\documentclass[11pt]{article}

\usepackage[T1]{fontenc}
\usepackage[utf8]{inputenc}
\usepackage[english]{babel}
\usepackage[margin=1in]{geometry}
\usepackage{lmodern}
\usepackage{microtype}
\usepackage{mathtools}
\usepackage{amsmath,amssymb,amsthm}
\usepackage{array}
\usepackage{booktabs}
\usepackage{tabularx}
\usepackage{longtable}
\usepackage{float}
\usepackage{graphicx}
\usepackage{tikz}
\usetikzlibrary{arrows.meta,positioning}
\usepackage{pdflscape}
\usepackage{caption}
\usepackage{enumitem}
\usepackage{fancyhdr}
\usepackage[numbers]{natbib}
\usepackage{doi}
\usepackage{xurl}
\usepackage{hyperref}
\usepackage{cleveref}
\usepackage{orcidlink}
\usepackage{titling}

\usepackage{ifluatex}
\ifluatex
  \usepackage{fontspec}
  \setmonofont{DejaVu Sans Mono}[Scale=MatchLowercase]
\fi



\providecommand{\BirdInt}{\mathbf{BirdInt}^{\mathrm{aud}}_{\mathrm{fin}}}

\providecommand{\factor}{\mathrel{\sqsubseteq}}

\providecommand{\Pdrive}{\mathrm{P6}_{\mathrm{drive}}}
\providecommand{\PdriveH}[1]{\mathrm{P6}_{\mathrm{drive}}^{#1}}

\providecommand{\tightlist}{}
\providecommand{\pandocbounded}[1]{#1}
\providecommand{\real}[1]{#1}
 
\newtheorem{theorem}{Theorem}[section]
\newtheorem{lemma}[theorem]{Lemma}
\newtheorem{proposition}[theorem]{Proposition}
\newtheorem{corollary}[theorem]{Corollary}
\theoremstyle{definition}
\newtheorem{definition}[theorem]{Definition}
\newtheorem{remark}[theorem]{Remark}
\newenvironment{claimtheorem}[1]{\par\medskip\noindent\textbf{#1.}\itshape}{\par\medskip}
\newenvironment{claimproposition}[1]{\par\medskip\noindent\textbf{#1.}\itshape}{\par\medskip}
\newenvironment{claimdefinition}[1]{\par\medskip\noindent\textbf{#1.}}{\par\medskip}
\floatstyle{ruled}
\newfloat{algorithm}{tbp}{loa}
\floatname{algorithm}{Algorithm}

\urlstyle{same}
\captionsetup{
  font=small,
  labelfont=bf,
  labelsep=period
}
\captionsetup[table]{skip=6pt,position=top}
\captionsetup[figure]{skip=6pt}
\captionsetup[algorithm]{skip=6pt,position=top}
\crefname{algorithm}{algorithm}{algorithms}
\Crefname{algorithm}{Algorithm}{Algorithms}
\crefname{theorem}{theorem}{theorems}
\Crefname{theorem}{Theorem}{Theorems}
\crefname{lemma}{lemma}{lemmas}
\Crefname{lemma}{Lemma}{Lemmas}
\crefname{proposition}{proposition}{propositions}
\Crefname{proposition}{Proposition}{Propositions}
\crefname{corollary}{corollary}{corollaries}
\Crefname{corollary}{Corollary}{Corollaries}
\crefname{definition}{definition}{definitions}
\Crefname{definition}{Definition}{Definitions}
\crefname{remark}{remark}{remarks}
\Crefname{remark}{Remark}{Remarks}
\setlist[enumerate]{leftmargin=*, itemsep=2pt, topsep=4pt}
\setlength{\droptitle}{-3.2em}
\pretitle{\begin{center}\LARGE\bfseries}
\posttitle{\par\end{center}\vspace{0.5em}}
\preauthor{\begin{center}\normalsize}
\postauthor{\par\end{center}\vspace{-0.4em}}
\predate{\begin{center}\small}
\postdate{\par\end{center}\vspace{0.6em}}
\setlength{\emergencystretch}{4em}
\tolerance=2000
\hbadness=2000
\clubpenalty=10000
\widowpenalty=10000
\displaywidowpenalty=10000
\interfootnotelinepenalty=10000
\allowdisplaybreaks[2]
\fancypagestyle{firstpage}{\fancyhf{}
  \fancyfoot[C]{\scriptsize Automorph Inc., Wilmington, DE, USA \quad \texttt{ioannis@automorph.io} \quad DOI: \href{https://doi.org/10.5281/zenodo.23096391}{10.5281/zenodo.23096391} \quad \textcopyright\ 2026 \quad CC-BY 4.0}
  \renewcommand{\headrulewidth}{0pt}
  \renewcommand{\footrulewidth}{0pt}
}

\title{Six Birds Foundations III:\\
A Finite Audited Interaction Calculus for SBT}
\author{Ioannis Tsiokos\,\orcidlink{0009-0009-7659-5964}}
\date{Version 2: October 2, 2026\\[2pt] \footnotesize (Version 1: May 11, 2026)}
\hypersetup{
  pdftitle={Six Birds Foundations III: A Finite Audited Interaction Calculus for SBT},
  pdfauthor={Ioannis Tsiokos},
  colorlinks=true,
  linkcolor=blue!45!black,
  citecolor=blue!45!black,
  urlcolor=blue!45!black
}

\begin{document}

\maketitle
\thispagestyle{firstpage}

\begin{abstract}
In Six Birds theory (SBT), Foundations I and II fixed six roles that a layer of description uses to close over what it tracks: descent, representability, route mismatch, refinement, packaging and audit ($P_1,\ldots,P_6$). This paper asks how the roles act on one another. It builds $\BirdInt$, a finite calculus in which an admissible interaction, one role acting on another, is a typed cell, classified under an explicit instrument from its records and measured defects. The calculus works with finite records; when numerical fields have decidable equality and order and every threshold-evaluated defect is computable in its declared type, its admissibility and classification judgments are finite checks.

The main results are as follows. No total algebra on the six roles recovers the status of a cell. An interaction may have distinct acting and informing roles. Five non-implications separate the verdicts of different judgment families: evidence that one role is active does not make an interaction cell active; an active cell does not make the corresponding unordered pair observable, which needs its own source record, count as real; and a strict promotion to a new layer implies neither an active cell, nor acceptance, one level higher, of the claim that a lower instrument's recorded verdicts are correct, nor acceptance of an instrument's soundness claim about itself. The last is never accepted, because the classifier's circularity rule applies to any soundness claim that cites a verdict at its own level (Theorem 24). Every well-formed cell receives exactly one status. Descent through a quotient, idempotent completion, strict extension of a theory and cycles on a finite graph have exact counterparts in the calculus, with their hypotheses made explicit. Further results say when promoting a construction to a new layer is sound, when a claim reads more than its instrument can see, and how an instrument at a higher level audits a stack of at most two lower instruments without certifying itself. Twelve countermodels mark what the calculus does not deliver. Four scoped model realizations exhibit fragments of the calculus on concrete hosts: abstract interpretation, graph cohomology, an audited class of states of a Cantor system, and record stores of the shape used by PICA (the Primitive Interaction Closure Algebra of a finite Markov chain). Each realization holds for records that meet a stated admissibility condition. For PICA we show that the condition can be met over every nonempty finite Markov chain, whatever raw interaction statuses are prescribed; the construction uses the chain only to check which finite trajectories have positive probability, and we do not check the condition on the records of the existing PICA simulator. Appendix F lists, result by result, what the accompanying Lean 4 artifact proves. The paper does not claim a universal exact-six theorem (that every description decomposes into exactly the six roles), a universal algebra of the six roles, a unique decomposition, an empirical realization, or complete self-certification at a single level.
\end{abstract}

\medskip
\noindent\textbf{Keywords.} finite typed calculus; audited interaction; emergence calculus; status semantics; defect calculus; theory packages; promotion bridge; model realization; rotating audit; nonfactorization.

\setcounter{tocdepth}{2}
\tableofcontents
\clearpage

\section{Introduction}\label{sn:1}

Six Birds theory (SBT) studies how a layer of description closes over what it tracks: which distinctions a coarse description can keep, which updates it can follow, and what it must record to justify its claims. Foundations I and II \citep{Tsiokos2026Foundations,Tsiokos2026FoundationsII} identified six roles that such closure uses. This paper asks how those roles act on one another, and builds a finite calculus, \(\mathbf{BirdInt}^{\mathrm{aud}}_{\mathrm{fin}}\), in which that question has exact answers. Subsection \hyperref[sn:1.1]{1.1} explains the question and the object, Subsection \hyperref[sn:1.2]{1.2} lists the results in words, Subsection \hyperref[sn:1.3]{1.3} states what the paper does not prove, Subsection \hyperref[sn:1.4]{1.4} gives a reading guide, and Subsection \hyperref[sn:1.5]{1.5} maps every numbered result to its place.

\subsection{The Question and the Object}\label{sn:1.1}

\subsubsection{Starting Point}\label{sn:1.1.1}

The six roles, written \(P_1,\ldots,P_6\), are:

\begin{itemize}
\tightlist
\item
  \(P_1\), \emph{descent}: whether an update respects a grouping of states;
\item
  \(P_2\), \emph{representability}: whether a specification can be realized, and which support it may use;
\item
  \(P_3\), \emph{route mismatch}: whether two ways of reaching the same result agree;
\item
  \(P_4\), \emph{refinement}: stages, refinements, and gluing local data into global data;
\item
  \(P_5\), \emph{packaging}: forming objects by quotients and completions;
\item
  \(P_6\), \emph{audit}: records, provenance and checks.
\end{itemize}

Foundations I derived these roles; Foundations II gave them exact meanings and showed that, on a domain \(\mathsf{FATCD}\) of finite audited typed closure descriptions, every description has a record with exactly six role channels \citep{Tsiokos2026FoundationsII}. Two earlier models show the roles at work: PICA (the Primitive Interaction Closure Algebra), built on a finite Markov chain whose dynamics come from six primitive mechanisms \citep{Tsiokos2026PICA}, and a Cantor system on which a packaged theory strictly extends a base theory \citep{Tsiokos2026Cantor}.

\subsubsection{The Question}\label{sn:1.1.2}

A role rarely acts alone. An audit (\(P_6\)) may record a route mismatch (\(P_3\)); a gate (\(P_2\)) may cut the cycles of a graph that carry a drive. The paper writes \(P_i\leftarrow P_j\) for such a \emph{directed cell}: the \emph{actor} \(P_i\) performs an update, using a witness supplied by the \emph{informant} \(P_j\). Whether such an interaction happened, failed, or was never attempted is a matter of record, and the paper asks what exact structure these records must have. One might hope for an algebra, a rule that combines two roles into a third or into a verdict. Two cells on the same pair of roles can have different statuses (countermodel \(\mathsf{CM}_7\)), and Theorem \hyperref[cl:thm:1]{1} shows that then no such rule recovers cell status. Interactions must therefore be judged cell by cell, from their records.

\subsubsection{The Object}\label{sn:1.1.3}

The calculus \(\mathbf{BirdInt}^{\mathrm{aud}}_{\mathrm{fin}}\), defined in Sections \hyperref[sn:3]{3} and \hyperref[sn:4]{4}, is a finite typed object with the following parts:

\begin{enumerate}
\def\labelenumi{\arabic{enumi}.}
\tightlist
\item
  the six role labels, with no operation on them;
\item
  directed cells, each a finite record of its actor and informant, witness, update, threshold, audit, visibility and defect data, together with a profile \(\lambda\) recording the levels, scale, regime and other conditions under which the cell is read;
\item
  seven finite families of statuses, for roles, cells, pairs of roles, promotions, claims, gates and squares, with priority orders for the role, cell, pair, promotion, claim and square classifiers; a few of the conditions that the classifiers test, including whether a role configuration collapses, whether a pair contract forbids a source class, whether partial evidence is present, and whether a pair's branches are compatible, are declared by the instance rather than computed from the record;
\item
  theory packages, instruments (the declared means by which content is seen or suppressed), source and audit records, and bridges between packages and between visibility regimes;
\item
  an admissibility check for each component;
\item
  defect records, which route a claim away from acceptance through the classifier of Section \hyperref[sn:5]{5};
\item
  a register of nonclaims, which records what each admissible claim does not assert.
\end{enumerate}

\emph{The directed cell in brief.} An instrument \(I\) splits the finite record set \(\mathsf{Cont}\) into visible, suppressed, out-of-scope and unknown records. A directed cell \(P_i\leftarrow P_j\) with profile \(\lambda\) is a record \((W_j,U_i,L,V,\Theta,A,\delta)\): a witness with signature in \(\mathsf{Wit}(P_j)\) and an update with signature in \(\mathsf{Upd}(P_i)\) (either may be an audited absence marker), a level and lens record \(L\), a visibility record \(V\) listing any bridges for suppressed fields, thresholds \(\Theta\), an audit record \(A\), and a defect record \(\delta\). It is well formed when the typing and recomputation conditions of Subsubsection \hyperref[sn:4.3.3]{4.3.3} hold, among them that \(\delta\) equals its value recomputed from the other fields; this check is decidable under the numerical decidability and computability hypotheses of Theorem \hyperref[cl:thm:3]{3}. Its status is that of the first row of the table of Subsubsection \hyperref[sn:5.6.6]{5.6.6} that holds, in the order \(\texttt{outside\_scope}\), \(\texttt{absent}\), \(\texttt{blocked}\), \(\texttt{undefined\_circular}\), \(\texttt{collapsed}\), \(\texttt{below\_threshold}\), \(\texttt{trivial}\), \(\texttt{implicit}\), \(\texttt{action}\), followed by a final default row \(\texttt{blocked}\) that makes the table total; Theorem \hyperref[cl:thm:4]{4} shows that no well-formed cell reaches it.

The calculus works with finite records. When numerical fields have decidable equality and order and every threshold-evaluated defect is computable in its declared type, its admissibility and classification judgments are finite checks: Theorem \hyperref[cl:thm:3]{3} decides cell well-formedness, Theorem \hyperref[cl:thm:4]{4} makes the cell classifier total, and the admissibility predicates of Subsubsections \hyperref[sn:3.4.1]{3.4.1}, \hyperref[sn:3.6.2]{3.6.2}, \hyperref[sn:10.2.1]{10.2.1} and \hyperref[sn:11.7.3]{11.7.3} are finite checks on declared fields. Some classifier inputs are declared Booleans rather than computed values, including role collapse, source prohibition, partial evidence and branch compatibility. The running example of Subsection \hyperref[sn:3.8]{3.8} shows the parts at work. On the subsets of \(\{a,b,c\}\), let \(E_1\) add \(b\) to every set containing \(a\), and \(E_2\) add \(c\) to every set containing \(b\). Writing \(E_1E_2\) for \(E_1\circ E_2\) (apply \(E_2\) first), \(E_1E_2(\{a\})=E_1(\{a\})=\{a,b\}\) while \(E_2E_1(\{a\})=E_2(\{a,b\})=\{a,b,c\}\): the two completions do not commute, which is a route mismatch (\(P_3\)). An audit that recomputes both sides and records the difference is the directed cell \(P_6\leftarrow P_3\), and the classifier gives it the status \(\texttt{action}\).

Figure \ref{fig:pipeline} shows how a cell is judged.

\begin{figure}[htbp]
\centering
\begin{tikzpicture}[>={Stealth[length=2mm]},node distance=6mm,
  box/.style={rectangle,rounded corners,draw,align=center,font=\footnotesize,minimum height=9mm,text width=2.3cm}]
\node[box] (rec) {cell record\\ $P_i\leftarrow P_j$};
\node[box,right=of rec] (wf) {well-formedness\\ (Theorem \hyperref[cl:thm:3]{3})};
\node[box,right=of wf] (def) {defect records\\ (Section \hyperref[sn:5]{5})};
\node[box,right=of def] (cls) {priority classifier\\ (Subsection \hyperref[sn:5.6]{5.6})};
\node[box,below=of cls,text width=3.4cm] (st) {one status for each\\ well-formed cell (Theorem \hyperref[cl:thm:4]{4})};
\node[box,above=of wf,text width=7.4cm] (ins) {instrument $I$: what is visible, which sources and bridges are admitted, which audits apply};
\draw[->] (rec) -- (wf); \draw[->] (wf) -- (def); \draw[->] (def) -- (cls); \draw[->] (cls) -- (st);
\draw[->,dashed] (ins.south) -- (wf.north);
\draw[->,dashed] (ins.east) -| (cls.north);
\end{tikzpicture}
\caption{How a directed cell is judged. Well-formedness (Theorem \hyperref[cl:thm:3]{3}) includes recomputing the cell's defect records (step 8 of Subsubsection \hyperref[sn:6.3.1]{6.3.1}); the classifier reads those defects and assigns each well-formed cell exactly one status; the instrument fixes visibility, admitted sources and bridges, and audit rules throughout.}\label{fig:pipeline}
\end{figure}

\subsection{What the Paper Proves}\label{sn:1.2}

The results fall into five groups. Table \ref{tab:map} gives the exact statement and location of each.

\subsubsection{The Core Calculus (Section 6 and Subsection 4.7)}\label{sn:1.2.1}

\begin{itemize}
\tightlist
\item
  \emph{No total algebra} (Theorem \hyperref[cl:thm:1]{1}): cell status does not factor through the pair of roles, so no binary operation on the six labels recovers it. This rules out reading a cell as a product of its two labels; the proof exhibits two records on the pair \(P_6\leftarrow P_3\), one classified \(\texttt{action}\) and one \(\texttt{blocked}\).
\item
  \emph{Typed non-collapse} (Theorem \hyperref[cl:thm:2]{2}): an active cell can have distinct actor and informant, and a cell judgment never forces them to be equal. This rules out identifying the acting role with the informing role; the proof uses the typing rule that takes the informant's data from \(\mathsf{Wit}(P_j)\) and the actor's data from \(\mathsf{Upd}(P_i)\).
\item
  \emph{Decidability} (Theorem \hyperref[cl:thm:3]{3}): a finite procedure decides whether a candidate cell is well formed, for numerical fields with decidable equality and order and computable threshold data.
\item
  \emph{Unique status} (Theorem \hyperref[cl:thm:4]{4}): the nine substantive rows of the classifier cover every well-formed cell, so the first matching row gives each well-formed cell exactly one status and the default \(\texttt{blocked}\) row is never used. The proof shows that a well-formed cell meeting none of the rows \(\texttt{outside\_scope}\) through \(\texttt{implicit}\) meets the \(\texttt{action}\) row; uniqueness is then the first-match rule.
\item
  \emph{Top-down channels} (Subsection \hyperref[sn:4.7]{4.7}): structural downward influence, such as a completion or a context index, is distinguished from a top-down causal channel, which must pass intervention gates; three propositions make the distinction precise.
\item
  \emph{Separation of status families} (Theorem \hyperref[cl:thm:5]{5}): an active role does not imply an active cell, an active cell does not imply a real pair observable, and a strict promotion implies neither an active cell, nor acceptance of a same-level soundness claim, nor acceptance of a lifted soundness claim under a higher instrument (the same-level case holds because, by Theorem \hyperref[cl:thm:24]{24}, a same-level soundness claim is never accepted).
\end{itemize}

\subsubsection{Descent, Packaging, Completion and Strictness (Section 7)}\label{sn:1.2.2}

\begin{itemize}
\tightlist
\item
  \emph{Descent} (Theorems \hyperref[cl:thm:6]{6} and \hyperref[cl:thm:7]{7}): an update descends through a finite map exactly when it respects the map's fibers, and the split pairs record every failure.
\item
  \emph{Closure deficit} (Theorem \hyperref[cl:thm:8]{8}): on a finite Markov host, when the admissible class of coarse-grained kernels contains a kernel that agrees with the conditional macro kernel on the support of the coarse state, its least Kullback--Leibler predictive loss equals a conditional mutual information.
\item
  \emph{Noncommuting completions} (Theorems \hyperref[cl:thm:9]{9} and \hyperref[cl:thm:10]{10}): two idempotent completions of a finite set need not commute, and their mismatch is a finite defect.
\item
  \emph{Saturation} (Theorem \hyperref[cl:thm:11]{11}): an idempotent map equals each of its positive powers.
\item
  \emph{Strict extension} (Theorems \hyperref[cl:thm:12]{12} and \hyperref[cl:thm:13]{13}): a new object map fails to factor through an old one exactly when some pair of points is identified by the old map and separated by the new one.
\item
  \emph{Definability} (Theorem \hyperref[cl:thm:14]{14}): a finite map whose image has \(k\) elements defines exactly \(2^k\) preimage sets.
\end{itemize}

\subsubsection{No-Go Results and Countermodels (Sections 8 and 9)}\label{sn:1.2.3}

\begin{itemize}
\tightlist
\item
  \emph{No drive on graphs without cycles} (Theorems \hyperref[cl:thm:15]{15}--\hyperref[cl:thm:18]{18}): a drive here is a finite-graph \(1\)-cochain with a nonzero integral around some cycle; for a Markov chain with positive transitions in both directions on each edge, the forward/backward log-ratio cochain is a motivating example when a cycle affinity is nonzero. Such a drive needs a nonzero cohomology class; a forest has none, an exact cochain integrates to zero on every cycle, and gating a graph down to a forest removes the drive.
\item
  \emph{Twelve countermodels} \(\mathsf{CM}_1\)--\(\mathsf{CM}_{12}\) (Section \hyperref[sn:9]{9} and Appendix \hyperref[sn:C]{C}): finite examples, each blocking one tempting overclaim, for instance that packaging implies closure or that holonomy implies drive.
\end{itemize}

\subsubsection{Promotion and Instruments (Sections 10 and 11)}\label{sn:1.2.4}

\begin{itemize}
\tightlist
\item
  \emph{Promotion} (Theorems \hyperref[cl:thm:19]{19}--\hyperref[cl:thm:22]{22}): an accepted, strict or non-strict promotion verdict requires every required core gate to pass; a strict promotion verdict implies that the new object map does not factor through the old one, and a non-strict verdict implies that it does, and a two-point example realizes the strict case.
\item
  \emph{No overreading} (Theorem \hyperref[cl:thm:23]{23}): an accepted claim uses suppressed content only through an admissible visibility bridge. The proof uses the claim classifier's priority order: a nonempty overreading defect prevents \(\texttt{accepted}\), because \(\texttt{blocked}\) or a higher-priority status is returned first.
\item
  \emph{No same-level self-certification} (Theorem \hyperref[cl:thm:24]{24}): the classifier does not accept an instrument's claim of its own complete soundness without a change of level. The proof uses the same priority order: by the definition of its use-set, the soundness claim refers to its own verdict, so in scope the classifier's rule for such claims gives \(\texttt{undefined\_circular}\) or a status ranked higher.
\item
  \emph{Finite rotating audit} (Theorem \hyperref[cl:thm:25]{25}): given a level above a finite stack of instruments, an instrument at that level can audit the stack without certifying its own soundness. With the three levels of the calculus, a stack has at most two instruments below the auditing one.
\end{itemize}

\subsubsection{Models (Sections 12--14)}\label{sn:1.2.5}

\begin{itemize}
\tightlist
\item
  \emph{Abstract interpretation} (Theorems \hyperref[cl:thm:26]{26}--\hyperref[cl:thm:28]{28} and Model Theorem \hyperref[cl:thm:29]{29}): the closure of an abstract-interpretation host is a packaging, sound transformers are lax descent squares, and exact ones are descents, and representable elements are fixed points; every nonempty finite family of sound hosts over a common concrete domain and transformer realizes \(P_1/P_2/P_5/P_6\), and also \(P_4\) when its refinement records contain a strict refinement.
\item
  \emph{PICA and Cantor} (Model Theorems \hyperref[cl:thm:33]{33} and \hyperref[cl:thm:34]{34}): an admissible PICA realization object realizes the cell, pair and provenance fragment of the calculus, and the Cantor shell realizes a strict promotion. The PICA admissibility predicate is shown satisfiable over every nonempty finite Markov chain by stores that reproduce the raw cell statuses; these stores read the chain only through traces of positive transition probability, and the predicate is not checked on the stores of the PICA simulator. The theorem therefore shows that PICA's record format is consistent with the calculus; it does not compute any cell status from the chain.
\item
  \emph{Squares} (Theorems \hyperref[cl:thm:30]{30}--\hyperref[cl:thm:32]{32}): in a finite fragment of decorated squares (double-category shape; the axioms are not assumed, Subsubsection \hyperref[sn:14.1.1]{14.1.1}), descent and completion appear as decorated squares, and strict extension appears as the absence of a filler.
\item
  \emph{Lean} (Proposition \hyperref[cl:prop:35]{35}): a dated validation run of the Lean artifact passed for every entry of its manifest.
\end{itemize}

\subsection{What the Paper Does Not Prove}\label{sn:1.3}

The negative scope of the paper is stated in Subsubsection \hyperref[sn:2.4.3]{2.4.3}, and its numbered nonclaims are recorded in Subsection \hyperref[sn:15.2]{15.2}. In brief, the paper does not construct a total algebra on the six primitive labels (Theorem \hyperref[cl:thm:1]{1} shows that none recovers cell status), does not assert an unrestricted exact-six decomposition, does not equate a role channel with an active projection or an active directed cell, does not claim that its four model families realize the whole calculus, and does not claim an instrument that certifies its own complete soundness (Theorems \hyperref[cl:thm:24]{24} and \hyperref[cl:thm:25]{25}).

\subsection{Reader's Map}\label{sn:1.4}

The paper can be entered by two independent routes. The finite mathematics of Theorems \hyperref[cl:thm:6]{6}--\hyperref[cl:thm:18]{18} (Sections \hyperref[sn:7]{7}--\hyperref[sn:8]{8}), 26--28 (Section \hyperref[sn:12]{12}) and 30--32 (Section \hyperref[sn:14]{14}) uses Subsections \hyperref[sn:2.1]{2.1}--\hyperref[sn:2.2]{2.2}, the local definitions of those sections and the definitions they cite, and can be read first. The calculus proper (Theorems \hyperref[cl:thm:1]{1}--\hyperref[cl:thm:5]{5} and \hyperref[cl:thm:19]{19}--\hyperref[cl:thm:25]{25}) uses the definitions of Sections \hyperref[sn:3]{3}--\hyperref[sn:5]{5} listed below. For a first reading of Theorems \hyperref[cl:thm:1]{1}--\hyperref[cl:thm:4]{4}, start with the running example of Subsection \hyperref[sn:3.8]{3.8}, the directed-cell table of Subsubsection \hyperref[sn:5.6.6]{5.6.6}, whose last column evaluates that example, and the first-match rule at the head of Section \hyperref[sn:6]{6}, then read Section \hyperref[sn:6]{6} directly. The list below supplies the definitions used to verify well-formedness and the classifier computations in the proofs. The proofs use four objects only: the running-example records (Subsection \hyperref[sn:3.8]{3.8}), the eight well-formedness conditions (Subsubsection \hyperref[sn:4.3.3]{4.3.3}), the defect sets and audit codes glossed in the key of Subsubsection \hyperref[sn:5.6.6]{5.6.6}, and the directed-cell table that follows that key; the list says where each ingredient of these four objects is defined. These four passages suffice to follow the directed-cell row computations of Theorems \hyperref[cl:thm:1]{1}--\hyperref[cl:thm:4]{4}; the list below says where each field they name is defined and is meant for lookup, not for reading in advance:

\begin{itemize}
\tightlist
\item
  labels and levels: Subsection \hyperref[sn:2.2]{2.2};
\item
  records and the running example: Subsubsections \hyperref[sn:3.1.1]{3.1.1}, \hyperref[sn:3.3.1]{3.3.1}, \hyperref[sn:3.3.3]{3.3.3}, \hyperref[sn:3.4.1]{3.4.1}, \hyperref[sn:3.5.2]{3.5.2}, \hyperref[sn:3.5.3]{3.5.3} and \hyperref[sn:3.6.2]{3.6.2}, and Subsections \hyperref[sn:3.7]{3.7} and \hyperref[sn:3.8]{3.8};
\item
  the cell record and its well-formedness: Subsubsections \hyperref[sn:4.3.1]{4.3.1}--\hyperref[sn:4.3.4]{4.3.4} and \hyperref[sn:4.9.2]{4.9.2}, and Subsection \hyperref[sn:4.10]{4.10};
\item
  the cell's defects: Subsubsections \hyperref[sn:5.1.2]{5.1.2}, \hyperref[sn:5.2.1]{5.2.1}, \hyperref[sn:5.2.2]{5.2.2}, \hyperref[sn:5.3.5]{5.3.5}, \hyperref[sn:5.3.6]{5.3.6} and \hyperref[sn:5.5.1]{5.5.1}--\hyperref[sn:5.5.3]{5.5.3};
\item
  the classifier: the reference-graph paragraph of Subsubsection \hyperref[sn:5.6.5]{5.6.5} and the first table of Subsubsection \hyperref[sn:5.6.6]{5.6.6}, whose last column evaluates the running example.
\end{itemize}

Later results cite the judgment families and host definitions they need. Subsections \hyperref[sn:3.3]{3.3}--\hyperref[sn:3.6]{3.6} and \hyperref[sn:4.9]{4.9}--\hyperref[sn:4.10]{4.10} are reference material; the parts that Theorems \hyperref[cl:thm:1]{1}--\hyperref[cl:thm:4]{4} need are listed above.

The paper is organized into sixteen sections, the last of which is the conclusion (Section \hyperref[sn:16]{16}), and ten appendices (A--J). The reader's map is:

\begin{itemize}
\tightlist
\item
  \emph{Section \hyperref[sn:2]{2}} records preliminaries and scope: finite host mathematics, primitive labels and levels, claim strengths, status discipline, and the global nonclaim conventions.
\item
  \emph{Sections \hyperref[sn:3]{3} and \hyperref[sn:4]{4}} define the finite audited calculus: the finite typed domain, the central object \(\mathbf{BirdInt}^{\mathrm{aud}}_{\mathrm{fin}}\) as a finite typed tuple, theory packages, instruments, source-of-truth records, audit records, visibility bridges, directed-cell records, pair-observable records, promotion bridges, and the per-cell profile schema, and the distinction between structural downward influence and top-down causal channels (Subsection \hyperref[sn:4.7]{4.7}).
\item
  \emph{Section \hyperref[sn:5]{5}} records the status and defect semantics: the seven status families, the per-component defect schemas, the visibility-defect family, and the classifier priority orders.
\item
  \emph{Sections \hyperref[sn:6]{6}--\hyperref[sn:8]{8}} prove the finite theorem spine: the no-total-algebra theorem and the four core finite-interaction theorems (Section \hyperref[sn:6]{6}); the nine descent/packaging/completion/strictness theorems (Section \hyperref[sn:7]{7}); and the four cohomological no-drive theorems on graph-supported drive (Section \hyperref[sn:8]{8}).
\item
  \emph{Section \hyperref[sn:9]{9}} records the countermodel atlas in main-text form, with full sheets for \(\mathsf{CM}_{1}\)--\(\mathsf{CM}_{6}\) and named pointers to \(\mathsf{CM}_{7}\)--\(\mathsf{CM}_{12}\). The full twelve-countermodel atlas is recorded in Appendix \hyperref[sn:C]{C}.
\item
  \emph{Section \hyperref[sn:10]{10}} records promotion and stacking: the promotion-bridge schema, the eight-gate machinery, the promotion-status consequences of the classifier, Theorems \hyperref[cl:thm:19]{19}--\hyperref[cl:thm:21]{21}, a two-point strict-bridge example (Theorem \hyperref[cl:thm:22]{22}), and the stacking claims not proved.
\item
  \emph{Section \hyperref[sn:11]{11}} records instruments, visibility, and rotating audit: the instrument-indexed claim semantics, the classifier priority, the visibility/suppression discipline, the no-overreading theorem (Theorem \hyperref[cl:thm:23]{23}), the same-level self-audit failure theorem (Theorem \hyperref[cl:thm:24]{24}), the \(P_{6}\leftarrow P_{6}\) status separation, the rotating-audit chain and bridge, the finite rotating-audit theorem (Theorem \hyperref[cl:thm:25]{25}), and the instrument-transfer rule.
\item
  \emph{Section \hyperref[sn:12]{12}} records the abstract-interpretation model family: the host data, the realized fragment \(P_{1}/P_{2}/P_{5}/P_{6}\), with \(P_{4}\) added when the refinement records contain a strict refinement, Theorems \hyperref[cl:thm:26]{26}--\hyperref[cl:thm:28]{28} (closure as packaging, descent as sound transformer, representability as fixedness), and Model Theorem \hyperref[cl:thm:29]{29}.
\item
  \emph{Section \hyperref[sn:13]{13}} records the four scoped model realizations (PICA, Cantor strict-bridge, abstract interpretation, graph/cohomology), Model Theorems \hyperref[cl:thm:33]{33} and \hyperref[cl:thm:34]{34}, the combined model-realization matrix, and the cross-model role-coverage observation.
\item
  \emph{Section \hyperref[sn:14]{14}} records the secondary high-structure semantics: the finite decorated double-category fragment and the recovery theorems 30--32 (descent-square recovery, completion-pasting recovery, strict extension as no-filler).
\item
  \emph{Section \hyperref[sn:15]{15}} records limitations, nonclaims, and future work, with the register of numbered nonclaims (Subsection \hyperref[sn:15.2]{15.2}).
\item
  \emph{Section \hyperref[sn:16]{16}} records the conclusion.
\item
  \emph{Appendices \hyperref[sn:A]{A}--\hyperref[sn:J]{J}} record the long-form supplementary material: full definitions and tables (A), long proofs (B), the full countermodel atlas (C), the four model-realization sheets (D), the high-structure semantics extension (E), the Lean mechanization appendix (F), the dependency graph (G), the deferred and excluded claims register (H), the claim strengths and instrument-indexed forms of the numbered results (I), and the version history (J).
\end{itemize}

The result numbers do not follow the section order at one point: Model Theorem \hyperref[cl:thm:29]{29} is in Section \hyperref[sn:12]{12}, Model Theorems \hyperref[cl:thm:33]{33} and \hyperref[cl:thm:34]{34} are in Section \hyperref[sn:13]{13}, and Theorems \hyperref[cl:thm:30]{30}--\hyperref[cl:thm:32]{32} are in Section \hyperref[sn:14]{14}. Subsection \hyperref[sn:1.5]{1.5} lists every result with its location.

Sections \hyperref[sn:2]{2}--\hyperref[sn:5]{5} supply the formal-development prerequisites for the later calculus results. For a first reading of the mathematical statements of Theorems \hyperref[cl:thm:6]{6}--\hyperref[cl:thm:18]{18} (Sections \hyperref[sn:7]{7} and \hyperref[sn:8]{8}), Subsections \hyperref[sn:2.1]{2.1}--\hyperref[sn:2.2]{2.2} and the local definitions in Sections \hyperref[sn:7]{7} and \hyperref[sn:8]{8} suffice: the statements use finite sets and maps, finite probability and information quantities, and finite graph/cohomology data. Likewise, the statements of Theorems \hyperref[cl:thm:26]{26}--\hyperref[cl:thm:28]{28} use Subsubsection \hyperref[sn:12.1.1]{12.1.1} together with their local definitions and the definitions they cite, and those of Theorems \hyperref[cl:thm:30]{30}--\hyperref[cl:thm:32]{32} use Subsection \hyperref[sn:14.1]{14.1} and the definitions they cite. Sections \hyperref[sn:3]{3}--\hyperref[sn:5]{5} supply the role, instrument, and status readings around them. Sections \hyperref[sn:6]{6}--\hyperref[sn:11]{11} are the theorem spine: they build on Sections \hyperref[sn:2]{2}--\hyperref[sn:5]{5} and, where stated, on earlier theorem sections. Section \hyperref[sn:12]{12} depends on Sections \hyperref[sn:2]{2}--\hyperref[sn:5]{5} and \hyperref[sn:7]{7} and on Subsection \hyperref[sn:13.1]{13.1}; Section \hyperref[sn:13]{13} depends on Sections \hyperref[sn:2]{2}--\hyperref[sn:8]{8} and \hyperref[sn:10]{10}--\hyperref[sn:12]{12}, and Subsections \hyperref[sn:13.5]{13.5}--\hyperref[sn:13.6]{13.6} use the family of Section \hyperref[sn:12]{12}; its model-realization claims are recorded under instruments. Section \hyperref[sn:14]{14} depends on Sections \hyperref[sn:2]{2}--\hyperref[sn:5]{5} and \hyperref[sn:7]{7}; its visibility and audit square types use Subsubsections \hyperref[sn:5.5.2]{5.5.2} and \hyperref[sn:11.7.3]{11.7.3}, and Theorem \hyperref[cl:thm:32]{32} cites Subsection \hyperref[sn:13.3]{13.3} only as an example. Section \hyperref[sn:15]{15} collects scope discipline from every preceding section. Appendix \hyperref[sn:F]{F} records the proof-assistant disclosure table, validation boundary, and audit-grade Proposition \hyperref[cl:prop:35]{35}. The reader interested in model realizations can add Sections \hyperref[sn:12]{12} and \hyperref[sn:13]{13} and Appendix \hyperref[sn:D]{D}; the reader interested in mechanization can read Appendix \hyperref[sn:F]{F}.

\subsection{Map of Results}\label{sn:1.5}

Table \ref{tab:map} lists the numbered results, together with the three unnumbered propositions of Subsection \hyperref[sn:4.7]{4.7}, and where each is stated and proved. Appendix \hyperref[sn:I]{I} records the claim strength and verdict of each numbered result, which is its logical status (theorem, model realization, or audit report), not a measure of importance; the three propositions of Subsection \hyperref[sn:4.7]{4.7} are proposition-grade with verdict accepted (Subsubsection \hyperref[sn:2.3.1]{2.3.1}). Appendix \hyperref[sn:F]{F} records the Lean coverage of every row.

\begingroup\small
\begin{longtable}{@{}>{\raggedright\arraybackslash}p{0.2\linewidth}>{\raggedright\arraybackslash}p{0.58\linewidth}>{\raggedright\arraybackslash}p{0.16\linewidth}@{}}
\caption{Map of results.}\label{tab:map}\\
\toprule
Result & What it establishes & Where \\
\midrule
\endfirsthead
\toprule
Result & What it establishes & Where \\
\midrule
\endhead
\bottomrule
\endfoot
Proposition (Subsection \hyperref[sn:4.7]{4.7}), top-down gate soundness & An accepted top-down channel claim requires every listed gate to pass, including at least two admissible macro interventions and a passing effect comparison. & §4.7 \\
Proposition (Subsection \hyperref[sn:4.7]{4.7}), top-down implies structural & An accepted top-down channel claim carries a structural downward dependency. & §4.7 \\
Proposition (Subsection \hyperref[sn:4.7]{4.7}), structural is insufficient & Structural downward influence alone does not yield an accepted top-down channel claim. & §4.7 \\
Theorem \hyperref[cl:thm:1]{1} & Two covered directed cells on the pair $(P_6,P_3)$ receive the statuses $\texttt{action}$ and $\texttt{blocked}$, so cell status does not factor through the pair; no binary operation on the six labels, followed by a decoder, recovers cell status. & §6.1 \\
Theorem \hyperref[cl:thm:2]{2} & Some well-formed active directed cell has distinct actor and informant primitives; a directed-cell judgment never forces the two to be equal. & §6.2 \\
Theorem \hyperref[cl:thm:3]{3} & Well-formedness of a candidate directed cell is decided by a finite eight-step procedure (for numerical fields with decidable equality and order and computable threshold data). & §6.3 \\
Theorem \hyperref[cl:thm:4]{4} & The priority classifier assigns exactly one status to each well-formed directed cell. & §6.4 \\
Theorem \hyperref[cl:thm:5]{5} & Five non-implications between the judgment-attached status families: an active role channel does not imply an active cell, an active cell does not imply a real pair observable, and a strict promotion implies neither an active cell, nor acceptance of a same-level soundness claim (never accepted, by Theorem \hyperref[cl:thm:24]{24}), nor accepted lifted soundness under a higher instrument. & §6.5 \\
Theorem \hyperref[cl:thm:6]{6} & A finite update descends through a finite map exactly when it respects the fibers of the map. & §7.1 \\
Theorem \hyperref[cl:thm:7]{7} & The descent defect, the set of split pairs, is empty exactly when the update descends. & §7.2 \\
Theorem \hyperref[cl:thm:8]{8} & On a finite Markov host, for any class $\mathcal K$ of macro kernels, $\inf_{\mathcal K}\mathcal L_\tau\ge I(X_t;\Pi(X_{t+\tau})\mid\Pi(X_t))$; the minimum over $\mathcal K$ exists and equals $I(X_t;\Pi(X_{t+\tau})\mid\Pi(X_t))$ exactly when $\mathcal K$ contains a kernel that agrees with the conditional macro kernel $K^\ast$ on the support of $\Pi(X_t)$. & §7.3; App.~B.1 \\
Theorem \hyperref[cl:thm:9]{9} & Two idempotent completion operators on a finite carrier need not commute; an explicit pair is given. & §7.4 \\
Theorem \hyperref[cl:thm:10]{10} & The route-mismatch defect of two completions is nonempty exactly when they do not commute. & §7.5 \\
Theorem \hyperref[cl:thm:11]{11} & An idempotent self-map of a finite set equals each of its positive powers. & §7.6 \\
Theorem \hyperref[cl:thm:12]{12} & A finite object map $\pi_1$ fails to factor through $\pi_0$ exactly when some pair of points is identified by $\pi_0$ and separated by $\pi_1$. & §7.7 \\
Theorem \hyperref[cl:thm:13]{13} & The factorization defect is nonempty exactly when the factorization fails. & §7.8 \\
Theorem \hyperref[cl:thm:14]{14} & A finite map whose image has $k$ elements defines exactly $2^k$ preimage sets. & §7.9; App.~B.4 \\
Theorem \hyperref[cl:thm:15]{15} & A finite graph with first Betti number zero has trivial first cohomology and admits no cycle-supported drive certificate. & §8.2; App.~B.2.1 \\
Theorem \hyperref[cl:thm:16]{16} & An exact $1$-cochain integrates to zero on every cycle. & §8.3; App.~B.2.2 \\
Theorem \hyperref[cl:thm:17]{17} & A $1$-cochain has a nonzero cohomology class exactly when some cycle integral is nonzero. & §8.4; App.~B.2.3 \\
Theorem \hyperref[cl:thm:18]{18} & A gate that turns a graph into a forest leaves trivial first cohomology, and no cycle-supported drive certificate exists on the forest. & §8.5; App.~B.2.4 \\
Theorem \hyperref[cl:thm:19]{19} & A promotion verdict of accepted, strict or non-strict requires every required core gate to exist and pass. & §10.3 \\
Theorem \hyperref[cl:thm:20]{20} & A strict promotion verdict implies that the new object map does not factor through the old one. & §10.4 \\
Theorem \hyperref[cl:thm:21]{21} & A non-strict promotion verdict implies that the new object map factors through the old one. & §10.5 \\
Theorem \hyperref[cl:thm:22]{22} & A promotion bridge from a one-point layer to a two-point layer over a two-element carrier is admissible and is classified strict. & §10.6 \\
Theorem \hyperref[cl:thm:23]{23} & Every suppressed record in the use-set of an accepted claim is bridged by an admissible visibility bridge. & §11.4 \\
Theorem \hyperref[cl:thm:24]{24} & Without a level shift, the claim classifier does not accept an instrument's claim of its own complete soundness. & §11.5 \\
Theorem \hyperref[cl:thm:25]{25} & Given a level above a finite instrument stack, an extension instrument at that level audits the stack, accepts its compliance exactly when the stack-audit defect is empty, and does not thereby certify its own soundness. & §11.8 \\
Theorem \hyperref[cl:thm:26]{26} & The closure $\rho=\gamma\alpha$ of a finite abstract-interpretation host is extensive, monotone and idempotent. & §12.2 \\
Theorem \hyperref[cl:thm:27]{27} & The two soundness inequalities for an abstract transformer, $\alpha F\le F^{\sharp}\alpha$ and $F\gamma\le\gamma F^{\sharp}$, are equivalent. & §12.3 \\
Theorem \hyperref[cl:thm:28]{28} & The representable concrete elements, the fixed points of $\rho$ and the image of $\gamma$ coincide. & §12.4 \\
Model Theorem \hyperref[cl:thm:29]{29} & A nonempty finite family of admissible (sound) abstract-interpretation hosts sharing one concrete domain and one concrete transformer realizes the role fragment $P_1/P_2/P_5/P_6$, and $P_1/P_2/P_4/P_5/P_6$ when its refinement records contain a strict refinement; it also realizes $P_3$ when two host closures fail to commute (for instance, the hosts built from the completions of Theorem \hyperref[cl:thm:9]{9}). & §12.5 \\
Theorem \hyperref[cl:thm:30]{30} & The descent square is exact exactly when its given descended map commutes with the quotient, and a commuting (set-map) filler exists exactly when the update descends through the quotient. & §14.2 \\
Theorem \hyperref[cl:thm:31]{31} & The completion square is exact exactly when the two completions commute; otherwise their noncommutation is recorded as the square defect $\Delta_3^{\mathrm{comp}}(E_1,E_2)$. & §14.3 \\
Theorem \hyperref[cl:thm:32]{32} & A factorization filler exists exactly when the new object map factors through the old one. & §14.4 \\
Model Theorem \hyperref[cl:thm:33]{33} & Stores meeting $\mathsf{AdmPICAReal}$ realize the finite stochastic cell, pair and provenance fragment; the predicate is met over every nonempty finite Markov chain for every raw status table; the satisfying stores read the chain only through positive-probability traces, and the status table is an input, not a consequence. & §13.2 \\
Model Theorem \hyperref[cl:thm:34]{34} & An admissible audited Cantor shell with a nonempty factorization defect and passing gates realizes a strict bridge; the two-point construction in its Moreover clause has the object-map pattern of Theorem \hyperref[cl:thm:22]{22} on a cited witness pair. & §13.3 \\
Proposition \hyperref[cl:prop:35]{35} & A dated strict validation run of the Lean artifact passed: each manifest entry has a checked Lean declaration. & App.~F.5 \\
\end{longtable}
\endgroup

\subsection{Relation to Existing Work}\label{sn:1.6}

Theorems \hyperref[cl:thm:6]{6} and \hyperref[cl:thm:7]{7} give a finite fiberwise descent criterion, related to descent along a quotient \citep{Grothendieck1959Descent}. Theorem \hyperref[cl:thm:8]{8} relates predictive closure loss to the question of Markov lumpability \citep{KemenySnell1976,Buchholz1994Lumpability}. Theorems \hyperref[cl:thm:9]{9}--\hyperref[cl:thm:11]{11} use finite completion operators; Theorems \hyperref[cl:thm:26]{26}--\hyperref[cl:thm:28]{28} use closure operators arising from Galois connections \citep{DaveyPriestley2002,Ore1944Galois,CousotCousot1977}. The \(P_3\) witness family includes critical pairs familiar from term rewriting \citep{BaaderNipkow1998}. Theorems \hyperref[cl:thm:24]{24} and \hyperref[cl:thm:25]{25} concern this paper's claim classifier: \(\mathrm{Sound}(I_0)\) is defined to reference its own verdict, and the classifier's circularity rule acts on that reference; no diagonal or arithmetization argument is involved, and the comparison with work on incompleteness and reflection \citep{Goedel1931Incompleteness,Feferman1962Progressions} is one of theme only. The stratification rule behind Theorems \hyperref[cl:thm:24]{24} and \hyperref[cl:thm:25]{25}, under which a soundness claim judged at level \(\ell\) may cite only verdicts of strictly lower levels, follows Tarski's hierarchy of object language and metalanguage \citep{Tarski1936Truth}; Theorem \hyperref[cl:thm:25]{25} is a finite three-level instance of that ascent. Theorems \hyperref[cl:thm:15]{15}--\hyperref[cl:thm:18]{18} are the finite-graph form of the cycle decomposition of Markov affinities in the network theory of \citet{Schnakenberg1976}; for an irreducible chain with \(P_{xy}>0\iff P_{yx}>0\), vanishing of every cycle integral of the log-ratio cochain is Kolmogorov's criterion for detailed balance \citep{Kelly1979}. Section \hyperref[sn:14]{14} organizes selected constructions as decorated squares in a finite double-category fragment \citep{Ehresmann1963Doubles,GrandisPare1999}.
 \section{Preliminaries and Scope}\label{sn:2}

This section fixes the host mathematics, the primitive labels, the level set, the claim-strength discipline, and the scope statement that the rest of the paper will assume. Every later definition, theorem, and model realization is interpreted inside the domain declared here, or inside an explicitly named finite extension host. No object, judgment, or proof in this paper depends on hosts that fall outside the inventory below. On a first reading, Subsection \hyperref[sn:2.2]{2.2} (the six roles and their levels) is enough; Subsections \hyperref[sn:2.1]{2.1}, \hyperref[sn:2.3]{2.3} and \hyperref[sn:2.4]{2.4} can be consulted when later sections cite them.

\subsection{Finite Host Mathematics}\label{sn:2.1}

The calculus represents its judgments and the data they use by finite typed records on declared hosts; each judgment is relative to a declared instrument. Later sections take these conventions for granted and do not repeat them.

\subsubsection{Base Host}\label{sn:2.1.1}

The base host of the calculus is

\[
\begin{aligned}
\mathsf H_{\mathrm{fin}} \;=\; \text{finite sets, finite maps, finite relations, finite records,}\\
\text{finite diagrams, finite graphs, finite audit records}.
\end{aligned}
\]

All carriers, lenses, packages, instruments, profiles, witnesses, updates, defects, gates, audits, statuses, and judgment instances introduced in Sections \hyperref[sn:3]{3} and \hyperref[sn:4]{4} are objects of \(\mathsf H_{\mathrm{fin}}\) unless an extension host is declared. In particular, every carrier, index set and record is finite, every diagram has finitely many vertices and edges, and every audit record is a finite tuple of finite fields. Fields used by a finite decision procedure have declared types with effective equality and the operations and order comparisons that procedure uses. The macro-kernel class used in Theorem \hyperref[cl:thm:8]{8} and the cochain and cycle spaces used in Theorems \hyperref[cl:thm:15]{15}--\hyperref[cl:thm:18]{18} may be infinite. The mathematical identities and equivalences of Theorems \hyperref[cl:thm:6]{6}, \hyperref[cl:thm:7]{7} and \hyperref[cl:thm:10]{10}--\hyperref[cl:thm:13]{13} hold for arbitrary sets and maps; the cardinality identity of Theorem \hyperref[cl:thm:14]{14} holds for an arbitrary map \(f\), with \(2^{|\mathrm{im}(f)|}\) read as a cardinal; and the closure, soundness, descent and representability conclusions of Theorems \hyperref[cl:thm:26]{26}--\hyperref[cl:thm:28]{28} hold between arbitrary posets under their stated Galois and monotonicity hypotheses. Finite-test, finite-record and audited-realization conclusions retain their finite-host hypotheses.

The base host carries the following structures, each used by a specific judgment family later in the paper:

\begin{longtable}[]{@{}ll@{}}
\toprule
\begin{minipage}[b]{0.47\columnwidth}\raggedright
Structure\strut
\end{minipage} & \begin{minipage}[b]{0.47\columnwidth}\raggedright
Use in the calculus\strut
\end{minipage}\tabularnewline
\midrule
\endhead
\begin{minipage}[t]{0.47\columnwidth}\raggedright
finite sets\strut
\end{minipage} & \begin{minipage}[t]{0.47\columnwidth}\raggedright
carriers \(Z\), primitive labels \(\mathbb P\), witness/update domains, status sets\strut
\end{minipage}\tabularnewline
\begin{minipage}[t]{0.47\columnwidth}\raggedright
finite maps\strut
\end{minipage} & \begin{minipage}[t]{0.47\columnwidth}\raggedright
lenses \(f:Z\to X\), quotients \(q\), completion endomaps \(E\), object maps \(\pi_j\), update maps \(U_i\)\strut
\end{minipage}\tabularnewline
\begin{minipage}[t]{0.47\columnwidth}\raggedright
finite relations\strut
\end{minipage} & \begin{minipage}[t]{0.47\columnwidth}\raggedright
equivalence on a carrier, dependency between primitives, gluing compatibility on a cover\strut
\end{minipage}\tabularnewline
\begin{minipage}[t]{0.47\columnwidth}\raggedright
finite records\strut
\end{minipage} & \begin{minipage}[t]{0.47\columnwidth}\raggedright
theory packages \(\mathcal T\), instruments \(I\), directed cells, promotion bridges, claim records, defect records, gate records\strut
\end{minipage}\tabularnewline
\begin{minipage}[t]{0.47\columnwidth}\raggedright
finite diagrams\strut
\end{minipage} & \begin{minipage}[t]{0.47\columnwidth}\raggedright
finite decorated squares, descent squares, refinement and gluing diagrams\strut
\end{minipage}\tabularnewline
\begin{minipage}[t]{0.47\columnwidth}\raggedright
finite graphs\strut
\end{minipage} & \begin{minipage}[t]{0.47\columnwidth}\raggedright
support graphs, finite-dimensional cycle/cochain complexes for the drive module\strut
\end{minipage}\tabularnewline
\begin{minipage}[t]{0.47\columnwidth}\raggedright
finite audit records\strut
\end{minipage} & \begin{minipage}[t]{0.47\columnwidth}\raggedright
provenance, replay traces, source identifiers, threshold checks, nonclaims\strut
\end{minipage}\tabularnewline
\bottomrule
\end{longtable}

\subsubsection{Permitted Extension Hosts}\label{sn:2.1.2}

The following finite or scoped extension hosts are permitted in the paper. Each is named and scoped under its own rules, and used only for the specific theorem groups indicated. No theorem in the paper uses an extension host implicitly; every appearance is gated by an explicit declaration on the surrounding judgment.

\[
\mathsf{Host}\;=\;\bigl\{\mathsf H_{\mathrm{fin}},\,\mathsf H_{\mathrm{prob}},\,\mathsf H_{\mathrm{graph}},\,\mathsf H_{\mathrm{AI}},\,\mathsf H_{\mathrm{PICA}},\,\mathsf H_{\mathrm{CantorShell}},\,\mathsf H_{\mathrm{DC}},\,\mathsf H_{\mathrm{instr}}\bigr\}.
\]

\begingroup
\setlength{\LTleft}{0pt}
\setlength{\LTright}{0pt}
\setlength{\tabcolsep}{4pt}
\renewcommand{\arraystretch}{1.08}
\begin{longtable}{@{}>{\raggedright\arraybackslash}p{0.21\linewidth}>{\raggedright\arraybackslash}p{0.35\linewidth}>{\raggedright\arraybackslash}p{0.34\linewidth}@{}}
\toprule
Host & Content & Scope \\
\midrule
\endfirsthead
\toprule
Host & Content & Scope \\
\midrule
\endhead
\bottomrule
\endfoot
\bottomrule
\endlastfoot
$\mathsf H_{\mathrm{fin}}$ & finite sets, maps, relations, records, diagrams, graphs, audit records & base host for typed-calculus theorems \\
$\mathsf H_{\mathrm{prob}}$ & finite stochastic kernels on finite state spaces, finite information quantities & Markov closure-deficit theorem, under explicit kernel/partition assumptions \\
$\mathsf H_{\mathrm{graph}}$ & finite graphs, finite cycle bases, finite-dimensional (finitely generated free) cochain/coboundary complexes & graph/cohomology drive module of Section \hyperref[sn:8]{8} \\
$\mathsf H_{\mathrm{AI}}$ & finite concrete domains, finite abstract domains, finite abstraction/concretization pairs, finite sound transformers & abstract-interpretation model family of Section \hyperref[sn:12]{12} \\
$\mathsf H_{\mathrm{PICA}}$ & finite PICA realization records (rows, columns, statuses, source records, audit restrictions) & PICA model-realization theorem of Section \hyperref[sn:13]{13} \\
$\mathsf H_{\mathrm{CantorShell}}$ & audited Cantor shell records (base/target packages, interfaces, audit) & Cantor strict-bridge realization of Section \hyperref[sn:13]{13} \\
$\mathsf H_{\mathrm{DC}}$ & finite decorated double-category fragment (objects, horizontal and vertical arrows, decorated squares, finite strict pasting) & secondary high-structure semantics of Section \hyperref[sn:14]{14} \\
$\mathsf H_{\mathrm{instr}}$ & instrument records, visibility maps, rotating-audit chains & instrument and self-reference theorems of Section \hyperref[sn:11]{11} \\
\end{longtable}
\endgroup

The core and finite extension hosts use finite carriers and records: \(\mathsf H_{\mathrm{prob}}\) uses finite state spaces and finite kernels, \(\mathsf H_{\mathrm{graph}}\) uses finite vertex and edge sets, \(\mathsf H_{\mathrm{AI}}\) uses finite abstract domains and finite sound transformers, and so on. An extension host never licenses an infinitary construction, except the quantification over macro-kernel, cochain and cycle spaces and the arbitrary-set and arbitrary-poset forms of Theorems \hyperref[cl:thm:6]{6}, \hyperref[cl:thm:7]{7}, \hyperref[cl:thm:10]{10}--\hyperref[cl:thm:14]{14} and \hyperref[cl:thm:26]{26}--\hyperref[cl:thm:28]{28} noted in Subsubsection \hyperref[sn:2.1.1]{2.1.1}. The Cantor host is the one host whose underlying system is not finite: it packages finite records over the audited shell of a continuous Cantor system, its strictness input comes from the cited Cantor theorem, and the calculus does not assert that the shell itself is finite.

\subsubsection{Excluded Hosts}\label{sn:2.1.3}

The following hosts are excluded. A claim that depends on any of them, however informally stated, falls outside the formal scope of this paper.

\begin{itemize}
\tightlist
\item
  \textbf{Unrestricted infinite reflection towers.} No theorem in the paper requires an infinite tower of theories, languages, or instruments closed under reflection. Every audit chain used in Section \hyperref[sn:11]{11} is finite, and the rotating-audit theorem is stated for chains with strictly increasing levels in the three-element \(\mathsf{Lev}\), so its hypothesis forces \(N\le1\).
\item
  \textbf{Unrestricted semantic universes.} No theorem assumes a model class closed under arbitrary semantic operations. The model families of Section \hyperref[sn:13]{13} are scoped finite hosts and are accepted only as model realizations of named fragments of the calculus.
\item
  \textbf{Empirical systems without bridge records.} No theorem appeals to a physical, biological, or computational system unless its content has been packaged into a finite record with explicit bridge data, source records, audit, and nonclaims. An empirical-bridge assertion in this paper is admissible only as a finite record on a named host, never as a bare claim about an external system.
\item
  \textbf{Untyped total six-symbol products.} No definition or theorem assumes a total binary operation \(*:\mathbb P^2\to\mathbb P\) on the primitive labels, and no later object is interpreted as a primitive product. The later no-total-algebra theorem is stated against the typed witness, update, profile, and status requirements of the calculus, so subsequent uses of pair notation \((P_i,P_j)\) or directed-cell notation \(P_i\leftarrow P_j\) are typed bridge judgments rather than algebraic products.
\end{itemize}

The host inventory above, together with these exclusions, fixes what counts as an admissible mathematical environment for the paper. Section \hyperref[sn:3]{3} builds the calculus over this inventory; Section \hyperref[sn:4]{4} defines its central judgment families; later sections invoke individual extension hosts only where named.

\subsection{Primitive Labels and Levels}\label{sn:2.2}

\subsubsection{Primitive Set}\label{sn:2.2.1}

The primitive set of the calculus is the fixed finite set

\[
\mathbb P \;=\; \{P_1,P_2,P_3,P_4,P_5,P_6\}.
\]

The labels \(P_1,\ldots,P_6\) are role names. They are not elements of an interaction algebra: there is no total binary operation on \(\mathbb P\) in this paper, and the later no-total-algebra theorem makes that exclusion formal. Pair notation \((P_i,P_j)\) names a primitive pair only; it does not determine status, effect, claim strength, visibility, or audit. Directed-cell notation \(P_i\leftarrow P_j\) refers to a typed judgment whose admissibility depends on the witness, update, profile, level, visibility, threshold, audit, and defect data attached to the judgment, not on the primitive labels alone.

The set \(\mathbb P\) is reserved for primitive labels for the duration of the paper. Where a probability kernel and a primitive label appear in the same local context, the kernel is written with an explicit subscript or as a named transition operator so that no expression \(P\cdot P'\) or \(Px\) is ambiguous between a primitive label and a kernel application.

\subsubsection{Informal Role Gloss}\label{sn:2.2.2}

Each primitive carries a stable role description, fixed for the entire paper. The descriptions are role labels, not propositions: they identify the kind of typed witness, update, defect, and audit data that a primitive participates in, and they are not theorems about what a primitive does in any particular host.

\begin{itemize}
\tightlist
\item
  \(P_1\): descent, closure of updates, rewrite compatibility, transport of content. Witness data record descent successes, descent obstructions, and rewrite/closure deficits; updates install descended maps or rewrites.
\item
  \(P_2\): representability, gating, admissibility, support control. Witness data record representability or its failure, support, capacity, and gate decisions; updates mark representable or non-representable status and adjust support.
\item
  \(P_3\): route mismatch, holonomy, critical pairs, pasting. Witness data record commutator failures, holonomy values, critical-pair mismatches, and pasting defects; updates record these as typed defect entries.
\item
  \(P_4\): refinement, staging, gluing, local-to-global structure. Witness data record stages, refinements, covers, and local/global obstructions; updates refine the lens, advance the stage, or record a gluing defect.
\item
  \(P_5\): packaging, completion, objecthood, idempotence. Witness data record package maps, quotients, completion outcomes, and fixed points; updates install a package, apply a completion, or record a fixed-point residual.
\item
  \(P_6\): audit, provenance, replay, source-of-truth, drive certification. Witness data record audit checks, provenance, source identity, threshold checks, and (where present) drive certificates; updates add an audit record, commit a source, or certify drive or no-drive.
\end{itemize}

The audit face \(P_6\) and the drive face \(\mathrm{P6}_{\mathrm{drive}}\) are distinguished: a generic \(P_6\) audit witness is not a drive certificate. Section \hyperref[sn:8]{8} isolates the drive face inside the graph/cohomology module, and the no-drive theorems there are stated for \(\mathrm{P6}_{\mathrm{drive}}\) specifically. The judgment discipline of the calculus preserves this separation throughout.

\subsubsection{Levels}\label{sn:2.2.3}

The finite level set used by the paper is

\[
\mathsf{Lev} \;=\; \{\mathsf{Beh},\,\mathsf{Ver},\,\mathsf{Str}\},
\]

with the following roles.

\begin{longtable}[]{@{}ll@{}}
\toprule
\begin{minipage}[b]{0.47\columnwidth}\raggedright
Level\strut
\end{minipage} & \begin{minipage}[b]{0.47\columnwidth}\raggedright
Role\strut
\end{minipage}\tabularnewline
\midrule
\endhead
\begin{minipage}[t]{0.47\columnwidth}\raggedright
\(\mathsf{Beh}\)\strut
\end{minipage} & \begin{minipage}[t]{0.47\columnwidth}\raggedright
behavioral level: model traces, runs, observed dynamics, raw data\strut
\end{minipage}\tabularnewline
\begin{minipage}[t]{0.47\columnwidth}\raggedright
\(\mathsf{Ver}\)\strut
\end{minipage} & \begin{minipage}[t]{0.47\columnwidth}\raggedright
verification level: instruments, audits, checks, replay, certification\strut
\end{minipage}\tabularnewline
\begin{minipage}[t]{0.47\columnwidth}\raggedright
\(\mathsf{Str}\)\strut
\end{minipage} & \begin{minipage}[t]{0.47\columnwidth}\raggedright
structural level: theorems, objects, packages, interfaces, definitional content\strut
\end{minipage}\tabularnewline
\bottomrule
\end{longtable}

The levels are ordered \(\mathsf{Beh}<\mathsf{Ver}<\mathsf{Str}\); the rotating-audit chains of Section \hyperref[sn:11]{11} compare instrument levels in this order, and ``above'' and ``highest level'' refer to it. Every directed-cell, pair-observable, promotion, and claim judgment in Sections \hyperref[sn:4]{4} and beyond carries an explicit level annotation drawn from \(\mathsf{Lev}\). A judgment whose compared level tags differ is admissible only when an explicit level bridge is recorded; for a directed cell these are the tags of \(\lambda\), \(L\), \(W_j\), \(U_i\) (step 5 of Subsubsection \hyperref[sn:6.3.1]{6.3.1}), for a claim those of its level profile and use-set (Subsubsection \hyperref[sn:11.1.2]{11.1.2}), and for other families those their schemas name. The admissible level bridges are

\[
\texttt{translation},\quad\texttt{lift},\quad\texttt{forgetting},\quad\texttt{reflection},
\]

each interpreted as a finite typed map between level-tagged records. A directed cell whose level tags differ and whose level record \(L\) carries no level bridge is malformed (Subsubsection \hyperref[sn:6.3.1]{6.3.1}, step 5) and receives no status. For other judgment families, the applicable profile or judgment schema checks any required level bridge; a missing bridge makes the judgment inadmissible. The claim classifier returns \(\texttt{outside\_scope}\) exactly when the claim's host tag is not in \(\mathsf{Hosts}_I\), a level tag of its \(\mathsf{LevelProfile}\) is not in \(\mathsf{Levels}_I\), \(\mathsf{Use}(\varphi)\not\subseteq\mathsf{Scope}_I\), or the claim type is not in \(\mathsf{ClaimTypes}_I\); the cell classifier returns it when the host tag of \(C\), a level tag of \(L\), or part of its content lies outside \(\mathsf{Scope}_I\).

Levels are not a substitute for the channel/activation distinction: a primitive may be associated with a level annotation without producing an active witness at that level, and no theorem in this paper infers an active witness from a level tag alone. Levels are not a substitute for instruments either: a level fixes which kind of content a judgment refers to, while an instrument fixes which content is visible, suppressed, or outside scope and what claim strengths are admissible. Both are required.

\subsection{Claim Strengths and Status Discipline}\label{sn:2.3}

\subsubsection{Claim Strengths}\label{sn:2.3.1}

Every numbered statement in this paper carries a claim strength, drawn from a finite taxonomy: theorem-grade (proved in the calculus under its stated assumptions), proposition-grade (following directly from definitions and earlier theorems), schema-grade (a definition, judgment form or rule, not a theorem), model-only (valid in a named model host), countermodel-grade (a finite construction together with the no-go it supports), future-work (plausible but not proved here), and excluded (not asserted). Claim strengths apply to theorems, propositions, definitions, classifier rules, model-realization statements, countermodels, and the contents of the limitations section. Countermodel-grade applies to the unnumbered sheets \(\mathsf{CM}_1\)--\(\mathsf{CM}_{12}\); numbered results proved by a finite countermodel are recorded as theorem-grade. Model-only statements are also called model-grade; the two terms name the same strength. Appendix \hyperref[sn:I]{I} also uses three refinements of these strengths: theorem-grade with assumptions (a theorem-grade statement whose assumptions include finite probabilistic data, Theorem \hyperref[cl:thm:8]{8}), theorem-grade example (a theorem-grade explicit instance, Theorem \hyperref[cl:thm:22]{22}) and audit-grade (the validation record of Proposition \hyperref[cl:prop:35]{35}). A result's recorded verdict is \(\texttt{accepted}\) for theorem-grade and proposition-grade results, \(\texttt{model\_realization}\) for model-grade results and \(\texttt{audit\_report}\) for the audit-grade record; only \(\texttt{accepted}\) is a value of the claim-status family. The meaning of each strength, and the strength and verdict of each numbered result, are recorded in Appendix \hyperref[sn:I]{I}.

A theorem-grade statement does not become more general by being labelled theorem-grade; it remains scoped by the host, instrument, package, profile, visibility, threshold, audit, and nonclaim assumptions on its judgment line (Appendix \hyperref[sn:I]{I}). A model-only statement does not gain universal force by being placed beside a theorem-grade statement; the surrounding model-realization discipline of Section \hyperref[sn:13]{13} keeps the two distinct.

Throughout, a statement ``\(X\not\Rightarrow Y\)'' means that some admissible instance satisfies \(X\) and fails \(Y\); the indices, profiles and records of the instance are named where the statement is proved. The countermodel-grade label is reserved for statements of this form supported by a finite explicit construction that exhibits an instance of \(X\) together with a failure of \(Y\). The countermodels of Section \hyperref[sn:9]{9} and Appendix \hyperref[sn:C]{C} have this strength (in \(\mathsf{CM}_6\) (\texttt{blocks}) and \(\mathsf{CM}_{12}\) (\texttt{destroys}) the universal clause of Subsubsection \hyperref[sn:4.6.1]{4.6.1} is supplied by Theorems \hyperref[cl:thm:24]{24} and \hyperref[cl:thm:18]{18}, and the construction exhibits the antecedent; in \(\mathsf{CM}_2\) the second conjunct of \texttt{requires\_bridge} comes from the bridge record of Subsubsection \hyperref[sn:9.3.5]{9.3.5}); the numbered no-go and separation results of Sections \hyperref[sn:6]{6}, \hyperref[sn:7]{7}, \hyperref[sn:8]{8} and \hyperref[sn:10]{10} are theorem-grade, several of them proved by a finite countermodel.

\subsubsection{Status Families}\label{sn:2.3.2}

The calculus uses several distinct status families, one per judgment family. They are introduced informally here and defined fully in Sections \hyperref[sn:4]{4} and \hyperref[sn:5]{5}. The families are:

\begin{itemize}
\tightlist
\item
  role-channel statuses, classifying the standing of a primitive role channel under an instrument, with values such as \(\texttt{active\_projection}\), \(\texttt{below\_threshold}\), \(\texttt{collapsed\_boundary\_case}\), \(\texttt{absent\_with\_record}\), and \(\texttt{inapplicable\_outside\_scope}\) inherited from the channel discipline of Foundations II;
\item
  directed-cell statuses, classifying typed directed-cell judgments \(P_i\leftarrow P_j\) with values such as \(\texttt{action}\), \(\texttt{implicit}\), \(\texttt{trivial}\), \(\texttt{blocked}\), \(\texttt{undefined\_circular}\), \(\texttt{below\_threshold}\), \(\texttt{collapsed}\), \(\texttt{absent}\), and \(\texttt{outside\_scope}\);
\item
  pair-observable statuses, classifying unordered pair observables, with values \(\texttt{real}\), \(\texttt{real\_provisional}\), \(\texttt{real\_duplicated}\), and \(\texttt{blocked}\);
\item
  promotion statuses, classifying promotion-bridge judgments, with values \(\texttt{candidate}\), \(\texttt{accepted}\), \(\texttt{strict}\), \(\texttt{non\_strict}\), \(\texttt{failed\_descent}\), \(\texttt{failed\_stability}\), \(\texttt{failed\_audit}\), \(\texttt{failed\_no\_smuggling}\), \(\texttt{local\_only}\), \(\texttt{globally\_obstructed}\), and \(\texttt{outside\_scope}\);
\item
  claim statuses, classifying instrument-indexed claim records, with values \(\texttt{accepted}\), \(\texttt{rejected}\), \(\texttt{provisional}\), \(\texttt{blocked}\), \(\texttt{outside\_scope}\), \(\texttt{absent\_with\_record}\), \(\texttt{below\_threshold}\), \(\texttt{failed\_audit}\), and \(\texttt{undefined\_circular}\);
\item
  gate statuses, classifying individual gate checks during promotion or claim acceptance, with values \(\texttt{pass}\), \(\texttt{fail}\), \(\texttt{not\_required}\), \(\texttt{not\_checked}\), and \(\texttt{outside\_scope}\), together with the strictness-gate refinements defined later;
\item
  square statuses, classifying decorated-square fillers in the secondary high-structure module of Section \hyperref[sn:14]{14}, with values \(\texttt{exact}\), \(\texttt{lax}\), \(\texttt{obstructed}\), \(\texttt{blocked}\), \(\texttt{outside\_scope}\), \(\texttt{failed\_audit}\), \(\texttt{undefined\_circular}\), and \(\texttt{provisional}\).
\end{itemize}

The discipline of the paper is that these families are kept apart. No theorem in the paper infers a value in one family from a value in another without an explicit bridge rule. In particular, role-channel activity does not by itself imply directed-cell action, directed-cell action does not imply pair-observable realness, claim acceptance under one instrument does not imply instrument-free truth or acceptance under another instrument, and gate passage does not imply theorem-grade strength of a surrounding statement. Section \hyperref[sn:6]{6} includes the formal status-family separation theorem that closes off the most common collapses.

\subsubsection{No-Overreading Rule}\label{sn:2.3.3}

The visibility discipline of the paper is summarized informally as the no-overreading rule:

\begin{quote}
An accepted claim under instrument \(I\) may use only the content that \(I\) marks visible, together with content that an admissible bridge under \(I\) has connected to visible content. Suppressed content not connected by an admissible bridge cannot support an accepted claim, even when the suppressed content is in fact correct. The test runs on \(\mathsf{Use}(\varphi)\): \(\mathsf{Refs}(\varphi)\) is included directly, while the evidence, threshold and audit fields contribute the one-step content references of Subsubsection \hyperref[sn:5.5.2]{5.5.2}. A suppressed dependency outside this use-set is not tested by the visibility rule; listing it in \(\mathsf{Refs}(\varphi)\) brings it within that test.
\end{quote}

The rule is not a metaphysical statement about truth. Suppressed content is not asserted to be false, and content outside the scope of \(I\) is not asserted to be unknowable. The rule is a record-level constraint: an instrument's acceptance verdict is admissible only when the use-set of the claim has been audited against the instrument's visibility map and bridge data, and any suppressed dependency without a bridge invalidates acceptance under \(I\). The formal statement, with the definition of the overreading defect and the precise acceptance rule, appears in Section \hyperref[sn:11]{11}; the rule is invoked in earlier sections only at this informal level.

\subsubsection{Three Vocabularies}\label{sn:2.3.4}

The paper uses three vocabularies that are easy to conflate. The table places them side by side.

\begingroup
\setlength{\tabcolsep}{4pt}
\renewcommand{\arraystretch}{1.08}
\begin{center}
\begin{tabularx}{0.96\linewidth}{@{}>{\raggedright\arraybackslash}p{0.17\linewidth}>{\raggedright\arraybackslash}X>{\raggedright\arraybackslash}p{0.2\linewidth}@{}}
\toprule
Vocabulary & What it records, and its values & Where \\
\midrule
Claim strength & The logical status of a numbered statement of the paper, not its importance: theorem-grade, proposition-grade, schema-grade, model-only, countermodel-grade, future-work, or excluded. Appendix \hyperref[sn:I]{I} lists the model theorems as model-grade and Proposition \hyperref[cl:prop:35]{35} as audit-grade. & Subsubsection \hyperref[sn:2.3.1]{2.3.1}; Appendix \hyperref[sn:I]{I} \\
Claim status & How an instrument's classifier adjudicates a finite claim record: $\texttt{accepted}$, $\texttt{rejected}$, $\texttt{provisional}$, $\texttt{blocked}$, $\texttt{outside\_scope}$, $\texttt{absent\_with\_record}$, $\texttt{below\_threshold}$, $\texttt{failed\_audit}$, or $\texttt{undefined\_circular}$. & Subsubsections \hyperref[sn:4.8.2]{4.8.2}, \hyperref[sn:5.1.5]{5.1.5} and \hyperref[sn:5.6.5]{5.6.5} \\
Verdict of a result & The verdict in the instrument-indexed form of a numbered result: $\texttt{accepted}$ for a theorem, $\texttt{model\_realization}$ for a model theorem, and $\texttt{audit\_report}$ for Proposition \hyperref[cl:prop:35]{35}, which is not a judgment verdict but names the recorded outcome of the validator run of that proposition. & Subsection \hyperref[sn:13.1]{13.1}; Appendix \hyperref[sn:I.2]{I.2} \\
\bottomrule
\end{tabularx}
\end{center}
\endgroup

\subsection{Paper Scope Statement}\label{sn:2.4}

\subsubsection{Formal Scope}\label{sn:2.4.1}

All judgments, theorems, propositions, schemas, model realizations, and countermodels in this paper are interpreted inside

\[
\mathbf{BirdInt}^{\mathrm{aud}}_{\mathrm{fin}}
\]

or inside one of the explicitly named extension hosts of Subsubsection \hyperref[sn:2.1.2]{2.1.2}. The body states each theorem in plain mathematical form, and that statement is the theorem: it holds under its mathematical hypotheses alone. Readers interested only in the mathematics can skip the rest of this subsection and Appendix \hyperref[sn:I]{I}: no proof in Sections \hyperref[sn:6]{6}--\hyperref[sn:14]{14} uses the instrument-indexed forms. The instrument-indexed form of each theorem, listed in Appendix \hyperref[sn:I.2]{I.2}, records the result under any admissible instrument \(I\) satisfying the following conditions:

\begin{enumerate}
\def\labelenumi{\arabic{enumi}.}
\tightlist
\item
  the associated claim record is well formed under admissible \(\Gamma,\mathcal T\);
\item
  every record referred to by its content is present in \(\Gamma\) or \(\mathcal T\) rather than represented by an audited absence;
\item
  the claim's host and level tags lie in \(\mathsf{Hosts}_I\) and \(\mathsf{Levels}_I\) and its type in \(\mathsf{ClaimTypes}_I\);
\item
  its use-set lies in \(\mathsf{Visible}_I\), its own verdict reference \(\ulcorner\varphi,I\urcorner\) (a record of \(\mathsf{Cont}\) whose two reference fields name \(\varphi\) and \(I\); Subsubsection \hyperref[sn:5.6.5]{5.6.5}) is not reachable from \(\mathsf{Use}(\varphi)\) by resolved references (reachability in the reference graph of Subsubsection \hyperref[sn:5.6.5]{5.6.5}, whose edges go from each record to every record named by one of its resolved reference fields), and, if its type is \(\texttt{soundness}\), it contains no verdict reference \(\ulcorner\psi,I'\urcorner\) with \(\mathsf{Level}(I')\ge\mathsf{Level}(I)\);
\item
  its evidence references are source records admitted by \(\mathsf{EvidencePolicy}_I\) and \(\mathsf{SourcePolicy}_I\), and each passes \(\Theta_{\mathrm{source}}\) under the default entry of \(\mathsf{SourcePolicy}_I\) (Subsubsection \hyperref[sn:5.3.6]{5.3.6});
\item
  the check rules of \(I\) for its type reject exactly when the recomputed finite witnesses falsify the conclusion;
\item
  its thresholds are exact tests on the recomputed witnesses and each evaluates to \(\texttt{pass}\) for this claim instance;
\item
  no partial evidence is declared;
\item
  its audit defect is \(\texttt{audit\_passes}\).
\end{enumerate}

The names \(I_{\mathrm{fin}}\), \(I_{\mathrm{prob}}\), \(I_{\mathrm{graph}}\), \(I_{\mathrm{instr}}\), \(I_{\mathrm{AI}}\) and \(I_{\mathrm{DC}}\) denote any such instrument for the corresponding host; \(I_{\mathrm{meta}}\) names the instrument under which the validator run of Proposition \hyperref[cl:prop:35]{35} is recorded, whose outcome \(\texttt{audit\_report}\) is not a claim status. Sections \hyperref[sn:9]{9}, \hyperref[sn:12]{12} and \hyperref[sn:13]{13} also use \(I_{\mathrm{fin}}\), \(I_{\mathrm{graph}}\), \(I_{\mathrm{instr}}\) and \(I_{\mathrm{AI}}\) for particular instruments, the judging instrument of a countermodel dependency (Subsubsection \hyperref[sn:9.1.1]{9.1.1}) or the instrument auditing a model host (Subsubsections \hyperref[sn:12.1.1]{12.1.1} and \hyperref[sn:13.4.1]{13.4.1}); those instruments are fixed by their sections and need not meet conditions 1--9.

For a result quantified over records, the instrument-indexed form is read per instance. Fix an admissible \(\mathfrak D\) and a tuple of its records meeting the hypotheses, and extend \(\mathfrak D\) by a claim record \(\varphi\), its audit record and its evidence sources (as for Theorem \hyperref[cl:thm:22]{22} in Appendix \hyperref[sn:I.2]{I.2}), where the content of \(\varphi\) is the conclusion for that tuple and \(\mathsf{Refs}(\varphi)\) is that tuple together with its recomputed witnesses (so that, by Subsubsection \hyperref[sn:3.2.1]{3.2.1}, \(\mathsf{Use}(\varphi)\) also contains \(\varphi\), its audit record and its evidence sources). The form asserts that every such instrument for this extension accepts \(\varphi\).

Such instruments require finite witnesses with decidable equality: for Theorem \hyperref[cl:thm:8]{8} the form is asserted for instances whose kernel and the law \(\mu\) of \(X_t\) are rational and whose class \(\mathcal K_{\mathrm{adm}}\) is a finite record, and for Theorems \hyperref[cl:thm:15]{15}--\hyperref[cl:thm:18]{18} for coefficients with decidable equality and cycles given on a declared finite basis; for other instances only the body statement is asserted.

For a result whose conclusion is not a decidable property of finitely many records of \(\mathfrak D\), the form records its per-instance content: for Theorem \hyperref[cl:thm:3]{3}, that \(\mathsf{WFCell}(C)=\texttt{wf}\) if and only if \(C\) meets conditions 1--8 of Subsubsection \hyperref[sn:4.3.3]{4.3.3}, for each candidate cell \(C\) over \(\mathfrak D\); for Theorem \hyperref[cl:thm:4]{4}, that \(\mathsf{CellClassify}(C)\) is the first row holding at \(C\) and is not the final row; for the existential results, Theorems \hyperref[cl:thm:1]{1}, \hyperref[cl:thm:2]{2}, \hyperref[cl:thm:5]{5}, \hyperref[cl:thm:9]{9}, \hyperref[cl:thm:22]{22} and \hyperref[cl:thm:25]{25}, whose witnesses are built in an admissible instance or extension \(\mathfrak D^{+}\) (for Theorem \hyperref[cl:thm:1]{1}, the instances named in its statement, one for each pair \((P_i,P_j)\) of its last clause; for Theorem \hyperref[cl:thm:2]{2}, one for each of its nine status variants; for Theorem \hyperref[cl:thm:5]{5}, one for each of items (a)--(e)), the form is asserted in each such instance with those witnesses in the use-set. Decidability and computability are statements of the body only.

The displays of Appendix \hyperref[sn:I.2]{I.2} abbreviate this family, and a displayed \(\forall\) ranges over records of the instance. The host, package, profile, visibility, threshold, audit and nonclaim assumptions of that form are explicit and are part of the recorded form rather than an ambient context that may be relaxed without comment. Where a proof closes by saying that the instrument-indexed verdict is \(\texttt{accepted}\), it restates the result in that form; the host, instrument and record conditions it names are the assumptions of the form, not further steps of the proof.

The shorthand \(I\vdash\varphi(\mathcal T)\) for \(\Gamma;\mathcal T;I\vdash\varphi:\texttt{accepted}\) is fixed in Subsubsection \hyperref[sn:3.4.2]{3.4.2}.

\subsubsection{Main Positive Scope}\label{sn:2.4.2}

Within the formal scope above, the paper proves and presents the results listed in Subsection \hyperref[sn:1.2]{1.2} and Table \ref{tab:map}.

\subsubsection{Main Negative Scope}\label{sn:2.4.3}

The paper does not prove, and does not assume as background, any of the following:

\begin{itemize}
\tightlist
\item
  a universal exact-six theorem stating that every phenomenon, every theory, or all of mathematics reduces to the six primitive labels;
\item
  a total binary algebra \(*:\mathbb P^2\to\mathbb P\) on the primitive labels, or any equivalent reduction of typed interaction to primitive-label arithmetic;
\item
  a rule by which the primitive pair \((P_i,P_j)\) alone determines directed-cell status;
\item
  collapse of role-channel, directed-cell, pair-observable, promotion, claim, gate, or square status families;
\item
  an implication from role-channel activity to directed-cell action, or from directed-cell action to pair-observable realness;
\item
  a unique decomposition theorem stating that every theory, package, or interaction admits a canonical decomposition into the six primitives;
\item
  package-implies-closure, idempotence-implies-commutation, repeated fixed completion as strict theory growth, local-implies-global packaging, or refinement-implies-improvement;
\item
  strictness-implies-closure, strictness-implies-drive, or route-mismatch-implies-drive;
\item
  audit-is-truth, instrument-free acceptance, accepted claims over unbridged suppressed content, or rotating-audit finality;
\item
  a same-level complete self-certification of any instrument or theory, including the interpreting instrument of this paper;
\item
  empirical realization of the calculus by any external physical, biological, or computational system; empirical-bridge assertions about such systems require a finite admissible record on a named host and do not become realization theorems;
\item
  model-realization-as-universal-theorem, PICA universality, or a PICA \(6\times6\) table as a universal interaction algebra;
\item
  complete model coverage by the model families of Section \hyperref[sn:13]{13}, that is, a theorem stating that every fragment of the calculus is realized by PICA, Cantor, abstract interpretation, or graph/cohomology data;
\item
  proof-assistant certification beyond the finite audited declarations explicitly disclosed in Appendix \hyperref[sn:F]{F}. That appendix records the validation track for the listed declarations, not a claim about unlisted statements, empirical implementations, or unrestricted host extensions.
\end{itemize}

This negative scope constrains every later section: no theorem of the paper asserts any of these excluded claims or relaxes its host, instrument or audit assumptions. Section \hyperref[sn:15]{15} collects the excluded and deferred statements.

\subsection{Where Symbols Are Defined}\label{sn:2.5}

The following table lists the main symbols of the calculus and the place where each is defined. Several of them appear before their full definition, in the overview of Section \hyperref[sn:3]{3}. The turnstile \(\vdash\) appears only in judgments \(\Gamma;\mathcal T;I\vdash\varphi:\chi\) and in the shorthand \(I\vdash\varphi(\mathcal T)\) of Subsubsection \hyperref[sn:3.4.2]{3.4.2} for such a judgment with verdict \(\texttt{accepted}\); factorization of one object map through another is written \(\pi_1\factor\pi_0\). About thirty symbols carry further local meanings; each is announced again where it arises:

\begin{itemize}
\tightlist
\item
  \(A\) is an audit record, but the abstract domain in Section \hyperref[sn:12]{12} (where the host audit record is \(A_{\mathrm{AI}}\)).
\item
  \(I\) is an instrument, but a conditional mutual information in Subsubsection \hyperref[sn:5.3.1]{5.3.1}, Subsection \hyperref[sn:7.3]{7.3} and the remark after Theorem \hyperref[cl:thm:3]{3}.
\item
  \(P\) is a role label, but a probability in Subsection \hyperref[sn:7.3]{7.3} and a pressure in Subsection \hyperref[sn:13.3]{13.3}.
\item
  \(C\) is a directed-cell record, the carrier \(\mathcal P(\{a,b,c\})\) in Subsection \hyperref[sn:7.4]{7.4}, and the concrete domain in Section \hyperref[sn:12]{12}.
\item
  \(\mathrm{Ver}\) names the vertical arrows in Section \hyperref[sn:14]{14}, distinct from the level \(\mathsf{Ver}\).
\item
  ``channel'' means a role channel, except for the top-down channels of Subsection \hyperref[sn:4.7]{4.7}.
\item
  \(L\) is also a visibility bridge in Subsubsections \hyperref[sn:3.5.3]{3.5.3} and \hyperref[sn:5.5.2]{5.5.2} and the language map \(L_{n\to n+1}\) in Subsection \hyperref[sn:11.7]{11.7}.
\item
  \(V\) is also the vertex set \(V(G)\) and a patch in \(\mathsf{CM}_3\).
\item
  \(K\) is a macro kernel in Subsection \hyperref[sn:7.3]{7.3}, the Markov kernel of PICA and of the Cantor system in Section \hyperref[sn:13]{13}.
\item
  \(T\) denotes a spanning forest in the proof of Theorem \hyperref[cl:thm:17]{17}, and \(T_0,T_1\) name the base and extended Cantor theories of Subsection \hyperref[sn:13.3]{13.3}, whose packages are \(\mathcal T^{\mathrm{Cantor}}_0,\mathcal T^{\mathrm{Cantor}}_1\).
\item
  \(E\) is also the edge set \(E(G)\).
\item
  \(W\) is also the Cantor carrier in Subsection \hyperref[sn:13.3]{13.3}.
\item
  \(\gamma\) is a cycle in Section \hyperref[sn:8]{8} and the concretization map in Section \hyperref[sn:12]{12}.
\item
  \(\mu\) is the law of \(X_t\) in Subsection \hyperref[sn:7.3]{7.3} and a pair-observable profile in Subsection \hyperref[sn:4.4]{4.4}.
\item
  \(\rho\) is the closure \(\gamma\alpha\) in Section \hyperref[sn:12]{12}, the forgetting map in Subsubsection \hyperref[sn:5.3.4]{5.3.4} and the role tag in Subsubsection \hyperref[sn:5.4.2]{5.4.2}, and, as \(\rho_i\), the role-channel status in Subsubsection \hyperref[sn:4.2.1]{4.2.1}.
\item
  \(\tau\) is the lag in Subsection \hyperref[sn:7.3]{7.3}, the origin tag in Subsubsection \hyperref[sn:5.4.2]{5.4.2} and a timescale in Subsection \hyperref[sn:13.3]{13.3}.
\item
  \(U\) is the set \(\{a,b,c\}\) in Subsections \hyperref[sn:3.8]{3.8} and \hyperref[sn:7.4]{7.4} and a patch in \(\mathsf{CM}_3\), besides the updates \(U_i\).
\item
  \(\Phi\) is a package map in Subsubsection \hyperref[sn:3.3.2]{3.3.2}, a fiber in Subsubsection \hyperref[sn:5.3.3]{5.3.3}, a target scope in Subsection \hyperref[sn:11.7]{11.7} and the bijection in the proof of Theorem \hyperref[cl:thm:14]{14}.
\item
  \(\mathcal L\) is also the KL predictive loss \(\mathcal L_\tau\) of Subsection \hyperref[sn:7.3]{7.3}.
\item
  \(\sigma\) is a cell status, and the holonomy swap of \(\{u,v\}\) in \(\mathsf{CM}_2\).
\item
  \(\Delta\) is also the simplex \(\Delta(Y)\) of Subsection \hyperref[sn:7.3]{7.3} and the pressure-gap consequence record of Subsection \hyperref[sn:13.3]{13.3}.
\item
  \(R\) is role evidence \(R_i\) in Subsubsection \hyperref[sn:5.6.1]{5.6.1}, a pair's source record in Subsubsections \hyperref[sn:5.6.3]{5.6.3} and \hyperref[sn:13.2.8]{13.2.8}, a top-down channel record in Subsection \hyperref[sn:4.7]{4.7}, and a refinement certificate \(R_{k\to\ell}\) in Subsection \hyperref[sn:12.5]{12.5}; \(\mathcal R\) is the rule language of Subsubsection \hyperref[sn:3.4.1]{3.4.1}.
\item
  \(Q\) is a pair record in Subsubsection \hyperref[sn:5.5.2]{5.5.2}, a query in Subsubsection \hyperref[sn:13.2.8]{13.2.8} and the forgetting map \(Q_\ell\) in Subsection \hyperref[sn:13.3]{13.3}.
\item
  \(D\) is the effect comparator of Subsection \hyperref[sn:4.7]{4.7} and a descended update in the proof of Theorem \hyperref[cl:thm:27]{27}.
\item
  \(a\) is an element of \(U\), a \(1\)-cochain in Section \hyperref[sn:8]{8}, an abstract element in Section \hyperref[sn:12]{12} and a vertex in \(\mathsf{CM}_{10}\).
\item
  \(\mathsf{Use}(\cdot)\) is defined for claims in Subsubsection \hyperref[sn:3.2.1]{3.2.1} and for cells, pair records, promotion bridges, visibility bridges and decorated squares in Subsubsection \hyperref[sn:5.5.2]{5.5.2}.
\item
  \(M\) is the stochastic matrix of Subsection \hyperref[sn:7.3]{7.3} and a macro record in Subsection \hyperref[sn:4.7]{4.7}.
\item
  \(S\) is the carrier of object maps (Subsubsection \hyperref[sn:4.5.2]{4.5.2}, Theorem \hyperref[cl:thm:12]{12}), a subset of \(U\) in Subsection \hyperref[sn:7.4]{7.4}, the set of \(\Theta_{\in S}\) in Subsubsection \hyperref[sn:5.2.2]{5.2.2}, and the source class of a dependency \(S\mathrel{\mathsf{dep}_r}X\) in Subsection \hyperref[sn:4.6]{4.6}.
\item
  \(X\) is the carrier of Theorems \hyperref[cl:thm:6]{6}--\hyperref[cl:thm:7]{7}, the state space of Theorem \hyperref[cl:thm:8]{8}, a substrate record in Subsection \hyperref[sn:4.7]{4.7} and the target class of a dependency in Subsection \hyperref[sn:4.6]{4.6}.
\item
  \(G\) is a graph in Section \hyperref[sn:8]{8} and, with a subscript, a gate \(G_{\mathrm{suff}},G_{\mathrm{strict}},\ldots\) (Subsection \hyperref[sn:10.2]{10.2}).
\item
  \emph{Bridge} names distinct records: visibility bridge \(L\) (Subsubsection \hyperref[sn:3.5.3]{3.5.3}), level bridge (Subsubsection \hyperref[sn:2.2.3]{2.2.3}), promotion bridge \(B_{j\to j+1}\) (Subsection \hyperref[sn:4.5]{4.5}), drive bridge \(B^{\mathrm{drive}}_{3\to6}\) (Subsubsection \hyperref[sn:9.3.5]{9.3.5}), source-upgrade and proxy source bridges (Subsubsection \hyperref[sn:5.3.6]{5.3.6}), gluing bridge \(B_{\mathrm{glue}}\) (Subsubsection \hyperref[sn:9.4.5]{9.4.5}), rotating-audit bridge \(B^{\mathrm{audit}}_{n\to n+1}\) (Subsubsection \hyperref[sn:11.7.2]{11.7.2}) and instrument-transfer bridge \(B_{I\to J}\) (Subsubsection \hyperref[sn:11.9.1]{11.9.1}); admissibility is defined separately for each.
\end{itemize}

\emph{Admissibility and well-formedness.} \(\mathsf{AdmPkg}\) (Subsubsection \hyperref[sn:3.3.3]{3.3.3}), \(\mathsf{AdmProfile}\) (Subsubsection \hyperref[sn:3.6.2]{3.6.2}), instrument admissibility (Subsubsections \hyperref[sn:3.4.1]{3.4.1} and \hyperref[sn:11.1.1]{11.1.1}), \(\mathsf{AdmDomain}\) (Subsubsection \hyperref[sn:4.10.1]{4.10.1}), acyclicity of the evaluation dependency graph (Subsubsection \hyperref[sn:5.3.6]{5.3.6}, Well-foundedness; required by \(\mathsf{AdmDomain}\)), cell well-formedness and \(\mathsf{WFCell}\) (Subsubsection \hyperref[sn:4.3.3]{4.3.3}, Subsection \hyperref[sn:6.3]{6.3}), pair records (Subsubsection \hyperref[sn:4.4.1]{4.4.1}), \(\mathsf{AdmVisBridge}\) and squares (Subsubsection \hyperref[sn:5.5.2]{5.5.2}), \(\mathrm{WF}(B)\) (Subsubsections \hyperref[sn:5.6.4]{5.6.4} and \hyperref[sn:10.2.1]{10.2.1}), claim records (Subsubsection \hyperref[sn:11.1.2]{11.1.2}), \(\mathsf{AdmRotAuditBridge}\) (Subsubsection \hyperref[sn:11.7.3]{11.7.3}), \(\mathsf{AdmInstrBridge}\) (Subsubsection \hyperref[sn:11.9.1]{11.9.1}), \(\mathsf{AdmAIHost}\) (Subsubsection \hyperref[sn:12.1.1]{12.1.1}), \(\mathsf{AdmPICAReal}\) (Subsubsection \hyperref[sn:13.2.2]{13.2.2}), \(\mathsf{AdmCantorShell}\) (Subsubsection \hyperref[sn:13.3.8]{13.3.8}), \(\mathsf{AdmCohGraph}\) (Subsubsection \hyperref[sn:13.4.1]{13.4.1}).

\begingroup\small
\begin{longtable}{@{}>{\raggedright\arraybackslash}p{0.26\linewidth}>{\raggedright\arraybackslash}p{0.5\linewidth}>{\raggedright\arraybackslash}p{0.18\linewidth}@{}}
\toprule
Symbol & Meaning & Defined in \\
\midrule
\endfirsthead
\toprule
Symbol & Meaning & Defined in \\
\midrule
\endhead
\bottomrule
\endfoot
$\mathbf{BirdInt}^{\mathrm{aud}}_{\mathrm{fin}}$ & the finite audited interaction calculus, a nineteen-component tuple & §3.1 \\
$\mathsf H_{\mathrm{fin}}$, $\mathsf H_{\mathrm{prob}}$, $\mathsf H_{\mathrm{graph}}$, \dots & the base host and the permitted extension hosts & §2.1 \\
$\mathbb P=\{P_1,\ldots,P_6\}$, $\mathsf{Lev}$ & the primitive labels and the ordered level set $\mathsf{Beh}<\mathsf{Ver}<\mathsf{Str}$ & §2.2 \\
$\mathsf{Cont}$ & the finite content universe & §3.2 \\
$\mathcal T$, $\mathsf{Pkg}$ & a theory package $(Z,f,\Sigma_f,E,\mathcal A)$ and the class of packages & §3.3 \\
$I$, $\mathsf{Instr}$ & an instrument and the class of instruments & §3.4, §11.1 \\
$\Gamma;\mathcal T;I\vdash\varphi:\chi$ & an instrument-relative judgment with verdict $\chi$ & §3.4.2, §4.8 \\
$\mathsf{Vis}$, $\mathsf{Suppressed}_I$, $\mathcal L$ & visibility maps, suppressed content, and a set of visibility bridges & §3.5 \\
$\lambda$, $\mathsf{Prof}$ & a profile and the profile family & §3.6 \\
$\mathsf{Wit}(P_j)$, $\mathsf{Upd}(P_i)$ & the witness family of an informant and the update family of an actor & §3.7 \\
$\mathsf{Cell}_{ij}^{\lambda}$ & a directed cell $P_i\leftarrow P_j$ with profile $\lambda$ & §4.3 \\
$L$; $L\in\mathcal L$ (instances $L_{\mathrm{drive}}$, $L_t$); $L_{n\to n+1}$ & the level/lens/interface record of a cell or bridge; a visibility bridge (also written $L$ in its defining schema of Subsubsection \hyperref[sn:3.5.3]{3.5.3}); a rotating-audit language map & §4.3.3, §3.5.3, §11.7.2 \\
$\mathsf{PairObs}_{\{i,j\}}^{\mu}$ & a pair observable on $\{P_i,P_j\}$ with profile $\mu$ & §4.4 \\
$B_{j\to j+1}$, $\mathrm{Promote}(B)$ & a promotion bridge and its promotion judgment & §4.5, §10.1 \\
$S\mathrel{\mathsf{dep}_r}X$ & a dependency judgment of type $r$ & §4.6 \\
$\mathsf{Source}$, $\mathsf{Audit}$ & source records and audit records & §4.9 \\
$\mathrm{RoleStatus}$, $\mathrm{CellStatus}$, \dots & the status families & §§4.2--4.8, §5.1 \\
$\delta$, $\mathsf{Def}$ & a defect record and the defect families & §5.2--5.5 \\
$\mathrm{Definable}(f)$ & the preimages under $f$ of subsets of its image & §7.9 \\
$\mathrm{CD}_{\tau}^{\mu}(\Pi)$ & the Markov closure deficit, a conditional mutual information & §5.3.1, §7.3 \\
$(C,A,\alpha,\gamma)$ & concrete domain, abstract domain and Galois connection of an abstract-interpretation host & §12.1 \\
$\delta_1^{\mathrm{split}}(q,F)$ & the descent defect: pairs identified by $q$ and separated by $qF$ & §7.2 \\
$\Delta_3^{\mathrm{comp}}(E_1,E_2)$ & the route-mismatch (completion commutator) defect of two completions & §7.5, §5.3.3 \\
$\pi_1\factor\pi_0$, $\pi_1\not\factor\pi_0$ & $\pi_1$ factors, or does not factor, through $\pi_0$ on its image & §7.7 \\
$\Delta_{\mathrm{fact}}(\pi_0,\pi_1)$ & the factorization defect & §7.8, §5.4.1 \\
$\delta_{\mathrm{overread}}(I,\varphi,\mathcal L)$ & the suppressed content that a claim uses without a bridge & §5.5.1 \\
$\mathsf{Use}(\varphi)$, $\mathsf{Refs}(\varphi)$, $\mathsf{Use}(C)$ & use-set and declared references of a claim; use-set of a cell & §3.2.1, §4.3.3, §5.5.2 \\
$\ulcorner\varphi,I\urcorner$ & verdict reference; reachability in the reference graph & §5.6.5 \\
$\delta_6^{\mathrm{audit}}$, $\texttt{audit\_passes}$, $\delta_6^{\mathrm{source}}$ & audit defect and its passing value; source defect & §5.3.6 \\
$\delta_{\mathrm{vis}}$, $\delta_{\mathrm{weakbridge}}$, $\delta_{\mathrm{reqbridge}}$, $\mathsf{AdmVisBridge}$ & visibility and required-bridge defects; bridge admissibility & §5.5 \\
$\mathsf{noopU}$, $\mathsf{implicitW}$ & computed no-op and attached-witness fields of a cell & §4.3.4, §5.6.6 \\
$\mathsf{ClaimLog}(I)$, $\mathrm{Sound}(I)$ & declared verdict log of $I$; its complete self-soundness claim & §11.5.1 \\
$H^1(G)$, $\oint_\gamma a$ & first cohomology of a finite graph and the integral of a cochain along a cycle & §8.1 \\
$\mathsf{Gate}$ & the promotion gates and bridge-admissibility gates & §10.2 \\
$\mathcal M\models\mathcal S$ & a model realization of a fragment $\mathcal S$ & §13.1 \\
$\mathcal N$, NC-$n$ & a nonclaim register and the numbered nonclaims & §15.2 \\
\end{longtable}
\endgroup
 \section{The Finite Audited Domain}\label{sn:3}

This section declares the ambient domain of the calculus, before the central judgment object is defined in Section \hyperref[sn:4]{4}. The domain is a finite tuple of finite typed components: a primitive set, a level set, a host inventory, a content universe, instruments, packages, profiles, witness and update families, defect records, judgment families, status families, gates, dependencies, audits, visibility data, source records, nonclaim records, and model-realization modules. Each component is introduced as a finite schema; the theorems and classifiers that depend on these components are stated in later sections and never quoted as already established here. The definitions needed for Theorems \hyperref[cl:thm:1]{1}--\hyperref[cl:thm:4]{4} are listed in the reader's map (Subsection \hyperref[sn:1.4]{1.4}); the rest of this section is reference material.

\subsection{Domain Declaration}\label{sn:3.1}

\subsubsection{Full Tuple}\label{sn:3.1.1}

\begin{claimdefinition}{Finite audited interaction domain}

The finite audited interaction calculus is the nineteen-component signature

\[
\mathbf{BirdInt}^{\mathrm{aud}}_{\mathrm{fin}}
\;=\;
\begin{aligned}
(\,&\mathbb P,\;\mathsf{Lev},\;\mathsf{Host},\;\mathsf{Cont},\;\mathsf{Instr},\;\mathsf{Pkg},\;\mathsf{Prof},\\
&\mathsf{Wit},\;\mathsf{Upd},\;\mathsf{Def},\;\mathsf{Judg},\;\mathsf{Status},\;\mathsf{Gate},\;\mathsf{Dep},\\
&\mathsf{Audit},\;\mathsf{Vis},\;\mathsf{Source},\;\mathsf{Nonclaim},\;\mathsf{Real}\,).
\end{aligned}
\]

An instance \(\mathfrak D\) assigns to each component a finite value, either a finite set, a finite family of finite records, or a finite typed map between such families, with \(\mathbb P\), \(\mathsf{Lev}\), \(\mathsf{Judg}\) and \(\mathsf{Status}\) fixed as in Subsubsection \hyperref[sn:4.10.1]{4.10.1} and \(\mathsf{Host}\) a subfamily of the inventory of Subsubsection \hyperref[sn:2.1.2]{2.1.2}; \(\mathfrak D\in\mathbf{BirdInt}^{\mathrm{aud}}_{\mathrm{fin}}\) means that \(\mathfrak D\) is such an assignment. The components are introduced in Subsections \hyperref[sn:3.2]{3.2} through \hyperref[sn:3.7]{3.7} and Section \hyperref[sn:4]{4}; full definitions of the status, defect, gate, audit, and model-realization components appear in Sections \hyperref[sn:4]{4}, \hyperref[sn:5]{5}, \hyperref[sn:10]{10}, \hyperref[sn:11]{11}, and \hyperref[sn:13]{13}, respectively. The signature is the schema-grade declaration of what counts as the ambient domain of the calculus.

The notation \(\mathbf{BirdInt}^{\mathrm{aud}}_{\mathrm{fin}}\) refers to this declared domain throughout the paper. It is distinct from the inherited Foundations II domain \(\mathsf{FATCD}\) of finite audited typed closure descriptions: \(\mathsf{FATCD}\) is the description domain for the scoped exact-six program of that paper, and \(\mathbf{BirdInt}^{\mathrm{aud}}_{\mathrm{fin}}\) is the interaction calculus of this paper. Where Foundations II structures are reused inside the calculus, the paper records the embedding rather than renaming \(\mathsf{FATCD}\).

\end{claimdefinition}

The notation \(\mathbf{BirdInt}^{\mathrm{aud}}_{\mathrm{fin}}\) names this schema. A particular description \(\mathfrak D\in\mathbf{BirdInt}^{\mathrm{aud}}_{\mathrm{fin}}\) is an instance of it; the instance is admissible when \(\mathsf{AdmDomain}(\mathfrak D)\) holds (Subsection \hyperref[sn:4.10]{4.10}), and the theorems of the paper quantify over admissible instances. \emph{Conventions.} The words domain, calculus and central object used for \(\mathbf{BirdInt}^{\mathrm{aud}}_{\mathrm{fin}}\) all refer to this schema; the definition in Subsubsection \hyperref[sn:4.1.1]{4.1.1} restates it together with its judgment families. An instance \(\mathfrak D\) is a particular choice of the nineteen components. A universally stated theorem is read as a statement about an arbitrary admissible instance \(\mathfrak D\), and \(\mathfrak D\) is left out of the notation: a record, judgment or status written without an instance belongs to \(\mathfrak D\). Existential conclusions have the hypotheses and quantifier scopes stated in their individual results; their proofs construct the required admissible instances or extensions.

\emph{Records and typed families.} A record signature is a finite set of field names, each with a finite value type or a declared numerical type, with effective equality and every operation or order comparison used by an applicable decision procedure. A finite typed record of signature \(\sigma\) assigns to each field of \(\sigma\) a value of its type, and \(\sigma\) is its type. A finite typed family, such as \(\mathsf{Wit}(P_i)\) or \(\mathsf{Upd}(P_i)\) of Subsection \hyperref[sn:3.7]{3.7}, is a finite set of named signatures, and \(W\in\mathsf{Wit}(P_i)\) means that the signature of \(W\) is one of them; membership is therefore decided by comparing finite signatures, and whether the values of \(W\) meet a threshold or pass an audit is decided separately by the threshold and audit records of the judgment that uses \(W\).

\emph{Signatures of the running example.} The following three signatures are used in the running example of Subsection \hyperref[sn:3.8]{3.8} and in later constructions; their descriptions can be skipped until then: \(\mathsf{RouteMismatchWit}\) has the fields first completion, second completion and point, with values maps on the finite carrier and an element of it; \(\mathsf{AddAuditRecord}\) has the fields log, entry, target and effect: the entry records the two recomputed images, the target is the log \(\Lambda\) and the effect is \(e_U(\Lambda)=\Lambda\cup\{\text{entry}\}\); and \(\mathsf{CertifyDrive}\) has the fields log, entry, drive check, target and effect, with the same target and effect, the drive-check field naming the candidate drive bridge checked for the drive face \(\mathrm{P6}_{\mathrm{drive}}\); it certifies that face only when the bridge passes the required strength and audit checks.

\emph{Record identifiers.} Distinct records of \(\mathsf{Cont}\) carry distinct identifiers; a reference field names the unique record of \(\mathsf{Cont}\) carrying that identifier, and a candidate record outside \(\mathsf{Cont}\) is named by no reference field. Every record of \(\mathsf{Cont}\) carries a record identifier: the displayed field \(\mathrm{id}\) where its schema has one (\(\mathrm{id}_I\), \(\mathrm{id}_\varphi\), \(\mathrm{id}_{\mathrm{src}}\), \(\mathrm{id}_L\), the bridge identifier), and otherwise (among them theory packages, directed cells, pair records, decorated squares, audit and threshold records, absence markers, verdict references, claim logs and context records) an undisplayed identifier field that is not a reference field. Reference fields, \(\mathsf{ClaimRef}\) and verdict references name records by identifier, so a cell can hold its audit record \(A\) while \(A\) names the cell, and records with equal displayed fields but distinct identifiers, such as \(Q\) and \(Q'\) of Subsubsection \hyperref[sn:4.4.1]{4.4.1}, are distinct.

\subsubsection{Interpretation of Components}\label{sn:3.1.2}

The components of the tuple have the following roles. Each description is schema-grade: it fixes what kind of object the component is, not what theorems it satisfies.

\begin{itemize}
\tightlist
\item
  \(\mathbb P\) is the fixed finite set of primitive labels \(\{P_1,\ldots,P_6\}\) defined in Subsubsection \hyperref[sn:2.2.1]{2.2.1}. It is the underlying label alphabet of every role-channel, directed-cell, pair-observable, promotion, and dependency judgment.
\item
  \(\mathsf{Lev}=\{\mathsf{Beh},\mathsf{Ver},\mathsf{Str}\}\) is the finite level set defined in Subsubsection \hyperref[sn:2.2.3]{2.2.3}, used to tag the kind of content a judgment refers to and to gate cross-level claims through translation, lift, forgetting, or reflection bridges.
\item
  \(\mathsf{Host}\) is a finite subfamily of the inventory of permitted hosts declared in Subsubsection \hyperref[sn:2.1.2]{2.1.2}: \(\mathsf H_{\mathrm{fin}}\), \(\mathsf H_{\mathrm{prob}}\), \(\mathsf H_{\mathrm{graph}}\), \(\mathsf H_{\mathrm{AI}}\), \(\mathsf H_{\mathrm{PICA}}\), \(\mathsf H_{\mathrm{CantorShell}}\), \(\mathsf H_{\mathrm{DC}}\), \(\mathsf H_{\mathrm{instr}}\). Every theorem statement names its host.
\item
  \(\mathsf{Cont}\) is the finite content universe of the domain, defined in Subsection \hyperref[sn:3.2]{3.2}. It collects the records that claims may use: theory packages, instruments, witnesses, updates, profiles, defect records, threshold records, audit records, visibility records, level/lens/interface records, package and promotion interfaces, quotients, lenses, completion endomaps, kernels, graphs, finite families, refinements, and other finite typed records introduced as the calculus is built.
\item
  \(\mathsf{Instr}\) is the finite family of instrument records, defined in Subsection \hyperref[sn:3.4]{3.4}. An instrument fixes the visible/suppressed map, scope, claim types, threshold and audit policies, level assignment, bridge policy, source policy, and nonclaim list under which a claim is to be evaluated.
\item
  \(\mathsf{Pkg}\) is the finite family of theory packages, defined in Subsection \hyperref[sn:3.3]{3.3}, each carrying the inherited Foundations I package data \((Z,f,\Sigma_f,E,\mathcal A)\) together with the additional Foundations III typed records (visibility, source, defect, promotion, model-realization) attached around it.
\item
  \(\mathsf{Prof}\) is the finite family of profile records, defined in Subsection \hyperref[sn:3.6]{3.6}. A profile fixes per-judgment data such as actor and informant level tags, scale, regime, bridge type, instrument mode, and locality, and is required by precisely the record families listed in Subsection \hyperref[sn:3.6]{3.6}.
\item
  \(\mathsf{Wit}\) assigns to each primitive \(P_i\in\mathbb P\) its typed witness family \(\mathsf{Wit}(P_i)\), defined in Subsection \hyperref[sn:3.7]{3.7}. The informant-side data of a directed cell is required to belong to one of these families.
\item
  \(\mathsf{Upd}\) assigns to each primitive \(P_i\in\mathbb P\) its typed update family \(\mathsf{Upd}(P_i)\), defined in Subsection \hyperref[sn:3.7]{3.7}. The actor-side data of a directed cell is required to belong to one of these families.
\item
  \(\mathsf{Def}\) is the finite family of defect records, defined in Section \hyperref[sn:5]{5}. A defect record carries a defect type, a domain, a witness, a threshold, a measured value, a status consequence, and an audit/source record, and feeds the defect-to-status rules of Subsection \hyperref[sn:5.6]{5.6}.
\item
  \(\mathsf{Judg}\) is the finite family of judgment forms used by the calculus: role-channel, directed-cell, pair-observable, promotion, dependency, and instrument-indexed claim judgments, all defined in Section \hyperref[sn:4]{4}, and the model-realization judgment of Subsubsection \hyperref[sn:13.1.1]{13.1.1}.
\item
  \(\mathsf{Status}\) is the finite family of status sets, one per judgment family: role-channel, directed-cell, pair-observable, promotion, claim, gate, and square statuses. The sets are listed in Subsubsection \hyperref[sn:2.3.2]{2.3.2} and defined fully in Section \hyperref[sn:5]{5}.
\item
  \(\mathsf{Gate}\) is the finite family of gate specifications used during promotion and claim acceptance, defined in Section \hyperref[sn:10]{10}. A gate is a finite check whose \(\texttt{pass}\), \(\texttt{fail}\), \(\texttt{not\_required}\), \(\texttt{not\_checked}\), or \(\texttt{outside\_scope}\) result feeds promotion-status and claim-status decisions.
\item
  \(\mathsf{Dep}\) is the finite family of dependency relations between primitives, including positive dependencies, blocked dependencies, bridge-required dependencies, and countermodel-supported no-go dependencies, defined in Subsection \hyperref[sn:4.6]{4.6}.
\item
  \(\mathsf{Audit}\) is the finite family of audit records, defined in Subsection \hyperref[sn:4.9]{4.9}. An audit record carries a target, an instrument, replay or check data, and admissibility data; it is distinct from the inherited audit functional \(\mathcal A\) that lives inside a theory package \(\mathcal T\).
\item
  \(\mathsf{Vis}\) is the finite family of visibility maps and visibility bridges, defined in Subsection \hyperref[sn:3.5]{3.5}. A visibility map sends content to one of \(\texttt{visible}\), \(\texttt{suppressed}\), \(\texttt{outside}\), \(\texttt{unknown}\) under an instrument, and a visibility bridge is the admissible record by which a claim may use suppressed content.
\item
  \(\mathsf{Source}\) is the finite family of source records, defined in Subsection \hyperref[sn:4.9]{4.9}. A source record fixes source identity, source type, trace, and source validity, and is consulted by pair-observable and claim acceptance rules.
\item
  \(\mathsf{Nonclaim}\) is the finite family of nonclaim records carried by an instrument or attached to a judgment. A nonclaim records what the surrounding statement does not assert, in the same finite typed format as the rest of the calculus.
\item
  \(\mathsf{Real}\) is the finite family of model-realization modules, defined in Section \hyperref[sn:13]{13}: the PICA, audited Cantor shell, abstract-interpretation, and graph/cohomology realizations of named fragments of the calculus.
\end{itemize}

The tuple is finite at every component: each of \(\mathbb P\), \(\mathsf{Lev}\), \(\mathsf{Host}\), the per-instrument and per-package records, the witness and update families per primitive, and the per-judgment status sets is finite. No component is defined by an unbounded reflection scheme or a closure under arbitrary semantic operations. A claim that requires an extension of any component is not interpreted in \(\mathbf{BirdInt}^{\mathrm{aud}}_{\mathrm{fin}}\) as defined here and lies outside the formal scope of this paper. The body statements also include quantification over the possibly infinite kernel, cochain and cycle spaces of Theorems \hyperref[cl:thm:8]{8} and \hyperref[cl:thm:15]{15}--\hyperref[cl:thm:18]{18}, and the arbitrary-set and arbitrary-poset forms of Theorems \hyperref[cl:thm:6]{6}, \hyperref[cl:thm:7]{7}, \hyperref[cl:thm:10]{10}--\hyperref[cl:thm:14]{14} and \hyperref[cl:thm:26]{26}--\hyperref[cl:thm:28]{28}. These mathematical assertions do not enlarge the finite record schema; their instrument-indexed forms concern finite recorded instances.

\subsection{Content Universe}\label{sn:3.2}

\subsubsection{Finite Content Objects}\label{sn:3.2.1}

The content universe of the calculus is the finite family

\[
\mathsf{Cont}
\]

of records that claims, judgments, and audits may refer to. Every member of \(\mathsf{Cont}\) is a finite typed record built from the base host \(\mathsf H_{\mathrm{fin}}\), possibly with components drawn from a declared extension host. The schema is extended only by the finite record types introduced in this paper's declared components; each such record type is required to be a finite typed object on a named host, and no record type is defined by closure under an unbounded operation.

The standard members of \(\mathsf{Cont}\) used by this paper are the following.

\begin{longtable}[]{@{}ll@{}}
\toprule
\begin{minipage}[b]{0.47\columnwidth}\raggedright
Record kind\strut
\end{minipage} & \begin{minipage}[b]{0.47\columnwidth}\raggedright
Reference\strut
\end{minipage}\tabularnewline
\midrule
\endhead
\begin{minipage}[t]{0.47\columnwidth}\raggedright
theory package \(\mathcal T=(Z,f,\Sigma_f,E,\mathcal A)\) together with attached Foundations III records\strut
\end{minipage} & \begin{minipage}[t]{0.47\columnwidth}\raggedright
Subsection \hyperref[sn:3.3]{3.3}\strut
\end{minipage}\tabularnewline
\begin{minipage}[t]{0.47\columnwidth}\raggedright
instrument record \(I\)\strut
\end{minipage} & \begin{minipage}[t]{0.47\columnwidth}\raggedright
Subsection \hyperref[sn:3.4]{3.4}\strut
\end{minipage}\tabularnewline
\begin{minipage}[t]{0.47\columnwidth}\raggedright
primitive label \(P_i\in\mathbb P\)\strut
\end{minipage} & \begin{minipage}[t]{0.47\columnwidth}\raggedright
Subsubsection \hyperref[sn:2.2.1]{2.2.1}\strut
\end{minipage}\tabularnewline
\begin{minipage}[t]{0.47\columnwidth}\raggedright
witness \(W_i\in\mathsf{Wit}(P_i)\)\strut
\end{minipage} & \begin{minipage}[t]{0.47\columnwidth}\raggedright
Subsection \hyperref[sn:3.7]{3.7}\strut
\end{minipage}\tabularnewline
\begin{minipage}[t]{0.47\columnwidth}\raggedright
update \(U_i\in\mathsf{Upd}(P_i)\)\strut
\end{minipage} & \begin{minipage}[t]{0.47\columnwidth}\raggedright
Subsection \hyperref[sn:3.7]{3.7}\strut
\end{minipage}\tabularnewline
\begin{minipage}[t]{0.47\columnwidth}\raggedright
profile \(\lambda\)\strut
\end{minipage} & \begin{minipage}[t]{0.47\columnwidth}\raggedright
Subsection \hyperref[sn:3.6]{3.6}\strut
\end{minipage}\tabularnewline
\begin{minipage}[t]{0.47\columnwidth}\raggedright
level/lens/interface record \(L\)\strut
\end{minipage} & \begin{minipage}[t]{0.47\columnwidth}\raggedright
Subsection \hyperref[sn:4.3]{4.3}\strut
\end{minipage}\tabularnewline
\begin{minipage}[t]{0.47\columnwidth}\raggedright
visibility record \(V\), suppressed-content record, visibility bridge\strut
\end{minipage} & \begin{minipage}[t]{0.47\columnwidth}\raggedright
Subsection \hyperref[sn:3.5]{3.5}\strut
\end{minipage}\tabularnewline
\begin{minipage}[t]{0.47\columnwidth}\raggedright
threshold record \(\Theta\)\strut
\end{minipage} & \begin{minipage}[t]{0.47\columnwidth}\raggedright
Subsection \hyperref[sn:5.2]{5.2}\strut
\end{minipage}\tabularnewline
\begin{minipage}[t]{0.47\columnwidth}\raggedright
audit record \(A\), source record, replay trace\strut
\end{minipage} & \begin{minipage}[t]{0.47\columnwidth}\raggedright
Subsection \hyperref[sn:4.9]{4.9}\strut
\end{minipage}\tabularnewline
\begin{minipage}[t]{0.47\columnwidth}\raggedright
defect record \(\delta\) with type, domain, witness, threshold, measured value, status consequence\strut
\end{minipage} & \begin{minipage}[t]{0.47\columnwidth}\raggedright
Section \hyperref[sn:5]{5}\strut
\end{minipage}\tabularnewline
\begin{minipage}[t]{0.47\columnwidth}\raggedright
object map or interface \(\pi_j\), \(\pi_{j+1}\) in a promotion bridge\strut
\end{minipage} & \begin{minipage}[t]{0.47\columnwidth}\raggedright
Subsection \hyperref[sn:4.5]{4.5}\strut
\end{minipage}\tabularnewline
\begin{minipage}[t]{0.47\columnwidth}\raggedright
quotient \(q\), finite map \(F\), descended map \(F^\sharp\), completion endomap \(E\)\strut
\end{minipage} & \begin{minipage}[t]{0.47\columnwidth}\raggedright
Sections \hyperref[sn:6]{6} and \hyperref[sn:7]{7}\strut
\end{minipage}\tabularnewline
\begin{minipage}[t]{0.47\columnwidth}\raggedright
finite stochastic kernel \(K\) on a finite state space, finite information quantities\strut
\end{minipage} & \begin{minipage}[t]{0.47\columnwidth}\raggedright
Subsection \hyperref[sn:7.3]{7.3} (host \(\mathsf H_{\mathrm{prob}}\))\strut
\end{minipage}\tabularnewline
\begin{minipage}[t]{0.47\columnwidth}\raggedright
finite graph \(G\), cycle basis, cochain family \(\mathcal F\)\strut
\end{minipage} & \begin{minipage}[t]{0.47\columnwidth}\raggedright
Section \hyperref[sn:8]{8} (host \(\mathsf H_{\mathrm{graph}}\))\strut
\end{minipage}\tabularnewline
\begin{minipage}[t]{0.47\columnwidth}\raggedright
refinement, staging, gluing record \(\Delta\)\strut
\end{minipage} & \begin{minipage}[t]{0.47\columnwidth}\raggedright
Subsection \hyperref[sn:5.3]{5.3} and Section \hyperref[sn:7]{7}\strut
\end{minipage}\tabularnewline
\begin{minipage}[t]{0.47\columnwidth}\raggedright
nonclaim record \(\mathcal N\)\strut
\end{minipage} & \begin{minipage}[t]{0.47\columnwidth}\raggedright
Subsection \hyperref[sn:15.2]{15.2}\strut
\end{minipage}\tabularnewline
\begin{minipage}[t]{0.47\columnwidth}\raggedright
verdict reference \(\ulcorner\varphi,I\urcorner\), naming the verdict of a claim \(\varphi\) (or of another audit target) under an instrument \(I\)\strut
\end{minipage} & \begin{minipage}[t]{0.47\columnwidth}\raggedright
Subsubsection \hyperref[sn:5.6.5]{5.6.5}\strut
\end{minipage}\tabularnewline
\bottomrule
\end{longtable}

Every claim \(\varphi\) in the calculus has a finite use-set

\[
\mathsf{Use}(\varphi)\;\subseteq\;\mathsf{Cont}
\]

defined as follows. The Content field of a claim record (Subsubsection \hyperref[sn:11.1.2]{11.1.2}) carries a finite declared set \(\mathsf{Refs}(\varphi)\) of record identifiers, and the record is well formed only if every record identifier occurring in the statement of \(\varphi\) lies in \(\mathsf{Refs}(\varphi)\). Then \(\mathsf{Use}(\varphi)\) consists of \(\{\varphi\}\cup\mathsf{Refs}(\varphi)\) together with, for a claim record, the one-step content references (Subsubsection \hyperref[sn:5.5.2]{5.5.2}) of its \(\mathsf{EvidenceRefs}\), threshold and audit fields: each named source record with the records named in its trace, and each threshold or audit record with the records named by its reference fields, except the instrument named in \(\mathsf{CheckRules}\). The claim record \(\varphi\) itself lies in \(\mathsf{Use}(\varphi)\) even when its audit slot holds an audited absence marker, so scope and visibility conditions apply to the claim record too. Visibility, suppression, source-of-truth, and overreading checks in later sections operate on \(\mathsf{Use}(\varphi)\), not on the entire content universe. A claim whose use-set contains a record that is not finite, not typed, or not on a permitted host is not admissible in \(\mathbf{BirdInt}^{\mathrm{aud}}_{\mathrm{fin}}\).

\subsubsection{Typed Content}\label{sn:3.2.2}

Every record in \(\mathsf{Cont}\) is typed in four respects.

\begin{itemize}
\tightlist
\item
  \textbf{Host.} Each record carries a host tag drawn from \(\mathsf{Host}\). The host tag fixes which finite mathematical universe the record lives in: \(\mathsf H_{\mathrm{fin}}\) for the base host, \(\mathsf H_{\mathrm{prob}}\) for finite stochastic content, \(\mathsf H_{\mathrm{graph}}\) for finite graph and cohomology content, \(\mathsf H_{\mathrm{AI}}\) for finite abstract-interpretation content, \(\mathsf H_{\mathrm{PICA}}\) and \(\mathsf H_{\mathrm{CantorShell}}\) for the model-realization hosts, \(\mathsf H_{\mathrm{DC}}\) for the high-structure fragment, and \(\mathsf H_{\mathrm{instr}}\) for instrument-side content.
\item
  \textbf{Level.} Each record carries a level tag drawn from \(\mathsf{Lev}=\{\mathsf{Beh},\mathsf{Ver},\mathsf{Str}\}\). The level tag is a declared field; the glosses of Subsubsection \hyperref[sn:2.2.3]{2.2.3} describe typical content, and the calculus uses level tags only through equality (which decides when a level bridge is needed) and the order \(\mathsf{Beh}<\mathsf{Ver}<\mathsf{Str}\). For an instrument, \(\mathsf{Level}_I\) ranks which instruments may audit which (Subsection \hyperref[sn:11.7]{11.7}); it does not assert that the instrument's records are behavioral, verification or structural content. A judgment that has differing compared level tags (Subsubsection \hyperref[sn:2.2.3]{2.2.3}) is admissible only with an explicit level bridge (translation, lift, forgetting, or reflection).
\item
  \textbf{Instrument.} Each record is interpreted under the instrument \(I\) that the surrounding judgment names. The instrument fixes whether the record is visible, suppressed, outside-scope, or unknown for the purpose of the claim, and what bridges may connect it to visible content. A record may be in \(\mathsf{Cont}\) globally and yet be suppressed under a particular instrument; suppression is a relation between a record and an instrument, not a property of the record alone, and it does not assert that the suppressed record is false.
\item
  \textbf{Profile.} Each directed-cell, pair-observable, promotion-bridge, rotating-audit-bridge, instrument-transfer and top-down channel record is tagged by a profile \(\lambda\in\mathsf{Prof}\), and each claim record by a level profile (Subsection \hyperref[sn:3.6]{3.6}); a profile fixes per-judgment data including actor and informant level tags, scale, regime, bridge type, instrument mode, and locality. The profile of a judgment together with its host, level, and instrument tags determines the admissibility of every record in its use-set.
\end{itemize}

Records carry host, level and instrument context, and profiles only where their record family requires one; without the required context, a statement is a schema rather than a judgment and cannot support a claim.

\subsection{Theory Packages (Reference)}\label{sn:3.3}

\emph{Used by:} Theorems \hyperref[cl:thm:3]{3} and \hyperref[cl:thm:23]{23} use theory packages in their statements, through the judgment context \(\Gamma;\mathcal T;I\).

\subsubsection{Package Record}\label{sn:3.3.1}

A theory package is the inherited Foundations I object

\[
\mathcal T \;=\; (\,Z,\;f,\;\Sigma_f,\;E,\;\mathcal A\,)
\]

together with the additional Foundations III typed records that the calculus attaches to it. The inherited components have the inherited meaning: \(Z\) is a finite carrier or scoped comparison domain, \(f\) is a finite lens or coarse-description map of \(Z\), \(\Sigma_f\) is the expressive content (finite definability) induced by \(f\), \(E\) is a completion or packaging endomap, and \(\mathcal A\) is the audit functional inherited from Foundations I. The inherited form is preserved verbatim throughout the paper; no Foundations III usage redefines \(\mathcal T\), \(f\), \(\Sigma_f\), \(E\), or \(\mathcal A\) as different objects.

In Foundations III, a package record carries the inherited five-tuple together with finite typed records for the additional content the interaction calculus depends on:

\begin{itemize}
\tightlist
\item
  a host tag \(h_{\mathcal T}\in\mathsf{Host}\) fixing the host on which the package lives;
\item
  a level tag \(\ell_{\mathcal T}\in\mathsf{Lev}\) fixing whether the package describes behavioral, verification, or structural content;
\item
  a finite visibility profile \(V_{\mathcal T}\) recording which fields of the package are designated visible, suppressed, outside-scope, or unknown by the instruments that consult it;
\item
  a finite source record \(\mathrm{src}_{\mathcal T}\) recording source identity, source type, trace, and source validity;
\item
  a finite audit record family \(A_{\mathcal T}\) supplying provenance, replay, and threshold-check entries used by the package's audit functional \(\mathcal A\);
\item
  a finite defect record family \(\delta_{\mathcal T}\) tracking package-level closure deficits, gluing obstructions, idempotence residuals, and other defects;
\item
  finite promotion data when the package participates in a promotion bridge: an old interface, a target package, a target interface, and the bridge gate record;
\item
  finite model-realization data when the package realizes a fragment of \(\mathbf{BirdInt}^{\mathrm{aud}}_{\mathrm{fin}}\) inside a named model host;
\item
  a finite family \(W_{\mathcal T}\) of witness records attached to the package, each with signature in some \(\mathsf{Wit}(P_i)\), visible, sourced and audited like the other attached records;
\item
  a finite nonclaim record \(\mathcal N_{\mathcal T}\) recording what claims about the package are explicitly not asserted.
\end{itemize}

The visibility, suppression, source, audit, defect, promotion, model-realization, and nonclaim records are extra structure attached around the inherited five-tuple. They never silently replace any inherited component, and a Foundations III theorem about a package always restates which inherited fields it uses.

The canonical lens-induced expressive content is the inherited family

\[
\Sigma_f \;=\; \{\,f^{-1}(B)\;:\;B\subseteq f(Z)\,\},
\]

and the canonical lens completion is the inherited endomap

\[
E_f(S)\;=\;f^{-1}(f(S)).
\]

Where Foundations III uses a different completion, the choice is named explicitly and the relationship to \(E_f\) is recorded.

The word ``closure'' is used only with its declared sense. The completion endomap \(E\) is a packaging/completion operation on a declared content space; an order-closure operator is used only on a poset host satisfying the relevant closure axioms; and macro closure refers to closed induced dynamics such as a descended update law. None of these follows from the others without an explicit theorem or bridge.

\subsubsection{Package Morphisms}\label{sn:3.3.2}

The calculus uses four kinds of finite morphism between theory packages, all of which are records on the base host (or on an extension host when a package lives there).

\begin{itemize}
\tightlist
\item
  \textbf{Package map.} A package map \(\Phi:\mathcal T\to\mathcal T'\) is a finite typed record consisting of a finite carrier map \(\phi_Z:Z\to Z'\) and the induced data on lenses, expressive content, completion endomaps, audit functionals, and Foundations III attached records. Its recorded lens data must satisfy a finite compatibility equality in the direction of the map: either \(f'\circ\phi_Z=\psi\circ f\) for a recorded \(\psi:f(Z)\to X'\), or \(f=r\circ f'\circ\phi_Z\) for a recorded \(r:f'(Z')\to X\). Admissibility also requires a visible image for each visible source field and extension of the source and audit records, including \(A_{\mathcal T}\subseteq A_{\mathcal T'}\). A refinement requires the second equality, as specified below.
\item
  \textbf{Refinement.} A package refinement is a package map \(\Phi:\mathcal T\to\mathcal T'\), with lenses \(f:Z\to X\) and \(f':Z'\to X'\), for which there is a map \(r:f'(Z')\to X\) with \(f=r\circ f'\circ\phi_Z\); then \(\Sigma_f\subseteq\{\phi_Z^{-1}(S'):S'\in\Sigma_{f'}\}\), so the codomain expressive content extends the domain expressive content, and the refinement is recorded in the level/lens/interface field of any judgment that uses it. A refinement is not assumed to improve any closure deficit; refinement-implies-improvement is excluded by the negative scope of Subsubsection \hyperref[sn:2.4.3]{2.4.3}.
\item
  \textbf{Quotient.} A package quotient is a package map \(\Phi\) whose carrier component \(\phi_Z\) is a finite surjection \(q:Z\to Z/\!\sim\) for a finite equivalence relation \(\sim\) on \(Z\), together with the matching quotient data on lens, expressive content, and completion. Quotients are the finite maps used by the descent theorems of Section \hyperref[sn:7]{7}.
\item
  \textbf{Completion-induced map.} A completion-induced map is a finite package map or package record induced by applying a named completion endomap \(E:\mathcal V\to\mathcal V\) on its declared content space. When the host supplies inclusions or fixed-point records, those maps are recorded explicitly; they are not inferred from the symbol \(E\) alone. Repeated application of an idempotent completion is not treated as strict package growth unless a new interface, bridge, refinement, or nonfactorization witness is recorded.
\end{itemize}

Package maps compose as records, but composition is not assumed to preserve every Foundations III attached record without an explicit compatibility condition: in particular, audit and source records compose only when the audit functionals of the codomain extend those of the domain, and a composition that violates an inherited audit condition is a defect, not a new package map.

\subsubsection{Package Admissibility}\label{sn:3.3.3}

A theory package \(\mathcal T\) is admissible in \(\mathbf{BirdInt}^{\mathrm{aud}}_{\mathrm{fin}}\), written

\[
\mathsf{AdmPkg}(\mathcal T),
\]

when all of the following hold.

\begin{itemize}
\tightlist
\item
  The inherited five-tuple \((Z,f,\Sigma_f,E,\mathcal A)\) is present and finite. The carrier \(Z\) is a finite set, the lens \(f\) is a finite map from \(Z\) to a finite codomain, the expressive content \(\Sigma_f\) is a finite family of subsets of \(Z\), the completion endomap \(E\) is a finite endomap on the relevant content space, and the audit functional \(\mathcal A\) is a finite map from the entries of the audit record family \(A_{\mathcal T}\) to \(\{\texttt{pass},\texttt{fail}\}\).
\item
  The host tag \(h_{\mathcal T}\) lies in \(\mathsf{Host}\), and every component of the package lives in or under that host.
\item
  The level tag \(\ell_{\mathcal T}\) lies in \(\mathsf{Lev}\). Cross-level use of the package requires a level bridge.
\item
  The visibility profile \(V_{\mathcal T}\) partitions the package's fields into visible, suppressed, outside-scope, and unknown classes for each instrument that may consult \(\mathcal T\), and the partition is finite; for each instrument \(I\) for which \(V_{\mathcal T}\) records a partition, that partition is the restriction of \(\mathsf{vis}_I\) to the package's fields.
\item
  The source record \(\mathrm{src}_{\mathcal T}\) is present, finite, well-formed, and has resolved references; whether it satisfies a consulting instrument's source policy is determined by the applicable judgment classifier.
\item
  The audit record family \(A_{\mathcal T}\) is finite and consistent with \(\mathcal A\), that is, \(\mathcal A(e)=\texttt{pass}\) for every entry \(e\) of \(A_{\mathcal T}\). The audit record family is distinct from the audit functional: \(A_{\mathcal T}\) supplies the finite checked entries that \(\mathcal A\) certifies.
\item
  The defect record family \(\delta_{\mathcal T}\) is finite, with each entry conforming to the defect record schema of Section \hyperref[sn:5]{5}.
\item
  Any attached promotion data, model-realization data, or level/lens/interface data is finite, hosted, assigned an explicit visibility status under the relevant instrument, and accompanied by its own source, audit, and nonclaim records.
\item
  The attached witness family \(W_{\mathcal T}\) is finite, and each of its records has signature in some \(\mathsf{Wit}(P_i)\) and its own visibility, source and audit records.
\item
  The nonclaim record \(\mathcal N_{\mathcal T}\) is present and finite.
\end{itemize}

A package that fails any of these conditions is not admissible. Subsequent sections only consider admissible packages, and a theorem statement that uses a package implicitly carries \(\mathsf{AdmPkg}(\mathcal T)\) as a hypothesis. The finite shape of these checks is part of the input to the later well-formedness results; this subsection only records the package schema and admissibility conditions.

\subsection{Instruments (Reference)}\label{sn:3.4}

\emph{Used by:} Theorems \hyperref[cl:thm:1]{1}--\hyperref[cl:thm:5]{5} (the running-example instrument of Subsection \hyperref[sn:3.8]{3.8} and its variants) and Theorems \hyperref[cl:thm:22]{22}--\hyperref[cl:thm:25]{25} use instruments in their statements; Theorem \hyperref[cl:thm:3]{3} decides well-formedness relative to one.

\subsubsection{Instrument Record}\label{sn:3.4.1}

An instrument is a finite typed record

\[
\begin{aligned}
I\;=\;(\,&\mathrm{id}_I,\;\mathsf{Scope}_I,\;\mathsf{ClaimTypes}_I,\;\mathsf{Visible}_I,\;\mathsf{Suppressed}_I,\\
&\mathsf{CheckRules}_I,\;\mathsf{EvidencePolicy}_I,\;\mathsf{ThresholdPolicy}_I,\;\mathsf{AuditPolicy}_I,\\
&\mathsf{Level}_I,\;\mathsf{BridgePolicy}_I,\;\mathsf{SourcePolicy}_I,\;\mathcal N_I\,).
\end{aligned}
\]

The components have the following roles, all finite:

\begin{itemize}
\tightlist
\item
  \(\mathrm{id}_I\): a finite identifier under which the instrument is registered. Two distinct instruments have distinct identifiers, and an instrument identifier is part of the data of every claim accepted under that instrument.
\item
  \(\mathsf{Scope}_I\): the finite scope record of the instrument. It consists of finite sets \(\mathsf{Hosts}_I\subseteq\mathsf{Host}\) and \(\mathsf{Levels}_I\subseteq\mathsf{Lev}\) and the in-scope content, also written \(\mathsf{Scope}_I\subseteq\mathsf{Cont}\), with \(\mathsf{Outside}_I=\mathsf{Cont}\setminus\mathsf{Scope}_I\). A host tag \(h\) (a level tag \(\ell\)) lies outside \(\mathsf{Scope}_I\) when \(h\notin\mathsf{Hosts}_I\) (\(\ell\notin\mathsf{Levels}_I\)). A claim whose use-set is not contained in \(\mathsf{Scope}_I\) has status \(\texttt{outside\_scope}\) under \(I\).
\item
  \(\mathsf{ClaimTypes}_I\): the finite family of claim types that \(I\) is configured to evaluate, including role-channel, directed-cell, pair-observable, promotion, dependency, and instrument-indexed claim types as appropriate.
\item
  \(\mathsf{Visible}_I,\;\mathsf{Suppressed}_I\): two disjoint subsets of \(\mathsf{Scope}_I\): the content \(I\) may use directly, and the in-scope content \(I\) may not use directly. The complement of \(\mathsf{Visible}_I\cup\mathsf{Suppressed}_I\) within \(\mathsf{Scope}_I\) is the unknown class; content outside \(\mathsf{Scope}_I\) is in the outside-scope class. The full four-way partition is given in Subsection \hyperref[sn:3.5]{3.5}.
\item
  \(\mathsf{CheckRules}_I\): the finite family of check rules used to classify a claim under \(I\). Check rules are typed finite records: each rule states which fields of a claim it consults, what comparison or test it performs, and what classifier verdict it contributes. This component also carries a finite mode map \(\mathsf{ModePolicy}_I\) on judgment families and the finite host/profile compatibility relation \(\mathsf{Compat}_I\); both are declarations of \(I\).
\item
  \(\mathsf{EvidencePolicy}_I\): the finite policy for what counts as admissible evidence under \(I\), distinguishing direct evidence, bridge-relayed evidence, and inadmissible evidence; it admits a source-type tag \(t\) of Subsubsection \hyperref[sn:4.9.1]{4.9.1} when it classifies source records of tag \(t\) as direct or bridge-relayed evidence.
\item
  \(\mathsf{ThresholdPolicy}_I\): the finite family of thresholds used by \(I\), including exact, approximate, tolerance, rank/dimension, visibility, and source thresholds, together with a finite set \(\mathsf{ProvisionalTypes}_I\) of claim and square types that may receive \(\texttt{provisional}\). The threshold record schema is fixed in Subsection \hyperref[sn:5.2]{5.2}.
\item
  \(\mathsf{AuditPolicy}_I\): the finite policy for which audit records \(I\) requires, in what order, and against which sources. The policy is checked by the gate semantics of Section \hyperref[sn:10]{10} and the rotating-audit semantics of Section \hyperref[sn:11]{11}. Passing an audit policy contributes to an instrument-indexed status; it is not instrument-free truth.
\item
  \(\mathsf{Level}_I\in\mathsf{Lev}\): the level annotation of \(I\). An instrument operates at a single level; cross-level operation requires an explicit level bridge inside the instrument's claim record.
\item
  \(\mathsf{BridgePolicy}_I\) (including a finite map to \(\mathsf{Strength}_I\) that is total on every admitted (bridge type, regime) pair, and a finite map \(\mathsf{ReqStrength}_I:\mathsf{ClaimTypes}_I\to\mathsf{Strength}_I\)): the finite policy for which bridges \(I\) admits, including visibility bridges, level bridges (among them the \(\texttt{reflection}\) level bridge of rotating audits), instrument-transfer bridges, drive bridges, source-upgrade bridges, gluing bridges, and model-realization bridges. The policy fixes which suppressed content can be brought into the use-set of an accepted claim. A generic audit bridge is not a drive certificate unless the bridge records the required \(\mathrm{P6}_{\mathrm{drive}}\) data. An instrument classifying decorated squares also declares a separate total map from the seven square types of Subsubsection \hyperref[sn:14.1.2]{14.1.2} to \(\mathsf{Strength}_I\), likewise written \(\mathsf{ReqStrength}_I\).
\item
  \(\mathsf{SourcePolicy}_I\): a finite relation between source-type tags, admission modes (\(\texttt{full}\), \(\texttt{fallback}\), \(\texttt{provisional}\)) and profiles, with a default entry; with the source slot it determines \(\delta_6^{\mathrm{source}}\) by the rules of Subsubsection \hyperref[sn:5.3.6]{5.3.6}, where the absence, explicit unknown/contradictory tags, dirtiness and proxy-bridge tests precede the source-policy admission tests; an unknown value can also result from failed full-mode admission.
\item
  \(\mathcal N_I\): the finite nonclaim record of \(I\). The nonclaims of \(I\) are part of the instrument and propagate to every claim accepted under \(I\); an accepted claim under \(I\) is interpreted with the nonclaims of \(I\) attached.
\end{itemize}

A field \(\mathsf F_I\) of \(I\) is also written \(\mathsf F(I)\) when \(I\) carries an index; thus \(\mathsf{Level}(I_0)=\mathsf{Level}_{I_0}\). Equality of instruments is a finite record comparison across all fields. An instrument is admissible when every component is finite and well typed: its level lies in \(\mathsf{Lev}\) and the hosts named in \(\mathsf{Scope}_I\) lie in \(\mathsf{Host}\), and its visibility, threshold, gate and audit components are records of the schemas of Subsection \hyperref[sn:3.5]{3.5}, Subsection \hyperref[sn:5.2]{5.2} and Sections \hyperref[sn:10]{10} and \hyperref[sn:11]{11}. These are finite typing checks. Instrument admissibility also requires every rule of \(\mathsf{CheckRules}_I\) to be a program of a declared rule language \(\mathcal R\) whose programs are total by construction (finite table lookups, the declared operations of the field types, equality and order comparisons on declared field types, and iteration bounded by the size of finite fields), so that each rule is a terminating decision procedure on its declared field types, and \(\mathcal R\) is closed under sequential composition and finite, acyclic calls to named \(\mathcal R\)-programs; membership of a rule in \(\mathcal R\) is a syntactic check. It also requires every audit record judged under \(I\), including the audit records \(A_L\) of visibility bridges, to draw its \(\mathsf{CheckRules}\) field from \(\mathsf{CheckRules}_I\). If \(\texttt{top\_down\_channel}\in\mathsf{ClaimTypes}_I\), admissibility also requires \(\mathsf{CheckRules}_I\) to contain the rule that assigns the rejection verdict to \(\mathsf{TopDownChannel}(R)\) whenever some gate of \(\mathsf{TDGate}(R)\) (Subsection \hyperref[sn:4.7]{4.7}) has a status other than \(\texttt{pass}\).

\subsubsection{Instrument-Relative Judgments}\label{sn:3.4.2}

Every nontrivial claim in the calculus is interpreted under an instrument. The general form of an instrument-indexed claim is

\[
\Gamma\,;\;\mathcal T\,;\;I\;\vdash\;\varphi:\chi,
\]

where \(\Gamma\) is a finite typed context of attached records (witnesses, updates, profiles, level/lens/interface records, visibility records, threshold records, audit records, defect records, and nonclaim records), \(\mathcal T\) is an admissible theory package, \(I\) is an admissible instrument, \(\varphi\) is a claim with a finite use-set \(\mathsf{Use}(\varphi)\subseteq\mathsf{Cont}\), and \(\chi\) is the claim status assigned to \(\varphi\) under \(\Gamma,\mathcal T,I\).

The middle slot may also hold a host \(\mathsf H\in\mathsf{Host}\) (Subsubsection \hyperref[sn:2.1.2]{2.1.2}) in place of a package, when the claim concerns the host as a whole rather than one theory package. Appendix \hyperref[sn:I]{I} records every numbered result in this form, and Theorem \hyperref[cl:thm:25]{25} states its rotating-audit report in it.

The shorthand

\[
I\;\vdash\;\varphi(\mathcal T)
\]

abbreviates

\[
\Gamma\,;\;\mathcal T\,;\;I\;\vdash\;\varphi:\texttt{accepted}.
\]

A claim that holds in this sense is accepted by \(I\) under the package \(\mathcal T\) and the surrounding context \(\Gamma\). The shorthand never abbreviates instrument-free truth, and a statement of the form \(I\vdash\varphi(\mathcal T)\) does not transfer to a different instrument \(I'\) except through an explicit instrument-transfer bridge admitted by both \(I\) and \(I'\).

The full classifier under which \(\chi\) is determined consults the visibility partition \(\mathsf{Visible}_I/\mathsf{Suppressed}_I\), the bridge policy \(\mathsf{BridgePolicy}_I\), the threshold policy \(\mathsf{ThresholdPolicy}_I\), the source policy \(\mathsf{SourcePolicy}_I\), the audit policy \(\mathsf{AuditPolicy}_I\), and the nonclaim record \(\mathcal N_I\). The classifier is defined by the claim-classifier table of Subsubsection \hyperref[sn:5.6.6]{5.6.6}; the no-overreading and same-level self-audit theorems that depend on it are stated and proved at theorem-grade in Subsections \hyperref[sn:11.4]{11.4} and \hyperref[sn:11.5]{11.5}. This subsection only fixes the form of the judgment and the meaning of the shorthand. No theorem about acceptance or overreading is asserted here.

A claim that names a context and package but does not name an instrument is not a judgment in the calculus. A claim that names an instrument but whose context, package, or use-set is incomplete fails the claim well-formedness conditions of Subsubsection \hyperref[sn:11.1.2]{11.1.2} and so receives no status from the claim classifier of Subsubsection \hyperref[sn:5.6.5]{5.6.5}.

\subsection{Visibility and Suppression (Reference)}\label{sn:3.5}

\emph{Used by:} Theorems \hyperref[cl:thm:1]{1}, \hyperref[cl:thm:2]{2} and \hyperref[cl:thm:4]{4} (the suppressed record \(x\) and the bridge \(L_{\mathrm{drive}}\)) and Theorems \hyperref[cl:thm:22]{22} and \hyperref[cl:thm:23]{23} use visibility records.

Each instrument \(I\) partitions the content universe according to what \(I\) adjudicates and how it adjudicates it. The partition has four classes; this subsection fixes the partition map, the two named classes (visible and suppressed), and the schema of admissible visibility bridges. No theorem about acceptance or overreading is asserted here; the formal no-overreading theorem and the bridge-aware acceptance rule are stated in Section \hyperref[sn:11]{11}.

Each instrument \(I\) carries a finite visibility map

\[
\mathsf{vis}_I:\;\mathsf{Cont}\;\to\;\{\,\texttt{visible},\;\texttt{suppressed},\;\texttt{outside},\;\texttt{unknown}\,\},
\]

inducing a finite four-way partition of the content universe of Subsection \hyperref[sn:3.2]{3.2}:

\[
\mathsf{Cont}\;=\;\mathsf{Visible}_I\;\sqcup\;\mathsf{Suppressed}_I\;\sqcup\;\mathsf{Outside}_I\;\sqcup\;\mathsf{Unknown}_I.
\]

The four classes correspond, respectively, to content \(I\) may use directly, content within \(I\)'s scope that \(I\) may not use directly, content outside \(I\)'s declared scope, and content within \(I\)'s scope whose visibility verdict has not yet been resolved. The classes are part of \(I\)'s record (Subsubsection \hyperref[sn:3.4.1]{3.4.1}); changing the partition is changing the instrument.

\subsubsection{Visible Content}\label{sn:3.5.1}

The visible class under \(I\) is

\[
\mathsf{Visible}_I\;=\;\{\,x\in\mathsf{Cont}\;:\;\mathsf{vis}_I(x)=\texttt{visible}\,\}.
\]

A claim whose use-set \(\mathsf{Use}(\varphi)\) lies entirely in \(\mathsf{Visible}_I\) is a visible claim under \(I\); its threshold, audit and evidence-source records are then visible, since they lie in \(\mathsf{Use}(\varphi)\) (Subsubsection \hyperref[sn:3.2.1]{3.2.1}). The classifier of Section \hyperref[sn:11]{11} may consult visible content directly. Visibility under \(I\) is not a property of \(x\) alone: the same record can be visible under one instrument and suppressed, outside, or unknown under another.

\subsubsection{Suppressed Content}\label{sn:3.5.2}

The suppressed class under \(I\) is

\[
\mathsf{Suppressed}_I\;=\;\{\,x\in\mathsf{Cont}\;:\;\mathsf{vis}_I(x)=\texttt{suppressed}\,\}.
\]

Suppression is a record-level relation between \(x\) and \(I\), not a truth verdict on \(x\):

\[
x\in\mathsf{Suppressed}_I\;\;\not\Rightarrow\;\;x\;\text{is false}.
\]

Suppressed content may be exactly correct under another instrument or in an external description. The relation \(x\in\mathsf{Suppressed}_I\) asserts only that \(I\) does not have admissible visibility into \(x\); under \(I\), a claim that depends on \(x\) without an admissible bridge over-reads \(I\) (in the informal sense fixed in Subsubsection \hyperref[sn:2.3.3]{2.3.3}).

The map \(\mathsf{vis}_I\) is determined by the fields of \(I\): \(\mathsf{vis}_I(x)\) is visible iff \(x\in\mathsf{Visible}_I\), suppressed iff \(x\in\mathsf{Suppressed}_I\), outside iff \(x\notin\mathsf{Scope}_I\), and unknown otherwise. The remaining two classes are recorded for completeness:

\[
\mathsf{Outside}_I\;=\;\{\,x\in\mathsf{Cont}\;:\;\mathsf{vis}_I(x)=\texttt{outside}\,\},\qquad \mathsf{Unknown}_I\;=\;\{\,x\in\mathsf{Cont}\;:\;\mathsf{vis}_I(x)=\texttt{unknown}\,\}.
\]

A claim whose use-set meets \(\mathsf{Outside}_I\) has status \(\texttt{outside\_scope}\) under \(I\); a claim whose use-set meets \(\mathsf{Unknown}_I\) requires the later classifier to consult the corresponding unknown-content rule. No acceptance status for unknown content is asserted in this subsection.

\subsubsection{Visibility Bridges}\label{sn:3.5.3}

Suppressed content can become admissibly usable by a claim only via a recorded visibility bridge. A visibility bridge under \(I\) is a finite typed record

\[
L\;=\;(\,\mathrm{id}_L,\;c,\;c',\;\mathsf{BridgeType}_L,\;\mathsf{Guarantee}_L,\;\mathsf{Strength}_L,\;A_L,\;V_L,\;\mathcal N_L\,),
\]

with \(c\in\mathsf{Suppressed}_I\), \(c'\in\mathsf{Visible}_I\), \(\mathsf{BridgeType}_L\) drawn from the bridge inventory of \(\mathsf{BridgePolicy}_I\) (Subsubsection \hyperref[sn:3.4.1]{3.4.1}), and \(\mathsf{Guarantee}_L\), \(\mathsf{Strength}_L\), \(A_L\), \(V_L\), \(\mathcal N_L\) the finite guarantee, strength, audit, visibility, and nonclaim records that the bridge carries; the audit record \(A_L\) has \(\mathsf{ClaimRef}=L\), so \(L\in\mathsf{Use}(L)\). Admissibility of a bridge for a claim or cell, \(\mathsf{AdmVisBridge}(L,I,X)\), is defined in Subsubsection \hyperref[sn:5.5.2]{5.5.2}.

\(\mathcal L\) is a finite set of registered, audited visibility-bridge candidates under \(I\). Registration checks the bridge's type, endpoints and audit record but does not assert that its strength suffices for every claim. A member \(L\in\mathcal L\) is admissible \emph{for \(\varphi\)} only when \(\mathsf{AdmVisBridge}(L,I,\varphi)\) holds, including the strength comparison. The schema-level use of \(\mathcal L\) is to record which suppressed content has been brought into admissible reach for which claims; the formal acceptance rule, including the constraint that every suppressed dependency of an accepted claim is bridged by some \(L\in\mathcal L\), is stated and proved in Section \hyperref[sn:11]{11}. This subsection only fixes the schema; no acceptance or no-overreading theorem is asserted here.

The witness and update records introduced in Subsection \hyperref[sn:3.7]{3.7} are content records in this sense: when a later directed-cell or pair-observable judgment uses a witness \(W_i\) or update \(U_i\), that record has a visibility status under the named instrument and may require a visibility bridge if it is suppressed.

\subsection{Profiles (Reference)}\label{sn:3.6}

\emph{Used by:} The directed-cell records in the statements of Theorems \hyperref[cl:thm:1]{1}--\hyperref[cl:thm:5]{5} carry a profile \(\lambda\).

A profile is the finite typed record attached to a directed-cell, pair-observable, promotion-bridge, rotating-audit-bridge, instrument-transfer or top-down channel record that fixes the per-judgment scale, regime, bridge type, instrument mode, and level data. Claim records carry only a level profile (Subsubsection \hyperref[sn:11.1.2]{11.1.2}), and role-channel and dependency records carry none. Profiles are attached to the six record families listed above; their finite admissibility schema and their dependence on judgment status are recorded here. Theorem \hyperref[cl:thm:1]{1} of Section \hyperref[sn:6]{6} uses the profile schema.

\subsubsection{Profile Grammar}\label{sn:3.6.1}

A profile is a finite typed record

\[
\lambda\;=\;(\,\ell_i,\;\ell_j,\;\mathsf{Scale}_\lambda,\;\mathsf{Regime}_\lambda,\;\mathsf{BridgeType}_\lambda,\;\mathsf{InstrumentMode}_\lambda,\;\mathsf{Locality}_\lambda\,)
\]

with the following finite components.

\begin{itemize}
\tightlist
\item
  \(\ell_i,\ell_j\in\mathsf{Lev}\) are the actor-side and informant-side level tags. For a directed cell \(P_i\leftarrow P_j\), \(\ell_i\) is the level on which the actor update \(U_i\) lives and \(\ell_j\) is the level on which the informant witness \(W_j\) lives. For a judgment that is not directed (e.g., a pair observable or a promotion bridge), the two tags coincide or are determined by the judgment family's convention.
\item
  \(\mathsf{Scale}_\lambda\) is a finite scale label, distinguishing the scale of the carriers, partitions, kernels, or graphs the judgment depends on (for example, microscale or macroscale on a stochastic host, or fine or coarse on a finite carrier).
\item
  \(\mathsf{Regime}_\lambda\) is a finite regime label, distinguishing dynamical, structural, and audit regimes; specific regimes used by the paper are named where the relevant judgment is introduced.
\item
  \(\mathsf{BridgeType}_\lambda\) is drawn from the finite bridge inventory of \(\mathsf{BridgePolicy}_I\) (Subsubsection \hyperref[sn:3.4.1]{3.4.1}): visibility, level (including the \(\texttt{reflection}\) level bridge of rotating audits), instrument-transfer, drive, source-upgrade, gluing, or model-realization. The profile records which bridge type the surrounding judgment uses.
\item
  \(\mathsf{InstrumentMode}_\lambda\) is a finite instrument-mode label, distinguishing the surrounding instrument's default, audit-only, transfer, and rotating-audit modes.
\item
  \(\mathsf{Locality}_\lambda\) is a finite locality label, distinguishing local, semi-global, and global use of the judgment data on a cover or staging.
\end{itemize}

Specific named profiles used elsewhere in the paper include \(\lambda_{\mathrm{audit}}\) and \(\lambda_{\mathrm{drive}}\) (Subsection \hyperref[sn:3.8]{3.8}, Subsubsection \hyperref[sn:6.1.1]{6.1.1}), \(\mu_{\mathrm{drive}}\) (Subsection \hyperref[sn:9.5]{9.5}), \(\lambda_{\mathrm{prom}}\) (Subsubsection \hyperref[sn:10.6.1]{10.6.1}), \(\lambda_{\mathrm{rot}}\) (Subsubsection \hyperref[sn:11.7.2]{11.7.2}), \(\lambda_{\mathrm{PICA}}\) and \(\mu_{\mathrm{PICA}}\) (Subsubsection \hyperref[sn:13.2.5]{13.2.5}), \(\lambda_{\mathrm{Cantor}}\) (Model Theorem \hyperref[cl:thm:34]{34}), \(\lambda_{\mathrm{obj}}\) and \(\lambda_{\mathrm{refresh}}\) (Subsection \hyperref[sn:11.6]{11.6}) and \(\lambda_{\mathrm{transfer}}\) (Subsubsection \hyperref[sn:11.9.1]{11.9.1}). Each is a finite record with the components above. The values of \(\lambda_{\mathrm{audit}}\), \(\lambda_{\mathrm{drive}}\), \(\mu_{\mathrm{drive}}\), \(\lambda_{\mathrm{prom}}\), \(\lambda_{\mathrm{PICA}}\) and \(\mu_{\mathrm{PICA}}\) are fixed where the corresponding judgment is introduced and are reused thereafter without redefinition; \(\lambda_{\mathrm{rot}}\) is fixed up to its two level tags, the levels of the bridge's instruments (Subsubsection \hyperref[sn:11.7.2]{11.7.2}); Model Theorem \hyperref[cl:thm:34]{34} constrains \(\lambda_{\mathrm{Cantor}}\) only through \(\mathsf{AdmProfile}\) and its level tags, its Moreover clause giving one admissible value; of \(\lambda_{\mathrm{obj}}\) and \(\lambda_{\mathrm{refresh}}\) only the audit-only instrument mode is fixed (Subsection \hyperref[sn:11.6]{11.6}), and of \(\lambda_{\mathrm{transfer}}\) only the instrument-transfer bridge type (Subsubsection \hyperref[sn:11.9.1]{11.9.1}), the other components being any values for which \(\mathsf{AdmProfile}\) holds.

\subsubsection{Profile Admissibility}\label{sn:3.6.2}

A profile \(\lambda\) is admissible under instrument \(I\) and theory package \(\mathcal T\), written

\[
\mathsf{AdmProfile}(\lambda,\,I,\,\mathcal T),
\]

when its components are finite and the following checks hold:

\begin{itemize}
\tightlist
\item
  \(\ell_i,\ell_j\in\mathsf{Lev}\) (a profile carries no level-bridge record; when \(\ell_i\neq\ell_j\), the level bridge of one of the four types from Subsubsection \hyperref[sn:2.2.3]{2.2.3} (translation, lift, forgetting, reflection) is required of the judgment's level/lens/interface record and checked with it, step 5 of Subsubsection \hyperref[sn:6.3.1]{6.3.1}, not by \(\mathsf{AdmProfile}\));
\item
  \(\mathsf{BridgeType}_\lambda\) admitted by \(\mathsf{BridgePolicy}_I\);
\item
  \(\mathsf{InstrumentMode}_\lambda=\mathsf{ModePolicy}_I(\text{surrounding judgment family})\);
\item
  \((\mathsf{Scale}_\lambda,\mathsf{Regime}_\lambda,\mathsf{Locality}_\lambda,h_{\mathcal T})\in\mathsf{Compat}_I\), where \(\mathsf{Compat}_I\) is a finite relation declared by the instrument, listing the scale, regime and locality values it accepts on each host.
\end{itemize}

Each check is a finite typing, equality, bridge-reference, or membership check on records carried by the instance, so \(\mathsf{AdmProfile}\) is decidable.

A profile failing a typing or declared-mode check is inadmissible, so the candidate judgment is malformed. A structurally admissible profile may accompany content outside \(\mathsf{Scope}_I\); the classifier then records \(\texttt{outside\_scope}\) for a directed cell, promotion bridge or claim, and \(\texttt{blocked}\) for a pair record, whose status family has no such value (Subsubsection \hyperref[sn:5.5.3]{5.5.3}). Profile admissibility is a finite check on finite records; like package admissibility (Subsubsection \hyperref[sn:3.3.3]{3.3.3}) and instrument admissibility (Subsubsection \hyperref[sn:3.4.1]{3.4.1}), it feeds the well-formedness classifier of Section \hyperref[sn:6]{6} rather than producing its own theorem here.

Profile mutation is not allowed silently. A judgment whose statement carries one profile and whose proof or downstream use depends on a different profile requires an explicit profile-update record, itself a typed entry in \(\mathsf{Cont}\) subject to the visibility, source, and audit discipline of Subsections \hyperref[sn:3.5]{3.5} and \hyperref[sn:4.9]{4.9} and Section \hyperref[sn:10]{10}.

\subsubsection{Profile Dependence}\label{sn:3.6.3}

The status assigned to a directed-cell, pair-observable or promotion record depends on the profile attached to the judgment. The same primitive labels and the same witness, update, visibility, threshold and audit data, attached to a different profile, can produce a different status. The regime variant below keeps the primitive labels, witness, update, visibility, threshold and audit data fixed and recomputes the required-bridge defect; in general, a profile change requires recomputing every affected record field, defect and audit check to satisfy well-formedness. Concretely:

\begin{itemize}
\tightlist
\item
  the running-example cell of Subsection \hyperref[sn:3.8]{3.8} is \(\texttt{action}\) under \(\lambda_{\mathrm{audit}}\); the cell \(C_{\mathrm{reg}}\) of Theorem \hyperref[cl:thm:1]{1} differs from it only in its regime, set to \(\texttt{drive-certification}\) (admitted by \(\mathsf{Compat}_I\)), in the target of its audit and in the defect record recomputed from these. Its profile gives \(\mathsf{ReqBridge}(C)=\{\texttt{drive}\}\); since \(\mathsf{AddAuditRecord}\) names no drive bridge, \(\delta_{\mathrm{reqbridge}}=\{\texttt{drive}\}\) and the first \(\texttt{blocked}\) row of the directed-cell table of Subsubsection \hyperref[sn:5.6.6]{5.6.6} applies. A cell whose profile has a level tag outside \(\mathsf{Levels}_I\) receives \(\texttt{outside\_scope}\), since step 5 of the well-formedness procedure requires \(L\) to carry the level tags of the profile. A cell whose profile fails \(\mathsf{AdmProfile}\) is malformed. The rows \(\texttt{below\_threshold}\) and \(\texttt{collapsed}\) read \(\Theta\), \(L\) and \(W\), not \(\lambda\);
\item
  a pair-observable judgment on the same primitive pair, with the same source data, may be classified \(\texttt{real}\) under one profile and \(\texttt{blocked}\) under another;
\item
  a promotion bridge with fixed packages and object maps may be \(\texttt{strict}\) under one admissible profile and not accepted under another. Take a suppressed record \(x\in\mathsf{Use}(B)\) whose only registered bridge is \(L\) of strength \(\texttt{weak}\). Assume every other use is visible, every bridge-admissibility condition except strength passes, and both profiles meet \(\mathsf{AdmProfile}\). Recompute the audits, defects and gate records under each profile. If \(\mathsf{BridgePolicy}_I\) assigns \(\texttt{weak}\) to the (bridge type, regime) pair of the first profile, so that \(\mathsf{RequiredStrength}(B)=\texttt{weak}\) (Subsubsection \hyperref[sn:5.5.2]{5.5.2}), the visibility defect is empty; with the other required gates passing, strict object maps and a required strictness verdict, the classifier returns \(\texttt{strict}\). If it assigns \(\texttt{strong}\) to that pair of the second profile, then \((x,L)\in\delta_{\mathrm{weakbridge}}(B)\) and \(G_{\mathrm{vis}}\) fails; for the resulting well-formed, in-scope bridge, the classifier returns \(\texttt{failed\_audit}\) if its audit gate fails, and \(\texttt{failed\_no\_smuggling}\) otherwise. The choice between \(\texttt{strict}\) and \(\texttt{non\_strict}\) is fixed by \(\Delta_{\mathrm{fact}}(\pi_j,\pi_{j+1})\) (Theorems \hyperref[cl:thm:20]{20} and \hyperref[cl:thm:21]{21}).
\end{itemize}

This is a schema-grade observation about which fields the later classifiers must inspect: status is not a function of primitive labels alone, because the classifier's input includes the profile (along with witness, update, threshold, audit, visibility, defect, and source records). Theorem \hyperref[cl:thm:1]{1} of Section \hyperref[sn:6]{6} is the theorem-grade statement that status does not factor through the primitive pair; its two cells also differ in update, visibility and audit records, and the regime variant above isolates the profile.

\subsection{Witness and Update Families}\label{sn:3.7}

\subsubsection{Witness Map}\label{sn:3.7.1}

The witness map of the calculus is the finite typed family

\[
\mathsf{Wit}\colon\mathbb P\;\to\;\mathrm{Type},\qquad P_i\;\mapsto\;\mathsf{Wit}(P_i),
\]

assigning to each primitive label a finite collection of admissible witness types. \(\mathsf{Wit}\) and \(\mathsf{Upd}\) are components of the instance: for each \(P_i\) the instance fixes a finite set of named record signatures, and for any host tag drawn from \(\mathsf{Host}\) the witness types specialize to finite records on that host. The tables of Subsubsections \hyperref[sn:3.7.4]{3.7.4}--\hyperref[sn:3.7.5]{3.7.5} list the kinds of signature this paper's instances use; the proofs use \(\mathsf{RouteMismatchWit}\in\mathsf{Wit}(P_3)\), the source-record signature of Subsubsection \hyperref[sn:4.9.1]{4.9.1} in \(\mathsf{Wit}(P_6)\) (the role evidence of Theorem \hyperref[cl:thm:5]{5}(a), audited by a record \(A_{R_6}\) with \(\mathsf{ClaimRef}=R_6\) and empty audit defect), and \(\mathsf{AddAuditRecord},\mathsf{CertifyDrive}\in\mathsf{Upd}(P_6)\) (Subsubsection \hyperref[sn:3.1.1]{3.1.1}), and the signatures named in Subsection \hyperref[sn:12.5]{12.5} and Subsubsection \hyperref[sn:13.2.5]{13.2.5}. A witness \(W_i\in\mathsf{Wit}(P_i)\) is the informant-side data of a directed cell, or one side of a pair-observable record, in which \(P_i\) supplies the witness family; its host, level, instrument, and profile tags are determined by the surrounding judgment.

\subsubsection{Update Map}\label{sn:3.7.2}

The update map of the calculus is the finite typed family

\[
\mathsf{Upd}\colon\mathbb P\;\to\;\mathrm{Type},\qquad P_i\;\mapsto\;\mathsf{Upd}(P_i),
\]

assigning to each primitive label a finite collection of admissible update types. Each \(\mathsf{Upd}(P_i)\) is a finite family of finite typed records describing the effect that the primitive \(P_i\) produces when it occupies the actor role in a directed cell or promotion bridge. As with witnesses, every update record carries host, level, instrument, and profile tags consistent with the surrounding judgment. Every update signature of \(\mathsf{Upd}(P_i)\) includes a target field, naming the record \(r\) of \(\mathsf{Cont}\) it modifies, and an effect field recording its action on \(r\) as a finite map \(e_U\) on the type of \(r\), each of which may hold the value none (for a completion, \(e_U=E\) and \(r\) is the recorded input); \(\mathsf{noopU}\) reads only these two fields and is false when either holds none.

\subsubsection{Witness/Update Constraint}\label{sn:3.7.3}

A directed-cell witness slot is either a present \(W_j\in\mathsf{Wit}(P_j)\) or a typed, audited absence marker for \(P_j\); its update slot is either a present \(U_i\in\mathsf{Upd}(P_i)\) or a typed, audited absence marker for \(P_i\). Each marker is a finite record in \(\mathsf{Cont}\) naming its primitive, slot, and passing absence audit. A wrong-typed present record or an unaudited absence marker is malformed. A typed, audited absence passes its slot check; if the remaining well-formedness checks pass, the cell classifier tests the marker. The rule is checked at the well-formedness step of Section \hyperref[sn:6]{6} together with the package, instrument, profile, level, visibility, threshold, audit, and defect conditions of the surrounding judgment.

The notation \(P_i\leftarrow P_j\) for a directed cell is a typed bridge judgment that records this rule. It is not an expression of a binary product on \(\mathbb P\), and a statement that uses the notation but violates the witness/update rule is a malformed cell, not a primitive product.

\subsubsection{Primitive Witness Table}\label{sn:3.7.4}

The primitive witness families used in this paper are listed below. Each entry is the finite witness family of the corresponding primitive: the entry names the kinds of records that count as witnesses, and each named record type is itself a finite typed record on the relevant host.

\begingroup
\setlength{\tabcolsep}{4pt}
\renewcommand{\arraystretch}{1.08}
\begin{center}
\begin{tabularx}{0.96\linewidth}{@{}>{\raggedright\arraybackslash}p{0.11\linewidth}>{\raggedright\arraybackslash}X@{}}
\toprule
Primitive & $\mathsf{Wit}(P_i)$ \\
\midrule
$P_1$ & descent success/failure records, rewrite obstruction records, closure-deficit records \\
$P_2$ & representability and nonrepresentability records, gate decisions, support and capacity records \\
$P_3$ & route-mismatch records, holonomy records, commutator records, critical-pair records, pasting-mismatch records \\
$P_4$ & stage records, refinement records, filtration records, cover records, local/global obstruction records \\
$P_5$ & package map records, quotient records, completion records, idempotence records, fixed-point records \\
$P_6$ & audit records, provenance records, source-of-truth records, threshold-check records, nonclaim records, drive certificates \\
\bottomrule
\end{tabularx}
\end{center}
\endgroup

The drive-certificate entry under \(\mathsf{Wit}(P_6)\) is reserved for \(\mathrm{P6}_{\mathrm{drive}}\) data, which is the drive face of \(P_6\) rather than the generic audit face. A drive certificate is admissible only when the surrounding host supports it, and Section \hyperref[sn:8]{8} lists the finite cycle-supported and gating-suppression cases in which the certificate may or may not be supported. The presence of generic audit records under \(\mathsf{Wit}(P_6)\) does not by itself supply a drive certificate.

\subsubsection{Primitive Update Table}\label{sn:3.7.5}

The primitive update families used in this paper are listed below.

\begingroup
\setlength{\tabcolsep}{4pt}
\renewcommand{\arraystretch}{1.08}
\begin{center}
\begin{tabularx}{0.96\linewidth}{@{}>{\raggedright\arraybackslash}p{0.11\linewidth}>{\raggedright\arraybackslash}X@{}}
\toprule
Primitive & $\mathsf{Upd}(P_i)$ \\
\midrule
$P_1$ & mark descent failure, install descended update, rewrite, repair \\
$P_2$ & mark representable or nonrepresentable, gate support, update admissibility record \\
$P_3$ & record route mismatch, holonomy, commutator, critical pair \\
$P_4$ & refine lens, advance stage, record local/global obstruction \\
$P_5$ & install package map, apply completion endomap, register fixed-point record, saturate \\
$P_6$ & add audit, check provenance, commit source, record nonclaim, certify drive or no-drive \\
\bottomrule
\end{tabularx}
\end{center}
\endgroup

Each entry above is a finite typed update record on the host of the surrounding judgment. The certify-drive-or-no-drive entry under \(\mathsf{Upd}(P_6)\) is the update counterpart to the drive certificate witness type; it is admissible only when the surrounding host and instrument support drive judgments, and the drive or no-drive records of Section \hyperref[sn:8]{8} are recorded through this update family.

The witness and update tables together fix the typed actor/informant data for directed cells and the related witness/update records used by pair observables, promotion bridges, and instrument-indexed claims. Section \hyperref[sn:4]{4} introduces the judgment families that consume this data; Section \hyperref[sn:5]{5} introduces the defect records that the data may produce; and Sections \hyperref[sn:6]{6} through \hyperref[sn:11]{11} use these records in the theorem-grade results.

\subsection{A Running Example}\label{sn:3.8}

The following directed cell serves as a running example; later sections return to it briefly. It is the active cell of the countermodels \(\mathsf{CM}_4\) and \(\mathsf{CM}_7\), with every field written out.

Take the finite carrier \(\mathcal P(U)\) with \(U=\{a,b,c\}\) and the two completion maps of Subsubsection \hyperref[sn:7.4.1]{7.4.1}: \(E_1\) adds \(b\) to every set that contains \(a\), and \(E_2\) adds \(c\) to every set that contains \(b\). Each map is idempotent, and, writing \(E_1E_2\) for \(E_1\circ E_2\) (apply \(E_2\) first), the two do not commute at \(\{a\}\):

\[
E_1E_2(\{a\})\;=\;\{a,b\},\qquad E_2E_1(\{a\})\;=\;\{a,b,c\}.
\]

In plain words, the informant is the route mismatch, shown by the set \(\{a\}\) on which the two orders of completion differ, and the actor is an audit, whose update adds a record of both recomputations to its log. The fields listed below say which witness and update the cell carries, at which level, what is visible, which threshold applies, which audit record checks it, and which defects it has.

Its status is computed by the directed-cell table of Subsubsection \hyperref[sn:5.6.6]{5.6.6}, whose rows are tested from the top and whose last column evaluates this cell; \(\delta_6^{\mathrm{audit}}\), \(\mathsf{implicitW}\) and \(\mathsf{noopU}\) are defined in Subsubsections \hyperref[sn:5.3.6]{5.3.6} and \hyperref[sn:4.3.4]{4.3.4}. In summary: \(W_3\) is the mismatch at \(\{a\}\); \(U_6\) logs both recomputed images; every field is visible, both levels are \(\mathsf{Ver}\), the threshold \(\Theta_{\ge1}\) (the images differ) is met, the audit passes and \(\delta=\varnothing\). Every row above \(\texttt{action}\) fails, so the status is \(\texttt{action}\). The instrument fields below, and the records \(U_6'\), \(x\) and \(L_{\mathrm{drive}}\), do not affect the status of this cell; Theorems \hyperref[cl:thm:1]{1}, \hyperref[cl:thm:2]{2} and \hyperref[cl:thm:5]{5} use them, and the item \emph{Bridge strengths} fixes the comparison that blocks the cell \(C_{\mathrm{blk}}\) of Theorem \hyperref[cl:thm:1]{1}.

The running example is the directed cell \(P_6\leftarrow P_3\) in which an audit (\(P_6\)) records this route mismatch (\(P_3\)):

\[
\mathsf{Cell}_{63}^{\lambda_{\mathrm{audit}}}\;=\;(\,P_6\leftarrow P_3,\;W_3,\;U_6,\;L,\;V,\;\Theta,\;A,\;\delta\,),
\]

with the following fields, defined in Subsubsection \hyperref[sn:4.3.3]{4.3.3}: \(W_3\) is the informant's witness, \(U_6\) the actor's update, \(L\) the level, lens and interface record, \(V\) the visibility record (which also lists any visibility bridges), \(\Theta\) the thresholds, \(A\) the audit record (Subsubsection \hyperref[sn:4.9.2]{4.9.2}), \(\delta\) the defect records (Section \hyperref[sn:5]{5}), and \(\lambda\) the profile (Subsection \hyperref[sn:3.6]{3.6}). \(\mathsf{AuditPolicy}_I\) and \(\mathsf{BridgePolicy}_I\) are policies of the instrument \(I\) (Subsection \hyperref[sn:3.4]{3.4}), and NC-12 is a nonclaim of the register in Subsection \hyperref[sn:15.2]{15.2}.

\begin{itemize}
\item
  \(W_3=\mathsf{RouteMismatchWit}\in\mathsf{Wit}(P_3)\), the route-mismatch record \((E_1,E_2,\{a\})\): the set \(\{a\}\) lies in the route-mismatch defect \(\Delta_3^{\mathrm{comp}}(E_1,E_2)\) of Theorem \hyperref[cl:thm:10]{10}.
\item
  \(U_6=\mathsf{AddAuditRecord}\in\mathsf{Upd}(P_6)\), the update that adds to the audit log \(\Lambda\) a fresh entry \(\varepsilon\notin\Lambda\) recording the recomputation of \(E_1E_2(\{a\})\) and \(E_2E_1(\{a\})\), so \(e_U(\Lambda)=\Lambda\cup\{\varepsilon\}\neq\Lambda\).
\item
  \(L\): both level tags are \(\mathsf{Ver}\), the lens is the identity of \(\mathcal P(U)\), and the interface is the pair \((E_1,E_2)\). The two tags coincide, so no level bridge is needed.
\item
  \(V\): every field of the cell and every record of \(\mathsf{Use}(C)\) is visible under the surrounding instrument \(I\), so no visibility bridge is used.
\item
  \(\Theta\): the threshold \(\Theta_{\ge\varepsilon}\) of Subsubsection \hyperref[sn:5.2.2]{5.2.2} with lower bound \(1\), written \(\Theta_{\ge1}\), on the datum \(|E_1E_2(\{a\})\,\triangle\,E_2E_1(\{a\})|\), a direct check on the witness values; the datum equals \(1\), so the threshold is met.
\item
  \(A\): an audit record whose target (\(\mathsf{ClaimRef}\)) is the cell, whose evidence references (\(\mathsf{EvidenceRefs}\)) are the tables of \(E_1\) and \(E_2\), whose check rules recompute the two composites at \(\{a\}\) under \(\mathsf{AuditPolicy}_I\), whose source references point to the committed record of \(E_1\) and \(E_2\), which satisfies the source policy of \(I\), and whose nonclaim field records that the cell asserts no drive (NC-12 of Subsection \hyperref[sn:15.2]{15.2}). Its check results and its visibility, threshold and circularity checks all pass.
\item
  \(\delta=\varnothing\): its visibility and required-bridge defect sets are empty, it has no package-collapse or \(\Theta\)-named primitive defect entry, \(\mathsf{noopU}\) and \(\mathsf{implicitW}\) are false, and its audit defect has the passing value \(\delta_6^{\mathrm{audit}}=\texttt{audit\_passes}\) (the convention of Subsubsection \hyperref[sn:5.3.6]{5.3.6}). The route-mismatch defect \(\Delta_3^{\mathrm{comp}}(E_1,E_2)\) is the content of the witness \(W_3\), not a defect of the cell.
\item
  \(\lambda_{\mathrm{audit}}\): both level tags \(\mathsf{Ver}\), the fine scale of the finite carrier, the audit regime, the bridge type \(\texttt{visibility}\) (admitted by \(\mathsf{BridgePolicy}_I\); no visibility bridge is needed, since every field is visible), the audit-only instrument mode, and local use of the data. For this instance, \(\mathcal T\) has host \(\mathsf H_{\mathrm{fin}}\), \(\mathsf{ModePolicy}_I(\texttt{directed-cell})=\texttt{audit-only}\), and \(\mathsf{Compat}_I\) contains \((\texttt{fine},\texttt{audit},\texttt{local},\mathsf H_{\mathrm{fin}})\) and \((\texttt{fine},\texttt{drive-certification},\texttt{local},\mathsf H_{\mathrm{fin}})\). The bridge policy admits both the visibility and the drive bridge types.
\item
  \emph{Records for the later variations (used only by the blocked cell of Theorem \hyperref[cl:thm:1]{1}, its reuse in Theorem \hyperref[cl:thm:2]{2}, and kept in \(\mathsf{Cont}\) by the instances of Theorem \hyperref[cl:thm:5]{5}(a) and (c); skip on first reading).} The instance also contains the update \(U_6'=\mathsf{CertifyDrive}\in\mathsf{Upd}(P_6)\), whose entry field names the candidate drive-certificate record \(x\), and a registered drive-type visibility bridge \(L_{\mathrm{drive}}\) from \(x\) to a visible record, of strength \(\texttt{weak}\); the cell above does not use them.
\item
  \emph{Package and instrument.} \(\mathcal T=(\mathcal P(U),\mathrm{id}_{\mathcal P(U)},\mathcal P(\mathcal P(U)),\mathrm{id},\mathcal A)\) on \(\mathsf H_{\mathrm{fin}}\) at level \(\mathsf{Ver}\), with every field visible; source record \((\texttt{e12},\texttt{committed\_state},\text{the tables of }E_1,E_2,\texttt{valid})\); an empty defect family; an empty attached witness family \(W_{\mathcal T}=\varnothing\); one audit entry recomputing both tables; and the nonclaim NC-12. Every record of the instance, in particular \(W_3\), \(U_6\), \(U_6'\), \(x\), \(L_{\mathrm{drive}}\), \(\Theta\), \(A\), \(V\) and \(\texttt{e12}\), carries host tag \(\mathsf H_{\mathrm{fin}}\) and level tag \(\mathsf{Ver}\). The instrument \(I\) has these fields:

  \begin{itemize}
  \tightlist
  \item
    \emph{Scope}: \(\mathsf{Scope}_I=\mathsf{Cont}\), the records of the instance; only this is needed for the status of this cell. \(\mathsf{Cont}\) contains the records of \(\mathcal T\), \(W_3\), \(U_6\), \(U_6'\) with its entry field \(x\), \(L_{\mathrm{drive}}\), every directed-cell record with its fields, and each cell's verdict reference \(\ulcorner\mathsf{Cell},I\urcorner\), so the circular variant of Theorem \hyperref[cl:thm:2]{2} stays in scope. Each variant of Subsections \hyperref[sn:6.1]{6.1}--\hyperref[sn:6.5]{6.5} takes its own \(\mathsf{Cont}\) as \(\mathsf{Scope}_I\). Its instrument keeps the identifier and all other fields of \(I\), lists the records of the variant in \(\mathsf{Scope}_I\) and \(\mathsf{Visible}_I\), and replaces \(I\) in that variant's instance. There are two exceptions: the \(\texttt{outside\_scope}\) variant of Theorem \hyperref[cl:thm:2]{2} uses the instrument \(I'\) described there, and in Theorem \hyperref[cl:thm:5]{5}(c)--(d) \(\mathsf{Scope}_I\) stays the \(\mathsf{Cont}\) of Theorem \hyperref[cl:thm:5]{5}(a), so the added records lie outside the scope of \(I\);
  \item
    \emph{Hosts and levels}: \(\mathsf{Hosts}_I=\{\mathsf H_{\mathrm{fin}}\}\) and \(\mathsf{Levels}_I=\{\mathsf{Ver}\}\);
  \item
    \emph{Claim types, modes and check rules}: claim types \(\texttt{directed-cell}\), \(\texttt{role-channel}\), \(\texttt{drive-check}\) and \(\texttt{pair-observable}\) with \(\mathsf{ModePolicy}_I(\texttt{pair-observable})=\texttt{audit-only}\), with terminating check rules for the role channel, the candidate drive/no-drive entry, the cell recomputation and threshold, the pair source and required-bridge checks, the audit of \(L_{\mathrm{drive}}\), and the replay of every audit that a variant of Subsections \hyperref[sn:6.1]{6.1}--\hyperref[sn:6.5]{6.5} or \(\mathsf{CM}_4\) adds (absence audits of typed absence markers, and audits of added cells, attached witnesses and branch-compatibility records); each audit under \(I\) draws its rules from this family, \(\mathsf H_{\mathrm{fin}}\) admitting that candidate check without asserting a positive drive certificate;
  \item
    \emph{Visibility}: \(\mathsf{Suppressed}_I=\{x\}\), the record named by the entry field of \(U_6'\) above, and \(\mathsf{Visible}_I=\mathsf{Scope}_I\setminus\{x\}\);
  \item
    \emph{Thresholds and level}: the threshold \(\Theta_{\ge1}\) on an exact datum; level \(\mathsf{Ver}\);
  \item
    \emph{Bridge strengths}: bridge strengths \(\mathsf{Strength}_I=\{\texttt{weak}<\texttt{strong}\}\), with \(\mathsf{Strength}(L_{\mathrm{drive}})=\texttt{weak}\) and required strength \(\texttt{strong}\) under drive certification; the strength map of \(\mathsf{BridgePolicy}_I\) sends \((\texttt{drive},\texttt{drive-certification})\) to \(\texttt{strong}\) and the other admitted (bridge type, regime) pairs to \(\texttt{weak}\), so a cell with profile \(\lambda_{\mathrm{drive}}\), which agrees with \(\lambda_{\mathrm{audit}}\) except that its regime is drive-certification and its bridge type is drive (Subsubsection \hyperref[sn:6.1.1]{6.1.1}), has \(\mathsf{RequiredStrength}=\texttt{strong}\); and \(\mathsf{ReqStrength}_I(\texttt{drive-check})=\texttt{strong}\), \(\mathsf{ReqStrength}_I(\chi)=\texttt{weak}\) for every other admitted claim type \(\chi\);
  \item
    \emph{Source policy}: a source policy admitting \(\texttt{committed\_state}\) in the sense of Subsubsection \hyperref[sn:5.3.6]{5.3.6};
  \item
    \emph{Evidence, audit and nonclaim fields}: \(\mathsf{EvidencePolicy}_I\) admits \(\texttt{committed\_state}\) as direct evidence; \(\mathsf{AuditPolicy}_I\) requires replay of every recorded check result under the rules above and lists NC-12 as the only required nonclaim for directed cells; \(\mathsf{ProvisionalTypes}_I=\varnothing\); \(\mathcal N_I=\{\text{NC-12}\}\). The tables of \(E_1\) and \(E_2\) are records of \(\mathsf{Cont}\), named in the trace of \(\texttt{e12}\).
  \end{itemize}

  Each condition of Subsubsections \hyperref[sn:3.3.3]{3.3.3} and \hyperref[sn:3.4.1]{3.4.1} and of the profile admissibility of Subsection \hyperref[sn:3.6]{3.6} holds by inspection of these finite records. The remaining components of the instance (hosts, audit records, source records, visibility records and bridges, defect records and nonclaims) consist of the records listed above, each written in its schema, and the dependency, gate and realization components are empty, so \(\mathsf{AdmDomain}\) of Subsubsection \hyperref[sn:4.10.1]{4.10.1} holds; each variant of Subsections \hyperref[sn:6.1]{6.1}--\hyperref[sn:6.5]{6.5} adds only records of these schemas.
\end{itemize}

This cell passes the checks of the well-formedness procedure of Subsubsection \hyperref[sn:6.3.1]{6.3.1}, and the classifier of Subsubsection \hyperref[sn:5.6.2]{5.6.2} assigns it the status \(\texttt{action}\). Two variations of it appear later. The same pair \(P_6\leftarrow P_3\) with a drive-certification update and a registered drive bridge too weak for the drive-certification profile is blocked (Theorem \hyperref[cl:thm:1]{1}). The pair observable on \(\{P_3,P_6\}\) with a drive profile is blocked while this cell is active (Theorem \hyperref[cl:thm:5]{5} and \(\mathsf{CM}_4\)): its source slot is an audited absence, so \(\delta_6^{\mathrm{source}}=\texttt{missing}\), and it has no drive bridge, so \(\delta_{\mathrm{reqbridge}}=\{\texttt{drive}\}\); either defect alone triggers the first \(\texttt{blocked}\) row.

The blocked cell of Theorem \hyperref[cl:thm:1]{1} keeps the same primitive pair and route-mismatch witness, but uses a drive-certification update under a profile requiring a drive bridge. The update's entry field \(x\), the candidate drive certificate, is suppressed and bridged only by the registered drive bridge \(L_{\mathrm{drive}}\), which is too weak for that profile; the audit record notes the failed strength comparison. The witness and the update are both present, so the row \(\texttt{absent}\) of the classifier table in Subsubsection \hyperref[sn:5.6.6]{5.6.6} does not apply. Both \(\delta_{\mathrm{weakbridge}}=\{(x,L_{\mathrm{drive}})\}\) and \(\delta_{\mathrm{reqbridge}}=\{\texttt{drive}\}\) are nonempty, so the row \(\texttt{blocked}\) applies, and the cell receives the status \(\texttt{blocked}\).
 \section{\texorpdfstring{The Central Object \(\mathbf{BirdInt}\)}{The Central Object \textbackslash mathbf\{BirdInt\}}}\label{sn:4}

Section \hyperref[sn:3]{3} declared the ambient domain \(\mathbf{BirdInt}^{\mathrm{aud}}_{\mathrm{fin}}\) as a finite tuple of finite typed components. As stated in Subsubsection \hyperref[sn:3.1.1]{3.1.1}, \(\mathbf{BirdInt}^{\mathrm{aud}}_{\mathrm{fin}}\) is a schema and the theorems quantify over its admissible instances. This section defines the judgment families that the rest of the paper consumes: role-channel, directed-cell, pair-observable, promotion, dependency, and instrument-indexed claim judgments, together with the audit, source, and admissibility records that the calculus carries on every judgment line. The judgment families are introduced at schema-grade; their classifier rules and theorem-grade results appear in Sections \hyperref[sn:5]{5} through \hyperref[sn:11]{11}. Subsection \hyperref[sn:4.7]{4.7} separates structural downward influence from a top-down causal channel, which must pass intervention gates. The definitions needed for Theorems \hyperref[cl:thm:1]{1}--\hyperref[cl:thm:4]{4} are listed in the reader's map (Subsection \hyperref[sn:1.4]{1.4}); the rest of this section is reference material.

\subsection{Central Object Definition}\label{sn:4.1}

\subsubsection{Object Form}\label{sn:4.1.1}

\begin{claimdefinition}{Central object $\mathbf{BirdInt}^{\mathrm{aud}}_{\mathrm{fin}}$}

The central object of this paper is the nineteen-component signature \(\mathbf{BirdInt}^{\mathrm{aud}}_{\mathrm{fin}}\) displayed in Subsubsection \hyperref[sn:3.1.1]{3.1.1}, whose instances are finite assignments of its components, equipped with the judgment families introduced in Subsections \hyperref[sn:4.2]{4.2} through \hyperref[sn:4.8]{4.8} and the source, audit, and admissibility records of Subsections \hyperref[sn:4.9]{4.9} and \hyperref[sn:4.10]{4.10}. The same notation \(\mathbf{BirdInt}^{\mathrm{aud}}_{\mathrm{fin}}\) refers to this object; no second name is introduced. It is distinct from the inherited Foundations II domain \(\mathsf{FATCD}\). The qualifier \(\mathrm{fin}\) records that every component is finite, and the qualifier \(\mathrm{aud}\) records that every accepted claim is interpreted under an explicit audit record.

The object is not a category, an algebra, or a closure operator. Each instance is a finite typed record over a finite host inventory, carrying finite typed records in each component. The judgment families of this section operate on the components directly, and a theorem about \(\mathbf{BirdInt}^{\mathrm{aud}}_{\mathrm{fin}}\) is a theorem about a finite check on those records, under explicit instrument, package, profile, level, visibility, threshold, audit, and nonclaim assumptions.

\end{claimdefinition}

\subsubsection{Internal Versus Model-Realization Components}\label{sn:4.1.2}

The components of \(\mathbf{BirdInt}^{\mathrm{aud}}_{\mathrm{fin}}\) split into two groups.

The \textbf{internal components} are \(\mathbb P\), \(\mathsf{Lev}\), \(\mathsf{Host}\), \(\mathsf{Cont}\), \(\mathsf{Instr}\), \(\mathsf{Pkg}\), \(\mathsf{Prof}\), \(\mathsf{Wit}\), \(\mathsf{Upd}\), \(\mathsf{Def}\), \(\mathsf{Judg}\), \(\mathsf{Status}\), \(\mathsf{Gate}\), \(\mathsf{Dep}\), \(\mathsf{Audit}\), \(\mathsf{Vis}\), \(\mathsf{Source}\), and \(\mathsf{Nonclaim}\). These are the components on which the calculus's syntax, classifier rules, defect-to-status rules, gate semantics, and instrument-indexed claim semantics operate. The theorems of Sections \hyperref[sn:6]{6}, \hyperref[sn:10]{10} and \hyperref[sn:11]{11} are theorems about the internal components; those of Sections \hyperref[sn:7]{7} and \hyperref[sn:8]{8} concern the finite maps, kernels and graphs that these components record.

The \textbf{model-realization component} is \(\mathsf{Real}\). This is the family of finite model-realization modules through which scoped finite hosts (PICA, the audited Cantor shell, finite abstract interpretation, finite graph/cohomology fragments) realize named fragments of the internal calculus. Statements about the model-realization component are tagged model-only when they assert a realization; they do not transfer to internal universal claims, and the negative scope of Subsubsection \hyperref[sn:2.4.3]{2.4.3} records this explicitly.

The split is enforced by the discipline of the paper: a theorem about an internal component does not silently use a model-realization fact, and a model-realization theorem is always tagged model-only and accompanied by the host, the realized fragment, the realization map, and the explicit nonclaims. Section \hyperref[sn:13]{13} carries the model-realization material; the present section defines only the internal judgment families.

\subsection{Role-Channel Judgments}\label{sn:4.2}

\subsubsection{Judgment Form}\label{sn:4.2.1}

A role-channel judgment fixes the standing of a primitive role channel under an instrument and a theory package. Its form is

\[
\Gamma\,;\;\mathcal T\,;\;I\;\vdash\;\mathrm{Chan}_i:\rho_i,
\]

where \(\Gamma\) is a finite typed context, \(\mathcal T\) is an admissible theory package, \(I\) is an admissible instrument, \(\mathrm{Chan}_i\) is the role channel for primitive \(P_i\in\mathbb P\), and \(\rho_i\) is the role-channel status assigned to \(\mathrm{Chan}_i\) under \(\Gamma,\mathcal T,I\). Unlike the directed-cell, pair-observable and promotion judgment forms, this form has no claim-type premise; scope enters only through rule 1 of Subsubsection \hyperref[sn:5.6.1]{5.6.1}.

The notation \(\mathrm{Chan}_i\) records that the subject of the judgment is the role channel of \(P_i\), not the primitive label \(P_i\) itself. Channels are typed slots through which role-aligned witness, threshold, level, visibility, and audit data are recorded; the primitive label only names the channel.

\subsubsection{Channel Statuses}\label{sn:4.2.2}

The role-channel status set is inherited from Foundations II:

\[
\begin{aligned}
\mathrm{RoleStatus}\;=\;\{\,&\texttt{active\_projection},\;\texttt{below\_threshold},\;\texttt{collapsed\_boundary\_case},\\
&\texttt{absent\_with\_record},\;\texttt{inapplicable\_outside\_scope}\,\}.
\end{aligned}
\]

The five values have the inherited meanings:

\begin{itemize}
\tightlist
\item
  \(\texttt{active\_projection}\): an active derived role projection exists with the required role-aligned witness, threshold, level, audit, and relevant visibility/support data.
\item
  \(\texttt{below\_threshold}\): candidate role evidence is present, but the threshold data required for activation is absent or not met.
\item
  \(\texttt{collapsed\_boundary\_case}\): the role-channel record's declared Boolean field \(\mathsf{collapse}_i\) is true.
\item
  \(\texttt{absent\_with\_record}\): no relevant raw features or witness data for the role are present, and the absence is explicitly recorded.
\item
  \(\texttt{inapplicable\_outside\_scope}\): the role-\(P_i\) content lies outside the scope covered by \(\mathcal T\) under \(I\).
\end{itemize}

Foundations III does not introduce additional role-channel statuses. The new status families introduced for directed cells, pair observables, promotions, claims, gates, and squares (Subsubsection \hyperref[sn:2.3.2]{2.3.2} and Subsection \hyperref[sn:5.1]{5.1}) extend the calculus; they do not replace or relabel the inherited five role-channel statuses.

\subsubsection{Channel/Activation Discipline}\label{sn:4.2.3}

A role channel is a typed slot. An active witness is a recorded value passing the threshold, level, witness, and audit conditions of the surrounding instrument and package. The two are distinct:

\begin{itemize}
\tightlist
\item
  a primitive may have an open role channel without an active witness, with status \(\texttt{absent\_with\_record}\) or \(\texttt{below\_threshold}\);
\item
  a primitive may have an active witness under one instrument and be inapplicable under another instrument whose scope excludes the role evidence, with status \(\texttt{inapplicable\_outside\_scope}\) there (rule 1 of Subsubsection \hyperref[sn:5.6.1]{5.6.1}); the role-channel classifier reads no profile;
\item
  a primitive may have status \(\texttt{active\_projection}\) on its role channel and yet not produce a directed cell with status \(\texttt{action}\), because directed-cell action additionally requires admissible actor/informant data, profile, visibility, threshold, audit, and defect records on the cell.
\end{itemize}

No theorem in this paper infers an active witness from the existence of a channel, and no theorem infers a directed-cell status from a role-channel status. The status-family separation theorem of Section \hyperref[sn:6]{6} makes the second of these formal.

\subsection{Directed-Cell Judgments}\label{sn:4.3}

Directed-cell judgments are the central typed bridge judgments of \(\mathbf{BirdInt}^{\mathrm{aud}}_{\mathrm{fin}}\). A directed cell records that, under an admissible context, package, and instrument, an actor primitive consumes typed witness data from an informant primitive through a typed update, with the level/lens/interface, visibility, threshold, audit, and defect data of the surrounding judgment. The cell is the working unit of typed interaction; the no-total-six-symbol-algebra theorem of Section \hyperref[sn:6]{6} records that cell status does not factor through the primitive pair.

\subsubsection{Judgment Form}\label{sn:4.3.1}

\begin{claimdefinition}{Directed-cell record}

The directed-cell judgment form is

\[
\Gamma\,;\;\mathcal T\,;\;I\;\vdash\;(\,P_i\leftarrow P_j,\;W_j,\;U_i,\;L,\;V,\;\Theta,\;A,\;\delta\,)^{\lambda}\;:\;\sigma,
\]

where \(\Gamma\) is a finite typed context, \(\mathcal T\) is an admissible theory package, \(I\) is an admissible instrument with \(\texttt{directed-cell}\in\mathsf{ClaimTypes}_I\) (a premise of the judgment form, as for promotion judgments, Subsubsection \hyperref[sn:4.5.1]{4.5.1}, not a classifier test; for claims, claim-type exclusion is instead a classifier outcome, Subsubsection \hyperref[sn:11.1.2]{11.1.2}), \(\lambda\in\mathsf{Prof}\) is an admissible profile, and \(\sigma\in\mathrm{CellStatus}\) is the directed-cell status assigned by the classifier of Subsubsection \hyperref[sn:5.6.2]{5.6.2}. The cell record itself is the finite typed tuple

\[
\mathsf{Cell}_{ij}^{\lambda}\;=\;(\,P_i\leftarrow P_j,\;W_j,\;U_i,\;L,\;V,\;\Theta,\;A,\;\delta\,),
\]

so the judgment may equivalently be written \(\Gamma;\mathcal T;I\vdash\mathsf{Cell}_{ij}^{\lambda}:\sigma\). The notation \(P_i\leftarrow P_j\) is a typed bridge designation, not a binary product on \(\mathbb P\); it records the actor and informant roles for the surrounding witness, update, and status data.

\end{claimdefinition}

\subsubsection{Actor and Informant Roles}\label{sn:4.3.2}

In the cell \(P_i\leftarrow P_j\), the actor primitive is \(P_i\) and the informant primitive is \(P_j\). When the slots are present, the actor receives a typed update \(U_i\in\mathsf{Upd}(P_i)\) and the informant supplies a typed witness \(W_j\in\mathsf{Wit}(P_j)\), with \(\mathsf{Wit}\) and \(\mathsf{Upd}\) as in Subsection \hyperref[sn:3.7]{3.7}. The arrow direction is part of the data: the cells \(P_i\leftarrow P_j\) and \(P_j\leftarrow P_i\) are distinct judgments with different witness/update typing requirements and may receive different statuses under the same package, instrument, and profile.

The slot rule of Subsubsection \hyperref[sn:3.7.3]{3.7.3} applies: each slot holds a record of the declared family or a typed, audited absence marker; a wrong-typed record or an unaudited marker is malformed. A typed, audited absence passes its slot check; if the remaining well-formedness checks pass, the cell classifier tests the marker.

\subsubsection{Cell Components}\label{sn:4.3.3}

The cell judgment uses the following finite typed records, each hosted and instrument-indexed as required by \(\mathcal T\), \(I\), and the profile \(\lambda\):

\begin{itemize}
\tightlist
\item
  \(W_j\): the informant-side witness data, a record with signature in \(\mathsf{Wit}(P_j)\) or a typed, audited absence marker for that slot.
\item
  \(U_i\): the actor-side update data, a record with signature in \(\mathsf{Upd}(P_i)\) or a typed, audited absence marker for that slot.
\item
  \(L\): the level/lens/interface record. \(L\) specifies the level tags on which the actor and informant operate, the lens through which the cell content is read, and the interface against which the cell is judged. Its level tags are the actor and informant level tags of \(\lambda\); if they differ, \(L\) contains a level-bridge record of one of the types of Subsubsection \hyperref[sn:2.2.3]{2.2.3}.
\item
  \(V\): the visibility record under \(I\) for the fields of the cell, drawn from the visibility partition of Subsection \hyperref[sn:3.5]{3.5}; it assigns each field of the cell, and each record of \(\mathsf{Use}(C)\), its class under \(\mathsf{vis}_I\), and a record whose marks differ from \(\mathsf{vis}_I\) fails condition (4) of well-formedness. \(V\) specifies which of \(W_j\), \(U_i\), \(L\), \(\Theta\), \(A\), \(\delta\) are visible, suppressed, outside-scope, or unknown under \(I\), and lists any visibility bridges admitted under \(\mathsf{BridgePolicy}_I\) that bring suppressed fields into admissible reach.
\item
  \(\Theta\): the threshold record. \(\Theta\) specifies the exact, approximate, tolerance, rank, visibility, and source thresholds the cell must meet, drawn from the eleven-element threshold family of Subsubsection \hyperref[sn:5.2.2]{5.2.2}. Each threshold names its evaluated datum: a direct check on witness or update values, a primitive defect of Subsection \hyperref[sn:5.3]{5.3}, a visibility or required-bridge defect of Subsection \hyperref[sn:5.5]{5.5}, the audit defect of a named audit record, or, for \(\Theta_{\mathrm{pass}}\), a recorded Boolean field of a record of \(\mathsf{Cont}\) named by \(\Theta\) (which then lies in \(\mathsf{Use}(C)\)).
\item
  \(A\): the audit slot, holding either a finite audit record in the schema of Subsubsection \hyperref[sn:4.9.2]{4.9.2} or a typed, audited absence marker. The marker has no check-result fields; its audit defect is \(\{\texttt{missing\_audit}\}\). \(A\) is distinct from the inherited audit functional \(\mathcal A\) of \(\mathcal T\): \(A\) is the finite audit record of the cell, and its relation to \(\mathsf{AuditPolicy}_I\) is read by the classifier through the audit defect of Subsubsection \hyperref[sn:5.3.6]{5.3.6}.
\item
  \(\delta\): the defect record family attached to the cell. Each entry of \(\delta\) is a typed defect record in the schema of Subsubsection \hyperref[sn:5.2.1]{5.2.1}, and \(\delta\) consists of (i) the visibility, required-bridge and audit defects of Subsection \hyperref[sn:5.5]{5.5} and Subsubsection \hyperref[sn:5.3.6]{5.3.6}, computed from the other fields; (ii) the package-collapse entries of Subsubsection \hyperref[sn:5.3.5]{5.3.5}, computed from \(L\), \(W\) and \(U\); (iii) for each primitive defect of Subsection \hyperref[sn:5.3]{5.3} explicitly named as the evaluated datum of a threshold in \(\Theta\), that defect computed on the witness and update data; and (iv) the Booleans \(\mathsf{noopU}\) and \(\mathsf{implicitW}\) of Subsubsection \hyperref[sn:4.3.4]{4.3.4}, recorded as computed, with passing value false. A payload defect to which \(\Theta\) attaches no threshold, such as \(\Delta_3^{\mathrm{comp}}\) of a \(P_3\) witness, is witness content, not an entry of \(\delta\).
\item
  \(\lambda\in\mathsf{Prof}\): the profile of the cell, fixing the actor and informant level tags, scale, regime, bridge type, instrument mode, and locality, in the schema of Subsubsection \hyperref[sn:3.6.1]{3.6.1}.
\item
  \(\sigma\in\mathrm{CellStatus}\): the directed-cell status assigned to the cell by the classifier. This is part of the judgment line, not an extra field of the cell tuple \(\mathsf{Cell}_{ij}^{\lambda}\).
\end{itemize}

The host tag of a cell is the host tag \(h_{\mathcal T}\) of its package; step 5 of Subsubsection \hyperref[sn:6.3.1]{6.3.1} checks that every component is a record on that host.

The use-set of the cell, \(\mathsf{Use}(C)\), consists of the records held or named by its fields \(W_j,U_i,L,\Theta,A\) together with every record named by a reference field of one of those records, except the instrument named in an audit record's \(\mathsf{CheckRules}\) field; the expansion is applied once (Subsubsection \hyperref[sn:5.5.2]{5.5.2}); the visibility defects of condition 7 are computed on it. The \(V\) and \(\delta\) slots do not themselves contribute references to \(\mathsf{Use}(C)\); references to these records occurring in the included fields are still counted. Bridge candidates listed in \(V\) are assessed through \(\mathsf{AdmVisBridge}\) where required, and \(\delta\) is recomputed in step 8.

A cell is \emph{well-formed} when the following eight conditions hold:

\begin{enumerate}
\def\labelenumi{\arabic{enumi}.}
\item
  its labels are primitives {[}decided by step 1 of Subsubsection \hyperref[sn:6.3.1]{6.3.1}{]};
\item
  its witness and update slots are typed or carry audited absence markers {[}decided by steps 2--3 of Subsubsection \hyperref[sn:6.3.1]{6.3.1}{]};
\item
  its profile is admissible {[}decided by step 4 of Subsubsection \hyperref[sn:6.3.1]{6.3.1}{]};
\item
  its \(L,V,\Theta,A,\delta,\lambda\) records are finite and typed with the level tags of \(L\) as above, and its present \(W_j\), \(U_i\) are records on the host of \(\mathcal T\) tagged \(\ell_j\), \(\ell_i\), and \(V\) assigns each field of the cell and each record of \(\mathsf{Use}(C)\) its class under \(\mathsf{vis}_I\) {[}decided by step 5 of Subsubsection \hyperref[sn:6.3.1]{6.3.1}{]};
\item
  its references resolve {[}decided by step 6 of Subsubsection \hyperref[sn:6.3.1]{6.3.1}{]};
\item
  its host and instrument tags are declared {[}decided by step 7 of Subsubsection \hyperref[sn:6.3.1]{6.3.1}{]};
\item
  its defect record \(\delta\) records the visibility and required-bridge defects of Subsection \hyperref[sn:5.5]{5.5}, the audit defect of Subsubsection \hyperref[sn:5.3.6]{5.3.6}, the package-collapse entries of Subsubsection \hyperref[sn:5.3.5]{5.3.5} and each primitive defect named by \(\Theta\), as computed from its other fields {[}decided by step 8 of Subsubsection \hyperref[sn:6.3.1]{6.3.1}{]};
\item
  \begin{enumerate}
  \def\labelenumii{(\alph{enumii})}
  \tightlist
  \item
    if the audit slot holds a present record \(A\): its \(\mathsf{ClaimRef}\) names \(C\) and its \(\mathsf{CheckRules}\) names the judging instrument \(I\), its \(\mathsf{VisibilityCheck}\) passes exactly when the visibility defect of condition 7 is empty, its \(\mathsf{ThresholdCheck}\) passes exactly when every threshold of \(\Theta\) evaluates to \(\texttt{pass}\), and its \(\mathsf{CircularityCheck}\) is the reachability test of Subsubsection \hyperref[sn:4.9.2]{4.9.2} with \(X=C\); (b) if it holds an audited absence marker: \(\delta_6^{\mathrm{audit}}=\{\texttt{missing\_audit}\}\); (c) in both cases: \(\mathsf{noopU}\) and \(\mathsf{implicitW}\) equal their computed values, and each type-(iii) defect entry carries exactly the threshold that \(\Theta\) assigns to it {[}decided by step 8 of Subsubsection \hyperref[sn:6.3.1]{6.3.1}{]}.
  \end{enumerate}
\end{enumerate}

Steps 1, 2--3, 4, 5, 6 and 7 decide conditions 1--6, and step 8 decides conditions 7 and 8. Consequently, on a well-formed cell, \(\delta\) and the check results of \(A\) are determined by the other fields: the classifier's status is a function of \(P_i,P_j,W,U,L\), the bridges listed in \(V\), \(\Theta\), the remaining fields of \(A\), \(\lambda\), \(\mathcal T\), \(I\) and the instance records they reference.

\emph{Covered cells.} A cell is \emph{covered} when it is well-formed and one of the nine rows \(\texttt{outside\_scope}\) through \(\texttt{action}\) of the table of Subsubsection \hyperref[sn:5.6.6]{5.6.6} holds at it; the final \(\texttt{blocked}\) row is a totality guard. The notion lets Theorem \hyperref[cl:thm:1]{1} be stated without Theorem \hyperref[cl:thm:4]{4}, which shows that every well-formed cell is covered, so that \(\mathsf{CoveredCell}\) is the set of well-formed cells. A missing witness or update, and content outside \(\mathsf{Scope}_I\), are classifier outcomes (\(\texttt{absent}\), \(\texttt{outside\_scope}\)), not well-formedness failures. The well-formedness step is a finite decidable check (Theorem \hyperref[cl:thm:3]{3}, Subsection \hyperref[sn:6.3]{6.3}); the resulting status \(\sigma\) is then determined by the rows of the directed-cell table of Subsubsection \hyperref[sn:5.6.6]{5.6.6} (glossed in Subsubsection \hyperref[sn:5.6.2]{5.6.2}) from the \(L,V,\Theta,A,\delta\) data and the profile. In the running example of Subsection \hyperref[sn:3.8]{3.8}, the witness is the route-mismatch record \((E_1,E_2,\{a\})\), the update adds an audit record, both level tags are \(\mathsf{Ver}\), and \(\delta\) is empty.

\subsubsection{Directed-Cell Statuses}\label{sn:4.3.4}

The directed-cell status set is

\[
\begin{aligned}
\mathrm{CellStatus}\;=\;\{\,&\texttt{action},\;\texttt{implicit},\;\texttt{trivial},\;\texttt{blocked},\;\texttt{undefined\_circular},\\
&\texttt{below\_threshold},\;\texttt{collapsed},\;\texttt{absent},\;\texttt{outside\_scope}\,\}.
\end{aligned}
\]

In words (Subsubsection \hyperref[sn:5.6.2]{5.6.2} gives the exact rules and their priority order): \(\texttt{action}\), the cell is well formed, its witness and update are present and every defect, visibility, audit and source check passes; \(\texttt{implicit}\), the witness is attached to the package and no threshold evaluates it; \(\texttt{trivial}\), the update changes nothing (\(\mathsf{noopU}\)) or re-installs the package map or completion already recorded for its target (\(\texttt{identity\_package}\), Subsubsection \hyperref[sn:5.3.5]{5.3.5}; an identity lens or completion of \(\mathcal T\) itself produces no entry); \(\texttt{below\_threshold}\), the witness and update are present and the audit's threshold check fails; \(\texttt{collapsed}\), the lens of \(L\) collapses the witness content; \(\texttt{undefined\_circular}\), the cell's own verdict is reachable from its audit's evidence or sources; \(\texttt{blocked}\), a bridge required by the profile is missing or too weak, a used record is suppressed without an admissible bridge or has unknown visibility, or the audit is missing or records a failed source, visibility, nonclaim or replay check; \(\texttt{absent}\), the witness or update is an audited absence marker and the audit defect contains neither \(\texttt{missing\_audit}\) nor \(\texttt{failed\_audit}\) (otherwise the cell is \(\texttt{blocked}\)); \(\texttt{outside\_scope}\), the instrument does not cover the cell's host, level or used records. The classifier assigns exactly one of these to a covered well-formed cell (Theorem \hyperref[cl:thm:4]{4}, Subsection \hyperref[sn:6.4]{6.4}). The classifier priority order, fixed for the rest of the paper, is displayed with the per-status rules in Subsubsection \hyperref[sn:5.6.2]{5.6.2}.

The field \(\mathsf{implicitW}\) holds when the witness \(W\) lies in the witness family \(W_{\mathcal T}\) attached to the package \(\mathcal T\) and no threshold of \(\Theta\) evaluates it, that is, has as its evaluated datum a field of \(W\), a direct check on values of \(W\), or a primitive defect computed from fields of \(W\) (status \(\texttt{implicit}\), Subsubsection \hyperref[sn:5.6.6]{5.6.6}); the well-formedness check computes \(\mathsf{noopU}\), which is \([e_U(r)=r]\) when the update \(U\) has a recorded finite effect map \(e_U\) on its target \(r\) and false otherwise, and which decides the status \(\texttt{trivial}\). Subsubsection \hyperref[sn:5.6.6]{5.6.6} lists the declared inputs of all the classifiers.

Under this priority, a cell whose data simultaneously matches more than one status condition is assigned the highest-priority match; the per-status conditions are the rows of the directed-cell table of Subsubsection \hyperref[sn:5.6.6]{5.6.6}, glossed in Subsubsection \hyperref[sn:5.6.2]{5.6.2}, and the uniqueness theorem is Theorem \hyperref[cl:thm:4]{4}. The priority ordering reflects the working discipline of the calculus: out-of-scope and absence verdicts are recorded before any further check on the cell's content; blocking, circularity, and collapse are recorded before threshold or activity verdicts; and the \(\texttt{action}\) verdict is reached only after every higher-priority condition has been ruled out.

\subsection{Pair-Observable Judgments}\label{sn:4.4}

A directed cell records one role acting on another. A pair observable instead records joint content for an unordered pair of roles, with a source-of-truth record. Its classifier uses that source together with the pair's compatibility, visibility, threshold and audit records to assign \(\texttt{real}\), \(\texttt{real\_provisional}\), \(\texttt{real\_duplicated}\) or \(\texttt{blocked}\). Pair-observable judgments record the standing of an unordered pair of primitives as a joint observable, distinct from the directed-cell judgments of Subsection \hyperref[sn:4.3]{4.3}. The pair-observable family has its own status set, its own classifier inputs (notably a source-of-truth record), and its own admissibility conditions. Directed-cell action does not by itself imply pair-observable realness; the schema-level distinction is recorded here, and the formal status-family separation appears in Theorem \hyperref[cl:thm:5]{5} of Section \hyperref[sn:6]{6}.

\subsubsection{Judgment Form}\label{sn:4.4.1}

\begin{claimdefinition}{Pair-observable record}

The pair-observable judgment form is

\[
\Gamma\,;\;\mathcal T\,;\;I\;\vdash\;\mathsf{PairObs}_{\{i,j\}}^{\mu}\;:\;\omega,
\]

where \(\{i,j\}\subseteq\{1,\ldots,6\}\) is an unordered pair of distinct primitive indices (\(i\neq j\); there are fifteen such pairs), \(\mu\in\mathsf{Prof}\) is a profile drawn from the family of profiles for pair-observable judgments, the instrument \(I\) has \(\texttt{pair-observable}\in\mathsf{ClaimTypes}_I\), and \(\omega\in\mathrm{PairStatus}\) is the pair-observable status assigned to the record by the classifier of Subsubsection \hyperref[sn:5.6.3]{5.6.3}. The unordered pair \(\{i,j\}\) equals \(\{j,i\}\); the judgment is the same for either ordering.

A pair-observable record is a finite typed tuple recording, in addition to \(\{i,j\}\) and \(\mu\), the joint witness data, the source-of-truth record, the branch-compatibility record, the visibility record, the threshold record, the audit record, the defect record, and the nonclaim record, together with declared Boolean fields \(\mathsf{forbidsSource}\) and \(\mathsf{partialEvidence}\) that the surrounding instrument and package require, and optionally a duplicate-pair record \(\mathsf{Dup}\), read by the \(\texttt{real\_duplicated}\) row of Subsubsection \hyperref[sn:5.6.3]{5.6.3}. The source-of-truth field is the record's \emph{source slot}: it holds a source record \(R\) of Subsubsection \hyperref[sn:4.9.1]{4.9.1} or a typed, audited absence marker. In the PICA display of Subsubsection \hyperref[sn:13.2.7]{13.2.7} the joint-witness, source and branch fields are named \(\mathsf{ObservableType}\), \(\mathsf{SourceOfTruth}\) and \(\mathsf{BranchData}\). A duplicate-pair record \(\mathsf{Dup}=(Q',A_{\mathsf{Dup}})\) names a second pair record \(Q'\neq Q\) of the instance, on the same unordered pair, that records the same joint content, with its own passing audit; the classifier reads only its presence. The full schema is registered in \(\mathsf{Cont}\). A pair-observable record is \emph{well formed} when its fields are present, finite and typed, \(\mathsf{AdmProfile}(\mu,I,\mathcal T)\) holds, every reference resolves or is an audited absence (a missing source is recorded as \(\delta_6^{\mathrm{source}}=\texttt{missing}\)), its audit slot holds an audit record whose \(\mathsf{ClaimRef}\) names the pair record, or a typed, audited absence marker, whose audit defect is \(\{\texttt{missing\_audit}\}\) (Subsubsection \hyperref[sn:5.3.6]{5.3.6}), and its visibility, required-bridge, source and audit defect entries equal the values computed as in Subsections \hyperref[sn:5.5]{5.5} and \hyperref[sn:5.3]{5.3}; policy failures are classifier outcomes, not well-formedness failures. The branch-compatibility record carries a Boolean field \(\mathsf{compatible}\) and its own audit record. Each threshold of a pair record names its evaluated datum as for a directed cell (Subsubsection \hyperref[sn:4.3.3]{4.3.3}), the source-of-truth and branch-compatibility records counting as witness data.

\end{claimdefinition}

\subsubsection{Pair Versus Directed Cell}\label{sn:4.4.2}

A pair observable is not a directed cell. The directed cell \(P_i\leftarrow P_j\) records typed actor/informant bridge data; the pair observable \(\mathsf{PairObs}_{\{i,j\}}^{\mu}\) records joint source-of-truth content for the unordered pair. The two are distinct judgment families with distinct status sets, distinct classifier inputs, and distinct admissibility schemas.

In particular,

\[
\Gamma\,;\;\mathcal T\,;\;I\;\vdash\;\mathsf{Cell}_{ij}^{\lambda}\;:\;\texttt{action}\;\;\not\Rightarrow\;\;\Gamma\,;\;\mathcal T\,;\;I\;\vdash\;\mathsf{PairObs}_{\{i,j\}}^{\mu}\;:\;\texttt{real}.
\]

A directed cell with status \(\texttt{action}\) records that the typed actor/informant bridge passes the well-formedness, threshold, visibility, and audit conditions of the cell record. A pair observable with status \(\texttt{real}\) additionally records that the unordered pair has a real source of truth, an admissible branch-compatibility record, and an admissible source policy under \(I\). The countermodel sheet \(\mathsf{CM}_4\) of Section \hyperref[sn:9]{9} supplies a finite construction in which a directed cell is \(\texttt{action}\) and the corresponding pair observable is \(\texttt{blocked}\). Section \hyperref[sn:6]{6} contains the theorem-grade form of this separation as part of Theorem \hyperref[cl:thm:5]{5}; this subsection records only the schema-level distinction.

\subsubsection{Source-of-Truth Conditions}\label{sn:4.4.3}

Pair-observable acceptance depends on a source-of-truth record drawn from a finite typed family. The source classes used by this paper, drawn from the source-policy values of Subsubsection \hyperref[sn:4.9.1]{4.9.1}, are:

\begin{itemize}
\tightlist
\item
  \textbf{real source.} A finite source record, certified under \(\mathsf{SourcePolicy}_I\) with admissible audit and bridge data, that supports the unordered pair under profile \(\mu\). Here ``real'' requires a source admitted by the instrument's source policy and passing the applicable checks; fallback and proxy sources are provisional by default, subject to the audited-upgrade rule for fallback stated in the next item.
\item
  \textbf{fallback source.} A finite source record certified as a designated fallback under \(\mathsf{SourcePolicy}_I\). A fallback source supports a pair observable at \(\texttt{real\_provisional}\) at most by default; it does not support \(\texttt{real}\) unless its record carries an audited source-upgrade bridge admitted by \(\mathsf{SourcePolicy}_I\) and \(\mathsf{BridgePolicy}_I\) (the \(\texttt{real}\) row of the pair classifier, Subsubsection \hyperref[sn:5.6.3]{5.6.3}).
\item
  \textbf{dirty source.} A finite source record whose trace records a failed audit or threshold check of the record itself, or whose validity record places it outside its validity condition (the \(\texttt{dirty}\) clause of Subsubsection \hyperref[sn:5.3.6]{5.3.6}); failed checks that the record only holds as data do not make it dirty. A dirty source does not support an accepted pair observable; the resulting status is \(\texttt{blocked}\).
\item
  \textbf{proxy source.} A finite source record obtained via a typed bridge from another source, admissible only when the bridge is recorded in \(\mathsf{BridgePolicy}_I\) and audited; a proxy source supports a pair observable at most at \(\texttt{real\_provisional}\).
\item
  \textbf{missing source.} No source record is supplied. The pair observable does not pass acceptance: the status is \(\texttt{blocked}\).
\end{itemize}

The pair-observable status set used by this paper is

\[
\mathrm{PairStatus}\;=\;\{\,\texttt{real},\;\texttt{real\_provisional},\;\texttt{real\_duplicated},\;\texttt{blocked}\,\}.
\]

The source-condition names above--fallback source, dirty source, proxy source, and missing source--are not additional pair-observable statuses. They are source-policy conditions that the classifier routes into one of the four values in \(\mathrm{PairStatus}\).

The classifier of Subsubsection \hyperref[sn:5.6.3]{5.6.3} consumes the pair record together with its source-of-truth record, branch-compatibility record, visibility record, threshold record, audit record, and defect record, and assigns a value in \(\mathrm{PairStatus}\) to every well-formed pair-observable record. As with directed cells, no theorem in this paper infers a pair-observable status from a directed-cell status, a role-channel status, or a primitive pair alone; the schema discipline of this subsection together with the status-family separation theorem of Section \hyperref[sn:6]{6} closes that gap.

\subsection{Promotion Judgments}\label{sn:4.5}

Promotion judgments record the standing of a finite promotion bridge from a source theory package to a target theory package under an instrument. They are the typed/statused machinery on which the strict-extension theorems of Section \hyperref[sn:7]{7} and the promotion-gate-soundness theorem of Section \hyperref[sn:10]{10} are built. Like directed-cell and pair-observable judgments, promotion judgments have their own status set and their own classifier inputs; in particular, the strict and non-strict refinements of an accepted promotion are properties of the bridge's object-map data and do not by themselves imply closure or drive on the target package.

\subsubsection{Judgment Form}\label{sn:4.5.1}

The promotion judgment form is

\[
\Gamma\,;\;\mathcal T_j\,;\;I\;\vdash\;\mathrm{Promote}(B_{j\to j+1})\;:\;\sigma_{\mathrm{promote}},
\]

where \(B_{j\to j+1}\) is the promotion bridge from source theory package \(\mathcal T_j\) to target theory package \(\mathcal T_{j+1}\), \(I\) is an admissible instrument with \(\mathsf{ClaimTypes}_I\) that includes promotion claims, and \(\sigma_{\mathrm{promote}}\in\mathrm{PromotionStatus}\) is the promotion status assigned by the classifier of Subsubsection \hyperref[sn:5.6.4]{5.6.4}.

\subsubsection{Promotion Bridge Object}\label{sn:4.5.2}

\begin{claimdefinition}{Promotion bridge record}

A promotion bridge is the finite typed tuple

\[
B_{j\to j+1}\;=\;(\,\mathcal T_j,\;\mathcal T_{j+1},\;\pi_j,\;\pi_{j+1},\;L,\;V,\;\Theta,\;A,\;\delta,\;\mathcal N\,),
\]

with the following components, all finite:

\begin{itemize}
\tightlist
\item
  \(\mathcal T_j\): the source theory package, admissible in the sense of Subsubsection \hyperref[sn:3.3.3]{3.3.3}.
\item
  \(\mathcal T_{j+1}\): the target theory package, admissible in the same sense.
\item
  \(\pi_j\): the old object map or interface, recording how content is read out of \(\mathcal T_j\).
\item
  \(\pi_{j+1}\): the new object map or interface, recording how content is read out of \(\mathcal T_{j+1}\).
\item
  \(L\): the level/lens/interface record of the bridge, fixing the levels at which \(\pi_j\) and \(\pi_{j+1}\) operate and, when drive is judged, a finite support graph \(G_B\) on the carrier \(S\) of the expanded form below.
\item
  \(V\): the visibility record of the bridge under \(I\), in the schema of Subsection \hyperref[sn:3.5]{3.5}.
\item
  \(\Theta\): the threshold record of the bridge, drawn from the threshold family of Subsubsection \hyperref[sn:5.2.2]{5.2.2}.
\item
  \(A\): the audit record of the bridge, distinct from the audit functionals \(\mathcal A_j\), \(\mathcal A_{j+1}\) of the source and target packages, whose \(\mathsf{ClaimRef}\) names \(B_{j\to j+1}\).
\item
  \(\delta\): the defect record family attached to the bridge, including the strictness, no-smuggling, stability, descent, and audit defects defined in Subsections \hyperref[sn:5.3]{5.3} and \hyperref[sn:5.4]{5.4}.
\item
  \(\mathcal N\): the nonclaim record attached to the bridge.
\end{itemize}

The expanded form of \(B_{j\to j+1}\) recorded in \(\mathsf{Cont}\) also carries a bridge identifier, the carrier \(S\) on which the object maps act, the codomains \(O_j\) and \(O_{j+1}\), the gate-results record, the host tag, and the bridge profile:

\[
\begin{aligned}
B_{j\to j+1}=\bigl(&\mathrm{id},\mathcal T_j,\mathcal T_{j+1},S,O_j,O_{j+1},\pi_j,\pi_{j+1},\\
&L,V,\Theta,A,\delta,\mathsf{GateResults},\mathcal N,\mathsf{HostTag},\lambda_{\mathrm{prom}}\bigr).
\end{aligned}
\]

The last field is the bridge profile, a variable of the schema; the named profile \(\lambda_{\mathrm{prom}}\) of Subsubsection \hyperref[sn:10.6.1]{10.6.1} is one value of it.

A bridge is admissible when its well-formedness gate \(G_{\mathrm{suff}}\) passes (Subsubsection \hyperref[sn:10.2.1]{10.2.1}); the other gates are inputs to the promotion classifier of Subsubsection \hyperref[sn:5.6.4]{5.6.4}, not admissibility conditions.

A promotion bridge has \textbf{strict object maps} when

\[
\pi_{j+1}\;\not\factor\;\pi_j,
\]

equivalently when there exist \(s,s'\in S\) with \(\pi_j(s)=\pi_j(s')\) and \(\pi_{j+1}(s)\neq\pi_{j+1}(s')\). A promotion bridge has \textbf{non-strict object maps} when \(\pi_{j+1}\) factors through \(\pi_j\), that is, when there exists a finite map \(\phi:\mathrm{im}(\pi_j)\to O_{j+1}\) with \(\pi_{j+1}=\phi\,\pi_j\). Strictness is a property of the typed object-map data. Strict object maps need not refine the old ones: \(\pi_{j+1}\) may also merge points that \(\pi_j\) separates, and strictness of two consecutive bridges then need not give strictness of the composite bridge (Subsubsection \hyperref[sn:10.7.1]{10.7.1}); when \(\pi_j=r\circ\pi_{j+1}\), strictness is strict refinement of the fibre partition (Subsubsection \hyperref[sn:10.1.2]{10.1.2}). Strictness is not by itself a closure or drive claim, and the negative scope of Subsubsection \hyperref[sn:2.4.3]{2.4.3} records the corresponding nonclaims (NC-10 and NC-11 of Subsection \hyperref[sn:15.2]{15.2}).

\end{claimdefinition}

\subsubsection{Promotion Statuses}\label{sn:4.5.3}

The promotion status set is

\[
\begin{aligned}
\mathrm{PromotionStatus}\;=\;\{\,&\texttt{candidate},\;\texttt{accepted},\;\texttt{strict},\;\texttt{non\_strict},\\
&\texttt{failed\_descent},\;\texttt{failed\_stability},\;\texttt{failed\_audit},\;\texttt{failed\_no\_smuggling},\\
&\texttt{local\_only},\;\texttt{globally\_obstructed},\;\texttt{outside\_scope}\,\}.
\end{aligned}
\]

The values \(\texttt{strict}\) and \(\texttt{non\_strict}\) are accepted-status refinements: a strict promotion is an accepted promotion in which the strictness condition \(\pi_{j+1}\not\factor\pi_j\) holds, and a non-strict promotion is an accepted promotion in which \(\pi_{j+1}\) factors through \(\pi_j\). The classifier of Subsubsection \hyperref[sn:5.6.4]{5.6.4} produces an accepted, strict, or non-strict verdict only when every required core gate in \(\mathsf{GateResults}(B)\) passes; the failure-tagged values record specific gate failures (descent, stability, audit, no-smuggling), and \(\texttt{local\_only}\) and \(\texttt{globally\_obstructed}\) record the optional local-to-global gate's verdict. The verdict \(\texttt{candidate}\) records a bridge whose required core gates have not yet been resolved, typically because at least one core gate has status \(\texttt{not\_checked}\).

A promotion bridge with status \(\texttt{strict}\) does not by that fact alone imply macro closure, descended dynamics, or drive: NC-10 and NC-11 of Subsection \hyperref[sn:15.2]{15.2} record the corresponding strictness nonclaims, and the later no-go results are stated against these would-be implications.

\subsection{Dependency and No-Go Judgments}\label{sn:4.6}

Dependency judgments record typed relationships between content classes under an instrument. They are the working format for the no-go and bridge-required claims that appear throughout the paper, including the route-mismatch-requires-bridge dependency that keeps the drive face \(\mathrm{P6}_{\mathrm{drive}}\) separate from generic \(P_3\) holonomy. Like the other judgment families, dependencies have their own admissibility schema; theorem-grade no-go statements supported by countermodels are recorded in Sections \hyperref[sn:7]{7} through \hyperref[sn:10]{10} and Appendix \hyperref[sn:C]{C}.

\subsubsection{Judgment Form}\label{sn:4.6.1}

The dependency judgment form is

\[
\Gamma\,;\;\mathcal C\,;\;I\;\vdash\;S\;\mathrel{\mathsf{dep}_r}\;X\;[\,H,\;A,\;V,\;\Theta\,],
\]

where \(\Gamma\) is a finite typed context, \(\mathcal C\) is a finite content context (typically a theory package, a model-realization module, or a content fragment of \(\mathsf{Cont}\)), \(I\) is an admissible instrument, \(S\) and \(X\) are content classes (primitive labels, primitive role channels, named structural records, drive certificates, gating data, or any record drawn from \(\mathsf{Cont}\)), \(r\) is the relation type, and \(H,A,V,\Theta\) are the host, audit, visibility, and threshold records under which the dependency is asserted.

The relation types admitted by this paper are

\[
\begin{aligned}
\{\,&\texttt{suffices},\;\texttt{insufficient},\;\texttt{destroys},\;\texttt{blocks},\\
&\texttt{enables},\;\texttt{equivalent},\;\texttt{requires\_bridge},\;\texttt{outside\_scope}\,\}.
\end{aligned}
\]

The relation type fixes the polarity and form of the dependency; the host, audit, visibility, and threshold records fix the scope under which the dependency is asserted.

\emph{Truth conditions.} Let \(\Sigma\) be the class of admissible instances that contain the records \(\Gamma,\mathcal C,I,H,A,V,\Theta\) and in which \(I\) admits the dependency, that is, \(\texttt{dependency}\in\mathsf{ClaimTypes}_I\), \(H\in\mathsf{Hosts}_I\), and the records \(\Gamma,\mathcal C,A,V,\Theta\) lie in \(\mathsf{Scope}_I\) with level tags in \(\mathsf{Levels}_I\), and read \(S,X\) as predicates on the instances of \(\Sigma\). \(S\mathrel{\mathsf{dep}_{\texttt{suffices}}}X\) holds when every instance in \(\Sigma\) satisfying \(S\) satisfies \(X\); \(\texttt{insufficient}\) when some instance satisfies \(S\) and not \(X\) (\(S\not\Rightarrow X\), Subsubsection \hyperref[sn:2.3.1]{2.3.1}); \(\texttt{blocks}\) when every instance satisfying \(S\) fails \(X\); \(\texttt{destroys}\) when \(S\) names an operation on instances and, for every instance \(\mathfrak D\) in \(\Sigma\) satisfying \(X\) to which it applies, the result \(S(\mathfrak D)\) fails \(X\); \(\texttt{equivalent}\) when \(S\Leftrightarrow X\) on \(\Sigma\); \(\texttt{enables}\) when some instance satisfies both; \(\texttt{requires\_bridge}\) when \(S\mathrel{\mathsf{dep}_{\texttt{insufficient}}}X\) and \(S\) together with an admissible drive bridge (Subsubsection \hyperref[sn:9.3.5]{9.3.5}) \(\mathrel{\mathsf{dep}_{\texttt{suffices}}}X\); this paper uses \(\texttt{requires\_bridge}\) only with \(X\) a drive predicate (Subsubsection \hyperref[sn:4.6.2]{4.6.2}, \(\mathsf{CM}_2\)). When a drive bridge by itself carries a certificate of \(X\), as \(B^{\mathrm{drive}}_{3\to6}\) does in Subsubsection \hyperref[sn:9.3.5]{9.3.5}, the second conjunct holds for every \(S\); the judgment then asserts that \(S\) is insufficient for \(X\) and that a drive bridge is one admissible route to \(X\). The judgment does not assert necessity: instances of \(\Sigma\) may satisfy \(X\) with no bridge of the named type (a cochain with a nonzero cycle integral certifies drive without \(P_3\) data). Beyond \(\texttt{insufficient}\), its additional requirement is that \(S\) together with an admissible bridge of the named type suffices for \(X\). Finally, \(S\mathrel{\mathsf{dep}_{\texttt{outside\_scope}}}X\) holds when \(I\) does not admit the dependency.

\subsubsection{Dependency Claims}\label{sn:4.6.2}

The dependency judgment family supports four categories of statement, each appearing throughout the paper.

\textbf{Positive dependency.} A statement of the form \(S\mathrel{\mathsf{dep}_{\texttt{suffices}}}X\) or \(S\mathrel{\mathsf{dep}_{\texttt{enables}}}X\) records that, under the named scope, \(S\) supplies or enables \(X\). Positive dependencies in this paper are typically theorem-grade results stated under explicit host, instrument, and audit assumptions: the descent theorem of Section \hyperref[sn:7]{7}, for example, gives a positive dependency from a vanishing descent defect to the existence of a descended map.

\textbf{Blocked dependency.} A statement of the form \(S\mathrel{\mathsf{dep}_{\texttt{insufficient}}}X\), \(S\mathrel{\mathsf{dep}_{\texttt{blocks}}}X\), or \(S\mathrel{\mathsf{dep}_{\texttt{destroys}}}X\) records that \(S\) does not support \(X\) under the named scope, or actively obstructs it. Blocked dependencies are the natural form of the paper's no-go results: the gating-affinity-suppression theorem of Section \hyperref[sn:8]{8}, for instance, records a \(\texttt{destroys}\) dependency from \(P_2\)-style support deletion to cohomological drive on the resulting forest.

\textbf{Bridge-required dependency.} A statement of the form \(S\mathrel{\mathsf{dep}_{\texttt{requires\_bridge}}}X\) records that \(S\) does not by itself yield \(X\), but does so via a typed bridge admitted by \(\mathsf{BridgePolicy}_I\). The canonical instance in this paper is the route-mismatch dependency

\[
P_3\;\mathrel{\mathsf{dep}_{\texttt{requires\_bridge}}}\;\mathrm{P6}_{\mathrm{drive}},
\]

recording that route mismatch or holonomy alone does not establish a drive certificate; an explicit drive bridge is required to infer a positive drive claim from route-mismatch data; a drive claim can also be certified directly by a cochain with a nonzero cycle integral (Subsubsection \hyperref[sn:8.1.2]{8.1.2}); a directed cell under the \(\texttt{drive-certification}\) regime needs the required bridge of Subsubsection \hyperref[sn:5.5.2]{5.5.2}, a drive-type visibility bridge when the certified datum is suppressed, or a drive bridge of Subsubsection \hyperref[sn:9.3.5]{9.3.5} named by its drive-check field when the certified datum is visible.

\textbf{Countermodel-supported no-go.} A statement of the form \(S\mathrel{\mathsf{dep}_{\texttt{insufficient}}}X\) or \(S\mathrel{\mathsf{dep}_{\texttt{blocks}}}X\) supported by a finite explicit countermodel construction. Sections \hyperref[sn:7]{7}, \hyperref[sn:8]{8}, \hyperref[sn:9]{9}, and \hyperref[sn:10]{10} use this category for their no-go and separation results; the countermodel atlas of Section \hyperref[sn:9]{9} and Appendix \hyperref[sn:C]{C} records the constructions in full.

A dependency judgment is admissible when every component is finite and consistent with the surrounding instrument and content context, and when its relation label \(r\) is one of the eight allowed values above. The label \(\texttt{outside\_scope}\) records the case in which the surrounding instrument does not adjudicate the dependency.

\subsection{Structural Downward Influence Versus Top-Down Causal Channels}\label{sn:4.7}

The calculus records several forms of downward influence across packages, profiles, levels, and hosts. A macro package may select a representative substrate record by a completion or lift; a macro admissibility condition may exclude substrate states or histories; a phase, profile, or protocol record may index which substrate update rule is active. These are structural downward influences: they are paths, restrictions, or context-indexing records from a macro record to a substrate record.

They are not, by themselves, top-down causal channels. In this subsection ``top-down channel'' means an intervention-respecting difference-making record; it is distinct from the role channels \(\mathrm{Chan}_i\) of Subsection \hyperref[sn:4.2]{4.2}, which are slots for the six roles. A completion, lift, representative map, feasibility restriction, or context index is a structural path unless the top-down channel gates pass. Such a record must exhibit an admissible finite family of macro interventions and an audited comparison showing that changing the macro intervention changes a substrate-future response under matched controls.

A macro record is a finite level-tagged or lens-related record \(M\) in the current package: for example a package object, profile parameter, phase, protocol, admissibility gate, policy, abstract value, or promoted object-map value. A substrate record \(X\) is a lower-level or finer-lens future observable, state predicate, update trace, or output record in the same host or a connected host. The macro/substrate relation is therefore relative to a package, lens, level profile, host, and bridge record, rather than an external micro/macro ontology.

\begin{claimdefinition}{Top-down channel record}

A top-down channel record is a finite tuple

\[
\mathsf{TopDownChannelRecord}_{M\to X}^{\lambda}
=
(M,X,\mathsf{StructPath},\mathsf{Interv}_M,\mathsf{Ctrl},\mathsf{Resp},D,\Theta,L,V,A,
\delta_{\mathrm{td}},\mathcal N).
\]

Here \(M\) is a macro record relative to \((\mathcal T,I,\lambda)\), \(X\) is a substrate-future record, \(\mathsf{StructPath}\) is a declared record from \(M\) to \(X\) of one of the five structural kinds (completion, lift, representative map, feasibility restriction, context index) or an audited absence marker, \(\mathsf{Interv}_M\) is a finite family of admissible macro interventions, \(\mathsf{Ctrl}\) records matched controls, \(\mathsf{Resp}\) records the substrate response under each intervention, \(D\) is the host-supplied effect comparator, \(\Theta\) is the effect threshold, \(L,V,A\) are bridge, visibility, and audit records, \(\delta_{\mathrm{td}}\) is the measured difference-making defect or certificate, and \(\mathcal N\) records the required nonclaims.

\end{claimdefinition}

The calculus does not fix a universal probability metric for top-down channels. The host supplies the comparator \(D_{\mathsf H}\) and threshold \(\Theta_{\mathrm{td}}\); \(\Theta_{\mathrm{td}}\) is an instance of \(\Theta_{\ge\varepsilon}\) or \(\Theta_{\in S}\) (Subsubsection \hyperref[sn:5.2.2]{5.2.2}) on \(\delta_{\mathrm{td}}\). A probabilistic host may instantiate the comparator by total variation, KL divergence, or a channel-capacity quantity; a deterministic host may instantiate it by finite output inequality. The present calculus requires only that the comparator, threshold, source records, and audit be finite and visible or bridged.

\emph{Reading the statements.} \(\mathsf{StructDown}(M,X)\) says that a macro record \(M\) constrains, indexes or maps to a substrate record \(X\) by one of the structural paths listed above: a completion, lift, representative map, feasibility restriction or context index. \(\mathsf{TopDownChannel}(R)\) is the claim that \(R\) is a top-down channel from \(M\) to \(X\). Its claim-type tag is \(\texttt{top\_down\_channel}\), its \(\mathsf{Refs}\) consists of \(R\) and all records held or named by its fields, and its audit field is the audit record \(A\) of \(R\), whose \(\mathsf{ClaimRef}\) names this claim; its use-set is that of Subsubsection \hyperref[sn:3.2.1]{3.2.1}, and \(G_{\mathrm{audit}}\), \(G_{\mathrm{source}}\) and \(G_{\mathrm{vis}}\) below read this \(A\) and this use-set. It is classified by the claim classifier of Subsubsection \hyperref[sn:5.6.5]{5.6.5}; for this tag the check rules of every admissible \(I\) that admits the tag assign the rejection verdict unless every gate of \(\mathsf{TDGate}(R)\), whose nine gates are macro record, intervention, matched controls, feasibility, source, visibility, audit, no-smuggling and effect (in particular \(G_{\mathrm{interv}}\) passes iff \(|\mathsf{Interv}_M|\ge2\), and \(G_{\mathrm{effect}}\) passes iff the set of comparator values of distinct matched interventions is nonempty and its maximum passes \(\Theta_{\mathrm{td}}\); full conditions after the third proposition), passes (Subsubsection \hyperref[sn:3.4.1]{3.4.1}). The display of \(\mathsf{TDGate}(R)\) and the paragraph after it, which follow the third proposition and the satisfiability example, give the pass conditions of all nine gates and define \(\mathcal M\) and \(\delta_{\mathrm{td}}\), which the three proofs use.

\begin{claimproposition}{Top-down channel gate soundness}

The claim

\[
\Gamma;\mathcal T;I
\vdash
\mathsf{TopDownChannel}(R):\texttt{accepted}
\]

is accepted only if the following gates pass: macro-record well-formedness, admissible intervention, matched controls, feasibility, source-of-truth, visibility, audit, no-smuggling, and effect threshold. In particular, the record must contain at least two admissible macro interventions and a host-supplied comparison satisfying

\[
\Theta_{\mathrm{td}}(\delta_{\mathrm{td}})=\texttt{pass}.
\]

\end{claimproposition}

\begin{proof}

The claim classifier of Subsubsection \hyperref[sn:5.6.5]{5.6.5} returns \(\texttt{accepted}\) only when no earlier row applies, in particular not the row \(\texttt{rejected}\). For the tag \(\texttt{top\_down\_channel}\), admissibility of \(I\) (Subsubsection \hyperref[sn:3.4.1]{3.4.1}) requires its check rules to assign the rejection verdict whenever some gate of \(\mathsf{TDGate}(R)\) does not pass; if the tag is not admitted, the claim is \(\texttt{outside\_scope}\). Hence an accepted top-down channel claim has every gate passing: macro-record well-formedness, intervention, matched controls, feasibility, source, visibility, audit, no-smuggling and effect. The intervention gate requires at least two admissible macro interventions, and the effect gate is \(\Theta_{\mathrm{td}}(\delta_{\mathrm{td}})=\texttt{pass}\).

\end{proof}

\begin{claimproposition}{Top-down channel implies structural downward influence}

An accepted top-down channel claim carries a positive structural dependency,

\[
\mathsf{TopDownChannel}(R):\texttt{accepted}
\;\mathrel{\mathsf{dep}_{\texttt{suffices}}}\;
\mathsf{StructDown}(M,X).
\]

\end{claimproposition}

\begin{proof}

By the previous proposition, acceptance requires the macro gate \(G_{\mathrm{macro}}\) to pass, and that gate requires the field \(\mathsf{StructPath}\) of \(R\) to be a declared record from \(M\) to \(X\) of one of the five structural kinds. That record witnesses \(\mathsf{StructDown}(M,X)\). Hence the accepted channel claim suffices for the structural-downward dependency.

\end{proof}

\begin{claimproposition}{Structural downward influence is insufficient for top-down channel acceptance}

Structural downward influence does not imply top-down causal-channel acceptance:

\[
\mathsf{StructDown}(M,X)
\;\mathrel{\mathsf{dep}_{\texttt{insufficient}}}\;
\mathsf{TopDownChannel}(R):\texttt{accepted}.
\]

\end{claimproposition}

\begin{proof}

A structural downward path records only that some macro-side record constrains, indexes, or maps to substrate-side data. Top-down channel acceptance requires strictly more data: at least two admissible macro interventions, matched controls, source-of-truth response records, visibility, audit, no-smuggling, and a passed effect threshold. The following finite instance provides the structural path without that data. The host is \(\mathsf H_{\mathrm{fin}}\), with deterministic substrate states \(\{0,1\}\), and the judging instrument \(I\) is admissible, admits \(\mathsf H_{\mathrm{fin}}\) and the claim types \(\texttt{top\_down\_channel}\) and \(\texttt{dependency}\), and has the records of the judgment in its scope with every recorded level tag in \(\mathsf{Levels}_I\), so the instance lies in the class \(\Sigma\) of Subsubsection \hyperref[sn:4.6.1]{4.6.1}; \(M\) is a phase record of \(\mathcal T\) with the single value \(\phi_0\); \(X\) is the next-state record; and \(\mathsf{StructPath}\) is the context index by which \(\phi_0\) selects the update rule \(s\mapsto s\). Then \(\mathsf{StructDown}(M,X)\) holds. Let \(R\) carry this path, the one-element intervention family \(\mathsf{Interv}_M=\{\phi_0\}\), and passing records for every gate except \(G_{\mathrm{interv}}\) and \(G_{\mathrm{effect}}\); record the audited absence of \(\delta_{\mathrm{td}}\): with the single intervention \(\phi_0\) there is no pair \(m\neq m'\), so the comparison set \(\mathcal M\) defined after the third proposition is empty and \(G_{\mathrm{effect}}\) fails. Its intervention gate fails, since it requires two admissible macro interventions, so the check rules assign the rejection verdict and the claim \(\mathsf{TopDownChannel}(R)\) is not \(\texttt{accepted}\): it receives \(\texttt{rejected}\), or a status of higher priority. So structural downward influence is insufficient for channel acceptance.

\end{proof}

The acceptance hypothesis of the first two propositions is satisfiable. Take the host and instrument of the last proof with a phase record \(M\) attached to \(\mathcal T\) (a record of \(\mathcal T\) that exists before \(R\)), whose two values \(\phi_0,\phi_1\) select the update rules \(s\mapsto 0\) and \(s\mapsto 1\), and let \(\mathsf{StructPath}\) be the context index; give \(R\) a profile \(\lambda\) with \(\mathsf{AdmProfile}(\lambda,I,\mathcal T)\). Let \(\mathsf{Interv}_M=\{\phi_0,\phi_1\}\), each with a feasibility record whose audit passes; let \(\mathsf{Ctrl}\) give both interventions initial state \(0\) and the same background, support, package, profile and source policy; let \(\mathsf{Resp}(\phi_k)\) be the recorded next state \(k\), each carrying a \(\texttt{committed\_state}\) source record listed in the \(\mathsf{SourceRefs}\) of \(A\); let \(D\) be the output-inequality indicator and \(\Theta_{\mathrm{td}}\) the threshold \(D=1\); and let every used record be visible, every origin tag be \(\texttt{target}\) and every role tag be \(\texttt{other}\), with all six no-smuggling tests passing, and every audit check pass. All nine gates pass and \(\delta_{\mathrm{td}}=1\), so under an admissible instrument whose host, levels, scope, visibility, evidence and source policies admit this claim and the committed response sources, whose claim types include \(\texttt{top\_down\_channel}\), and whose check rules reject only when a gate fails, the claim \(\mathsf{TopDownChannel}(R)\) is \(\texttt{accepted}\).

A completion, lift, feasibility gate, or context index is a structural path. It becomes a top-down causal channel only when the intervention and matched-response gates pass. The required top-down gates are

\[
\mathsf{TDGate}(R)
=
(G_{\mathrm{macro}},G_{\mathrm{interv}},G_{\mathrm{matched}},G_{\mathrm{feas}},
G_{\mathrm{source}},G_{\mathrm{vis}},G_{\mathrm{audit}},G_{\mathrm{nosmuggle}},
G_{\mathrm{effect}}).
\]

The macro gate requires that \(M\) be a package/profile record rather than a post hoc label, the intervention gate at least two macro choices, the matched-control gate a common background, and the feasibility gate feasible interventions; the source, visibility, audit and no-smuggling gates apply the source-of-truth, no-overreading, audit and anti-smuggling discipline to the records used by the top-down claim; and the effect gate is the host-specific threshold check \(\Theta_{\mathrm{td}}\). Write \(\varphi_R=\mathsf{TopDownChannel}(R)\) and \(\mathcal L_R\) for the bridges recorded in \(L\) and \(V\). The pass conditions are as follows. \(G_{\mathrm{macro}}\) passes iff \(M\) is a record of \(\mathcal T\) or of \(\lambda\) whose identifier is not introduced by \(R\), and \(\mathsf{StructPath}\) is a record, not an absence marker, of one of the five structural kinds with source \(M\) and target \(X\). \(G_{\mathrm{feas}}\) passes iff each \(m\in\mathsf{Interv}_M\) has a feasibility record in the declared layer whose audit defect is \(\texttt{audit\_passes}\). \(G_{\mathrm{source}}\) passes iff each \(\mathsf{Resp}(m)\) carries a source-of-truth record \(R_m\) listed in the \(\mathsf{SourceRefs}\) of \(A\), and \(\delta_6^{\mathrm{source}}(R_m;I)\) is \(\texttt{committed}\), \(\texttt{fallback\_admitted}\) or \(\texttt{proxy\_admitted}\). \(G_{\mathrm{vis}}\) passes iff \(\delta_{\mathrm{vis}}(I,\varphi_R,\mathcal L_R)=\varnothing\) (Subsection \hyperref[sn:5.5]{5.5}). \(G_{\mathrm{audit}}\) passes iff \(\delta_6^{\mathrm{audit}}(A)\) is \(\texttt{audit\_passes}\). \(G_{\mathrm{nosmuggle}}\) passes iff the six violation tests of Subsubsection \hyperref[sn:5.4.2]{5.4.2} are empty on the audited origin and role tags of \(\mathsf{Use}(\varphi_R)\), using the visibility record of \(R\) and \(\delta_{\mathrm{vis}}(I,\varphi_R,\mathcal L_R)\); in the first test, an independent substrate-effect assertion used in place of the recorded response comparison is the forbidden premise. For the effect gate, let the values of \(D\) have a declared decidable total order, and let \(\mathcal M=\{D(\mathsf{Resp}(m),\mathsf{Resp}(m')):m\neq m'\in\mathsf{Interv}_M,\ \mathsf{Ctrl}(m)=\mathsf{Ctrl}(m')\}\). When \(\mathcal M\neq\varnothing\), set \(\delta_{\mathrm{td}}=\max\mathcal M\); otherwise record its audited absence and set \(G_{\mathrm{effect}}=\texttt{fail}\). \(G_{\mathrm{interv}}\) passes iff \(|\mathsf{Interv}_M|\ge2\); \(G_{\mathrm{matched}}\) passes iff \(\mathsf{Ctrl}\) assigns all interventions the same background, initial class, support, package/lens, profile and source-policy record; and \(G_{\mathrm{effect}}\) passes exactly when \(\mathcal M\neq\varnothing\) and \(\Theta_{\mathrm{td}}(\delta_{\mathrm{td}})=\texttt{pass}\), with \(\delta_{\mathrm{td}}\) recomputed from \(D\), \(\mathsf{Resp}\) and \(\mathsf{Ctrl}\).

The record is not assigned to a single primitive. Its macro package or representative path may use \(P_5\), its feasible intervention family may use \(P_2\), its layer/profile bridge may use \(P_4\), its response or transition-law comparison may use \(P_1\) when an induced operator is at issue, and its provenance, source, threshold, visibility, and no-smuggling checks are \(P_6\)-audited. \(P_3\) route data may enter the record when protocol or order matters, but it is not by itself a channel certificate.

\subsection{Instrument-Indexed Claims}\label{sn:4.8}

Instrument-indexed claims are the most general judgment family in the calculus. Every nontrivial assertion in the paper --- directed-cell, pair-observable, promotion, dependency, theorem, model-realization, or countermodel --- is interpreted as an instrument-indexed claim under an admissible instrument. This subsection fixes the claim form at the level of generality required throughout the paper, and records the per-claim status set used by the classifier of Subsubsection \hyperref[sn:5.6.5]{5.6.5}; the classifier rules and the no-overreading and same-level self-audit theorems are stated in Section \hyperref[sn:11]{11}.

\subsubsection{Claim Form}\label{sn:4.8.1}

The general instrument-indexed claim form is

\[
\Gamma\,;\;\mathcal T\,;\;I\;\vdash\;\varphi\;:\;\chi,
\]

introduced at schema-grade in Subsubsection \hyperref[sn:3.4.2]{3.4.2} and reused throughout the paper. Here \(\Gamma\) is a finite typed context, \(\mathcal T\) is an admissible theory package, \(I\) is an admissible instrument, \(\varphi\) is a claim with a finite use-set \(\mathsf{Use}(\varphi)\subseteq\mathsf{Cont}\), and \(\chi\in\mathrm{ClaimStatus}\) is the claim status assigned by the classifier of Subsubsection \hyperref[sn:5.6.5]{5.6.5}.

\begin{claimdefinition}{Instrument-indexed claim record}

A candidate claim record under \(I\) is the finite typed tuple \(\varphi(\mathcal T)=(\mathrm{id}_\varphi,\mathcal T,\mathsf{ClaimType},\mathsf{Content},\mathsf{LevelProfile},\mathsf{HostTag},\mathsf{EvidenceRefs},\Theta,A,V,\mathcal L,\mathcal N_\varphi,\mathsf{partialEvidence})\), whose fields are glossed in Subsubsection \hyperref[sn:11.1.2]{11.1.2}; its use-set is defined in Subsubsection \hyperref[sn:3.2.1]{3.2.1}, the reference graph of its circularity conditions in the Reference-graph paragraph of Subsubsection \hyperref[sn:5.6.5]{5.6.5}, the profile governing its audit's sources in Subsubsection \hyperref[sn:5.3.6]{5.3.6}, and \(\mathsf{ClaimLog}(I)\) in Subsubsection \hyperref[sn:11.5.1]{11.5.1}: it carries an identifier, the target package \(\mathcal T\), a finite declared claim-type tag, the claim content \(\varphi\) with its use-set, a level profile and host tag, the source records consulted (\(\mathsf{EvidenceRefs}\)), threshold and audit records, a visibility record with its visibility bridges, and nonclaim records. The surrounding \(\Gamma\) is judgment context, not a field of the tuple. The record is well formed exactly when it satisfies the conditions of Subsubsection \hyperref[sn:11.1.2]{11.1.2}. Membership of the claim-type tag in \(\mathsf{ClaimTypes}_I\) is decided by the claim classifier. The shorthand

\[
I\;\vdash\;\varphi(\mathcal T)
\]

abbreviates \(\Gamma;\mathcal T;I\vdash\varphi:\texttt{accepted}\) and is used in cross-references where the surrounding context is fixed; the shorthand never abbreviates instrument-free truth, and a statement of the form \(I\vdash\varphi(\mathcal T)\) does not transfer to a different instrument \(I'\) except through an explicit instrument-transfer bridge.

\end{claimdefinition}

\subsubsection{Claim Status}\label{sn:4.8.2}

The claim status set used for instrument-indexed claims is

\[
\begin{aligned}
\mathrm{ClaimStatus}\;=\;\{\,&\texttt{accepted},\;\texttt{rejected},\;\texttt{provisional},\;\texttt{blocked},\;\texttt{outside\_scope},\\
&\texttt{absent\_with\_record},\;\texttt{below\_threshold},\;\texttt{failed\_audit},\;\texttt{undefined\_circular}\,\}.
\end{aligned}
\]

The classifier defined by the claim-classifier table of Subsubsection \hyperref[sn:5.6.6]{5.6.6} assigns one of these to every well-formed claim record. The values record specific verdicts:

\begin{itemize}
\tightlist
\item
  \(\texttt{accepted}\): the claim is accepted by \(I\) under \(\Gamma,\mathcal T\) with admissible visibility, threshold, source, audit, and nonclaim data.
\item
  \(\texttt{rejected}\): the check rules of \(I\) assign the rejection verdict, for example because a recomputed defect contradicts the claim or a required gate of the claim type fails.
\item
  \(\texttt{provisional}\): \(\mathsf{partialEvidence}\) is true and the claim type lies in \(\mathsf{ProvisionalTypes}_I\).
\item
  \(\texttt{blocked}\): the claim is in scope and a used record is suppressed without an admissible bridge or has unknown visibility, an entry of \(\mathsf{EvidenceRefs}\) is not an admitted source record or fails \(\Theta_{\mathrm{source}}\), or its audit records a failed source or visibility check; it is also the final default, for example when \(\mathsf{partialEvidence}\) is true and the claim type is not in \(\mathsf{ProvisionalTypes}_I\).
\item
  \(\texttt{outside\_scope}\): the claim's host tag is not in \(\mathsf{Hosts}_I\), a level tag of its \(\mathsf{LevelProfile}\) is not in \(\mathsf{Levels}_I\), \(\mathsf{Use}(\varphi)\not\subseteq\mathsf{Scope}_I\), or its claim type is not in \(\mathsf{ClaimTypes}_I\).
\item
  \(\texttt{absent\_with\_record}\): a record that the content of \(\varphi\) refers to is absent from \(\Gamma\) or \(\mathcal T\), and the absence is audited. An absent evidence source is not of this kind: an audited absence marker in \(\mathsf{EvidenceRefs}\) does not resolve to a source record, so, unless the earlier \(\texttt{outside\_scope}\) row applies, the claim receives \(\texttt{blocked}\).
\item
  \(\texttt{below\_threshold}\): the audit's \(\mathsf{ThresholdCheck}\) fails (\(\texttt{threshold\_failure}\in\delta_6^{\mathrm{audit}}\)), that is, a threshold of the family tested by \(\mathsf{ThresholdCheck}\) (the claim's field \(\Theta\), Subsubsection \hyperref[sn:4.9.2]{4.9.2}) does not evaluate to \(\texttt{pass}\).
\item
  \(\texttt{failed\_audit}\): the audit record is missing, a recorded check result disagrees with its replay, or a nonclaim check fails (failure codes of the audit defect, Subsubsection \hyperref[sn:5.3.6]{5.3.6}).
\item
  \(\texttt{undefined\_circular}\): the claim depends on an unstratified circular support or audit loop, including the same-level self-audit case of Subsection \hyperref[sn:11.5]{11.5}.
\end{itemize}

The paper-facing claim strengths of Subsubsection \hyperref[sn:2.3.1]{2.3.1} (theorem-grade, proposition-grade, schema-grade, model-only, countermodel-grade, future-work, excluded) and the per-claim instrument statuses above are different layers. Claim strength records how a numbered statement is presented in the manuscript; claim status records how an instrument classifies a finite claim record. A theorem-grade or proposition-grade statement, under the hypotheses recorded in its statement, is represented by an accepted instrument-indexed claim; a schema-grade record is a definition or rule that is not adjudicated as a claim unless separately wrapped as one; a model-only statement is adjudicated only inside its declared model host; a countermodel-grade statement is adjudicated with its finite countermodel record attached; future-work and excluded statements are not asserted as accepted claims in this paper.

\subsection{Source, Provenance, and Audit Records (Reference)}\label{sn:4.9}

\emph{Used by:} The audit field \(A\) of the directed-cell records used in Theorems \hyperref[cl:thm:1]{1}--\hyperref[cl:thm:5]{5} is an audit record of this subsection.

Source and audit records are the working format for the calculus's reproducibility and scope discipline. They are not themselves proofs, and admissibility under an instrument is not the same as instrument-free truth. This subsection records the source-record and audit-record schemas used throughout the paper, with explicit separation between the audit record \(A\) and the inherited audit functional \(\mathcal A\) of Foundations I.

\subsubsection{Source Records}\label{sn:4.9.1}

\begin{claimdefinition}{Source record}

A source record is a finite typed tuple

\[
\mathrm{src}\;=\;(\,\mathrm{id}_{\mathrm{src}},\;\mathrm{type}_{\mathrm{src}},\;\mathrm{trace}_{\mathrm{src}},\;\mathrm{valid}_{\mathrm{src}}\,),
\]

with the following finite components.

\begin{itemize}
\item
  \(\mathrm{id}_{\mathrm{src}}\): a finite identifier under which the source is registered. Distinct source records carry distinct identifiers (Subsubsection \hyperref[sn:3.1.1]{3.1.1}).
\item
  \(\mathrm{type}_{\mathrm{src}}\): the source-type tag, drawn from the finite source-of-truth family

  \[
  \begin{aligned}
  \mathsf{SourceOfTruth}\;=\;\{\,&\texttt{committed\_state},\;\texttt{audited\_cell\_records},\;\texttt{independent\_pair\_witness},\\
  &\texttt{simulation\_trace},\;\texttt{ablation\_record},\;\texttt{fallback},\;\texttt{unknown},\;\texttt{contradictory}\,\},
  \end{aligned}
  \]

  The first five tags (\(\texttt{committed\_state}\), \(\texttt{audited\_cell\_records}\), \(\texttt{independent\_pair\_witness}\), \(\texttt{simulation\_trace}\), \(\texttt{ablation\_record}\)) are candidate source-of-truth tags whose force is determined by \(\mathsf{SourcePolicy}_I\) and the surrounding profile; \(\texttt{fallback}\) supports at most provisional pair-observable acceptance unless the record carries an audited source-upgrade bridge admitted by \(\mathsf{SourcePolicy}_I\) and \(\mathsf{BridgePolicy}_I\), and \(\texttt{unknown}\) or \(\texttt{contradictory}\) does not supply a full source-of-truth record.
\item
  \(\mathrm{trace}_{\mathrm{src}}\): a finite trace recording how the source was constructed and how it has been audited; entries are typed records consistent with \(\mathsf{AuditPolicy}_I\) for the surrounding instrument.
\item
  \(\mathrm{valid}_{\mathrm{src}}\): a finite validity record, recording whether the source is currently within its declared validity window or condition.
\end{itemize}

A claim that depends on a source records the identifier of that source in its use-set. A claim whose source is \(\texttt{unknown}\) or \(\texttt{contradictory}\) does not supply a full source-of-truth record under any instrument that requires one; the resulting judgment is routed by the relevant source policy, typically to \(\texttt{blocked}\) or \(\texttt{outside\_scope}\). Provisional pair support comes only from an admitted fallback source or an admitted proxy source (row 4 of Subsubsection \hyperref[sn:5.6.3]{5.6.3}), never from an unknown or contradictory source.

\end{claimdefinition}

\subsubsection{Audit Records}\label{sn:4.9.2}

\begin{claimdefinition}{Audit record}

An audit record is a finite typed tuple

\[
\begin{aligned}
A\;=\;(\,&\mathsf{ClaimRef},\;\mathsf{EvidenceRefs},\;\mathsf{CheckRules},\;\mathsf{CheckResults},\;\mathsf{SourceRefs},\\
&\mathsf{VisibilityCheck},\;\mathsf{ThresholdCheck},\;\mathsf{CircularityCheck},\;\mathsf{Nonclaims},\;\mathsf{Provenance}\,),
\end{aligned}
\]

with the following components.

\begin{itemize}
\tightlist
\item
  \(\mathsf{ClaimRef}\): a finite reference to the object or claim being audited. Targets include theory packages, directed cells, pair observables, promotion bridges, instrument-indexed claim records, dependency records, source records, visibility bridges (\(A_L\)), rotating-audit bridges, decorated squares (\(A_\Xi\)), absence markers, and branch-compatibility, attached-witness and feasibility records.
\item
  \(\mathsf{EvidenceRefs}\): the finite evidence references consulted by the audit. This field may name records of any type; it is distinct from the claim field \(\mathsf{EvidenceRefs}\) (Subsubsection \hyperref[sn:11.1.2]{11.1.2}), which must name source records. Only \(\mathsf{SourceRefs}\) is tested against the source policy.
\item
  \(\mathsf{CheckRules}\): the finite check rules, including the instrument and audit policy under which the audit is performed. An audit record without an instrumented check policy is malformed. An audit record filling the audit slot of a judgment under \(I\), or the audit \(A_L\) of a bridge registered under \(I\), names \(I\) in this field; so in the checks below the named instrument is the judging instrument.
\item
  \(\mathsf{CheckResults}\): the finite replay or check results produced by applying the check rules.
\item
  \(\mathsf{SourceRefs}\): the source records used by the audited claim or object.
\item
  \(\mathsf{VisibilityCheck}\): the finite visibility check under the named instrument. For the audit of a directed cell it is condition 8(a) of Subsubsection \hyperref[sn:4.3.3]{4.3.3}. For the audit of any other target \(X\) it passes exactly when every record named by \(\mathsf{EvidenceRefs}\) or \(\mathsf{SourceRefs}\) lies in \(\mathsf{Visible}_I\) or is bridged by an admissible bridge listed in the visibility field of \(X\). For targets \(X\) for which \(\mathsf{AdmVisBridge}(L',I,X)\) is not defined, the bridged alternative is unavailable. For the audit \(A_L\) of a visibility bridge \(L=(\mathrm{id}_L,c,c',\ldots)\), its suppressed endpoint \(c\) is exempt, but every other record named by \(\mathsf{EvidenceRefs}\) or \(\mathsf{SourceRefs}\) must be visible. The use-set of \(X\) is tested by its classifier or gate (\(\delta_{\mathrm{vis}}\), \(G_{\mathrm{vis}}\)), not by this check.
\item
  \(\mathsf{ThresholdCheck}\): the finite threshold check against the target's field \(\Theta\) when present, and otherwise against the finite threshold family prescribed by \(\mathsf{AuditPolicy}_I\), empty when none is prescribed. It passes exactly when \(\mathsf{Eval}_\theta\) (Subsubsection \hyperref[sn:5.2.1]{5.2.1}) returns \(\texttt{pass}\) for every threshold in that family; any other verdict (\(\texttt{fail}\), \(\texttt{not\_checked}\) or \(\texttt{outside\_scope}\)) makes it fail with code \(\texttt{threshold\_failure}\).
\item
  \(\mathsf{CircularityCheck}\): the finite check that the audit does not rest on the verdict it supports. Let \(X\) be the target of \(\mathsf{ClaimRef}\) and \(I\) the instrument named in \(\mathsf{CheckRules}\). The check passes exactly when no record reachable from \(\mathsf{EvidenceRefs}\) or \(\mathsf{SourceRefs}\) (the named records themselves included) by resolved references (reachability in the reference graph of Subsubsection \hyperref[sn:5.6.5]{5.6.5}, whose edges go from each record to every record named by one of its resolved reference fields) is a verdict reference \(\ulcorner X,I\urcorner\) naming the verdict of \(X\) under \(I\); its result is computed from these fields. A verdict reference to \(X\) under an instrument other than \(I\), in particular a higher one, does not by itself make the check fail.
\item
  \(\mathsf{Nonclaims}\): the nonclaim records required by the audited claim or object, namely those \(\mathsf{AuditPolicy}_I\) lists for the target's record family or claim type (none if it lists none) and those its schema names (Subsubsections \hyperref[sn:11.7.3]{11.7.3}, \hyperref[sn:13.1.1]{13.1.1}). The check passes exactly when each required record is present in this field; its failure is the code \(\texttt{nonclaim\_failure}\).
\item
  \(\mathsf{Provenance}\): the finite provenance trace of the audit record; every record it names is also named by \(\mathsf{EvidenceRefs}\) or \(\mathsf{SourceRefs}\), so the visibility and circularity checks cover it.
\end{itemize}

The audit passes only when its target resolves, its evidence resolves, the source-of-truth policy passes, visibility passes, thresholds pass, the circularity check passes, and the required nonclaims are present. The audit record family \(\mathsf{Audit}\) of \(\mathbf{BirdInt}^{\mathrm{aud}}_{\mathrm{fin}}\) is the finite collection of all such records used by judgments in the paper.

The audit record \(A\) is distinct from the inherited audit functional \(\mathcal A\) that lives inside a theory package \(\mathcal T=(Z,f,\Sigma_f,E,\mathcal A)\): \(\mathcal A\) is the package-internal audit functional, while \(A\) is a finite audit record, certificate, trace, or evidence object. Compatibility with \(\mathcal A\) is a separate typed condition when the inherited package audit functional is used. The two are not interchangeable, and a theorem that uses one does not by that fact use the other.

A drive certificate is a specialization of an audit record whose check data records the drive face \(\mathrm{P6}_{\mathrm{drive}}\) --- a nonzero cycle affinity, a nontrivial cohomology class, or another declared drive bridge --- rather than a generic \(P_6\) audit. A generic admissible audit record does not by itself supply a drive certificate; Subsubsection \hyperref[sn:5.3.6]{5.3.6} records the drive-certificate schema and Section \hyperref[sn:8]{8} states the no-drive theorems that constrain it.

\end{claimdefinition}

\subsubsection{Evidence Hygiene}\label{sn:4.9.3}

The discipline of the paper for using source and audit records is:

\begin{itemize}
\tightlist
\item
  a source record fixes which source a claim depends on, so that the claim's reliance on a real, fallback, dirty, proxy, unknown, or contradictory source is identifiable at the classifier step;
\item
  an audit record fixes which checks a claim has passed or failed under a specific instrument, so that the claim is reproducible by anyone with access to the same records and the same audit policy;
\item
  neither record asserts instrument-free truth. An audited record may be exact under one instrument and suppressed or outside-scope under another, and the instrument-transfer rule of Subsection \hyperref[sn:11.9]{11.9} requires an admissible instrument-transfer bridge and a fresh verdict of the target instrument's classifier before a claim accepted under one instrument is classified under another; the no-overreading theorem of Section \hyperref[sn:11]{11} separately ensures that an accepted claim uses suppressed content only through admissible bridges.
\end{itemize}

NC-15 of Subsection \hyperref[sn:15.2]{15.2} records the corresponding nonclaim: \(A\) passing under \(I\), \(\mathcal T\), and the surrounding scope yields acceptance under that scope, not instrument-free truth. Audits and source records record the hygiene of the calculus, not its conclusions; the conclusions are theorems whose proofs use the records together with the typed witness, update, profile, level, visibility, and threshold data of the surrounding judgments.

\subsection{Domain Instance Admissibility (Reference)}\label{sn:4.10}

\emph{Used by:} the theorems of Sections \hyperref[sn:5]{5} through \hyperref[sn:14]{14}, each of which assumes an admissible instance.

This subsection records the admissibility conditions on a full domain instance \(\mathfrak D\in\mathbf{BirdInt}^{\mathrm{aud}}_{\mathrm{fin}}\). Sections \hyperref[sn:5]{5} through \hyperref[sn:14]{14} work with admissible instances throughout, and the hypothesis ``let \(\mathfrak D\) be admissible'' is implicit in every universally stated theorem.

\subsubsection{Admissible Instance}\label{sn:4.10.1}

A full domain instance \(\mathfrak D\in\mathbf{BirdInt}^{\mathrm{aud}}_{\mathrm{fin}}\) is admissible, written

\[
\mathsf{AdmDomain}(\mathfrak D),
\]

when every component of the instance is finite and every record in every component meets its admissibility schema:

\begin{itemize}
\tightlist
\item
  \(\mathbb P=\{P_1,\ldots,P_6\}\) is the fixed finite primitive set of Subsubsection \hyperref[sn:2.2.1]{2.2.1}.
\item
  \(\mathsf{Lev}=\{\mathsf{Beh},\mathsf{Ver},\mathsf{Str}\}\) is the fixed finite level set of Subsubsection \hyperref[sn:2.2.3]{2.2.3}.
\item
  \(\mathsf{Host}\) is a finite subfamily of the host inventory of Subsubsection \hyperref[sn:2.1.2]{2.1.2}; only hosts named in \(\mathsf{Host}\) may appear on the host tag of any record in the instance.
\item
  \(\mathsf{Cont}\) is the finite content universe of the instance (Subsection \hyperref[sn:3.2]{3.2}): finite typed records, each with a host tag from \(\mathsf{Host}\), a level tag from \(\mathsf{Lev}\), an instrument-relative visibility class drawn from Subsection \hyperref[sn:3.5]{3.5}, and a profile attachment when required.
\item
  All references among content records resolve, or else the missing reference is represented by an explicit absence record (a typed, audited absence record, Subsubsection \hyperref[sn:3.7.3]{3.7.3}).
\item
  \(\mathsf{Instr}\) is a finite family of admissible instruments, each meeting the conditions of Subsubsection \hyperref[sn:3.4.1]{3.4.1}.
\item
  \(\mathsf{Pkg}\) is a finite family of admissible theory packages, each meeting the conditions of Subsubsection \hyperref[sn:3.3.3]{3.3.3}.
\item
  \(\mathsf{Prof}\) is a finite family of finite typed profile records in the schema of Subsubsection \hyperref[sn:3.6.1]{3.6.1}; \(\mathsf{AdmProfile}(\lambda,I,\mathcal T)\) is relative to the instrument and package of the judgment carrying \(\lambda\) and is checked with that judgment (step 4 of Subsubsection \hyperref[sn:6.3.1]{6.3.1}), not as a condition on the instance.
\item
  \(\mathsf{Wit}\) and \(\mathsf{Upd}\) are the witness and update maps of Subsubsections \hyperref[sn:3.7.1]{3.7.1}--\hyperref[sn:3.7.2]{3.7.2}; their values on each \(P_i\) are finite typed families, with the present-or-audited-absence slot rule of Subsubsection \hyperref[sn:3.7.3]{3.7.3} enforced on every directed-cell record.
\item
  \(\mathsf{Def}\) is a finite family of admissible defect records in the schema of Subsubsection \hyperref[sn:5.2.1]{5.2.1}.
\item
  \(\mathsf{Judg}\) is the finite family of judgment forms introduced in Subsections \hyperref[sn:4.2]{4.2} through \hyperref[sn:4.8]{4.8} (role-channel, directed-cell, pair-observable, promotion, dependency, instrument-indexed claim), together with the model-realization judgment of Subsubsection \hyperref[sn:13.1.1]{13.1.1}.
\item
  \(\mathsf{Status}\) is the finite family of the seven status sets of Subsection \hyperref[sn:5.1]{5.1}: those of the role-channel, directed-cell, pair-observable, promotion and claim judgments, and those of gate and square records.
\item
  Threshold records are finite and typed in the threshold schema of Subsection \hyperref[sn:5.2]{5.2}.
\item
  The evaluation dependency graph of Subsubsection \hyperref[sn:5.3.6]{5.3.6} is acyclic.
\item
  \(\mathsf{Gate}\) is the finite family of gate specifications used in Sections \hyperref[sn:10]{10} and \hyperref[sn:11]{11}, with the gate status sets \(\mathrm{GateStatus}\), \(\mathrm{StrictGateStatus}\) and \(\mathrm{LocGlobGateStatus}\) from Subsubsection \hyperref[sn:5.1.6]{5.1.6}; promotion gates are explicit whenever a promotion judgment is claimed.
\item
  \(\mathsf{Dep}\) is a finite family of admissible dependency records in the schema of Subsection \hyperref[sn:4.6]{4.6}.
\item
  \(\mathsf{Audit}\) and \(\mathsf{Source}\) are finite families of admissible audit and source records in the schemas of Subsubsections \hyperref[sn:4.9.2]{4.9.2} and \hyperref[sn:4.9.1]{4.9.1}.
\item
  \(\mathsf{Vis}\) is the finite family of visibility maps and visibility bridges of Subsection \hyperref[sn:3.5]{3.5}.
\item
  \(\mathsf{Nonclaim}\) is a finite family of admissible nonclaim records, attached either to instruments (as \(\mathcal N_I\)) or to individual judgments.
\item
  \(\mathsf{Real}\) is a finite family of admissible model-realization modules; admissibility of a model-realization module is the schema fixed in Section \hyperref[sn:13]{13}.
\end{itemize}

A full instance that satisfies all of the above is the working object of every theorem in Sections \hyperref[sn:5]{5} through \hyperref[sn:11]{11} and of every model-only realization statement in Section \hyperref[sn:13]{13}. Admissibility is a finite check on finite records: it composes the per-component admissibility checks of Sections \hyperref[sn:3]{3} and \hyperref[sn:4]{4} into a single instance-level condition and feeds the well-formedness, classifier, gate, and audit machinery of Sections \hyperref[sn:5]{5} through \hyperref[sn:11]{11}.

\subsubsection{Exclusion Conditions}\label{sn:4.10.2}

A candidate instance is rejected as outside the scope of \(\mathbf{BirdInt}^{\mathrm{aud}}_{\mathrm{fin}}\) when any of the following holds:

\begin{itemize}
\tightlist
\item
  a component of the tuple is missing or infinite;
\item
  a host tag on any record points to a host outside \(\mathsf{Host}\) (in particular, outside the host inventory of Subsubsection \hyperref[sn:2.1.2]{2.1.2});
\item
  a level tag on any record points outside \(\mathsf{Lev}\);
\item
  a theory package fails \(\mathsf{AdmPkg}\) from Subsubsection \hyperref[sn:3.3.3]{3.3.3}, including missing inherited Foundations I data \((Z,f,\Sigma_f,E,\mathcal A)\) or missing Foundations III attached records;
\item
  an instrument fails the instrument admissibility conditions of Subsubsection \hyperref[sn:3.4.1]{3.4.1};
\item
  a profile record is not a finite typed record of the schema of Subsubsection \hyperref[sn:3.6.1]{3.6.1};
\item
  a present witness or update has the wrong primitive signature, or an absence marker lacks its required typed, passing audit;
\item
  a directed-cell, pair-observable, promotion, dependency, or claim record has unresolved references, missing required fields, or status-family collapse where the record uses a status from one family in place of another;
\item
  an audit record is malformed, has unresolved evidence/source references, or has no \(\mathsf{Nonclaims}\) field;
\item
  a source record is malformed or has unresolved references;
\item
  a model-realization module lacks a declared fragment \(\mathcal S\), or its nonclaim register lacks the model-realization nonclaims of Subsection \hyperref[sn:13.6]{13.6} (clause 5 of Subsubsection \hyperref[sn:13.1.1]{13.1.1}; NC-19 of Subsection \hyperref[sn:15.2]{15.2}).
\end{itemize}

The exclusions concern the form of a record, not the outcome of its checks. A well-formed record whose visibility, threshold, audit or source check fails stays in the instance; where the applicable classifier has a matching condition, it assigns the status prescribed by Subsection \hyperref[sn:5.6]{5.6}. In particular, claim failures can receive \(\texttt{blocked}\), \(\texttt{below\_threshold}\), \(\texttt{undefined\_circular}\) or \(\texttt{failed\_audit}\), and every well-formed cell receives a cell status from one of the rows \(\texttt{outside\_scope}\) through \(\texttt{action}\) (Theorem \hyperref[cl:thm:4]{4}). Role-channel records and theory packages are the exceptions: for in-scope role evidence, visibility and audit passage are well-formedness conditions (Subsubsection \hyperref[sn:5.6.1]{5.6.1}), and a package is admissible only when its audit functional passes every entry of its audit record family (Subsubsection \hyperref[sn:3.3.3]{3.3.3}), so a package with a failing audit entry falls under the fourth exclusion above. A candidate directed cell over \(\mathfrak D\) (Theorem \hyperref[cl:thm:3]{3}) is a tuple of the form of Subsubsection \hyperref[sn:4.3.1]{4.3.1} whose fields are records of \(\mathfrak D\) or audited absence markers; it is not assumed to be a record of \(\mathfrak D\), and the exclusions above do not apply to it. Because \(\mathsf{AdmDomain}(\mathfrak D)\) requires the \(\mathsf{ClaimRef}\) of every audit record of \(\mathfrak D\) to resolve in \(\mathsf{Cont}\), a candidate whose audit slot holds a present record of \(\mathfrak D\) meets condition 8(a) only if it is itself a record of \(\mathfrak D\); a candidate outside \(\mathfrak D\) can be well formed only with an audited absence marker in its audit slot. A malformed candidate receives no status and does not make \(\mathfrak D\) inadmissible.

A claim, judgment, or theorem stated against a rejected candidate instance is not interpreted in \(\mathbf{BirdInt}^{\mathrm{aud}}_{\mathrm{fin}}\) and is excluded from the formal scope of this paper. Sections \hyperref[sn:5]{5} through \hyperref[sn:14]{14} work with admissible instances throughout, and a hypothesis ``let \(\mathfrak D\) be admissible'' is implicit in every theorem statement.
 \section{Status and Defect Semantics}\label{sn:5}

This section gathers the full status families used by the calculus and the defect record schemas that feed the defect-to-status rules consumed by Sections \hyperref[sn:6]{6} through \hyperref[sn:11]{11}. The role-channel, directed-cell, pair-observable, promotion, and claim status families were introduced per judgment in Section \hyperref[sn:4]{4}; the gate and square families were referenced there. This section consolidates all seven families, fixes the defect record schema, lists the primitive and promotion-specific defect families used by the paper, and states the defect-to-status rules that drive the classifiers. The definitions needed for Theorems \hyperref[cl:thm:1]{1}--\hyperref[cl:thm:4]{4} are listed in the reader's map (Subsection \hyperref[sn:1.4]{1.4}); the rest of this section is reference material.

\subsection{Status-Family Separation}\label{sn:5.1}

\begin{claimdefinition}{Status families}

The seven status families used in this paper are

\[
\begin{aligned}&\mathrm{RoleStatus},\quad\mathrm{CellStatus},\quad\mathrm{PairStatus},\quad\mathrm{PromotionStatus},\\&\mathrm{ClaimStatus},\quad\mathrm{GateStatus},\quad\mathrm{SqStatus}.\end{aligned}
\]

They are kept apart at the schema level: a value in one family is not by definition a value in another, and a status assigned to a judgment in one family does not by itself fix a status in another. Theorem \hyperref[cl:thm:5]{5} (Subsection \hyperref[sn:6.5]{6.5}) proves five non-implications between these families: role to cell, cell to pair, strict promotion to cell, strict promotion to same-level soundness (a consequence of Theorem \hyperref[cl:thm:24]{24}, since that claim is never accepted), and strict promotion to accepted lifted soundness under a higher instrument. NC-4 of Subsection \hyperref[sn:15.2]{15.2} records the general family separation.

\end{claimdefinition}

The table below sets the seven families side by side. Each row is one value name; an entry gives the value's place in the priority order of that family's classifier (1 is checked first), \(\bullet\) marks a value of a family whose rules fix no priority order here, and -- marks a value the family does not have. The Gate column includes the refined values of \(\mathrm{StrictGateStatus}\) (Subsubsection \hyperref[sn:5.1.6]{5.1.6}); the values \(\texttt{glue\_pass}\), \(\texttt{glue\_fail}\) and \(\texttt{local\_only}\) of \(\mathrm{LocGlobGateStatus}\) are not marked in that column, and the \(\texttt{local\_only}\) row records the promotion status only. The same name can mean different things in different families: \(\texttt{accepted}\) is a verdict on a claim in \(\mathrm{ClaimStatus}\) and a promotion whose core gates pass in \(\mathrm{PromotionStatus}\).

\begingroup\footnotesize\setlength{\tabcolsep}{3pt}
\begin{longtable}{@{}lccccccc@{}}
\toprule
Value & Role & Cell & Pair & Prom. & Claim & Gate & Square \\
\midrule
\endfirsthead
\toprule
Value & Role & Cell & Pair & Prom. & Claim & Gate & Square \\
\midrule
\endhead
\bottomrule
\endfoot
$\texttt{outside\_scope}$ & -- & 1 & -- & 1 & 1 & $\bullet$ & 1 \\
$\texttt{inapplicable\_outside\_scope}$ & 1 & -- & -- & -- & -- & -- & -- \\
$\texttt{absent}$ & -- & 2 & -- & -- & -- & -- & -- \\
$\texttt{absent\_with\_record}$ & 2 & -- & -- & -- & 3 & -- & -- \\
$\texttt{blocked}$ & -- & 3, 10 (final) & 1, 5 (final) & -- & 2, 10 (final) & -- & 2, 9 (final) \\
$\texttt{failed\_audit}$ & -- & -- & -- & 2 & 4 & -- & 3 \\
$\texttt{failed\_no\_smuggling}$ & -- & -- & -- & 3 & -- & -- & -- \\
$\texttt{failed\_descent}$ & -- & -- & -- & 4 & -- & -- & -- \\
$\texttt{failed\_stability}$ & -- & -- & -- & 5 & -- & -- & -- \\
$\texttt{undefined\_circular}$ & -- & 4 & -- & -- & 5 & -- & 4 \\
$\texttt{collapsed}$ & -- & 5 & -- & -- & -- & -- & -- \\
$\texttt{collapsed\_boundary\_case}$ & 3 & -- & -- & -- & -- & -- & -- \\
$\texttt{below\_threshold}$ & 4 & 6 & -- & -- & 6 & -- & -- \\
$\texttt{globally\_obstructed}$ & -- & -- & -- & 6 & -- & -- & -- \\
$\texttt{local\_only}$ & -- & -- & -- & 7 & -- & -- & -- \\
$\texttt{obstructed}$ & -- & -- & -- & -- & -- & -- & 5 \\
$\texttt{rejected}$ & -- & -- & -- & -- & 7 & -- & -- \\
$\texttt{candidate}$ & -- & -- & -- & 8, 12 (final) & -- & -- & -- \\
$\texttt{trivial}$ & -- & 7 & -- & -- & -- & -- & -- \\
$\texttt{implicit}$ & -- & 8 & -- & -- & -- & -- & -- \\
$\texttt{fail}$ & -- & -- & -- & -- & -- & $\bullet$ & -- \\
$\texttt{not\_checked}$ & -- & -- & -- & -- & -- & $\bullet$ & -- \\
$\texttt{not\_required}$ & -- & -- & -- & -- & -- & $\bullet$ & -- \\
$\texttt{provisional}$ & -- & -- & -- & -- & 8 & -- & 7 \\
$\texttt{real\_provisional}$ & -- & -- & 4 & -- & -- & -- & -- \\
$\texttt{real\_duplicated}$ & -- & -- & 2 & -- & -- & -- & -- \\
$\texttt{lax}$ & -- & -- & -- & -- & -- & -- & 6 \\
$\texttt{action}$ & -- & 9 & -- & -- & -- & -- & -- \\
$\texttt{active\_projection}$ & 5 & -- & -- & -- & -- & -- & -- \\
$\texttt{real}$ & -- & -- & 3 & -- & -- & -- & -- \\
$\texttt{strict}$ & -- & -- & -- & 9 & -- & -- & -- \\
$\texttt{non\_strict}$ & -- & -- & -- & 10 & -- & -- & -- \\
$\texttt{accepted}$ & -- & -- & -- & 11 & 9 & -- & -- \\
$\texttt{exact}$ & -- & -- & -- & -- & -- & -- & 8 \\
$\texttt{pass}$ & -- & -- & -- & -- & -- & $\bullet$ & -- \\
$\texttt{strict\_pass}$ & -- & -- & -- & -- & -- & $\bullet$ & -- \\
$\texttt{non\_strict\_pass}$ & -- & -- & -- & -- & -- & $\bullet$ & -- \\
$\texttt{bad\_audit}$ & -- & -- & -- & -- & -- & $\bullet$ & -- \\
\end{longtable}
\endgroup

\subsubsection{Role-Channel Statuses}\label{sn:5.1.1}

The role-channel status set \(\mathrm{RoleStatus}\) is the inherited Foundations II set of five values, displayed with their meanings in Subsubsection \hyperref[sn:4.2.2]{4.2.2}.

\subsubsection{Directed-Cell Statuses}\label{sn:5.1.2}

The directed-cell status set \(\mathrm{CellStatus}\), with nine values from \(\texttt{action}\) to \(\texttt{outside\_scope}\), is displayed in Subsubsection \hyperref[sn:4.3.4]{4.3.4}.

Its classifier priority order is displayed, with the per-status conditions that the classifier checks, in Subsubsection \hyperref[sn:5.6.2]{5.6.2}.

\subsubsection{Pair-Observable Statuses}\label{sn:5.1.3}

The pair-observable status set \(\mathrm{PairStatus}\) is displayed, with its meanings and source-of-truth conditions, in Subsection \hyperref[sn:4.4]{4.4}.

\subsubsection{Promotion Statuses}\label{sn:5.1.4}

The promotion status set \(\mathrm{PromotionStatus}\) is displayed in Subsubsection \hyperref[sn:4.5.3]{4.5.3}.

The values \(\texttt{strict}\) and \(\texttt{non\_strict}\) are accepted-status refinements; the failure-tagged values record specific gate failures, as in Subsubsection \hyperref[sn:4.5.3]{4.5.3}.

\subsubsection{Claim Statuses}\label{sn:5.1.5}

The instrument-indexed claim status set \(\mathrm{ClaimStatus}\) is displayed in Subsubsection \hyperref[sn:4.8.2]{4.8.2}.

Subsubsection \hyperref[sn:4.8.2]{4.8.2} records the meanings; Subsubsection \hyperref[sn:5.6.5]{5.6.5} specifies the classifier, and Section \hyperref[sn:11]{11} proves the no-overreading and same-level self-audit theorems against it.

\subsubsection{Gate Statuses}\label{sn:5.1.6}

A promotion or claim gate has status drawn from the finite set

\[
\mathrm{GateStatus}\;=\;\{\,\texttt{pass},\;\texttt{fail},\;\texttt{not\_required},\;\texttt{not\_checked},\;\texttt{outside\_scope}\,\}.
\]

A strictness gate has the refined set

\[
\begin{aligned}
\mathrm{StrictGateStatus}\;=\;\{\,&\texttt{strict\_pass},\;\texttt{non\_strict\_pass},\;\texttt{not\_required},\\
&\texttt{not\_checked},\;\texttt{bad\_audit},\;\texttt{outside\_scope}\,\}.
\end{aligned}
\]

The local-to-global gate has the refined set

\[
\begin{aligned}\mathrm{LocGlobGateStatus}\;=\;\{\,&\texttt{glue\_pass},\;\texttt{local\_only},\;\texttt{glue\_fail},\\&\texttt{not\_required},\;\texttt{not\_checked},\;\texttt{outside\_scope}\,\}:\end{aligned}
\]

\(\texttt{glue\_fail}\) when a recorded gluing record has \(\delta_4^{\mathrm{glue}}\neq\varnothing\); \(\texttt{local\_only}\) when every local package on the recorded cover is admissible and no gluing record is present; and \(\texttt{glue\_pass}\) when a gluing record with \(\delta_4^{\mathrm{glue}}=\varnothing\) is present. \(G_{\mathrm{suff}}\) checks the recorded value.

The core promotion gates used by the paper are

\[
\mathsf{CoreGates}\;=\;\{\,G_{\mathrm{suff}},\;G_{\mathrm{desc}},\;G_{\mathrm{stab}},\;G_{\mathrm{ctrl}},\;G_{\mathrm{nosmuggle}},\;G_{\mathrm{vis}},\;G_{\mathrm{audit}}\,\},
\]

with the additional strictness gate \(G_{\mathrm{strict}}\) and the optional local-to-global gate \(G_{\mathrm{locglob}}\). Their roles are fixed in Section \hyperref[sn:10]{10}.

\subsubsection{Square Statuses}\label{sn:5.1.7}

The high-structure module of Section \hyperref[sn:14]{14} uses the square status set

\[
\begin{aligned}
\mathrm{SqStatus}\;=\;\{\,&\texttt{exact},\;\texttt{lax},\;\texttt{obstructed},\;\texttt{blocked},\;\texttt{outside\_scope},\\
&\texttt{failed\_audit},\;\texttt{undefined\_circular},\;\texttt{provisional}\,\}.
\end{aligned}
\]

The values have the meanings of an exact square (\(\texttt{exact}\)), an ordered or approximate square comparison (\(\texttt{lax}\)), a square in scope but with nonzero defect (\(\texttt{obstructed}\)), and the standard outside, failed-audit, circular, and provisional cases. Subsubsection \hyperref[sn:5.6.6]{5.6.6} gives the square classifier.

\subsection{General Defect Record (Reference)}\label{sn:5.2}

\emph{Used by:} The defect field \(\delta\) of the directed-cell records used in Theorems \hyperref[cl:thm:1]{1}--\hyperref[cl:thm:5]{5} uses this record form.

\subsubsection{Record Form}\label{sn:5.2.1}

\begin{claimdefinition}{Typed defect record}

A typed defect record is the finite tuple

\[
\mathsf{Def}_i^{\lambda}(\mathcal T,I)\;=\;(\,\delta_i,\;\Omega_i,\;\Theta_i,\;\mathrm{polarity}_i,\;W_i,\;A_i,\;V_i,\;\sigma_i\,)
\]

attached to a primitive index \(i\in\{1,\ldots,6\}\) (or to a non-primitive defect family in Subsections \hyperref[sn:5.4]{5.4} and \hyperref[sn:5.5]{5.5}), under profile \(\lambda\), theory package \(\mathcal T\), and instrument \(I\). The fields are finite and have the following roles:

\begin{itemize}
\tightlist
\item
  \(\delta_i\): the defect value, a finite typed value drawn from the value type of \(\Omega_i\).
\item
  \(\Omega_i\): the value object or defect type, naming the kind of value \(\delta_i\) takes (numeric, Boolean, categorical, set-valued, cohomological, audit-valued).
\item
  \(\Theta_i\): the threshold or pass condition against which \(\delta_i\) is evaluated.
\item
  \(\mathrm{polarity}_i\): the polarity tag, naming whether the defect is an obstruction, certificate, gap, residual, capacity loss, or other typed polarity drawn from a finite family fixed by the surrounding family.
\item
  \(W_i\): the witness, certificate, or counterexample supporting \(\delta_i\).
\item
  \(A_i\): the audit record of the defect, in the schema of Subsubsection \hyperref[sn:4.9.2]{4.9.2}.
\item
  \(V_i\): the visibility record of the defect under \(I\), in the schema of Subsection \hyperref[sn:3.5]{3.5}.
\item
  \(\sigma_i\): the induced status or status annotation, drawn from the status family appropriate to the surrounding judgment.
\end{itemize}

The general threshold rule for evaluating a defect is

\[
\mathsf{Eval}_{\Theta_i}(\delta_i)\;\in\;\{\,\texttt{pass},\;\texttt{fail},\;\texttt{not\_checked},\;\texttt{outside\_scope}\,\},
\]

producing one of four verdicts. The verdict feeds the defect-to-status rules of Subsection \hyperref[sn:5.6]{5.6}.

\end{claimdefinition}

\subsubsection{Threshold Types}\label{sn:5.2.2}

The threshold field \(\Theta_i\) is drawn from the finite family

\[
\Theta_i\;\in\;\{\,\Theta_{=0},\;\Theta_{\neq0},\;\Theta_{\le\varepsilon},\;\Theta_{\ge\varepsilon},\;\Theta_{\in S},\;\Theta_{\notin S},\;\Theta_{\mathrm{pass}},\;\Theta_{\mathrm{source}},\;\Theta_{\mathrm{bridge}},\;\Theta_{\mathrm{audit}},\;\Theta_{\mathrm{vis}}\,\}.
\]

The eleven types have the following meanings, used by the per-defect rules of Subsections \hyperref[sn:5.3]{5.3} through \hyperref[sn:5.5]{5.5}:

\begin{longtable}[]{@{}ll@{}}
\toprule
Threshold & Meaning\tabularnewline
\midrule
\endhead
\(\Theta_{=0}\) & exact defect vanishes\tabularnewline
\(\Theta_{\neq0}\) & nontrivial witness exists\tabularnewline
\(\Theta_{\le\varepsilon}\) & approximate defect tolerated below \(\varepsilon\)\tabularnewline
\(\Theta_{\ge\varepsilon}\) & activation passes a lower bound \(\varepsilon\)\tabularnewline
\(\Theta_{\in S}\) & representability or membership condition holds\tabularnewline
\(\Theta_{\notin S}\) & exclusion or nonrepresentability holds\tabularnewline
\(\Theta_{\mathrm{pass}}\) & a Boolean audit or gate pass\tabularnewline
\(\Theta_{\mathrm{source}}\) & the source-of-truth requirement is met\tabularnewline
\(\Theta_{\mathrm{bridge}}\) & a required bridge exists and is strong enough\tabularnewline
\(\Theta_{\mathrm{audit}}\) & the audit policy passes\tabularnewline
\(\Theta_{\mathrm{vis}}\) & the visibility check passes\tabularnewline
\bottomrule
\end{longtable}

Each threshold type has a fixed evaluation rule on values of the appropriate kind: \(\Theta_{=0}\) checks \(\delta=0\) for numeric defects and \(\delta=\varnothing\) for set-valued defects, \(\Theta_{\le\varepsilon}\) checks \(|\delta|\le\varepsilon\), \(\Theta_{\neq0}\) checks \(\delta\neq0\) (\(\delta\neq\varnothing\) for a set-valued defect), \(\Theta_{\ge\varepsilon}\) checks \(\delta\ge\varepsilon\), \(\Theta_{\in S}\) and \(\Theta_{\notin S}\) check \(\delta\in S\) and \(\delta\notin S\), \(\Theta_{\mathrm{pass}}\) checks that the recorded Boolean is true, \(\Theta_{\mathrm{source}}\) checks \(\delta_6^{\mathrm{source}}\in\{\texttt{committed},\texttt{fallback\_admitted},\texttt{proxy\_admitted}\}\), \(\Theta_{\mathrm{bridge}}\) checks \(\delta_{\mathrm{reqbridge}}=\varnothing\), \(\Theta_{\mathrm{audit}}\) checks that the evaluated audit record \(A\) has its \(\mathsf{CheckRules}\) field drawn from \(\mathsf{CheckRules}_I\) and \(\delta_6^{\mathrm{audit}}(A)=\texttt{audit\_passes}\), and \(\Theta_{\mathrm{vis}}\) checks \(\delta_{\mathrm{vis}}=\varnothing\). \(\mathsf{Eval}\) tests scope on the record supplying the evaluated datum, using its host, level and content membership as in Subsubsection \hyperref[sn:3.4.1]{3.4.1}. It returns \(\texttt{outside\_scope}\) when that record is out of scope; otherwise it returns \(\texttt{not\_checked}\) when an audited absence marker replaces the datum, and \(\texttt{pass}\) or \(\texttt{fail}\) by the rules above when the datum is present. The threshold types are the finite vocabulary used throughout Sections \hyperref[sn:5]{5} through \hyperref[sn:11]{11}; no theorem in the paper introduces a threshold type outside this list without explicitly extending the family. The named thresholds of later sections are instances: \(\Theta_1^{\mathrm{desc}}\), \(\Theta_5^{\mathrm{idem}}\), \(\Theta_{\mathrm{nosmuggle}}\), \(\Theta_{\mathrm{nonstrict}}\) and \(\Theta_2^{\mathrm{forest}}\) (on the datum \(\beta_1(G')\)) of \(\Theta_{=0}\); \(\Theta_{\mathrm{strict}}\) of \(\Theta_{\neq0}\); \(\Theta_{\mathrm{rot}}\) of \(\Theta_{\in S}\); and \(\Theta_1^{\mathrm{closure}}\) of \(\Theta_{=0}\) or \(\Theta_{\le\varepsilon}\).

\subsection{Primitive Defect Families (Reference)}\label{sn:5.3}

\emph{Used by:} Theorems \hyperref[cl:thm:1]{1}--\hyperref[cl:thm:5]{5} use the package-collapse and audit defects of Subsubsections \hyperref[sn:5.3.5]{5.3.5}--\hyperref[sn:5.3.6]{5.3.6}; Theorems \hyperref[cl:thm:7]{7} and \hyperref[cl:thm:10]{10} state their results with defects of these families.

For each primitive \(P_i\), the calculus declares a finite family of defect record types. Each entry below uses a named defect of the calculus. Full constructions and proofs that depend on these defects appear in Sections \hyperref[sn:6]{6}, \hyperref[sn:7]{7}, and \hyperref[sn:8]{8}.

\subsubsection{\texorpdfstring{\(P_1\): Descent, Closure, Rewrite Defects}{P\_1: Descent, Closure, Rewrite Defects}}\label{sn:5.3.1}

The \(P_1\) defect family contains three named entries.

\begin{itemize}
\item
  \textbf{Exact quotient descent defect.} Given a finite quotient \(q:X\to Y\) and a finite map \(F:X\to X\), the split-pair set is

  \[
  \delta_1^{\mathrm{split}}(q,F)\;=\;\{\,(x,x')\;:\;q(x)=q(x')\;\wedge\;qF(x)\neq qF(x')\,\}.
  \]

  Threshold: \(\Theta_1^{\mathrm{desc}}:\delta_1^{\mathrm{split}}=\varnothing\). Polarity: obstruction. The status rule is \(\delta_1^{\mathrm{split}}=\varnothing\Rightarrow\) descent passes, while \(\delta_1^{\mathrm{split}}\neq\varnothing\) records an active \(P_1\) obstruction.
\item
  \textbf{Candidate descended-map defect.} Let \(q_{\mathrm{im}}:X\to\mathrm{im}(q)\) be \(q\) with its codomain restricted to its image. For a candidate \(F^\sharp:\mathrm{im}(q)\to\mathrm{im}(q)\), the defect is

  \[
  \delta_1^{\mathrm{cand}}(q,F,F^\sharp)\;=\;\{\,x:F^\sharp(q_{\mathrm{im}}(x))\neq q_{\mathrm{im}}F(x)\,\}.
  \]

  Threshold: \(\Theta_{=0}\). Polarity: obstruction.
\item
  \textbf{Markov closure deficit.} For finite Markov abstraction, the closure deficit is

  \[
  \mathrm{CD}_{\tau}^{\mu}(\Pi)\;=\;I(X_t;\Pi(X_{t+\tau})\mid\Pi(X_t)).
  \]

  Threshold: \(\Theta_1^{\mathrm{closure}}:\mathrm{CD}_{\tau}^{\mu}(\Pi)=0\) for exact closure, or \(\mathrm{CD}_{\tau}^{\mu}(\Pi)\le\varepsilon\) with \(\varepsilon=\log\eta\) for a declared rational \(\eta>1\) for approximate closure. Polarity: gap.
\end{itemize}

\subsubsection{\texorpdfstring{\(P_2\): Representability and Gating Defects}{P\_2: Representability and Gating Defects}}\label{sn:5.3.2}

The \(P_2\) defect family contains two named entries.

\begin{itemize}
\item
  \textbf{Representability defect.} For a lens \(f:Z\to O\) and predicate \(s\subseteq Z\),

  \[
  \delta_2^{\mathrm{rep}}(s;f)\;=\;
  \begin{cases}
  0, & s\in\Sigma_f,\\
  1, & s\notin\Sigma_f.
  \end{cases}
  \]

  By definition of \(\Sigma_f\) (Subsubsection \hyperref[sn:3.3.1]{3.3.1}), \(s\in\Sigma_f\) iff \(s=f^{-1}(B)\) for some \(B\subseteq\mathrm{im}(f)\). A nonrepresentability witness is a pair \(z,z'\) with \(f(z)=f(z')\) and \(\mathbf 1_s(z)\neq\mathbf 1_s(z')\). Threshold: \(\Theta_{\in S}\) for representable, \(\Theta_{\notin S}\) for nonrepresentable. Polarity: obstruction.
\item
  \textbf{Gating / support-loss defect.} For graph/support gating \(G\leadsto G'\),

  \[
  \delta_2^{\mathrm{cycle}}(G,G')\;=\;\beta_1(G)-\beta_1(G').
  \]

  The drive-suppression threshold is \(\Theta_2^{\mathrm{forest}}\), the instance of \(\Theta_{=0}\) whose evaluated datum is \(\beta_1(G')\), not \(\delta_2^{\mathrm{cycle}}\): it passes exactly when \(\beta_1(G')=0\), and \(\beta_1(G')=0\Rightarrow H^1(G')=0\). Polarity: capacity loss.
\end{itemize}

\subsubsection{\texorpdfstring{\(P_3\): Route-Mismatch, Holonomy, Pasting Defects}{P\_3: Route-Mismatch, Holonomy, Pasting Defects}}\label{sn:5.3.3}

The \(P_3\) defect family contains the following formula-level entries.

\begin{itemize}
\item
  \textbf{Completion commutator defect.} For completions \(E_1,E_2:C\to C\),

  \[
  \Delta_3^{\mathrm{comp}}(E_1,E_2)\;=\;\{\,c:E_1E_2(c)\neq E_2E_1(c)\,\}.
  \]

  Route mismatch is active when \(\Delta_3^{\mathrm{comp}}\neq\varnothing\), and route coherence holds when \(\Delta_3^{\mathrm{comp}}=\varnothing\). Polarity: obstruction.
\item
  \textbf{Holonomy defect.} A \(P_3\) transport record on a finite graph consists of a finite fiber \(\Phi\) and a permutation \(g_e\) of \(\Phi\) for each oriented edge \(e\), with \(g_{\bar e}=g_e^{-1}\) for the reversed edge. The holonomy of a closed path \(\gamma=e_1\cdots e_n\) is \(\operatorname{Hol}(\gamma)=g_{e_n}\circ\cdots\circ g_{e_1}\), and the holonomy defect is

  \[
  \delta_3^{\mathrm{hol}}(\gamma)\;=\;\{\,\phi\in\Phi:\operatorname{Hol}(\gamma)(\phi)\neq\phi\,\}.
  \]

  It is active when it is nonempty, that is, when \(\operatorname{Hol}(\gamma)\neq\mathrm{id}\). The associated nonclaim is

  \[
  P_3^{\mathrm{holonomy}}\not\Rightarrow \mathrm{P6}_{\mathrm{drive}}
  \]

  without an explicit bridge.
\end{itemize}

Pasting and critical-pair material belong to the \(P_3\) role vocabulary and to the high-structure square module, but no additional primitive \(P_3\) pasting-defect formula is introduced here beyond the named entries.

\subsubsection{\texorpdfstring{\(P_4\): Refinement, Staging, Gluing Defects}{P\_4: Refinement, Staging, Gluing Defects}}\label{sn:5.3.4}

The \(P_4\) defect family contains two named entries.

\begin{itemize}
\item
  \textbf{Refinement factor defect.} For \(f_0:Z\to O_0\), \(f_1:Z\to O_1\), and forgetting map \(\rho:O_1\to O_0\),

  \[
  \delta_4^{\mathrm{ref}}(f_0,f_1,\rho)\;=\;\{\,z:f_0(z)\neq\rho f_1(z)\,\}.
  \]

  Refinement passes iff \(\delta_4^{\mathrm{ref}}=\varnothing\). Polarity: obstruction.
\item
  \textbf{Gluing defect.} For local packages over an overlap,

  \[
  \delta_4^{\mathrm{glue}}(q_U,q_V)\;=\;\{\,x\in U\cap V:\mathrm{res}_{U,U\cap V}(q_U(x))\neq\mathrm{res}_{V,U\cap V}(q_V(x))\,\},
  \]

  for local packages \(q_U:U\to O_U\) and \(q_V:V\to O_V\) with restriction maps \(\mathrm{res}_{U,U\cap V}:O_U\to O_{U\cap V}\) and \(\mathrm{res}_{V,U\cap V}:O_V\to O_{U\cap V}\). With identity restrictions it is empty exactly when some \(q:U\cup V\to O\) has \(q|_U=q_U\) and \(q|_V=q_V\). For a finite cover \(\{U_\alpha\}\) with local packages \(q_\alpha\) and restrictions \(\mathrm{res}_{\alpha,\alpha\beta}\), set \(\delta_4^{\mathrm{glue}}(\{q_\alpha\})=\{(\alpha,\beta,x):x\in U_\alpha\cap U_\beta,\ \mathrm{res}_{\alpha,\alpha\beta}(q_\alpha(x))\neq\mathrm{res}_{\beta,\alpha\beta}(q_\beta(x))\}\); with identity restrictions it is empty exactly when some \(q:\bigcup_\alpha U_\alpha\to O\) has \(q|_{U_\alpha}=q_\alpha\) for every \(\alpha\), since pairwise agreement makes \(q(x):=q_\alpha(x)\) independent of \(\alpha\).

  Gluing passes iff \(\delta_4^{\mathrm{glue}}=\varnothing\). If local data exist but gluing fails, the induced promotion status is \(\texttt{globally\_obstructed}\).
\end{itemize}

\subsubsection{\texorpdfstring{\(P_5\): Packaging, Completion, Objecthood Defects}{P\_5: Packaging, Completion, Objecthood Defects}}\label{sn:5.3.5}

The \(P_5\) defect family contains three named entries.

\begin{itemize}
\item
  \textbf{Idempotence defect.} Given an endomap \(E:X\to X\), the idempotence defect is

  \[
  \delta_5^{\mathrm{idem}}(E)\;=\;\{\,x:E^2(x)\neq E(x)\,\}.
  \]

  Threshold: \(\Theta_5^{\mathrm{idem}}:\delta_5^{\mathrm{idem}}=\varnothing\). If \(E^2=E\), then \(E^n=E\) for \(n\ge1\); an update that applies \(E\) to a recorded input \(s\in\mathrm{im}(E)\) is then a no-op, and the cell classifier assigns \(\texttt{trivial}\) to every cell \(P_5\leftarrow P_j\) with such an update, for any informant \(P_j\), unless a row above it applies.
\item
  \textbf{Fixed-point residual.} The fixed-point residual is

  \[
  \delta_5^{\mathrm{fix}}(E)\;=\;\{\,o:E(o)\neq o\,\},
  \]

  On a host with a declared norm, \(\|E(o)-o\|\) measures a pointwise residual. The finite set \(\delta_5^{\mathrm{fix}}(E)\) passes \(\Theta_{=0}\) exactly when empty; an approximate threshold \(\Theta_{\le\varepsilon}\) uses a separately declared scalar residual.
\item
  \textbf{Package collapse defect.} The package-collapse values are \(\texttt{identity\_package}\), \(\texttt{total\_collapse}\) and \(\texttt{singleton\_collapse}\). For a cell whose level/lens/interface record \(L\) has lens \(\ell_L\), \(\texttt{total\_collapse}\in\delta\) exactly when \(\ell_L\) is constant on a carrier with at least two elements; \(\texttt{singleton\_collapse}\in\delta\) exactly when \(\ell_L\) is not constant, the fields of \(W\) whose declared type is the domain of \(\ell_L\) record at least two distinct values, and \(\ell_L\) maps every such recorded value to the same value; and \(\texttt{identity\_package}\in\delta\) exactly when \(U\in\mathsf{Upd}(P_5)\) installs a package map or completion equal to the one already recorded for its target. An identity lens or completion of the ambient package \(\mathcal T\) produces no entry. The values \(\texttt{total\_collapse}\) and \(\texttt{singleton\_collapse}\) induce \(\texttt{collapsed}\), and \(\texttt{identity\_package}\) induces \(\texttt{trivial}\) (Subsubsection \hyperref[sn:5.6.6]{5.6.6}).
\end{itemize}

\subsubsection{\texorpdfstring{\(P_6\): Audit, Provenance, Drive Defects}{P\_6: Audit, Provenance, Drive Defects}}\label{sn:5.3.6}

The \(P_6\) defect family contains three named entries.

\begin{itemize}
\item
  \textbf{Audit defect.} The audit defect is a finite set of failure codes,

  \[
  \begin{aligned}
  \delta_6^{\mathrm{audit}}\subseteq\{\,&\texttt{missing\_audit},\;\texttt{failed\_audit},\;\texttt{source\_failure},\\
  &\texttt{visibility\_failure},\;\texttt{threshold\_failure},\;\texttt{nonclaim\_failure},\;\texttt{circular\_audit}\,\},
  \end{aligned}
  \]

  and we write \(\delta_6^{\mathrm{audit}}=\texttt{audit\_passes}\) for \(\delta_6^{\mathrm{audit}}=\varnothing\). The code \(\texttt{failed\_audit}\) (a failed \(\mathsf{CheckResults}\) check) is distinct from the status \(\texttt{failed\_audit}\); the classifier rows of Subsection \hyperref[sn:5.6]{5.6} state which codes trigger each status. The audit threshold \(\Theta_6^{\mathrm{audit}}\) is the instance of \(\Theta_{\mathrm{audit}}\) (Subsubsection \hyperref[sn:5.2.2]{5.2.2}): it passes exactly when the \(\mathsf{CheckRules}\) field of \(A\) is drawn from \(\mathsf{CheckRules}_I\) and \(\delta_6^{\mathrm{audit}}(A)=\texttt{audit\_passes}\). Circular audit induces \(\texttt{undefined\_circular}\). A code \(c\) is present when \(c\in\delta_6^{\mathrm{audit}}\), and a set \(S\) of codes is met when \(\delta_6^{\mathrm{audit}}\cap S\neq\varnothing\); the classifiers below read the audit defect only through these two tests, and when several codes are present their priority orders decide which row applies. A defect record written \(\delta=\varnothing\), as in the running example of Subsection \hyperref[sn:3.8]{3.8}, has every defect set empty and every defect code at its passing value.

  The value of \(\delta_6^{\mathrm{audit}}\) is computed from the audit record \(A\) of Subsubsection \hyperref[sn:4.9.2]{4.9.2}: it is \(\{\texttt{missing\_audit}\}\) when the audit slot carries an absence marker; otherwise it is the set of codes of the checks of \(A\) that fail, the codes being \(\mathsf{SourceRefs}\) against the source policy (\(\texttt{source\_failure}\); this check fails exactly when some record \(R\) of \(\mathsf{SourceRefs}\) has \(\delta_6^{\mathrm{source}}(R;I)\notin\{\texttt{committed},\texttt{fallback\_admitted},\texttt{proxy\_admitted}\}\), the values at which \(\Theta_{\mathrm{source}}\) passes), \(\mathsf{VisibilityCheck}\) (\(\texttt{visibility\_failure}\)), \(\mathsf{CircularityCheck}\) (\(\texttt{circular\_audit}\)), \(\mathsf{ThresholdCheck}\) (\(\texttt{threshold\_failure}\)), \(\mathsf{Nonclaims}\) (\(\texttt{nonclaim\_failure}\)), and \(\mathsf{CheckResults}\) (\(\texttt{failed\_audit}\)); it is empty, that is \(\texttt{audit\_passes}\), when every check passes. An audit record \(A\) \emph{passes} \(\mathsf{AuditPolicy}_I\) when its \(\mathsf{CheckRules}\) field is drawn from \(\mathsf{CheckRules}_I\) and \(\delta_6^{\mathrm{audit}}(A)=\texttt{audit\_passes}\); this is the meaning of the phrase in Subsubsections \hyperref[sn:5.6.1]{5.6.1} and \hyperref[sn:10.2.1]{10.2.1}, Theorem \hyperref[cl:thm:22]{22} and Model Theorem \hyperref[cl:thm:34]{34}. The \(\mathsf{CheckResults}\) check includes resolution of \(\mathsf{ClaimRef}\) and every \(\mathsf{EvidenceRefs}\) entry, validation of \(\mathsf{Provenance}\) (resolution of each record named in the provenance trace, by a program of \(\mathsf{CheckRules}_I\)), and replay of \(\mathsf{CheckRules}\), which fails exactly when a replayed output differs from the recorded \(\mathsf{CheckResults}\) entry; a correctly recorded rejection is not a failure and is read by the \(\texttt{rejected}\) row. Failure of any of these gives \(\texttt{failed\_audit}\).
\item
  \textbf{Provenance / source defect.} The source defect has values

  \[
  \begin{aligned}
  \delta_6^{\mathrm{source}}\in\{\,&\texttt{committed},\;\texttt{fallback\_admitted},\;\texttt{proxy\_admitted},\;\texttt{dirty},\\
  &\texttt{fallback\_forbidden},\;\texttt{unknown},\;\texttt{contradictory},\;\texttt{missing}\,\}.
  \end{aligned}
  \]

  In outline: \(\delta_6^{\mathrm{source}}(R;I)\) is the first value in the eight-item list below that applies to the record \(R\) in the source slot. Reading the list needs three conventions, fixed in the next two paragraphs: which values pass \(\Theta_{\mathrm{source}}\), the governing profile \(\lambda\) at which \(\mathsf{SourcePolicy}_I\) is read, and what ``admits'' means.

  \emph{Passing values.} The source threshold \(\Theta_{\mathrm{source}}\) passes at \(\texttt{committed}\), \(\texttt{fallback\_admitted}\) and \(\texttt{proxy\_admitted}\) and fails at the other values. Real pair claims require \(\texttt{committed}\), or \(\texttt{fallback\_admitted}\) with an audited source-upgrade bridge; \(\texttt{fallback\_admitted}\) and \(\texttt{proxy\_admitted}\) otherwise support at most \(\texttt{real\_provisional}\).
\end{itemize}

\emph{Governing profile.} The value is computed from the source slot, \(\mathsf{SourcePolicy}_I\) and the governing profile \(\lambda\). The governing profile is that of the judgment whose source slot holds \(R\); for the role evidence \(R_i\) of a role-channel record, which carries no profile (Subsection \hyperref[sn:3.6]{3.6}), it is the default entry; for \(R\) in the \(\mathsf{EvidenceRefs}\) of a claim record, whose level profile is not in \(\mathsf{Prof}\), it is the default entry; for \(R\) in the \(\mathsf{SourceRefs}\) of an audit record it is the profile of the record named by \(\mathsf{ClaimRef}\), or, if that record carries no profile of \(\mathsf{Prof}\) (a claim record, whose level profile is not in \(\mathsf{Prof}\), or a package, source or host record), the default entry that \(\mathsf{SourcePolicy}_I\) declares; for \(R\) named as the evaluated datum of a source threshold \(\Theta_{\mathrm{source}}\) in the threshold record of a judgment, it is the profile of that judgment, or the default entry when that judgment carries no profile of \(\mathsf{Prof}\).

\emph{Fallback designation.} \(\mathsf{SourcePolicy}_I\) is a finite relation between source-type tags, admission modes (\(\texttt{full}\), \(\texttt{fallback}\), \(\texttt{provisional}\)) and profiles. A record of type \(\texttt{fallback}\) is a \emph{designated fallback} under \(\lambda\) when \((\texttt{fallback},m,\lambda)\in\mathsf{SourcePolicy}_I\) for some mode \(m\), and then receives \(\texttt{fallback\_admitted}\) when no earlier clause of the first-match list applies; the rules below do not distinguish the mode \(\texttt{fallback}\) from the other modes. The \(\texttt{real\_provisional}\) row needs \((\mathrm{type}_{\mathrm{src}}(R),\texttt{provisional},\mu)\in\mathsf{SourcePolicy}_I\), that is \((\texttt{fallback},\texttt{provisional},\mu)\) for a fallback source; a source-upgrade bridge needs \((\texttt{fallback},\texttt{full},\lambda)\).

\emph{Admission.} Throughout the paper, ``\(\mathsf{SourcePolicy}_I\) admits \(t\)'' (or ``admits a source of type \(t\)'') means \((t,\texttt{full},\lambda)\in\mathsf{SourcePolicy}_I\) for every profile \(\lambda\) of the instance and for the default entry, which every such instrument declares; for a profile \(\lambda'\) that \(\mathsf{SourcePolicy}_I\) does not list, \((t,m,\lambda')\in\mathsf{SourcePolicy}_I\) is read as \((t,m,\text{default})\in\mathsf{SourcePolicy}_I\), so extending or joining an instance changes no instrument record; admission for provisional use means \((t,\texttt{provisional},\lambda)\in\mathsf{SourcePolicy}_I\). A \emph{source-upgrade bridge} for \(R\) is a record named in \(\mathrm{trace}_{\mathrm{src}}\), of type \(\texttt{source-upgrade}\), both of whose endpoints are \(R\); it relays no other source, so it is not in the set \(B_R\) below. \(\mathsf{BridgePolicy}_I\) admits it when it admits that type, and \(\mathsf{SourcePolicy}_I\) admits it when \((\texttt{fallback},\texttt{full},\lambda)\in\mathsf{SourcePolicy}_I\). The value is the first of the following that holds:

\begin{enumerate}
\def\labelenumi{\arabic{enumi}.}
\tightlist
\item
  \(\texttt{missing}\): the slot is an audited absence marker;
\item
  \(\texttt{contradictory}\) or \(\texttt{unknown}\): \(\mathrm{type}_{\mathrm{src}}\) is that tag;
\item
  \(\texttt{dirty}\): \(\mathrm{trace}_{\mathrm{src}}\) records a failed audit or threshold check of \(R\) itself (one whose target is \(R\) or its construction; records that \(R\) holds as data do not make \(R\) dirty), or \(\mathrm{valid}_{\mathrm{src}}\) records that \(R\) is outside its validity condition;
\item
  if \(\mathrm{trace}_{\mathrm{src}}\) names a nonempty set \(B_R\) of typed source-bridge records with target \(R\) and another source record as source endpoint (\(R\) is then a proxy source): \(\texttt{proxy\_admitted}\) when some member of \(B_R\) is admitted by \(\mathsf{BridgePolicy}_I\) and has a passing audit, and \(\texttt{unknown}\) otherwise;
\item
  \(\texttt{fallback\_forbidden}\): \(\mathrm{type}_{\mathrm{src}}=\texttt{fallback}\) and \((\texttt{fallback},m,\lambda)\notin\mathsf{SourcePolicy}_I\) for every mode \(m\);
\item
  \(\texttt{fallback\_admitted}\): \(\mathrm{type}_{\mathrm{src}}=\texttt{fallback}\);
\item
  \(\texttt{committed}\): \(\mathrm{type}_{\mathrm{src}}\) is one of the five candidate tags of Subsubsection \hyperref[sn:4.9.1]{4.9.1} and \((\mathrm{type}_{\mathrm{src}},\texttt{full},\lambda)\in\mathsf{SourcePolicy}_I\);
\item
  \(\texttt{unknown}\) otherwise (in particular for a candidate-tag record not admitted with mode \(\texttt{full}\)).
\end{enumerate}

\emph{Well-foundedness.} A source record \(R'\) is a \emph{source predecessor} of a proxy source \(R\) when \(R'\) lies in the \(\mathsf{SourceRefs}\) of the audit of a member of \(B_R\). Rule 4 terminates because the evaluations it actually queries follow the finite acyclic evaluation dependency graph below; cycles in the unqualified source-predecessor relation need not be evaluation cycles. Form also the finite \emph{evaluation dependency graph}: its vertices are defect and threshold evaluations, including their record, instrument, profile and judgment context; an edge points from an evaluation to every evaluation queried by its defining checks, including source, audit, visibility and threshold checks. Admissibility (Subsubsection \hyperref[sn:4.10.1]{4.10.1}) requires this graph to be acyclic, and evaluations are defined by recursion along it. Reference reachability tests inspect records without evaluating the verdicts they reference.
- \textbf{Drive certificate.} The drive certificate is the separate typed record

\[
  \kappa_6^{\mathrm{drive}},
  \]

kept distinct from \(\delta_6^{\mathrm{audit}}\). In the finite graph/cohomology host, the cohomological drive-certificate schema is

\[
  [a]\neq0\in H^1(G)\;\Longleftrightarrow_{\mathrm{schema}}\;\exists\gamma:\oint_\gamma a\neq0.
  \]

This schema is available only in the declared finite graph/cohomology host with the relevant cohomology-drive bridge. It is not supplied by a generic audit record.

The separation between the audit face \(P_6\) and the drive face \(\mathrm{P6}_{\mathrm{drive}}\) is preserved at the defect level: a generic audit defect is in the \(P_6\) audit subfamily and does not by itself imply a drive certificate failure, and a missing drive certificate does not by itself imply an audit defect. The route-mismatch dependency \(P_3\mathrel{\mathsf{dep}_{\texttt{requires\_bridge}}}\mathrm{P6}_{\mathrm{drive}}\) of Subsubsection \hyperref[sn:4.6.2]{4.6.2} records this separation at the dependency level.

\subsection{Promotion-Specific Defects (Reference)}\label{sn:5.4}

\emph{Used by:} Theorem \hyperref[cl:thm:13]{13} defines the factorization defect; the proofs of Theorems \hyperref[cl:thm:20]{20} and \hyperref[cl:thm:21]{21}, the construction of Theorem \hyperref[cl:thm:22]{22}, and Model Theorem \hyperref[cl:thm:34]{34} use it.

Promotion bridges carry a finite family of defect records that are not attached to any single primitive but to the bridge as a whole. Each entry below names the defect, the threshold it is checked against, and the promotion status it induces on failure, in the schema of Subsections \hyperref[sn:5.2]{5.2} and \hyperref[sn:4.5]{4.5}.

\subsubsection{Strictness / Factorization Defect}\label{sn:5.4.1}

The strictness defect of a promotion bridge \(B_{j\to j+1}\) with object maps \(\pi_j,\pi_{j+1}\) on a finite carrier \(S\) is

\[
\Delta_{\mathrm{fact}}(\pi_j,\pi_{j+1})\;=\;\{\,(s,s')\in S\times S\;:\;\pi_j(s)=\pi_j(s')\;\wedge\;\pi_{j+1}(s)\neq\pi_{j+1}(s')\,\}.
\]

The strict threshold \(\Theta_{\mathrm{strict}}\) asks \(\Delta_{\mathrm{fact}}\neq\varnothing\); the non-strict threshold \(\Theta_{\mathrm{nonstrict}}\) asks \(\Delta_{\mathrm{fact}}=\varnothing\). When the strictness gate \(G_{\mathrm{strict}}\) of Subsubsection \hyperref[sn:10.2.1]{10.2.1} is checked,

\[
G_{\mathrm{strict}}=\texttt{strict\_pass}\iff\Delta_{\mathrm{fact}}\neq\varnothing,\qquad G_{\mathrm{strict}}=\texttt{non\_strict\_pass}\iff\Delta_{\mathrm{fact}}=\varnothing;
\]

the other four values of \(\mathrm{StrictGateStatus}\) record an unchecked or unauditable gate. The well-formedness gate \(G_{\mathrm{suff}}\) verifies that a recorded checked value agrees with the recomputed defect (Subsection \hyperref[sn:10.2]{10.2}).

Strictness is a property of the typed object-map data \(\pi_j,\pi_{j+1}\) alone; it is not by itself a closure or drive claim, and the nonclaim register of Subsection \hyperref[sn:15.2]{15.2} records NC-10 and NC-11 against any such inference. Theorem \hyperref[cl:thm:13]{13} of Section \hyperref[sn:7]{7} (factorization defect equivalence) states the equivalence between \(\Delta_{\mathrm{fact}}=\varnothing\) and the existence of a factorization \(\pi_{j+1}=\phi\,\pi_j\), and is the theorem-grade form of this defect.

\subsubsection{No-Smuggling Defect}\label{sn:5.4.2}

The no-smuggling defect \(\delta_{\mathrm{nosmuggle}}(B)\) records violations of the bridge's declared scope. The violation kinds form the following finite list:

\begin{enumerate}
\def\labelenumi{\arabic{enumi}.}
\tightlist
\item
  target-layer data used to prove target novelty,
\item
  suppressed content used without an admissible bridge,
\item
  simulation used as theorem,
\item
  local evidence used as global without a gluing record,
\item
  a \(P_3\)-witness used as a \(\mathrm{P6}_{\mathrm{drive}}\) certificate without a drive bridge,
\item
  a lossy summary used as exact data.
\end{enumerate}

Each record \(r\) of \(\mathsf{Use}(B)\) carries an audited origin tag \(\tau(r)\in\{\texttt{target},\texttt{simulation},\texttt{local},\texttt{P3\_witness},\texttt{lossy\_summary},\texttt{other}\}\), and each use of \(r\) an audited role tag \(\rho(r)\in\{\texttt{novelty},\texttt{theorem},\texttt{global},\texttt{drive},\texttt{exact},\texttt{other}\}\). Then \(\delta_{\mathrm{nosmuggle}}(B)\) is the set of pairs \((k,r)\) with: \(k=1\) if an independent target-layer assertion or certificate, rather than a defining object map evaluated on the common carrier, is used as a premise for target novelty, that is, \(\tau(r)=\texttt{target}\), \(\rho(r)=\texttt{novelty}\) and \(r\) is not the table of \(\pi_{j+1}\); \(k=2\) if \(r\) lies in \(\delta_{\mathrm{overread}}\) or \(\delta_{\mathrm{unknown}}\), or is the first component of an entry of \(\delta_{\mathrm{weakbridge}}\), for the bridge; \(k=3\) if \(\tau(r)=\texttt{simulation}\) and \(\rho(r)=\texttt{theorem}\); \(k=4\) if \(\tau(r)=\texttt{local}\), \(\rho(r)=\texttt{global}\), and \(\mathsf{Use}(B)\) contains no gluing record with \(\delta_4^{\mathrm{glue}}=\varnothing\); \(k=5\) if \(\tau(r)=\texttt{P3\_witness}\), \(\rho(r)=\texttt{drive}\), and \(V\) lists no admissible drive bridge; and \(k=6\) if \(\tau(r)=\texttt{lossy\_summary}\) and \(\rho(r)=\texttt{exact}\).

The threshold is \(\Theta_{\mathrm{nosmuggle}}:\delta_{\mathrm{nosmuggle}}=\varnothing\). Failure induces

\[
\mathrm{Promote}(B)\;:\;\texttt{failed\_no\_smuggling}
\]

in the promotion classifier of Subsubsection \hyperref[sn:5.6.4]{5.6.4}, and the no-smuggling gate \(G_{\mathrm{nosmuggle}}\) of Subsubsection \hyperref[sn:10.2.1]{10.2.1} records the gate verdict.

\subsubsection{Stability Defect}\label{sn:5.4.3}

Let \(E\) be the completion endomap of the target package on its content space. When the content space is \(\mathcal P(S)\) for the carrier \(S\) of the target object map \(\pi_{j+1}:S\to O_{j+1}\), the stability defect is

\[
\delta_{\mathrm{stab}}\;=\;\{\,o\in\mathrm{im}(\pi_{j+1})\;:\;E\bigl(\pi_{j+1}^{-1}(o)\bigr)\neq\pi_{j+1}^{-1}(o)\,\};
\]

otherwise the bridge records an embedding \(\iota:O_{j+1}\to\) content space, and \(\delta_{\mathrm{stab}}=\{\,o\in O_{j+1}:E(\iota o)\neq\iota o\,\}\),

or its approximate-residual analogue under threshold \(\Theta_{\le\varepsilon}\) when an approximate completion is in scope. The stability gate \(G_{\mathrm{stab}}\) is required when the bridge claims that the target objects are fixed points of \(E\). Failure induces

\[
\mathrm{Promote}(B)\;:\;\texttt{failed\_stability}.
\]

When the bridge does not claim fixed-point structure on the target, the stability gate is recorded as \(\texttt{not\_required}\) rather than enforced.

\subsubsection{Descent Defect Inside Promotion}\label{sn:5.4.4}

When a promotion bridge claims macro closure on the target package --- that is, when the bridge claims a descended dynamics \(F^\sharp\) on the target's quotient or interface --- the descent gate \(G_{\mathrm{desc}}\) is required. The gate uses the descent defect of Subsubsection \hyperref[sn:5.3.1]{5.3.1} (\(\delta_1^{\mathrm{split}}(q,F)\)) on the relevant quotient and update data of the target package. Failure induces

\[
\mathrm{Promote}(B)\;:\;\texttt{failed\_descent}.
\]

If no closure or descended-dynamics claim is made on the bridge, the descent gate is recorded as \(G_{\mathrm{desc}}=\texttt{not\_required}\) rather than enforced; in particular, the strictness gate alone does not require descent, since \(\pi_{j+1}\not\factor\pi_j\) does not by itself entail any claim about a descended dynamics on the target. NC-10 of Subsection \hyperref[sn:15.2]{15.2} is the corresponding nonclaim.

\subsection{Claim-Level Visibility Defects (Reference)}\label{sn:5.5}

The use-sets on which these defects are computed, \(\mathsf{Use}(C)\), \(\mathsf{Use}(B)\), \(\mathsf{Use}(Q)\), \(\mathsf{Use}(L)\) and \(\mathsf{Use}(\Xi)\), are defined in Subsubsection \hyperref[sn:5.5.2]{5.5.2} (paragraph \emph{Use-sets of cells, bridges and squares}); \(\mathsf{Use}(\varphi)\) in Subsubsection \hyperref[sn:3.2.1]{3.2.1}.

\emph{Used by:} Theorems \hyperref[cl:thm:1]{1}, \hyperref[cl:thm:2]{2} and \hyperref[cl:thm:4]{4} use the weak-bridge and required-bridge defects of Subsubsection \hyperref[sn:5.5.2]{5.5.2}; Theorem \hyperref[cl:thm:23]{23} states its result with the overreading defect of this subsection.

Instrument-indexed claim records carry a finite family of visibility defects, all aggregated into a single visibility defect record. The threshold and induced statuses are recorded here at schema-grade; the formal no-overreading theorem is Theorem \hyperref[cl:thm:23]{23} in Section \hyperref[sn:11]{11}.

The aggregated visibility defect of a claim \(\varphi\) under instrument \(I\) relative to a finite family \(\mathcal L\) of registered, audited bridge candidates is

\[
\delta_{\mathrm{vis}}(I,\varphi,\mathcal L)\;=\;\delta_{\mathrm{overread}}\;\cup\;\delta_{\mathrm{weakbridge}}\;\cup\;\delta_{\mathrm{outside}}\;\cup\;\delta_{\mathrm{unknown}}.
\]

The threshold is \(\Theta_{\mathrm{vis}}:\delta_{\mathrm{vis}}=\varnothing\). The components \(\delta_{\mathrm{overread}}\), \(\delta_{\mathrm{outside}}\) and \(\delta_{\mathrm{unknown}}\) are finite sets of use-set records, and \(\delta_{\mathrm{weakbridge}}\) is a finite set of pairs \((x,L)\); the union is taken as a disjoint union, and only its emptiness is used.

\subsubsection{Hidden-Content Overread}\label{sn:5.5.1}

The hidden-content overread defect is

\[
\delta_{\mathrm{overread}}(I,\varphi,\mathcal L)\;=\;\{\,x\in\mathsf{Use}(\varphi)\cap\mathsf{Suppressed}_I\;:\;\text{no }L\in\mathcal L\text{ bridges }x\,\}
\]

(a bridge \(L\) \emph{bridges} \(x\) when its suppressed endpoint \(c\) equals \(x\), Subsubsection \hyperref[sn:5.5.2]{5.5.2}).

A nonempty \(\delta_{\mathrm{overread}}\) records that the claim depends on suppressed content for which no bridge has been registered or audited under \(\mathsf{BridgePolicy}_I\). The no-overreading theorem (Theorem \hyperref[cl:thm:23]{23} of Section \hyperref[sn:11]{11}) later uses the condition \(\delta_{\mathrm{overread}}=\varnothing\) in its acceptance argument; this subsection only records the defect schema.

\subsubsection{Weak Bridge}\label{sn:5.5.2}

The weak-bridge defect \(\delta_{\mathrm{weakbridge}}(I,\varphi,\mathcal L)\) records suppressed content that has a registered bridge in \(\mathcal L\) but no bridge in \(\mathcal L\) that is admissible for the claim, for instance because every registered bridge is weaker than the claim requires (\(\mathsf{Strength}(L)<\mathsf{RequiredStrength}(\varphi)\)), unaudited, or circular:

\[
\begin{aligned}
\delta_{\mathrm{weakbridge}}(I,\varphi,\mathcal L)\;=\;\{\,(x,L)\;:\;&x\in\mathsf{Use}(\varphi)\cap\mathsf{Suppressed}_I,\;L\in\mathcal L\text{ bridges }x,\\
&\text{and no }L'\in\mathcal L\text{ with }\mathsf{AdmVisBridge}(L',I,\varphi)\text{ bridges }x\,\}.
\end{aligned}
\]

Here a bridge \(L=(\mathrm{id}_L,c,c',\ldots)\) of Subsubsection \hyperref[sn:3.5.3]{3.5.3} \emph{bridges} \(x\) when its suppressed endpoint is \(c=x\). Bridge strengths are values in a finite totally ordered set \((\mathsf{Strength}_I,\le)\) declared by \(\mathsf{BridgePolicy}_I\); \(\mathsf{Strength}(L)\in\mathsf{Strength}_I\) is a field of \(L\), \(\mathsf{RequiredStrength}(\varphi)=\mathsf{ReqStrength}_I(\mathsf{ClaimType}(\varphi))\in\mathsf{Strength}_I\) is the strength that the policy requires for the claim type of \(\varphi\), and \(<\) is the strict order of \(\mathsf{Strength}_I\).

\emph{Use-sets of cells, bridges and squares.} The same defects are defined for a directed cell \(C\), with \(\varphi\) replaced by \(C\): the use-set \(\mathsf{Use}(C)\) is the set of content references in the fields \(W_j,U_i,L,\Theta,A\) of \(C\) (a content reference in these fields is the record the field holds or names together with each record named by a reference field of that record, except the instrument named in an audit record's \(\mathsf{CheckRules}\) field; this expansion is applied once, not iterated, and a claim's declared \(\mathsf{Refs}(\varphi)\) is included directly, without expanding those records' reference fields), and \(\mathsf{RequiredStrength}(C)\in\mathsf{Strength}_I\) is the value that \(\mathsf{BridgePolicy}_I\) assigns to the pair \((\mathsf{BridgeType}_\lambda,\mathsf{Regime}_\lambda)\) of the profile of \(C\). Reachability tests, unlike use-sets, follow the full reference graph of Subsubsection \hyperref[sn:5.6.5]{5.6.5}. The same defects are defined for a promotion bridge \(B\): \(\mathsf{Use}(B)\) is the set of content references in the fields \(\mathcal T_j,\mathcal T_{j+1},S,O_j,O_{j+1},\pi_j,\pi_{j+1},L,\Theta,A,\mathcal N\) of its expanded record, \(\mathcal L_B\) is the set of bridges listed in its visibility field \(V\), and \(\mathsf{RequiredStrength}(B)\) is the value that \(\mathsf{BridgePolicy}_I\) assigns to \((\mathsf{BridgeType},\mathsf{Regime})\) of the profile in the last field of its expanded record (\(\lambda_{\mathrm{prom}}\) of Subsubsection \hyperref[sn:10.6.1]{10.6.1} in Theorem \hyperref[cl:thm:22]{22}, \(\lambda_{\mathrm{Cantor}}\) in Model Theorem \hyperref[cl:thm:34]{34}). For a visibility bridge \(L\), \(\mathsf{Use}(L)\) is the set of content references in \(c,c',\mathsf{Guarantee}_L,A_L,V_L,\mathcal N_L\). For a decorated square \(\Xi\) of Section \hyperref[sn:14]{14}, \(\mathsf{Use}(\Xi)\) is the set of content references in its fields \(X,F,q,F^\sharp,\mathrm{src}_\Xi,\mathrm{tgt}_\Xi,\mathsf{Decor}_\Xi,A_\Xi\); \(\mathcal L_\Xi\) is the set of bridges listed in \(V_\Xi\); \(\mathsf{RequiredStrength}(\Xi)=\mathsf{ReqStrength}_{I_{\mathrm{DC}}}(\mathrm{type}_\Xi)\), declared on square types as \(\mathsf{ProvisionalTypes}_I\) is; the host tag of \(\Xi\) is that of \(X\).

\emph{Square well-formedness.} \(\Xi\) is well formed when its fields are finite and typed, its boundary edges match (Subsubsection \hyperref[sn:14.1.2]{14.1.2}), its references resolve, and \(\delta_\Xi\) and its audit defect equal their recomputed values.

For a claim \(\varphi\), let \(\mathcal L_\varphi\) be its visibility-bridge field; for a directed cell \(C\), let \(\mathcal L_C\) be the registered bridges listed in its visibility field \(V\). The same defects are defined for a pair record \(Q=\mathsf{PairObs}^{\mu}_{\{i,j\}}\): \(\mathsf{Use}(Q)\) is the set of content references in its joint-witness, source-of-truth, branch-compatibility, threshold and audit fields, \(\mathcal L_Q\) is the set of bridges listed in its visibility record, and \(\mathsf{RequiredStrength}(Q)\) is the value \(\mathsf{BridgePolicy}_I\) assigns to \((\mathsf{BridgeType}_\mu,\mathsf{Regime}_\mu)\). For a claim, directed cell, pair record, promotion bridge or decorated square \(X\) (with the corresponding registered bridge family \(\mathcal L_X\) defined above), the bridge-admissibility predicate \(\mathsf{AdmVisBridge}(L,I,X)\) holds if and only if (i) \(L\in\mathcal L_X\), the registered bridges of \(X\); (ii) \(c\in\mathsf{Use}(X)\cap\mathsf{Suppressed}_I\) and \(c'\in\mathsf{Visible}_I\); (iii) \(\mathsf{BridgeType}_L\) is admitted by \(\mathsf{BridgePolicy}_I\) (for a cell or pair record with profile \(\lambda\), the pair \((\mathsf{BridgeType}_L,\mathsf{Regime}_\lambda)\) is an admitted (bridge type, regime) pair of \(\mathsf{BridgePolicy}_I\), so that its strength value is defined, Subsubsection \hyperref[sn:3.4.1]{3.4.1}); (iv) \(\mathsf{Strength}_L\ge\mathsf{RequiredStrength}(X)\) in \(\mathsf{Strength}_I\); (v) the audit defect of \(A_L\) is \(\texttt{audit\_passes}\); and (vi) \(\ulcorner X,I\urcorner\) is not reachable from \(\mathsf{Use}(L)\) by resolved references.

Each entry pairs a suppressed use-set element \(x\) with a registered bridge \(L\) at \(x\). It is recorded when no registered bridge at \(x\) meets conditions (i)--(vi) of \(\mathsf{AdmVisBridge}\) above: for instance, every bridge at \(x\) is weaker than required, has a failing audit, has a type the policy does not admit, or reaches the verdict reference of the claim. The promotion-gate semantics of Section \hyperref[sn:10]{10} give the later promotion-specific use of this comparison. A nonempty \(\delta_{\mathrm{weakbridge}}\) is reported separately from \(\delta_{\mathrm{overread}}\) so that ``no bridge'' and ``bridge failed'' failures are distinguishable in the classifier output and the audit record.

A directed cell may also require a bridge because of its profile, whether or not any of its content is suppressed. The required bridge types of a cell \(C\) with profile \(\lambda\) are

\[
\mathsf{ReqBridge}(C)\;=\;\begin{cases}\{\texttt{drive}\}&\text{if }\mathsf{Regime}_\lambda=\texttt{drive-certification},\\ \varnothing&\text{otherwise,}\end{cases}
\]

A required bridge of type \(\beta\) for \(C\) is a bridge \(L\) listed in \(V\) with \(\mathsf{BridgeType}_L=\beta\), admitted by \(\mathsf{BridgePolicy}_I\) for \(\lambda\), named by the bridge-check field of \(U\) (for \(\mathsf{CertifyDrive}\), its drive-check field), bridging the record named by the field of \(U\) that carries the certified datum (for \(\mathsf{CertifyDrive}\), its entry field), with \(A_L\) at \(\texttt{audit\_passes}\) and \(\mathsf{Strength}_L\ge\mathsf{RequiredStrength}(C)\), and, for \(\beta=\texttt{drive}\), \(L\) is admissible as a drive bridge in the sense of Subsubsection \hyperref[sn:9.3.5]{9.3.5}: \(\mathsf{BridgePolicy}_I\) admits the type \(\texttt{drive}\), and the guarantee record of \(L\) holds a declared finite graph \(G\), a cochain \(a\in C^1(G)\) and a cycle \(\gamma\in Z_1(G)\) with \(\oint_\gamma a\neq0\). When that record is visible under \(I\), a required bridge of type \(\texttt{drive}\) is instead a drive bridge of Subsubsection \hyperref[sn:9.3.5]{9.3.5} named by the drive-check field of \(U\) whose guarantee record is that record, admitted by \(\mathsf{BridgePolicy}_I\) for \(\lambda\), with \(A_L\) at \(\texttt{audit\_passes}\) and \(\mathsf{Strength}_L\ge\mathsf{RequiredStrength}(C)\). The required-bridge defect is

\[
\delta_{\mathrm{reqbridge}}(C)\;=\;\{\,\beta\in\mathsf{ReqBridge}(C):\text{no such bridge of type }\beta\text{ exists}\,\}.
\]

For a pair record \(Q\) with profile \(\mu\), \(\mathsf{ReqBridge}(Q)=\{\texttt{drive}\}\) when \(\mathsf{Regime}_\mu=\texttt{drive-certification}\), and is empty otherwise. A required bridge is listed in \(Q\)'s visibility record, admitted for \(\mu\), referenced by its threshold or audit check, and audited as certifying its declared datum at the strength required by \(\mu\), and, for the type \(\texttt{drive}\), admissible in the sense of Subsubsection \hyperref[sn:9.3.5]{9.3.5}. The set \(\delta_{\mathrm{reqbridge}}(Q)\) contains each required type for which no such bridge exists.

\subsubsection{Outside and Unknown Content}\label{sn:5.5.3}

The outside and unknown components of \(\delta_{\mathrm{vis}}\) record use-set entries that lie outside \(I\)'s scope or in the \(\texttt{unknown}\) visibility class:

\[
\delta_{\mathrm{outside}}(I,\varphi)\;=\;\mathsf{Use}(\varphi)\cap\mathsf{Outside}_I,\qquad \delta_{\mathrm{unknown}}(I,\varphi)\;=\;\mathsf{Use}(\varphi)\cap\mathsf{Unknown}_I.
\]

A nonempty \(\delta_{\mathrm{outside}}\) records that the claim refers to content the surrounding instrument does not adjudicate at all; a nonempty \(\delta_{\mathrm{unknown}}\) records that the claim refers to content within \(I\)'s scope whose visibility verdict has not yet been resolved. Both record use-set entries that \(I\) cannot treat as visible, independently of any bridge strength; the classifier of Subsubsection \hyperref[sn:5.6.5]{5.6.5} routes such records away from \(\texttt{accepted}\).

The table illustrates visibility-related outcomes; the first applicable row of each classifier determines the final status, and a nonempty \(\delta_{\mathrm{outside}}\) routes claims and cells to \(\texttt{outside\_scope}\) first:

\begin{longtable}[]{@{}ll@{}}
\toprule
\begin{minipage}[b]{0.47\columnwidth}\raggedright
Claim family\strut
\end{minipage} & \begin{minipage}[b]{0.47\columnwidth}\raggedright
Induced status\strut
\end{minipage}\tabularnewline
\midrule
\endhead
\begin{minipage}[t]{0.47\columnwidth}\raggedright
ordinary instrument-indexed claim\strut
\end{minipage} & \begin{minipage}[t]{0.47\columnwidth}\raggedright
\(\texttt{blocked}\)\strut
\end{minipage}\tabularnewline
\begin{minipage}[t]{0.47\columnwidth}\raggedright
directed cell\strut
\end{minipage} & \begin{minipage}[t]{0.47\columnwidth}\raggedright
\(\texttt{blocked}\)\strut
\end{minipage}\tabularnewline
\begin{minipage}[t]{0.47\columnwidth}\raggedright
pair observable\strut
\end{minipage} & \begin{minipage}[t]{0.47\columnwidth}\raggedright
\(\texttt{blocked}\)\strut
\end{minipage}\tabularnewline
\begin{minipage}[t]{0.47\columnwidth}\raggedright
promotion bridge\strut
\end{minipage} & \begin{minipage}[t]{0.47\columnwidth}\raggedright
the first applicable promotion-classifier status, including \(\texttt{outside\_scope}\), \(\texttt{failed\_audit}\), \(\texttt{failed\_no\_smuggling}\) or \(\texttt{candidate}\) as its gates require\strut
\end{minipage}\tabularnewline
\begin{minipage}[t]{0.47\columnwidth}\raggedright
same-level audit\strut
\end{minipage} & \begin{minipage}[t]{0.47\columnwidth}\raggedright
\(\texttt{undefined\_circular}\) when the failure is circular, unless a status of higher priority in Subsubsection \hyperref[sn:5.6.5]{5.6.5} applies first\strut
\end{minipage}\tabularnewline
\bottomrule
\end{longtable}

The promotion case routes through the no-smuggling defect of Subsubsection \hyperref[sn:5.4.2]{5.4.2} because a visibility failure on a promotion bridge is recorded as a no-smuggling violation in the gate semantics of Section \hyperref[sn:10]{10} (content outside scope is caught earlier, by the \(\texttt{outside\_scope}\) row of the promotion classifier); the same-level audit case is the special-case overreading of self-soundness recorded as circularity in Subsection \hyperref[sn:11.5]{11.5}.

\subsection{Defect-to-Status Rules}\label{sn:5.6}

The classifier rules below convert finite defect-record verdicts into status assignments per judgment family. Each rule is a finite check on the surrounding judgment, and each rule is consumed by the corresponding theorem in Sections \hyperref[sn:6]{6}, \hyperref[sn:10]{10}, or 11. The rules are stated as priority orders: when more than one condition matches, the highest-priority match is recorded.

\subsubsection{Role-Status Rules}\label{sn:5.6.1}

For a designated role-channel record \(C_i=(\mathrm{Chan}_i,R_i,\Theta_i^{\mathrm{act}},A_{R_i},V_{R_i},\mathsf{collapse}_i)\) in \(\Gamma\), the role-channel classifier \(\mathsf{RoleClassify}(C_i;\Gamma,\mathcal T,I)\) assigns a value in \(\mathrm{RoleStatus}\) in the following priority order

\begin{enumerate}
\def\labelenumi{\arabic{enumi}.}
\tightlist
\item
  \(\texttt{inapplicable\_outside\_scope}\), when the role evidence \(R_i\) lies outside \(\mathsf{Scope}_I\) or the host of \(\mathcal T\);
\item
  \(\texttt{absent\_with\_record}\), when a required witness for the role is absent and the absence is audited;
\item
  \(\texttt{collapsed\_boundary\_case}\), when the role-channel record's declared Boolean field \(\mathsf{collapse}_i\) is true;
\item
  \(\texttt{below\_threshold}\), when role evidence is present but the activation threshold is not met;
\item
  \(\texttt{active\_projection}\), otherwise, when role evidence is visible and audited.
\end{enumerate}

A role record whose required witness is absent without an audited absence record is malformed, as for directed cells (Subsubsection \hyperref[sn:3.7.3]{3.7.3}). A role-channel judgment with present evidence in \(\mathsf{Scope}_I\) and in the host of \(\mathcal T\) is well formed only when that evidence is visible under \(I\) and its audit passes \(\mathsf{AuditPolicy}_I\); a record failing either check receives no \(\mathrm{RoleStatus}\). Present evidence outside \(\mathsf{Scope}_I\) or outside the host of \(\mathcal T\) is well formed and receives \(\texttt{inapplicable\_outside\_scope}\) by rule 1. For well-formed role-channel judgments, rules 1--5 are exhaustive.

A role-channel record is a tuple \((\mathrm{Chan}_i,R_i,\Theta_i^{\mathrm{act}},A_{R_i},V_{R_i},\mathsf{collapse}_i)\), where \(R_i\) is a record with signature in \(\mathsf{Wit}(P_i)\) or an audited absence marker, \(\Theta_i^{\mathrm{act}}\) is a declared activation threshold, a finite check on \(R_i\), \(A_{R_i}\) is the audit record of \(R_i\) and \(V_{R_i}\) its visibility under \(I\); and \(\mathsf{collapse}_i\) is a declared Boolean field, which decides rule 3; the classifier reads only these fields. Role evidence is \emph{present} when \(R_i\) is a witness record; a required witness is \emph{absent} when \(R_i\) is an absence marker; the evidence is \emph{visible and audited} when \(R_i\in\mathsf{Visible}_I\) and its audit record passes \(\mathsf{AuditPolicy}_I\); and the activation threshold is \emph{not met} when the check \(\Theta_i^{\mathrm{act}}\) fails on \(R_i\).

\subsubsection{Cell-Status Rules}\label{sn:5.6.2}

The classifier is defined by the table of Subsubsection \hyperref[sn:5.6.6]{5.6.6}; the list below glosses its rows. The directed-cell classifier consumes the cell record together with the per-cell defect records of Subsection \hyperref[sn:5.3]{5.3} and the visibility defect of Subsection \hyperref[sn:5.5]{5.5}. The priority order is

\[
\begin{aligned}
&\texttt{outside\_scope}\;\succ\;\texttt{absent}\;\succ\;\texttt{blocked}\;\succ\;\texttt{undefined\_circular}\;\succ\\
&\texttt{collapsed}\;\succ\;\texttt{below\_threshold}\;\succ\;\texttt{trivial}\;\succ\;\texttt{implicit}\;\succ\;\texttt{action}.
\end{aligned}
\]

The per-status conditions are:

\begin{itemize}
\tightlist
\item
  \(\texttt{outside\_scope}\): the cell's host tag is not in \(\mathsf{Hosts}_I\), a level tag is not in \(\mathsf{Levels}_I\), or \(\delta_{\mathrm{outside}}(I,C)=\mathsf{Use}(C)\cap\mathsf{Outside}_I\neq\varnothing\).
\item
  \(\texttt{absent}\): \(W\) or \(U\) is an audited absence marker and \(\delta_6^{\mathrm{audit}}\) contains neither \(\texttt{missing\_audit}\) nor \(\texttt{failed\_audit}\); a cell whose audit is missing or failed is caught by the next row.
\item
  \(\texttt{blocked}\): the cell uses a missing or weak bridge, hidden suppressed content or unknown content, or its audit records a failed source or visibility check, is missing, failed, or fails its nonclaim check.
\item
  \(\texttt{undefined\_circular}\): the verdict reference \(\ulcorner C,I\urcorner\) of the cell under the judging instrument is reachable from the evidence or source references of its audit (\(\texttt{circular\_audit}\in\delta_6^{\mathrm{audit}}\)).
\item
  \(\texttt{collapsed}\): \(\delta\) contains \(\texttt{total\_collapse}\) or \(\texttt{singleton\_collapse}\), computed from the lens of \(L\) and \(W\) (Subsubsection \hyperref[sn:5.3.5]{5.3.5}).
\item
  \(\texttt{below\_threshold}\): \(W\) and \(U\) are present and the audit's \(\mathsf{ThresholdCheck}\) fails (\(\texttt{threshold\_failure}\in\delta_6^{\mathrm{audit}}\)).
\item
  \(\texttt{trivial}\): \(\mathsf{noopU}\) holds, or \(\delta\) contains \(\texttt{identity\_package}\); in particular, if the update applies an idempotent completion \(E\) to a recorded input \(s\in\mathrm{im}(E)\), then \(\mathsf{noopU}\) holds and the cell is \(\texttt{trivial}\) unless a row above it applies.
\item
  \(\texttt{implicit}\): \(\mathsf{implicitW}\) holds: \(W\) is attached to \(\mathcal T\) and no threshold of \(\Theta\) evaluates \(W\) (a threshold evaluates \(W\) when its evaluated datum is a field of \(W\), a direct check on values of \(W\) in the sense of Subsubsection \hyperref[sn:4.3.3]{4.3.3}, or a primitive defect computed from fields of \(W\)).
\item
  \(\texttt{action}\): the cell is well-formed, its witness and update are present, all required defects pass their thresholds, the visibility defect is empty (with admissible bridges in the bridge family \(\mathcal L\) of Subsection \hyperref[sn:5.5]{5.5} for any suppressed content used), and the audit and source policies pass.
\item
  \(\texttt{blocked}\) (final rule): included so that the table is total by construction; by Theorem \hyperref[cl:thm:4]{4} no well-formed cell reaches it.
\end{itemize}

For the running example of Subsection \hyperref[sn:3.8]{3.8}, no rule above \(\texttt{action}\) applies: the cell is in scope, its witness and update are present, it uses no bridge and no suppressed content, its audit is not circular, its content does not collapse, its threshold is met, its update is not the identity, and \(\mathsf{implicitW}\) is false, since \(W_3\notin W_{\mathcal T}=\varnothing\) and the threshold of \(\Theta\) evaluates \(W_3\). The classifier therefore assigns \(\texttt{action}\).

The final rule makes the classifier total on well-formed cells; the covered-cell uniqueness theorem (Theorem \hyperref[cl:thm:4]{4} of Section \hyperref[sn:6]{6}) proves that it assigns each of them exactly one status using this priority order.

\subsubsection{Pair-Status Rules}\label{sn:5.6.3}

The pair-observable classifier consumes the pair record together with its source-of-truth record and visibility defect, and assigns a value in the set \(\mathrm{PairStatus}\) of Subsection \hyperref[sn:4.4]{4.4}. The priority order is

\begin{enumerate}
\def\labelenumi{\arabic{enumi}.}
\tightlist
\item
  \(\texttt{blocked}\), when any of the source, visibility, threshold, or audit defects fails its threshold; or \(\delta_{\mathrm{reqbridge}}\neq\varnothing\); or when the source-of-truth record has type \(\texttt{unknown}\) or \(\texttt{contradictory}\) (then \(\delta_6^{\mathrm{source}}\) fails \(\Theta_{\mathrm{source}}\)); or when the pair contract forbids the available source class entirely; or when the branch-compatibility record has \(\mathsf{compatible}=\mathrm{false}\) or an audit defect other than \(\texttt{audit\_passes}\); or the host tag of the pair record, which is the host tag \(h_{\mathcal T}\) of its package as for a cell, is not in \(\mathsf{Hosts}_I\), or a level tag of \(\mu\) is not in \(\mathsf{Levels}_I\);
\item
  \(\texttt{real\_duplicated}\), when the pair record carries an explicit duplicate-pair record and the conditions of \(\texttt{real}\) below hold;
\item
  \(\texttt{real}\), when both (a) \(\delta_6^{\mathrm{source}}=\texttt{committed}\) (so \((\mathrm{type}_{\mathrm{src}},\texttt{full},\mu)\in\mathsf{SourcePolicy}_I\)), or \(\delta_6^{\mathrm{source}}=\texttt{fallback\_admitted}\) and \(R\) carries an audited source-upgrade bridge admitted by \(\mathsf{SourcePolicy}_I\) and \(\mathsf{BridgePolicy}_I\), and (b) the visibility defect is empty, \(\mathsf{partialEvidence}\) is false, and the audit and threshold policies pass;
\item
  \(\texttt{real\_provisional}\), when \(\delta_6^{\mathrm{source}}\in\{\texttt{fallback\_admitted},\texttt{proxy\_admitted}\}\), \(\mathsf{SourcePolicy}_I\) admits \(R\) for provisional use under \(\mu\), the visibility defect is empty, \(\mathsf{partialEvidence}\) is false, and every bridge this row needs is admissible (the first row excludes \(\delta_{\mathrm{reqbridge}}\neq\varnothing\), and \(\texttt{proxy\_admitted}\) already requires an admitted, audited source bridge, Subsubsection \hyperref[sn:5.3.6]{5.3.6});
\item
  \(\texttt{blocked}\), as a final rule, when none of the conditions above applies; with it the pair classifier assigns every well-formed pair record a status.
\end{enumerate}

\subsubsection{Promotion-Status Rules}\label{sn:5.6.4}

The gates are defined in Subsubsection \hyperref[sn:10.2.1]{10.2.1}. Here \(\mathsf{ReqCoreGates}(B)=\{G_{\mathrm{suff}},G_{\mathrm{ctrl}},G_{\mathrm{nosmuggle}},G_{\mathrm{vis}},G_{\mathrm{audit}}\}\), together with \(G_{\mathrm{desc}}\) when \(B\) claims macro closure or descended dynamics and \(G_{\mathrm{stab}}\) when it claims fixed-point targets; \(G_{\mathrm{locglob}}\) is present only when \(B\) claims a global package built from local packages. In brief: \(G_{\mathrm{suff}}\) checks well-formedness and recorded gate values; \(G_{\mathrm{desc}}\) and \(G_{\mathrm{stab}}\) are required only for closure or fixed-point claims; \(G_{\mathrm{ctrl}}\) and \(G_{\mathrm{nosmuggle}}\) test the defect of Subsubsection \hyperref[sn:5.4.2]{5.4.2}; \(G_{\mathrm{vis}}\) tests \(\delta_{\mathrm{vis}}\) on \(\mathsf{Use}(B)\); \(G_{\mathrm{audit}}\) tests the bridge audit; \(G_{\mathrm{strict}}\) records whether \(\Delta_{\mathrm{fact}}\) is nonempty; \(G_{\mathrm{locglob}}\) records gluing.

The promotion classifier consumes the bridge \(B_{j\to j+1}\), its gate record \(\mathsf{GateResults}(B)\), its strictness defect, no-smuggling defect, stability defect, and descent defect, and assigns a value in \(\mathrm{PromotionStatus}\). It applies to bridges \(B\) whose \(G_{\mathrm{suff}}\) checks (Subsubsection \hyperref[sn:10.2.1]{10.2.1}), computed from the bridge data, all hold, written \(\mathrm{WF}(B)\); \(\mathsf{GateResults}(B)\) records \(G_{\mathrm{suff}}=\texttt{pass}\) exactly when \(\mathrm{WF}(B)\). A bridge with \(\neg\mathrm{WF}(B)\) is not a well-formed promotion bridge and receives no promotion status, as a malformed directed-cell record receives no cell status (Subsubsection \hyperref[sn:4.3.3]{4.3.3}). The priority order is

\begin{enumerate}
\def\labelenumi{\arabic{enumi}.}
\tightlist
\item
  \(\texttt{outside\_scope}\), when the host tag \(h_B\) of \(B\) is not in \(\mathsf{Hosts}_I\), a level tag of its profile or of \(L\) is not in \(\mathsf{Levels}_I\), or \(\mathsf{Use}(B)\not\subseteq\mathsf{Scope}_I\) (\(B\in\mathsf{Use}(B)\), since the \(\mathsf{ClaimRef}\) of its audit names \(B\));
\item
  \(\texttt{failed\_audit}\), when the audit gate fails;
\item
  \(\texttt{failed\_no\_smuggling}\), when the no-smuggling gate, the visibility gate, or the control gate \(G_{\mathrm{ctrl}}\) fails (a failed control gate is the first violation kind of the no-smuggling defect of Subsubsection \hyperref[sn:5.4.2]{5.4.2});
\item
  \(\texttt{failed\_descent}\), when the descent gate fails on a bridge that requires descent;
\item
  \(\texttt{failed\_stability}\), when the stability gate fails on a bridge that requires stability;
\item
  \(\texttt{globally\_obstructed}\), when the optional local-to-global gate has \(G_{\mathrm{locglob}}=\texttt{glue\_fail}\);
\item
  \(\texttt{local\_only}\), when \(G_{\mathrm{locglob}}=\texttt{local\_only}\);
\item
  \(\texttt{candidate}\), when at least one required core gate has status \(\texttt{not\_checked}\), or a required \(G_{\mathrm{locglob}}\) is \(\texttt{not\_checked}\);
\item
  \(\texttt{strict}\), when every \(G\in\mathsf{ReqCoreGates}(B)\) passes and \(G_{\mathrm{strict}}=\texttt{strict\_pass}\), that is \(\Delta_{\mathrm{fact}}\neq\varnothing\), and \(G_{\mathrm{locglob}}\in\{\texttt{glue\_pass},\texttt{not\_required}\}\);
\item
  \(\texttt{non\_strict}\), when every \(G\in\mathsf{ReqCoreGates}(B)\) passes, \(G_{\mathrm{strict}}=\texttt{non\_strict\_pass}\), that is \(\Delta_{\mathrm{fact}}=\varnothing\), and the factorization is witnessed: the bridge records a map \(\phi\) with \(\pi_{j+1}=\phi\circ\pi_j\), and \(G_{\mathrm{locglob}}\in\{\texttt{glue\_pass},\texttt{not\_required}\}\);
\item
  \(\texttt{accepted}\), when every required core gate passes and the bridge does not require a strictness verdict, and \(G_{\mathrm{locglob}}\in\{\texttt{glue\_pass},\texttt{not\_required}\}\);
\item
  \(\texttt{candidate}\), as a final rule, when none of the conditions above applies; with it the promotion classifier assigns every well-formed bridge a status.
\end{enumerate}

The promotion-gate-soundness theorem (Theorem \hyperref[cl:thm:19]{19} of Section \hyperref[sn:10]{10}) later proves the accepted-status consequence for verdicts in \(\{\texttt{accepted},\texttt{strict},\texttt{non\_strict}\}\); this rule supplies only the classifier schema.

\subsubsection{Claim-Status Rules}\label{sn:5.6.5}

The claim classifier is defined by the table of Subsubsection \hyperref[sn:5.6.6]{5.6.6}; the list below glosses its rows. It consumes the claim record, its use-set, the visibility defect of Subsection \hyperref[sn:5.5]{5.5} computed with the claim's own bridge field \(\mathcal L_\varphi\) (Subsubsection \hyperref[sn:11.1.2]{11.1.2}), the threshold and source records, the audit record, and the bridge policy of \(I\). The priority order is

\[
\begin{aligned}
&\texttt{outside\_scope}\;\succ\;\texttt{blocked}\;\succ\;\texttt{absent\_with\_record}\;\succ\;\texttt{failed\_audit}\;\succ\\
&\texttt{undefined\_circular}\;\succ\;\texttt{below\_threshold}\;\succ\;\texttt{rejected}\;\succ\;\texttt{provisional}\;\succ\;\texttt{accepted}.
\end{aligned}
\]

\emph{Reference graph.} For reachability, form the finite directed graph whose vertices are the records of \(\mathsf{Cont}\) in the instance, with an edge from a record to each record named by one of its resolved reference fields; every reachability condition of the paper (the conditions on claims in Section \hyperref[sn:2]{2}, \(\mathsf{CircularityCheck}\) of Subsubsection \hyperref[sn:4.9.2]{4.9.2}, the visibility defects of Subsubsection \hyperref[sn:5.5.2]{5.5.2}, the rules below and the stack defects of Section \hyperref[sn:11]{11}) refers to this graph. Reachability from a set of records, in particular from \(\mathsf{Use}(\varphi)\) and from the records named by the \(\mathsf{EvidenceRefs}\) or \(\mathsf{SourceRefs}\) of an audit record, includes paths of length zero. A field of a record is a reference field exactly when its declared type is a record identifier or a finite set of record identifiers, unless the schema lists its reference fields explicitly; in either case, the record's own identifier field is not a reference field. The following schemas declare their reference fields explicitly; every other schema (packages, pair records, promotion bridges, decorated squares, threshold, defect, context and claim-log records, absence markers) follows the default rule of the previous sentence: those of an audit record are \(\mathsf{ClaimRef}\), \(\mathsf{EvidenceRefs}\), \(\mathsf{SourceRefs}\) and the instrument named in \(\mathsf{CheckRules}\); of a claim record, \(\mathcal T\), \(\mathsf{EvidenceRefs}\), \(\Theta\), \(A\), \(V\), \(\mathcal L\) and the identifiers of \(\mathsf{Refs}(\varphi)\); of a directed cell, the records in \(W_j,U_i,L,V,\Theta,A\); of a visibility bridge, \(c,c',\mathsf{Guarantee}_L,A_L,V_L,\mathcal N_L\); of a source record, the records named in \(\mathrm{trace}_{\mathrm{src}}\). For an instrument record they are the named programs of \(\mathsf{CheckRules}_I\); the sets \(\mathsf{Scope}_I\), \(\mathsf{Visible}_I\), \(\mathsf{Suppressed}_I\) and the policy tables are data fields, not reference fields, and \(\mathsf{ClaimLog}(I)\) is a separate record that \(I\) does not reference. A verdict reference \(\ulcorner\varphi,I\urcorner\) has exactly two reference fields, naming the record \(\varphi\) and the instrument \(I\). Cont may hold several records with these two reference fields; a condition that \(\ulcorner\varphi,I\urcorner\) is reachable means that some such record is reachable, and a condition that it is not reachable means that none is.

The per-status conditions are:

\begin{itemize}
\tightlist
\item
  \(\texttt{outside\_scope}\): the claim's host tag is not in \(\mathsf{Hosts}_I\), a level tag of its \(\mathsf{LevelProfile}\) is not in \(\mathsf{Levels}_I\), \(\mathsf{Use}(\varphi)\not\subseteq\mathsf{Scope}_I\), or its claim type is not in \(\mathsf{ClaimTypes}_I\).
\item
  \(\texttt{blocked}\): a visibility check, missing bridge, wrong evidence type, an evidence source failing \(\Theta_{\mathrm{source}}\), or a failed audit source check; threshold failures are routed separately to \(\texttt{below\_threshold}\).
\item
  \(\texttt{absent\_with\_record}\): the claim refers to records that are absent in \(\Gamma\) or \(\mathcal T\), with the absence audited (an absent entry of \(\mathsf{EvidenceRefs}\) is caught earlier by the \(\texttt{blocked}\) row).
\item
  \(\texttt{failed\_audit}\): the audit defect of the claim's audit record (Subsubsection \hyperref[sn:5.3.6]{5.3.6}) contains \(\texttt{missing\_audit}\), \(\texttt{failed\_audit}\) or \(\texttt{nonclaim\_failure}\); the other failing codes are routed to \(\texttt{blocked}\), \(\texttt{undefined\_circular}\) and \(\texttt{below\_threshold}\).
\item
  \(\texttt{undefined\_circular}\): the claim's audit defect contains \(\texttt{circular\_audit}\), or the claim refers to its own verdict under the judging instrument \(I\), that is, the verdict reference \(\ulcorner\varphi,I\urcorner\) is reachable from its use-set \(\mathsf{Use}(\varphi)\) by resolved references (as in \(\mathsf{CircularityCheck}\), Subsubsection \hyperref[sn:4.9.2]{4.9.2}), \(\ulcorner\varphi,I\urcorner\) being a record of \(\mathsf{Cont}\) that names the verdict of \(\varphi\) under \(I\); or the claim type is \(\texttt{soundness}\) and \(\mathsf{Use}(\varphi)\) contains a verdict reference \(\ulcorner\psi,I'\urcorner\) with \(\mathsf{Level}(I')\ge\mathsf{Level}(I)\). Reachability is a finite graph check and does not require evaluating the referenced verdict. The third condition is a stratification rule, not a cycle test: it applies whenever a soundness claim uses the verdict of any instrument \(I'\) whose level is not below \(\mathsf{Level}(I)\), including \(I'\ne I\) with no reference cycle.
\item
  \(\texttt{below\_threshold}\): a threshold check fails.
\item
  \(\texttt{rejected}\): the claim is well-formed but the check rules of \(I\) assign a rejection verdict.
\item
  \(\texttt{provisional}\): \(\mathsf{partialEvidence}\) is true and \(\mathsf{ClaimType}\in\mathsf{ProvisionalTypes}_I\).
\item
  \(\texttt{accepted}\): the claim is well-formed, \(\mathsf{partialEvidence}\) is false, the visibility defect is empty (with all required bridges admissible), the threshold and source policies pass, the audit policy passes, and the claim is not circular at the same level.
\item
  \(\texttt{blocked}\), as a final rule, when none of the conditions above applies; with it the claim classifier assigns every well-formed claim record a status.
\end{itemize}

If the check rules of \(I\) for the claim type of \(\varphi\) reject exactly when the recomputed finite witnesses falsify the content of \(\varphi\) (condition 6 of Subsubsection \hyperref[sn:2.4.1]{2.4.1}), then acceptance implies that those witnesses do not falsify it, and the content holds when it is a decidable property of them; without that condition, \(\texttt{accepted}\) records only that the listed checks passed.

The no-overreading theorem (Theorem \hyperref[cl:thm:23]{23} of Section \hyperref[sn:11]{11}) later uses the constraint \(\delta_{\mathrm{overread}}=\varnothing\) in the acceptance argument, and the same-level self-audit failure theorem (Theorem \hyperref[cl:thm:24]{24} of Section \hyperref[sn:11]{11}) later applies the circularity rule to complete same-level self-audit attempts without a level shift.

\subsubsection{Classifier Conditions on Record Fields}\label{sn:5.6.6}

The tables below restate the conditions of Subsubsections \hyperref[sn:5.6.1]{5.6.1} through \hyperref[sn:5.6.5]{5.6.5} as conditions on named fields of the classified record and on the defect values of Subsections \hyperref[sn:5.3]{5.3} through \hyperref[sn:5.5]{5.5}. Each classifier returns the first status in its table whose condition holds. For the directed-cell classifier, the last column evaluates each condition on the running example of Subsection \hyperref[sn:3.8]{3.8}.

\emph{Key.} \(\delta_{\mathrm{reqbridge}}\): profile-required bridge types with no admissible bridge (5.5.2); \(\delta_{\mathrm{weakbridge}}\): suppressed uses whose registered bridges are all inadmissible (5.5.2); \(\delta_{\mathrm{overread}},\delta_{\mathrm{unknown}},\delta_{\mathrm{outside}}\): unbridged suppressed, unknown-class and out-of-scope uses (5.5.1, 5.5.3); \(\delta_6^{\mathrm{audit}}\): codes of failed audit checks (5.3.6); \(\texttt{total\_collapse}\), \(\texttt{singleton\_collapse}\), \(\texttt{identity\_package}\): package-collapse entries (5.3.5); \(\mathsf{noopU}\): the update fixes its target; \(\mathsf{implicitW}\): the witness is attached to \(\mathcal T\) and evaluated by no threshold (4.3.4).

\emph{Directed-cell classifier}, for a well-formed cell record \(C=(P_i\leftarrow P_j,W,U,L,V,\Theta,A,\delta)^{\lambda}\) under \(I\) and \(\mathcal T\) (Subsubsection \hyperref[sn:4.3.3]{4.3.3}); a malformed record receives no status:

\begingroup\small
\begin{longtable}{@{}>{\raggedright\arraybackslash}p{0.2\linewidth}>{\raggedright\arraybackslash}p{0.5\linewidth}>{\raggedright\arraybackslash}p{0.24\linewidth}@{}}
\toprule
Status & Condition & Running example \\
\midrule
\endfirsthead
\toprule
Status & Condition & Running example \\
\midrule
\endhead
\bottomrule
\endfoot
$\texttt{outside\_scope}$ & the host tag of $C$ is not in $\mathsf{Hosts}_I$, a level tag of $L$ is not in $\mathsf{Levels}_I$, or $\delta_{\mathrm{outside}}(I,C)=\mathsf{Use}(C)\cap\mathsf{Outside}_I\neq\varnothing$ & false \\
$\texttt{absent}$ & $W$ or $U$ is an audited absence marker, and $\delta_6^{\mathrm{audit}}$ contains neither $\texttt{missing\_audit}$ nor $\texttt{failed\_audit}$ & false: $W$, $U$ present \\
$\texttt{blocked}$ & a bridge required by $\lambda$ is missing or weak ($\delta_{\mathrm{reqbridge}}\neq\varnothing$), or a suppressed use has registered bridges but none admissible ($\delta_{\mathrm{weakbridge}}\neq\varnothing$); or $\delta_{\mathrm{overread}}\neq\varnothing$ or $\delta_{\mathrm{unknown}}\neq\varnothing$; or $\delta_6^{\mathrm{audit}}$ contains one of $\texttt{source\_failure}$, $\texttt{visibility\_failure}$, $\texttt{missing\_audit}$, $\texttt{failed\_audit}$, $\texttt{nonclaim\_failure}$ & false \\
$\texttt{undefined\_circular}$ & $\texttt{circular\_audit}\in\delta_6^{\mathrm{audit}}$ & false \\
$\texttt{collapsed}$ & $\delta$ contains $\texttt{total\_collapse}$ or $\texttt{singleton\_collapse}$ & false: $\delta=\varnothing$ \\
$\texttt{below\_threshold}$ & $W$ and $U$ are present and $\texttt{threshold\_failure}\in\delta_6^{\mathrm{audit}}$, that is (condition 8 of Subsubsection \hyperref[sn:4.3.3]{4.3.3}), some threshold of $\Theta$ does not evaluate to $\texttt{pass}$ & false: $\Theta$ met \\
$\texttt{trivial}$ & $\mathsf{noopU}$ holds, or $\delta$ contains $\texttt{identity\_package}$ & false \\
$\texttt{implicit}$ & $\mathsf{implicitW}$ holds: $W$ is attached to $\mathcal T$ and no threshold of $\Theta$ evaluates $W$ & false \\
$\texttt{action}$ & none of the above; $W$ and $U$ are present; $C$ is well-formed; every threshold of $\Theta$ returns $\texttt{pass}$; the visibility defect is empty; $\delta_6^{\mathrm{audit}}=\texttt{audit\_passes}$ & true \\
$\texttt{blocked}$ & none of the above (unreachable for well-formed cells, Theorem \hyperref[cl:thm:4]{4}) & false \\
\end{longtable}
\endgroup

A well-formed cell whose audit defect contains \(\texttt{missing\_audit}\), \(\texttt{failed\_audit}\) or \(\texttt{nonclaim\_failure}\) is caught by the first \(\texttt{blocked}\) row unless an earlier row applies. By Theorem \hyperref[cl:thm:4]{4} no well-formed cell reaches the final row; it makes the classifier total by construction. Every well-formed cell therefore meets some row.

\(\mathsf{implicitW}\) holds exactly when \(W\in W_{\mathcal T}\), the witness family attached to \(\mathcal T\) (Subsubsection \hyperref[sn:3.3.1]{3.3.1}), and no threshold of \(\Theta\) evaluates \(W\); step 8 of Subsubsection \hyperref[sn:6.3.1]{6.3.1} recomputes it. The field \(\mathsf{noopU}\) is computed: it is \([e_U(r)=r]\) when the update \(U\) records its effect on its target record \(r\) as a finite map \(e_U\), including a declared completion, and false otherwise; step 8 of Subsubsection \hyperref[sn:6.3.1]{6.3.1} rechecks it. Some conditions in these tables are not reduced to other record fields: whether content collapses to a degenerate record, whether the pair contract forbids an available source class, whether partial evidence is present, and whether the branches of a pair are compatible. Each is a Boolean field of the classified record, set by the instance: \(\mathsf{collapse}\) on a role-channel record (deciding \(\texttt{collapsed\_boundary\_case}\); on a directed cell the package-collapse values of \(\delta\) play this part and are computed entries of the defect record, rechecked in step 8 of Subsubsection \hyperref[sn:6.3.1]{6.3.1}), \(\mathsf{forbidsSource}\) on a pair record, \(\mathsf{partialEvidence}\) on a claim, pair or square record, and \(\mathsf{compatible}\) on a branch-compatibility record. The classifiers read these fields like any other, so every classifier is a finite computation on the record's fields; what the fields assert about the instance is the instance's declaration, recorded and audited with the record.

\emph{Pair-observable classifier}, for a pair record \(\mathsf{PairObs}_{\{i,j\}}^{\mu}\) with source-of-truth record \(R\):

\begingroup\small
\begin{longtable}{@{}>{\raggedright\arraybackslash}p{0.22\linewidth}>{\raggedright\arraybackslash}p{0.72\linewidth}@{}}
\toprule
Status & Condition \\
\midrule
\endfirsthead
\toprule
Status & Condition \\
\midrule
\endhead
\bottomrule
\endfoot
$\texttt{blocked}$ & a source, visibility, threshold or audit defect fails its threshold; or $\delta_{\mathrm{reqbridge}}\neq\varnothing$; or $R$ has type $\texttt{unknown}$ or $\texttt{contradictory}$ (then $\delta_6^{\mathrm{source}}$ fails $\Theta_{\mathrm{source}}$); or the pair contract forbids the available source class; or the branch-compatibility record has $\mathsf{compatible}=\mathrm{false}$ or an audit defect other than $\texttt{audit\_passes}$; or the host tag of the pair record, which is the host tag $h_{\mathcal T}$ of its package as for a cell, is not in $\mathsf{Hosts}_I$, or a level tag of $\mu$ is not in $\mathsf{Levels}_I$ \\
$\texttt{real\_duplicated}$ & the record carries an explicit duplicate-pair record and the conditions of the $\texttt{real}$ row hold \\
$\texttt{real}$ & both (a) $\delta_6^{\mathrm{source}}=\texttt{committed}$ (so $(\mathrm{type}_{\mathrm{src}},\texttt{full},\mu)\in\mathsf{SourcePolicy}_I$), or $\delta_6^{\mathrm{source}}=\texttt{fallback\_admitted}$ and $R$ carries an audited source-upgrade bridge admitted by $\mathsf{SourcePolicy}_I$ and $\mathsf{BridgePolicy}_I$, and (b) the visibility defect is empty, $\mathsf{partialEvidence}$ is false, and the audit and threshold policies pass \\
$\texttt{real\_provisional}$ & $\delta_6^{\mathrm{source}}\in\{\texttt{fallback\_admitted},\texttt{proxy\_admitted}\}$, $\mathsf{SourcePolicy}_I$ admits $R$ for provisional use under $\mu$, the visibility defect is empty, $\mathsf{partialEvidence}$ is false, and every bridge this row needs is admissible (the first row excludes $\delta_{\mathrm{reqbridge}}\neq\varnothing$, and $\texttt{proxy\_admitted}$ already requires an admitted, audited source bridge, Subsubsection \hyperref[sn:5.3.6]{5.3.6}) \\
$\texttt{blocked}$ & none of the above: no row above assigns a status, for instance any source with $\mathsf{partialEvidence}$ true, or a $\texttt{fallback\_admitted}$ source that $\mathsf{SourcePolicy}_I$ does not admit for provisional use under $\mu$ \\
\end{longtable}
\endgroup

\emph{Promotion classifier}, for a well-formed bridge \(B\) (its gate \(G_{\mathrm{suff}}\) passes) with gate record \(\mathsf{GateResults}(B)\):

\begingroup\small
\begin{longtable}{@{}>{\raggedright\arraybackslash}p{0.22\linewidth}>{\raggedright\arraybackslash}p{0.72\linewidth}@{}}
\toprule
Status & Condition \\
\midrule
\endfirsthead
\toprule
Status & Condition \\
\midrule
\endhead
\bottomrule
\endfoot
$\texttt{outside\_scope}$ & the host tag $h_B$ of $B$ is not in $\mathsf{Hosts}_I$, a level tag of its profile or of $L$ is not in $\mathsf{Levels}_I$, or $\mathsf{Use}(B)\not\subseteq\mathsf{Scope}_I$ ($B\in\mathsf{Use}(B)$, since the $\mathsf{ClaimRef}$ of its audit names $B$) \\
$\texttt{failed\_audit}$ & $G_{\mathrm{audit}}=\texttt{fail}$ \\
$\texttt{failed\_no\_smuggling}$ & $G_{\mathrm{nosmuggle}}$, $G_{\mathrm{vis}}$ or $G_{\mathrm{ctrl}}$ has status $\texttt{fail}$ \\
$\texttt{failed\_descent}$ & $G_{\mathrm{desc}}$ is required and $G_{\mathrm{desc}}=\texttt{fail}$ \\
$\texttt{failed\_stability}$ & $G_{\mathrm{stab}}$ is required and $G_{\mathrm{stab}}=\texttt{fail}$ \\
$\texttt{globally\_obstructed}$ & $G_{\mathrm{locglob}}=\texttt{glue\_fail}$ \\
$\texttt{local\_only}$ & $G_{\mathrm{locglob}}=\texttt{local\_only}$ \\
$\texttt{candidate}$ & some $G\in\mathsf{ReqCoreGates}(B)$ has status $\texttt{not\_checked}$, or a required $G_{\mathrm{locglob}}$ is $\texttt{not\_checked}$ \\
$\texttt{strict}$ & every $G\in\mathsf{ReqCoreGates}(B)$ passes and $G_{\mathrm{strict}}=\texttt{strict\_pass}$ ($\Delta_{\mathrm{fact}}\neq\varnothing$), and $G_{\mathrm{locglob}}\in\{\texttt{glue\_pass},\texttt{not\_required}\}$ \\
$\texttt{non\_strict}$ & every $G\in\mathsf{ReqCoreGates}(B)$ passes, $G_{\mathrm{strict}}=\texttt{non\_strict\_pass}$ ($\Delta_{\mathrm{fact}}=\varnothing$), and $B$ records a map $\phi$ with $\pi_{j+1}=\phi\circ\pi_j$, and $G_{\mathrm{locglob}}\in\{\texttt{glue\_pass},\texttt{not\_required}\}$ \\
$\texttt{accepted}$ & every $G\in\mathsf{ReqCoreGates}(B)$ passes and the bridge does not require a strictness verdict, and $G_{\mathrm{locglob}}\in\{\texttt{glue\_pass},\texttt{not\_required}\}$ \\
$\texttt{candidate}$ & none of the above: no row above assigns a status, for instance strictness is required and $G_{\mathrm{strict}}$ is $\texttt{not\_checked}$ or $\texttt{bad\_audit}$, or a core gate or $G_{\mathrm{locglob}}$ is $\texttt{outside\_scope}$ \\
\end{longtable}
\endgroup

\emph{Claim classifier}, for a claim record \(\varphi\) under \(I\) with use-set \(\mathsf{Use}(\varphi)\):

\begingroup\small
\begin{longtable}{@{}>{\raggedright\arraybackslash}p{0.22\linewidth}>{\raggedright\arraybackslash}p{0.72\linewidth}@{}}
\toprule
Status & Condition \\
\midrule
\endfirsthead
\toprule
Status & Condition \\
\midrule
\endhead
\bottomrule
\endfoot
$\texttt{outside\_scope}$ & the claim's host tag is not in $\mathsf{Hosts}_I$, a level tag of its $\mathsf{LevelProfile}$ is not in $\mathsf{Levels}_I$, $\mathsf{Use}(\varphi)\not\subseteq\mathsf{Scope}_I$, or its claim type is not in $\mathsf{ClaimTypes}_I$ \\
$\texttt{blocked}$ & a visibility check fails ($\delta_{\mathrm{overread}}\neq\varnothing$, $\delta_{\mathrm{unknown}}\neq\varnothing$, or $\texttt{visibility\_failure}\in\delta_6^{\mathrm{audit}}$), a required bridge is missing or weak ($\delta_{\mathrm{weakbridge}}\neq\varnothing$), some entry $R$ of $\mathsf{EvidenceRefs}$ does not resolve to a source record of Subsubsection \hyperref[sn:4.9.1]{4.9.1} whose type tag $\mathsf{EvidencePolicy}_I$ admits, or has $\delta_6^{\mathrm{source}}(R;I)\notin\{\texttt{committed},\texttt{fallback\_admitted},\texttt{proxy\_admitted}\}$ (governing profile: the default entry, Subsubsection \hyperref[sn:5.3.6]{5.3.6}), or the source check fails ($\texttt{source\_failure}\in\delta_6^{\mathrm{audit}}$) \\
$\texttt{absent\_with\_record}$ & $\varphi$ refers to records absent from $\Gamma$ or $\mathcal T$, and the absence is audited \\
$\texttt{failed\_audit}$ & $\delta_6^{\mathrm{audit}}\cap\{\texttt{missing\_audit},\texttt{failed\_audit},\texttt{nonclaim\_failure}\}\neq\varnothing$ \\
$\texttt{undefined\_circular}$ & $\ulcorner\varphi,I\urcorner$ is reachable from $\mathsf{Use}(\varphi)$ by resolved references, or $\texttt{circular\_audit}\in\delta_6^{\mathrm{audit}}$, or the claim type is $\texttt{soundness}$ and some $\ulcorner\psi,I'\urcorner\in\mathsf{Use}(\varphi)$ has $\mathsf{Level}(I')\ge\mathsf{Level}(I)$ \\
$\texttt{below\_threshold}$ & a threshold check fails ($\texttt{threshold\_failure}\in\delta_6^{\mathrm{audit}}$) \\
$\texttt{rejected}$ & $\varphi$ is well-formed and the check rules of $I$ assign a rejection \\
$\texttt{provisional}$ & $\mathsf{partialEvidence}$ is true and $\mathsf{ClaimType}\in\mathsf{ProvisionalTypes}_I$ \\
$\texttt{accepted}$ & none of the above; $\varphi$ is well-formed, $\mathsf{partialEvidence}$ is false, the visibility defect is empty with all required bridges admissible, and the threshold, source and audit policies pass \\
$\texttt{blocked}$ & none of the above: no row above assigns a status \\
\end{longtable}
\endgroup

\emph{Role-channel classifier}, for a primitive \(P_i\) under \(\mathcal T\) and \(I\):

\begingroup\small
\begin{longtable}{@{}>{\raggedright\arraybackslash}p{0.3\linewidth}>{\raggedright\arraybackslash}p{0.64\linewidth}@{}}
\toprule
Status & Condition \\
\midrule
\endfirsthead
\toprule
Status & Condition \\
\midrule
\endhead
\bottomrule
\endfoot
$\texttt{inapplicable\_outside\_scope}$ & the role evidence $R_i$ lies outside $\mathsf{Scope}_I$ or the host of $\mathcal T$ \\
$\texttt{absent\_with\_record}$ & a required witness for the role is absent and the absence is audited \\
$\texttt{collapsed\_boundary\_case}$ & the role-channel record's declared Boolean field $\mathsf{collapse}_i$ is true \\
$\texttt{below\_threshold}$ & role evidence is present and the activation threshold is not met \\
$\texttt{active\_projection}$ & none of the above; role evidence is visible and audited \\
\end{longtable}
\endgroup

\emph{Square classifier}, for a well-formed decorated square \(\Xi\) of Section \hyperref[sn:14]{14} under \(I_{\mathrm{DC}}\). For a square whose decoration declares an order comparison, \(\delta_\Xi\) records the failures of that comparison; otherwise it records the failures of the exactness condition of its type (for a descent square, the candidate defect \(\delta_1^{\mathrm{cand}}\)).

\begingroup\small
\begin{longtable}{@{}>{\raggedright\arraybackslash}p{0.22\linewidth}>{\raggedright\arraybackslash}p{0.72\linewidth}@{}}
\toprule
Status & Condition \\
\midrule
\endfirsthead
\toprule
Status & Condition \\
\midrule
\endhead
\bottomrule
\endfoot
$\texttt{outside\_scope}$ & the host tag or part of the content of $\Xi$ lies outside the scope of the instrument \\
$\texttt{blocked}$ & the visibility defect of $\Xi$ is nonempty, or the audit defect of $A_\Xi$ contains $\texttt{source\_failure}$ or $\texttt{visibility\_failure}$ \\
$\texttt{failed\_audit}$ & $\delta_6^{\mathrm{audit}}(A_\Xi)\cap\{\texttt{missing\_audit},\texttt{failed\_audit},\texttt{nonclaim\_failure},\texttt{threshold\_failure}\}\neq\varnothing$ \\
$\texttt{undefined\_circular}$ & the audit defect of $A_\Xi$ contains $\texttt{circular\_audit}$ \\
$\texttt{obstructed}$ & $\delta_\Xi\neq\varnothing$ \\
$\texttt{lax}$ & the decoration declares an order comparison, $\mathsf{partialEvidence}_\Xi$ is false, $\delta_\Xi=\varnothing$, and the equality of the two boundary composites fails \\
$\texttt{provisional}$ & $\mathsf{partialEvidence}$ is true and $\mathrm{type}_\Xi\in\mathsf{ProvisionalTypes}_I$ \\
$\texttt{exact}$ & none of the above: $\mathsf{partialEvidence}$ is false, the exactness condition of the square's type holds and $\delta_\Xi=\varnothing$ \\
$\texttt{blocked}$ & none of the above: the recorded decoration and defect do not complete a per-type exactness or order-comparison check \\
\end{longtable}
\endgroup

Theorems \hyperref[cl:thm:30]{30}--\hyperref[cl:thm:32]{32} use \emph{mathematical exactness} as defined in Subsubsection \hyperref[sn:14.1.2]{14.1.2}. A mathematically exact square receives classifier status \(\texttt{exact}\) only after the preceding scope, visibility, audit and evidence rows have been ruled out.
 \section{Core Finite Interaction Theorems}\label{sn:6}

\emph{Reading key.} The status computations below use the directed-cell table of Subsubsection \hyperref[sn:5.6.6]{5.6.6}: rows are tested from top to bottom and the first that holds gives the status. Theorem \hyperref[cl:thm:1]{1} needs two records on one pair that trigger the \(\texttt{action}\) row and the \(\texttt{blocked}\) row (a too-weak bridge). Theorem \hyperref[cl:thm:2]{2} needs one \(\texttt{action}\) cell with \(i\neq j\). Theorem \hyperref[cl:thm:4]{4} is the first-match rule plus exhaustiveness of the rows \(\texttt{outside\_scope}\) through \(\texttt{action}\). Theorem \hyperref[cl:thm:5]{5} compares first-match outputs of different tables on records that share data. In order, the rows test: \(\texttt{outside\_scope}\) (a host or level tag outside \(I\), or a use outside \(\mathsf{Scope}_I\)); \(\texttt{absent}\) (an audited absence marker in the \(W\) or \(U\) slot, with neither \(\texttt{missing\_audit}\) nor \(\texttt{failed\_audit}\)); \(\texttt{blocked}\) (a nonempty required-bridge, weak-bridge, overread or unknown defect, or one of the codes \(\texttt{source\_failure}\), \(\texttt{visibility\_failure}\), \(\texttt{missing\_audit}\), \(\texttt{failed\_audit}\), \(\texttt{nonclaim\_failure}\)); \(\texttt{undefined\_circular}\) (\(\texttt{circular\_audit}\)); \(\texttt{collapsed}\) (\(\texttt{total\_collapse}\) or \(\texttt{singleton\_collapse}\)); \(\texttt{below\_threshold}\) (\(W\) and \(U\) present and \(\texttt{threshold\_failure}\)); \(\texttt{trivial}\) (\(\mathsf{noopU}\) or \(\texttt{identity\_package}\)); \(\texttt{implicit}\) (\(\mathsf{implicitW}\)); \(\texttt{action}\) (none of the above, with the remaining conditions of its row); and a final default row \(\texttt{blocked}\) that no well-formed cell reaches.

Sections \hyperref[sn:2]{2} through \hyperref[sn:5]{5} declared the finite audited interaction calculus, its judgment families, and the status and defect schemas. This section proves the core theorems of \(\mathbf{BirdInt}^{\mathrm{aud}}_{\mathrm{fin}}\): directed-cell status does not factor through the primitive pair, active typed interaction does not identify primitive roles, well-formedness of directed cells is decidable, every covered well-formed directed cell receives a unique status, and the five judgment-attached status families are not collapsed into one. All five theorems are theorem-grade, stated under explicit host, instrument, and admissibility assumptions, and proved finitely.

Theorems \hyperref[cl:thm:1]{1}--\hyperref[cl:thm:4]{4} use the definitions listed in the reader's map of Subsection \hyperref[sn:1.4]{1.4}, together with the profile grammar of Subsubsection \hyperref[sn:3.6.1]{3.6.1}, the source-record schema of Subsubsection \hyperref[sn:4.9.1]{4.9.1} and \(\delta_{\mathrm{vis}}\) at the head of Subsection \hyperref[sn:5.5]{5.5}; a reader who starts here needs no other part of Sections \hyperref[sn:3]{3}--\hyperref[sn:5]{5}. Theorem \hyperref[cl:thm:5]{5} also uses Subsubsections \hyperref[sn:5.6.1]{5.6.1}, \hyperref[sn:5.6.3]{5.6.3} and \hyperref[sn:5.6.4]{5.6.4}, Subsubsection \hyperref[sn:11.5.1]{11.5.1}, and one fact from each of three later theorems: the bridge \(B^{\mathrm{Toy22}}_{0\to1}\) of Theorem \hyperref[cl:thm:22]{22} is admissible and classified \(\texttt{strict}\) under \(I_{\mathrm{prom}}\), by Theorem \hyperref[cl:thm:24]{24} no admissible instrument \(I\) accepts the claim \(\mathrm{Sound}(I)\) about itself, and clause 4 of Theorem \hyperref[cl:thm:25]{25} applies to the one-instrument chain \((I_{\mathrm{prom}})\) with \(N=0\). Theorems \hyperref[cl:thm:22]{22}, \hyperref[cl:thm:24]{24} and \hyperref[cl:thm:25]{25} do not use Theorem \hyperref[cl:thm:5]{5} (Appendix \hyperref[sn:G]{G} lists these as cross-stage edges), so there is no circularity; cases (a)--(c) can be read now and cases (d)--(e) after Subsections \hyperref[sn:11.5]{11.5} and \hyperref[sn:11.8]{11.8}.

\subsection{Theorem 1: No Total Six-Symbol Algebra}\label{sn:6.1}

The two records of Theorem \hyperref[cl:thm:1]{1} are the running example of Subsection \hyperref[sn:3.8]{3.8}, the cell \(P_6\leftarrow P_3\) classified \(\texttt{action}\), and the cell described at the end of that subsection, with the same pair, a drive-certification update and a registered drive bridge too weak for the drive-certification profile, classified \(\texttt{blocked}\). Subsubsection \hyperref[sn:6.1.1]{6.1.1} lists the fields of both.

\begin{claimtheorem}{Theorem 1 (NoTotalAlgebra)}\phantomsection\label{cl:thm:1}

Across admissible instances of \(\mathbf{BirdInt}^{\mathrm{aud}}_{\mathrm{fin}}\), let \(\mathsf{CellData}\) be the class of typed directed-cell records together with their instance and judgment context, with the projection

\[
\mathsf{pair}\colon\mathsf{CellData}\to\mathbb P\times\mathbb P,\qquad C\mapsto(P_i,P_j)
\]

extracting the actor and informant primitive labels of \(C\), and the status assignment

\[
\mathsf{status}\colon\mathsf{CoveredCell}\to\mathrm{CellStatus}
\]

returning, on the covered records \(\mathsf{CoveredCell}\subseteq\mathsf{CellData}\) of Subsubsection \hyperref[sn:4.3.3]{4.3.3}, the directed-cell status of \(C\) given by the rows of the directed-cell table of Subsubsection \hyperref[sn:5.6.6]{5.6.6} (glossed in Subsubsection \hyperref[sn:5.6.2]{5.6.2}). There exist covered records \(C_{\mathrm{act}},C_{\mathrm{blk}}\in\mathsf{CoveredCell}\) with

\[
\mathsf{pair}(C_{\mathrm{act}})\;=\;\mathsf{pair}(C_{\mathrm{blk}})\;=\;(P_6,P_3),\qquad\mathsf{status}(C_{\mathrm{act}})\;=\;\texttt{action},\qquad\mathsf{status}(C_{\mathrm{blk}})\;=\;\texttt{blocked}.
\]

Consequently \(\mathsf{status}\) does not factor through the restriction of \(\mathsf{pair}\) to \(\mathsf{CoveredCell}\):

\[
\mathsf{status}\;\not\factor\;\mathsf{pair}|_{\mathsf{CoveredCell}}.
\]

In particular, there are no maps \(*\colon\mathbb P\times\mathbb P\to\mathbb P\) and \(d\colon\mathbb P\to\mathrm{CellStatus}\) with \(\mathsf{status}=d\circ *\circ\mathsf{pair}\) on \(\mathsf{CoveredCell}\). Both records lie in one admissible instance: the variant of the instance of Subsection \hyperref[sn:3.8]{3.8} that adds \(C_{\mathrm{blk}}\), \(V'\), \(A'\) and the verdict reference of \(C_{\mathrm{blk}}\), with these records added to \(\mathsf{Scope}_I\) and \(\mathsf{Visible}_I\) (item \emph{Scope} of Subsection \hyperref[sn:3.8]{3.8}). Moreover, \(C_{\mathrm{blk}}\) can be replaced by the record \(C_{\mathrm{reg}}\) that differs from \(C_{\mathrm{act}}\) only in the regime of its profile, set to \(\texttt{drive-certification}\), in the target of its audit, and in its recomputed defect record, which contains \(\delta_{\mathrm{reqbridge}}=\{\texttt{drive}\}\) (Subsubsection \hyperref[sn:3.6.3]{3.6.3}); \(C_{\mathrm{reg}}\) is covered with status \(\texttt{blocked}\). Finally, status is non-constant on every fibre of \(\mathsf{pair}|_{\mathsf{CoveredCell}}\): for every \((P_i,P_j)\in\mathbb P^2\), including \(i=j\), the instance of Subsection \hyperref[sn:3.8]{3.8} with two cells on \((P_i,P_j)\) added, together with their visibility, audit and verdict-reference records, in the same way, contains covered cells of statuses \(\texttt{absent}\) and \(\texttt{blocked}\).

\end{claimtheorem}

\subsubsection{Witness Records}\label{sn:6.1.1}

The theorem is stated under the following finite admissibility assumptions:

\begin{itemize}
\tightlist
\item
  \(\mathbb P=\{P_1,\ldots,P_6\}\) is the fixed finite primitive set of Subsubsection \hyperref[sn:2.2.1]{2.2.1}.
\item
  \(\mathsf{CellData}\) is the class of typed directed-cell records of Subsection \hyperref[sn:4.3]{4.3}, finite typed tuples \((P_i\leftarrow P_j,W_j,U_i,L,V,\Theta,A,\delta)^{\lambda}\) together with \(\Gamma,\mathcal T,I\); its subclass \(\mathsf{CoveredCell}\) consists of the covered records of Subsubsection \hyperref[sn:4.3.3]{4.3.3} (by Theorem \hyperref[cl:thm:4]{4}, exactly the well-formed records; the proof below checks the rows of both records directly and does not use Theorem \hyperref[cl:thm:4]{4}).
\item
  \(\mathrm{CellStatus}\) is the directed-cell status set fixed in Subsubsection \hyperref[sn:4.3.4]{4.3.4} and Subsubsection \hyperref[sn:5.1.2]{5.1.2}.
\item
  \(\mathsf{status}\) is the classifier of Subsubsection \hyperref[sn:5.6.2]{5.6.2}, applied to covered records.
\item
  \(C_{\mathrm{act}}\) is the running example \(\mathsf{Cell}_{63}^{\lambda_{\mathrm{audit}}}\) of Subsection \hyperref[sn:3.8]{3.8}, with the fields listed there. Both witnesses occur in the running-example instance extended by \(C_{\mathrm{blk}}\); the counterexample shows that no pair decoder works across the calculus.
\item
  \(C_{\mathrm{blk}}=\mathsf{Cell}_{63}^{\lambda_{\mathrm{drive}}}=(P_6\leftarrow P_3,W_3,U_6',L,V',\Theta,A',\delta')\) has the following fields. \(W_3=\mathsf{RouteMismatchWit}\), \(L\) and \(\Theta\) are those of \(C_{\mathrm{act}}\). The update is \(U_6'=\mathsf{CertifyDrive}\in\mathsf{Upd}(P_6)\). The profile \(\lambda_{\mathrm{drive}}\) agrees with \(\lambda_{\mathrm{audit}}\) except that its regime is drive certification and its bridge type is \(\texttt{drive}\), admitted by \(\mathsf{BridgePolicy}_I\). Let \(x\in\mathsf{Cont}\) be the candidate drive-certificate record named by the entry field of \(U'_6\) (a record-identifier field, so \(x\in\mathsf{Use}(C_{\mathrm{blk}})\)), with \(x\in\mathsf{Suppressed}_I\). The visibility record \(V'\) marks \(x\) suppressed and every other field of the cell and record of \(\mathsf{Use}(C_{\mathrm{blk}})\) visible, and lists \(L_{\mathrm{drive}}\) as the sole member of \(\mathcal L=\{L_{\mathrm{drive}}\}\) bridging \(x\), an audited bridge of the drive type registered by \(\mathsf{BridgePolicy}_I\). The bridge is registered but has strength below that required by \(\lambda_{\mathrm{drive}}\): \(\mathsf{Strength}(L_{\mathrm{drive}})=\texttt{weak}<\texttt{strong}=\mathsf{RequiredStrength}(C_{\mathrm{blk}})\) (item \emph{Bridge strengths}, Subsection \hyperref[sn:3.8]{3.8}). The drive-check field of \(U'_6\) names \(L_{\mathrm{drive}}\). The audit record \(A'\) has the cell as target, recomputes the two composites at \(\{a\}\) as in \(C_{\mathrm{act}}\), records the failed strength comparison and that no drive certificate is asserted, and its source references and nonclaim field are as in \(C_{\mathrm{act}}\); its \(\mathsf{VisibilityCheck}\) fails, because of that strength comparison, and its other checks pass.
\item
  \emph{Computed defects of \(C_{\mathrm{blk}}\).} The defect record \(\delta'\) records the computed values of Subsection \hyperref[sn:5.5]{5.5} and Subsubsection \hyperref[sn:5.3.6]{5.3.6}. Since \(x\in\mathsf{Use}(C_{\mathrm{blk}})\) is suppressed, \(L_{\mathrm{drive}}\) bridges \(x\), \(\mathsf{Strength}(L_{\mathrm{drive}})<\mathsf{RequiredStrength}(C_{\mathrm{blk}})\), and \(L_{\mathrm{drive}}\), the only bridge at \(x\), is therefore not admissible, the computed weak-bridge defect is \(\delta'_{\mathrm{weakbridge}}=\{(x,L_{\mathrm{drive}})\}\); since \(x\) has a registered bridge, \(\delta'_{\mathrm{overread}}=\varnothing\); the other use-set elements are visible, so the remaining visibility defects are empty; the regime of \(\lambda_{\mathrm{drive}}\) is drive certification and \(L_{\mathrm{drive}}\) is too weak, so \(\delta'_{\mathrm{reqbridge}}=\{\texttt{drive}\}\); and the visibility check of \(A'\) fails while its other checks pass, so \(\delta_6^{\mathrm{audit}}=\{\texttt{visibility\_failure}\}\).
\end{itemize}

\begin{proof}

Both records pass the checks of Subsubsection \hyperref[sn:6.3.1]{6.3.1}: their labels are in \(\mathbb P\), \(W_3\in\mathsf{Wit}(P_3)\), \(\mathsf{AddAuditRecord},\mathsf{CertifyDrive}\in\mathsf{Upd}(P_6)\), the profiles are admissible and every reference resolves. Both project to \((P_6,P_3)\). By Subsubsection \hyperref[sn:5.6.6]{5.6.6}, \(C_{\mathrm{act}}\) meets no row above \(\texttt{action}\) and meets the \(\texttt{action}\) row, so it is covered with status \(\texttt{action}\). For \(C_{\mathrm{blk}}\), the row \(\texttt{outside\_scope}\) fails since its host, levels and content are those of \(C_{\mathrm{act}}\); the row \(\texttt{absent}\) fails since \(W_3\) and \(U_6'\) are present; and the row \(\texttt{blocked}\) holds since the registered bridge is too weak, as witnessed by \((x,L_{\mathrm{drive}})\in\delta'_{\mathrm{weakbridge}}\). So \(C_{\mathrm{blk}}\) is covered with status \(\texttt{blocked}\).

Suppose for contradiction that \(\mathsf{status}\) factored through \(\mathsf{pair}\): that is, suppose there were a map

\[
g\colon\mathbb P\times\mathbb P\to\mathrm{CellStatus}
\]

such that \(\mathsf{status}(C)=g(\mathsf{pair}(C))\) for every \(C\in\mathsf{CoveredCell}\). Apply this to the two records \(C_{\mathrm{act}},C_{\mathrm{blk}}\) constructed in the first paragraph. Since \(\mathsf{pair}(C_{\mathrm{act}})=\mathsf{pair}(C_{\mathrm{blk}})\), evaluating \(g\) at this common pair gives

\[
\mathsf{status}(C_{\mathrm{act}})\;=\;g\bigl(\mathsf{pair}(C_{\mathrm{act}})\bigr)\;=\;g\bigl(\mathsf{pair}(C_{\mathrm{blk}})\bigr)\;=\;\mathsf{status}(C_{\mathrm{blk}}),
\]

contradicting \(\mathsf{status}(C_{\mathrm{act}})=\texttt{action}\neq\texttt{blocked}=\mathsf{status}(C_{\mathrm{blk}})\). Hence no such \(g\) exists, and \(\mathsf{status}\) does not factor through \(\mathsf{pair}\).

A total operation \(*\colon\mathbb P\times\mathbb P\to\mathbb P\) followed by a decoder \(d\colon\mathbb P\to\mathrm{CellStatus}\) is the special case \(g=d\circ *\), which is itself a map \(\mathbb P\times\mathbb P\to\mathrm{CellStatus}\). The same argument applies, and no such \((*,d)\) pair can encode \(\mathsf{status}\) faithfully.

For \(C_{\mathrm{reg}}\): its profile is admissible, since \((\texttt{fine},\texttt{drive-certification},\texttt{local},\mathsf H_{\mathrm{fin}})\in\mathsf{Compat}_I\) and its bridge type and mode are those of \(\lambda_{\mathrm{audit}}\). It uses no suppressed record, so its visibility defects are empty. \(\delta_{\mathrm{reqbridge}}\) is not part of the visibility defect, so its audit checks are those of \(C_{\mathrm{act}}\), retargeted to \(C_{\mathrm{reg}}\), and it passes steps 1--8 of Subsubsection \hyperref[sn:6.3.1]{6.3.1}. The rows \(\texttt{outside\_scope}\) and \(\texttt{absent}\) fail as for \(C_{\mathrm{act}}\). \(\mathsf{AddAuditRecord}\) has no bridge-check field, so \(\delta_{\mathrm{reqbridge}}(C_{\mathrm{reg}})=\{\texttt{drive}\}\) and the first \(\texttt{blocked}\) row holds. For the last clause, fix \((P_i,P_j)\), including \(i=j\). Extend the running instance by two cells and their supporting records, all typed on \(\mathsf H_{\mathrm{fin}}\) at \(\mathsf{Ver}\), and include every new record in the instrument's scope and visible set. Both cells use the \(L\) and \(\lambda_{\mathrm{audit}}\) of \(C_{\mathrm{act}}\), typed audited absence markers for \(P_j\) and \(P_i\) in their witness and update slots, and \(\Theta_{\mathrm{audit}}\) on the witness marker's distinct passing absence audit. Give the first cell a passing audit naming it and the second a typed, audited absence marker in its audit slot. For each cell construct a visibility record marking every field and every record of its use-set visible, and recompute its defect record and audit checks as required by Subsubsection \hyperref[sn:6.3.1]{6.3.1}. Both cells are well formed and in scope. The first meets the \(\texttt{absent}\) row. The second has \(\delta_6^{\mathrm{audit}}=\{\texttt{missing\_audit}\}\) by condition 8(b), so the \(\texttt{absent}\) row fails and the first \(\texttt{blocked}\) row holds.

\end{proof}

\subsubsection{Consequence}\label{sn:6.1.2}

The running example of Subsection \hyperref[sn:3.8]{3.8} supplies the hypothesis of Theorem \hyperref[cl:thm:1]{1}. Its cell on \(P_6\leftarrow P_3\) has status \(\texttt{action}\), and the blocked cell constructed in Subsubsection \hyperref[sn:6.1.1]{6.1.1} and restated in \(\mathsf{CM}_7\) keeps the same pair but replaces the update by \(\mathsf{CertifyDrive}\) under a profile that requires a drive bridge; its only registered drive bridge is too weak for that profile, so that cell is \(\texttt{blocked}\).

Directed cells \(P_i\leftarrow P_j\) are typed bridge judgments, not primitive products. The notation is closed against algebraic interpretation: there is no underlying multiplication of primitives whose values can determine cell status, and the typed witness, update, profile, level, visibility, threshold, audit, and defect data carried by every cell record are essential for the classifier's verdict. NC-1 of Subsection \hyperref[sn:15.2]{15.2} is the corresponding nonclaim. Subsequent sections of this paper use directed-cell records, not primitive-pair labels, as the working unit of typed interaction; in particular, Section \hyperref[sn:13]{13}'s PICA model-realization sheet does not interpret the \(6\times 6\) table as a universal interaction algebra.

\subsection{Theorem 2: Typed Non-Collapse}\label{sn:6.2}

The running example of Subsection \hyperref[sn:3.8]{3.8} is a well-formed active cell whose actor \(P_6\) and informant \(P_3\) differ.

\begin{claimtheorem}{Theorem 2 (TypedNonCollapse)}\phantomsection\label{cl:thm:2}

There exists a well-formed active directed cell

\[
\mathsf{Cell}_{ij}^{\lambda}\;=\;(\,P_i\leftarrow P_j,\;W_j,\;U_i,\;L,\;V,\;\Theta,\;A,\;\delta\,)
\]

in \(\mathbf{BirdInt}^{\mathrm{aud}}_{\mathrm{fin}}\) with \(i\neq j\) and status \(\texttt{action}\). Therefore active typed interaction does not imply equality of the actor and informant primitives:

\[
\mathsf{Cell}_{ij}^{\lambda}\;:\;\texttt{action}\;\;\not\Rightarrow\;\;P_i\;=\;P_j.
\]

More generally, for each \(\sigma\in\mathrm{CellStatus}\) there is an admissible instance containing a well-formed cell on the pair \((P_6,P_3)\) with status \(\sigma\); so no cell status forces actor and informant to coincide.

\end{claimtheorem}

\begin{proof}

The proof uses the type separation of the witness and update families of Subsection \hyperref[sn:3.7]{3.7} and the finite active-cell record supplied by the calculus.

A directed cell is a record

\[
(P_i\leftarrow P_j,\;W_j,\;U_i,\;L,\;V,\;\Theta,\;A,\;\delta)
\]

with the typing rules of Subsubsection \hyperref[sn:3.7.3]{3.7.3}:

\[
W_j\;\in\;\mathsf{Wit}(P_j),\qquad U_i\;\in\;\mathsf{Upd}(P_i).
\]

The rules carry both primitive labels as part of the data of the cell; they do not assert any equality relation between \(P_i\) and \(P_j\). The record constructor permits distinct primitive labels, for example \(P_1\leftarrow P_5\) for descent through a package or \(P_3\leftarrow P_5\) for noncommuting completions, with witness data drawn from the informant family and update data drawn from the actor family. The running example \(\mathsf{Cell}_{63}^{\lambda_{\mathrm{audit}}}\) of Subsection \hyperref[sn:3.8]{3.8} is such a cell: it is well formed, the table of Subsubsection \hyperref[sn:5.6.6]{5.6.6} gives it \(\texttt{action}\) (last column), and \(6\neq3\). This single admissible active cross-primitive cell refutes the implication from active directed-cell status to primitive equality.

The typing rule constrains \(W_j\) and \(U_i\) separately and states no relation between \(i\) and \(j\), so the judgment form does not entail \(P_i=P_j\); the nine variants below exhibit, for every status, a well-formed cell whose actor \(P_6\) differs from its informant \(P_3\).

The typing rule also admits diagonal cells \(P_i\leftarrow P_i\), whose witness and update are drawn from the distinct families \(\mathsf{Wit}(P_i)\) and \(\mathsf{Upd}(P_i)\); the last clause of Theorem \hyperref[cl:thm:1]{1} exhibits, on every diagonal pair, covered cells of statuses \(\texttt{absent}\) and \(\texttt{blocked}\).

For each of the nine cell statuses there is moreover a well-formed cell on the pair \((6,3)\), with \(P_6\neq P_3\), that receives it; each is a coherent finite variant of the running-example template of Subsection \hyperref[sn:3.8]{3.8}: change the indicated status-determining datum, then recompute its dependent audit, threshold, visibility, defect, and instance records as required by Subsubsection \hyperref[sn:6.3.1]{6.3.1}, step 8, and check the result against the table of Subsubsection \hyperref[sn:5.6.6]{5.6.6}. Except where stated, each variant keeps \(W_3\), \(U_6\), \(L\), \(V\) and \(\Theta\), has empty visibility and required-bridge defects, no package-collapse entry, \(\mathsf{implicitW}=\mathsf{noopU}=\mathrm{false}\) and \(\delta_6^{\mathrm{audit}}=\varnothing\), so every row before the one named fails.

\begin{itemize}
\tightlist
\item
  \(\texttt{outside\_scope}\): judge the cell under a new instrument \(I'\) with a fresh identifier that otherwise differs from \(I\) only in \(\mathsf{Scope}\), \(\mathsf{Visible}\) and \(\mathsf{vis}\); \(I\) stays in the instance, so the records naming \(I\) still resolve: its scope keeps the host and level tags of \(I\) and the record \(x\) but omits every record of \(\mathsf{Use}(C)\), including \(C\) itself, which the \(\mathsf{ClaimRef}\) of \(A\) names, \(\mathsf{Suppressed}_{I'}=\{x\}\) and \(\mathsf{Visible}_{I'}=\mathsf{Scope}_{I'}\setminus\{x\}\). \(I'\) is admissible, since its check rules and the audit's rules are those of \(I\). Reissue \(V\) and \(A\) under \(I'\), so that \(A\) names \(I'\) in its \(\mathsf{CheckRules}\) field. Step 8 then gives \(\delta_{\mathrm{outside}}=\mathsf{Use}(C)\) and \(\delta_6^{\mathrm{audit}}=\{\texttt{visibility\_failure},\texttt{threshold\_failure}\}\): the threshold datum lies in \(W_3\notin\mathsf{Scope}_{I'}\), so \(\mathsf{Eval}\) returns \(\texttt{outside\_scope}\) and the \(\mathsf{ThresholdCheck}\) fails (Subsubsections \hyperref[sn:5.2.2]{5.2.2} and \hyperref[sn:4.9.2]{4.9.2}). The first row holds and takes priority.
\item
  \(\texttt{absent}\): the witness slot is replaced by an audited absence marker, \(\Theta\) is replaced by \(\Theta_{\mathrm{audit}}\) on the audit defect of the marker's absence audit (a record distinct from \(A\), at \(\texttt{audit\_passes}\)), and \(A\) records the absence with passing check results; the cell is in scope, so the second row holds.
\item
  \(\texttt{blocked}\): the cell \(C_{\mathrm{blk}}\) of Theorem \hyperref[cl:thm:1]{1}, with \(\delta'_{\mathrm{reqbridge}}=\{\texttt{drive}\}\) and \(\delta_6^{\mathrm{audit}}=\{\texttt{visibility\_failure}\}\).
\item
  \(\texttt{undefined\_circular}\): the \(\mathsf{EvidenceRefs}\) of \(A\) are extended by the verdict reference \(\ulcorner C,I\urcorner\) of this variant cell \(C\), so the recomputed \(\mathsf{CircularityCheck}\) fails while its other checks pass, and \(\delta_6^{\mathrm{audit}}=\{\texttt{circular\_audit}\}\); the blocked row fails because no bridge or visibility defect occurs and neither \(\texttt{source\_failure}\) nor \(\texttt{visibility\_failure}\) is in \(\delta_6^{\mathrm{audit}}\).
\item
  \(\texttt{collapsed}\): the lens of \(L\) is replaced by the constant map \(\mathcal P(\{a,b,c\})\to\{*\}\), so step 8 computes \(\texttt{total\_collapse}\in\delta\).
\item
  \(\texttt{below\_threshold}\): replace the threshold by \(\Theta_{\ge2}\) on the datum \(|E_1E_2(\{a\})\,\triangle\,E_2E_1(\{a\})|\) (the two recomputed images differ in at least two elements). The images \(\{a,b\}\) and \(\{a,b,c\}\) differ only in \(c\), so the \(\mathsf{ThresholdCheck}\) fails and \(\delta_6^{\mathrm{audit}}=\{\texttt{threshold\_failure}\}\). This is the cell \(C_{\mathrm{thr}}\) reused in Theorem \hyperref[cl:thm:5]{5}.
\item
  \(\texttt{trivial}\): the audit log is recorded as a finite set \(\Lambda\) with an entry \(\epsilon_0\in\Lambda\), and \(\mathsf{AddAuditRecord}\) carries the effect map \(e_U(\Lambda)=\Lambda\cup\{\epsilon_0\}=\Lambda\), so step 8 computes \(\mathsf{noopU}=\mathrm{true}\).
\item
  \(\texttt{implicit}\): \(W_3\) is added to the witness family \(W_{\mathcal T}\) attached to \(\mathcal T\), with its own visibility, source and audit records (Subsubsection \hyperref[sn:3.3.3]{3.3.3}), and \(\Theta\) is replaced by \(\Theta_{\mathrm{pass}}\) on the recorded Boolean of the package audit entry of \(\mathcal T\) that recomputes the tables of \(E_1\) and \(E_2\), a record distinct from \(A\) whose value is true. No threshold then evaluates \(W_3\), so step 8 computes \(\mathsf{implicitW}=\mathrm{true}\).
\item
  \(\texttt{action}\): the running example itself.
\end{itemize}

So the non-implication holds for every status.

\end{proof}

The running example of Subsection \hyperref[sn:3.8]{3.8} is an instance: the cell \(P_6\leftarrow P_3\) has status \(\texttt{action}\) and distinct actor and informant.

\subsubsection{Corollary}\label{sn:6.2.1}

Pair labels alone do not determine cell status. The same primitive pair \((P_i,P_j)\), with different witness, update, profile, threshold, visibility, audit, or defect data, can produce different cell statuses; this was the working content of Theorem \hyperref[cl:thm:1]{1}, and Theorem \hyperref[cl:thm:2]{2} records the underlying typed-record discipline that licenses it. No theorem in this paper infers the equality of two primitive labels from the existence of a directed-cell judgment between them; in particular, no theorem in Sections \hyperref[sn:7]{7}, \hyperref[sn:10]{10}, or 13 does so even when the cell receives status \(\texttt{action}\).

\subsection{Theorem 3: Cell Well-Formedness Decidability}\label{sn:6.3}

\begin{claimtheorem}{Theorem 3 (CellWellFormednessDecidability)}\phantomsection\label{cl:thm:3}

In \(\mathbf{BirdInt}^{\mathrm{aud}}_{\mathrm{fin}}\), well-formedness of a candidate directed cell is decidable, provided (a) its numerical fields take values in declared types with decidable equality and order (the record-signature convention of Subsubsection \hyperref[sn:3.1.1]{3.1.1} supplies effective equality and every order comparison used by the procedure), and (b) every threshold that step 8 evaluates, directly or through an audit or source check it replays, has an evaluated datum computable in its type; these are all threshold vertices reachable from step 8 in the evaluation dependency graph (Subsubsection \hyperref[sn:5.3.6]{5.3.6}, paragraph Well-foundedness), including thresholds of \(\Theta\) and thresholds of audits queried by visibility, required-bridge, source or audit-defect checks. For each candidate cell record (a tuple of the form of Subsubsection \hyperref[sn:4.3.1]{4.3.1} whose fields are records of the instance or audited absence markers; last paragraph of Subsubsection \hyperref[sn:4.10.2]{4.10.2})

\[
C\;=\;(\,P_i\leftarrow P_j,\;W,\;U,\;L,\;V,\;\Theta,\;A,\;\delta\,)^{\lambda},
\]

together with \(\Gamma,\mathcal T,I\), there is a finite decision procedure \(\mathsf{WFCell}(C)\in\{\texttt{wf},\texttt{not\_wf}\}\) that returns \(\texttt{wf}\) exactly when \(C\) is well-formed in the sense of Subsubsection \hyperref[sn:4.3.3]{4.3.3}.

\end{claimtheorem}

\subsubsection{Algorithm}\label{sn:6.3.1}

The procedure runs the following finite checks in order; the cell is well-formed iff every step returns \(\texttt{wf}\).

\begin{enumerate}
\def\labelenumi{\arabic{enumi}.}
\tightlist
\item
  \textbf{Primitive labels.} Verify \(P_i,P_j\in\mathbb P\). This is a finite membership test on the six-element set.
\item
  \textbf{Witness slot.} Verify either a witness with signature in \(\mathsf{Wit}(P_j)\) or a typed, audited absence marker for that slot. The family \(\mathsf{Wit}(P_j)\) is a finite typed family (Subsubsection \hyperref[sn:3.7.1]{3.7.1}); the test is a finite membership check.
\item
  \textbf{Update slot.} Verify either an update with signature in \(\mathsf{Upd}(P_i)\) or a typed, audited absence marker for that slot. The family \(\mathsf{Upd}(P_i)\) is finite (Subsubsection \hyperref[sn:3.7.2]{3.7.2}); the test is a finite membership check.
\item
  \textbf{Profile admissibility.} Verify \(\mathsf{AdmProfile}(\lambda,I,\mathcal T)\). The check is the finite per-field comparison of Subsubsection \hyperref[sn:3.6.2]{3.6.2}.
\item
  \textbf{Resolution of \(L,V,\Theta,A,\delta,\lambda\).} Verify that each component is a finite typed record on the host of \(\mathcal T\), with all references in its fields resolved against \(\mathsf{Cont}\) or recorded as audited absences. In particular, verify that the level tags of \(L\) equal the actor and informant level tags of \(\lambda\), that a present \(U\) carries the tag \(\ell_i\) and a present \(W\) the tag \(\ell_j\), and that \(W\) and \(U\) are records on the host of \(\mathcal T\), and, if the actor and informant level tags of \(\lambda\) differ, that \(L\) contains a level-bridge record of one of the types of Subsubsection \hyperref[sn:2.2.3]{2.2.3}. Verify also that \(V\) assigns each field \(x\) of the cell, and each record \(x\) of the finite set \(\mathsf{Use}(C)\) (Subsubsection \hyperref[sn:5.5.2]{5.5.2}), the class \(\mathsf{vis}_I(x)\) determined by \(\mathsf{Visible}_I\), \(\mathsf{Suppressed}_I\) and \(\mathsf{Scope}_I\) (Subsubsection \hyperref[sn:3.5.2]{3.5.2}); this is a finite membership test.
\item
  \textbf{Reference resolution.} Verify that every reference in \(W,U,\delta,A\) --- to a quotient, a kernel, a defect record, an audit entry, a source identifier, or any other element of \(\mathsf{Cont}\) --- resolves to a record of \(\mathsf{Cont}\) of the declared type, or to an audited absence record. (Admissibility of the records of \(\mathsf{Cont}\) is the standing hypothesis \(\mathsf{AdmDomain}(\mathfrak D)\) of Subsection \hyperref[sn:4.10]{4.10}, not a step of \(\mathsf{WFCell}\).)
\item
  \textbf{Host and instrument typing.} Verify that the host and instrument tags are declared and every content reference is resolved or recorded as an audited absence. Whether the content is in \(\mathsf{Scope}_I\) is decided by the cell classifier.
\item
  \textbf{Defect agreement.} Compute the visibility defects and the required-bridge defect of Subsection \hyperref[sn:5.5]{5.5} from the finite use-set \(\mathsf{Use}(C)\), the profile \(\lambda\), the classes \(\mathsf{vis}_I(x)\) and the bridges listed in \(V\), and the finite bridge family, and the audit defect of Subsubsection \hyperref[sn:5.3.6]{5.3.6} from \(A\). Verify that the visibility, required-bridge and audit entries of \(\delta\) equal these values. Verify that \(\delta\) contains exactly the type-(iii) defects of Subsubsection \hyperref[sn:4.3.3]{4.3.3} named by \(\Theta\), each with its recomputed value and carrying exactly the threshold that \(\Theta\) assigns to it, and that its package-collapse entries equal the values computed from \(L\), \(W\) and \(U\) as in Subsubsection \hyperref[sn:5.3.5]{5.3.5}. When \(A\) is present, verify that its \(\mathsf{ClaimRef}\) names \(C\), that its \(\mathsf{CheckRules}\) names \(I\), and that its \(\mathsf{VisibilityCheck}\), \(\mathsf{ThresholdCheck}\) and \(\mathsf{CircularityCheck}\) results agree with the computed visibility defect, with the evaluation of \(\Theta\), and with the finite reachability test of Subsubsection \hyperref[sn:4.9.2]{4.9.2}; when \(A\) is an audited absence marker, verify that \(\delta_6^{\mathrm{audit}}=\{\texttt{missing\_audit}\}\). This is a finite computation. Verify also that the field \(\mathsf{noopU}\) equals its computed value (Subsubsection \hyperref[sn:5.6.6]{5.6.6}): \([e_U(r)=r]\) for an update whose effect on its target record \(r\) is recorded as a finite map \(e_U\), and false otherwise, and that \(\mathsf{implicitW}\) equals its value computed from the records attached to \(\mathcal T\) and from \(\Theta\); this includes an update that applies a declared completion \(E\) to a recorded input \(s\), with \(e_U=E\) and \(r=s\).
\end{enumerate}

\begin{proof}

Each step of the procedure is a finite decidable check: a membership test in a finite set, a finite per-field record comparison, or a finite policy comparison against a finite policy record. Step 8 also evaluates every threshold reachable from its evaluations in that graph, which terminate by hypothesis (the defect and threshold evaluations follow the acyclic evaluation dependency graph of Subsubsection \hyperref[sn:5.3.6]{5.3.6}; memoization keyed by record, instrument, profile and judgment context evaluates each reachable vertex at most once), and replays the check rules of each present \(A\) and \(A_L\), which terminate by admissibility of \(I\), including the reference-resolution and provenance checks of Subsubsection \hyperref[sn:5.3.6]{5.3.6}, which are finite lookups in \(\mathsf{Cont}\). The composition of finitely many decidable checks is decidable, hence \(\mathsf{WFCell}\) is decidable. It returns \(\texttt{wf}\) exactly when the cell is well-formed in the sense of Subsubsection \hyperref[sn:4.3.3]{4.3.3}, because steps 1, 2--3, 4, 5, 6 and 7 decide conditions 1--6 of that definition and step 8 decides conditions 7 and 8, and each step terminates.

The procedure is the well-formedness step that the classifier of Subsubsection \hyperref[sn:5.6.2]{5.6.2} invokes before assigning a cell status. Theorem \hyperref[cl:thm:4]{4} will use this decidability step when proving uniqueness of the status assigned to a covered well-formed cell.

\end{proof}

The hypothesis of Theorem \hyperref[cl:thm:3]{3} holds for the closure-deficit threshold of Subsubsection \hyperref[sn:5.3.1]{5.3.1} when the kernel and the law \(\mu\) of \(X_t\) are rational, \(\varepsilon=\log\eta\) with \(\eta\) rational, and logarithms are natural. Record \(\mathrm{CD}\) in the type of finite lists \(((r_k,s_k))_k\) of positive rationals, read as \(\sum_k r_k\log s_k\); two such values \(a,b\) satisfy \(a=b\) (resp. \(a\le b\)) exactly when \(\prod s_k^{Nr_k}=\prod {s'_k}^{Nr'_k}\) (resp. \(\le\)) for a common denominator \(N\), a rational test, so equality and order in this type are decidable. Explicitly, \(\mathrm{CD}=\sum\mu(x)p_x(y')\log\bigl(p_x(y')/\bar p_{\Pi(x)}(y')\bigr)\) over \(\mu(x)>0\), \(p_x(y')>0\), where \(p_x(y')=\sum_{\Pi(x')=y'}M^\tau(x,x')\) and \(\bar p_y(y')=\sum_{\Pi(x)=y}\mu(x)p_x(y')/\sum_{\Pi(x)=y}\mu(x)\), which is at least \(\mu(x)p_x(y')/\sum_{\Pi(x)=y}\mu(x)>0\); for rational \(M,\mu\) each coefficient and ratio is a positive rational. Then \(\mathrm{CD}=\sum_k r_k\log s_k\) with \(r_k,s_k\in\mathbb Q_{>0}\), so \(e^{\mathrm{CD}}=\prod_k s_k^{r_k}\) is algebraic. For \(\varepsilon=0\), \(\mathrm{CD}=0\) exactly when \(\prod_k s_k^{Nr_k}=1\) for a common denominator \(N\) of the \(r_k\), an exact rational test. For \(\varepsilon=\log\eta\) with \(\eta\in\mathbb Q_{>1}\), \(\mathrm{CD}\le\log\eta\) exactly when \(\prod_k s_k^{Nr_k}\le\eta^N\) for a common denominator \(N\), a rational comparison; with multiplication and natural-number powers declared among the operations of the rational type, it is a rule of \(\mathcal R\).

\subsection{Theorem 4: Covered Cell Status Uniqueness}\label{sn:6.4}

A cell is \emph{covered} when it is well formed and one of the nine rows \(\texttt{outside\_scope}\) through \(\texttt{action}\) of Subsubsection \hyperref[sn:5.6.6]{5.6.6} holds at it (Subsubsection \hyperref[sn:4.3.3]{4.3.3}). The coverage clause of Theorem \hyperref[cl:thm:4]{4} shows that every well-formed cell is covered, so the final \(\texttt{blocked}\) row is never used; uniqueness of the status follows from the first-match rule.

\begin{claimtheorem}{Theorem 4 (CoveredCellUniqueStatus)}\phantomsection\label{cl:thm:4}

Every well-formed directed cell satisfies at least one of the nine rows \(\texttt{outside\_scope}\) through \(\texttt{action}\) of the table of Subsubsection \hyperref[sn:5.6.6]{5.6.6}; a well-formed cell meeting none of the rows \(\texttt{outside\_scope}\) through \(\texttt{implicit}\) meets the \(\texttt{action}\) row. Hence \(\mathsf{CoveredCell}\) is the set of well-formed cells, decidable under the hypotheses of Theorem \hyperref[cl:thm:3]{3}, and the first-match classifier \(\mathsf{CellClassify}\) is a total function from it to \(\mathrm{CellStatus}\), computable under those hypotheses and never using the final row:

\[
\forall\,C\;\;\bigl(\mathsf{WFCell}(C)=\texttt{wf}\;\Rightarrow\;\exists\,r\in\{\texttt{outside\_scope},\ldots,\texttt{action}\}:\ \text{row }r\text{ holds at }C\bigr);
\]

\(\mathsf{CellClassify}(C)=\sigma\) exactly when the first row holding at \(C\) is a row of \(\sigma\). In particular, a well-formed cell whose audit defect contains \(\texttt{missing\_audit}\) or \(\texttt{failed\_audit}\) receives \(\texttt{outside\_scope}\) or \(\texttt{blocked}\), since the \(\texttt{absent}\) row excludes \(\texttt{missing\_audit}\) and \(\texttt{failed\_audit}\), and one whose audit defect contains \(\texttt{nonclaim\_failure}\) receives \(\texttt{outside\_scope}\), \(\texttt{absent}\) or \(\texttt{blocked}\).

\end{claimtheorem}

\subsubsection{Classifier Priority}\label{sn:6.4.1}

The status set is fixed in Subsubsection \hyperref[sn:5.1.2]{5.1.2}, and the priority order is frozen in Subsubsections \hyperref[sn:4.3.4]{4.3.4} and \hyperref[sn:5.6.2]{5.6.2}:

\[
\begin{aligned}
&\texttt{outside\_scope}\;\succ\;\texttt{absent}\;\succ\;\texttt{blocked}\;\succ\;\texttt{undefined\_circular}\;\succ\\
&\texttt{collapsed}\;\succ\;\texttt{below\_threshold}\;\succ\;\texttt{trivial}\;\succ\;\texttt{implicit}\;\succ\;\texttt{action}.
\end{aligned}
\]

A cell whose data simultaneously matches more than one per-status condition is assigned the highest-priority match; the per-status conditions are the rows of the directed-cell table of Subsubsection \hyperref[sn:5.6.6]{5.6.6}, glossed in Subsubsection \hyperref[sn:5.6.2]{5.6.2}.

\begin{proof}

The classifier is a finite first-match procedure over the fixed priority order. The proof has two parts: coverage (every well-formed cell meets one of the rows \(\texttt{outside\_scope}\) through \(\texttt{action}\)), proved in the coverage paragraph below, and uniqueness, which follows from the first-match rule; by coverage the final row is never used.

Existence of a status follows: the classifier walks the priority list, evaluates each per-status condition in order, and stops at the first match. Since, by the coverage paragraph below, one of the nine rows holds at every well-formed cell, the classifier returns some \(\sigma\in\mathrm{CellStatus}\) without reaching the final row.

Uniqueness follows from the priority discipline: if more than one condition matches a cell, the classifier records only the highest-priority one. The output is therefore a single value in \(\mathrm{CellStatus}\). Each per-status check in Subsubsection \hyperref[sn:5.6.2]{5.6.2} is a finite test on finite records; the entire classifier is a finite first-match procedure. Hence \(\mathsf{CellClassify}\) is a total function on \(\mathsf{CoveredCell}\) with codomain \(\mathrm{CellStatus}\), and assigns exactly one status to each covered well-formed cell.

By the coverage argument of the coverage paragraph below, one of those nine rows holds at every well-formed cell. Hence \(\mathsf{CoveredCell}\) is exactly the set of well-formed cells and is decidable under the hypotheses of Theorem \hyperref[cl:thm:3]{3}. The first \(\texttt{blocked}\) row tests the codes \(\texttt{missing\_audit}\), \(\texttt{failed\_audit}\) and \(\texttt{nonclaim\_failure}\), so a cell whose audit defect contains \(\texttt{missing\_audit}\) or \(\texttt{failed\_audit}\) receives \(\texttt{outside\_scope}\) if that row holds and \(\texttt{blocked}\) otherwise, since the \(\texttt{absent}\) row excludes these codes; a cell whose audit defect contains \(\texttt{nonclaim\_failure}\) receives \(\texttt{outside\_scope}\) or \(\texttt{absent}\) if one of those rows holds, and \(\texttt{blocked}\) otherwise.

For the coverage clause, let \(C\) be well formed and meet none of the rows \(\texttt{outside\_scope}\) through \(\texttt{implicit}\). If \(W\) or \(U\) were an audited absence marker, \(C\) would meet the \(\texttt{absent}\) row, or, when \(\delta_6^{\mathrm{audit}}\) contains \(\texttt{missing\_audit}\) or \(\texttt{failed\_audit}\), the first \(\texttt{blocked}\) row; so \(W\) and \(U\) are present. The first \(\texttt{blocked}\) row excludes the audit codes \(\texttt{missing\_audit}\), \(\texttt{failed\_audit}\), \(\texttt{nonclaim\_failure}\), \(\texttt{source\_failure}\) and \(\texttt{visibility\_failure}\), a nonempty required-bridge defect and a nonempty \(\delta_{\mathrm{overread}}\), \(\delta_{\mathrm{weakbridge}}\) or \(\delta_{\mathrm{unknown}}\); the \(\texttt{outside\_scope}\) row excludes a nonempty \(\delta_{\mathrm{outside}}\); the \(\texttt{undefined\_circular}\) row excludes \(\texttt{circular\_audit}\); the \(\texttt{collapsed}\) row excludes the collapse entries; and the \(\texttt{below\_threshold}\) row excludes \(\texttt{threshold\_failure}\). So the audit record is present and \(\delta_6^{\mathrm{audit}}=\texttt{audit\_passes}\), and by well-formedness condition 8 its \(\mathsf{ThresholdCheck}\) equals the evaluation of \(\Theta\), which therefore passes, so every threshold of \(\Theta\) returns \(\texttt{pass}\). The visibility defect is empty. Hence \(C\) meets the \(\texttt{action}\) row.

The proof uses the priority order of Subsubsection \hyperref[sn:5.6.2]{5.6.2} and introduces no other ordering. Subsubsection \hyperref[sn:5.6.2]{5.6.2} supplies the per-status conditions; this subsection only records that the priority assignment is well-defined and produces a unique value on each covered cell.

\end{proof}

\subsection{Theorem 5: Status-Family Separation}\label{sn:6.5}

\begin{claimtheorem}{Theorem 5 (StatusFamilySeparation)}\phantomsection\label{cl:thm:5}

In each of the five displayed non-implications below, some admissible instance satisfies the left judgment and not the right one; items (a)--(e) name the witnessing instances, and (d)--(e) use \(\mathrm{Sound}(I)\) and \(\mathsf{ClaimLog}(I)\) of Subsubsection \hyperref[sn:11.5.1]{11.5.1} and the constructions of Theorems \hyperref[cl:thm:22]{22} and \hyperref[cl:thm:25]{25}. The five judgment-attached status families

\[
\mathrm{RoleStatus},\quad\mathrm{CellStatus},\quad\mathrm{PairStatus},\quad\mathrm{PromotionStatus},\quad\mathrm{ClaimStatus}
\]

satisfy the following five non-implications in \(\mathbf{BirdInt}^{\mathrm{aud}}_{\mathrm{fin}}\). In each, some admissible instance satisfies the left judgment and not the right one. The fourth holds in the stronger form \(\mathrm{Promote}(B):\texttt{strict}\mathrel{\mathsf{dep}_{\texttt{blocks}}}\mathrm{Sound}(I^{+}_{\mathrm{prom}}):\texttt{accepted}\) of Subsubsection \hyperref[sn:4.6.1]{4.6.1}: by Theorem \hyperref[cl:thm:24]{24} the right judgment fails in every admissible instance, and the first instance of (d) shows that it fails with the claim in scope, visible, sourced and audited, at exactly \(\texttt{undefined\_circular}\); in the fifth, the right judgment holds in another instance:

\[
\mathrm{Chan}_i\;:\;\texttt{active\_projection}\;\wedge\;\mathsf{Cell}_{ij}^{\lambda}\text{ well formed with }W_j,U_i\text{ present}\;\;\not\Rightarrow\;\;\mathsf{Cell}_{ij}^{\lambda}\;:\;\texttt{action},
\]

\[
\mathsf{Cell}_{ij}^{\lambda}\;:\;\texttt{action}\;\;\not\Rightarrow\;\;\mathsf{PairObs}_{\{i,j\}}^{\mu}\;:\;\texttt{real},
\]

\[
\mathrm{Promote}(B)\;:\;\texttt{strict}\;\;\not\Rightarrow\;\;\exists i,j,\lambda\;\mathsf{Cell}_{ij}^{\lambda}\;:\;\texttt{action},
\]

\[
\mathrm{Promote}(B)\;:\;\texttt{strict}\;\;\not\Rightarrow\;\;\Gamma;\mathcal T_0;I^{+}_{\mathrm{prom}}\vdash\mathrm{Sound}(I^{+}_{\mathrm{prom}})\;:\;\texttt{accepted},
\]

\[
\mathrm{Promote}(B)\;:\;\texttt{strict}\;\;\not\Rightarrow\;\;\Gamma;\mathcal T^{\mathrm{audit}}_{\le 0};I^{\mathrm{stack}}_1\vdash\mathrm{Sound}^{\uparrow}_0\;:\;\texttt{accepted}.
\]

Here \(\mathrm{Sound}(I)\) is the claim, of type \(\texttt{soundness}\), that every verdict \(I\) has recorded (including its verdict on this claim) equals the status its claim classifier assigns (Subsubsection \hyperref[sn:11.5.1]{11.5.1}); \(\mathrm{Sound}^{\uparrow}_0\) is the same assertion about the recorded verdicts of \(I_{\mathrm{prom}}\), judged by an instrument one level higher (Theorem \hyperref[cl:thm:25]{25}); \(I_{\mathrm{prom}}^{+}\) is \(I_{\mathrm{prom}}\) with a fresh identifier, \(\texttt{soundness}\) added to its claim types, the use-set of \(\mathrm{Sound}(I^{+}_{\mathrm{prom}})\) added to its scope and visible set, and its required-strength map extended by \(\mathsf{ReqStrength}_{I_{\mathrm{prom}}^+}(\texttt{soundness})=\texttt{weak}\), retaining its old values; and \(\mathrm{Sound}^{\uparrow}_0\), \(\mathcal T^{\mathrm{audit}}_{\le 0}\) and \(I^{\mathrm{stack}}_1\) are the lifted soundness claim, audit package and extension instrument of Theorem \hyperref[cl:thm:25]{25} for the one-instrument chain \((I_{\mathrm{prom}})\), at the level \(\mathsf{Str}\) above \(\mathsf{Level}(I_{\mathrm{prom}})=\mathsf{Ver}\). Each non-implication asserts the existence of an admissible instance of \(\mathbf{BirdInt}^{\mathrm{aud}}_{\mathrm{fin}}\) in which the left judgment holds and the right one fails for the records listed:

\begin{itemize}
\item
  \begin{enumerate}
  \def\labelenumi{(\alph{enumi})}
  \tightlist
  \item
    for \(i=6\), \(j=3\), the instance obtained from the running example of Subsection \hyperref[sn:3.8]{3.8} by replacing its cell with the cell \(C_{\mathrm{thr}}\), the running-example cell with threshold \(\Theta_{\ge2}\) on the datum \(|E_1E_2(\{a\})\,\triangle\,E_2E_1(\{a\})|\) (the two recomputed images differ in at least two elements), keeping \(U_6'\), \(x\) and \(L_{\mathrm{drive}}\) in \(\mathsf{Cont}\) and adding the role-channel record \((\mathrm{Chan}_6,R_6,\Theta_6^{\mathrm{act}},A_{R_6},V_{R_6},\mathrm{false})\) with its audit record \(A_{R_6}\) (Subsubsection \hyperref[sn:3.7.1]{3.7.1}); there the role channel of \(P_6\) is \(\texttt{active\_projection}\), and \(C_{\mathrm{thr}}\), the only directed cell of the instance, is \(\texttt{below\_threshold}\), so no directed cell on the pair \((6,3)\) is \(\texttt{action}\);
  \end{enumerate}
\item
  \begin{enumerate}
  \def\labelenumi{(\alph{enumi})}
  \setcounter{enumi}{1}
  \tightlist
  \item
    for \(i=6\), \(j=3\), the instance of countermodel \(\mathsf{CM}_4\), in which the running-example cell \(\mathsf{Cell}_{63}^{\lambda_{\mathrm{audit}}}\) is \(\texttt{action}\) and the pair observable \(\mathsf{PairObs}_{\{3,6\}}^{\mu_{\mathrm{drive}}}\) on the same two primitives is \(\texttt{blocked}\);
  \end{enumerate}
\item
  \begin{enumerate}
  \def\labelenumi{(\alph{enumi})}
  \setcounter{enumi}{2}
  \tightlist
  \item
    the instance obtained by joining the records of the instance of Theorem \hyperref[cl:thm:22]{22} and of the instance of (a): its bridge is \(\texttt{strict}\), and its only directed cell, \(C_{\mathrm{thr}}\) on the pair \((6,3)\), is \(\texttt{below\_threshold}\), so that no directed cell of the instance is \(\texttt{action}\);
  \end{enumerate}
\item
  \begin{enumerate}
  \def\labelenumi{(\alph{enumi})}
  \setcounter{enumi}{3}
  \tightlist
  \item
    for the fourth display, the instance of (c) extended, with fresh identifiers, by \(I^{+}_{\mathrm{prom}}\), by the visibility, audit and gate records of the bridge of Theorem \hyperref[cl:thm:22]{22} reissued under \(I^{+}_{\mathrm{prom}}\), and by \(\mathrm{Sound}(I^{+}_{\mathrm{prom}})\), with host tag \(\mathsf H_{\mathrm{fin}}\) and level profile \((\mathsf{Ver},\mathsf{Ver})\), \(\mathsf{ClaimLog}(I^{+}_{\mathrm{prom}})\) and the verdict references of its use-set, all in the scope and visible set of \(I^{+}_{\mathrm{prom}}\); there the copied bridge is again \(\texttt{strict}\) and \(\mathrm{Sound}(I_{\mathrm{prom}}^{+})\), with a fresh \(\texttt{committed\_state}\) evidence source admitted by \(I_{\mathrm{prom}}^{+}\) and an audit record whose checks pass, receives exactly \(\texttt{undefined\_circular}\). For the variant with \(I_{\mathrm{prom}}\) in place of \(I^{+}_{\mathrm{prom}}\), extend the instance of (c) instead, with fresh identifiers, by \(\mathrm{Sound}(I_{\mathrm{prom}})\) of Subsubsection \hyperref[sn:11.5.1]{11.5.1}, \(\mathsf{ClaimLog}(I_{\mathrm{prom}})\) and the verdict references of its use-set, each outside the scope of \(I_{\mathrm{prom}}\); the bridge stays \(\texttt{strict}\) and \(\mathrm{Sound}(I_{\mathrm{prom}})\) receives \(\texttt{outside\_scope}\);
  \end{enumerate}
\item
  \begin{enumerate}
  \def\labelenumi{(\alph{enumi})}
  \setcounter{enumi}{4}
  \tightlist
  \item
    the instance of Theorem \hyperref[cl:thm:22]{22} extended, with fresh identifiers, by one claim record of type promotion in \(\mathsf{ClaimLog}(I_{\mathrm{prom}})\) whose recorded verdict differs from the status the claim classifier of \(I_{\mathrm{prom}}\) assigns it, and by the records of Theorem \hyperref[cl:thm:25]{25} for the chain \((I_{\mathrm{prom}})\) with \(N=0\); the bridge stays \(\texttt{strict}\), \(\delta_{\mathrm{status}}\) has an entry with index \(0\), and by clause 4 of Theorem \hyperref[cl:thm:25]{25} \(\mathrm{Sound}^{\uparrow}_0\) is not \(\texttt{accepted}\). With an empty claim log the same clause accepts \(\mathrm{Sound}^{\uparrow}_0\), so the right judgment of the fifth case can hold.
  \end{enumerate}
\end{itemize}

By Theorem \hyperref[cl:thm:24]{24}, in every admissible instance containing a strict bridge and a candidate record \(\mathrm{Sound}(I)\) for the instrument \(I\) that classifies it, that record is not \(\texttt{accepted}\) under \(I\).

\end{claimtheorem}

\begin{proof}

Each non-implication is established by exhibiting finite admissible records, or a finite admissible record schema, satisfying the antecedent without satisfying the consequent. The channel and promotion separations use the classifier-input distinctions of Sections \hyperref[sn:4]{4} and \hyperref[sn:5]{5}; the cell-to-pair separation is the countermodel-backed case, with the explicit construction recorded as \(\mathsf{CM}_4\) in Section \hyperref[sn:9]{9} and Appendix \hyperref[sn:C]{C}.

\textbf{Channel does not imply cell.} Let \(C_{\mathrm{thr}}\) be the running-example cell with its exact threshold replaced by the threshold \(\Theta_{\ge2}\) on the datum \(|E_1E_2(\{a\})\,\triangle\,E_2E_1(\{a\})|\) (the two recomputed images differ in at least two elements). The images \(\{a,b\}\) and \(\{a,b,c\}\) differ only in \(c\), so this threshold is not met. Keep the underlying maps, witness, update, profile and visibility data. The audit's threshold check, recomputed for the new threshold, fails, while its source, visibility, circularity, nonclaim and other checks pass as before; so the computed audit defect of Subsubsection \hyperref[sn:5.3.6]{5.3.6} is \(\delta_6^{\mathrm{audit}}=\{\texttt{threshold\_failure}\}\), and the visibility defects remain empty. This is a well-formed directed cell \(P_6\leftarrow P_3\) whose informant witness \(W_3\in\mathsf{Wit}(P_3)\) is present but fails a threshold check of its threshold record. The rows of Subsubsection \hyperref[sn:5.6.6]{5.6.6} above \(\texttt{below\_threshold}\) do not apply: the cell is in scope, as the running example is; \(W_3\) and \(U_6\) are present; it uses no bridge, has empty overread and weak-bridge defects, and its audit defect contains neither \(\texttt{source\_failure}\) nor \(\texttt{visibility\_failure}\); its audit defect does not contain \(\texttt{circular\_audit}\); and \(\delta\) contains no package-collapse value. The \(\texttt{below\_threshold}\) row holds, since \(W_3\) and \(U_6\) are present and \(\texttt{threshold\_failure}\in\delta_6^{\mathrm{audit}}\); so \(C_{\mathrm{thr}}\) is \(\texttt{below\_threshold}\), not \(\texttt{action}\). In the instance whose only directed cell is \(C_{\mathrm{thr}}\), take the role-channel record \((\mathrm{Chan}_6,R_6,\Theta_6^{\mathrm{act}},A_{R_6},V_{R_6},\mathrm{false})\) of Subsubsection \hyperref[sn:5.6.1]{5.6.1}, whose role evidence is the source-of-truth record \(R_6\in\mathsf{Wit}(P_6)\) of \(E_1\) and \(E_2\) to which the audit record refers, with activation threshold \(\Theta_6^{\mathrm{act}}\) the instance of \(\Theta_{\mathrm{source}}\) on \(R_6\), which passes since \(\delta_6^{\mathrm{source}}(R_6;I)=\texttt{committed}\) under the default entry. This record is unchanged by the replacement: it is in scope, visible, audited under \(\mathsf{AuditPolicy}_I\), its declared field \(\mathsf{collapse}_6\) is false, and it satisfies the source policy. By Subsubsection \hyperref[sn:5.6.1]{5.6.1} the role channel of \(P_6\) is \(\texttt{active\_projection}\). Thus an active role channel does not by itself produce a directed cell with status \(\texttt{action}\); the channel and the cell are adjudicated by separate classifiers with separate inputs.

\textbf{Cell does not imply pair.} Take a directed cell \(\mathsf{Cell}_{ij}^{\lambda}\) that is well-formed and classified \(\texttt{action}\) under \(\Gamma,\mathcal T,I\). Now consider the unordered pair observable \(\mathsf{PairObs}_{\{i,j\}}^{\mu}\) on the same primitives, with a drive profile, whose source slot is an audited absence (\(\delta_6^{\mathrm{source}}=\texttt{missing}\)) and which has no drive bridge (\(\delta_{\mathrm{reqbridge}}=\{\texttt{drive}\}\)) (either defect alone triggers the first \(\texttt{blocked}\) row), as in \(\mathsf{CM}_4\). The cell is \(\texttt{action}\), but the source defect of the pair record fails its threshold, so the first condition of the \(\texttt{blocked}\) row of the pair-observable table (Subsubsection \hyperref[sn:5.6.6]{5.6.6}) holds and the pair classifier assigns the pair record status \(\texttt{blocked}\). Countermodel \(\mathsf{CM}_4\) of Section \hyperref[sn:9]{9} supplies the explicit construction; the cell-to-pair non-implication is therefore witnessed by a finite admissible pair of records.

\textbf{Strict promotion does not imply cell action.} Take a promotion bridge \(B_{j\to j+1}\) classified \(\texttt{strict}\) under \(\Gamma,\mathcal T_j,I\) --- that is, an admissible bridge whose strictness gate \(G_{\mathrm{strict}}\) is \(\texttt{strict\_pass}\), equivalently \(\Delta_{\mathrm{fact}}(\pi_j,\pi_{j+1})\neq\varnothing\). Strictness is a property of the typed object-map data alone (Subsubsection \hyperref[sn:5.4.1]{5.4.1}) and does not by itself produce a directed cell of any specific primitive pair. A directed-cell record on the target package, with its own witness, update, profile, threshold, audit, visibility, and defect data, is required to obtain a \(\texttt{action}\) verdict on a cell. For case (c), join the records of the instance of Theorem \hyperref[cl:thm:22]{22} and of the instance of case (a); the two record sets are disjoint and each judgment keeps its own instrument. In the joined instance, give the records of the two instances distinct identifiers and leave \(I\) and \(I_{\mathrm{prom}}\) unchanged; since visibility is determined by scope (Subsubsection \hyperref[sn:3.5.2]{3.5.2}), the Theorem \hyperref[cl:thm:22]{22} records are \(\texttt{outside}\) for \(I\) and the records of case (a) are \(\texttt{outside}\) for \(I_{\mathrm{prom}}\); no record used by either judgment changes class, so every row computation is unchanged, and the bridge remains \(\texttt{strict}\) by Theorem \hyperref[cl:thm:22]{22} and the cell \(C_{\mathrm{thr}}\) remains \(\texttt{below\_threshold}\) by case (a). Hence strict promotion does not imply directed-cell action.

\textbf{Strict promotion does not imply acceptance of a same-level soundness claim.} By Theorem \hyperref[cl:thm:24]{24}, which applies to every admissible instrument, neither \(\mathrm{Sound}(I_{\mathrm{prom}})\) nor \(\mathrm{Sound}(I^{+}_{\mathrm{prom}})\) is accepted. The same claim-classifier rows are checked directly for the Toy22 instruments below. These constructions and Theorem \hyperref[cl:thm:22]{22} do not use Theorem \hyperref[cl:thm:5]{5}. In the Theorem \hyperref[cl:thm:22]{22} instance, \(I_{\mathrm{prom}}\) classifies the bridge as \(\texttt{strict}\), while the candidate same-level claim \(\mathrm{Sound}(I_{\mathrm{prom}})\) receives \(\texttt{outside\_scope}\): its claim-type tag \(\texttt{soundness}\) is excluded from \(\mathsf{ClaimTypes}_{I_{\mathrm{prom}}}\). This follows directly from the first row of the claim classifier. Take \(I_{\mathrm{prom}}^+\) with a fresh identifier, copying the finite promotion policies of \(I_{\mathrm{prom}}\) and extending the required-strength map by \(\mathsf{ReqStrength}_{I_{\mathrm{prom}}^+}(\texttt{soundness})=\texttt{weak}\), adding \(\texttt{soundness}\) to its claim types, and the self-soundness use-set together with the reissued bridge record and its visibility, audit and gate records to its scope and visible set. Reissue, with a fresh bridge identifier, the Toy22 bridge's instrument-indexed visibility, audit and gate records under \(I_{\mathrm{prom}}^+\), and recompute their checks; the same finite tables give passing core gates and \(\texttt{strict\_pass}\), so the copied bridge is \(\texttt{strict}\) under \(I_{\mathrm{prom}}^+\). The soundness claim contains its own verdict reference and receives \(\texttt{undefined\_circular}\) unless a higher-priority row applies. In neither case is the claim \(\texttt{accepted}\).

For \(I_{\mathrm{prom}}^{+}\) in (d), the host tag \(\mathsf H_{\mathrm{fin}}\) and level tag \(\mathsf{Ver}\) of \(\mathrm{Sound}(I_{\mathrm{prom}}^{+})\) lie in \(\mathsf{Hosts}\) and \(\mathsf{Levels}\) of \(I_{\mathrm{prom}}^{+}\), which are those of \(I_{\mathrm{prom}}\), the use-set of \(\mathrm{Sound}(I_{\mathrm{prom}}^{+})\) is in scope and visible, its evidence resolves to an admitted source, no record is absent and its audit defect is empty, so the rows above \(\texttt{undefined\_circular}\) fail; the third condition of the \(\texttt{undefined\_circular}\) row holds by membership. For the stronger form, the dependency is judged, as in Subsubsection \hyperref[sn:9.1.1]{9.1.1}, by a fresh local copy of \(I_{\mathrm{prom}}^{+}\) that also admits the claim type \(\texttt{dependency}\); since \(\mathrm{Sound}(I_{\mathrm{prom}}^{+})\) is accepted in no admissible instance (Theorem \hyperref[cl:thm:24]{24}), every instance of its class \(\Sigma\) satisfying the left side fails the right side.

\textbf{Strict promotion does not imply accepted lifted soundness.} For (e), give the promotion-type claim a fresh \(\texttt{committed\_state}\) evidence source outside the scope of \(I_{\mathrm{prom}}\), complete typed fields and correctly recomputed audit defects. Its classifier status under \(I_{\mathrm{prom}}\) is then \(\texttt{outside\_scope}\); record \(\texttt{accepted}\) for it in \(\mathsf{ClaimLog}(I_{\mathrm{prom}})\). The bridge records are unchanged, so the bridge stays \(\texttt{strict}\). The one-instrument chain \((I_{\mathrm{prom}})\) satisfies the hypotheses of Theorem \hyperref[cl:thm:25]{25} with \(N=0\), since \(\mathsf{Level}(I_{\mathrm{prom}})=\mathsf{Ver}\) lies below \(\mathsf{Str}\); the log entry puts \((0,\varphi)\) in \(\delta_{\mathrm{status}}\), so clause 4 rejects acceptance of \(\mathrm{Sound}^{\uparrow}_0\). With an empty log, \(\delta_{\mathrm{status}}\) has no entry with index \(0\) and the same clause accepts \(\mathrm{Sound}^{\uparrow}_0\).

The gate and square families of Subsubsections \hyperref[sn:5.1.6]{5.1.6} and \hyperref[sn:5.1.7]{5.1.7} are likewise distinct from these five: their classifier inputs are gate records (per the gate semantics of Section \hyperref[sn:10]{10}) and square fillers (per the high-structure semantics of Section \hyperref[sn:14]{14}), not the cell, pair, promotion, or claim records.

\end{proof}

\subsubsection{Consequence}\label{sn:6.5.1}

In the running example of Subsection \hyperref[sn:3.8]{3.8}, the cell \(P_6\leftarrow P_3\) is \(\texttt{action}\), while the pair observable on \(\{P_3,P_6\}\) with a drive profile, whose source slot is an audited absence (\(\delta_6^{\mathrm{source}}=\texttt{missing}\)) and which has no drive bridge (\(\delta_{\mathrm{reqbridge}}=\{\texttt{drive}\}\)) (either defect alone triggers the first \(\texttt{blocked}\) row), is \(\texttt{blocked}\); this is the cell-to-pair separation, constructed as \(\mathsf{CM}_4\). For the separation of strict promotion from cell activity, the bridge of Theorem \hyperref[cl:thm:22]{22} is classified \(\texttt{strict}\) (Subsubsection \hyperref[sn:10.6.2]{10.6.2}), and nothing in its record asserts an active directed cell.

The paper does not infer pair-observable realness from directed-cell action, promotion strictness from package or channel activity, or claim acceptance from raw visibility alone. The status discipline of the calculus is the working content of the separation theorem: each status family is consulted only for its own judgments, and a status-level inference between families is recorded only through an explicit bridge or theorem. NC-3, NC-4, NC-5, and NC-6 of Subsection \hyperref[sn:15.2]{15.2} are the corresponding nonclaims, and the no-overreading theorem of Section \hyperref[sn:11]{11} gives the claim-level analogue at theorem-grade.
 \section{Descent, Packaging, Completion, and Strictness}\label{sn:7}

This section proves the finite theorems on descent, packaging, completion and strictness. Theorem \hyperref[cl:thm:6]{6} gives the descent condition for a finite update through a quotient or finite map: an update induces a map on classes exactly when it respects the fibres. This is the universal property of a finite quotient, the set-level analogue of the descent question of \citet{Grothendieck1959Descent}, not an instance of that theory. Theorem \hyperref[cl:thm:7]{7} shows that the update descends exactly when its descent defect, the set of split pairs, is empty. Theorem \hyperref[cl:thm:8]{8} is the Markov closure-deficit theorem, the only theorem in this paper that lives on the stochastic host \(\mathsf H_{\mathrm{prob}}\); its information-theoretic quantities are those of \citet{CoverThomas2006}. Theorems \hyperref[cl:thm:9]{9} and \hyperref[cl:thm:10]{10} record noncommuting completions and the resulting pasting defect. Theorem \hyperref[cl:thm:11]{11} is idempotent saturation. Theorems \hyperref[cl:thm:12]{12} and \hyperref[cl:thm:13]{13} record strict-extension nonfactorization and its defect equivalence. Theorem \hyperref[cl:thm:14]{14} is the fixed-interface definability bound.

\subsection{Theorem 6: Descent Through Quotient / Finite Map}\label{sn:7.1}

\begin{claimtheorem}{Theorem 6 (DescentThroughQuotient)}\phantomsection\label{cl:thm:6}

Let \(X\) and \(Y\) be finite sets on the host \(\mathsf H_{\mathrm{fin}}\), let \(q:X\to Y\) be a finite package, quotient, or abstraction map, and let \(F:X\to X\) be a finite update. Let \(q_{\mathrm{im}}:X\to\mathrm{im}(q)\) denote the effective-image map. There exists a finite descended update \(F^\sharp:\mathrm{im}(q)\to\mathrm{im}(q)\) with

\[
F^\sharp\,q_{\mathrm{im}}\;=\;q_{\mathrm{im}}\,F
\]

if and only if

\[
\forall\,x,x'\in X,\quad q(x)=q(x')\;\Longrightarrow\;qF(x)=qF(x').
\]

Equivalently, for a quotient \(X/\!\sim\) with quotient map \(q\),

\[
\exists\,F^\sharp\colon\;qF=F^\sharp q\;\;\Longleftrightarrow\;\;\bigl(x\sim x'\;\Longrightarrow\;F(x)\sim F(x')\bigr).
\]
The equivalence, with the same proof, holds for arbitrary sets \(X,Y\) and arbitrary maps \(q,F\); finiteness only makes the descent condition a finite check.

\end{claimtheorem}

\subsubsection{Data}\label{sn:7.1.1}

The hypotheses are: a finite set \(X\) on the host \(\mathsf H_{\mathrm{fin}}\), a finite map \(q:X\to Y\) with effective image \(\mathrm{im}(q)\subseteq Y\), a finite update \(F:X\to X\). The conclusion quantifies over candidate descended updates \(F^\sharp:\mathrm{im}(q)\to\mathrm{im}(q)\); none is given in advance. The judgment is registered under an admissible finite profile whose lens data include the effective-image map \(q_{\mathrm{im}}\).

\begin{proof}

For the forward direction, suppose \(F^\sharp\) exists with \(F^\sharp q_{\mathrm{im}}=q_{\mathrm{im}}F\). Take \(x,x'\in X\) with \(q(x)=q(x')\). Then \(q_{\mathrm{im}}(x)=q_{\mathrm{im}}(x')\), so

\[
q_{\mathrm{im}}F(x)\;=\;F^\sharp q_{\mathrm{im}}(x)\;=\;F^\sharp q_{\mathrm{im}}(x')\;=\;q_{\mathrm{im}}F(x').
\]

Taking the underlying values in the subset \(\mathrm{im}(q)\subseteq Y\) gives \(qF(x)=qF(x')\), establishing the fiber-constancy condition.

For the reverse direction, suppose \(q(x)=q(x')\Rightarrow qF(x)=qF(x')\). Define \(F^\sharp\) on \(\mathrm{im}(q)\) by

\[
F^\sharp(q_{\mathrm{im}}(x))\;=\;q_{\mathrm{im}}(F(x)),\qquad x\in X.
\]

The fiber-constancy hypothesis ensures this is well-defined: if \(q(x)=q(x')\), then \(qF(x)=qF(x')\), so the value does not depend on the chosen representative in the fiber. Restricting the codomain to \(\mathrm{im}(q)\) gives a finite map \(F^\sharp:\mathrm{im}(q)\to\mathrm{im}(q)\), and the relation \(F^\sharp q_{\mathrm{im}}=q_{\mathrm{im}}F\) holds by construction.

Both directions are finite checks on finite sets, so the proof is finitary throughout.

\end{proof}

\subsubsection{Role Interpretation}\label{sn:7.1.2}

The theorem records the \(P_1/P_5/P_6\) descent condition: \(P_1\) supplies the update \(F\), \(P_5\) supplies the package \(q\), and \(P_6\) records the audit trail under which the descended map \(F^\sharp\) is accepted. Descent through \(q\) is admissible exactly when the descent defect \(\delta_1^{\mathrm{split}}(q,F)\) of Subsubsection \hyperref[sn:5.3.1]{5.3.1} vanishes; Theorem \hyperref[cl:thm:7]{7} states this equivalence at the defect level.

\subsection{Theorem 7: Descent Defect Equivalence}\label{sn:7.2}

\begin{claimtheorem}{Theorem 7 (DescentDefectEquivalence)}\phantomsection\label{cl:thm:7}

With \(X,Y,q,F\) as in Theorem \hyperref[cl:thm:6]{6}, define the descent defect

\[
\delta_1^{\mathrm{split}}(q,F)\;=\;\{\,(x,x')\in X\times X\;:\;q(x)=q(x')\;\wedge\;qF(x)\neq qF(x')\,\}.
\]

Then

\[
\delta_1^{\mathrm{split}}(q,F)\;=\;\varnothing\;\;\Longleftrightarrow\;\;F\;\text{descends through}\;q,
\]

and equivalently

\[
\delta_1^{\mathrm{split}}(q,F)\;\neq\;\varnothing\;\;\Longleftrightarrow\;\;\nexists\,F^\sharp\colon F^\sharp\,q_{\mathrm{im}}=q_{\mathrm{im}}\,F.
\]

With the same proof, this holds for arbitrary sets and maps; finiteness only makes it a finite check.

\end{claimtheorem}

\begin{proof}

By definition, \(\delta_1^{\mathrm{split}}(q,F)\) collects exactly the pairs \((x,x')\) that violate the fiber-constancy condition of Theorem \hyperref[cl:thm:6]{6}. The set is empty iff the fiber-constancy condition holds for every pair, which by Theorem \hyperref[cl:thm:6]{6} is equivalent to the existence of a descended update \(F^\sharp\). The set is nonempty iff there exists a violating pair, which by Theorem \hyperref[cl:thm:6]{6} is equivalent to the nonexistence of any \(F^\sharp\).

Both directions are finite checks on the finite set \(X\times X\).

The theorem makes the descent obstruction a finite typed record: a non-empty \(\delta_1^{\mathrm{split}}\) is the witness used by the descent gate of Subsubsection \hyperref[sn:5.4.4]{5.4.4} in promotion records that claim macro closure on the target package, and a vanishing \(\delta_1^{\mathrm{split}}\) is the corresponding passing condition for that gate.

\end{proof}

\subsection{Theorem 8: Markov Closure-Deficit Theorem}\label{sn:7.3}

In this subsection only, \(P\) denotes a probability and \(I(\cdot\,;\cdot\mid\cdot)\) a conditional mutual information; neither refers to a role label or an instrument.

\begin{claimtheorem}{Theorem 8 (FiniteMarkovClosureDeficit)}\phantomsection\label{cl:thm:8}

On the finite stochastic host \(\mathsf H_{\mathrm{prob}}\), fix a finite state space \(X\), a stochastic matrix \(M\) on \(X\), a distribution \(\mu\) on \(X\), a finite abstraction \(\Pi:X\to Y\) to a finite set \(Y\), and a lag \(\tau\in\mathbb N_{\ge1}\). Let the pair \((X_t,X_{t+\tau})\) have the joint law \(P(X_t=x,\ X_{t+\tau}=x')=\mu(x)\,M^\tau(x,x')\), and write \(Y_t=\Pi(X_t)\) and \(Y_{t+\tau}=\Pi(X_{t+\tau})\); all probabilities and expectations below are taken under this law. For a stochastic kernel \(K:Y\to\Delta(Y)\), the KL predictive closure loss is
\[
\mathcal L_\tau(K)=\mathbb E\Bigl[\,D_{\mathrm{KL}}\bigl(P(Y_{t+\tau}\mid X_t)\,\big\|\,K(\,\cdot\,\mid Y_t)\bigr)\Bigr],
\]
where \(D_{\mathrm{KL}}(p\,\|\,q)=\sum_{y}p(y)\log\bigl(p(y)/q(y)\bigr)\), a term with \(p(y)=0\) is \(0\), and a term with \(p(y)>0=q(y)\) is \(+\infty\); the expectation is \(\sum_{x:\mu(x)>0}\mu(x)\,D_{\mathrm{KL}}\bigl(P(Y_{t+\tau}\mid X_t=x)\,\big\|\,K(\,\cdot\,\mid\Pi(x))\bigr)\), so states of zero \(\mu\)-mass do not contribute. Let \(\mathcal K_{\mathrm{adm}}\) be a declared class of such kernels containing a kernel \(K^\ast\) with \(K^\ast(\cdot\mid y)=P(Y_{t+\tau}\mid Y_t=y)\) for every \(y\) in the support of \(Y_t\). Then the minimum predictive closure loss over admissible macro kernels exists, is attained at \(K^\ast\), and is

\[
\min_{K\in\mathcal K_{\mathrm{adm}}}\;\mathcal L_\tau(K)\;=\;I\bigl(X_t\,;\,\Pi(X_{t+\tau})\,\big|\,\Pi(X_t)\bigr),
\]

where \(I(\,\cdot\,;\,\cdot\,|\,\cdot\,)\) is the conditional mutual information. Moreover, for every kernel \(K:Y\to\Delta(Y)\),

\[
\mathcal L_\tau(K)\;=\;I\bigl(X_t\,;\,\Pi(X_{t+\tau})\,\big|\,\Pi(X_t)\bigr)\;+\;\sum_{y\in\operatorname{supp}(Y_t)}P(Y_t=y)\,D_{\mathrm{KL}}\bigl(P(Y_{t+\tau}\mid Y_t=y)\,\big\|\,K(\,\cdot\,\mid y)\bigr),
\]

where the equality is in the extended nonnegative reals: if \(K\) gives zero mass to an outcome of positive supported conditional probability, both sides are \(+\infty\). So for any class \(\mathcal K\) of kernels \(\inf_{K\in\mathcal K}\mathcal L_\tau(K)\ge I\bigl(X_t\,;\,\Pi(X_{t+\tau})\,\big|\,\Pi(X_t)\bigr)\), and the minimum over \(\mathcal K\) exists and equals this conditional mutual information exactly when \(\mathcal K\) contains a kernel that agrees with \(K^\ast\) on the support of \(Y_t\). In particular \(I(X_t;\Pi(X_{t+\tau})\mid\Pi(X_t))=0\) if and only if, for every \(x\) with \(\mu(x)>0\), the law \(\sum_{\Pi(x')=\cdot}M^\tau(x,x')\) depends on \(x\) only through \(\Pi(x)\); that is, the \(\tau\)-step kernel \(M^\tau\) satisfies the Kemeny--Snell lumpability condition for \(\Pi\) at every state of positive \(\mu\)-mass (implied by Kemeny--Snell lumpability of \(M\) for \(\Pi\), and in general strictly weaker, since it constrains only \(M^\tau\) and only states of positive \(\mu\)-mass), and then \(\mathcal L_\tau(K^\ast)=0\).

\end{claimtheorem}

\subsubsection{Assumptions}\label{sn:7.3.1}

The theorem requires the following finite assumptions, all carried explicitly on the judgment line:

\begin{itemize}
\tightlist
\item
  \(\mathsf H_{\mathrm{prob}}\) is the host: a finite state space \(X\) with a stochastic matrix \(M\) and a distribution \(\mu\), which together fix the joint law of \((X_t,X_{t+\tau})\).
\item
  \(\Pi:X\to Y\) is a finite abstraction to a finite macro space \(Y\).
\item
  \(\tau\) is a positive integer lag, and \(P(Y_{t+\tau}=y'\mid X_t=x)=\sum_{x':\Pi(x')=y'}M^\tau(x,x')\).
\item
  \(\mathcal L_\tau\) is the declared closure loss, the KL predictive loss of the statement.
\item
  The admissible macro-kernel class \(\mathcal K_{\mathrm{adm}}\) is a declared class of stochastic kernels \(K:Y\to\Delta(Y)\) over which the minimization is taken, and it contains \(K^\ast\); values of \(K^\ast(\cdot\mid y)\) for \(y\) outside the support of \(Y_t\) do not affect \(\mathcal L_\tau\).
\end{itemize}

The theorem is theorem-grade only under these declared assumptions; it is not a universal stochastic-emergence statement.

\begin{proof}

Write \(Y_t=\Pi(X_t)\) and \(Y_{t+\tau}=\Pi(X_{t+\tau})\). For each macro state \(y\in Y\) in the support of \(Y_t\), decompose the conditional KL loss into a sum of two non-negative terms:

\[
\begin{aligned}
&\mathbb E\Bigl[\,D_{\mathrm{KL}}\bigl(P(Y_{t+\tau}\mid X_t)\,\big\|\,K(\,\cdot\,\mid y)\bigr)\;\Big|\;Y_t=y\,\Bigr]\\
&\qquad=\;\mathbb E\Bigl[\,D_{\mathrm{KL}}\bigl(P(Y_{t+\tau}\mid X_t)\,\big\|\,P(Y_{t+\tau}\mid Y_t=y)\bigr)\;\Big|\;Y_t=y\,\Bigr]\\
&\qquad\quad+\;D_{\mathrm{KL}}\bigl(P(Y_{t+\tau}\mid Y_t=y)\,\big\|\,K(\,\cdot\,\mid y)\bigr).
\end{aligned}
\]

The decomposition is the chain rule for relative entropy: since \(Y_t\) is a function of \(X_t\), \(\mathbb E[P(Y_{t+\tau}\mid X_t)\mid Y_t=y]=P(Y_{t+\tau}\mid Y_t=y)\), and the cross term is the expectation of \(\log\bigl(P(Y_{t+\tau}\mid Y_t=y)/K(\cdot\mid y)\bigr)\) under that averaged law. When \(K(\cdot\mid y)\) vanishes where \(P(Y_{t+\tau}\mid Y_t=y)\) does not, both sides are \(+\infty\). The first term does not depend on the choice of \(K\). The second term is non-negative and is minimized by the admissible choice \(K^\ast(\,\cdot\,\mid y)\;=\;P(Y_{t+\tau}\mid Y_t=y)\), at which it vanishes. Taking expectation over \(Y_t=y\) and minimizing over \(K\in\mathcal K_{\mathrm{adm}}\),

\[
\min_{K\in\mathcal K_{\mathrm{adm}}}\;\mathcal L_\tau(K)\;=\;\mathbb E\Bigl[\,D_{\mathrm{KL}}\bigl(P(Y_{t+\tau}\mid X_t)\,\big\|\,P(Y_{t+\tau}\mid Y_t)\bigr)\,\Bigr]\;=\;I\bigl(X_t\,;\,Y_{t+\tau}\,\big|\,Y_t\bigr).
\]

Substituting back \(Y_t=\Pi(X_t)\) and \(Y_{t+\tau}=\Pi(X_{t+\tau})\) gives the stated equality. Taking the expectation of the decomposition without minimizing gives the displayed identity for every kernel \(K\), since the expectation of the first term is the conditional mutual information. The second term is nonnegative, which gives the lower bound for any class; it vanishes exactly when \(K(\cdot\mid y)=P(Y_{t+\tau}\mid Y_t=y)\) for every \(y\) in the support of \(Y_t\), which gives the characterization of attainment. Support-zero cases are handled by restricting to the support of the finite distribution; the minimizer is \(K^\ast\) on that support, and no convexity of \(\mathcal K_{\mathrm{adm}}\) is used.

For the lumpability clause, the conditional mutual information equals \(\sum_{\mu(x)>0}\mu(x)D_{\mathrm{KL}}\bigl(P(Y_{t+\tau}\mid X_t=x)\,\|\,P(Y_{t+\tau}\mid Y_t=\Pi(x))\bigr)\), a sum of nonnegative terms; it vanishes if and only if each conditional law equals its fibre average, which is the lumpability condition.

\end{proof}

\subsubsection{Caution}\label{sn:7.3.2}

The theorem is assumption-scoped: the equality holds only when the host is \(\mathsf H_{\mathrm{prob}}\) and the assumptions of Subsubsection \hyperref[sn:7.3.1]{7.3.1} are met. It is not a universal stochastic-emergence theorem and does not by itself imply package-implies-closure or strictness-implies-closure; NC-7 of Subsection \hyperref[sn:15.2]{15.2} records the corresponding nonclaim. The closure deficit is a finite information-theoretic quantity on a declared finite Markov host; its interpretation outside that host requires a separate bridge.

\subsection{Theorem 9: Noncommuting Completions}\label{sn:7.4}

\begin{claimtheorem}{Theorem 9 (NoncommutingCompletions)}\phantomsection\label{cl:thm:9}

There exist finite idempotent completion operators \(E_1,E_2:C\to C\) on the finite power-set lattice \(C=\mathcal P(U)\) of a finite set \(U\), which are closure operators on \((C,\subseteq)\) (extensive, monotone and idempotent), with

\[
E_1^2\;=\;E_1,\qquad E_2^2\;=\;E_2,
\]

but

\[
E_1\,E_2\;\neq\;E_2\,E_1.
\]

Equivalently, there exists \(S\in C\) with \(E_1E_2(S)\neq E_2E_1(S)\).

\end{claimtheorem}

\subsubsection{Construction}\label{sn:7.4.1}

Take \(U=\{a,b,c\}\) and \(C=\mathcal P(U)\), the finite power set of \(U\). Define

\[
E_1(S)\;=\;\begin{cases}\,S\cup\{b\}, & a\in S,\\[2pt] S, & a\notin S,\end{cases}\qquad E_2(S)\;=\;\begin{cases}\,S\cup\{c\}, & b\in S,\\[2pt] S, & b\notin S.\end{cases}
\]

Each map is a finite endomap of \(C\).

\textbf{Extensivity and monotonicity.} \(S\subseteq E_1(S)\) is clear. If \(S\subseteq T\) and \(a\in S\), then \(a\in T\) and \(E_1(S)=S\cup\{b\}\subseteq T\cup\{b\}=E_1(T)\); if \(a\notin S\), then \(E_1(S)=S\subseteq T\subseteq E_1(T)\). The argument for \(E_2\) is the same with \(b\) in place of \(a\) and \(c\) in place of \(b\).

\textbf{Idempotence.} For \(E_1\): if \(a\in S\), then \(E_1(S)=S\cup\{b\}\) still contains \(a\), so \(E_1(E_1(S))=(S\cup\{b\})\cup\{b\}=S\cup\{b\}=E_1(S)\). If \(a\notin S\), then \(E_1(S)=S\), hence \(E_1(E_1(S))=E_1(S)\). So \(E_1^2=E_1\). The argument for \(E_2\) is symmetric (with \(b\) in place of \(a\) and \(c\) in place of \(b\)): \(E_2^2=E_2\).

\textbf{Noncommutation.} Take \(S_0=\{a\}\). Then

\[
E_1\,E_2(\{a\})\;=\;E_1(\{a\})\;=\;\{a,b\},
\]

since \(E_2\) leaves \(\{a\}\) fixed (because \(b\notin\{a\}\)). On the other hand,

\[
E_2\,E_1(\{a\})\;=\;E_2(\{a,b\})\;=\;\{a,b,c\},
\]

since \(E_1\) adds \(b\) (because \(a\in\{a\}\)) and then \(E_2\) adds \(c\) (because \(b\in\{a,b\}\)). Hence \(E_1E_2(\{a\})\neq E_2E_1(\{a\})\).

\begin{proof}

The construction above is a finite witness: \(E_1\) and \(E_2\) are finite idempotent endomaps of a finite carrier \(C\) with \(E_1E_2(\{a\})\neq E_2E_1(\{a\})\). The theorem is established.

\end{proof}

\subsubsection{Interpretation}\label{sn:7.4.2}

Route mismatch is a real defect even when each individual completion is exact and idempotent. The result rules out the overclaim targeted by NC-8 (no idempotence-implies-commutation) of Subsection \hyperref[sn:15.2]{15.2} at theorem-grade: \(E_1^2=E_1\wedge E_2^2=E_2\) does not imply \(E_1E_2=E_2E_1\). Theorem \hyperref[cl:thm:10]{10} will record the equivalence between non-vanishing route-mismatch defect \(\Delta_3^{\mathrm{comp}}(E_1,E_2)\) and noncommutation; the construction here is the finite witness behind that equivalence.

\subsection{Theorem 10: Completion-Pasting Defect}\label{sn:7.5}

\begin{claimtheorem}{Theorem 10 (CompletionPastingDefect)}\phantomsection\label{cl:thm:10}

For any endomaps \(E_1,E_2:C\to C\) of a finite carrier \(C\) (in particular finite completion endomaps), define the route-mismatch defect

\[
\Delta_3^{\mathrm{comp}}(E_1,E_2)\;=\;\{\,c\in C\;:\;E_1E_2(c)\;\neq\;E_2E_1(c)\,\}.
\]

Then

\[
\Delta_3^{\mathrm{comp}}(E_1,E_2)\;\neq\;\varnothing\;\;\Longleftrightarrow\;\;E_1E_2\;\neq\;E_2E_1.
\]

The defect \(\Delta_3^{\mathrm{comp}}\) is a \(P_3\)-route/pasting defect (Subsubsection \hyperref[sn:5.3.3]{5.3.3}); when \(E_1,E_2\) are recorded completion endomaps, their completion records supply the \(P_5\) witnesses (Subsubsection \hyperref[sn:5.3.5]{5.3.5}).

With the same proof, this holds for arbitrary sets and maps; finiteness only makes it a finite check.

\end{claimtheorem}

\begin{proof}

By definition, \(\Delta_3^{\mathrm{comp}}(E_1,E_2)\) collects exactly the points \(c\in C\) at which the two composites disagree.

If \(\Delta_3^{\mathrm{comp}}=\varnothing\), then for every \(c\in C\), \(E_1E_2(c)=E_2E_1(c)\), so \(E_1E_2=E_2E_1\) as maps on \(C\).

If \(\Delta_3^{\mathrm{comp}}\neq\varnothing\), then there exists some \(c\in C\) with \(E_1E_2(c)\neq E_2E_1(c)\), so \(E_1E_2\neq E_2E_1\) as maps on \(C\).

Both directions are finite checks on the finite carrier \(C\).

Together with Theorem \hyperref[cl:thm:9]{9}, the result records that the route-mismatch defect is a finite typed record: any noncommutation of finite completions is witnessed by a non-empty \(\Delta_3^{\mathrm{comp}}\), and the construction of Theorem \hyperref[cl:thm:9]{9} is a finite explicit instance with \(\{a\}\in\Delta_3^{\mathrm{comp}}(E_1,E_2)\). The high-structure pasting-square fragment of Section \hyperref[sn:14]{14} later uses this route-mismatch record as the finite input for decorated-square pasting defects.

\end{proof}

\subsection{Theorem 11: Idempotent Saturation}\label{sn:7.6}

\begin{claimtheorem}{Theorem 11 (IdempotentSaturation)}\phantomsection\label{cl:thm:11}

Let \(X\) be a finite set on \(\mathsf H_{\mathrm{fin}}\), and let \(E:X\to X\) be a finite self-map. If \(E^2=E\), then for every \(n\ge 1\),

\[
E^n\;=\;E.
\]

Pointwise: \(\forall n\ge 1\;\forall x\in X,\;E^n(x)=E(x)\).

With the same proof, this holds for arbitrary sets and maps; finiteness only makes it a finite check.

\end{claimtheorem}

\begin{proof}

Induct on \(n\).

\textbf{Base case} (\(n=1\)). \(E^1=E\) by definition.

\textbf{Inductive step.} Assume \(E^n=E\) for some \(n\ge 1\). Then

\[
E^{n+1}\;=\;E\circ E^n\;=\;E\circ E\;=\;E^2\;=\;E.
\]

At each induction step, the equality \(E^n=E\) is an equality of finite functions on the finite carrier \(X\); the argument is the ordinary induction on positive integers using the idempotence equation \(E^2=E\).

\end{proof}

\subsubsection{Consequence}\label{sn:7.6.1}

Repeated application of a fixed exact completion does not produce theory growth: the saturation \(E,E^2,E^3,\ldots\) collapses to \(E\) once \(E^2=E\). Strict theory growth requires a new object map, bridge, refinement, or nonfactorization witness, not iteration of a fixed completion; NC-9 of Subsection \hyperref[sn:15.2]{15.2} is the corresponding nonclaim. For the directed cell, if \(U\) applies \(E\) to a recorded input \(s\in\mathrm{im}(E)\), then \(E(s)=s\) by idempotence (\(s=E(t)\) gives \(E(s)=E^2(t)=E(t)=s\)), so \(\mathsf{noopU}\) holds, and every cell \(P_5\leftarrow P_j\) carrying this update, whatever its informant, is \(\texttt{trivial}\) in the classifier of Subsubsection \hyperref[sn:5.6.2]{5.6.2} unless a row above it applies.

\subsection{Theorem 12: Strict-Extension Nonfactorization}\label{sn:7.7}

\begin{claimtheorem}{Theorem 12 (StrictExtensionNonfactorization)}\phantomsection\label{cl:thm:12}

Let \(S\) be a finite set on \(\mathsf H_{\mathrm{fin}}\), and let

\[
\pi_0\colon S\to O_0,\qquad \pi_1\colon S\to O_1
\]

be finite object maps with effective codomain \(O_0^{\mathrm{eff}}=\mathrm{im}(\pi_0)\subseteq O_0\). Then \(\pi_1\) does not factor through \(\pi_0\), written \(\pi_1\not\factor\pi_0\), meaning

\[
\nexists\,\phi\colon O_0^{\mathrm{eff}}\to O_1\;\;\text{with}\;\;\pi_1\;=\;\phi\,\pi_0,
\]

if and only if

\[
\exists\,s,s'\in S\;:\;\pi_0(s)=\pi_0(s')\;\wedge\;\pi_1(s)\neq\pi_1(s').
\]

The equivalence, with the same proof, holds for an arbitrary set \(S\) and arbitrary maps \(\pi_0,\pi_1\); finiteness only makes the witness search a finite check.

\end{claimtheorem}

\subsubsection{Data}\label{sn:7.7.1}

The hypotheses are: a finite carrier \(S\), finite object maps \(\pi_0,\pi_1\), the effective codomain \(O_0^{\mathrm{eff}}=\mathrm{im}(\pi_0)\); the conclusion quantifies over candidate factorizations \(\phi:O_0^{\mathrm{eff}}\to O_1\). The theorem treats \(\pi_0,\pi_1\) as the old and new interface of a promotion bridge or extension; the strict-extension condition is the absence of \(\phi\).

\begin{proof}

For the reverse direction, suppose there exist \(s,s'\in S\) with \(\pi_0(s)=\pi_0(s')\) and \(\pi_1(s)\neq\pi_1(s')\). Suppose for contradiction that some \(\phi:O_0^{\mathrm{eff}}\to O_1\) satisfies \(\pi_1=\phi\,\pi_0\). Then

\[
\phi(\pi_0(s))\;=\;\pi_1(s)\;\neq\;\pi_1(s')\;=\;\phi(\pi_0(s')),
\]

contradicting that \(\phi\) is well-defined as a function on the common value \(\pi_0(s)=\pi_0(s')\). Hence no \(\phi\) exists, and \(\pi_1\not\factor\pi_0\).

For the forward direction, argue by contraposition: suppose no such pair \((s,s')\) exists, that is, \(\pi_0(s)=\pi_0(s')\Rightarrow\pi_1(s)=\pi_1(s')\) for every \(s,s'\in S\). Define \(\phi\) on \(O_0^{\mathrm{eff}}\) by \(\phi(\pi_0(s))=\pi_1(s)\) for \(s\in S\). The fiber-constancy condition just stated ensures that \(\phi\) does not depend on the choice of \(s\) in the fiber, so \(\phi\) is well-defined. Thus there is a finite map \(\phi:O_0^{\mathrm{eff}}\to O_1\) with \(\pi_1=\phi\,\pi_0\). Therefore, if no witness pair exists, \(\pi_1\) factors through \(\pi_0\); equivalently, nonfactorization forces such a witness pair.

Both directions are finite checks on the finite set \(S\times S\).

\end{proof}

\subsubsection{Consequence}\label{sn:7.7.2}

The strict-extension condition is exactly nonfactorization of the new interface through the old interface; this is the strictness-gate comparison of Subsubsection \hyperref[sn:5.4.1]{5.4.1} at theorem grade. Strictness does not imply closure (NC-10) or drive (NC-11); the theorem records only the nonfactorization equivalence on the typed object-map data, and the no-go countermodels of Section \hyperref[sn:9]{9} supply finite explicit instances in which strictness coexists with absent macro closure or absent drive certificates.

\subsection{Theorem 13: Factorization Defect Equivalence}\label{sn:7.8}

\begin{claimtheorem}{Theorem 13 (FactorizationDefectEquivalence)}\phantomsection\label{cl:thm:13}

With \(S,\pi_0,\pi_1\) as in Theorem \hyperref[cl:thm:12]{12}, define the factorization defect

\[
\Delta_{\mathrm{fact}}(\pi_0,\pi_1)\;=\;\{\,(s,s')\in S\times S\;:\;\pi_0(s)=\pi_0(s')\;\wedge\;\pi_1(s)\neq\pi_1(s')\,\}.
\]

Then

\[
\Delta_{\mathrm{fact}}(\pi_0,\pi_1)\;\neq\;\varnothing\;\;\Longleftrightarrow\;\;\pi_1\;\not\factor\;\pi_0,
\]

and equivalently

\[
\Delta_{\mathrm{fact}}(\pi_0,\pi_1)\;=\;\varnothing\;\;\Longleftrightarrow\;\;\pi_1\;\factor\;\pi_0.
\]

With the same proof, this holds for arbitrary sets and maps; finiteness only makes it a finite check.

\end{claimtheorem}

\begin{proof}

By definition, \(\Delta_{\mathrm{fact}}(\pi_0,\pi_1)\) collects exactly the witness pairs of Theorem \hyperref[cl:thm:12]{12}. The equivalence then follows immediately: nonfactorization holds iff a witness pair exists iff \(\Delta_{\mathrm{fact}}\) is nonempty; factorization holds iff no witness pair exists iff \(\Delta_{\mathrm{fact}}\) is empty.

Both directions are finite checks on the finite set \(S\times S\).

The theorem makes the strictness condition a finite typed record: \(\Delta_{\mathrm{fact}}\) is the strictness defect of Subsubsection \hyperref[sn:5.4.1]{5.4.1} read in the present setting, and when the strictness gate \(G_{\mathrm{strict}}\) is checked, it is \(\texttt{strict\_pass}\) exactly when \(\Delta_{\mathrm{fact}}\neq\varnothing\) (Subsubsection \hyperref[sn:5.4.1]{5.4.1}).

\end{proof}

\subsection{Theorem 14: Fixed-Interface Definability Bound}\label{sn:7.9}

\begin{claimtheorem}{Theorem 14 (FixedInterfaceDefinabilityBound)}\phantomsection\label{cl:thm:14}

Let \(f:X\to Y\) be a finite map on \(\mathsf H_{\mathrm{fin}}\), and define the full finite-lens definability of \(f\) as

\[
\mathrm{Definable}(f)\;=\;\{\,f^{-1}(B)\;:\;B\subseteq\mathrm{im}(f)\,\}.
\]

Then

\[
|\,\mathrm{Definable}(f)\,|\;=\;2^{|\mathrm{im}(f)|}.
\]

For an arbitrary map \(f\), \(B\mapsto f^{-1}(B)\) is still a bijection from the subsets of \(\mathrm{im}(f)\) onto \(\mathrm{Definable}(f)\), so the identity holds with \(2^{|\mathrm{im}(f)|}\) read as a cardinal.

\end{claimtheorem}

\begin{proof}

Consider the map \(\Phi:\mathcal P(\mathrm{im}(f))\to\mathrm{Definable}(f)\) defined by \(\Phi(B)=f^{-1}(B)\). The map is well-defined and surjective onto \(\mathrm{Definable}(f)\) by construction.

\textbf{Injectivity.} Suppose \(B,B'\subseteq\mathrm{im}(f)\) with \(B\neq B'\). Without loss of generality, take \(y\in B\setminus B'\). Since \(y\in\mathrm{im}(f)\), there is some \(x\in X\) with \(f(x)=y\). Then \(x\in f^{-1}(B)\) but \(x\notin f^{-1}(B')\), so \(f^{-1}(B)\neq f^{-1}(B')\), that is, \(\Phi(B)\neq\Phi(B')\).

\textbf{Surjectivity.} By definition, every element of \(\mathrm{Definable}(f)\) is of the form \(f^{-1}(B)\) for some \(B\subseteq\mathrm{im}(f)\), hence in the image of \(\Phi\).

Therefore \(\Phi\) is a bijection \(\mathcal P(\mathrm{im}(f))\to\mathrm{Definable}(f)\), and

\[
|\,\mathrm{Definable}(f)\,|\;=\;|\mathcal P(\mathrm{im}(f))|\;=\;2^{|\mathrm{im}(f)|}.
\]

\end{proof}

\subsubsection{Consequence}\label{sn:7.9.1}

For a fixed finite interface \(f\), the number of definable package distinctions is exactly \(2^{|\mathrm{im}(f)|}\). Open-ended definability requires changing the interface, the package, the instrument, or the host; it does not arise from repeated application of a fixed completion (Theorem \hyperref[cl:thm:11]{11}) or from the same lens at a deeper level. Combined with the strict-extension theorems (Theorems \hyperref[cl:thm:12]{12} and \hyperref[cl:thm:13]{13}), the bound makes the discipline of Section \hyperref[sn:10]{10} explicit: strict promotion is recorded as an interface change rather than as theory growth at fixed interface.
 \section{Graph/Cohomology Drive Module}\label{sn:8}

This section proves the scoped no-drive and drive-separation results that the calculus uses to keep the drive face \(\mathrm{P6}_{\mathrm{drive}}\) distinct from generic \(P_6\) audit. The graph cohomology used is the elementary cycle/cochain theory of a finite graph \citep{Hatcher2002AT}. A drive is a circulation that no potential explains: a \(1\)-cochain whose integral around some cycle is nonzero, that is, a nonzero class in \(H^1\). The motivating case is a Markov chain on the graph whose transitions have positive probability in both directions: the cochain \(a(x\to y)=\log(P_{xy}/P_{yx})\) integrates around a cycle to that cycle's affinity, and a nonzero cycle affinity rules out detailed balance on that bidirectional Markov support. The theorems below say when such a drive can and cannot exist. All theorems in this section live on the host \(\mathsf H_{\mathrm{graph}}\) of finite graphs and finite cycle/cochain data, and all theorems are scoped to cycle-supported cohomological drive on a declared finite graph. The module is not a universal theory of drive: it does not classify all thermodynamic force, all entropy production, or all directionality. The negative scope of Subsubsection \hyperref[sn:2.4.3]{2.4.3} records the corresponding nonclaims, NC-12 and NC-32 of Subsection \hyperref[sn:15.2]{15.2} in particular.

\subsection{Finite Graph/Cohomology Support Model}\label{sn:8.1}

This subsection fixes the finite graph/cohomology support model used by Theorems \hyperref[cl:thm:15]{15} through \hyperref[cl:thm:18]{18}. It is a schema-grade declaration of the host data, the drive certificate format, and the gating operations admitted on the host. The no-drive and gating results are stated and proved in Subsections \hyperref[sn:8.2]{8.2} through \hyperref[sn:8.5]{8.5}.

\subsubsection{Graph Data}\label{sn:8.1.1}

The host \(\mathsf H_{\mathrm{graph}}\) provides finite graph data:

\begin{itemize}
\tightlist
\item
  a finite vertex set \(V(G)\);
\item
  a finite set \(E(G)\) of oriented edges (an unoriented graph is given a fixed orientation, with one oriented edge for each unordered adjacent pair and \(a(\bar e)=-a(e)\) for the reversed edge \(\bar e\); self-loops \(x\to x\) are omitted (the log-ratio cochain vanishes on them); Theorems \hyperref[cl:thm:15]{15}--\hyperref[cl:thm:18]{18} are stated for graphs without self-loops);
\item
  the finite first Betti number \(\beta_1(G)\;=\;|E(G)|\;-\;|V(G)|\;+\;c(G)\), where \(c(G)\) is the number of connected components of \(G\);
\item
  the finite-dimensional (finitely generated free) cycle complex \(Z_*(G)\) and cochain complex \(C^*(G)\) with coboundary \(d:C^0\to C^1\), \((d\phi)(e)=\phi(\mathrm{head}(e))-\phi(\mathrm{tail}(e))\), and boundary \(\partial:k^{E(G)}\to k^{V(G)}\), \(\partial e=\mathrm{head}(e)-\mathrm{tail}(e)\), with \(Z_1(G)=\ker\partial\), taken over a fixed coefficient field or, more generally, a declared commutative ring \(k\) with identity such as \(\mathbb Z\) (a \emph{free coefficient host}: \(C^0\), \(C^1\) and \(Z_1\) are then free \(k\)-modules, with bases the vertices, the edges and the fundamental cycles of a spanning forest), with dimensions read as free ranks;
\item
  the first cohomology \(H^1(G)\;=\;C^1(G)/\mathrm{im}(d)\), a free module of rank \(\beta_1(G)\) over the declared coefficients (Theorem \hyperref[cl:thm:17]{17}).
\end{itemize}

The cycle integral of a finite \(1\)-cochain \(a\in C^1(G)\) around a cycle \(\gamma=\sum_e\gamma(e)\,e\in Z_1(G)\) is

\[
\oint_\gamma a\;:=\;\sum_{e\in E(G)}\gamma(e)\,a(e).
\]

For a closed walk \(e_1\cdots e_n\) with signs \(\varepsilon_k=\pm1\) (\(+1\) when \(e_k\) is traversed from tail to head), this equals \(\sum_k\varepsilon_k\,a(e_k)\). The map \(\gamma\mapsto\oint_\gamma a\) is linear. Every record in this section carries a host tag \(\mathsf H_{\mathrm{graph}}\), a level tag drawn from \(\mathsf{Lev}\), and an instrument annotation. Cycle-supported drive is the only drive notion adjudicated here.

\subsubsection{Drive Certificate}\label{sn:8.1.2}

The drive face \(\mathrm{P6}_{\mathrm{drive}}^{H^1}\) on \(\mathsf H_{\mathrm{graph}}\) is a finite typed record stating that some cohomology class \([a]\in H^1(G)\) is nonzero. A positive cycle-supported cohomological drive certificate consists of

\begin{itemize}
\tightlist
\item
  a finite \(1\)-cochain \(a\in C^1(G)\);
\item
  a finite cycle \(\gamma\in Z_1(G)\) with \(\oint_\gamma a\neq 0\);
\item
  the audit and source records required by \(\mathsf{AuditPolicy}_I\) and \(\mathsf{SourcePolicy}_I\) for the surrounding instrument.
\end{itemize}

A no-drive certificate is a \(P_6\) threshold-check witness record of Subsubsection \hyperref[sn:3.7.4]{3.7.4} containing the graph \(G\), the cochain \(a\), a declared finite basis of \(Z_1(G)\), and the checked equalities \(\oint_\gamma a=0\) for every basis cycle, together with the same audit and source data. By linearity these equalities hold on every cycle.

A generic \(P_6\) audit record is not a drive certificate: it does not record cycle integrals or cohomology data. The separation between the audit face \(P_6\) and the drive face \(\mathrm{P6}_{\mathrm{drive}}^{H^1}\) is preserved at the certificate level throughout this section.

\subsubsection{Gating}\label{sn:8.1.3}

A \(P_2\)-style gate on \(\mathsf H_{\mathrm{graph}}\) is a finite admissible map \(G\mapsto G'\) obtained by deleting a finite set of edges or by restricting to a finite induced subgraph. A forest-producing gate is a gate whose output \(G'\) has \(\beta_1(G')=0\). Gates are recorded as finite typed records on the surrounding judgment; an admissible gate is one that conforms to the scope and bridge policy of the surrounding instrument.

\subsection{Theorem 15: Forest No-Drive}\label{sn:8.2}

\begin{claimtheorem}{Theorem 15 (ForestNoDrive)}\phantomsection\label{cl:thm:15}

For a finite graph \(G\) on \(\mathsf H_{\mathrm{graph}}\), if \(\beta_1(G)=0\), then

\[
Z_1(G)\;=\;0\quad\text{and}\quad H^1(G)\;=\;0.
\]

In particular, no cycle-supported cohomological drive certificate \(\mathrm{P6}_{\mathrm{drive}}^{H^1}\) is admissible on \(G\).

\end{claimtheorem}

\begin{proof}

If \(\beta_1(G)=0\) then \(|E(G)|=|V(G)|-c(G)\), the edge count of any spanning forest, so \(G\) is a forest. Then \(Z_1(G)=0\) over any coefficient ring: for \(\gamma\neq0\) the edges with \(\gamma(e)\neq0\) form a nonempty forest, which has a vertex \(v\) on exactly one such edge \(e\), and \((\partial\gamma)(v)=\pm\gamma(e)\neq0\). The first cohomology

\[
H^1(G)\;=\;C^1(G)/\mathrm{im}(d)
\]

vanishes over any coefficient ring. Indeed, \(G\) is a forest, since \(\beta_1(G)=0\). Fix a root in each component and, for \(a\in C^1(G)\), define \(\phi(v)\) as the sum of \(\pm a(e)\) over the edges \(e\) of the unique path from the root of \(v\)'s component to \(v\), with sign \(+\) when \(e\) is traversed from tail to head and \(-\) otherwise; the path is unique, so \(\phi\) is well defined. For each edge \(e\), the path to \(\mathrm{head}(e)\) is the path to \(\mathrm{tail}(e)\) extended by \(e\), or the path to \(\mathrm{tail}(e)\) is the path to \(\mathrm{head}(e)\) extended by \(e\) reversed, so \(d\phi(e)=\phi(\mathrm{head}(e))-\phi(\mathrm{tail}(e))=a(e)\). Hence \(a=d\phi\) for every \(a\in C^1(G)\), and \(H^1(G)=0\).

With \(H^1(G)=0\) and \(Z_1(G)=0\), the only cycle is \(\gamma=0\), and \(\oint_0a=0\) by linearity, so no nonzero cycle affinity is recorded. Hence no cycle-supported drive certificate is admissible: a positive certificate would require some \(\gamma\in Z_1(G)\) with \(\oint_\gamma a\neq 0\), but no such cycle exists.

\end{proof}

\subsection{Theorem 16: Exact-Form Null-Drive}\label{sn:8.3}

\begin{claimtheorem}{Theorem 16 (ExactFormNullDrive)}\phantomsection\label{cl:thm:16}

Let \(G\) be a finite graph on \(\mathsf H_{\mathrm{graph}}\). If a finite \(1\)-cochain \(a\in C^1(G)\) is exact, that is,

\[
a\;=\;d\phi\quad\text{for some }0\text{-cochain }\phi\in C^0(G),
\]

then for every cycle \(\gamma\in Z_1(G)\),

\[
\oint_\gamma a\;=\;0.
\]

\end{claimtheorem}

\begin{proof}

Let \(\gamma\) be a closed walk \(v_0,\ldots,v_n=v_0\) whose \(k\)-th step traverses the edge \(e_k\) with sign \(\varepsilon_k=+1\) along its orientation and \(\varepsilon_k=-1\) against it. By definition of the coboundary \(d\phi(e)=\phi(\mathrm{head}(e))-\phi(\mathrm{tail}(e))\), in either direction \(\varepsilon_k\,d\phi(e_k)=\phi(v_k)-\phi(v_{k-1})\), so the cycle integral telescopes:

\[
\oint_\gamma a\;=\;\sum_{k=1}^{n}\,\varepsilon_k\,a(e_k)\;=\;\sum_{k=1}^{n}\bigl(\phi(v_k)\;-\;\phi(v_{k-1})\bigr)\;=\;\phi(v_n)\;-\;\phi(v_0)\;=\;0,
\]

since \(v_n=v_0\) on the walk. For a general \(\gamma\in Z_1(G)\), \(\oint_\gamma d\phi=\sum_e\gamma(e)\bigl(\phi(\mathrm{head}(e))-\phi(\mathrm{tail}(e))\bigr)=\sum_v\phi(v)(\partial\gamma)(v)=0\), since \(\partial\gamma=0\) (as in Appendix \hyperref[sn:B.2.2]{B.2.2}), so \(\oint_\gamma a=0\) for every \(\gamma\in Z_1(G)\).

The result records that exact \(1\)-cochains carry no cycle drive: any cohomology class supported on an exact cochain is zero, and the corresponding drive certificate is empty on every cycle.

\end{proof}

\subsection{Theorem 17: Nonzero-Affinity Equivalence}\label{sn:8.4}

\begin{claimtheorem}{Theorem 17 (NonzeroAffinityEquivalence)}\phantomsection\label{cl:thm:17}

Let \(G\) be a finite graph on \(\mathsf H_{\mathrm{graph}}\) with the fixed coefficient field or declared free coefficient host of Subsubsection \hyperref[sn:8.1.1]{8.1.1}, and let \(a\in C^1(G)\) with cohomology class \([a]\in H^1(G)\). Then

\[
[a]\;\neq\;0\;\;\Longleftrightarrow\;\;\exists\,\gamma\in Z_1(G)\;:\;\oint_\gamma a\;\neq\;0.
\]

Moreover, for any spanning forest \(T\subseteq E(G)\), \([a]\neq0\) if and only if \(\oint_{\gamma_e}a\neq0\) for some \(e\notin T\), where \(\gamma_e\) is the fundamental cycle of \(e\); so testing the \(\beta_1(G)\) fundamental cycles suffices; indeed \(a\mapsto(\oint_{\gamma_e}a)_{e\notin T}\) induces an isomorphism \(H^1(G)\cong k^{\beta_1(G)}\), where \(k\) is the coefficient field or free coefficient host. In particular, when \(1\neq0\) in \(k\), \(H^1(G)=0\) if and only if \(\beta_1(G)=0\); for \(e\notin T\) the cochain \(\mathbf 1_e\) has \(\oint_{\gamma_e}\mathbf 1_e=1\neq0\), so a graph, or the output of a gate, admits a cochain and a cycle with nonzero integral, the mathematical part of a positive drive certificate of Subsubsection \hyperref[sn:8.1.2]{8.1.2}, exactly when its first Betti number is positive.

\end{claimtheorem}

\begin{proof}

For the reverse direction, suppose some \(\gamma\in Z_1(G)\) satisfies \(\oint_\gamma a\neq 0\). By the contrapositive of Theorem \hyperref[cl:thm:16]{16}, \(a\) is not exact: if \(a=d\phi\) for some \(\phi\), then \(\oint_\gamma a=0\), contradicting \(\oint_\gamma a\neq 0\). Hence \([a]\neq 0\in H^1(G)\).

For the forward direction, we prove the contrapositive: if \(\oint_\gamma a=0\) for every \(\gamma\in Z_1(G)\), then \(a\) is exact. Choose a spanning forest \(T\subseteq E(G)\) and a root vertex in each connected component. For a vertex \(v\), let \(\phi(v)\) be the sum of \(a\) along the unique path in \(T\) from the root of its component to \(v\), each edge counted with sign \(+1\) when traversed from tail to head and \(-1\) otherwise; so \(\phi(\text{root})=0\). For an edge \(e\in T\) from \(u\) to \(v\), the \(T\)-path to \(v\) is the \(T\)-path to \(u\) followed by \(e\), or the \(T\)-path to \(u\) passes through \(v\) and ends with \(e\) traversed backwards; in both cases \(\phi(v)-\phi(u)=a(e)\), that is, \(a(e)=d\phi(e)\). For an edge \(e\notin T\) from \(u\) to \(v\), let \(\gamma_e\in Z_1(G)\) be its fundamental cycle: \(e\) followed by the \(T\)-path from \(v\) back to \(u\), with the same sign convention. Along \(T\), \(a=d\phi\), so the path part contributes \(\phi(u)-\phi(v)\), and

\[
0\;=\;\oint_{\gamma_e}a\;=\;a(e)\;+\;\phi(u)\;-\;\phi(v),
\]

hence again \(a(e)=d\phi(e)\). So \(a=d\phi\) and \([a]=0\). The argument uses only addition and subtraction of coefficients, so it holds over the coefficient field and over any declared free coefficient host. Therefore \([a]\neq 0\) implies that some cycle \(\gamma\) has \(\oint_\gamma a\neq 0\). The contrapositive used only the cycles \(\gamma_e\) with \(e\notin T\), which gives the Moreover clause; the reverse direction of that clause is the reverse direction above. For the isomorphism, the map \(a\mapsto(\oint_{\gamma_e}a)_{e\notin T}\) is linear, and its kernel is \(\mathrm{im}(d)\) by the argument above and Theorem \hyperref[cl:thm:16]{16}. It is surjective: given \((c_e)_{e\notin T}\), set \(a(e)=c_e\) off \(T\) and \(a=0\) on \(T\); \(\gamma_e\) contains exactly one edge outside \(T\), namely \(e\) with coefficient \(1\), so \(\oint_{\gamma_e}a=c_e\). Since \(|E(G)\setminus T|=|E(G)|-|V(G)|+c(G)=\beta_1(G)\), \(H^1(G)\cong k^{\beta_1(G)}\).

The theorem records the equivalence between the cohomological drive certificate \([a]\neq 0\) and the cycle-affinity drive certificate \(\oint_\gamma a\neq 0\) on \(\mathsf H_{\mathrm{graph}}\) with the declared coefficients; both carry the same cycle-supported cohomological drive face \(\mathrm{P6}_{\mathrm{drive}}^{H^1}\).

\end{proof}

\subsection{Theorem 18: Gating-Affinity Suppression}\label{sn:8.5}

\begin{claimtheorem}{Theorem 18 (GatingAffinitySuppression)}\phantomsection\label{cl:thm:18}

Let \(G\) be a finite graph on \(\mathsf H_{\mathrm{graph}}\), and let \(G\mapsto G'\) be a \(P_2\)-style gate that produces a forest, that is,

\[
\beta_1(G')\;=\;0.
\]

Then

\[
H^1(G')\;=\;0,
\]

and no positive cycle-supported cohomological drive certificate \(\mathrm{P6}_{\mathrm{drive}}^{H^1}\) exists on \(G'\). For every \(a\in C^1(G')\), every cycle integral vanishes. Since \(\beta_1(G')=0\), the declared finite basis of \(Z_1(G')\) is empty, and the vanishing record on it is a no-drive certificate whenever it is accompanied by the source and audit records required by the surrounding instrument.

\end{claimtheorem}

\begin{proof}

By Theorem \hyperref[cl:thm:15]{15} applied to \(G'\), the cohomology \(H^1(G')\) vanishes when \(\beta_1(G')=0\). Since the drive certificate \(\mathrm{P6}_{\mathrm{drive}}^{H^1}\) requires a cycle \(\gamma\in Z_1(G')\) with \(\oint_\gamma a\neq 0\), and \(Z_1(G')=0\) by Theorem \hyperref[cl:thm:15]{15}, no such certificate exists on \(G'\). For the same reason \(\oint_\gamma a=0\) holds for every cycle of \(G'\) and every \(a\in C^1(G')\). Thus the vanishing-integral condition holds for every \(a\). Attaching the source and audit records required by the surrounding instrument gives the no-drive certificate of Subsubsection \hyperref[sn:8.1.2]{8.1.2}.

\end{proof}

\subsubsection{Nonclaim}\label{sn:8.5.1}

The theorem concerns cycle-supported cohomological drive on a declared finite graph and gate. It does not classify all drive, all thermodynamic force, all entropy production, or all directionality. A drive notion that does not reduce to a cycle-cohomology certificate on a finite graph is outside the scope of this theorem; its adjudication requires a separate drive bridge.

The dependency

\[
P_3\;\mathrel{\mathsf{dep}_{\texttt{requires\_bridge}}}\;\mathrm{P6}_{\mathrm{drive}}
\]

of Subsubsection \hyperref[sn:4.6.2]{4.6.2} remains active: route mismatch alone does not establish a drive certificate, even when route mismatch is recorded as a \(P_3\) defect. NC-12 of Subsection \hyperref[sn:15.2]{15.2} records the corresponding nonclaim.
 \section{Countermodel Atlas}\label{sn:9}

This section presents six of the twelve countermodels. Each countermodel is a finite explicit construction together with the no-go statement it supports, the overclaim it blocks, and any positive bridge that would have to be supplied for a positive claim. The full atlas of twelve countermodels is recorded in Appendix \hyperref[sn:C]{C}; Sections \hyperref[sn:6]{6} through \hyperref[sn:11]{11} also refer to appendix countermodels, for example \(\mathsf{CM}_7\) for Theorem \hyperref[cl:thm:1]{1}, and \(\mathsf{CM}_8\) records the instance of Theorem \hyperref[cl:thm:9]{9}.

\subsection{Atlas Convention}\label{sn:9.1}

This subsection fixes the sheet form for countermodels and selects which countermodels appear in the main text. It is a schema-grade declaration, with the actual constructions in Subsections \hyperref[sn:9.2]{9.2} through \hyperref[sn:9.7]{9.7}.

\subsubsection{Countermodel Sheet Form}\label{sn:9.1.1}

Unless a sheet lists them, its packages and instruments are completed as in Subsubsection \hyperref[sn:10.6.1]{10.6.1}: host \(\mathsf H_{\mathrm{fin}}\) (\(\mathsf H_{\mathrm{graph}}\) for graph data), level \(\mathsf{Ver}\), every field visible, a \(\texttt{committed\_state}\) source holding the tables, one audit entry recomputing them, an empty defect family and a nonclaim record; the conditions of Subsubsection \hyperref[sn:3.3.3]{3.3.3} then hold by inspection. For every displayed dependency judgment, its judging instrument admits the claim type \(\texttt{dependency}\) and the displayed host; its scope contains \(\Gamma,\mathcal C,A,V,\Theta\) and their referenced records, with their level tags admitted. The sheet supplies the finite audit, source and check-rule records it uses. When a sheet borrows an instrument from another construction, it uses a fresh local copy with these additions, extending its mode and required-strength policies to \(\texttt{dependency}\) and preserving its existing policies and verdicts; in particular, the local \(I_{\mathrm{prom}}\) of \(\mathsf{CM}_9\) admits both promotion and dependency judgments.

Each countermodel sheet contains the following five fields, in this order:

\begin{itemize}
\tightlist
\item
  \textbf{Construction.} A finite explicit construction on a named host, with the data of the calculus (carriers, maps, packages, profiles, witnesses, updates, defects, audit and source records) made explicit. The construction is finite at every step.
\item
  \textbf{Supported no-go.} The dependency judgment \(S\mathrel{\mathsf{dep}_r}X\) that the construction supports, in the form fixed by Subsubsection \hyperref[sn:4.6.1]{4.6.1}, with relation type \(r\) drawn from the eight admissible types \(\texttt{suffices}\), \(\texttt{insufficient}\), \(\texttt{destroys}\), \(\texttt{blocks}\), \(\texttt{enables}\), \(\texttt{equivalent}\), \(\texttt{requires\_bridge}\), and \(\texttt{outside\_scope}\). \emph{Reading the predicates.} In a sheet, except for \(\texttt{destroys}\), \(S\) and \(X\) are read as predicates on instances containing the displayed records; for \(\texttt{destroys}\), \(S\) names the operation and \(X\) is the predicate tested before and after it. For example, \(P_5(q)\) holds when the instance contains an admissible package with lens \(q\), \(P_1(qF=F^\sharp q)\) when some map \(F^\sharp:\mathrm{im}(q)\to\mathrm{im}(q)\) satisfies \(F^\sharp q_{\mathrm{im}}=q_{\mathrm{im}}F\), with \(q_{\mathrm{im}}\) the corestriction of \(q\) to its image, recorded or not; \(P_5(q_{\mathrm{global}})\) in \(\mathsf{CM}_3\) when some \(q:Z\to O\) with \(q|_U=q_U\) and \(q|_V=q_V\) exists; and \(\mathrm{P6}^{H^1}_{\mathrm{drive}}\) holds when the instance records a graph \(G\) and a cochain \(a\in C^1(G)\) with \([a]\neq0\) in \(H^1(G)\); a cochain recorded by a drive bridge \(B^{\mathrm{drive}}_{3\to6}\) counts on the graph that bridge declares (Subsubsection \hyperref[sn:9.3.5]{9.3.5}), where its class is nonzero by Theorem \hyperref[cl:thm:17]{17}. In \(\mathsf{CM}_2\) and \(\mathsf{CM}_{10}\) a cochain counts when it is recorded on the triangle \(G\), respectively on the support graph \(G_B\), or by a drive bridge on the graph that bridge declares; in \(\mathsf{CM}_{12}\) only cochains recorded on the graph the gate acts on count, the triangle \(G\) before the gate and its output \(G'\) after it. An instance that records no cochain fails the predicate. Likewise \(P_3^{\mathrm{holonomy}}\) holds when the instance records a \(P_3\) transport record and a cycle \(\gamma\) with \(\delta_3^{\mathrm{hol}}(\gamma)\neq\varnothing\); \(P_4(q_0\rightsquigarrow q_1)\) when it contains a package refinement (Subsubsection \hyperref[sn:3.3.2]{3.3.2}) from the package with lens \(q_0\) to the package with lens \(q_1\); ``\(I_0\) audits \(\mathcal T_0\)'' when it contains an audit record whose \(\mathsf{ClaimRef}\) is \(\mathcal T_0\) and whose check rules are drawn from \(\mathsf{CheckRules}_{I_0}\); and \(\mathrm{src}:\texttt{fallback}\) when the source slot of \(Q\) holds a record of type \(\texttt{fallback}\).
\item
  \textbf{What it blocks.} The overclaim that the no-go rules out, named together with the corresponding nonclaim NC-\(n\) from Subsection \hyperref[sn:15.2]{15.2}.
\item
  \textbf{What it does not block.} The positive content that the construction is consistent with, so that the no-go is not over-read into a stronger negative statement.
\item
  \textbf{Required positive bridge.} When a positive claim of the form \(S\Rightarrow X\) would need an additional bridge, the bridge type is named (for example, a drive bridge \(B_{3\to 6}^{\mathrm{drive}}\) or a gluing bridge \(B_{\mathrm{glue}}\)).
\end{itemize}

Formally, a countermodel sheet is a finite tuple

\[
(\mathrm{id},H,\mathsf{Data},\mathsf{ActiveClaim},\mathsf{FailedClaim},W,A,V,\mathcal N),
\]

with the five prose fields above unpacking that tuple for the reader.

\subsubsection{Main-Text Selection}\label{sn:9.1.2}

The six countermodels of this section are:

\begin{itemize}
\tightlist
\item
  \(\mathsf{CM}_1\) --- Active package but no macro closure (Subsection \hyperref[sn:9.2]{9.2});
\item
  \(\mathsf{CM}_2\) --- Active route/holonomy but no drive (Subsection \hyperref[sn:9.3]{9.3});
\item
  \(\mathsf{CM}_3\) --- Local packages fail globally (Subsection \hyperref[sn:9.4]{9.4});
\item
  \(\mathsf{CM}_4\) --- Directed cell active but pair observable blocked (Subsection \hyperref[sn:9.5]{9.5});
\item
  \(\mathsf{CM}_5\) --- Refinement worsens closure (Subsection \hyperref[sn:9.6]{9.6});
\item
  \(\mathsf{CM}_6\) --- Same-level self-audit circularity (Subsection \hyperref[sn:9.7]{9.7}).
\end{itemize}

The full twelve-countermodel atlas, including same-primitive-pair-different-statuses (\(\mathsf{CM}_7\)), idempotent-completions-do-not-commute (\(\mathsf{CM}_8\)), strict-extension-without-closure (\(\mathsf{CM}_9\)), strict-extension-without-drive (\(\mathsf{CM}_{10}\)), fallback-source (\(\mathsf{CM}_{11}\)), and gating-destroys-cycle-supported-drive (\(\mathsf{CM}_{12}\)), is recorded with full constructions in Appendix \hyperref[sn:C]{C}. Subsection \hyperref[sn:9.8]{9.8} names these appendix-only countermodels and points to Appendix \hyperref[sn:C]{C} for the constructions.

\subsection{Countermodel 1: Active Package but No Macro Closure}\label{sn:9.2}

\subsubsection{Construction}\label{sn:9.2.1}

Take \(X=\{a,b,c\}\) and \(Y=\{A,B\}\), both finite sets on \(\mathsf H_{\mathrm{fin}}\). Define a finite package map

\[
q:X\to Y,\qquad q(a)=A,\quad q(b)=A,\quad q(c)=B.
\]

The map \(q\) is admissible as a \(P_5\)-package on \(X\) in the schema of Subsection \hyperref[sn:3.3]{3.3}: it is a finite map on a finite carrier with a finite codomain.

Define a finite update

\[
F:X\to X,\qquad F(a)=a,\quad F(b)=c,\quad F(c)=c.
\]

Then \(q(a)=q(b)=A\), while \(qF(a)=q(a)=A\) and \(qF(b)=q(c)=B\), so

\[
qF(a)\;\neq\;qF(b).
\]

By Theorem \hyperref[cl:thm:7]{7}, \((a,b)\in\delta_1^{\mathrm{split}}(q,F)\), so

\[
\delta_1^{\mathrm{split}}(q,F)\;\neq\;\varnothing,
\]

and no finite descended update \(F^\sharp:Y\to Y\) with \(qF=F^\sharp q\) exists.

\subsubsection{Supported No-Go}\label{sn:9.2.2}

The construction supports the dependency judgment

\[
\Gamma\,;\;\mathsf H_{\mathrm{fin}}\,;\;I_{\mathrm{fin}}\;\vdash\;P_5(q)\;\mathrel{\mathsf{dep}_{\texttt{insufficient}}}\;P_1(qF=F^\sharp q),
\]

that is, \(P_5\not\Rightarrow P_1\text{ macro closure}\) on the host \(\mathsf H_{\mathrm{fin}}\) under the declared admissible instrument.

\subsubsection{What It Blocks}\label{sn:9.2.3}

The overclaim ``a package automatically closes the dynamics'' is rejected. NC-7 of Subsection \hyperref[sn:15.2]{15.2} records the corresponding nonclaim. The construction is the working content of NC-7: package admissibility on a finite carrier does not by itself produce a descended update.

\subsubsection{What It Does Not Block}\label{sn:9.2.4}

The construction does not say packages are useless or that descent is generically impossible. It says only that package existence alone does not guarantee descent of a given update, and a positive descent claim requires an explicit witness: in this construction, replacing \(F\) with a different finite update consistent with the fibers of \(q\) would produce an admissible descended map.

\subsubsection{Required Positive Bridge}\label{sn:9.2.5}

A positive descent claim \(P_5\Rightarrow P_1\) on a specific package \(q\) and update \(F\) requires an explicit descent record: the empty defect \(\delta_1^{\mathrm{split}}(q,F)=\varnothing\) together with the audit and visibility data required by the surrounding judgment. Theorem \hyperref[cl:thm:7]{7} makes the vanishing of \(\delta_1^{\mathrm{split}}(q,F)\) necessary and sufficient for the descended update, while the audit and visibility records remain separate admissibility requirements.

\subsection{Countermodel 2: Active Route / Holonomy but No Drive}\label{sn:9.3}

\subsubsection{Construction}\label{sn:9.3.1}

Take a finite cycle graph \(G\) on \(\mathsf H_{\mathrm{graph}}\), for example the triangle \(0\to 1\to 2\to 0\) with \(\beta_1(G)=1\), so that \(Z_1(G)\) is spanned by the triangle cycle \(\gamma\). Equip the surrounding judgment with the \(P_3\) transport record of Subsubsection \hyperref[sn:5.3.3]{5.3.3} with fiber \(\Phi=\{u,v\}\), \(g_{0\to1}=g_{1\to2}=\mathrm{id}\) and \(g_{2\to0}=\sigma\), the swap of \(u\) and \(v\). Then

\[
\operatorname{Hol}(\gamma)\;=\;\sigma\;\neq\;\mathrm{id},\qquad \delta_3^{\mathrm{hol}}(\gamma)=\{u,v\},
\]

so that a role-channel record for \(P_3\) with evidence \(\delta_3^{\mathrm{hol}}(\gamma)\neq\varnothing\) can receive \(\texttt{active\_projection}\).

For the drive face, take the symmetric random walk on \(G\), \(M(x,y)=1/2\) for each edge \(\{x,y\}\), and its log-ratio cochain \(a(x\to y)=\log\bigl(M(x,y)/M(y,x)\bigr)\); then \(a=0\in C^1(G)\). By Theorem \hyperref[cl:thm:16]{16} (or directly from \(a=0\)), every cycle integral vanishes:

\[
\oint_\gamma a\;=\;0\quad\text{for every cycle }\gamma\in Z_1(G).
\]

Hence \([a]=0\in H^1(G)\), and no positive cycle-supported cohomological drive certificate \(\mathrm{P6}_{\mathrm{drive}}^{H^1}\) is admissible on \(G\) for this \(a\). The route/holonomy data is active, but the drive certificate is empty: the affinity cochain of a reversible chain on the same graph has no drive, while the holonomy is \(\sigma\).

\subsubsection{Supported No-Go}\label{sn:9.3.2}

The construction supports the dependency judgment

\[
\Gamma\,;\;\mathsf H_{\mathrm{graph}}\,;\;I_{\mathrm{graph}}\;\vdash\;P_3^{\mathrm{holonomy}}\;\mathrel{\mathsf{dep}_{\texttt{requires\_bridge}}}\;\mathrm{P6}_{\mathrm{drive}}^{H^1},
\]

recording that route mismatch or holonomy alone does not establish a cohomological drive certificate. The construction establishes the first conjunct, \(P_3^{\mathrm{holonomy}}\mathrel{\mathsf{dep}_{\texttt{insufficient}}}\mathrm{P6}^{H^1}_{\mathrm{drive}}\); the second conjunct of the truth condition of \(\texttt{requires\_bridge}\) (Subsubsection \hyperref[sn:4.6.1]{4.6.1}) holds trivially: by Subsubsection \hyperref[sn:9.3.5]{9.3.5}, \(B^{\mathrm{drive}}_{3\to6}\) itself records a cochain \(a\) and a cycle \(\gamma\) with \(\oint_\gamma a\neq0\), a complete positive certificate (Subsubsection \hyperref[sn:8.1.2]{8.1.2}) whatever the route data. The content of the judgment is its insufficiency half.

\subsubsection{What It Blocks}\label{sn:9.3.3}

The overclaim ``route mismatch or holonomy automatically certifies drive'' is rejected. NC-12 of Subsection \hyperref[sn:15.2]{15.2} records the corresponding nonclaim. The construction is consistent with active \(P_3\) data and absent \(\mathrm{P6}_{\mathrm{drive}}^{H^1}\) data on the same graph.

\subsubsection{What It Does Not Block}\label{sn:9.3.4}

The construction does not say drive certificates are unattainable on \(G\) or that route mismatch is incompatible with drive. A different choice of \(1\)-cochain \(a\) on \(G\) --- one with \(\oint_\gamma a\neq 0\) --- would supply a positive drive certificate; the construction shows only that the choice of \(a\) is independent data from the route/holonomy record.

\subsubsection{Required Positive Bridge}\label{sn:9.3.5}

A positive inference of the form \(P_3\Rightarrow\mathrm{P6}_{\mathrm{drive}}^{H^1}\) requires an explicit drive bridge

\[
B_{3\to 6}^{\mathrm{drive}},
\]

a finite typed record in \(\mathsf{BridgePolicy}_I\) that connects the route/holonomy data to a specific cochain \(a\) and cycle \(\gamma\) with \(\oint_\gamma a\neq 0\), together with the audit and source records required by \(\mathsf{AuditPolicy}_I\) and \(\mathsf{SourcePolicy}_I\) (Subsubsection \hyperref[sn:8.1.2]{8.1.2}). The bridge supplies the missing data: the route witness alone does not produce \(a\), and the choice of \(a\) is the substantive content of the bridge. Such a drive bridge is \emph{admissible} under \(I\) when \(\mathsf{BridgePolicy}_I\) admits the type \(\texttt{drive}\), the bridge records a declared finite graph \(G\), a cochain \(a\in C^1(G)\) and a cycle \(\gamma\in Z_1(G)\) with \(\oint_\gamma a\neq0\), and its audit draws its \(\mathsf{CheckRules}\) from \(\mathsf{CheckRules}_I\) with audit defect \(\texttt{audit\_passes}\); this is the sense used in Subsubsections \hyperref[sn:4.6.1]{4.6.1} and \hyperref[sn:5.4.2]{5.4.2}.

\subsection{Countermodel 3: Local Packages Fail Globally}\label{sn:9.4}

\subsubsection{Construction}\label{sn:9.4.1}

Take a finite carrier \(Z=\{a,b,c\}\) on \(\mathsf H_{\mathrm{fin}}\) with a finite cover \(\{U,V\}\) given by

\[
U\;=\;\{a,b\},\qquad V\;=\;\{b,c\},
\]

so that the overlap is \(U\cap V=\{b\}\). Let local packages be finite maps

\[
q_U:U\to O_U,\qquad q_V:V\to O_V,
\]

each admissible on its own patch. The restriction system is part of the data. All package targets are the two-element set,

\[
O_U\;=\;O_V\;=\;O_{U\cap V}\;=\;O\;=\;\{0,1\},
\]

every restriction map on targets is the identity of \(\{0,1\}\), and a global package is a map \(q:Z\to O\) whose restrictions to the patches are the local packages, \(q|_U=q_U\) and \(q|_V=q_V\). Take \(q_U(a)=q_U(b)=0\) and \(q_V(b)=q_V(c)=1\), so that on the overlap

\[
\mathrm{res}_{U,U\cap V}(q_U(b))\;=\;0,\qquad \mathrm{res}_{V,U\cap V}(q_V(b))\;=\;1,
\]

and \(0\neq1\) in the overlap target. Each local package is admissible on its own patch, but no global package \(q:Z\to O\) extends both: \(q(b)\) would have to equal \(q_U(b)=0\) and \(q_V(b)=1\).

The gluing defect \(\delta_4^{\mathrm{glue}}\) on the cover \(\{U,V\}\) is therefore nonempty, with \(b\) as a witness.

\subsubsection{Supported No-Go}\label{sn:9.4.2}

The construction supports the dependency judgment

\[
\Gamma\,;\;\mathsf H_{\mathrm{fin}}\,;\;I_{\mathrm{fin}}\;\vdash\;\bigl\{P_5(q_U),\,P_5(q_V)\bigr\}\;\mathrel{\mathsf{dep}_{\texttt{insufficient}}}\;P_5(q_{\mathrm{global}}),
\]

recording that local-\(P_5\)-validity does not by itself imply global package existence on the cover.

\subsubsection{What It Blocks}\label{sn:9.4.3}

The overclaim ``local validity implies global validity'' is rejected. NC-14 of Subsection \hyperref[sn:15.2]{15.2} records the corresponding nonclaim. The construction is consistent with two admissible local packages and an absent global package on the same finite carrier.

\subsubsection{What It Does Not Block}\label{sn:9.4.4}

The construction does not say global packaging is generically impossible or that covers are useless. Different overlap data --- with admissibly identifiable values on \(U\cap V\) --- would produce a global package; the construction shows only that the overlap data is part of the global packaging input. The no-go is relative to the restriction system fixed above. A different global target with different restriction maps, for example one that records the pair of local values at \(b\), poses a different gluing problem, which this construction does not address.

\subsubsection{Required Positive Repair}\label{sn:9.4.5}

A positive global-packaging claim from local data requires a \(P_4\)-gluing bridge

\[
B_{\mathrm{glue}}\quad\text{with}\quad \delta_4^{\mathrm{glue}}\;=\;\varnothing,
\]

together with a \(P_6\)-audit compatibility record on the overlaps. Absent these, the gluing defect blocks the global claim, and Subsubsection \hyperref[sn:5.3.4]{5.3.4} records the gluing-defect schema used by \(B_{\mathrm{glue}}\).

\subsection{Countermodel 4: Directed Cell Active but Pair Observable Blocked}\label{sn:9.5}

\subsubsection{Construction}\label{sn:9.5.1}

The active cell of this construction is the running example of Subsection \hyperref[sn:3.8]{3.8}.

Take the primitive pair \(P_6\leftarrow P_3\) in \(\mathbf{BirdInt}^{\mathrm{aud}}_{\mathrm{fin}}\), with the directed-cell schema of Subsection \hyperref[sn:4.3]{4.3}. Equip the cell with witness and update data

\[
W_3\;=\;\mathsf{RouteMismatchWit},\qquad U_6\;=\;\mathsf{AddAuditRecord},
\]

drawn from \(\mathsf{Wit}(P_3)\) and \(\mathsf{Upd}(P_6)\) respectively. Suppose the witness is visible under the surrounding instrument \(I\) and audited under \(\mathsf{AuditPolicy}_I\), and that the cell's profile \(\lambda_{\mathrm{audit}}\) records the audit-only instrument mode for an audit update. The classifier of Subsubsection \hyperref[sn:5.6.2]{5.6.2} then assigns

\[
\mathsf{Cell}_{63}^{\lambda_{\mathrm{audit}}}\;:\;\texttt{action}.
\]

Now consider the unordered pair-observable judgment

\[
\mathsf{PairObs}_{\{3,6\}}^{\mu_{\mathrm{drive}}}
\]

with profile \(\mu_{\mathrm{drive}}=(\mathsf{Ver},\mathsf{Ver},\texttt{fine},\texttt{drive-certification},\texttt{drive},\texttt{audit-only},\texttt{local})\), whose compatibility tuple \((\texttt{fine},\texttt{drive-certification},\texttt{local},\mathsf H_{\mathrm{fin}})\) lies in \(\mathsf{Compat}_I\), and a pair record whose joint witness is \(W_3\), whose source slot carries an audited absence, and whose visibility and audit checks pass and whose defect record correctly contains \(\delta_6^{\mathrm{source}}=\texttt{missing}\) and \(\delta_{\mathrm{reqbridge}}=\{\texttt{drive}\}\); its regime \(\texttt{drive-certification}\) gives \(\mathsf{ReqBridge}(Q)=\{\texttt{drive}\}\) (Subsubsection \hyperref[sn:5.5.2]{5.5.2}), so the pair can be real only with an admissible drive bridge. Give the pair a branch-compatibility record with \(\mathsf{compatible}=\mathrm{true}\) and its own passing audit. Suppose no drive bridge is present: \(B_{3\to 6}^{\mathrm{drive}}\) is absent from \(\mathcal L\), and no admissible drive source-of-truth certifying \([a]\neq0\) is recorded under \(\mu_{\mathrm{drive}}\). The cell is well-formed and active, but the pair record's source defect is \(\delta_6^{\mathrm{source}}=\texttt{missing}\), which fails the source threshold; so the first condition of the \(\texttt{blocked}\) row of the pair-observable table in Subsubsection \hyperref[sn:5.6.6]{5.6.6} holds, and the pair classifier of Subsubsection \hyperref[sn:5.6.3]{5.6.3} assigns

\[
\mathsf{PairObs}_{\{3,6\}}^{\mu_{\mathrm{drive}}}\;:\;\texttt{blocked}.
\]

The same primitive pair \(\{3,6\}\) thus carries an active directed cell and a blocked pair observable simultaneously.

\subsubsection{Supported No-Go}\label{sn:9.5.2}

The construction supports the dependency judgment

\[
\Gamma\,;\;\mathsf H_{\mathrm{fin}}\,;\;I_{\mathrm{fin}}\;\vdash\;\mathsf{Cell}_{63}^{\lambda_{\mathrm{audit}}}:\texttt{action}\;\mathrel{\mathsf{dep}_{\texttt{insufficient}}}\;\mathsf{PairObs}_{\{3,6\}}^{\mu_{\mathrm{drive}}}:\texttt{real},
\]

recording that directed-cell action does not by itself imply pair-observable realness on the same primitive pair.

\subsubsection{What It Blocks}\label{sn:9.5.3}

The overclaim ``directed-cell action implies pair-observable realness'' is rejected. NC-5 of Subsection \hyperref[sn:15.2]{15.2} records the corresponding nonclaim, and Theorem \hyperref[cl:thm:5]{5} of Section \hyperref[sn:6]{6} makes the cell-to-pair non-implication formal as part of the status-family separation theorem.

\subsubsection{What It Does Not Block}\label{sn:9.5.4}

The construction does not say pair observables are unattainable on \(\{3,6\}\) or that audit cells are incompatible with pair realness. A different pair record on the same primitives, with an admissible source and the required drive bridge, may classify \(\texttt{real}\) or \(\texttt{real\_provisional}\) in the schema of Subsubsection \hyperref[sn:4.4.3]{4.4.3}.

\subsubsection{Required Positive Bridge}\label{sn:9.5.5}

A positive pair-observable claim requires its own typed records: a real or fallback source under \(\mathsf{SourcePolicy}_I\), an admissible branch-compatibility record, the visibility, threshold, and audit data of the pair record, and (when the pair claims drive content) the drive bridge \(B_{3\to 6}^{\mathrm{drive}}\) under \(\mathsf{BridgePolicy}_I\). The directed-cell record contributes none of these; the pair record carries them as separate data.

\subsection{Countermodel 5: Refinement Worsens Closure}\label{sn:9.6}

\subsubsection{Construction}\label{sn:9.6.1}

Take a finite carrier \(X=\{a,b,c,d\}\) on \(\mathsf H_{\mathrm{fin}}\). Define a coarse abstraction

\[
q_0\colon X\to\{*\},\qquad q_0(x)\;=\;*\;\text{ for every }x\in X,
\]

so that \(q_0\) has a single fiber and every finite update on \(X\) descends through \(q_0\) trivially.

Now define a refinement

\[
q_1\colon X\to\{A,B,C\},\qquad q_1(a)=A,\quad q_1(b)=A,\quad q_1(c)=B,\quad q_1(d)=C.
\]

The coarse abstraction \(q_0\) factors through \(q_1\) via the unique forgetting map \(r:\{A,B,C\}\to\{*\}\), so \(q_0=rq_1\). Thus the package data of \(q_1\) is more refined than that of \(q_0\) in the sense of Subsection \hyperref[sn:3.3]{3.3}. Record the package refinement \(\Phi=(\mathrm{id}_X,r)\) from the package with lens \(q_0\) to the package with lens \(q_1\). The latter has a source record holding the tables of \(q_0\) and \(q_1\), and an audit family containing the audit entry of the package with lens \(q_0\), so that \(A_{\mathcal T_0}\subseteq A_{\mathcal T_1}\) and \(\Phi\) is admissible (Subsubsection \hyperref[sn:3.3.2]{3.3.2}).

Now choose a finite update \(F:X\to X\) that creates a \(q_1\)-split. For instance, take

\[
F(a)\;=\;c,\quad F(b)\;=\;d,\quad F(c)\;=\;c,\quad F(d)\;=\;d.
\]

Then \(q_1(a)=q_1(b)=A\), but \(q_1F(a)=q_1(c)=B\) and \(q_1F(b)=q_1(d)=C\), so

\[
q_1F(a)\;\neq\;q_1F(b),
\]

and \((a,b)\in\delta_1^{\mathrm{split}}(q_1,F)\neq\varnothing\). By Theorem \hyperref[cl:thm:7]{7}, no descended update through \(q_1\) exists.

By contrast, every update on \(X\) --- including this one --- descends through \(q_0\), since \(q_0\) has a single fiber and the descended map on \(\{*\}\) is the identity. Thus the descent defect \(\delta_1^{\mathrm{split}}(q_0,F)\) is empty, while \(\delta_1^{\mathrm{split}}(q_1,F)\) is non-empty: the refinement from \(q_0\) to \(q_1\) has worsened the closure status of the same update.

\subsubsection{Supported No-Go}\label{sn:9.6.2}

The construction supports the dependency judgment

\[
\Gamma\,;\;\mathsf H_{\mathrm{fin}}\,;\;I_{\mathrm{fin}}\;\vdash\;P_4(q_0\leadsto q_1)\;\mathrel{\mathsf{dep}_{\texttt{insufficient}}}\;\bigl(\delta_1^{\mathrm{split}}(q_1,F)\;\subseteq\;\delta_1^{\mathrm{split}}(q_0,F)\bigr),
\]

recording that \(P_4\)-refinement does not by itself improve closure: a refinement can introduce a non-empty \(\delta_1^{\mathrm{split}}\) where there was none before.

\subsubsection{What It Blocks}\label{sn:9.6.3}

The overclaim ``finer is always better'' is rejected. NC-13 of Subsection \hyperref[sn:15.2]{15.2} records the corresponding nonclaim. The construction is the working content of NC-13: refinement can introduce closure defects rather than improve them.

\subsubsection{What It Does Not Block}\label{sn:9.6.4}

The construction does not say refinement is always harmful or that closure improvement under refinement is impossible. A different finite update \(F'\) --- one whose action is consistent with the fibers of \(q_1\) --- would descend through \(q_1\) as well as \(q_0\). The construction shows only that compatibility of \(F\) with \(q_1\) is independent of compatibility with \(q_0\).

\subsubsection{Required Positive Bridge}\label{sn:9.6.5}

A positive refinement-improvement theorem requires monotonicity hypotheses on the refinement, the update, and the closure functional, recorded as a \(P_4\)-monotonicity bridge. Such hypotheses are not asserted in this paper.

\subsection{Countermodel 6: Same-Level Self-Audit Circularity}\label{sn:9.7}

\subsubsection{Construction}\label{sn:9.7.1}

Let \(I_0\) be a finite admissible instrument under an admissible package \(\mathcal T_0\) on the host \(\mathsf H_{\mathrm{instr}}\), with the instrument record schema of Subsubsection \hyperref[sn:3.4.1]{3.4.1}. The instance also contains an audit record \(A_0\) with \(\mathsf{ClaimRef}=\mathcal T_0\), \(\mathsf{CheckRules}\) drawn from \(\mathsf{CheckRules}_{I_0}\) and \(\delta_6^{\mathrm{audit}}(A_0)=\texttt{audit\_passes}\), replaying the package audit entry of \(\mathcal T_0\); so ``\(I_0\) audits \(\mathcal T_0\)'' holds in the sense of Subsubsection \hyperref[sn:9.1.1]{9.1.1}. Consider the candidate same-level self-soundness claim

\[
\Gamma\,;\;\mathcal T_0\,;\;I_0\;\vdash\;\mathrm{Sound}(I_0)\;:\;?,
\]

asserting that \(I_0\) certifies its own soundness without a level shift. Two cases arise from the classifier of Subsubsection \hyperref[sn:5.6.5]{5.6.5}.

\subsubsection{Cases}\label{sn:9.7.2}

\textbf{Outside-scope case.} Suppose the claim's host tag is not in \(\mathsf{Hosts}_{I_0}\), a level tag of its profile is not in \(\mathsf{Levels}_{I_0}\), \(\texttt{soundness}\notin\mathsf{ClaimTypes}(I_0)\), or \(\mathsf{Use}(\mathrm{Sound}(I_0))\not\subseteq\mathsf{Scope}_{I_0}\). The first classifier row assigns

\[
\Gamma\,;\;\mathcal T_0\,;\;I_0\;\vdash\;\mathrm{Sound}(I_0)\;:\;\texttt{outside\_scope}.
\]

This is the case for instruments configured to evaluate object-level claims only.

\textbf{Circular case.} Suppose the claim's host tag is in \(\mathsf{Hosts}_{I_0}\), its level tags are in \(\mathsf{Levels}_{I_0}\), \(\texttt{soundness}\in\mathsf{ClaimTypes}(I_0)\) and the use-set of \(\mathrm{Sound}(I_0)\) is contained in \(\mathsf{Scope}_{I_0}\); the use-set then contains the verdict reference \(\ulcorner\mathrm{Sound}(I_0),I_0\urcorner\) (Subsubsection \hyperref[sn:11.5.1]{11.5.1}). When no higher-priority \(\texttt{blocked}\), \(\texttt{absent\_with\_record}\) or \(\texttt{failed\_audit}\) condition is present, the classifier assigns

\[
\Gamma\,;\;\mathcal T_0\,;\;I_0\;\vdash\;\mathrm{Sound}(I_0)\;:\;\texttt{undefined\_circular}.
\]

This is the case when an instrument is asked to certify a property of itself whose verification cannot be carried out without already knowing the verdict.

In neither case does \(I_0\) certify its own soundness at \(\texttt{accepted}\) without a level shift.

\subsubsection{Supported No-Go}\label{sn:9.7.3}

The construction supports the dependency judgment

\[
\Gamma\,;\;\mathsf H_{\mathrm{instr}}\,;\;I_{\mathrm{instr}}\;\vdash\;I_0\text{ audits }\mathcal T_0\;\mathrel{\mathsf{dep}_{\texttt{blocks}}}\;\bigl(I_0\vdash\mathrm{Sound}(I_0):\texttt{accepted}\bigr),
\]

recording, by Theorem \hyperref[cl:thm:24]{24}, that in every admissible instance an instrument's audit capacity for object-level content is accompanied by a non-\(\texttt{accepted}\) verdict on its own same-level soundness.

\subsubsection{What It Blocks}\label{sn:9.7.4}

The overclaim ``an instrument can certify its own complete soundness without a level shift'' is rejected within this calculus: the claim classifier routes that claim to \(\texttt{outside\_scope}\), \(\texttt{undefined\_circular}\), or a status of higher priority, and never to \(\texttt{accepted}\). NC-16 of Subsection \hyperref[sn:15.2]{15.2} records the corresponding nonclaim, and Theorem \hyperref[cl:thm:24]{24} of Section \hyperref[sn:11]{11} states this classifier routing formally. The construction concerns this classifier; it is not a general impossibility result about self-certification.

\subsubsection{What It Does Not Block}\label{sn:9.7.5}

The construction does not say self-audit data is unattainable or that instruments cannot record any self-referential information. Object-level audit updates within \(P_6\leftarrow P_6\) (auditing a specific object-level record) and fixed audit refreshes (re-running an existing audit policy on existing records) remain admissible: Subsection \hyperref[sn:11.6]{11.6} records the three-case separation between object-level audit, fixed audit refresh, and complete same-level self-audit. Only in the third case does the complete self-soundness claim itself force a status that precedes \(\texttt{accepted}\) in the classifier's priority order; the first two cases receive whatever status the cell classifier assigns to their records.

\subsubsection{Required Positive Repair}\label{sn:9.7.6}

A positive same-level soundness claim requires a level shift: a soundness claim whose use-set contains a verdict reference of an instrument at the judging instrument's level or above receives \(\texttt{undefined\_circular}\) or a status of higher priority, never \(\texttt{accepted}\) (Subsubsection \hyperref[sn:5.6.5]{5.6.5}; Theorem \hyperref[cl:thm:24]{24}), so the claim must be judged by a higher instrument \(I_1\); an admissible audit bridge \(B^{\mathrm{audit}}_{0\to 1}\) in the rotating-audit chain of Subsection \hyperref[sn:11.7]{11.7} records such an audit. The chain is finite and admissible; Theorem \hyperref[cl:thm:25]{25} of Section \hyperref[sn:11]{11} records that an extension instrument can audit the lower stack, and by its clause 4 accept an upper-level soundness claim for a lower instrument exactly when that instrument's recorded verdicts are correct, but does not certify its own soundness, so the rotating audit is non-final (NC-17 of Subsection \hyperref[sn:15.2]{15.2}).

\subsection{Additional Countermodels for Appendix}\label{sn:9.8}

The remaining six countermodels of the full atlas are recorded with their full constructions and dependency judgments in Appendix \hyperref[sn:C]{C}; they are named here for cross-reference. Each uses the sheet form of Subsubsection \hyperref[sn:9.1.1]{9.1.1}.

\subsubsection{Same Primitive Pair, Different Statuses}\label{sn:9.8.1}

\(\mathsf{CM}_7\) exhibits two directed-cell records on the same primitive pair \(P_6\leftarrow P_3\): one with audit-only profile \(\lambda_{\mathrm{audit}}\) classified \(\texttt{action}\), the other with drive-certification profile and a registered drive bridge too weak for it, classified \(\texttt{blocked}\). Supports NC-3 (no primitive-pair status determination); the construction is the working content of Theorem \hyperref[cl:thm:1]{1} of Section \hyperref[sn:6]{6}. Full sheet in Appendix \hyperref[sn:C]{C}.

\subsubsection{Idempotent Completions Do Not Commute}\label{sn:9.8.2}

\(\mathsf{CM}_8\) exhibits two finite idempotent completion operators \(E_1,E_2\) on the power set \(\mathcal P(\{a,b,c\})\) that fail to commute. Supports NC-8 (no idempotence-implies-commutation); the construction is the working content of Theorem \hyperref[cl:thm:9]{9} of Section \hyperref[sn:7]{7}. Full sheet in Appendix \hyperref[sn:C]{C}.

\subsubsection{Strict Extension Without Macro Closure}\label{sn:9.8.3}

\(\mathsf{CM}_9\) exhibits a finite strict object map, with nonfactorization \(\pi_{j+1}\not\factor\pi_j\), and a finite update with no descended dynamics on the target package. Supports NC-10 (no strictness-implies-closure). Full sheet in Appendix \hyperref[sn:C]{C}.

\subsubsection{Strict Extension Without Drive}\label{sn:9.8.4}

\(\mathsf{CM}_{10}\) exhibits a finite strict object map whose support graph is a forest, so \(H^1=0\) and no admissible cohomological drive certificate exists. Supports NC-11 (no strictness-implies-drive). Full sheet in Appendix \hyperref[sn:C]{C}.

\subsubsection{Fallback Source Cannot Support Real Pair}\label{sn:9.8.5}

\(\mathsf{CM}_{11}\) exhibits a pair-observable record whose only available source has type \(\texttt{fallback}\) and which carries no audited source-upgrade bridge admitted by \(\mathsf{SourcePolicy}_I\) and \(\mathsf{BridgePolicy}_I\). The classifier of Subsubsection \hyperref[sn:5.6.3]{5.6.3} assigns at most \(\texttt{real\_provisional}\), never \(\texttt{real}\). Supports the source-of-truth discipline of Subsubsection \hyperref[sn:4.4.3]{4.4.3}. Full sheet in Appendix \hyperref[sn:C]{C}.

\subsubsection{Gating Destroys Cycle-Supported Drive}\label{sn:9.8.6}

\(\mathsf{CM}_{12}\) exhibits a triangle graph with a nonzero drive cochain and a \(P_2\)-style gate that deletes one edge, leaving a path with \(\beta_1=0\) and hence \(H^1=0\). Supports the gating-suppression discipline of Section \hyperref[sn:8]{8} and is the working content of Theorem \hyperref[cl:thm:18]{18}. Full sheet in Appendix \hyperref[sn:C]{C}.
 \section{Promotion and Stacking}\label{sn:10}

This section gives the finite promotion-bridge theory in scoped form. Subsection \hyperref[sn:10.1]{10.1} recalls the promotion-bridge object of Subsection \hyperref[sn:4.5]{4.5} and adds the gate record. Subsection \hyperref[sn:10.2]{10.2} specifies the gate semantics. Subsections \hyperref[sn:10.3]{10.3}, \hyperref[sn:10.4]{10.4}, and \hyperref[sn:10.5]{10.5} prove the three promotion theorems: gate soundness, strict-implies-nonfactorization, and non-strict-implies-factorization. Subsection \hyperref[sn:10.6]{10.6} proves Theorem \hyperref[cl:thm:22]{22}, the admissibility of a small bridge that is then classified as strict. Subsection \hyperref[sn:10.7]{10.7} lists what this paper does not prove about stacking.

\subsection{Promotion Bridge Object}\label{sn:10.1}

\subsubsection{Bridge Data}\label{sn:10.1.1}

A promotion bridge \(B_{j\to j+1}\) from a source theory package \(\mathcal T_j\) to a target theory package \(\mathcal T_{j+1}\) is the finite typed record of Subsubsection \hyperref[sn:4.5.2]{4.5.2}, in its expanded form with identifier, carrier \(S\), codomains \(O_j,O_{j+1}\), object maps \(\pi_j:S\to O_j\) and \(\pi_{j+1}:S\to O_{j+1}\), level/lens/interface record \(L\), visibility record \(V\), threshold record \(\Theta\), audit record \(A\), defect record \(\delta\), gate-results record \(\mathsf{GateResults}(B)\), nonclaim record \(\mathcal N\), host tag \(h_B\in\mathsf{Host}\) and bridge profile \(\lambda_{\mathrm{prom}}\). Every component is finite; the defect record \(\delta\) collects the strictness, no-smuggling, stability, and descent defects of Subsection \hyperref[sn:5.4]{5.4}.

\subsubsection{Strict Versus Non-Strict Promotion}\label{sn:10.1.2}

A promotion bridge has \textbf{strict object maps} when

\[
\pi_{j+1}\;\not\factor\;\pi_j,
\]

equivalently when there exist \(s,s'\in S\) with \(\pi_j(s)=\pi_j(s')\) and \(\pi_{j+1}(s)\neq\pi_{j+1}(s')\). A promotion bridge has \textbf{non-strict object maps} when \(\pi_{j+1}\) factors through \(\pi_j\), that is, when there exists a finite map \(\phi:\mathrm{im}(\pi_j)\to O_{j+1}\) with \(\pi_{j+1}=\phi\,\pi_j\). For any object maps \(\pi_j,\pi_{j+1}\) on \(S\), the factorization dichotomy is exhaustive: by Theorem \hyperref[cl:thm:13]{13} of Section \hyperref[sn:7]{7}, the factorization defect \(\Delta_{\mathrm{fact}}(\pi_j,\pi_{j+1})\) is empty exactly when \(\pi_{j+1}\) factors through \(\pi_j\), and non-empty exactly when it does not. Strictness in this sense does not require \(\pi_j\) to factor through \(\pi_{j+1}\). When \(\pi_j=r\circ\pi_{j+1}\) for some map \(r\), strictness is equivalent to \(\{(s,s'):\pi_{j+1}(s)=\pi_{j+1}(s')\}\subsetneq\{(s,s'):\pi_j(s)=\pi_j(s')\}\): the new object map strictly refines the old one. Theorem \hyperref[cl:thm:22]{22}, \(\mathsf{CM}_9\), and the two-point witness realization \(W=\{x,x'\}\) in Model Theorem \hyperref[cl:thm:34]{34} have this property, since \(\pi_0\) is constant on those carriers.

Strictness is a property of the typed object-map data \(\pi_j,\pi_{j+1}\) alone; it is not by itself a closure or drive claim. The nonclaim register of Subsection \hyperref[sn:15.2]{15.2} records NC-10 (no strictness-implies-closure) and NC-11 (no strictness-implies-drive); the no-go countermodels \(\mathsf{CM}_9\) and \(\mathsf{CM}_{10}\) of Appendix \hyperref[sn:C]{C} (named in Section \hyperref[sn:9]{9}) supply finite explicit instances in which strictness coexists with absent closure or absent drive.

\subsection{Promotion Gates}\label{sn:10.2}

\subsubsection{Gate List}\label{sn:10.2.1}

The required gate family of a promotion bridge is

\[
\mathsf{PromGate}(B)\;=\;(\,G_{\mathrm{suff}},\;G_{\mathrm{desc}},\;G_{\mathrm{stab}},\;G_{\mathrm{ctrl}},\;G_{\mathrm{nosmuggle}},\;G_{\mathrm{vis}},\;G_{\mathrm{audit}},\;G_{\mathrm{strict}}\,),
\]

with the optional local-to-global gate

\[
G_{\mathrm{locglob}}.
\]

Each gate has a fixed role:

\begin{itemize}
\tightlist
\item
  \(G_{\mathrm{suff}}\): well-formedness gate, checking that \(B\) as a record satisfies the typed schema of Subsubsection \hyperref[sn:10.1.1]{10.1.1}, that each gate-feeding entry of \(\delta\) (the strictness, descent, stability and no-smuggling defects of Subsection \hyperref[sn:5.4]{5.4}) equals its value recomputed from the bridge data as in Subsection \hyperref[sn:5.4]{5.4}, the strictness defect being \(\Delta_{\mathrm{fact}}(\pi_j,\pi_{j+1})\) computed from \(\pi_j,\pi_{j+1}\) on \(S\), and that every recorded value of a defect-fed gate (\(G_{\mathrm{strict}}\), \(G_{\mathrm{desc}}\), \(G_{\mathrm{stab}}\), \(G_{\mathrm{nosmuggle}}\), \(G_{\mathrm{ctrl}}\)) other than \(\texttt{not\_checked}\), \(\texttt{not\_required}\) or an unauditable value equals the value recomputed from its defect; and that a recorded \(G_{\mathrm{vis}}=\texttt{pass}\) holds exactly when the visibility defect \(\delta_{\mathrm{vis}}\) of the bridge's use-set, computed from \(V\) and its bridge family as in Subsection \hyperref[sn:5.5]{5.5}, is empty, and a recorded \(G_{\mathrm{audit}}=\texttt{pass}\) exactly when the audit defect of \(A\) computed as in Subsubsection \hyperref[sn:5.3.6]{5.3.6} is \(\texttt{audit\_passes}\). \(G_{\mathrm{suff}}\) also checks that \(G_{\mathrm{locglob}}=\texttt{not\_required}\) exactly when no global-package claim is made, and checks each recorded \(\texttt{glue\_pass}\), \(\texttt{local\_only}\) or \(\texttt{glue\_fail}\) against the cover, local-package admissibility and gluing defect.
\item
  \(G_{\mathrm{desc}}\): descent gate, required when the bridge claims macro closure or descended dynamics on the target package; it is fed by the descent defect of Subsubsection \hyperref[sn:5.3.1]{5.3.1} (\(\delta_1^{\mathrm{split}}\)) and Theorem \hyperref[cl:thm:7]{7}.
\item
  \(G_{\mathrm{stab}}\): stability gate, required when the bridge claims target objects are fixed points of a completion endomap; it is fed by the stability defect \(\delta_{\mathrm{stab}}\) of Subsubsection \hyperref[sn:5.4.3]{5.4.3}.
\item
  \(G_{\mathrm{ctrl}}\): control-of-content gate, checking that no target-layer data has been used to certify target novelty: \(G_{\mathrm{ctrl}}=\texttt{pass}\) iff \(\delta_{\mathrm{nosmuggle}}(B)\) has no entry with \(k=1\) (Subsubsection \hyperref[sn:5.4.2]{5.4.2}); \(G_{\mathrm{ctrl}}\) is implied by \(G_{\mathrm{nosmuggle}}\) and is listed separately so that \(\texttt{failed\_no\_smuggling}\) records its cause.
\item
  \(G_{\mathrm{nosmuggle}}\): no-smuggling gate, fed by the no-smuggling defect of Subsubsection \hyperref[sn:5.4.2]{5.4.2}.
\item
  \(G_{\mathrm{vis}}\): visibility gate, checking that the bridge's content is admissible under \(\mathsf{Visible}_I\) and the bridge policy of \(I\).
\item
  \(G_{\mathrm{audit}}\): audit gate, checking that the audit record \(A\) passes \(\mathsf{AuditPolicy}_I\).
\item
  \(G_{\mathrm{strict}}\): strictness gate, with refined status set \(\mathrm{StrictGateStatus}\) from Subsubsection \hyperref[sn:5.1.6]{5.1.6}, fed by the factorization defect \(\Delta_{\mathrm{fact}}\) of Subsubsection \hyperref[sn:5.4.1]{5.4.1}.
\item
  \(G_{\mathrm{locglob}}\): optional local-to-global gate, present when the bridge claims a global package built from local packages on a cover.
\end{itemize}

The required core gates of a bridge are the subset of

\[
\mathsf{CoreGates}\;=\;\{\,G_{\mathrm{suff}},\,G_{\mathrm{desc}},\,G_{\mathrm{stab}},\,G_{\mathrm{ctrl}},\,G_{\mathrm{nosmuggle}},\,G_{\mathrm{vis}},\,G_{\mathrm{audit}}\,\}
\]

given by

\[
\begin{aligned}
\mathsf{ReqCoreGates}(B)\;=\;&\{\,G_{\mathrm{suff}},\,G_{\mathrm{ctrl}},\,G_{\mathrm{nosmuggle}},\,G_{\mathrm{vis}},\,G_{\mathrm{audit}}\,\}\\
&\cup\;\{\,G_{\mathrm{desc}}:B\text{ claims macro closure or descended dynamics}\,\}\\
&\cup\;\{\,G_{\mathrm{stab}}:B\text{ claims fixed-point targets}\,\}.
\end{aligned}
\]

The content and target fields declare whether \(B\) claims macro closure or descended dynamics or fixed-point targets. Such a declaration requires, respectively, a threshold on \(\delta_1^{\mathrm{split}}\) (Subsubsection \hyperref[sn:5.3.1]{5.3.1}) or on \(\delta_{\mathrm{stab}}\) (Subsubsection \hyperref[sn:5.4.3]{5.4.3}) in \(\Theta\); a missing threshold fails \(G_{\mathrm{suff}}\). The presence of \(\Theta_{\mathrm{strict}}\) or \(\Theta_{\mathrm{nonstrict}}\) in \(\Theta\) (Subsubsection \hyperref[sn:5.4.1]{5.4.1}) determines whether a strictness verdict is required. Since the bridge audit's \(\mathsf{ThresholdCheck}\) tests \(\Theta\), a bridge whose \(\Theta\) contains \(\Theta_{\mathrm{strict}}\) (resp. \(\Theta_{\mathrm{nonstrict}}\)) while \(\Delta_{\mathrm{fact}}=\varnothing\) (resp. \(\neq\varnothing\)) has \(\texttt{threshold\_failure}\) in its audit defect, so \(G_{\mathrm{audit}}\) fails and the classifier returns \(\texttt{failed\_audit}\) or a status of higher priority; the threshold declares which refinement the bridge claims. \(\mathsf{GateResults}(B)\) must mark the corresponding gate as required, and \(G_{\mathrm{suff}}\) checks that agreement. The five gates of the first set are required of every bridge. The strictness gate \(G_{\mathrm{strict}}\) is part of the promotion-gate record but sits outside \(\mathsf{CoreGates}\); when strictness is required and classified, it records the strict/non-strict refinement, while a bridge that does not require a strictness verdict may remain at the unrefined status \(\texttt{accepted}\).

\emph{Decidability of the well-formedness gate.} \(G_{\mathrm{suff}}\) performs finitely many checks: the typing of the seventeen fields of the expanded bridge record, \(\mathsf{AdmPkg}\) of both packages, \(\mathsf{AdmProfile}(\lambda,I,\mathcal T_j)\) for the bridge profile, finiteness of \(S\), \(O_j\) and \(O_{j+1}\), the agreement of each defect-fed gate value, including \(G_{\mathrm{ctrl}}\), and of the recorded values of \(G_{\mathrm{vis}}\) and \(G_{\mathrm{audit}}\), with its recomputed visibility or audit defect, and the agreement of the required marks of \(G_{\mathrm{desc}}\), \(G_{\mathrm{stab}}\) and \(G_{\mathrm{strict}}\) with the claims determined by the bridge's content, target and threshold fields. Each is a membership or comparison test on finite records, so, under hypotheses (a)--(b) of Theorem \hyperref[cl:thm:3]{3} for the thresholds these checks evaluate, whether a promotion bridge is well formed is decidable by the argument of Theorem \hyperref[cl:thm:3]{3}.

\subsubsection{Gate Statuses}\label{sn:10.2.2}

Each ordinary gate carries a status from the gate status set \(\mathrm{GateStatus}\) of Subsubsection \hyperref[sn:5.1.6]{5.1.6}.

The strictness gate carries a refined status from the set \(\mathrm{StrictGateStatus}\) of Subsubsection \hyperref[sn:5.1.6]{5.1.6}.

The gate-results record \(\mathsf{GateResults}(B)\) is the finite assignment of these statuses to the gates in the bridge record. A gate may be marked \(\texttt{not\_required}\) when the bridge does not claim the corresponding content; the descent gate, for example, is \(\texttt{not\_required}\) for a bridge that records only an interface change with no closure claim. A required core gate marked \(\texttt{not\_checked}\) does not yield an accepted promotion; the promotion classifier of Subsubsection \hyperref[sn:5.6.4]{5.6.4} routes such a bridge to \(\texttt{candidate}\).

\subsection{Theorem 19: Promotion Gate Soundness}\label{sn:10.3}

\begin{claimtheorem}{Theorem 19 (PromotionGateSoundness)}\phantomsection\label{cl:thm:19}

Let \(B_{j\to j+1}\) be a promotion bridge. If the promotion classifier of Subsubsection \hyperref[sn:5.6.4]{5.6.4} assigns

\[
\mathrm{Promote}(B):\sigma\quad\text{with}\quad\sigma\in\{\,\texttt{accepted},\,\texttt{strict},\,\texttt{non\_strict}\,\},
\]

then every required core gate exists in the gate-results record and passes:

\[
\forall\,G\in\mathsf{ReqCoreGates}(B),\quad\mathsf{GateExists}(B,G)\;\wedge\;\mathsf{GatePass}(B,G),
\]

that is, \(G\) has an entry in \(\mathsf{GateResults}(B)\) with value \(\texttt{pass}\); moreover \(G_{\mathrm{locglob}}\in\{\texttt{glue\_pass},\texttt{not\_required}\}\), so a verdict in this family never rests on an unchecked local-to-global claim.

In particular \(G_{\mathrm{suff}}\), \(G_{\mathrm{ctrl}}\), \(G_{\mathrm{nosmuggle}}\), \(G_{\mathrm{vis}}\) and \(G_{\mathrm{audit}}\) pass, and, since \(G_{\mathrm{suff}}\) checks the recorded values, \(\delta_{\mathrm{vis}}(B)=\varnothing\) and the audit defect of \(A\) is \(\texttt{audit\_passes}\).

\end{claimtheorem}

\begin{proof}

By the definition of the promotion classifier in Subsubsection \hyperref[sn:5.6.4]{5.6.4}, each accepted-family row requires a well-formed bridge, every required core gate to exist and pass, and \(G_{\mathrm{locglob}}\in\{\texttt{glue\_pass},\texttt{not\_required}\}\). The verdicts \(\texttt{strict}\) and \(\texttt{non\_strict}\) are refinements of the accepted core by the strictness gate: \(\texttt{strict}\) adds \(G_{\mathrm{strict}}=\texttt{strict\_pass}\), and \(\texttt{non\_strict}\) adds \(G_{\mathrm{strict}}=\texttt{non\_strict\_pass}\). The remaining promotion-status values record specific failure or unresolved conditions and are routed by the classifier to non-accepted statuses such as \(\texttt{failed\_audit}\), \(\texttt{failed\_no\_smuggling}\), \(\texttt{failed\_descent}\), \(\texttt{failed\_stability}\), \(\texttt{globally\_obstructed}\), \(\texttt{local\_only}\), \(\texttt{candidate}\), or \(\texttt{outside\_scope}\). Rows 9--11 of the classifier also require \(G_{\mathrm{locglob}}\in\{\texttt{glue\_pass},\texttt{not\_required}\}\). Hence whenever the verdict lies in the accepted family, every required core gate exists and passes and the local-to-global gate is passed or not required; the five gates listed last belong to \(\mathsf{ReqCoreGates}(B)\) for every bridge, and the passing \(G_{\mathrm{suff}}\) makes the recorded values of \(G_{\mathrm{vis}}\) and \(G_{\mathrm{audit}}\) agree with the computed visibility and audit defects.

The result records the basic discipline of promotion: an accepted, strict or non-strict promotion verdict requires every required core gate to pass its finite admissibility check and the local-to-global gate to pass or not be required.

\end{proof}

\subsection{Theorem 20: Strict Promotion Implies Nonfactorization}\label{sn:10.4}

\begin{claimtheorem}{Theorem 20 (StrictPromotionImpliesNonfactorization)}\phantomsection\label{cl:thm:20}

Let \(B_{j\to j+1}\) be a promotion bridge. If the promotion classifier assigns

\[
\mathrm{Promote}(B_{j\to j+1})\;:\;\texttt{strict},
\]

then \(\pi_{j+1}\not\factor\pi_j\). Equivalently, there exist \(s,s'\in S\) with \(\pi_j(s)=\pi_j(s')\) and \(\pi_{j+1}(s)\neq\pi_{j+1}(s')\).

\end{claimtheorem}

\begin{proof}

The verdict \(\texttt{strict}\) requires the accepted core, in which \(G_{\mathrm{suff}}\) passes, that is, \(\mathrm{WF}(B)\) holds, so the recorded value of \(G_{\mathrm{strict}}\) agrees with its recomputed value; by the strictness-gate semantics of Subsubsection \hyperref[sn:5.4.1]{5.4.1}, the strict-pass refinement \(G_{\mathrm{strict}}=\texttt{strict\_pass}\) is then recorded exactly when \(\Delta_{\mathrm{fact}}(\pi_j,\pi_{j+1})\neq\varnothing\). The promotion classifier of Subsubsection \hyperref[sn:5.6.4]{5.6.4} produces the verdict \(\texttt{strict}\) only when \(G_{\mathrm{strict}}=\texttt{strict\_pass}\), so \(\Delta_{\mathrm{fact}}\neq\varnothing\). By Theorem \hyperref[cl:thm:13]{13} of Section \hyperref[sn:7]{7}, \(\Delta_{\mathrm{fact}}\neq\varnothing\) is equivalent to \(\pi_{j+1}\not\factor\pi_j\), and a witness pair \((s,s')\) exists by definition of \(\Delta_{\mathrm{fact}}\).

\end{proof}

\subsection{Theorem 21: Non-Strict Promotion Implies Factorization}\label{sn:10.5}

\begin{claimtheorem}{Theorem 21 (NonStrictPromotionImpliesFactorization)}\phantomsection\label{cl:thm:21}

Let \(B_{j\to j+1}\) be a promotion bridge. If the promotion classifier assigns

\[
\mathrm{Promote}(B_{j\to j+1})\;:\;\texttt{non\_strict},
\]

then \(\pi_{j+1}\factor\pi_j\). There exists a finite map \(\phi:\mathrm{im}(\pi_j)\to O_{j+1}\) with \(\pi_{j+1}=\phi\,\pi_j\).

\end{claimtheorem}

\begin{proof}

As in the proof of Theorem \hyperref[cl:thm:20]{20}, \(G_{\mathrm{suff}}\) passes in the accepted core, that is, \(\mathrm{WF}(B)\) holds, so by the strictness-gate semantics of Subsubsection \hyperref[sn:5.4.1]{5.4.1}, the non-strict-pass refinement \(G_{\mathrm{strict}}=\texttt{non\_strict\_pass}\) is recorded exactly when \(\Delta_{\mathrm{fact}}(\pi_j,\pi_{j+1})=\varnothing\). The promotion classifier of Subsubsection \hyperref[sn:5.6.4]{5.6.4} produces the verdict \(\texttt{non\_strict}\) only when \(G_{\mathrm{strict}}=\texttt{non\_strict\_pass}\), so \(\Delta_{\mathrm{fact}}=\varnothing\). By Theorem \hyperref[cl:thm:13]{13} of Section \hyperref[sn:7]{7}, \(\Delta_{\mathrm{fact}}=\varnothing\) is equivalent to \(\pi_{j+1}\factor\pi_j\), and Theorem \hyperref[cl:thm:12]{12} supplies the explicit factorization \(\phi\) on \(\mathrm{im}(\pi_j)\).

\end{proof}

\subsection{Theorem 22: A Two-Point Strict Bridge}\label{sn:10.6}

\subsubsection{Data}\label{sn:10.6.1}

\emph{Setting.} The host is \(\mathsf H_{\mathrm{fin}}\) and the instrument is \(I_{\mathrm{prom}}\); the bridge schema is that of Subsection \hyperref[sn:10.1]{10.1} and the gates are those of Subsection \hyperref[sn:10.2]{10.2}. The gate results of the bridge, and its classification as \(\texttt{strict}\) by the promotion classifier, are recorded in Subsubsection \hyperref[sn:10.6.2]{10.6.2}.

\begin{claimtheorem}{Theorem 22 (Toy22StrictBridge)}\phantomsection\label{cl:thm:22}

There exists a finite two-layer promotion bridge \(B_{0\to 1}^{\mathrm{Toy22}}\) on the host \(\mathsf H_{\mathrm{fin}}\) with carrier \(S=\{a,b\}\) and object maps

\[
\pi_0\colon S\to O_0=\{*\},\qquad \pi_1\colon S\to O_1=\{u,v\}
\]

defined by \(\pi_0(a)=\pi_0(b)=*\) and \(\pi_1(a)=u\), \(\pi_1(b)=v\). The bridge is admissible under a finite instrument \(I_{\mathrm{prom}}\) whose claim types include promotion claims and exclude the claim-type tag \(\texttt{soundness}\), with audit and visibility records that pass the required core gates and with no closure or drive claim attached to the bridge, and the promotion classifier assigns \(\Gamma;\mathcal T_0;I_{\mathrm{prom}}\vdash\mathrm{Promote}(B_{0\to 1}^{\mathrm{Toy22}}):\texttt{strict}\).

\end{claimtheorem}

\emph{Records of the instance.} All records live on the host \(\mathsf H_{\mathrm{fin}}\) at the level \(\mathsf{Ver}\).

\begin{itemize}
\tightlist
\item
  \emph{Packages.} For \(k=0,1\), \(\mathcal T_k=(S,\pi_k,\Sigma_{\pi_k},\mathrm{id}_{\mathcal P(S)},\mathcal A_k)\), whose completion is the identity of \(\mathcal P(S)\) rather than \(E_{\pi_k}\); both fix every member of \(\Sigma_{\pi_k}\), since \(E_{\pi_k}(\pi_k^{-1}(Y))=\pi_k^{-1}(Y)\), and the bridge claims no fixed-point targets; here \(\Sigma_{\pi_k}=\{\pi_k^{-1}(Y):Y\subseteq O_k\}\) and \(\mathcal A_k\) certifies the single audit entry of \(A_{\mathcal T_k}\), which recomputes \(\pi_k(a)\) and \(\pi_k(b)\) from the table of \(\pi_k\). The visibility profile makes every field visible to \(I_{\mathrm{prom}}\); the source record is \(\mathrm{src}_{\mathcal T_k}=(\texttt{toy22-}k,\texttt{committed\_state},\text{the table of }\pi_k,\texttt{valid})\); the defect family is empty; the nonclaim record states that no closure or drive claim is made. \(\mathcal T_0\) carries the promotion data \((\pi_0,\mathcal T_1,\pi_1,\texttt{toy22})\), naming the bridge by its identifier, accompanied by \(\mathrm{src}_{\mathcal T_0}\), the bridge audit \(A\) and \(\mathcal N\). Each condition of Subsubsection \hyperref[sn:3.3.3]{3.3.3} holds by inspection of these finite records, and the remaining components of the instance consist of the records listed in this construction, each written in its schema, so \(\mathsf{AdmDomain}\) of Subsubsection \hyperref[sn:4.10.1]{4.10.1} holds.
\item
  \emph{Instrument.} \(I_{\mathrm{prom}}\) has scope and visible set the records of \(\mathcal T_0\), \(\mathcal T_1\) and the bridge, \(\mathsf{Hosts}_{I_{\mathrm{prom}}}=\{\mathsf H_{\mathrm{fin}}\}\), \(\mathsf{Levels}_{I_{\mathrm{prom}}}=\{\mathsf{Ver}\}\), strengths \(\{\texttt{weak}<\texttt{strong}\}\) with every admitted pair and claim type mapped to \(\texttt{weak}\), and \(\mathsf{Suppressed}_{I_{\mathrm{prom}}}=\varnothing\); claim types \(\{\texttt{promotion}\}\); check rules that recompute \(\pi_0\), \(\pi_1\) on \(S\) and the factorization defect \(\Delta_{\mathrm{fact}}(\pi_0,\pi_1)\); exact thresholds and \(\mathsf{ProvisionalTypes}_{I_{\mathrm{prom}}}=\varnothing\); an audit policy requiring every recomputed table entry to agree with the recorded one; level \(\mathsf{Ver}\); a bridge policy admitting the visibility bridge type and requiring no visibility bridge; \(\mathsf{ModePolicy}_{I_{\mathrm{prom}}}(\texttt{promotion})=\texttt{default}\) and \((\texttt{fine},\texttt{structural},\texttt{local},\mathsf H_{\mathrm{fin}})\in\mathsf{Compat}_{I_{\mathrm{prom}}}\), so that \(\mathsf{AdmProfile}(\lambda_{\mathrm{prom}},I_{\mathrm{prom}},\mathcal T_0)\) holds; a source policy admitting \(\texttt{committed\_state}\); an evidence policy admitting \(\texttt{committed\_state}\) as direct evidence; and a nonclaim record excluding soundness claims.
\item
  \emph{Bridge.} \(B_{0\to1}^{\mathrm{Toy22}}=(\texttt{toy22},\mathcal T_0,\mathcal T_1,S,O_0,O_1,\pi_0,\pi_1,L,V,\Theta,A,\delta,\mathsf{GateResults},\mathcal N,\mathsf H_{\mathrm{fin}},\lambda_{\mathrm{prom}})\) with \(L\) the level pair \((\mathsf{Ver},\mathsf{Ver})\), identity lens on \(S\) and interface \((\pi_0,\pi_1)\); \(V\) marking every field visible; \(\Theta\) the exact strictness threshold \(\Delta_{\mathrm{fact}}(\pi_0,\pi_1)\neq\varnothing\); \(A\) the audit record whose target is the bridge, whose evidence is the two tables, whose check rules recompute them and \(\Delta_{\mathrm{fact}}\), whose source references are \(\mathrm{src}_{\mathcal T_0}\) and \(\mathrm{src}_{\mathcal T_1}\), and whose visibility, threshold and circularity checks pass; \(\delta\) with strictness defect \(\{(a,b),(b,a)\}\) and empty no-smuggling, stability and descent defects; \(\mathcal N\) stating that no closure or drive claim is made; and \(\lambda_{\mathrm{prom}}=(\mathsf{Ver},\mathsf{Ver},\texttt{fine},\texttt{structural},\texttt{visibility},\texttt{default},\texttt{local})\), the promotion profile at the fine scale with local use. The bridge audit tags every record of \(\mathsf{Use}(B)\) \(\texttt{other}/\texttt{other}\), except the table of \(\pi_1\), tagged \(\texttt{target}/\texttt{novelty}\), which kind 1 of Subsubsection \hyperref[sn:5.4.2]{5.4.2} exempts; its strictness comparison uses the defining object maps on the common carrier and uses no independent target-layer novelty assertion.
\end{itemize}

\subsubsection{Gate Results}\label{sn:10.6.2}

The gate-results record \(\mathsf{GateResults}(B_{0\to 1}^{\mathrm{Toy22}})\) assigns:

\begin{itemize}
\tightlist
\item
  \(G_{\mathrm{suff}}=\texttt{pass}\), since the bridge data conforms to the schema of Subsubsection \hyperref[sn:10.1.1]{10.1.1};
\item
  \(G_{\mathrm{vis}}=\texttt{pass}\), since the bridge content lies in \(\mathsf{Visible}_{I_{\mathrm{prom}}}\);
\item
  \(G_{\mathrm{audit}}=\texttt{pass}\), since the audit record passes \(\mathsf{AuditPolicy}_{I_{\mathrm{prom}}}\);
\item
  \(G_{\mathrm{ctrl}}=\texttt{pass}\), since no target-layer data has been used to certify target novelty;
\item
  \(G_{\mathrm{nosmuggle}}=\texttt{pass}\), since the no-smuggling defect of Subsubsection \hyperref[sn:5.4.2]{5.4.2} is empty;
\item
  \(G_{\mathrm{desc}}=\texttt{not\_required}\), since no closure or descended-dynamics claim is made on the bridge;
\item
  \(G_{\mathrm{stab}}=\texttt{not\_required}\), since no fixed-point structure is claimed;
\item
  \(G_{\mathrm{locglob}}=\texttt{not\_required}\), since no local-to-global packaging claim is made.
\end{itemize}

The strictness gate is fed by the factorization defect \(\Delta_{\mathrm{fact}}(\pi_0,\pi_1)\). We have \(\pi_0(a)=\pi_0(b)=*\) but \(\pi_1(a)=u\neq v=\pi_1(b)\), so \((a,b)\in\Delta_{\mathrm{fact}}(\pi_0,\pi_1)\neq\varnothing\). By Theorems \hyperref[cl:thm:12]{12} and \hyperref[cl:thm:13]{13} of Section \hyperref[sn:7]{7} and the strictness-gate semantics, \(G_{\mathrm{strict}}=\texttt{strict\_pass}\).

The promotion classifier of Subsubsection \hyperref[sn:5.6.4]{5.6.4} then assigns

\[
\mathrm{Promote}(B_{0\to 1}^{\mathrm{Toy22}})\;:\;\texttt{strict}.
\]

\begin{proof}

Let \(S=\{a,b\}\), \(O_0=\{*\}\), and \(O_1=\{u,v\}\), and define \(\pi_0,\pi_1\) as in the statement. The bridge data are finite because all three carriers are finite, and its records are those listed under \emph{Records of the instance}. The sufficiency gate passes: \(B\) satisfies the typed schema of Subsubsection \hyperref[sn:10.1.1]{10.1.1} and both packages are admissible; the recorded strictness, descent, stability and no-smuggling defects equal their values recomputed from the bridge data; the recorded values of the defect-fed gates agree with them; and \(\Theta\) contains \(\Theta_{\mathrm{strict}}\) and no threshold on \(\delta_1^{\mathrm{split}}\) or \(\delta_{\mathrm{stab}}\), so a strictness verdict is required, \(G_{\mathrm{desc}}\) and \(G_{\mathrm{stab}}\) are not, and \(\mathsf{GateResults}(B)\) marks the gates accordingly. The visibility gate passes because \(\mathsf{Suppressed}_{I_{\mathrm{prom}}}=\varnothing\) and every field is visible. The audit gate passes because the audit record recomputes the four table entries, each agreeing with the recorded value, and its source references point to committed-state sources admitted by the source policy. The control gate passes because novelty is certified by the recomputed pair \((a,b)\) of the common carrier \(S\), not by a target-layer record asserting novelty. The no-smuggling defect is empty item by item in the list of Subsubsection \hyperref[sn:5.4.2]{5.4.2}: novelty is not taken from a target-layer assertion; no content is suppressed; no simulation, local-to-global, \(P_3\)-witness or lossy-summary record is used; and the tables are exact. The descent, stability, and local-to-global gates are not required because the bridge makes no descended-dynamics, fixed-point, or local-to-global packaging claim.

It remains to verify strictness. Since \(\pi_0(a)=\pi_0(b)=*\) while \(\pi_1(a)=u\neq v=\pi_1(b)\), the pair \((a,b)\) lies in the factorization defect \(\Delta_{\mathrm{fact}}(\pi_0,\pi_1)\). Hence this defect is nonempty, so Theorems \hyperref[cl:thm:12]{12} and \hyperref[cl:thm:13]{13} identify the bridge as nonfactorizing. The strictness gate therefore returns \(\texttt{strict\_pass}\). The row \(\texttt{outside\_scope}\) does not apply: \(\mathsf H_{\mathrm{fin}}\in\mathsf{Hosts}_{I_{\mathrm{prom}}}\), the level tags of \(\lambda_{\mathrm{prom}}\) and \(L\) are \(\mathsf{Ver}\in\mathsf{Levels}_{I_{\mathrm{prom}}}\), and \(\mathsf{Use}(B)\) lies in its scope. Applying the promotion classifier to the finite gate record gives \(\mathrm{Promote}(B_{0\to 1}^{\mathrm{Toy22}}):\texttt{strict}\).

\end{proof}

\subsubsection{Interpretation}\label{sn:10.6.3}

Theorem \hyperref[cl:thm:22]{22} is a finite instance of strict promotion, not a universal stacking theorem. It records that strictness is the working content of object-map nonfactorization on a finite carrier; it does not by itself imply closure, drive, or any deeper structure on the target package. Section \hyperref[sn:13]{13}'s Cantor model-realization sheet records a scoped strict-bridge realization on the audited Cantor shell, and Subsection \hyperref[sn:10.7]{10.7} records that no general stacking composition theorem follows from the toy example.

\subsection{Stacking Claims Not Proved}\label{sn:10.7}

This subsection records two claims that this paper does not prove. The first fails in general; Subsubsection \hyperref[sn:10.7.1]{10.7.1} gives a counterexample and proves its object-map form when the middle interface refines the source; the gate and nonclaim semantics of composed bridges, and the second claim, are deferred to future work, and Subsubsection \hyperref[sn:15.3.2]{15.3.2} lists them in the formal future-work register.

\subsubsection{No General Stacking Composition Theorem}\label{sn:10.7.1}

The paper does not prove a general stacking composition theorem of the form

\[
\mathrm{Promote}(B_{j\to j+1}):\texttt{strict}\;\wedge\;\mathrm{Promote}(B_{j+1\to j+2}):\texttt{strict}\;\Longrightarrow\;\mathrm{Promote}(B_{j\to j+2}):\texttt{strict},
\]

with \(B_{j\to j+2}\) a composed bridge derived from \(B_{j\to j+1}\) and \(B_{j+1\to j+2}\). The implication fails in general: for \(S=\{a,b,c\}\), \(\pi_0\) with fibres \(\{a,b\},\{c\}\), \(\pi_1\) with fibres \(\{a\},\{b,c\}\) and \(\pi_2=\pi_0\), we have \((a,b)\in\Delta_{\mathrm{fact}}(\pi_0,\pi_1)\) and \((b,c)\in\Delta_{\mathrm{fact}}(\pi_1,\pi_2)\), while \(\Delta_{\mathrm{fact}}(\pi_0,\pi_2)=\varnothing\), so a composed bridge with passing gates recording \(\phi=\mathrm{id}\) is \(\texttt{non\_strict}\). Strictness composes when the middle interface refines the source: if \(\pi_j=r\circ\pi_{j+1}\) for a map \(r\), then \(\Delta_{\mathrm{fact}}(\pi_{j+1},\pi_{j+2})\subseteq\Delta_{\mathrm{fact}}(\pi_j,\pi_{j+2})\), since \(\pi_{j+1}(s)=\pi_{j+1}(s')\) implies \(\pi_j(s)=\pi_j(s')\); so every witness pair of the second step witnesses strictness of the composite, and \(\pi_{j+2}\not\factor\pi_{j+1}\) implies \(\pi_{j+2}\not\factor\pi_j\). What remains open is the joint gate and nonclaim semantics of composed bridges; Subsubsection \hyperref[sn:15.3.2]{15.3.2} records the corresponding future-work item.

\subsubsection{No Local-to-Global Promotion Theorem}\label{sn:10.7.2}

The paper does not prove a local-to-global promotion theorem of the form

\[
\bigl(\,\forall\,U_\alpha\in\text{cover},\;\mathrm{Promote}(B_\alpha\text{ on }U_\alpha):\texttt{strict}\,\bigr)\;\Longrightarrow\;\mathrm{Promote}(B_{\mathrm{global}}):\texttt{strict}.
\]

If a global bridge has object maps restricting to those of a strict local bridge, the local witness belongs to its global \(\Delta_{\mathrm{fact}}\). Its computed strictness test therefore passes; a recorded \(\texttt{strict\_pass}\) and a strict promotion verdict also require the global gate checks. Open are the existence of global maps (gluing, compare \(\mathsf{CM}_3\)) and the other gates. The local-to-global gate \(G_{\mathrm{locglob}}\) of Subsection \hyperref[sn:10.2]{10.2} is optional precisely because the conditions under which local strict promotions glue into a global strict promotion are not classified at theorem-grade in this paper. The countermodel \(\mathsf{CM}_3\) of Section \hyperref[sn:9]{9} records the failure case for ordinary packages without a gluing bridge on a finite cover; a positive gluing theorem for compatible local promotion bridges is recorded as future work in Subsubsection \hyperref[sn:15.3.2]{15.3.2}.
 \section{Instruments, Visibility, and Rotating Audit}\label{sn:11}

This section specifies the instrument-indexed claim semantics, whose finite self-certification limits invite comparison with work on incompleteness and reflection \citep{Goedel1931Incompleteness,Feferman1962Progressions}, the visibility discipline that the calculus uses to prevent overreading, and the rotating audits by which one instrument audits another without certifying itself. Subsections \hyperref[sn:11.1]{11.1} and \hyperref[sn:11.2]{11.2} fix the claim semantics and the classifier priority. Subsection \hyperref[sn:11.3]{11.3} states the visibility and suppression rules that the proofs use. Subsections \hyperref[sn:11.4]{11.4} and \hyperref[sn:11.5]{11.5} prove Theorems \hyperref[cl:thm:23]{23} and \hyperref[cl:thm:24]{24}. Subsection \hyperref[sn:11.6]{11.6} isolates the \(P_6\leftarrow P_6\) status separation. Subsection \hyperref[sn:11.7]{11.7} records the rotating-audit chain. Subsection \hyperref[sn:11.8]{11.8} proves Theorem \hyperref[cl:thm:25]{25}. Subsection \hyperref[sn:11.9]{11.9} discusses instrument transfer and marks its full classifier as future work.

\subsection{Instrument-Indexed Claim Semantics}\label{sn:11.1}

This subsection records the instrument record schema, the claim record schema, and the form of the acceptance judgment. The classifier rules and the no-overreading and same-level self-audit theorems that depend on them are stated in Subsections \hyperref[sn:11.4]{11.4} and \hyperref[sn:11.5]{11.5}.

\subsubsection{Instrument Record}\label{sn:11.1.1}

The instrument record \(I\) is the thirteen-field finite record of Subsubsection \hyperref[sn:3.4.1]{3.4.1}. Every component is finite. In a fixed finite instance, admissibility of \(I\) is decidable from its fields and the audit records judged under it: rule-language membership and inclusion of each audit's check rules in \(\mathsf{CheckRules}_I\) are finite syntactic checks; write \(\mathsf{AdmInstrument}(I)\) for this instance-relative predicate.

\subsubsection{Claim Record}\label{sn:11.1.2}

A claim record under instrument \(I\) is a finite typed tuple

\[
\begin{aligned}
\varphi(\mathcal T)\;=\;(\,&\mathrm{id}_{\varphi},\;\mathcal T,\;\mathsf{ClaimType},\;\mathsf{Content},\;\mathsf{LevelProfile},\\
&\mathsf{HostTag},\;\mathsf{EvidenceRefs},\;\Theta,\;A,\;V,\;\mathcal L,\;\mathcal N_{\varphi},\;\mathsf{partialEvidence}\,).
\end{aligned}
\]

The tuple fields record:

\begin{itemize}
\tightlist
\item
  \(\mathrm{id}_{\varphi}\) and \(\mathsf{Content}\): the claim's identifier and its content \(\varphi\), with finite use-set \(\mathsf{Use}(\varphi)\subseteq\mathsf{Cont}\);
\item
  \(\mathsf{ClaimType}\): a finite declared claim-type tag; membership in \(\mathsf{ClaimTypes}_I\) is tested by the claim classifier;
\item
  \(\mathcal T\): the target package (the surrounding context \(\Gamma\) is not a field of the record);
\item
  \(\mathsf{LevelProfile}\) and \(\mathsf{HostTag}\): the level profile, a pair \((\ell,\ell')\) of level tags with an optional level-bridge record of a type of Subsubsection \hyperref[sn:2.2.3]{2.2.3}, and the host tag of the claim; the level tags of the claim are \(\ell\) and \(\ell'\);
\item
  \(\mathsf{EvidenceRefs}\): the source field of the claim, references to the source records consulted, drawn from the source-of-truth schema of Subsubsection \hyperref[sn:4.9.1]{4.9.1}; the classifier's source check reads this field;
\item
  the threshold and audit records, drawn from the threshold and audit schemas of Subsection \hyperref[sn:5.2]{5.2} and Subsubsection \hyperref[sn:4.9.2]{4.9.2}; each threshold of \(\Theta\) names its evaluated datum as for a directed cell (Subsubsection \hyperref[sn:4.3.3]{4.3.3}), a recorded field or computed defect of a record named by the threshold record, which then lies in \(\mathsf{Use}(\varphi)\) (Subsubsection \hyperref[sn:3.2.1]{3.2.1});
\item
  the visibility record \(V\) and the visibility bridges \(\mathcal L\) used to bring suppressed content into admissible reach, drawn from \(\mathsf{BridgePolicy}_I\);
\item
  the nonclaim record \(\mathcal N_\varphi\) attached to the claim;
\item
  \(\mathsf{partialEvidence}\in\{\mathrm{true},\mathrm{false}\}\): the instance's declaration that the evidence is incomplete.
\end{itemize}

A candidate claim record is well formed when its fields are present, finite, and correctly typed, its audit slot holds an audit record whose \(\mathsf{ClaimRef}\) names the claim, or a typed, audited absence marker, whose audit defect is \(\{\texttt{missing\_audit}\}\) (Subsubsection \hyperref[sn:5.3.6]{5.3.6}), and every record identifier occurring in its statement lies in \(\mathsf{Refs}(\varphi)\) (Subsubsection \hyperref[sn:3.2.1]{3.2.1}), and, when \(\ell\neq\ell'\) or a record of \(\mathsf{Use}(\varphi)\) carries a level tag other than \(\ell,\ell'\), its \(\mathsf{LevelProfile}\) carries a level-bridge record of a type of Subsubsection \hyperref[sn:2.2.3]{2.2.3} (the stack claims of Subsection \hyperref[sn:11.8]{11.8} carry the reflection bridge); policy failures, including claim-type exclusion, are classifier outcomes. Well-formedness is a finite check on finite records.

\subsubsection{Acceptance Rule}\label{sn:11.1.3}

The general claim form is

\[
\Gamma\,;\;\mathcal T\,;\;I\;\vdash\;\varphi\;:\;\chi,
\]

with \(\chi\) in the claim status set \(\mathrm{ClaimStatus}\) of Subsubsection \hyperref[sn:4.8.2]{4.8.2}.

The classifier of Subsubsection \hyperref[sn:5.6.5]{5.6.5} consumes the claim record, the visibility partition \(\mathsf{Visible}_I/\mathsf{Suppressed}_I/\mathsf{Outside}_I/\mathsf{Unknown}_I\), the threshold and source records, the audit record, the bridge family \(\mathcal L\), and the nonclaim records, and assigns a value \(\chi\in\mathrm{ClaimStatus}\) to the claim. The verdict \(\chi=\texttt{accepted}\) records that the required in-scope, visibility, evidence, threshold, audit, circularity, and nonclaim checks have passed under \(I\); the other statuses record the specific failure or scope condition that prevented acceptance.

The shorthand \(I\vdash\varphi(\mathcal T)\) abbreviates \(\Gamma;\mathcal T;I\vdash\varphi:\texttt{accepted}\) and never abbreviates instrument-free truth, in the sense fixed in Subsubsection \hyperref[sn:3.4.2]{3.4.2}.

\subsection{Claim Classifier Priority}\label{sn:11.2}

\subsubsection{Priority Order}\label{sn:11.2.1}

The claim classifier of Subsubsection \hyperref[sn:5.6.5]{5.6.5} is a finite first-match procedure with the priority order displayed there. Theorems \hyperref[cl:thm:23]{23}--\hyperref[cl:thm:25]{25} use only these inputs: the claim record and its well-formedness (Subsubsection \hyperref[sn:11.1.2]{11.1.2}); \(\mathsf{Use}(\varphi)\) (Subsubsection \hyperref[sn:3.2.1]{3.2.1}); the reference graph (paragraph \emph{Reference graph}, Subsubsection \hyperref[sn:5.6.5]{5.6.5}); the audit record and its visibility, threshold and circularity checks (Subsubsection \hyperref[sn:4.9.2]{4.9.2}); the audit-defect codes (Subsubsection \hyperref[sn:5.3.6]{5.3.6}); the four visibility defects (Subsubsections \hyperref[sn:5.5.1]{5.5.1}--\hyperref[sn:5.5.3]{5.5.3}); and the claim table of Subsubsection \hyperref[sn:5.6.6]{5.6.6}.

A claim record matching multiple per-status conditions is assigned the highest-priority match. The order encodes the intended discipline: an outside-scope claim is recorded outside scope rather than rejected; a claim whose use-set contains suppressed unbridged content is blocked rather than rejected; audited absence and failed audits are recorded before threshold or refutation outcomes; circularity is recorded once the earlier scope, visibility, source, and audit checks do not already determine the status; and the accepted verdict is the lowest-priority match, reached only when every higher condition has been ruled out.

\subsubsection{Soundness Role}\label{sn:11.2.2}

The classifier prevents suppressed content from supporting overstrong claims. Its priority order first routes claims with outside-scope use-set entries to \(\texttt{outside\_scope}\); for in-scope visibility failures, a nonempty visibility defect \(\delta_{\mathrm{vis}}\) of Subsection \hyperref[sn:5.5]{5.5} routes the claim to \(\texttt{blocked}\) before \(\texttt{accepted}\) can be reached. Likewise, a same-level circular self-reference routes to \(\texttt{undefined\_circular}\) once the earlier scope, visibility, source, and audit checks have not already determined the status. The no-overreading and same-level self-audit failure theorems of Subsections \hyperref[sn:11.4]{11.4} and \hyperref[sn:11.5]{11.5} are the theorem-grade form of these routing rules.

\subsection{Visibility and Suppression}\label{sn:11.3}

The visibility map \(\mathsf{vis}_I\), the classes \(\mathsf{Visible}_I\) and \(\mathsf{Suppressed}_I\), visibility bridges and the bridge family \(\mathcal L\) are those of Subsection \hyperref[sn:3.5]{3.5}; the predicate \(\mathsf{AdmVisBridge}\) and the four visibility defects \(\delta_{\mathrm{overread}}\), \(\delta_{\mathrm{weakbridge}}\), \(\delta_{\mathrm{outside}}\) and \(\delta_{\mathrm{unknown}}\), with \(\delta_{\mathrm{vis}}\) their union and threshold \(\Theta_{\mathrm{vis}}:\delta_{\mathrm{vis}}=\varnothing\), are those of Subsection \hyperref[sn:5.5]{5.5}. The proofs below use \(\delta_{\mathrm{overread}}\) and \(\delta_{\mathrm{weakbridge}}\) (Subsubsections \hyperref[sn:5.5.1]{5.5.1} and \hyperref[sn:5.5.2]{5.5.2}).

\subsection{Theorem 23: No-Overreading Suppression}\label{sn:11.4}

\begin{claimtheorem}{Theorem 23 (NoOverreadingSuppression)}\phantomsection\label{cl:thm:23}

Let \(I\) be an admissible instrument, \(\varphi\) a finite claim, and \(\mathcal L_\varphi\) the visibility-bridge field of the claim record of \(\varphi\) (Subsubsection \hyperref[sn:11.1.2]{11.1.2}), with which the classifier computes \(\delta_{\mathrm{vis}}\). If

\[
\Gamma\,;\;\mathcal T\,;\;I\;\vdash\;\varphi\;:\;\texttt{accepted},
\]

then

\[
\delta_{\mathrm{vis}}(I,\varphi,\mathcal L)\;=\;\varnothing,\quad\text{in particular}\quad\delta_{\mathrm{overread}}(I,\varphi,\mathcal L)\;=\;\varnothing,
\]

for every finite \(\mathcal L\supseteq\mathcal L_\varphi\) of registered, audited candidates under \(I\).

Moreover, under the accepted-status hypothesis, every \(x\in\mathsf{Use}(\varphi)\cap\mathsf{Suppressed}_I\) is bridged by some \(L\in\mathcal L_\varphi\) for which \(\mathsf{AdmVisBridge}(L,I,\varphi)\) holds.

\end{claimtheorem}

\begin{proof}

By the classifier priority order of Subsubsection \hyperref[sn:11.2.1]{11.2.1}, the verdict \(\texttt{accepted}\) is reached only after every higher-priority condition has been ruled out. The \(\texttt{accepted}\) row is reached only if no earlier row holds: a nonempty \(\delta_{\mathrm{outside}}\) gives \(\mathsf{Use}(\varphi)\not\subseteq\mathsf{Scope}_I\), the \(\texttt{outside\_scope}\) row; a nonempty \(\delta_{\mathrm{overread}}\), \(\delta_{\mathrm{weakbridge}}\) or \(\delta_{\mathrm{unknown}}\) gives the \(\texttt{blocked}\) row. Hence \(\texttt{accepted}\) implies \(\delta_{\mathrm{vis}}(I,\varphi,\mathcal L_\varphi)=\varnothing\), and in particular \(\delta_{\mathrm{overread}}(I,\varphi,\mathcal L_\varphi)=\varnothing\). The overread defect is antitone in the bridge family: enlarging \(\mathcal L\) can only supply more bridges, so \(\delta_{\mathrm{overread}}(I,\varphi,\mathcal L)\subseteq\delta_{\mathrm{overread}}(I,\varphi,\mathcal L_\varphi)=\varnothing\) for every \(\mathcal L\supseteq\mathcal L_\varphi\).

Hence also \(\delta_{\mathrm{weakbridge}}(I,\varphi,\mathcal L)=\varnothing\) for \(\mathcal L\supseteq\mathcal L_\varphi\), since each suppressed \(x\) in the use-set is bridged by an admissible \(L'\in\mathcal L_\varphi\subseteq\mathcal L\), and \(\delta_{\mathrm{outside}}\), \(\delta_{\mathrm{unknown}}\) do not depend on \(\mathcal L\).

For the second statement, let \(x\in\mathsf{Use}(\varphi)\cap\mathsf{Suppressed}_I\). Since \(\delta_{\mathrm{overread}}(I,\varphi,\mathcal L_\varphi)=\varnothing\), some \(L\in\mathcal L_\varphi\) bridges \(x\); since \(\delta_{\mathrm{weakbridge}}(I,\varphi,\mathcal L_\varphi)=\varnothing\), the pair \((x,L)\) is not in the weak-bridge defect, so by its definition some \(L'\in\mathcal L_\varphi\) with \(\mathsf{AdmVisBridge}(L',I,\varphi)\) bridges \(x\).

\end{proof}

The second statement applies to accepted claims that use suppressed content. Take \(I'\), with a fresh identifier, in the instance of Theorem \hyperref[cl:thm:22]{22} extended by the records below, by extending \(I_{\mathrm{prom}}\) of Subsubsection \hyperref[sn:10.6.1]{10.6.1} to admit the claim type \(\texttt{injectivity}\), declaring \(\mathsf{ReqStrength}_{I'}(\texttt{injectivity})=\texttt{strong}\) and \(\mathsf{Strength}(L_t)=\texttt{strong}\) in the ordered strengths \(\{\texttt{weak}<\texttt{strong}\}\), a \(\texttt{committed\_state}\) evidence source, and a rule checking the two recorded values of \(\pi_1\). Put the recorded table \(t\) of \(\pi_1\), a recomputed copy \(t'\), a visibility bridge \(L_t\), a claim \(\varphi\) and their evidence, threshold and audit records in \(\mathsf{Scope}_{I'}\). Make all these records visible except \(t\), and set \(\mathsf{Suppressed}_{I'}=\{t\}\). Let an audited \(L_t\) certify \(t'=t\) at the required strength without using \(\ulcorner\varphi,I'\urcorner\). Give \(\varphi\) host \(\mathsf H_{\mathrm{fin}}\), level \(\mathsf{Ver}\), claim type \(\texttt{injectivity}\), the statement that \(\pi_1\), recorded by its table \(t\), is injective, so that \(t\in\mathsf{Use}(\varphi)\cap\mathsf{Suppressed}_{I'}\), passing checks, bridge field \(\{L_t\}\) and \(\mathsf{partialEvidence}=\mathrm{false}\). Since \(t'\) records \(\pi_1(a)=u\ne v=\pi_1(b)\), the claim reaches \(\texttt{accepted}\) with \(t\) bridged by \(L_t\).

\subsubsection{Corollaries}\label{sn:11.4.1}

The no-overreading constraint specializes to three corollaries on the judgment families of Section \hyperref[sn:4]{4}.

\textbf{Directed-cell corollary.} A directed-cell judgment with status \(\texttt{action}\) has empty visibility defect, computed with the bridges listed in its visibility record: by the same argument, every suppressed field of the cell record (every element of \(\mathsf{Use}(C)\cap\mathsf{Suppressed}_I\), Subsubsection \hyperref[sn:5.5.2]{5.5.2}) is bridged by a bridge of that list that is admissible for the cell.

\textbf{Pair-observable corollary.} A pair-observable judgment with status \(\texttt{real}\) has empty overread defect on the source-of-truth record and the branch-compatibility record. By the classifier of Subsubsection \hyperref[sn:5.6.3]{5.6.3}, the source is either a real-source class admitted by \(\mathsf{SourcePolicy}_I\) or a fallback source with an audited source-upgrade bridge admitted by \(\mathsf{SourcePolicy}_I\) and \(\mathsf{BridgePolicy}_I\); any suppressed dependency must be bridged.

\textbf{Promotion corollary.} For a promotion bridge with status \(\texttt{strict}\), the visibility gate \(G_{\mathrm{vis}}\) and the no-smuggling gate \(G_{\mathrm{nosmuggle}}\) pass, and by Theorem \hyperref[cl:thm:19]{19} the visibility defect of the bridge's use-set is empty: no suppressed content is used without an admissible bridge.

\subsection{Theorem 24: Same-Level Self-Audit Failure}\label{sn:11.5}

\subsubsection{Target}\label{sn:11.5.1}

Let \(I_0\) be an admissible instrument under an admissible package \(\mathcal T_0\) on any host of Subsubsection \hyperref[sn:2.1.2]{2.1.2}. Consider the candidate same-level self-soundness claim \(\mathrm{Sound}(I_0)\), presented as the attempted judgment

\[
\Gamma\,;\;\mathcal T_0\,;\;I_0\;\vdash\;\mathrm{Sound}(I_0):\chi.
\]

If \(\chi=\texttt{accepted}\), this would be written in shorthand as \(I_0\vdash\mathrm{Sound}(I_0)\). The target claim asserts that \(I_0\) certifies its own soundness without a level shift to a higher instrument.

\emph{Claim log.} Write \(\mathsf{ClaimLog}(I_0)\) for the finite set of claim records, of claim types in \(\mathsf{ClaimTypes}(I_0)\), for which the instance records a verdict of \(I_0\) (the instance declares this set), each paired with the verdict that \(I_0\) recorded for it. Each entry also records the finite judgment context \(\Gamma_\psi\), stored as a single context record whose reference fields name the members of the context, and the admissible package or host \(\mathcal T_\psi\) under which \(\psi\) was judged; we keep writing \((\psi,\chi)\) for the entry, and ``the status the claim classifier of \(I_0\) assigns to \(\psi\)'' means the status it assigns under \(\Gamma_\psi;\mathcal T_\psi;I_0\).

\emph{Complete soundness.} Complete soundness of \(I_0\) means here that every verdict in \(\mathsf{ClaimLog}(I_0)\) equals the status the claim classifier assigns under the corresponding recorded context; when \(\mathrm{Sound}(I_0)\) is logged this includes \(I_0\)'s recorded verdict on it; it is not instrument-free truth. The claim log, including these context records, is a finite record of the instance, kept with \(I_0\) under its identifier \(\mathrm{id}_{I_0}\); it is not a component of the instrument tuple of Subsubsection \hyperref[sn:3.4.1]{3.4.1}, so equality of instruments does not depend on it, and \(\mathsf{ClaimLog}(I)\) is defined in the same way for every instrument \(I\) of the instance. Unlike the lifted claims of Theorem \hyperref[cl:thm:25]{25} (Subsection \hyperref[sn:11.8]{11.8}), which concern only the declared logs of lower instruments, \(\mathrm{Sound}(I_0)\) names its own verdict through \(\ulcorner\mathrm{Sound}(I_0),I_0\urcorner\), which lies in its use-set whether or not that verdict is logged; unlogged verdicts are not records and are not audited by these claims.

\emph{The record \(\mathrm{Sound}(I_0)\).} As a record, \(\mathrm{Sound}(I_0)\) is the claim record with claim-type tag \(\texttt{soundness}\) whose statement is that the verdict of \(I_0\) on each claim in \(\mathsf{ClaimLog}(I_0)\), and, if logged, on \(\mathrm{Sound}(I_0)\) itself, is the status that the claim classifier of Subsubsection \hyperref[sn:5.6.5]{5.6.5} assigns to that claim; the designated verdict-reference subset \(\mathsf{Use}_{\mathrm{verdict}}(\mathrm{Sound}(I_0))\) of its use-set is the finite set

\[
\mathsf{Use}_{\mathrm{verdict}}(\mathrm{Sound}(I_0))\;=\;\{\ulcorner\psi,I_0\urcorner:\exists\chi\,(\psi,\chi)\in\mathsf{ClaimLog}(I_0)\}\;\cup\;\{\ulcorner\mathrm{Sound}(I_0),I_0\urcorner\}
\]

of verdict references. Its declared reference set \(\mathsf{Refs}(\mathrm{Sound}(I_0))\) consists of \(\mathsf{Use}_{\mathrm{verdict}}(\mathrm{Sound}(I_0))\), the records \(I_0\), \(\mathrm{Sound}(I_0)\) and \(\mathsf{ClaimLog}(I_0)\), and each logged \(\psi\) with \(\Gamma_\psi\) and \(\mathcal T_\psi\) (the statement names each of them, so Subsubsection \hyperref[sn:3.2.1]{3.2.1} requires this), and by Subsubsection \hyperref[sn:3.2.1]{3.2.1} the use-set is \(\mathsf{Refs}(\mathrm{Sound}(I_0))\) together with the source records of \(\mathsf{EvidenceRefs}\) and the records referenced by the claim's threshold and audit fields. In particular the use-set contains \(\ulcorner\mathrm{Sound}(I_0),I_0\urcorner\). For a newly constructed candidate, adjoin the claim and its own verdict-reference record together, using fresh identifiers and mutually resolved references, along with its required audit, source and threshold records. For an existing well-formed candidate of this specified form, the verdict-reference record is already present. Evaluate the candidate in the resulting instance. The theorem below follows from membership of its own verdict reference in its use-set and the circularity rule of the claim classifier.

\subsubsection{Outside-Scope Case}\label{sn:11.5.2}

Suppose \(\texttt{soundness}\notin\mathsf{ClaimTypes}(I_0)\), that is, \(I_0\) is not configured to evaluate claims about its own soundness. By the classifier of Subsubsection \hyperref[sn:5.6.5]{5.6.5}, the \(\texttt{outside\_scope}\) condition holds because the claim type is not in \(\mathsf{ClaimTypes}(I_0)\) (the other condition of that row, a use-set not contained in \(\mathsf{Scope}_{I_0}\), would give the same status), and the priority routing of Subsection \hyperref[sn:11.2]{11.2} returns

\[
\Gamma\,;\;\mathcal T_0\,;\;I_0\;\vdash\;\mathrm{Sound}(I_0)\;:\;\texttt{outside\_scope}.
\]

This is the case for instruments configured to evaluate object-level claims only.

\subsubsection{Circular Case}\label{sn:11.5.3}

Suppose the claim's host tag is in \(\mathsf{Hosts}_{I_0}\), its level tags are in \(\mathsf{Levels}_{I_0}\), \(\texttt{soundness}\in\mathsf{ClaimTypes}(I_0)\), and the use-set of \(\mathrm{Sound}(I_0)\) is contained in \(\mathsf{Scope}_{I_0}\). By Subsubsection \hyperref[sn:11.5.1]{11.5.1} that use-set contains \(\ulcorner\mathrm{Sound}(I_0),I_0\urcorner\). The classifier records this circular self-reference at the same level and, when no higher-priority \(\texttt{blocked}\), \(\texttt{absent\_with\_record}\) or \(\texttt{failed\_audit}\) condition is present, returns

\[
\Gamma\,;\;\mathcal T_0\,;\;I_0\;\vdash\;\mathrm{Sound}(I_0)\;:\;\texttt{undefined\_circular}.
\]

This verdict follows from the same-level verdict reference explicitly included in \(\mathsf{Use}(\mathrm{Sound}(I_0))\).

\subsubsection{Statement and Proof}\label{sn:11.5.4}

\emph{Setting.} \(I_0\) is an admissible instrument under an admissible package \(\mathcal T_0\) on any host of Subsubsection \hyperref[sn:2.1.2]{2.1.2}, and \(\mathrm{Sound}(I_0)\) is the complete same-level self-soundness claim of Subsubsection \hyperref[sn:11.5.1]{11.5.1}, judged by the claim classifier of Subsubsection \hyperref[sn:5.6.5]{5.6.5}.

\begin{claimtheorem}{Theorem 24 (SameLevelSelfAuditFailure)}\phantomsection\label{cl:thm:24}

Let \(\mathrm{Sound}(I_0)\) be the claim record of Subsubsection \hyperref[sn:11.5.1]{11.5.1}, whose use-set contains its own verdict reference \(\ulcorner\mathrm{Sound}(I_0),I_0\urcorner\). The claim classifier of Subsubsection \hyperref[sn:5.6.5]{5.6.5} does not assign this same-level self-soundness claim the verdict \(\texttt{accepted}\) under \(I_0\). The content of \(\mathrm{Sound}(I_0)\) can nevertheless be true: if every verdict in \(\mathsf{ClaimLog}(I_0)\) equals the classifier's status and \(I_0\) records for \(\mathrm{Sound}(I_0)\) the status the classifier assigns it, its statement holds, and it is still not \(\texttt{accepted}\); so for this claim truth does not imply acceptance, whereas for claims whose check rules meet condition 6 of Subsubsection \hyperref[sn:2.4.1]{2.4.1} and whose content is a decidable property of the recomputed witnesses, acceptance implies truth (Subsubsection \hyperref[sn:5.6.5]{5.6.5}). If \(\mathrm{Sound}(I_0)\) is well formed (Subsubsection \hyperref[sn:11.1.2]{11.1.2}), the classifier returns \(\texttt{outside\_scope}\) if the claim's host tag is not in \(\mathsf{Hosts}_{I_0}\), a level tag is not in \(\mathsf{Levels}_{I_0}\), its type is not in \(I_0\)'s claim types or its use-set is not contained in \(I_0\)'s scope, and otherwise returns \(\texttt{undefined\_circular}\) or a status of higher priority; if \(\mathrm{Sound}(I_0)\) is not well formed, the classifier assigns it no status. In either case it is not accepted. More generally, under \(I_0\) the classifier accepts no well-formed claim of type \(\texttt{soundness}\) whose use-set contains a verdict reference \(\ulcorner\psi,I'\urcorner\) with \(\mathsf{Level}(I')\ge\mathsf{Level}(I_0)\) (by that stratification condition; no reference cycle is required). Whatever its claim-type tag, no well-formed claim whose statement includes \(I_0\)'s verdict on that claim is accepted under \(I_0\), since \(\ulcorner\varphi,I_0\urcorner\in\mathsf{Use}(\varphi)\) and the first condition of the \(\texttt{undefined\_circular}\) row applies.

\end{claimtheorem}

The result follows from the membership of the verdict reference \(\ulcorner\mathrm{Sound}(I_0),I_0\urcorner\) in the use-set of the claim and the stratification condition of the \(\texttt{undefined\_circular}\) row of Subsubsection \hyperref[sn:5.6.5]{5.6.5}; it uses no other property of \(I_0\).

\begin{proof}

For the complete same-level self-audit target, the classifier has the two cases of Subsubsections \hyperref[sn:11.5.2]{11.5.2} and \hyperref[sn:11.5.3]{11.5.3}. If the target claim's host tag is not in \(\mathsf{Hosts}_{I_0}\), a level tag is not in \(\mathsf{Levels}_{I_0}\), its type is not in \(\mathsf{ClaimTypes}(I_0)\), or its use-set is not contained in \(\mathsf{Scope}_{I_0}\), the \(\texttt{outside\_scope}\) rule applies, and \(\texttt{outside\_scope}\) has the highest priority. Otherwise the claim is in scope. By Subsubsection \hyperref[sn:11.5.1]{11.5.1}, the use-set of \(\mathrm{Sound}(I_0)\) then contains \(\ulcorner\mathrm{Sound}(I_0),I_0\urcorner\); since the claim type is \(\texttt{soundness}\) and \(\mathsf{Level}(I_0)\ge\mathsf{Level}(I_0)\), the third condition of the \(\texttt{undefined\_circular}\) rule holds, a condition that tests membership rather than reachability. The classifier then returns \(\texttt{undefined\_circular}\), or \(\texttt{blocked}\), \(\texttt{absent\_with\_record}\) or \(\texttt{failed\_audit}\) if one of those higher-priority rules also applies. Each of these statuses precedes \(\texttt{accepted}\) in the priority order of Subsubsection \hyperref[sn:5.6.5]{5.6.5}, so \(\texttt{accepted}\) is not reached. Hence the classifier does not accept the same-level self-soundness claim under \(I_0\). The same argument applies verbatim to the general statement, since the third condition tests membership: a well-formed soundness claim whose use-set contains such a verdict reference receives \(\texttt{outside\_scope}\), \(\texttt{undefined\_circular}\) or a status of higher priority. A level shift does not change this verdict: it judges the claim under a higher instrument instead of \(I_0\).

The theorem is a statement about the routing of the claim classifier of this calculus: its priority rules send a complete same-level self-soundness claim to a status that precedes \(\texttt{accepted}\) in the classifier's priority order. It is not a general impossibility result about self-certification for instruments outside this classifier.

NC-16 of Subsection \hyperref[sn:15.2]{15.2} is the corresponding nonclaim. A positive soundness claim about a lower instrument requires a level shift to a higher instrument \(I_1\) (Subsubsection \hyperref[sn:5.6.5]{5.6.5}), and an admissible audit bridge through which \(I_1\) audits \(I_0\) records it; Subsection \hyperref[sn:11.7]{11.7} records the rotating-audit chain that supplies such bridges, and Subsection \hyperref[sn:11.8]{11.8} records that a rotating audit can certify the lower stack but does not certify itself, so the rotating audit is non-final (NC-17).

\end{proof}

A same-level claim of a type other than \(\texttt{soundness}\) from whose use-set its own verdict reference is not reachable is not covered by Theorem \hyperref[cl:thm:24]{24}; one from whose use-set it is reachable receives \(\texttt{undefined\_circular}\) or a status of higher priority by the same rule of the claim classifier.

\subsection{\texorpdfstring{\(P_6\leftarrow P_6\) Status Separation}{P\_6\textbackslash leftarrow P\_6 Status Separation}}\label{sn:11.6}

This subsection records a schema-level distinction between three judgments that all carry the same primitive pair \(P_6\leftarrow P_6\): the object-level audit update, the fixed audit refresh, and the complete same-level self-audit. They differ in profile and update. The first two can differ in directed-cell status (\(\texttt{action}\) for an object-level audit whose records pass, \(\texttt{trivial}\) for a refresh, since \(\mathsf{noopU}\) holds); the third is separated at the level of its claim record, which Theorem \hyperref[cl:thm:24]{24} keeps away from \(\texttt{accepted}\). Theorem \hyperref[cl:thm:1]{1} of Section \hyperref[sn:6]{6} covers the formal non-factorization of cell status through the primitive pair.

\subsubsection{Object-Level Audit Update}\label{sn:11.6.1}

A directed cell \(P_6\leftarrow P_6\) in which the actor update \(U_6\) records an audit on a specific object-level record --- a directed cell, a pair observable, a promotion bridge, or a claim record other than \(I\)'s own soundness --- is admissible. Its profile \(\lambda_{\mathrm{obj}}\) has the audit-only instrument mode of Subsubsection \hyperref[sn:3.6.1]{3.6.1}, at object level. The cell may receive directed-cell status \(\texttt{action}\) in the classifier of Subsubsection \hyperref[sn:5.6.2]{5.6.2} when its witness, threshold, audit, and visibility records pass the standard cell admissibility checks.

\subsubsection{Fixed Audit Refresh}\label{sn:11.6.2}

A directed cell \(P_6\leftarrow P_6\) in which the actor update is a fixed-instrument audit refresh --- re-running an existing audit policy on an already present, unchanged audit record under a profile \(\lambda_{\mathrm{refresh}}\) with the audit-only instrument mode --- records the identity effect map on the audit record, since the refresh leaves that record unchanged, so \(\mathsf{noopU}\) is computed true, and is classified \(\texttt{trivial}\) in the priority order of Subsubsection \hyperref[sn:5.6.2]{5.6.2} unless a row above it applies. If the audit step produces new admissible audit content, it is no longer this fixed-refresh case and must be classified under its own object-level audit-update profile.

\subsubsection{Complete Same-Level Self-Audit}\label{sn:11.6.3}

A directed cell \(P_6\leftarrow P_6\) in which the actor update asserts a complete same-level self-soundness verdict on the surrounding instrument \(I_0\) is the case covered by Theorem \hyperref[cl:thm:24]{24} of Subsection \hyperref[sn:11.5]{11.5}. By that theorem, the claim record \(\mathrm{Sound}(I_0)\) that such an update asserts receives \(\texttt{outside\_scope}\), \(\texttt{undefined\_circular}\) or a claim status of higher priority, and never \(\texttt{accepted}\). The status of the directed cell itself is assigned separately, by the cell classifier of Subsubsection \hyperref[sn:5.6.2]{5.6.2} from the cell record.

The three cases together separate object-level audit, fixed audit refresh, and complete same-level self-certification within the same primitive pair \((P_6,P_6)\). They are an instance of the typed/profile-relative discipline of the calculus: when two of these cells on the same pair receive different statuses, Theorem \hyperref[cl:thm:1]{1} shows that cell status does not factor through the pair.

\subsection{Rotating Audit}\label{sn:11.7}

This subsection records the rotating-audit discipline that replaces forbidden complete same-level self-audit with a finite, level-shifted, bounded audit of one instrument by a strictly higher instrument. In a rotating-audit chain the auditing role passes upward: each instrument is audited by the next, strictly higher one, and the chain never returns to a lower level. It records the rotating-audit chain, the rotating-audit bridge, the admissibility predicate for that bridge, and the rotating-audit defect schema. Theorem \hyperref[cl:thm:25]{25} of Subsection \hyperref[sn:11.8]{11.8} then states the finite rotating-audit theorem and its non-finality.

\subsubsection{Rotating-Audit Chain}\label{sn:11.7.1}

A rotating-audit chain is a finite ordered family of admissible instruments

\[
\mathbb I_{\le N}\;=\;(I_0,I_1,\ldots,I_N),
\]

with strictly increasing instrument levels \(\mathsf{Level}(I_0)<\mathsf{Level}(I_1)<\cdots<\mathsf{Level}(I_N)\) in the order \(\mathsf{Beh}<\mathsf{Ver}<\mathsf{Str}\) of Subsubsection \hyperref[sn:2.2.3]{2.2.3}, together with a finite ordered family of rotating-audit bridges

\[
B^{\mathrm{audit}}_{n\to n+1}\quad\text{for each } 0\le n<N,
\]

each pointing to a strictly higher level. Each \(I_n\) is an instrument in the sense of Subsubsection \hyperref[sn:11.1.1]{11.1.1}, with finite logs, declared visibility policy, declared source policy, declared audit policy, declared nonclaim register, and declared scope. The chain is finite by construction: \(N\) is a fixed natural number and there is no implicit limit instrument at the top. The chain has no final self-certifying instrument; Subsection \hyperref[sn:11.8]{11.8} records this non-finality formally.

With the three levels of \(\mathsf{Lev}\), such a chain has at most three instruments, and Theorem \hyperref[cl:thm:25]{25} applies exactly to chains whose highest level is \(\mathsf{Beh}\) or \(\mathsf{Ver}\); such chains have \(N\le1\) (Subsubsection \hyperref[sn:11.8.1]{11.8.1}).

\subsubsection{Rotating-Audit Bridge}\label{sn:11.7.2}

For each adjacent pair \((I_n,I_{n+1})\) with \(\mathsf{Level}(I_{n+1})>\mathsf{Level}(I_n)\), the rotating-audit bridge

\[
B^{\mathrm{audit}}_{n\to n+1}\;=\;(I_n,\;I_{n+1},\;\Phi_n,\;L,\;V,\;\Theta,\;A,\;\delta,\;\mathcal N)
\]

with expanded finite record

\[
\begin{aligned}
B^{\mathrm{audit}}_{n\to n+1}\;=\;(&\mathrm{id},\;I_n,\;I_{n+1},\;\Phi_n,\;\mathsf{TargetRecords}_n,\;L_{n\to n+1},\;V_{n\to n+1},\\
&\Theta_{n\to n+1},\;A_{n\to n+1},\;\delta_{n\to n+1},\;\mathcal N_{n\to n+1},\;\lambda_{\mathrm{rot}}).
\end{aligned}
\]

It is a directed bridge from \(I_n\) into \(I_{n+1}\) whose bridge type is the level bridge \(\texttt{reflection}\), used for audit. Its components are:

\begin{itemize}
\tightlist
\item
  \(\Phi_n\), the rotating-audit target scope: a finite, declared subset of \(I_n\)'s records (the instrument record \(I_n\), the named programs of \(\mathsf{CheckRules}_{I_n}\), the log \(\mathsf{ClaimLog}(I_n)\) with the claim records, contexts, packages and verdict references of its entries, and the source, audit and nonclaim records these name) on which \(I_{n+1}\) undertakes to compute a finite audit defect, in the sense of Subsubsection \hyperref[sn:11.7.4]{11.7.4};
\item
  \(L_{n\to n+1}\), the language map that translates \(I_n\)-records into \(I_{n+1}\)'s claim language;
\item
  \(V_{n\to n+1}\), the visibility map that exposes the audited records to \(I_{n+1}\) either directly or through visibility bridges of Subsubsection \hyperref[sn:3.5.3]{3.5.3};
\item
  \(\Theta_{n\to n+1}\), the audit thresholds (witness, audit, visibility) that \(I_{n+1}\) applies when classifying the bridge;
\item
  \(A_{n\to n+1}\), the audit record of the bridge itself, recording that the bridge is auditable in \(I_{n+1}\);
\item
  \(\delta_{n\to n+1}\), the rotating-audit defect of Subsubsection \hyperref[sn:11.7.4]{11.7.4};
\item
  \(\mathcal N_{n\to n+1}\), the bridge-level nonclaim register that records, at minimum, \(\mathrm{NC}\text{-}16\), \(\mathrm{NC}\text{-}17\), and \(\mathrm{NC}\text{-}18\);
\item
  \(\lambda_{\mathrm{rot}}=(\mathsf{Level}(I_n),\mathsf{Level}(I_{n+1}),\texttt{fine},\texttt{audit},\texttt{reflection},\texttt{rotating-audit},\texttt{global})\), the bridge profile recording rotating-audit mode; the reflection level bridge its distinct tags require is the bridge itself.
\end{itemize}

\emph{Stack form.} For a finite chain \(I_0,\ldots,I_N\) and an instrument \(I_{N+1}\) with \(\mathsf{Level}(I_{N+1})>\max_{n\le N}\mathsf{Level}(I_n)\), a stack rotating-audit bridge \(B^{\mathrm{audit}}_{\le N\to N+1}\) has the same components, with the chain as its source, a finite declared target scope \(\Phi_{\le N}\) among the records of the chain and its bridges, and a language map on all of them. Its admissibility (Subsubsection \hyperref[sn:11.7.3]{11.7.3}) reads clause (1) as admissibility of every \(I_n\), \(n\le N\), and of \(I_{N+1}\), and clause (2) as \(\mathsf{Level}(I_{N+1})>\max_{n\le N}\mathsf{Level}(I_n)\); the other clauses are unchanged. The stack profile has source tag \(\max_{n\le N}\mathsf{Level}(I_n)\) and target tag \(\mathsf{Level}(I_{N+1})\); it retains the reflection, mode and compatibility requirements of clause (1).

The rotating-audit bridge is the instrument-level vehicle by which \(I_{n+1}\) reflects \(I_n\)'s declared records into its own claim language under a declared, finite scope \(\Phi_n\). It is a \(P_6\leftarrow P_6\) object at the bridge layer, distinct from the directed-cell \(P_6\leftarrow P_6\) judgments separated in Subsection \hyperref[sn:11.6]{11.6}: a rotating-audit bridge audits an instrument across a level shift; the cell-level judgments of Subsection \hyperref[sn:11.6]{11.6} audit object-level records under a single instrument.

\subsubsection{Rotating-Audit Admissibility}\label{sn:11.7.3}

The admissibility predicate

\[
\mathsf{AdmRotAuditBridge}(B^{\mathrm{audit}}_{n\to n+1})
\]

holds iff every clause below holds:

\begin{enumerate}
\def\labelenumi{\arabic{enumi}.}
\tightlist
\item
  \(I_n\) and \(I_{n+1}\) are admissible instruments in the sense of Subsubsection \hyperref[sn:11.1.1]{11.1.1}; \(\mathsf{BridgePolicy}_{I_{n+1}}\) admits reflection, \(\mathsf{ModePolicy}_{I_{n+1}}\) maps the rotating-audit-bridge family to rotating-audit, and \((\texttt{fine},\texttt{audit},\texttt{global},\mathsf H_{\mathrm{instr}})\in\mathsf{Compat}_{I_{n+1}}\);
\item
  \(\mathsf{Level}(I_{n+1})>\mathsf{Level}(I_n)\);
\item
  \(\Phi_n\) is finite and declared as part of \(B^{\mathrm{audit}}_{n\to n+1}\);
\item
  every record in \(\Phi_n\) lies in \(\mathsf{Scope}_{I_{n+1}}\) and is visible to \(I_{n+1}\) directly or through a visibility bridge of Subsubsection \hyperref[sn:3.5.3]{3.5.3};
\item
  \(L_{n\to n+1}\) is a finite table sending each record of \(\Phi_n\) to a claim record, of a type in \(\mathsf{ClaimTypes}_{I_{n+1}}\), that is well formed in the sense of Subsubsection \hyperref[sn:11.1.2]{11.1.2};
\item
  \(V_{n\to n+1}\) exposes the records required by the audit thresholds, and \(\delta_{\mathrm{vis}}(I_{n+1},B,\mathcal L_B)=\varnothing\) (Subsection \hyperref[sn:5.5]{5.5}), with \(\mathsf{AdmVisBridge}(L,I_{n+1},B)\) read as in Subsubsection \hyperref[sn:5.5.2]{5.5.2} for \(X=B\) and \(\mathsf{RequiredStrength}(B)\) the value \(\mathsf{BridgePolicy}_{I_{n+1}}\) assigns to \((\texttt{reflection},\texttt{audit})\), where \(\mathsf{Use}(B)\) is \(\Phi_n\) together with the content references (one step, Subsubsection \hyperref[sn:5.5.2]{5.5.2}) in \(L_{n\to n+1},\Theta_{n\to n+1},A_{n\to n+1}\), and \(\mathcal L_B\) is the set of bridges listed in \(V_{n\to n+1}\);
\item
  every threshold of \(\Theta_{n\to n+1}\) is a record of the schema of Subsubsection \hyperref[sn:5.2.2]{5.2.2} and belongs to \(\mathsf{ThresholdPolicy}_{I_{n+1}}\);
\item
  \(A_{n\to n+1}\) is an audit record in the schema of Subsubsection \hyperref[sn:4.9.2]{4.9.2} with \(\mathsf{ClaimRef}\) the bridge and \(\mathsf{CheckRules}\) drawn from \(\mathsf{CheckRules}_{I_{n+1}}\);
\item
  \(\delta_{n\to n+1}\) is typed by the rotating-audit defect schema of Subsubsection \hyperref[sn:11.7.4]{11.7.4};
\item
  no verdict reference \(\ulcorner\psi,I_{n+1}\urcorner\), for any \(\psi\), is reachable from a record of \(\Phi_n\) in the reference graph of Subsubsection \hyperref[sn:5.6.5]{5.6.5}, so no record of \(\Phi_n\) depends on a verdict of \(I_{n+1}\);
\item
  \(\mathcal N_{n\to n+1}\) is present and contains the rotating-audit nonclaims of the calculus, in particular \(\mathrm{NC}\text{-}16\), \(\mathrm{NC}\text{-}17\), and \(\mathrm{NC}\text{-}18\).
\end{enumerate}

The clauses are finite checks on the records of \(B^{\mathrm{audit}}_{n\to n+1}\), decidable under hypotheses (a)--(b) of Theorem \hyperref[cl:thm:3]{3}; admissibility is a predicate on the bridge object, not an external soundness condition.

\subsubsection{Rotating-Audit Defect}\label{sn:11.7.4}

The rotating-audit defect schema for a rotating-audit bridge \(B^{\mathrm{audit}}_{n\to n+1}\) is

\[
\mathsf{Def}^{\mathrm{rot}}_{n\to n+1}\;=\;(\delta_{\mathrm{rot}},\;\Omega_{\mathrm{rot}},\;\Theta_{\mathrm{rot}},\;\mathrm{polarity}_{\mathrm{rot}},\;W_{\mathrm{rot}},\;A_{\mathrm{rot}},\;V_{\mathrm{rot}},\;\sigma_{\mathrm{rot}}),
\]

following the defect-record schema of Subsection \hyperref[sn:5.2]{5.2}. The defect value \(\delta_{\mathrm{rot}}\) ranges over the rotating-audit defect codes

\[
\begin{aligned}
\Omega_{\mathrm{rot}}\;=\;\{\,&\texttt{audit\_passes},\;\texttt{missing\_target\_scope},\;\texttt{missing\_instrument\_record},\\
&\texttt{visibility\_failure},\;\texttt{translation\_failure},\;\texttt{source\_failure},\\
&\texttt{threshold\_failure},\;\texttt{nonclaim\_failure},\;\texttt{same\_level\_cycle},\;\texttt{higher\_audit\_failure}\,\},
\end{aligned}
\]

with intended readings:

\begin{itemize}
\tightlist
\item
  \(\texttt{audit\_passes}\): the bridge audit succeeded under its declared thresholds;
\item
  \(\texttt{missing\_target\_scope}\): \(\Phi_n\) is not finite or not declared;
\item
  \(\texttt{missing\_instrument\_record}\): a record required by \(\Phi_n\) is not present in \(I_n\) or not visible in \(I_{n+1}\);
\item
  \(\texttt{visibility\_failure}\): a visibility-policy clause of \(V_{n\to n+1}\) fails;
\item
  \(\texttt{translation\_failure}\): \(L_{n\to n+1}\) fails to type a record in \(\Phi_n\);
\item
  \(\texttt{source\_failure}\): \(I_{n+1}\) fails to check a source-of-truth clause on a record of \(\Phi_n\) as its bridge policy requires (a clause the check finds failing is report data);
\item
  \(\texttt{threshold\_failure}\): an audit threshold in \(\Theta_{n\to n+1}\) is violated;
\item
  \(\texttt{nonclaim\_failure}\): a clause in \(\mathcal N_{n\to n+1}\) is violated;
\item
  \(\texttt{same\_level\_cycle}\): an unstratified same-level audit cycle is detected in the bridge dependency graph;
\item
  \(\texttt{higher\_audit\_failure}\): the bridge-level audit record \(A_{n\to n+1}\) fails.
\end{itemize}

A defect found in a lower-stack record is report data. It is a bridge-level failure code only if the extension instrument fails to inspect, translate, source, or audit that record as required by the bridge policy. A correctly detected lower-stack threshold violation is likewise report data, not the bridge code \(\texttt{threshold\_failure}\).

The acceptance threshold is

\[
\Theta_{\mathrm{rot}}:\;\delta_{\mathrm{rot}}=\texttt{audit\_passes}.
\]

The value \(\texttt{audit\_passes}\) is the threshold value later consumed by the rotating-audit acceptance rule of Subsection \hyperref[sn:11.8]{11.8}.

\subsection{Theorem 25: Finite Rotating-Audit Theorem}\label{sn:11.8}

This subsection records the finite rotating-audit theorem. It states that, given a finite stack of admissible instruments and a level of \(\mathsf{Lev}\) above the highest level of the stack, there is an extension instrument at that level and an audit bridge that can finitely audit the lower stack and accept its compliance iff the stack-audit defect is empty, while the higher instrument does not thereby certify itself. The statement and proof follow, and Subsubsection \hyperref[sn:11.8.1]{11.8.1} records the non-finality of the construction. The theorem supports \(\mathrm{NC}\text{-}17\) of Subsection \hyperref[sn:15.2]{15.2}.

In words: the lower stack is audited by recomputation. Each \(I_n\) keeps a log of the claims it judged and the verdicts it recorded; the new instrument \(I^{\mathrm{stack}}_{N+1}\), at a strictly higher declared level, reruns \(I_n\)'s classifier on each logged claim and compares. Clause 1 says the report of all defect values is accepted whatever they are; clause 2 says the stack is compliant exactly when no logged verdict differs from the recomputed one (\(\delta_{\mathrm{status}}=\varnothing\)); clause 4 says the same instrument by instrument; clause 3 says none of this lets \(I^{\mathrm{stack}}_{N+1}\) accept its own soundness claim. Since \(\mathsf{Lev}\) has three levels, the chain has one or two instruments. Clauses 1, 2 (first equivalence) and 4 follow from the choice of the check rules of \(I^{\mathrm{stack}}_{N+1}\) once the other classifier rows are ruled out. The proof therefore consists of (i) admissibility of the extension instrument, the stack bridge and the extended instance; (ii) exclusion of the scope, visibility, absence, audit, circularity and threshold rows; (iii) the second equivalence of clause 2, which holds for every admissible chain; and (iv) clause 3, by Theorem \hyperref[cl:thm:24]{24}.

\emph{Setting.} Let

\[
\mathbb I_{\le N}\;=\;(I_0,\,I_1,\,\ldots,\,I_N)
\]

be a finite admissible rotating-audit chain in the sense of Subsubsection \hyperref[sn:11.7.1]{11.7.1}, with strictly increasing levels and a finite family of admissible rotating-audit bridges \(\{B^{\mathrm{audit}}_{n\to n+1}\}_{0\le n<N}\). Let the lower-stack audit target be the finite union

\[
\begin{aligned}
\mathsf{Target}_{\le N}\;=\;&\bigcup_{n=0}^{N}\mathsf{InstrRecord}(I_n)\;\cup\;\bigcup_{n=0}^{N-1}\mathsf{BridgeRecord}\bigl(B^{\mathrm{audit}}_{n\to n+1}\bigr)\\
&\cup\;\bigcup_{n=0}^{N}\mathsf{ClaimLog}(I_n)\;\cup\;\bigcup_{n=0}^{N}\{\,\ulcorner\psi,I_n\urcorner:\exists\chi\,(\psi,\chi)\in\mathsf{ClaimLog}(I_n)\,\}.
\end{aligned}
\]

Here \(\mathsf{InstrRecord}(I_n)\) is the set consisting of the instrument record \(I_n\) and the records named by its reference fields (Subsubsection \hyperref[sn:5.6.5]{5.6.5}); its data fields are read as fields of \(I_n\); \(\mathsf{BridgeRecord}(B)\) is the set of field records of the expanded record of \(B\) in Subsubsection \hyperref[sn:11.7.2]{11.7.2}; and each entry of \(\mathsf{ClaimLog}(I_n)\) contributes its claim record, \(\Gamma_\psi\), \(\mathcal T_\psi\), and every source, audit and bridge record they reference. The target also contains each instrument record \(I_n\) and each log record \(\mathsf{ClaimLog}(I_n)\), together with all records reachable from them in the original instance. In the displayed definition, each record family includes all records reachable from its listed members in the original instance's graph of Subsubsection \hyperref[sn:5.6.5]{5.6.5}. Thus \(\mathsf{Target}_{\le N}\) is finite and closed under resolved references, and its recorded induced graph suffices for every reachability test in the stack-defect computation. If \(\mathfrak D\) contains no verdict reference \(\ulcorner\psi,I_n\urcorner\) for some logged \((\psi,\chi)\), adjoin one with a fresh identifier. No record of \(\mathfrak D\) names it, so no reachability test, classifier verdict or claim log of \(\mathfrak D\) changes, and \(\mathfrak D\) stays admissible.

Let the finite stack-audit defect be the finite typed union

\[
\delta_{\le N}^{\mathrm{stack}}\;=\;\delta_{\mathrm{instr}}\,\cup\,\delta_{\mathrm{bridge}}\,\cup\,\delta_{\mathrm{visibility}}\,\cup\,\delta_{\mathrm{source}}\,\cup\,\delta_{\mathrm{audit}}\,\cup\,\delta_{\mathrm{circular}}\,\cup\,\delta_{\mathrm{nonclaim}}\,\cup\,\delta_{\mathrm{status}},
\]

where, writing \((\psi,\chi)\in\mathsf{ClaimLog}(I_n)\) for a claim record with its recorded verdict and \(\delta_6^{\mathrm{audit}}(\psi)\) for the audit defect of its audit record (Subsubsection \hyperref[sn:5.3.6]{5.3.6}),

\[
\begin{aligned}
\delta_{\mathrm{instr}}&=\{\,I_n:\neg\mathsf{AdmInstrument}(I_n)\,\},\\
\delta_{\mathrm{bridge}}&=\{\,B^{\mathrm{audit}}_{n\to n+1}:\neg\mathsf{AdmRotAuditBridge}(B^{\mathrm{audit}}_{n\to n+1})\,\},\\
\delta_{\mathrm{visibility}}&=\{\,(n,\psi):\delta_{\mathrm{vis}}(I_n,\psi,\mathcal L_\psi)\neq\varnothing\text{ and }\chi=\texttt{accepted}\,\},\\
\delta_{\mathrm{source}}&=\{\,(n,\psi):\texttt{source\_failure}\in\delta_6^{\mathrm{audit}}(\psi)\text{ and }\chi=\texttt{accepted}\,\},\\
\delta_{\mathrm{audit}}&=\{\,(n,\psi):\delta_6^{\mathrm{audit}}(\psi)\cap\{\texttt{missing\_audit},\texttt{failed\_audit},\texttt{visibility\_failure},\\
&\qquad\qquad\texttt{threshold\_failure},\texttt{circular\_audit}\}\neq\varnothing\text{ and }\chi=\texttt{accepted}\,\},\\
\delta_{\mathrm{nonclaim}}&=\{\,(n,\psi):\texttt{nonclaim\_failure}\in\delta_6^{\mathrm{audit}}(\psi)\text{ and }\chi=\texttt{accepted}\,\},\\
\delta_{\mathrm{circular}}&=\{\,(n,\psi):\ulcorner\psi,I_n\urcorner\text{ is reachable from }\mathsf{Use}(\psi)\text{ and }\chi=\texttt{accepted}\,\},\\
\delta_{\mathrm{status}}&=\{\,(n,\psi):\chi\text{ differs from the value the claim classifier of }I_n\\
&\qquad\qquad\text{assigns to }\psi\text{ under }\Gamma_\psi;\mathcal T_\psi;I_n\,\},
\end{aligned}
\]

with \(n\) ranging over \(0,\ldots,N\) and \((\psi,\chi)\) over \(\mathsf{ClaimLog}(I_n)\). If \(\psi\) is not well formed (Subsubsection \hyperref[sn:11.1.2]{11.1.2}), the classifier value is ``no status'', and \((n,\psi)\in\delta_{\mathrm{status}}\) for every recorded \(\chi\). Each is a finite set computed from the finite target under the computability hypothesis of Theorem \hyperref[cl:thm:25]{25}. The five summands that concern claims record an accepted claim that carries a defect, so a claim that an instrument correctly declined to accept makes the stack defect nonempty only through \(\delta_{\mathrm{status}}\), if its recorded verdict is wrong. Since the chain and its bridges are admissible by hypothesis, \(\delta_{\mathrm{instr}}=\delta_{\mathrm{bridge}}=\varnothing\). Under these hypotheses compliance is decided by \(\delta_{\mathrm{status}}\) alone (clause 2 of Theorem \hyperref[cl:thm:25]{25}). The five claim-level summands record sufficient record-level reasons why a logged acceptance is wrong; their conditions may coexist with rejection, audited absence or partial evidence. Every wrong non-acceptance appears only in \(\delta_{\mathrm{status}}\). A wrong acceptance appears only there exactly when none of the five claim-level defect conditions holds.

If \(\mathsf{Lev}\) contains a level above \(\max_{n\le N}\mathsf{Level}(I_n)\), as Theorem \hyperref[cl:thm:25]{25} assumes, let \(\ell^\star\) be such a level; the claims below are defined with it.

Fix a fresh identifier \(\mathrm{id}^\star\) for the extension instrument; below, a verdict reference under \(I^{\mathrm{stack}}_{N+1}\) means one carrying \(\mathrm{id}^\star\), and the instrument itself is constructed in the proof.

\emph{Reading guide.} The statement of Theorem \hyperref[cl:thm:25]{25} needs only the display defining \(\delta^{\mathrm{stack}}_{\le N}\) and \(\delta_{\mathrm{status}}\) and the bullets Report, Compliance, Stack audit, Audit package, Common fields of the three claims (which defines \(\mathrm{src}_{\le N}\)), Lifted soundness claims and Self-soundness; the other bullets fix the bookkeeping (use-set closure, sources, separate audits) used by the proof's row-by-row verification and may be skipped on a first reading.

Where the descriptions below identify a use-set with lower-stack records and the snapshot source, they specify its data portion. The full use-set additionally contains all statement, source, threshold and audit references required by Subsubsection \hyperref[sn:3.2.1]{3.2.1}; in particular, \(\mathsf{Refs}(\mathrm{Audit})\) contains \(A_{\mathrm{Audit}}\). Include these additional records in the reference-closed support family of Separate audits and in the extension instrument's scope and visible set. All classifier checks apply to the full use-set.

\begin{itemize}
\tightlist
\item
  \textbf{Report.} The claim \(\mathrm{Report}_{\Phi_{\le N}}(\mathbb I_{\le N})\) is the claim record with claim-type tag \(\texttt{stack\_report}\) whose statement lists the value of each component defect of \(\delta_{\le N}^{\mathrm{stack}}\) on \(\mathsf{Target}_{\le N}\), and whose use-set consists of \(\mathsf{Target}_{\le N}\cup\{s_{\le N}\}\) and the records referenced by the claim's threshold and audit fields (Subsubsection \hyperref[sn:3.2.1]{3.2.1}), where \(s_{\le N}\) denotes the source record \(\mathrm{src}_{\le N}\) below; the eight summands are computed from the records of \(\mathsf{Target}_{\le N}\) alone (they read the lower instruments \(I_n\) only through their records there) and do not depend on the extension instrument.
\item
  \textbf{Compliance.} The claim \(\mathrm{Compliant}_{\Phi_{\le N}}(\mathbb I_{\le N})\) has claim-type tag \(\texttt{stack\_compliance}\), the same use-set, and states \(\delta_{\le N}^{\mathrm{stack}}=\varnothing\).
\item
  \textbf{Stack audit.} The claim \(\mathrm{Audit}(\mathbb I_{\le N})\) has claim-type tag \(\texttt{stack\_audit}\) and the same use-set. Where \(\mathrm{Report}\) asserts the values of the eight component defects, \(\mathrm{Audit}\) asserts how they were obtained: its own audit record \(A_{\mathrm{Audit}}\) records a value for each component defect of \(\delta^{\mathrm{stack}}_{\le N}\) on \(\mathsf{Target}_{\le N}\), each computed by a rule named in the \(\mathsf{CheckRules}\) field of \(A_{\mathrm{Audit}}\), whatever that value is.
\item
  \textbf{Audit package.} Let \(\mathcal T_{\le N}^{\mathrm{audit}}\) be an admissible finite identity-lens package on \(Z=\mathsf{Target}_{\le N}\), hosted on \(\mathsf H_{\mathrm{instr}}\), with its source, package-audit and nonclaim records as required by Subsubsection \hyperref[sn:3.3.3]{3.3.3}.
\item
  \textbf{Common fields of the three claims.} Each targets \(\mathcal T_{\le N}^{\mathrm{audit}}\), has host tag \(\mathsf H_{\mathrm{instr}}\) and has level profile \((\ell^\star,\ell^\star)\), with the stack bridge recorded as a \(\texttt{reflection}\) level bridge in its \(\mathsf{LevelProfile}\), since its use-set consists of lower-level records and \(s_{\le N}\). Its evidence references are the single source record \(\mathrm{src}_{\le N}=(\mathrm{id},\texttt{committed\_state},\text{the finite snapshot of }\mathsf{Target}_{\le N},\text{valid})\), which is also the source record of \(\mathcal T_{\le N}^{\mathrm{audit}}\). Its trace holds the snapshot as data together with one passing audit of \(\mathrm{src}_{\le N}\) itself, so \(\mathrm{src}_{\le N}\) is not dirty whatever the lower-stack audits contain (Subsubsection \hyperref[sn:5.3.6]{5.3.6}); lower-stack source records enter only the use-set, as audited data. Its thresholds are \(\Theta_{\mathrm{rot}}\) and the \(\Theta_{\mathrm{pass}}\) check of the threshold policy below (every component defect is computed by its named rule), and its visibility record marks its use-set visible.
\item
  \textbf{Separate audits.} Give each claim \(\psi\) its own audit record \(A_\psi\) (here \(\psi\) ranges over the new claims Report, Compliant, Audit and \(\mathrm{Sound}^{\uparrow}_n\), never over the logged claims of \(\mathsf{ClaimLog}(I_n)\), whose audits are unchanged), with \(\mathsf{ClaimRef}=\psi\) and the bridge audit \(A_{\le N\to N+1}\) among its evidence references. These audits check that the reported computation and the check-rule verdict (rejection or not, computed from \(\delta^{\mathrm{stack}}_{\le N}\) alone) are correct; they do not evaluate the claim classifier on \(\psi\); the audit of \(\mathrm{Compliant}\) passes even when its correctly computed verdict is rejection. The bridge audit \(A_{\le N\to N+1}\) takes its \(\mathsf{EvidenceRefs}\) from \(\mathsf{Target}_{\le N}\cup\{\mathrm{src}_{\le N}\}\) and has \(\mathsf{SourceRefs}=\{\mathrm{src}_{\le N}\}\); lower-stack source records are evidence data, never source references of a new audit. The new records of \(\mathfrak D^{+}\) (the snapshot \(\mathrm{src}_{\le N}\) and its trace, the audit package \(\mathcal T^{\mathrm{audit}}_{\le N}\) and its package audits, the stack bridge and its audit, the instrument \(I^{\mathrm{stack}}_{N+1}\), its thresholds, the separate audits and the stack claims) are constructed so that, together with \(\mathsf{Target}_{\le N}\), they form a reference-closed support family containing no verdict reference under \(I^{\mathrm{stack}}_{N+1}\); the candidate \(\mathrm{Sound}(I^{\mathrm{stack}}_{N+1})\) and its self-verdict reference lie outside this family and are named by none of its records. All references of the original records are unchanged. Consequently no verdict reference under \(I^{\mathrm{stack}}_{N+1}\) is reachable from the evidence or source references of any audit \(A_\psi\).
\item
  \textbf{Lifted soundness claims.} For each \(n\le N\), let \(\mathrm{Sound}^{\uparrow}_n\) be a separate claim record under \(\mathcal T_{\le N}^{\mathrm{audit}}\). It has claim type \(\texttt{soundness}\), asserts that every verdict in \(\mathsf{ClaimLog}(I_n)\) agrees with \(I_n\)'s claim classifier, and has \(\mathsf{Refs}(\mathrm{Sound}^{\uparrow}_n)=\displaystyle\bigcup_{(\psi,\chi)\in\mathsf{ClaimLog}(I_n)}\{\ulcorner\psi,I_n\urcorner,\psi,\Gamma_\psi,\mathcal T_\psi\}\cup\{I_n,\mathsf{ClaimLog}(I_n),s_{\le N}\}\subseteq\mathsf{Target}_{\le N}\cup\{s_{\le N}\}\), so that its use-set (Subsubsection \hyperref[sn:3.2.1]{3.2.1}) is this set together with the source, threshold and audit records required in every claim use-set. Here \(\Gamma_\psi\) is the single context record of the log entry (Subsubsection \hyperref[sn:11.5.1]{11.5.1}), whose reference fields name the context's members; these are reachable from, but not members of, \(\mathsf{Use}(\mathrm{Sound}^{\uparrow}_n)\), so its only verdict references are the \(\ulcorner\psi,I_n\urcorner\). Give it the host tag \(\mathsf H_{\mathrm{instr}}\), the level profile \((\ell^\star,\ell^\star)\) with its upper-level reflection bridge, source, evidence, visibility, threshold, nonclaim and separate audit fields specified above for the three stack claims; its audit checks the computation even when the claim is rejected. Each separate audit \(A_\psi\) has \(\mathsf{SourceRefs}=\{\mathrm{src}_{\le N}\}\), and its evidence references include the bridge audit; the source trace records the snapshot and its package audit. The extension instrument asserted by Theorem \hyperref[cl:thm:25]{25} is required to have \(s_{\le N}\) in its scope and visible set and to admit \(\mathrm{src}_{\le N}\) under its source and evidence policies; the proof constructs one. No record reachable from this source or from the lower-stack part of a use-set is a verdict reference for a new stack claim, since the stack claims have fresh identifiers. Set \(\mathsf{partialEvidence}=\mathrm{false}\) on Report, Compliant, Audit and every lifted soundness claim.
\item
  \textbf{Self-soundness.} The claim \(\mathrm{Sound}(I)\) is the complete self-soundness claim of Subsubsection \hyperref[sn:11.5.1]{11.5.1}.
\end{itemize}

\begin{claimtheorem}{Theorem 25 (FiniteRotatingAudit)}\phantomsection\label{cl:thm:25}

Let \(N\ge0\), and let \(\mathsf{Target}_{\le N}\), \(\delta^{\mathrm{stack}}_{\le N}\), \(\delta_{\mathrm{status}}\), \(\mathrm{src}_{\le N}\), \(\mathcal T^{\mathrm{audit}}_{\le N}\) and the claims \(\mathrm{Report}\), \(\mathrm{Compliant}\), \(\mathrm{Audit}\), \(\mathrm{Sound}^{\uparrow}_n\) be as in the Setting above (in particular \(\delta_{\mathrm{status}}\) is the set of \((n,\psi)\), \((\psi,\chi)\in\mathsf{ClaimLog}(I_n)\), whose recorded verdict \(\chi\) differs from the status \(I_n\)'s claim classifier assigns \(\psi\) under its recorded context, and \(\mathrm{Sound}^{\uparrow}_n\) asserts that \(\delta_{\mathrm{status}}\) has no entry with index \(n\)). Suppose that \(\mathsf{Lev}\) contains a level above \(\max_{0\le n\le N}\mathsf{Level}(I_n)\) in the order of Subsubsection \hyperref[sn:2.2.3]{2.2.3}, and let \(\ell^{\star}\) be the level fixed above. Suppose also that every threshold evaluation reachable in the evaluation dependency graph of Subsubsection \hyperref[sn:5.3.6]{5.3.6} from the computation of any of the eight summands of \(\delta^{\mathrm{stack}}_{\le N}\), including logged-claim classification, audit-defect computation and bridge admissibility checks, is computed by a program of the rule language \(\mathcal R\) of Subsubsection \hyperref[sn:3.4.1]{3.4.1} (as for the exact tests on finite tables used in Theorem \hyperref[cl:thm:5]{5}(e)). (The levels of the chain strictly increase, so the hypothesis that \(\mathsf{Lev}\) contains a level above the chain requires \(N\le|\mathsf{Lev}|-2\), that is \(N\le1\) for the three levels of Subsubsection \hyperref[sn:2.2.3]{2.2.3}. The proof uses only that \(\ell^\star\) lies above every level of the chain; it applies verbatim when \(\mathsf{Lev}\) is replaced by any finite linear order of levels, where the stack may contain up to \(|\mathsf{Lev}|-1\) instruments.) Then, in an admissible extension \(\mathfrak D^{+}\) of the instance, there exists an extension instrument

\[
I_{N+1}^{\mathrm{stack}},
\]

whose scope and visible set contain \(\mathsf{Target}_{\le N}\cup\{s_{\le N}\}\) and whose source and evidence policies admit \(\mathrm{src}_{\le N}\), and an admissible audit bridge

\[
B^{\mathrm{audit}}_{\le N\to N+1}
\]

with \(\mathsf{Level}\bigl(I_{N+1}^{\mathrm{stack}}\bigr)=\ell^{\star}>\max_{0\le n\le N}\mathsf{Level}(I_n)\) and target scope \(\Phi_{\le N}=\mathsf{Target}_{\le N}\) such that the following hold:

\begin{enumerate}
\def\labelenumi{\arabic{enumi}.}
\tightlist
\item
  (Report acceptance.) The claim listing the computed values of the eight component defects is accepted:
  \[
  \Gamma\,;\;\mathcal T_{\le N}^{\mathrm{audit}}\,;\;I_{N+1}^{\mathrm{stack}}\;\vdash\;\mathrm{Report}_{\Phi_{\le N}}(\mathbb I_{\le N})\;:\;\texttt{accepted}.
  \]
\item
  (Compliance characterization.)
  \[
  \Gamma\,;\;\mathcal T_{\le N}^{\mathrm{audit}}\,;\;I_{N+1}^{\mathrm{stack}}\;\vdash\;\mathrm{Compliant}_{\Phi_{\le N}}(\mathbb I_{\le N})\;:\;\texttt{accepted}\quad\Longleftrightarrow\quad\delta_{\le N}^{\mathrm{stack}}=\varnothing\quad\Longleftrightarrow\quad\delta_{\mathrm{status}}=\varnothing.
  \]
\item
  (Non-finality.) In \(\mathfrak D^+\),
  \[
  \Gamma\,;\;\mathcal T_{\le N}^{\mathrm{audit}}\,;\;I_{N+1}^{\mathrm{stack}}\;\vdash\;\mathrm{Audit}(\mathbb I_{\le N})\;:\;\texttt{accepted},
  \]
  while \(\mathrm{Sound}\bigl(I_{N+1}^{\mathrm{stack}}\bigr)\) does not receive \(\texttt{accepted}\) under \(I_{N+1}^{\mathrm{stack}}\); in particular the acceptance of \(\mathrm{Audit}(\mathbb I_{\le N})\) does not imply that of \(\mathrm{Sound}\bigl(I_{N+1}^{\mathrm{stack}}\bigr)\).
\item
  (Soundness of lower instruments.) \(\texttt{soundness}\in\mathsf{ClaimTypes}(I_{N+1}^{\mathrm{stack}})\), and for each \(n\le N\), the lifted claim \(\mathrm{Sound}^{\uparrow}_n\) (that every verdict recorded in \(\mathsf{ClaimLog}(I_n)\) equals the status \(I_n\)'s claim classifier assigns under the recorded context) satisfies
  \[
  \Gamma\,;\;\mathcal T_{\le N}^{\mathrm{audit}}\,;\;I_{N+1}^{\mathrm{stack}}\;\vdash\;\mathrm{Sound}^{\uparrow}_n\;:\;\texttt{accepted}\quad\Longleftrightarrow\quad\delta_{\mathrm{status}}\ \text{has no entry with index}\ n.
  \]
\end{enumerate}

Consequently, by clause 2, \(\mathrm{Compliant}_{\Phi_{\le N}}(\mathbb I_{\le N})\) is \(\texttt{accepted}\) if and only if \(\mathrm{Sound}^{\uparrow}_n\) is \(\texttt{accepted}\) for every \(n\le N\).

\end{claimtheorem}

\begin{proof}

The proof constructs the extension instrument, verifies the audit bridge is admissible, and reads off the four clauses from the rotating-audit acceptance condition and the rotating-audit non-finality statement.

\emph{Construction of \(I_{N+1}^{\mathrm{stack}}\).} Define \(I_{N+1}^{\mathrm{stack}}\) as a finite instrument record in the form of Subsubsection \hyperref[sn:11.1.1]{11.1.1}, with components fixed as follows:

\begin{itemize}
\tightlist
\item
  \(\mathsf{ProvisionalTypes}\bigl(I_{N+1}^{\mathrm{stack}}\bigr)=\varnothing\);
\item
  \(\mathsf{Scope}\bigl(I_{N+1}^{\mathrm{stack}}\bigr)\supseteq\mathsf{Target}_{\le N}\), including the verdict references \(\ulcorner\psi,I_n\urcorner\) of the claim logs, the full record \(\mathrm{src}_{\le N}\) and every record referenced by the stack claims' threshold and audit fields, with \(\mathsf{Hosts}(I_{N+1}^{\mathrm{stack}})\ni\mathsf H_{\mathrm{instr}}\) and \(\mathsf{Levels}(I_{N+1}^{\mathrm{stack}})\ni\ell^\star\);
\item
  \(\mathsf{ClaimTypes}\bigl(I_{N+1}^{\mathrm{stack}}\bigr)\supseteq\{\,\texttt{stack\_report},\;\texttt{stack\_compliance},\;\texttt{stack\_audit},\;\texttt{soundness}\,\}\);
\item
  \(\mathsf{Visible}\bigl(I_{N+1}^{\mathrm{stack}}\bigr)\) contains \(\mathsf{Target}_{\le N}\), the full record \(\mathrm{src}_{\le N}\) and every record referenced by the stack claims' threshold and audit fields, and \(\mathsf{Suppressed}\bigl(I_{N+1}^{\mathrm{stack}}\bigr)\cap\mathsf{Target}_{\le N}=\varnothing\), so every record of the audit target is visible to the extension instrument and none of the three claims requires a visibility bridge of Subsubsection \hyperref[sn:3.5.3]{3.5.3};
\item
  \(\mathsf{EvidencePolicy}\) admits the finite lower-stack records as audit evidence; \(\mathsf{SourcePolicy}\) contains \((\texttt{committed\_state},\texttt{full},\lambda)\) for every profile \(\lambda\) of \(\mathfrak D^{+}\), including \(\lambda_{\mathrm{rot}}\), and for its default entry, so it admits \(\mathrm{src}_{\le N}\), while the recorded provenance of the lower stack is data to be checked, without requiring the lower stack to be compliant; and \(\mathsf{BridgePolicy}\) admits the declared upward rotating-audit bridge and its level shift, with \(\mathsf{Strength}=\{\texttt{weak}\}\), every admitted (bridge type, regime) pair, among them \((\texttt{reflection},\texttt{audit})\), and every claim type sent to \(\texttt{weak}\); \(\mathsf{EvidencePolicy}\) also admits \(\texttt{committed\_state}\) as direct evidence;
\item
  Set the extension instrument's mode for the rotating-audit-bridge family to \(\texttt{rotating-audit}\) and include \((\texttt{fine},\texttt{audit},\texttt{global},\mathsf H_{\mathrm{instr}})\) in its compatibility relation; its bridge policy admits \(\texttt{reflection}\).
\item
  \(\mathsf{CheckRules}\bigl(I_{N+1}^{\mathrm{stack}}\bigr)\) contains the finite rules that compute each of \(\delta_{\mathrm{instr}}\), \(\delta_{\mathrm{bridge}}\), \(\delta_{\mathrm{visibility}}\), \(\delta_{\mathrm{source}}\), \(\delta_{\mathrm{audit}}\), \(\delta_{\mathrm{circular}}\), \(\delta_{\mathrm{nonclaim}}\), \(\delta_{\mathrm{status}}\) on the finite records of \(\mathsf{Target}_{\le N}\), and a rule for \(\mathrm{Compliant}_{\Phi_{\le N}}(\mathbb I_{\le N})\) that returns the rejection verdict exactly when the computed \(\delta_{\le N}^{\mathrm{stack}}\) is non-empty, and rules for \(\mathrm{Report}_{\Phi_{\le N}}(\mathbb I_{\le N})\) and \(\mathrm{Audit}(\mathbb I_{\le N})\) that return the rejection verdict exactly when a listed value differs from the computed one or a component defect is not computed; and a rule for \(\mathrm{Sound}^{\uparrow}_n\) that returns the rejection verdict exactly when \(\delta_{\mathrm{status}}\) has an entry with index \(n\);
\item
  \(\mathsf{ThresholdPolicy}\bigl(I_{N+1}^{\mathrm{stack}}\bigr)\) applies the source instruments' per-component thresholds when computing the reported lower-stack defects. For the \(\mathrm{Report}\), \(\mathrm{Compliant}\), and \(\mathrm{Audit}\) claims themselves it requires \(\Theta_{\mathrm{pass}}\) on the recorded Boolean that every component defect of \(\delta^{\mathrm{stack}}_{\le N}\) is computed on \(\mathsf{Target}_{\le N}\) by its named rule, and the bridge threshold \(\Theta_{\mathrm{rot}}\), the instance of \(\Theta_{\in S}\) with \(S=\{\texttt{audit\_passes}\}\), on \(\delta_{\mathrm{rot}}\). A correctly reported failure of a lower threshold does not fail any of these claims' own threshold check;
\item
  \(\mathsf{AuditPolicy}\bigl(I_{N+1}^{\mathrm{stack}}\bigr)\) contains the audit rules for the bridge \(B^{\mathrm{audit}}_{\le N\to N+1}\);
\item
  \(\mathsf{Level}\bigl(I_{N+1}^{\mathrm{stack}}\bigr)=\ell^{\star}\), which exceeds \(\max_{0\le n\le N}\mathsf{Level}(I_n)\) by hypothesis;
\item
  \(\mathcal N\bigl(I_{N+1}^{\mathrm{stack}}\bigr)\) contains \(\mathrm{NC}\text{-}16\), \(\mathrm{NC}\text{-}17\), and \(\mathrm{NC}\text{-}18\).
\end{itemize}

Each component is finite and of the schema of Subsubsection \hyperref[sn:3.4.1]{3.4.1}: \(\mathsf{Target}_{\le N}\) is a finite union of finite records, the check rules, thresholds, audit rules and nonclaim register are finite, \(\ell^\star\in\mathsf{Lev}\) and \(\mathsf H_{\mathrm{instr}}\in\mathsf{Host}\). Each listed rule is an \(\mathcal R\)-program, by the computability hypothesis: the rules for \(\delta_{\mathrm{status}}\) and \(\delta_{\mathrm{instr}}\) call the finitely many programs of \(\mathsf{CheckRules}(I_n)\), \(n\le N\), copied into \(\mathsf{CheckRules}(I_{N+1}^{\mathrm{stack}})\), and otherwise use table lookups, comparisons and iteration over the finite target. Include every rule replayed by an audit judged under \(I_{N+1}^{\mathrm{stack}}\) in its \(\mathsf{CheckRules}\), give each such audit rules from that family, and put the audit records \(A_\psi\) and \(A_{\le N\to N+1}\) in its scope and visible set. Thus \(\mathsf{AdmInstrument}\bigl(I_{N+1}^{\mathrm{stack}}\bigr)\) holds.

\emph{The extended instance.} Every new record of \(\mathfrak D^{+}\) below carries host tag \(\mathsf H_{\mathrm{instr}}\) and level tag \(\ell^\star\), which lie in \(\mathsf{Hosts}\) and \(\mathsf{Levels}\) of \(I^{\mathrm{stack}}_{N+1}\), so \(\mathsf{Eval}\) finds the data of \(\Theta_{\mathrm{rot}}\) and \(\Theta_{\mathrm{pass}}\) in scope. Let \(\mathfrak D^{+}\) extend \(\mathfrak D\), with fresh identifiers, by \(I_{N+1}^{\mathrm{stack}}\), \(B^{\mathrm{audit}}_{\le N\to N+1}\), \(\mathcal T^{\mathrm{audit}}_{\le N}\) and the claim records \(\mathrm{Report}\), \(\mathrm{Compliant}\), \(\mathrm{Audit}\) and \(\mathrm{Sound}^{\uparrow}_n\), a declared log \(\mathsf{ClaimLog}(I_{N+1}^{\mathrm{stack}})\) (any finite set of claims of \(\mathfrak D^{+}\) with recorded verdicts, for instance the stack claims with the verdicts of clauses 1--4), with the verdict references of its entries, none named by a record of the support family, and the candidate \(\mathrm{Sound}(I_{N+1}^{\mathrm{stack}})\) with its self-verdict reference, with their typed audit and source records, which lie outside the scope of every instrument \(I\) of \(\mathfrak D\), so \(\mathsf{vis}_I\) is \(\texttt{outside}\) on them with \(I\) unchanged, and choosing \(\mathsf{vis}_{I_{N+1}^{\mathrm{stack}}}\) on the records of \(\mathfrak D\) outside \(\mathsf{Target}_{\le N}\). No record of \(\mathsf{Target}_{\le N}\) changes class and no \(\mathsf{ClaimLog}(I_n)\) changes, so \(\delta^{\mathrm{stack}}_{\le N}\) is unchanged. No evaluation of \(\mathfrak D\) queries a new record, and each new evaluation queries only evaluations of \(\mathfrak D\) (through the classifiers of the \(I_n\)), of \(\mathrm{src}_{\le N}\), or of the stack bridge and its audit; so the evaluation dependency graph of \(\mathfrak D^{+}\) is acyclic and \(\mathfrak D^{+}\) is admissible, and clauses 1--4 are judgments of \(\mathfrak D^{+}\).

\emph{Construction of the audit bridge.} Define \(B^{\mathrm{audit}}_{\le N\to N+1}\) to be the stack rotating-audit bridge of Subsubsection \hyperref[sn:11.7.2]{11.7.2} with target scope \(\Phi_{\le N}=\mathsf{Target}_{\le N}\), language map \(L\) that translates every record in \(\Phi_{\le N}\), including instrument records, bridge records, and claim-log records, into well-typed claims of \(I_{N+1}^{\mathrm{stack}}\), visibility map \(V\) exposing each record in \(\Phi_{\le N}\) either directly or through visibility bridges, audit thresholds \(\Theta\) that use the lower component thresholds to compute reported defects, while testing the bridge itself for complete, correct reporting and \(\delta_{\mathrm{rot}}=\texttt{audit\_passes}\), audit record \(A\) that records the finite checks above are auditable in \(I_{N+1}^{\mathrm{stack}}\), a bridge defect record with \(\delta_{\mathrm{rot}}=\texttt{audit\_passes}\), and a nonclaim register containing \(\mathrm{NC}\text{-}16\), \(\mathrm{NC}\text{-}17\), and \(\mathrm{NC}\text{-}18\). The bridge defect is \(\texttt{audit\_passes}\) because the check rules of \(I_{N+1}^{\mathrm{stack}}\) are the deterministic computations of the eight summands on the finite target, so the reported values are the computed ones; the computed \(\delta_{\le N}^{\mathrm{stack}}\) is stored separately as the report's result. Record \(L\) as a finite table mapping each \(r\in\Phi_{\le N}\) to a well-formed \(\texttt{stack\_report}\) claim reporting its recomputed audit data, with all statement identifiers in \(\mathsf{Refs}\). Include these claims and their separate audits in the reference-closed support family, scope and visible set. Register every bridge threshold in \(\mathsf{ThresholdPolicy}_{I_{N+1}^{\mathrm{stack}}}\). Set the bridge audit's \(\mathsf{ClaimRef}\) to the bridge and draw its rules from \(\mathsf{CheckRules}_{I_{N+1}^{\mathrm{stack}}}\).

\emph{Admissibility of the bridge.} Each clause of \(\mathsf{AdmRotAuditBridge}\bigl(B^{\mathrm{audit}}_{\le N\to N+1}\bigr)\) of Subsubsection \hyperref[sn:11.7.3]{11.7.3} is verified on the finite records: clause (1) follows from instrument admissibility and the reflection, mode and compatibility entries just specified; clause (2) is the strict level inequality fixed by construction; clause (3) is the finite declaration of \(\Phi_{\le N}\); clause (4) is satisfied because \(I_{N+1}^{\mathrm{stack}}\)'s scope contains \(\mathsf{Target}_{\le N}\); clause (5) is the language map; clauses (6)--(9) are the visibility, threshold, audit, and defect-typing requirements built into the construction; clause (10) holds because \(\Phi_{\le N}\) lies in the reference-closed support family of the separate audits, which contains no verdict reference under the fresh extension instrument; clause (11) is the nonclaim presence. Hence the bridge is admissible.

\emph{Clause 1 (Report acceptance).} We go through the rows of the claim classifier of Subsubsection \hyperref[sn:5.6.6]{5.6.6} for the report claim. The row \(\texttt{outside\_scope}\) does not apply: its tag \(\texttt{stack\_report}\) is in \(\mathsf{ClaimTypes}(I_{N+1}^{\mathrm{stack}})\), its host tag \(\mathsf H_{\mathrm{instr}}\) and level tag \(\ell^\star\) are in \(\mathsf{Hosts}\) and \(\mathsf{Levels}\) of \(I_{N+1}^{\mathrm{stack}}\), and its full use-set is contained in the scope. The row \(\texttt{blocked}\) does not apply: every record of the use-set is visible, so \(\delta_{\mathrm{overread}}=\delta_{\mathrm{unknown}}=\varnothing\) and no bridge is required, so \(\delta_{\mathrm{weakbridge}}=\varnothing\); the check rules read the records of \(\mathsf{Target}_{\le N}\) in the use-set, which have the types the rules read; and the source check passes because the evidence is the single source record \(\mathrm{src}_{\le N}\), whose type \(\texttt{committed\_state}\) is admitted by \(\mathsf{SourcePolicy}(I_{N+1}^{\mathrm{stack}})\) and \(\mathsf{EvidencePolicy}(I_{N+1}^{\mathrm{stack}})\), whatever the source records of the lower stack contain. The row \(\texttt{absent\_with\_record}\) does not apply, since every record of the use-set is present in the reference-closed support family. The claim's audit record is \(A_{\mathrm{Report}}\): its source references resolve to the admissible source records, its visibility check passes because every record named by its \(\mathsf{EvidenceRefs}\) and \(\mathsf{SourceRefs}\) (among them \(A_{\le N\to N+1}\) and \(\mathrm{src}_{\le N}\)) lies in the reference-closed support family, which is in the visible set of \(I^{\mathrm{stack}}_{N+1}\) (Subsubsection \hyperref[sn:4.9.2]{4.9.2}), its circularity check passes by the reference-closed support construction in \emph{Separate audits}, since no verdict reference under \(I^{\mathrm{stack}}_{N+1}\) is reachable from its evidence or source references, its threshold check passes because the data of \(\Theta_{\mathrm{rot}}\) and \(\Theta_{\mathrm{pass}}\) lie on in-scope new records and the reported values are the computed ones, its nonclaims are present, and its check results pass. By Subsubsection \hyperref[sn:5.3.6]{5.3.6} its audit defect is empty, that is \(\texttt{audit\_passes}\). So the row \(\texttt{failed\_audit}\) does not apply; the row \(\texttt{undefined\_circular}\) does not apply, since the audit defect is empty, the use-set lies in the reference-closed support family of \emph{Separate audits}, which contains no verdict reference under \(I_{N+1}^{\mathrm{stack}}\), so none is reachable from it, and the claim type is not \(\texttt{soundness}\); and the row \(\texttt{below\_threshold}\) does not apply, since the audit defect is empty. The row \(\texttt{rejected}\) does not apply, since the statement lists exactly the computed values, so the check rule does not reject. The row \(\texttt{provisional}\) does not apply, since the evidence is complete. Hence the \(\texttt{accepted}\) row holds:

\[
\Gamma\,;\;\mathcal T_{\le N}^{\mathrm{audit}}\,;\;I_{N+1}^{\mathrm{stack}}\;\vdash\;\mathrm{Report}_{\Phi_{\le N}}(\mathbb I_{\le N})\;:\;\texttt{accepted}.
\]

The report claim records the value of \(\delta_{\le N}^{\mathrm{stack}}\) together with the per-component defect values; it is accepted whether or not \(\delta_{\le N}^{\mathrm{stack}}\) is empty. The same row-by-row argument applies verbatim to \(\mathrm{Audit}(\mathbb I_{\le N})\), whose statement holds by construction of the check rules, so \(\mathrm{Audit}(\mathbb I_{\le N})\) is \(\texttt{accepted}\) under \(I_{N+1}^{\mathrm{stack}}\) as well.

\emph{Clause 2 (Compliance characterization).} The compliance claim \(\mathrm{Compliant}_{\Phi_{\le N}}(\mathbb I_{\le N})\) is in the claim types and scope of \(I_{N+1}^{\mathrm{stack}}\), its full use-set is visible through the bridge, its records are present, and its audit and threshold policies pass on the finite records, as for the report claim in clause 1. Its use-set lies in the reference-closed support family, which contains no verdict reference under \(I_{N+1}^{\mathrm{stack}}\), so none is reachable from it; its type is not \(\texttt{soundness}\) and its audit defect is empty, so the \(\texttt{undefined\_circular}\) row of Subsubsection \hyperref[sn:5.6.5]{5.6.5} does not apply. So the rows above \(\texttt{rejected}\) do not apply, and the verdict is decided by the \(\texttt{rejected}\) row, whose condition is the rejection verdict of the check rules. By construction that rule rejects exactly when \(\delta_{\le N}^{\mathrm{stack}}\neq\varnothing\). If \(\delta_{\le N}^{\mathrm{stack}}\neq\varnothing\), the claim is therefore \(\texttt{rejected}\); if \(\delta_{\le N}^{\mathrm{stack}}=\varnothing\), neither \(\texttt{rejected}\) nor \(\texttt{provisional}\) applies (the evidence is complete), and the \(\texttt{accepted}\) row holds. Hence the equivalence

\[
\Gamma\,;\;\mathcal T_{\le N}^{\mathrm{audit}}\,;\;I_{N+1}^{\mathrm{stack}}\;\vdash\;\mathrm{Compliant}_{\Phi_{\le N}}(\mathbb I_{\le N})\;:\;\texttt{accepted}\quad\Longleftrightarrow\quad\delta_{\le N}^{\mathrm{stack}}=\varnothing.
\]

For the second equivalence, \(\delta_{\mathrm{instr}}=\delta_{\mathrm{bridge}}=\varnothing\) by hypothesis, and \(\delta_{\mathrm{status}}\subseteq\delta_{\le N}^{\mathrm{stack}}\). Each entry \((n,\psi,\ldots)\) of the five remaining claim-level summands (\(\delta_{\mathrm{visibility}}\), \(\delta_{\mathrm{source}}\), \(\delta_{\mathrm{audit}}\), \(\delta_{\mathrm{nonclaim}}\), \(\delta_{\mathrm{circular}}\)) has recorded verdict \(\chi=\texttt{accepted}\) and a record condition under which the claim classifier of Subsubsection \hyperref[sn:5.6.5]{5.6.5} does not return \(\texttt{accepted}\): a nonempty visibility defect or a source or visibility failure gives the \(\texttt{blocked}\) row, or, for \(\delta_{\mathrm{outside}}\), the \(\texttt{outside\_scope}\) row, a missing, failed or nonclaim audit gives \(\texttt{failed\_audit}\), a circular audit or a self-reference gives \(\texttt{undefined\_circular}\), and a threshold failure gives \(\texttt{below\_threshold}\) unless an earlier row applies. So \((n,\psi)\in\delta_{\mathrm{status}}\), and \(\delta_{\le N}^{\mathrm{stack}}=\varnothing\) exactly when \(\delta_{\mathrm{status}}=\varnothing\).

\emph{Clause 3 (Non-finality).} By Theorem \hyperref[cl:thm:24]{24} of Subsection \hyperref[sn:11.5]{11.5}, the same-level self-soundness claim \(\mathrm{Sound}\bigl(I_{N+1}^{\mathrm{stack}}\bigr)\) does not receive verdict \(\texttt{accepted}\) under \(I_{N+1}^{\mathrm{stack}}\) in the absence of an admissible level shift to a strictly higher instrument. The rotating-audit acceptance of \(\mathrm{Audit}(\mathbb I_{\le N})\) concerns the lower stack; it is not an audit of \(I_{N+1}^{\mathrm{stack}}\) by a still higher instrument, and so does not supply a level shift over \(I_{N+1}^{\mathrm{stack}}\). By clause 1, \(\mathrm{Audit}(\mathbb I_{\le N})\) is accepted, while by Theorem \hyperref[cl:thm:24]{24} \(\mathrm{Sound}(I_{N+1}^{\mathrm{stack}})\) is not; this one instance refutes the implication. Hence

\[
\begin{aligned}
&\Gamma\,;\;\mathcal T_{\le N}^{\mathrm{audit}}\,;\;I_{N+1}^{\mathrm{stack}}\;\vdash\;\mathrm{Audit}(\mathbb I_{\le N})\;:\;\texttt{accepted}\\
&\quad\not\Longrightarrow\quad\Gamma\,;\;\mathcal T_{\le N}^{\mathrm{audit}}\,;\;I_{N+1}^{\mathrm{stack}}\;\vdash\;\mathrm{Sound}\bigl(I_{N+1}^{\mathrm{stack}}\bigr)\;:\;\texttt{accepted}.
\end{aligned}
\]

\emph{Clause 4 (Soundness of lower instruments).} The lower-stack portion of the use-set of \(\mathrm{Sound}^{\uparrow}_n\) consists of the recorded, visible lower-instrument verdict references and the claim records, contexts and packages of \(\mathsf{ClaimLog}(I_n)\), all in \(\mathsf{Target}_{\le N}\). The row-by-row argument of clause 1 applies: \(\texttt{soundness}\in\mathsf{ClaimTypes}(I_{N+1}^{\mathrm{stack}})\), and the use-set, contained in the reference-closed support family, is in scope and visible; the separate audit of \(\mathrm{Sound}^{\uparrow}_n\) has empty audit defect; and no threshold is violated, since correctly reported lower failures are report data. So the scope, visibility, absence, audit and threshold rows do not apply. Its use-set contains no verdict reference for \(\mathrm{Sound}^{\uparrow}_n\) under \(I_{N+1}^{\mathrm{stack}}\), and every verdict reference in it belongs to an \(I_n\) with \(\mathsf{Level}(I_n)<\mathsf{Level}(I_{N+1}^{\mathrm{stack}})\); hence none of the three conditions of the \(\texttt{undefined\_circular}\) row applies (the audit defect is empty, as above). The check rule rejects exactly when \(\delta_{\mathrm{status}}\) has an entry indexed by \(n\); otherwise the complete claim reaches \(\texttt{accepted}\).

This completes the proof.

\end{proof}

\subsubsection{Nonfinality}\label{sn:11.8.1}

Theorem \hyperref[cl:thm:25]{25} leaves the higher instrument \(I_{N+1}^{\mathrm{stack}}\) uncertified by itself. To audit \(I_{N+1}^{\mathrm{stack}}\), one treats

\[
\mathbb I_{\le N+1}\;=\;(I_0,\,I_1,\,\ldots,\,I_N,\,I_{N+1}^{\mathrm{stack}})
\]

as the next finite lower stack. This is possible only when \(N=0\) and \(\ell^\star=\mathsf{Ver}\): if \(\mathsf{Level}(I_0)=\mathsf{Beh}\) and \(\mathsf{Level}(I_1^{\mathrm{stack}})=\mathsf{Ver}\), then \((I_0,I_1^{\mathrm{stack}})\) with the stack bridge \(B^{\mathrm{audit}}_{\le0\to1}\) of that proof as its bridge \(B^{\mathrm{audit}}_{0\to1}\) (its scope \(\mathsf{Target}_{\le0}\) contains \(\mathsf{InstrRecord}(I_0)\cup\mathsf{ClaimLog}(I_0)\); clauses (1)--(11) were verified there) and the log of \(I_1^{\mathrm{stack}}\) declared in \(\mathfrak D^{+}\), is a chain with \(N=1\), and Theorem \hyperref[cl:thm:25]{25} yields an instrument \(I_2^{\mathrm{stack}}\) at level \(\mathsf{Str}\). No further step exists, since \(\mathsf{Lev}\) has no level above \(\mathsf{Str}\), and that top instrument is not audited by any higher instrument. With the three levels of \(\mathsf{Lev}\), a chain with strictly increasing levels has at most three instruments, and the hypothesis of Theorem \hyperref[cl:thm:25]{25} holds exactly for chains whose highest level is \(\mathsf{Beh}\) or \(\mathsf{Ver}\); such chains have \(N\le1\). No tower has a final self-certifying top, because each instrument of the tower inherits Theorem \hyperref[cl:thm:24]{24}'s classifier routing for its own soundness claim. This is the rotating-audit non-finality clause \(\mathrm{NC}\text{-}17\) of Subsection \hyperref[sn:15.2]{15.2}.

The instrument-family completeness claim, in the sense of an instrument that audits every higher level of itself, is therefore not within the scope of \(\mathbf{BirdInt}^{\mathrm{aud}}_{\mathrm{fin}}\). Any positive soundness statement about an instrument \(I_n\) requires an admissible level-shifted audit by a strictly higher instrument \(I_{n+1}\) (Subsubsection \hyperref[sn:5.6.5]{5.6.5}); the rotating-audit bridge of Subsection \hyperref[sn:11.7]{11.7} records such an audit; the framework records the chain, the bridge, the admissibility predicate, and the defect, and accepts \(\mathrm{Compliant}_{\Phi_{\le N}}(\mathbb I_{\le N})\) iff \(\delta_{\le N}^{\mathrm{stack}}=\varnothing\), but never accepts \(\mathrm{Sound}(I_n)\) from the same instrument \(I_n\) itself.

\subsection{Instrument Transfer}\label{sn:11.9}

This subsection records the instrument-transfer rule that governs whether a claim accepted under one instrument may be transported to another, and marks the corresponding full classifier as future work. The rule supports the exclusion of instrument-free acceptance in Subsubsection \hyperref[sn:2.4.3]{2.4.3}: \(I\vdash\varphi\) does not imply \(J\vdash\varphi\).

\subsubsection{Transfer Rule}\label{sn:11.9.1}

Let \(I\) and \(J\) be admissible instruments in the sense of Subsubsection \hyperref[sn:11.1.1]{11.1.1}. Suppose

\[
\Gamma\,;\;\mathcal T\,;\;I\;\vdash\;\varphi\;:\;\texttt{accepted}.
\]

Acceptance under \(I\) does not, by itself, license any claim status under \(J\): the default is

\[
\Gamma\,;\;\mathcal T\,;\;I\;\vdash\;\varphi\;:\;\texttt{accepted}\quad\not\Longrightarrow\quad\Gamma\,;\;\mathcal T\,;\;J\;\vdash\;\varphi\;:\;\texttt{accepted}.
\]

A non-default transfer requires an instrument-transfer bridge

\[
B_{I\to J}\;=\;(\,\mathrm{id}_B,\;I,\;J,\;\varphi,\;L_{I\to J},\;V_{I\to J},\;\Theta_{I\to J},\;A_{I\to J},\;\delta_{I\to J},\;\mathcal N_{I\to J},\;\lambda_{\mathrm{transfer}}\,),
\]

a directed bridge from \(I\) to \(J\) admitted by the surrounding bridge policies, with bridge type \(\texttt{instrument-transfer}\) and components recording, respectively: the source claim, the language map that translates \(\varphi\) into \(J\)'s claim language, the visibility map exposing the records that supported \(\varphi\) under \(I\) to \(J\), the audit thresholds \(J\) applies to the transfer, the audit record of the bridge, the transfer defect, the bridge nonclaim register, and the bridge profile. Admissibility of \(B_{I\to J}\) for the source claim \(\varphi\), written

\[
\mathsf{AdmInstrBridge}(B_{I\to J},\varphi),
\]

requires \(I\) and \(J\) admissible, \(\varphi\) in the claim types of both instruments via \(L_{I\to J}\), the supporting records visible to \(J\) either directly or through visibility bridges of Subsubsection \hyperref[sn:3.5.3]{3.5.3}, audit thresholds declared, the bridge audited, no unstratified circular dependency between \(\varphi\) in \(I\) and the records that support its transfer to \(J\), and the nonclaim register present.

The instrument-transfer rule is then

\[
\frac{\Gamma\,;\;\mathcal T\,;\;I\;\vdash\;\varphi\;:\;\texttt{accepted}\quad\quad\mathsf{AdmInstrBridge}(B_{I\to J},\varphi)}{\Gamma\,;\;\mathcal T\,;\;J\;\vdash\;B_{I\to J}(\varphi)\;:\;\chi},
\]

where \(\chi\in\mathrm{ClaimStatus}\) is the value that \(J\)'s claim classifier (Subsubsection \hyperref[sn:5.6.5]{5.6.5}) assigns to the transferred claim record. Acceptance under \(I\) thus enters the transfer as a premise; the verdict under \(J\) is decided by \(J\)'s own classifier on the bridged claim, and \(\texttt{accepted}\) requires \(J\)'s visibility, threshold, source and audit checks to pass on the bridged content. The rule preserves the exclusion of instrument-free acceptance in Subsubsection \hyperref[sn:2.4.3]{2.4.3}: the source acceptance and the bridge admissibility are the premises of the rule; they are not sufficient for \(\texttt{accepted}\) under \(J\), since \(J\)'s classifier decides the verdict on the bridged record.

\subsubsection{Future-Work Boundary}\label{sn:11.9.2}

A complete instrument-transfer classifier --- one that fixes, for every admissible \(B_{I\to J}\) and every claim type, the exact mapping from the source-side acceptance and the bridge defect record to the value \(\chi\in\mathrm{ClaimStatus}\) on the target side --- is not stated as a theorem of Foundations III. Subsubsection \hyperref[sn:11.9.1]{11.9.1} fixes the rule schema and the admissibility predicate; the per-claim-type classifier and its associated soundness theorem (the analogue of Theorem \hyperref[cl:thm:23]{23} for transferred claims) are deferred. This is consistent with the scope nonclaims of Subsection \hyperref[sn:2.4]{2.4} and with the deferred-claims register of Appendix \hyperref[sn:H]{H}: the calculus accepts only what is supported by an admissible bridge and a passing classifier verdict, and does not commit to a closed-form transfer-completeness statement.

The framework's default position is therefore the negative one: \(I\vdash\varphi\) does not imply \(J\vdash\varphi\). Positive transfer is recorded only by an admissible instrument-transfer bridge whose verdict under \(J\) is independently classified. Future-work strengthenings of Subsubsection \hyperref[sn:11.9.1]{11.9.1} would supply the per-claim-type completeness theorem and the corresponding countermodel atlas; both are outside the scope of the present paper.
 \section{Abstract Interpretation Model Family}\label{sn:12}

This section treats abstract interpretation \citep{CousotCousot1977,CousotCousot1979}, with its Galois connections \citep{Ore1944Galois}, as a finite model family for \(\mathbf{BirdInt}^{\mathrm{aud}}_{\mathrm{fin}}\) and proves the four model-family theorems that establish its scope. Subsection \hyperref[sn:12.1]{12.1} fixes the finite host data and the realized role fragment. Subsections \hyperref[sn:12.2]{12.2}--\hyperref[sn:12.4]{12.4} prove Theorems \hyperref[cl:thm:26]{26}, \hyperref[cl:thm:27]{27}, and \hyperref[cl:thm:28]{28}: closure as packaging, descent as sound transformer, and representability as fixedness. Subsection \hyperref[sn:12.5]{12.5} proves Model Theorem \hyperref[cl:thm:29]{29}, that a finite family of admissible abstract-interpretation hosts realizes the \(P_1/P_2/P_5/P_6\) fragment of the calculus, and \(P_1/P_2/P_4/P_5/P_6\) when its refinement records contain a strict refinement. Subsection \hyperref[sn:12.6]{12.6} records the abstract-interpretation nonclaims that fix the boundaries of this realization.

\emph{The model-realization convention.} A model-realization statement \(\Gamma;\mathcal T;I\vdash\mathcal M\models\mathcal S:\texttt{model\_realization}\) says that a finite realization object \(\mathcal M\), which carries host data, an instrument \(I\), and source, audit, visibility and nonclaim policies, realizes a declared finite fragment \(\mathcal S\) of the calculus. A finite family of assignments sends host data of \(\mathcal M\) to every component of \(\mathcal S\); each assignment is audited by \(I\); the source and visibility policies of \(\mathcal M\) meet those that \(\mathcal S\) requires; and the nonclaim register of \(\mathcal M\) contains the model-realization nonclaims. The verdict \(\texttt{model\_realization}\) is not the verdict \(\texttt{accepted}\) of a theorem, and it says nothing about fragments larger than \(\mathcal S\). Subsection \hyperref[sn:13.1]{13.1} gives the full definition; the proof of Model Theorem \hyperref[cl:thm:29]{29} checks its seven clauses: (1) every component assigned, (2) each assignment audited, (3) sources admitted, (4) required records visible, (5) model-realization nonclaims present, (6) witnesses and updates typed by \(\mathsf{Wit}\) and \(\mathsf{Upd}\), (7) every realized directed cell, pair observable or promotion bridge carries the status its classifier assigns (Subsection \hyperref[sn:5.6]{5.6}); role-channel statuses enter through clause 1.

The discipline of the section is the model-realization discipline of Subsection \hyperref[sn:13.1]{13.1}: abstract interpretation realizes a declared fragment of the calculus and not the whole calculus. In particular, the family additionally realizes \(P_3\) when two host closures fail to commute, and the section does not claim that sound abstraction is complete representation, that closure implies transformer soundness, or that the AI fragment is a foundation for the calculus.

\subsection{Abstract-Interpretation Host}\label{sn:12.1}

This subsection fixes the finite host data --- concrete and abstract domains, abstraction and concretization maps, the closure operator, sound transformers, and audit records --- and identifies the role fragment that the host realizes inside \(\mathbf{BirdInt}^{\mathrm{aud}}_{\mathrm{fin}}\).

\subsubsection{Host Data}\label{sn:12.1.1}

A finite abstract-interpretation host is a tuple

\[
H\;=\;(\,C,\;A,\;\alpha,\;\gamma,\;\rho,\;F,\;F^{\sharp},\;A_{\mathrm{AI}}\,)
\]

with the following components:

\begin{itemize}
\tightlist
\item
  \((C,\le_C)\) is a finite concrete domain: a finite set \(C\) together with a partial order \(\le_C\). Concreteness and finiteness are imposed throughout: every supremum, infimum, and fixed point used by the host is taken in a finite poset and is decidable on finite records.
\item
  \((A,\le_A)\) is a finite abstract domain: a finite set \(A\) together with a partial order \(\le_A\). In this section \(A\) always denotes the abstract domain; the audit record of the host is written \(A_{\mathrm{AI}}\).
\item
  \(\alpha:C\to A\) is a monotone abstraction map, and \(\gamma:A\to C\) is a monotone concretization map. The pair forms a Galois connection
  \[
  \alpha(c)\;\le_A\;a\quad\Longleftrightarrow\quad c\;\le_C\;\gamma(a),
  \]
  written \(\alpha\dashv\gamma\). Equivalently, \(\alpha\) and \(\gamma\) satisfy unit \(c\le_C\gamma\alpha(c)\) and counit \(\alpha\gamma(a)\le_A a\).
\item
  \(\rho:C\to C\) is the host closure operator, defined by \(\rho=\gamma\circ\alpha\). The fixedness set is \(\mathrm{Fix}(\rho)=\{c\in C:\rho(c)=c\}\).
\item
  \(F:C\to C\) is a monotone concrete transformer, and \(F^{\sharp}:A\to A\) is a monotone abstract transformer; transformer soundness is not assumed.
\item
  \(A_{\mathrm{AI}}\) is a finite audit record on the host data: it verifies the structural clauses above (\(C\) and \(A\) are finite ordered domains, \(\alpha\) and \(\gamma\) are monotone, \(\alpha\dashv\gamma\), \(\rho=\gamma\alpha\), and \(F\) and \(F^{\sharp}\) are monotone) and records the Boolean result of checking \(\alpha F\le_A F^\sharp\alpha\), which may be false. The audit record is finite and indexed by a host-level instrument \(I_{\mathrm{AI}}\).
\end{itemize}

These data form an \emph{AI pre-host}. An AI host is admissible, written \(\mathsf{AdmAIHost}(H)\), exactly when \(A_{\mathrm{AI}}\), with \(\mathsf{CheckRules}\) drawn from \(\mathsf{CheckRules}_{I_{\mathrm{AI}}}\), has \(\delta_6^{\mathrm{audit}}(A_{\mathrm{AI}})=\texttt{audit\_passes}\) (in particular its structural checks pass) and the recorded soundness Boolean is true, that is, \(\alpha F\le_A F^\sharp\alpha\) holds. Theorems \hyperref[cl:thm:26]{26}--\hyperref[cl:thm:28]{28} use only the pre-host clauses; soundness enters as the \(P_1\) witness of an admissible host.

The instrument-indexed claim form for any AI-host theorem in this section is

\[
\Gamma\,;\;\mathsf H_{\mathrm{AI}}\,;\;I_{\mathrm{AI}}\;\vdash\;\varphi\;:\;\chi,
\]

with \(\mathsf H_{\mathrm{AI}}\) the AI-host context recording \(H\) and \(I_{\mathrm{AI}}\) the instrument that audits the host.

\subsubsection{BirdInt Realized Fragment}\label{sn:12.1.2}

Every nonempty admissible finite family of abstract-interpretation hosts sharing one concrete domain \(C\) and one concrete transformer \(F\), equipped with the witness, source and instrument records constructed in the proof of Model Theorem \hyperref[cl:thm:29]{29} (Subsection \hyperref[sn:12.5]{12.5}), realizes the role fragment \(P_1/P_2/P_5/P_6\) of \(\mathbf{BirdInt}^{\mathrm{aud}}_{\mathrm{fin}}\); it also realizes \(P_4\), giving the fragment

\[
P_1\;/\;P_2\;/\;P_4\;/\;P_5\;/\;P_6,
\]

when its refinement records contain a pair with \(\rho_k\ne\rho_\ell\). The realization map is fixed as follows:

\begin{itemize}
\tightlist
\item
  \(P_5\) (packaging) is realized by the closure operator \(\rho=\gamma\alpha\). Theorem \hyperref[cl:thm:26]{26} of Subsection \hyperref[sn:12.2]{12.2} records that \(\rho\) is extensive, monotone, and idempotent, hence supplies the finite packaging/objecthood record for the AI host.
\item
  \(P_1\) (descent) is realized by the soundness inequality \(\alpha F\le_A F^{\sharp}\alpha\), equivalently \(F\gamma\le_C\gamma F^{\sharp}\). Theorem \hyperref[cl:thm:27]{27} of Subsection \hyperref[sn:12.3]{12.3} records that the equivalence is the descent-square soundness condition for the AI host.
\item
  \(P_2\) (representability) is realized by exact representability \(p\,\text{representable}\iff\exists a\in A:\gamma(a)=p\). Theorem \hyperref[cl:thm:28]{28} of Subsection \hyperref[sn:12.4]{12.4} records the equivalence with fixedness \(\rho(p)=p\) and with \(\mathrm{im}(\gamma)\).
\item
  \(P_4\) (refinement) is realized by audited refinement records between abstract domains within a finite family of hosts: a refinement record \(A_k\leadsto A_\ell\) is a finite audited bridge comparing the precision of two abstract domains within \(\mathsf H_{\mathrm{AI}}\), and it realizes \(P_4\) when \(\rho_k\ne\rho_\ell\).
\item
  \(P_6\) (audit) is realized by the host audit record \(A_{\mathrm{AI}}\) of Subsubsection \hyperref[sn:12.1.1]{12.1.1}, indexed by \(I_{\mathrm{AI}}\).
\end{itemize}

The role \(P_3\) (route mismatch) is not realized by the bare AI host. The route mismatch primitive of \(\mathbf{BirdInt}^{\mathrm{aud}}_{\mathrm{fin}}\) requires a route record and a route-mismatch defect, neither of which is supplied by a single Galois pair. The realized fragment of a family is \(P_1/P_2/P_5/P_6\), extended to \(P_1/P_2/P_4/P_5/P_6\) when the refinement records contain a strict refinement. The family additionally realizes \(P_3\) when two host closures fail to commute.

The fragment discipline of Subsection \hyperref[sn:13.1]{13.1} then applies: realization of the fragment \(P_1/P_2/P_4/P_5/P_6\) does not imply realization of the whole calculus, and the abstract-interpretation host is not a foundation for \(\mathbf{BirdInt}^{\mathrm{aud}}_{\mathrm{fin}}\) but a model family within it.

\subsection{Theorem 26: Closure as Packaging}\label{sn:12.2}

This subsection proves Theorem \hyperref[cl:thm:26]{26}: in any finite AI pre-host, the composite \(\rho=\gamma\alpha\) is a closure operator and so realizes the \(P_5\) packaging role of \(\mathbf{BirdInt}^{\mathrm{aud}}_{\mathrm{fin}}\).

Let \(H=(C,A,\alpha,\gamma,\rho,F,F^{\sharp},A_{\mathrm{AI}})\) be a finite AI pre-host of Subsubsection \hyperref[sn:12.1.1]{12.1.1}, not assumed sound, with Galois connection \(\alpha\dashv\gamma\) and \(\rho=\gamma\alpha\).

\begin{claimtheorem}{Theorem 26 (AIClosureAsPackaging)}\phantomsection\label{cl:thm:26}

The map \(\rho:C\to C\) is a closure operator on \((C,\le_C)\): it is extensive, monotone, and idempotent, that is,

\[
c\;\le_C\;\rho(c)\quad\text{(extensive)},
\]

\[
c\;\le_C\;c'\;\Longrightarrow\;\rho(c)\;\le_C\;\rho(c')\quad\text{(monotone)},
\]

\[
\rho^{2}\;=\;\rho\quad\text{(idempotent)}.
\]

In particular, \(\rho\) is a finite packaging operator on \(C\) in the sense required by the \(P_5\) role of \(\mathbf{BirdInt}^{\mathrm{aud}}_{\mathrm{fin}}\).

The three closure-operator clauses, with the same proof, hold for arbitrary posets and a Galois connection between them.

\end{claimtheorem}

\begin{proof}

The three closure-operator clauses follow directly from the Galois unit, the monotonicity of \(\alpha\) and \(\gamma\), and the Galois counit, applied to the finite ordered domains \((C,\le_C)\) and \((A,\le_A)\).

\emph{Extensivity.} The Galois unit asserts \(c\le_C\gamma\alpha(c)\) for every \(c\in C\). By definition \(\rho(c)=\gamma\alpha(c)\), so \(c\le_C\rho(c)\).

\emph{Monotonicity.} Suppose \(c\le_C c'\). Monotonicity of \(\alpha\) gives \(\alpha(c)\le_A\alpha(c')\). Monotonicity of \(\gamma\) then gives \(\gamma\alpha(c)\le_C\gamma\alpha(c')\), that is, \(\rho(c)\le_C\rho(c')\).

\emph{Idempotence.} Apply the Galois counit at \(a=\alpha(c)\): \(\alpha\gamma\alpha(c)\le_A\alpha(c)\). Monotonicity of \(\gamma\) gives \(\gamma\alpha\gamma\alpha(c)\le_C\gamma\alpha(c)\), that is, \(\rho^{2}(c)\le_C\rho(c)\). Conversely, extensivity applied at \(\rho(c)\) gives \(\rho(c)\le_C\rho(\rho(c))=\rho^{2}(c)\). Antisymmetry of \(\le_C\) on the finite poset \((C,\le_C)\) yields \(\rho^{2}(c)=\rho(c)\) for every \(c\in C\), hence \(\rho^{2}=\rho\).

Each clause is verified by a finite check on the finite host data \(H\) and recorded in the host audit record \(A_{\mathrm{AI}}\). The mathematical conclusion uses only the pre-host clauses.

The closure operator \(\rho\) is the host's \(P_5\) packaging witness: it supplies the finite extensive, monotone, idempotent map on \(C\) that the realization map of Subsubsection \hyperref[sn:12.1.2]{12.1.2} requires. Theorems \hyperref[cl:thm:27]{27} and \hyperref[cl:thm:28]{28} of Subsections \hyperref[sn:12.3]{12.3} and \hyperref[sn:12.4]{12.4} then state how the host's \(P_1\) and \(P_2\) roles are read off from \(\alpha\), \(\gamma\), \(F\), \(F^{\sharp}\), and \(\rho\).

\end{proof}

\subsection{Theorem 27: Descent as Sound Transformer}\label{sn:12.3}

This subsection proves Theorem \hyperref[cl:thm:27]{27}: in any finite AI pre-host, the abstract-side soundness inequality \(\alpha F\le_A F^{\sharp}\alpha\) is equivalent to the concrete-side soundness inequality \(F\gamma\le_C\gamma F^{\sharp}\). The equivalence gives a finite test for the \(P_1\) soundness condition; a \(P_1\) descent witness is supplied when either, hence both, inequalities hold.

Let \(H=(C,A,\alpha,\gamma,\rho,F,F^{\sharp},A_{\mathrm{AI}})\) be a finite AI pre-host of Subsubsection \hyperref[sn:12.1.1]{12.1.1}, not assumed sound, with Galois connection \(\alpha\dashv\gamma\) and monotone transformers \(F:C\to C\) and \(F^{\sharp}:A\to A\).

\begin{claimtheorem}{Theorem 27 (AIDescentAsSoundTransformer)}\phantomsection\label{cl:thm:27}

The two soundness inequalities

\[
\alpha F\;\le_A\;F^{\sharp}\alpha\qquad\text{and}\qquad F\gamma\;\le_C\;\gamma F^{\sharp}
\]

are equivalent: each holds iff the other does. Equivalently,

\[
\alpha F(c)\;\le_A\;F^{\sharp}\alpha(c)\;\;\text{for every}\;c\in C\quad\Longleftrightarrow\quad F\gamma(a)\;\le_C\;\gamma F^{\sharp}(a)\;\;\text{for every}\;a\in A.
\]

The equivalence gives a finite test for the \(P_1\) soundness condition; a \(P_1\) descent witness is supplied when either, hence both, inequalities hold.

Moreover, \(\alpha F=F^{\sharp}\alpha\) holds if and only if \(F\) descends through \(\alpha\) in the sense of Theorem \hyperref[cl:thm:6]{6} and \(F^{\sharp}|_{\mathrm{im}\,\alpha}\) is its descended update. The soundness inequality is thus the lax form of the descent square of Theorem \hyperref[cl:thm:6]{6}.

The soundness-inequality equivalence and the equality/descent characterization hold, with the same proof, for arbitrary posets, a Galois connection, and monotone transformers \(F,F^\sharp\).

\end{claimtheorem}

\begin{proof}

Each direction follows from the Galois adjunction together with monotonicity of \(\alpha\), \(\gamma\), \(F\), and \(F^{\sharp}\).

\emph{Forward direction.} Assume \(\alpha F\le_A F^{\sharp}\alpha\). Fix \(a\in A\); we show \(F\gamma(a)\le_C\gamma F^{\sharp}(a)\). By the Galois adjunction \(\alpha\dashv\gamma\), this is equivalent to \(\alpha F\gamma(a)\le_A F^{\sharp}(a)\). Apply the assumption at \(c=\gamma(a)\):

\[
\alpha F\gamma(a)\;\le_A\;F^{\sharp}\alpha\gamma(a).
\]

The Galois counit gives \(\alpha\gamma(a)\le_A a\), and monotonicity of \(F^{\sharp}\) gives \(F^{\sharp}\alpha\gamma(a)\le_A F^{\sharp}(a)\). Combining the two inequalities yields \(\alpha F\gamma(a)\le_A F^{\sharp}(a)\). Re-applying the adjunction returns \(F\gamma(a)\le_C\gamma F^{\sharp}(a)\), as required.

\emph{Reverse direction.} Assume \(F\gamma\le_C\gamma F^{\sharp}\). Fix \(c\in C\); we show \(\alpha F(c)\le_A F^{\sharp}\alpha(c)\). By the Galois adjunction, this is equivalent to \(F(c)\le_C\gamma F^{\sharp}\alpha(c)\). The Galois unit gives \(c\le_C\gamma\alpha(c)\), and monotonicity of \(F\) gives \(F(c)\le_C F\gamma\alpha(c)\). Apply the assumption at \(a=\alpha(c)\):

\[
F\gamma\alpha(c)\;\le_C\;\gamma F^{\sharp}\alpha(c).
\]

Combining the two inequalities yields \(F(c)\le_C\gamma F^{\sharp}\alpha(c)\). Re-applying the adjunction returns \(\alpha F(c)\le_A F^{\sharp}\alpha(c)\), as required.

\emph{Equality and descent.} If \(\alpha F=F^{\sharp}\alpha\) and \(\alpha(c)=\alpha(c')\), then \(\alpha F(c)=F^{\sharp}\alpha(c)=F^{\sharp}\alpha(c')=\alpha F(c')\), so \(F\) descends through \(\alpha\) by Theorem \hyperref[cl:thm:6]{6}; since \(F^{\sharp}(\alpha(c))=\alpha(F(c))\), the map \(F^{\sharp}\) sends \(\mathrm{im}(\alpha)\) into itself, and its restriction satisfies the descent equation, so it is the descended update. Conversely, if \(F\) descends with descended update \(D:\mathrm{im}(\alpha)\to\mathrm{im}(\alpha)\) and \(F^{\sharp}\) agrees with \(D\) on \(\mathrm{im}(\alpha)\), then \(F^{\sharp}\alpha(c)=D(\alpha(c))=\alpha F(c)\) for every \(c\).

Each step uses only the Galois unit, the Galois counit, monotonicity of \(\alpha\), \(\gamma\), \(F\), \(F^{\sharp}\), and antisymmetry/transitivity of \(\le_C\) and \(\le_A\) on the finite posets. Every check is finite and is recorded in the host audit record \(A_{\mathrm{AI}}\). The mathematical conclusion uses only the pre-host clauses.

The equivalence gives a finite test for the \(P_1\) soundness condition, and a \(P_1\) descent witness is supplied when either, hence both, inequalities hold, in either form: the abstract-side \(\alpha F\le_A F^{\sharp}\alpha\) is the upper-soundness inequality used when reasoning about \(F^{\sharp}\) as an over-approximation of \(F\) on the abstract side, and the concrete-side \(F\gamma\le_C\gamma F^{\sharp}\) is the descent inequality used when reasoning about the concretized abstract step as a finite over-approximation of the concrete step. The two inequalities are interchangeable under the Galois adjunction, and the host carries either form indifferently as its \(P_1\) record.

\end{proof}

\subsection{Theorem 28: Representability as Fixedness}\label{sn:12.4}

This subsection proves Theorem \hyperref[cl:thm:28]{28}: in any finite AI pre-host, exact representability of a concrete element coincides with fixedness under the closure operator \(\rho\) and with membership in the image of the concretization map \(\gamma\). The equivalence supplies the host's \(P_2\) representability witness.

Let \(H=(C,A,\alpha,\gamma,\rho,F,F^{\sharp},A_{\mathrm{AI}})\) be a finite AI pre-host of Subsubsection \hyperref[sn:12.1.1]{12.1.1}, not assumed sound, with Galois connection \(\alpha\dashv\gamma\) and \(\rho=\gamma\alpha\). For \(p\in C\), define exact representability by

\[
p\;\text{representable}\quad\Longleftrightarrow\quad\exists\,a\in A:\;\gamma(a)=p.
\]

\begin{claimtheorem}{Theorem 28 (AIRepresentabilityAsFixedness)}\phantomsection\label{cl:thm:28}

For every \(p\in C\),

\[
p\;\text{representable}\quad\Longleftrightarrow\quad\rho(p)=p.
\]

Equivalently,

\[
\mathrm{Rep}_{\gamma}(C)\;=\;\mathrm{Fix}(\rho)\;=\;\mathrm{im}(\gamma),
\]

where \(\mathrm{Rep}_{\gamma}(C)=\{p\in C:p\text{ representable}\}\), \(\mathrm{Fix}(\rho)=\{p\in C:\rho(p)=p\}\), and \(\mathrm{im}(\gamma)=\{p\in C:\exists\,a\in A,\;\gamma(a)=p\}\).

The theorem identifies three finite, equally decidable subsets of \(C\) --- the representable elements, the \(\rho\)-fixed points, and the image of \(\gamma\) --- and so supplies the host's \(P_2\) representability witness as a finite, audit-checkable predicate on \(C\).

The equivalences, with the same proof, hold for arbitrary posets \((C,\le_C)\), \((A,\le_A)\) and a Galois connection between them; finiteness only makes the three subsets decidable.

\end{claimtheorem}

\begin{proof}

The two implications follow from the Galois adjunction \(\alpha\dashv\gamma\) and the definition \(\rho=\gamma\alpha\).

\emph{Forward direction.} Suppose \(p\) is representable, so there exists \(a\in A\) with \(\gamma(a)=p\). Then

\[
\rho(p)\;=\;\gamma\alpha\gamma(a).
\]

The standard Galois-adjunction identity \(\gamma\alpha\gamma=\gamma\) follows from the unit and counit. The counit gives \(\alpha\gamma(a)\le_A a\); applying the monotone map \(\gamma\) gives

\[
\gamma\alpha\gamma(a)\;\le_C\;\gamma(a).
\]

The unit at the concrete element \(\gamma(a)\) gives

\[
\gamma(a)\;\le_C\;\gamma\alpha\gamma(a).
\]

By antisymmetry, \(\gamma\alpha\gamma(a)=\gamma(a)\). Hence \(\rho(p)=\gamma\alpha\gamma(a)=\gamma(a)=p\).

\emph{Reverse direction.} Suppose \(\rho(p)=p\), so \(\gamma\alpha(p)=p\). Set \(a=\alpha(p)\in A\). Then \(\gamma(a)=\gamma\alpha(p)=p\), so \(p\) is representable.

The two directions together yield \(\mathrm{Rep}_{\gamma}(C)=\mathrm{Fix}(\rho)\). The remaining equality \(\mathrm{Fix}(\rho)=\mathrm{im}(\gamma)\) is now immediate: if \(p=\rho(p)\), then \(p=\gamma\alpha(p)\in\mathrm{im}(\gamma)\); conversely, if \(p=\gamma(a)\) for some \(a\in A\), then \(\rho(p)=\gamma\alpha\gamma(a)=\gamma(a)=p\), so \(p\in\mathrm{Fix}(\rho)\).

Each step uses only the Galois unit and counit and the identity \(\gamma\alpha\gamma=\gamma\) that is read off from them on the finite posets \((C,\le_C)\) and \((A,\le_A)\). Every check is finite and is recorded in the host audit record \(A_{\mathrm{AI}}\). The mathematical conclusion uses only the pre-host clauses.

The fixed set \(\mathrm{Fix}(\rho)\) is the host's \(P_2\) representability witness. Together with Theorem \hyperref[cl:thm:26]{26} (which makes \(\rho\) a packaging operator) and Theorem \hyperref[cl:thm:27]{27} (which makes the soundness inequalities equivalent), Theorem \hyperref[cl:thm:28]{28} closes the local circle: representable elements are exactly the \(\rho\)-fixed elements, which are exactly the \(\gamma\)-images, and packaging, descent, and representability are read off from \(\rho\), \(F^{\sharp}\), and \(\gamma\) via the Galois adjunction.

\end{proof}

\subsection{Model Theorem 29: AI Model-Family Theorem}\label{sn:12.5}

This subsection proves Model Theorem \hyperref[cl:thm:29]{29}: every nonempty admissible finite family of abstract-interpretation hosts sharing one concrete domain \(C\) and one concrete transformer \(F\) realizes \(P_1/P_2/P_5/P_6\), and it also realizes \(P_4\) when its refinement records contain a pair with \(\rho_k\ne\rho_\ell\), giving the \(P_1/P_2/P_4/P_5/P_6\) fragment of \(\mathbf{BirdInt}^{\mathrm{aud}}_{\mathrm{fin}}\). The theorem assembles the per-host realizations of \(P_5\), \(P_1\), and \(P_2\) (Theorems \hyperref[cl:thm:26]{26}, \hyperref[cl:thm:27]{27}, \hyperref[cl:thm:28]{28}) with the audited refinement records between hosts and the per-host audit records. The theorem is model-grade rather than theorem-grade: its acceptance verdict is \(\texttt{model\_realization}\), not \(\texttt{accepted}\), in the sense fixed in Subsection \hyperref[sn:13.1]{13.1}.

Let

\[
\mathsf{AIFam}_{\mathrm{fin}}\;=\;\{\,H_k\,\}_{k\in\mathcal J}
\]

be a finite family, indexed by a finite set \(\mathcal J\), of admissible abstract-interpretation hosts in the sense of Subsubsection \hyperref[sn:12.1.1]{12.1.1}:

\[
H_k\;=\;(\,C,\;A_k,\;\alpha_k,\;\gamma_k,\;\rho_k,\;F,\;F_k^{\sharp},\;A_{\mathrm{AI},k}\,),
\]

with \(\alpha_k\dashv\gamma_k\) and \(\rho_k=\gamma_k\alpha_k\) for each \(k\in\mathcal J\), and with \(\mathsf{AdmAIHost}(H_k)\) holding for each \(k\). For \(k,\ell\in\mathcal J\) with \(\rho_k\le_C\rho_\ell\) pointwise, equivalently \(\mathrm{im}(\gamma_\ell)\subseteq\mathrm{im}(\gamma_k)\) (for closure operators, \(\rho_k\le\rho_\ell\) iff \(\mathrm{Fix}(\rho_\ell)\subseteq\mathrm{Fix}(\rho_k)\), and \(\mathrm{Fix}(\rho)=\mathrm{im}(\gamma)\) by Theorem \hyperref[cl:thm:28]{28}), let \(R_{k\to\ell}\) be the audited finite certificate of the inequalities \(\rho_k(c)\le_C\rho_\ell(c)\), \(c\in C\): the abstract domain \(A_k\) is at least as precise as \(A_\ell\). Let \(\mathsf{Refine}=\{(k,\ell,R_{k\to\ell}):\rho_k\le_C\rho_\ell\}\), and let \(\mathcal J\ne\varnothing\). The \(P_4\) assignment below uses a pair \((k,\ell,R_{k\to\ell})\in\mathsf{Refine}\) with \(\rho_k\neq\rho_\ell\), a strict refinement; the trivial records \((k,k,R_{k\to k})\) do not count: the activation threshold \(\Theta_4^{\mathrm{act}}\) of the \(P_4\) role-channel record checks that \(\rho_k(c)\neq\rho_\ell(c)\) for some \(c\in C\), so a record with \(\rho_k=\rho_\ell\), including every trivial record, receives \(\texttt{below\_threshold}\) and supplies no \(P_4\) evidence. If \(C\) has a top element and some \(\rho_k\) is not constantly \(\top_C\), adjoining the one-point host with \(\alpha(c)=*\), \(\gamma(*)=\top_C\), and \(F^\sharp(*)=*\) gives such a pair.

For each \(k\), the realization uses finite tagged witness records: a \(P_5\) package-map record carrying \(\rho_k\); a \(P_1\) descent record carrying the two compared maps and the verified inequality; a \(P_2\) representability record carrying the finite fixed-point comparison; a \(P_6\) audit record carrying \(A_{\mathrm{AI},k}\); and, for each selected \((k,\ell)\), a \(P_4\) refinement record carrying \(R_{k\to\ell}\). For each pair of hosts whose closures fail to commute, choose a point \(c_{k\ell}\in\Delta_3^{\mathrm{comp}}(\rho_k,\rho_\ell)\) by finite enumeration, declare the corresponding \(\mathsf{RouteMismatchWit}\) record and its defect record visible under \(I_{\mathrm{AI}}\), reference the admitted host sources, and add a family audit recomputing both composites at every \(c\in C\). Their named signatures belong, respectively, to \(\mathsf{Wit}(P_5)\), \(\mathsf{Wit}(P_1)\), \(\mathsf{Wit}(P_2)\), \(\mathsf{Wit}(P_6)\), and \(\mathsf{Wit}(P_4)\). Each selected tagged witness \(R_i\) is placed in \(\Gamma\), \(\mathsf{Scope}_{I_{\mathrm{AI}}}\) and \(\mathsf{Visible}_{I_{\mathrm{AI}}}\), and supplies a role-channel record \((\mathrm{Chan}_i,R_i,\Theta_i^{\mathrm{act}},A_{R_i},V_{R_i},\mathrm{false})\) of Subsubsection \hyperref[sn:5.6.1]{5.6.1}, judged as \(\Gamma;\mathcal T_{\mathrm{AI},k};I_{\mathrm{AI}}\vdash\mathrm{Chan}_i:\rho_i\) under the canonical AI package \(\mathcal T_{\mathrm{AI},k}=(C,\alpha_k,\Sigma_{\alpha_k},\rho_k,\mathcal A_{\mathrm{AI},k})\) of host \(H_k\), with \(\Sigma_{\alpha_k}=\{\alpha_k^{-1}(B):B\subseteq A_k\}\) and \(\mathcal A_{\mathrm{AI},k}\) assembled from \(A_{\mathrm{AI},k}\) (Subsubsection \hyperref[sn:D.3.2]{D.3.2}), hosted on \(\mathsf H_{\mathrm{AI}}\). Its passing audit checks its carried witness data: the refinement audit for \(P_4\), the family composite audit for \(P_3\), and finite audits of the selected per-host records otherwise. The activation thresholds are: for \(P_1\), \(\alpha_kF(c)\le F_k^{\sharp}\alpha_k(c)\) rechecked at every \(c\in C\); for \(P_2\), the recorded representable set equals the recomputed \(\mathrm{Fix}(\rho_k)\) and \(\mathrm{im}(\gamma_k)\); for \(P_5\), extensivity, monotonicity and idempotence of \(\rho_k\) rechecked on \(C\); for \(P_6\), \(\delta_6^{\mathrm{audit}}(A_{\mathrm{AI},k})=\texttt{audit\_passes}\); for \(P_3\), \(\rho_k\rho_\ell(c_{k\ell})\neq\rho_\ell\rho_k(c_{k\ell})\); for \(P_4\), \(\Theta_4^{\mathrm{act}}\) above. By Theorems \hyperref[cl:thm:26]{26} and \hyperref[cl:thm:28]{28} the \(P_2\) and \(P_5\) thresholds pass for every admissible host. The displayed role map denotes these records through their mathematical payloads. The family record is \((\mathsf{AIFam}_{\mathrm{fin}},\mathsf{Refine},\mathcal N_{\mathrm{AI}})\), where \(\mathcal N_{\mathrm{AI}}\) contains the nonclaims of Subsection \hyperref[sn:12.6]{12.6} and of Subsubsection \hyperref[sn:13.6.2]{13.6.2}; \(\mathsf{Visible}_{I_{\mathrm{AI}}}\) contains every host record and every \(R_{k\to\ell}\), and for each \(k\) a source record \(\mathrm{src}_{\mathrm{AI},k}=(\mathrm{id}_k,\texttt{committed\_state},\text{the recorded host data of }H_k,\texttt{valid})\) is audited by \(A_{\mathrm{AI},k}\); \(\mathsf{SourcePolicy}_{I_{\mathrm{AI}}}\) admits these source records.

\begin{claimtheorem}{Model Theorem 29 (AIModelFamily)}\phantomsection\label{cl:thm:29}

Fix \(k\in\mathcal J\). The four assignments displayed below for \(P_5\), \(P_1\), \(P_2\) and \(P_6\), taken at this \(k\), realize the role fragment \(P_1/P_2/P_5/P_6\) of \(\mathbf{BirdInt}^{\mathrm{aud}}_{\mathrm{fin}}\). If moreover \(\mathsf{Refine}\) contains a triple \((k,\ell,R_{k\to\ell})\) with this first index and \(\rho_k\neq\rho_\ell\), the five displayed assignments, with \(P_4\) taken at that triple, realize \(P_1/P_2/P_4/P_5/P_6\) of \(\mathbf{BirdInt}^{\mathrm{aud}}_{\mathrm{fin}}\). Such a \(k\) exists whenever \(\mathsf{Refine}\) contains any triple with \(\rho_k\ne\rho_\ell\). The assignments are:

\[
P_5\;\mapsto\;\rho_k,\qquad P_1\;\mapsto\;\alpha_k F\;\le_{A_k}\;F_k^{\sharp}\alpha_k,\qquad P_2\;\mapsto\;\bigl(p\;\text{representable}\;\Longleftrightarrow\;\rho_k(p)=p\bigr),
\]

\[
P_4\;\mapsto\;A_k\leadsto A_\ell\;\text{via}\;R_{k\to\ell},\qquad P_6\;\mapsto\;A_{\mathrm{AI},k}.
\]

In schematic form, when \(\mathsf{Refine}\) contains a pair with \(\rho_k\neq\rho_\ell\),

\[
\mathsf{AIFam}_{\mathrm{fin}}\;\models\;\mathbf{BirdInt}^{\mathrm{aud}}_{\mathrm{fin}}{}^{P_1/P_2/P_4/P_5/P_6}.
\]

Here the \(P_4\) assignment is taken at such a pair. If \(C\) has a top element and some \(\rho_k\) is not constantly \(\top_C\), then adjoining the one-point host gives a family whose \(\mathsf{Refine}\) contains such a pair, so that family realizes \(P_1/P_2/P_4/P_5/P_6\): \(A_k\) is strictly more precise than the adjoined one-point host. If \(\rho_k\rho_\ell\neq\rho_\ell\rho_k\) for some \(\ell\), the assignments at \(k\) together with \(P_3\mapsto(\rho_k,\rho_\ell,c_{k\ell})\), the route-mismatch record with \(c_{k\ell}\in\Delta_3^{\mathrm{comp}}(\rho_k,\rho_\ell)\) (the defect of Theorem \hyperref[cl:thm:10]{10}), realize \(P_1/P_2/P_3/P_5/P_6\), and all of \(P_1,\ldots,P_6\) when \(\mathsf{Refine}\) also contains \((k,\ell',R_{k\to\ell'})\) with \(\rho_k\neq\rho_{\ell'}\); such \(k,\ell\) are incomparable, since comparable closures commute.

The realization also holds for host-dependent monotone concrete transformers \(F_k\), with \(P_1\mapsto\alpha_kF_k\le_{A_k}F_k^{\sharp}\alpha_k\).

\end{claimtheorem}

\begin{proof}

The proof verifies, host by host and bridge by bridge, that the components required by the fragment are present, finite, and audited.

\emph{\(P_5\) packaging.} For each \(k\in\mathcal J\), Theorem \hyperref[cl:thm:26]{26} of Subsection \hyperref[sn:12.2]{12.2} applied to \(H_k\) asserts that \(\rho_k=\gamma_k\alpha_k\) is a finite closure operator on \((C,\le_C)\). Hence the assignment \(P_5\mapsto\rho_k\) supplies a per-host packaging witness recorded in \(A_{\mathrm{AI},k}\).

\emph{\(P_1\) descent.} For each \(k\in\mathcal J\), \(\mathsf{AdmAIHost}(H_k)\) gives \(\alpha_k F\le_{A_k}F_k^{\sharp}\alpha_k\), and by Theorem \hyperref[cl:thm:27]{27} of Subsection \hyperref[sn:12.3]{12.3} so does \(F\gamma_k\le_C\gamma_k F_k^{\sharp}\); this is the host's \(P_1\) witness in either form. The assignment \(P_1\mapsto(\alpha_k F\le_{A_k}F_k^{\sharp}\alpha_k)\) records the abstract-side form on \(H_k\), recorded in \(A_{\mathrm{AI},k}\).

\emph{\(P_2\) representability.} For each \(k\in\mathcal J\), Theorem \hyperref[cl:thm:28]{28} of Subsection \hyperref[sn:12.4]{12.4} applied to \(H_k\) asserts \(p\,\text{representable}\iff\rho_k(p)=p\) for every \(p\in C\), with \(\mathrm{Rep}_{\gamma_k}(C)=\mathrm{Fix}(\rho_k)=\mathrm{im}(\gamma_k)\). The assignment \(P_2\mapsto(p\,\text{representable}\iff\rho_k(p)=p)\) records the host's \(P_2\) representability witness, recorded in \(A_{\mathrm{AI},k}\).

\emph{\(P_4\) refinement.} Each refinement record \(R_{k\to\ell}\in\mathsf{Refine}\) is a finite audited bridge of the form fixed in Subsubsection \hyperref[sn:12.1.2]{12.1.2}: it certifies \(\rho_k(c)\le_C\rho_\ell(c)\) for every \(c\) in the same concrete domain \(C\), with audit conditions internal to \(\mathsf H_{\mathrm{AI}}\) and audit record contributing to the family-level audit aggregation. The assignment \(P_4\mapsto(A_k\leadsto A_\ell)\) thus supplies a \(P_4\) refinement witness at every recorded pair with \(\rho_k\neq\rho_\ell\), where it records a strict precision comparison. For the one-point host, \(\rho(c)=\top_C\) for every \(c\), so \(\rho_k\le_C\rho\) pointwise and \(\rho_k\neq\rho\) when \(\rho_k\) is not constantly \(\top_C\). Without such a pair no \(P_4\) assignment is made, and the other four assignments realize \(P_1/P_2/P_5/P_6\).

\emph{\(P_6\) audit.} Each per-host audit record \(A_{\mathrm{AI},k}\) supplies the \(P_6\) audit witness for \(H_k\), indexed by the family-level instrument \(I_{\mathrm{AI}}\). The aggregation \(\bigcup_{k\in\mathcal J}A_{\mathrm{AI},k}\) together with the bridge audit records of \(\mathsf{Refine}\) is the family-level audit witness.

\emph{\(P_3\) route mismatch.} If \(\rho_k\rho_\ell\neq\rho_\ell\rho_k\) for hosts \(k,\ell\), then \(\rho_k,\rho_\ell\) are idempotent completions on the finite carrier \(C\) (Theorem \hyperref[cl:thm:26]{26}), and \(\Delta_3^{\mathrm{comp}}(\rho_k,\rho_\ell)=\{c:\rho_k\rho_\ell(c)\neq\rho_\ell\rho_k(c)\}\) is nonempty; the record \((\rho_k,\rho_\ell,c_{k\ell})\), where \(c_{k\ell}\in\Delta_3^{\mathrm{comp}}(\rho_k,\rho_\ell)\), is a route-mismatch record in \(\mathsf{Wit}(P_3)\), as \(W_3\) is in the running example of Subsection \hyperref[sn:3.8]{3.8}, audited by the family audit of the construction. If \(\rho_k\le_C\rho_\ell\) pointwise, then \(\rho_\ell\rho_k=\rho_\ell\) and \(\rho_k\rho_\ell=\rho_\ell\) (since \(\rho_\ell(c)\le\rho_k\rho_\ell(c)\le\rho_\ell\rho_\ell(c)=\rho_\ell(c)\), and \(\rho_\ell(c)\le\rho_\ell\rho_k(c)\le\rho_\ell\rho_\ell(c)=\rho_\ell(c)\) by monotonicity and \(\rho_k\le\rho_\ell\)), so the two closures commute; hence the pair is incomparable. The noncommuting closures \(E_1,E_2\) of Theorem \hyperref[cl:thm:9]{9} arise in this way from the hosts \(A_i=\mathrm{Fix}(E_i)\), \(\alpha_i=E_i\), \(\gamma_i\) the inclusion, \(F=\mathrm{id}\) and \(F_i^\sharp=\mathrm{id}\).

\emph{Closure of the realization map.} The five assignments above are typed: each maps a role of the declared fragment to a finite, audited datum on \(\mathsf{AIFam}_{\mathrm{fin}}\) and \(\mathsf{Refine}\), recorded as a witness of that role, so clause 6 of Subsubsection \hyperref[sn:13.1.1]{13.1.1} holds. Clause 1 holds because each role-channel record \((\mathrm{Chan}_i,R_i,\Theta_i^{\mathrm{act}},A_{R_i},V_{R_i},\mathrm{false})\) built above receives \(\texttt{active\_projection}\): \(R_i\) is in scope on \(\mathsf H_{\mathrm{AI}}\) and present, \(\mathsf{collapse}_i\) is false, \(\Theta_i^{\mathrm{act}}\) passes (by \(\mathsf{AdmAIHost}\) for \(P_1\), by Theorems \hyperref[cl:thm:28]{28} and \hyperref[cl:thm:26]{26} for \(P_2\) and \(P_5\), because \(A_{\mathrm{AI},k}\) passes for \(P_6\), because \(\rho_k\ne\rho_\ell\) for \(P_4\), and because \(c_{k\ell}\in\Delta_3^{\mathrm{comp}}\) for \(P_3\)), and \(R_i\) is visible with a passing audit; clause 2 because each datum is audited in \(A_{\mathrm{AI},k}\) or in the audit of its refinement record; clause 3 because \(\mathsf{SourcePolicy}_{I_{\mathrm{AI}}}\) admits the audited host source records; clause 4 because every host record and refinement record is visible; and clause 5 because \(\mathcal N_{\mathrm{AI}}\) contains the nonclaims of Subsubsection \hyperref[sn:13.6.2]{13.6.2}. Clause 7 is vacuous for this realization, since its refinement records are not declared promotion bridges and it records no directed cell or pair observable. The family additionally realizes \(P_3\) when two host closures fail to commute, by the assignment of the \(P_3\) paragraph above. The instrument \(I_{\mathrm{AI}}\) records the model-realization judgment, which holds when the realization map satisfies clauses 1--7 of Subsubsection \hyperref[sn:13.1.1]{13.1.1}, once the per-host admissibility, the per-host theorems, the bridge admissibility, and the family-level audit records have all been recorded.

The verdict is therefore \(\texttt{model\_realization}\), in the sense fixed in Subsection \hyperref[sn:13.1]{13.1}: the family realizes the declared fragment, and not the whole calculus.

The model-realization verdict is the model-grade analogue of the \(\texttt{accepted}\) verdict of Theorems \hyperref[cl:thm:26]{26}, \hyperref[cl:thm:27]{27}, and \hyperref[cl:thm:28]{28}: it records that the host family supplies a finite, audited realization of the declared role fragment, while no realization datum is supplied for the unrealized role \(P_3\) on a bare host. The nonclaims of Subsection \hyperref[sn:12.6]{12.6} record the boundaries of this realization: the family supplies no cell, pair or promotion realization, sound abstraction is not complete representation, and closure does not imply transformer soundness.

For the last clause, the one-point host has \(A=\{*\}\), \(\alpha(c)=*\) and \(\gamma(*)=\top_C\); this is a Galois connection because \(c\le_C\top_C\) for every \(c\), and \(F^\sharp(*)=*\) satisfies the soundness inequality trivially. Equip the one-point host with the finite audit record checking its order, maps, adjunction, closure and soundness; it then satisfies \(\mathsf{AdmAIHost}\). Its closure \(\rho\) is constantly \(\top_C\), so \(\rho_k\le_C\rho\) for every \(k\), with \(\rho_k\neq\rho\) for the \(\rho_k\) that is not constantly \(\top_C\).

\emph{Host-dependent transformers.} The assignments at \(k\), the certificates \(R_{k\to\ell}\) and the defects \(\Delta_3^{\mathrm{comp}}(\rho_k,\rho_\ell)\) read only \(C\) and the data of the hosts involved; the one-point host satisfies \(\alpha F_k(c)=*=F^\sharp\alpha(c)\) for every \(F_k\).

\end{proof}

\subsection{AI Nonclaims}\label{sn:12.6}

This subsection records the abstract-interpretation nonclaims that fix the boundaries of the realization established in Subsections \hyperref[sn:12.2]{12.2}--\hyperref[sn:12.5]{12.5}. Each nonclaim is a scoped negative statement attached to the AI model family: it is a record of what the family does not establish, in the sense of the nonclaim discipline of Subsection \hyperref[sn:2.4]{2.4}. The nonclaims here support the global nonclaim discipline of Section \hyperref[sn:2]{2} and record the three AI-specific exclusions: no AI totality, no soundness-as-completeness, and no closure-implies-soundness.

\subsubsection{No AI Totality}\label{sn:12.6.1}

Abstract interpretation does not realize the whole calculus \(\mathbf{BirdInt}^{\mathrm{aud}}_{\mathrm{fin}}\): a family may realize all six role channels (Model Theorem \hyperref[cl:thm:29]{29}), but no directed cell, pair observable or promotion bridge. Theorem \hyperref[cl:thm:29]{29} of Subsection \hyperref[sn:12.5]{12.5} asserts the realization of the fragment \(P_1/P_2/P_5/P_6\), extended by \(P_4\) when the refinement records contain a strict refinement and by \(P_3\) when two host closures fail to commute. The bare AI host of Subsubsection \hyperref[sn:12.1.1]{12.1.1} does not supply a route record or a route-mismatch defect, and so does not realize the route-mismatch primitive \(P_3\) described in Subsubsection \hyperref[sn:2.2.2]{2.2.2}. A single bare host supplies no \(P_3\) data; a family supplies \(P_3\) data when two of its closures fail to commute (Model Theorem \hyperref[cl:thm:29]{29}). The realization map of Model Theorem \hyperref[cl:thm:29]{29} assigns no directed cell, pair observable or promotion bridge, so it does not realize the whole calculus; such extensions are not asserted here.

In particular, the family-level claim

\[
\Gamma\,;\;\mathsf H_{\mathrm{AI}}\,;\;I_{\mathrm{AI}}\;\vdash\;\mathsf{AIFam}_{\mathrm{fin}}\;\models\;\mathbf{BirdInt}^{\mathrm{aud}}_{\mathrm{fin}}\;:\;\chi
\]

does not receive verdict \(\texttt{model\_realization}\) for the whole calculus on the bare host. The verdict \(\texttt{model\_realization}\) is assigned only to the declared fragment: \(P_1/P_2/P_5/P_6\), and \(P_1/P_2/P_4/P_5/P_6\) when the refinement records contain a strict refinement. The family additionally realizes \(P_3\) when two host closures fail to commute.

\subsubsection{No Soundness-as-Completeness}\label{sn:12.6.2}

Sound abstraction does not imply complete representation. The \(P_1\) descent witness of Theorem \hyperref[cl:thm:27]{27} records the inequality \(\alpha F\le_A F^{\sharp}\alpha\) (equivalently \(F\gamma\le_C\gamma F^{\sharp}\)): the abstract transformer \(F^{\sharp}\) over-approximates the concrete transformer \(F\). The corresponding equality \(\alpha F=F^{\sharp}\alpha\), which would assert that \(F^{\sharp}\) is a complete (exact) representation of \(F\) through the Galois package \((\alpha,\gamma,\rho)\), is not implied by Theorem \hyperref[cl:thm:27]{27} and is not part of the host's \(P_1\) realization datum.

The formal distinction is the inequality/equality distinction itself: \(\alpha F\le_A F^{\sharp}\alpha\) is the soundness condition, while \(\alpha F=F^{\sharp}\alpha\) is the exactness condition. The host therefore records the descent inequality, not the descent equality, and the AI fragment does not commit to completeness.

\subsubsection{No Closure-Implies-Soundness}\label{sn:12.6.3}

The closure-operator status of \(\rho=\gamma\alpha\) from Theorem \hyperref[cl:thm:26]{26} does not imply the soundness of the abstract transformer \(F^{\sharp}\) from Theorem \hyperref[cl:thm:27]{27}. The Galois adjunction \(\alpha\dashv\gamma\) is sufficient to establish that \(\rho\) is extensive, monotone, and idempotent (Theorem \hyperref[cl:thm:26]{26}), and in particular this establishes the \(P_5\) packaging role of the host. It does not, by itself, establish \(\alpha F\le_A F^{\sharp}\alpha\), which is an additional, separate condition on the pair \((F,F^{\sharp})\).

The host records the closure operator \(\rho\) and the soundness inequality \(\alpha F\le_A F^{\sharp}\alpha\) as independent data. Theorem \hyperref[cl:thm:26]{26} commits only to the closure-operator properties of \(\rho\), and Theorem \hyperref[cl:thm:27]{27} commits only to the equivalence of the two soundness inequalities; neither theorem implies the other. The audit record \(A_{\mathrm{AI}}\) of Subsubsection \hyperref[sn:12.1.1]{12.1.1} records both checks separately, and the realization map of Subsubsection \hyperref[sn:12.1.2]{12.1.2} maps \(P_5\) to \(\rho\) and \(P_1\) to the soundness inequality, not from one to the other.

The three nonclaims above together fix the local boundary of the AI model family within \(\mathbf{BirdInt}^{\mathrm{aud}}_{\mathrm{fin}}\): the family does not realize the whole calculus, sound abstraction is not exact representation, and the packaging role does not imply the descent role. The model-realization verdict of Theorem \hyperref[cl:thm:29]{29} is therefore scoped: it records realization of the fragment \(P_1/P_2/P_5/P_6\), extended by \(P_4\) when the refinement records contain a strict refinement and by \(P_3\) when two host closures fail to commute, and is silent on the rest.
 \section{Model Realizations}\label{sn:13}

This section presents the four model realizations of \(\mathbf{BirdInt}^{\mathrm{aud}}_{\mathrm{fin}}\): PICA, Cantor strict-bridge, the abstract-interpretation host of Section \hyperref[sn:12]{12}, and the graph/cohomology drive support of Section \hyperref[sn:8]{8}. Subsection \hyperref[sn:13.1]{13.1} fixes the model-realization convention and the fragment discipline that governs every model in the family. Subsections \hyperref[sn:13.2]{13.2}, \hyperref[sn:13.3]{13.3}, and \hyperref[sn:13.4]{13.4} give the PICA, Cantor, and graph/cohomology realizations and their model theorems (Theorems \hyperref[cl:thm:33]{33} and \hyperref[cl:thm:34]{34} plus the model-grade reading of Section \hyperref[sn:8]{8}). Subsection \hyperref[sn:13.5]{13.5} organizes the four realizations into a combined model-realization matrix. Subsection \hyperref[sn:13.6]{13.6} records the cross-model role-coverage observation and the corresponding nonclaim.

The four model families are located as follows.

\begingroup\small
\begin{longtable}{@{}>{\raggedright\arraybackslash}p{0.24\linewidth}>{\raggedright\arraybackslash}p{0.22\linewidth}>{\raggedright\arraybackslash}p{0.32\linewidth}>{\raggedright\arraybackslash}p{0.14\linewidth}@{}}
\toprule
Model family & Host & Results & Sheet \\
\midrule
PICA & Subsection \hyperref[sn:13.2]{13.2} & Model Theorem \hyperref[cl:thm:33]{33} & Appendix \hyperref[sn:D.1]{D.1} \\
Cantor strict bridge & Subsection \hyperref[sn:13.3]{13.3} & Model Theorem \hyperref[cl:thm:34]{34} & Appendix \hyperref[sn:D.2]{D.2} \\
Abstract interpretation & Subsection \hyperref[sn:12.1]{12.1} & Theorems \hyperref[cl:thm:26]{26}--\hyperref[cl:thm:28]{28}, Model Theorem \hyperref[cl:thm:29]{29} & Appendix \hyperref[sn:D.3]{D.3} \\
Graph/cohomology & Subsection \hyperref[sn:8.1]{8.1} & Theorems \hyperref[cl:thm:15]{15}--\hyperref[cl:thm:18]{18}; realization in Subsection \hyperref[sn:13.4]{13.4} & Appendix \hyperref[sn:D.4]{D.4} \\
\bottomrule
\end{longtable}
\endgroup

The discipline of the section is uniform: each realization is scoped, each declares its host, instrument, source-of-truth policy, audit policy, visibility policy, and nonclaim register, and each realizes a declared fragment of the calculus. None of the four realizations is claimed to be the whole calculus, and the section does not state a global model-coverage theorem.

\subsection{Model-Realization Convention}\label{sn:13.1}

This subsection fixes the form of a model-realization statement and the fragment discipline imposed on every model realization in the section.

\subsubsection{Realization Map}\label{sn:13.1.1}

A model-realization statement has the form

\[
\Gamma\,;\;\mathcal T\,;\;I\;\vdash\;\mathcal M\;\models\;\mathcal S\;:\;\texttt{model\_realization},
\]

where:

\begin{itemize}
\tightlist
\item
  \(\mathcal M\) is the realization object: a finite typed record carrying the host data, the instrument, the source-of-truth policy, the audit policy, the visibility policy, and the nonclaim register required by the host;
\item
  \(\mathcal S\) is the declared structure: a finite, named fragment of \(\mathbf{BirdInt}^{\mathrm{aud}}_{\mathrm{fin}}\), written either as a role fragment (for example, \(P_1/P_2/P_4/P_5/P_6\)) or as a named substructure (for example, \(\mathrm{cell/pair/prov}\) for the finite stochastic cell/pair/provenance fragment, or \(0\to1\) for an audited strict-bridge passage); the PICA fragment is written \(\mathbf{BirdInt}^{\mathrm{aud}}_{\mathrm{fin\text{-}stoch}}{}^{\mathrm{cell/pair/prov}}\), the subscript recording the stochastic host, and the Cantor fragment \(\mathrm{StrictBridge}^{\mathrm{audited\ shell}}_{0\to1}\), the \(0\to1\) fragment on the audited shell;
\item
  \(I\) is the instrument under which the realization claim is recorded, and \(\mathcal T\) is the surrounding target package, in the sense of Subsection \hyperref[sn:11.1]{11.1}.
\end{itemize}

The verdict \(\texttt{model\_realization}\) records that the realization map of \(\mathcal M\) supplies finite host data satisfying clauses 1--7 below, with each assigned audit replayed and its computed outcome recorded; every assigned audit passes, except for failure codes that determine the classifier status carried under clause 7 (for example \(\texttt{circular\_audit}\) on a cell of status \(\texttt{undefined\_circular}\)). The verdict \(\texttt{model\_realization}\) is the verdict of a model-grade result (Subsubsection \hyperref[sn:2.3.4]{2.3.4}) and is distinct from the per-claim status \(\texttt{accepted}\) in \(\mathrm{ClaimStatus}\). The model-realization judgment is a judgment form of its own, separate from the instrument-indexed claim judgment of Subsection \hyperref[sn:4.8]{4.8}: \(\Gamma;\mathcal T;I\vdash\mathcal M\models\mathcal S:\texttt{model\_realization}\) holds exactly when the realization map of \(\mathcal M\) to \(\mathcal S\) defined below satisfies clauses 1--7.

The realization map of \(\mathcal M\) to \(\mathcal S\) is a finite typed family of assignments

\[
\mathrm{real}_{\mathcal M\to\mathcal S}\;=\;\bigl\{\,\mathrm{real}_X\;:\;X\in\mathrm{Components}(\mathcal S)\,\bigr\},
\]

with each \(\mathrm{real}_X\) mapping a host datum from \(\mathcal M\) to the corresponding component \(X\) of \(\mathcal S\). The components of the fragments used in this section are as follows. For \(\mathrm{cell/pair/prov}\): the thirty-six directed-cell records \(\mathsf{Cell}_{ij}\), \(1\le i,j\le 6\), each with its fields \(W_j,U_i,L,V,\Theta,A,\delta\) and profile; the fifteen pair-observable records \(\mathsf{PairObs}_{\{i,j\}}\), \(i\ne j\), each with its source-of-truth record; the provenance and audit records they reference; and the finite dependency and ablation records. For a role fragment such as \(P_1/P_2/P_4/P_5/P_6\): for each listed role \(P_i\), a role-channel record \((\mathrm{Chan}_i,R_i,\Theta_i^{\mathrm{act}},A_{R_i},V_{R_i},\mathsf{collapse}_i)\) of Subsubsection \hyperref[sn:5.6.1]{5.6.1} to which the role-channel classifier assigns \(\texttt{active\_projection}\) under \(I\), and the witness and update records it uses. For \(0\to 1\): the expanded bridge tuple of Subsubsection \hyperref[sn:10.1.1]{10.1.1}, including its two packages and gate record. For \(P_2/P_6^{H^1}\text{-drive}/P_3\text{-separation}\): a \(P_2\) gate record \(G\leadsto G'\) with its cycle-support defect, the \(P_6\) audit record \(A\), a drive record \(\kappa_6^{\mathrm{drive}}(a)\) or a no-drive certificate of Subsubsection \hyperref[sn:8.1.2]{8.1.2}, and NC-12 in the nonclaim register. The realization map is admissible when:

\begin{enumerate}
\def\labelenumi{\arabic{enumi}.}
\tightlist
\item
  every component of \(\mathcal S\) is assigned a finite host datum from \(\mathcal M\);
\item
  each assignment carries an audit record judged under the host instrument \(I\) whose checks are replayed and whose audit defect is \(\texttt{audit\_passes}\), up to the exception stated above;
\item
  every source record referenced by a realized record is admitted by \(\mathsf{SourcePolicy}_I\) of the host instrument \(I\);
\item
  the visibility policy of \(\mathcal M\) exposes the records required by \(\mathcal S\), with visibility bridges of Subsubsection \hyperref[sn:3.5.3]{3.5.3} supplied where required;
\item
  the nonclaim register of \(\mathcal M\) contains the model-realization nonclaims of Subsection \hyperref[sn:13.6]{13.6} and any host-specific nonclaims of the realization;
\item
  every realized witness or update has a signature in \(\mathsf{Wit}(P_j)\) or \(\mathsf{Upd}(P_i)\) for the role it realizes (Subsection \hyperref[sn:3.7]{3.7} and the convention of Subsubsection \hyperref[sn:3.1.1]{3.1.1});
\item
  every realized record that carries a status (a directed cell, a pair observable or a promotion bridge) carries the status that the classifier of Subsection \hyperref[sn:5.6]{5.6} assigns to the realized record; a status reported by the host that differs from it is recorded as a discrepancy and is not the status of the record.
\end{enumerate}

Admissibility of the realization map is a finite predicate on \((\mathcal M,\mathcal S,I)\). Clauses 6 and 7 are what make the realization respect the calculus: the host data are typed as the calculus types them, and the verdicts on realized records are the calculus's own verdicts.

\subsubsection{Fragment Discipline}\label{sn:13.1.2}

Every model realization in the section is scoped to a declared fragment \(\mathcal S\) of \(\mathbf{BirdInt}^{\mathrm{aud}}_{\mathrm{fin}}\). The fragment discipline imposes the following rules on every realization claim:

\begin{itemize}
\tightlist
\item
  \emph{Fragment-only realization.} A model-realization statement \(\mathcal M\models\mathcal S:\texttt{model\_realization}\) asserts realization of \(\mathcal S\) and not of the whole calculus. The verdict \(\texttt{model\_realization}\) for \(\mathcal S\) does not imply a \(\texttt{model\_realization}\) verdict for any larger fragment \(\mathcal S'\supsetneq\mathcal S\).
\item
  \emph{No theorem grade by realization.} A model-realization verdict is not a theorem-grade verdict. In the absence of a separately admissible theorem-grade bridge, \(\mathcal M\models\mathcal S:\texttt{model\_realization}\) does not yield a \(\texttt{accepted}\) verdict on a theorem-grade claim about \(\mathcal S\). This is the model-realization nonclaim of Subsection \hyperref[sn:2.4]{2.4}.
\item
  \emph{Host-bound realization.} A model-realization verdict for \(\mathcal S\) under \(\mathcal M\) is bound to the host of \(\mathcal M\): it does not transfer to other hosts without an admissible instrument-transfer bridge in the sense of Subsection \hyperref[sn:11.9]{11.9}.
\item
  \emph{No model-coverage theorem.} The four realizations recorded in Subsections \hyperref[sn:13.2]{13.2}--\hyperref[sn:13.4]{13.4} and Section \hyperref[sn:12]{12} do not, taken together, support a global model-coverage theorem for \(\mathbf{BirdInt}^{\mathrm{aud}}_{\mathrm{fin}}\). This is the model-coverage nonclaim of Subsection \hyperref[sn:13.6]{13.6}.
\end{itemize}

The fragment discipline is the central structural constraint of the section: it fixes the meaning of \(\models\), prevents over-reading of the model-realization verdict, and defines the boundary between model-grade and theorem-grade claims throughout Sections \hyperref[sn:13]{13} and \hyperref[sn:12]{12}. The four realizations of the section are each scoped to a declared fragment, and the four fragments are recorded in the combined matrix of Subsection \hyperref[sn:13.5]{13.5}.

\subsection{PICA Realization}\label{sn:13.2}

This subsection records the PICA model-realization sheet and Model Theorem \hyperref[cl:thm:33]{33}: any store family meeting \(\mathsf{AdmPICAReal}\) induces a model of the finite stochastic cell/pair/provenance fragment of \(\mathbf{BirdInt}^{\mathrm{aud}}_{\mathrm{fin}}\); the stores constructed to satisfy the predicate read the chain only through positive-probability traces; the predicate is met over every nonempty finite Markov chain for every raw status table, and is not checked here on the stores of the PICA simulator. The PICA realization is host-bound and fragment-scoped: it does not realize the whole calculus, the recovered \(6\times 6\) table is a PICA-profile fact rather than a universal law, and raw PICA statuses recover into formal statuses only via the host classifier.

\emph{Orientation.} PICA, the Primitive Interaction Closure Algebra, is the interaction algebra of the minimal stochastic substrate of \citet{Tsiokos2026PICA}: a finite Markov chain whose dynamics are built from six primitive mechanisms, simulated with audit records for every run. That work records a raw status (Action, Implicit, Trivial or Undefined) for each of the thirty-six ordered pairs of primitives, one for each directed interaction; Subsubsection \hyperref[sn:13.2.6]{13.2.6} maps these raw statuses to the formal cell statuses of the calculus. This subsection reads the simulator's finite stores as host data of the calculus and states the conditions (\(\mathsf{AdmPICAReal}\) and those of Subsubsection \hyperref[sn:13.2.5]{13.2.5}) under which they realize its cell, pair and provenance fragment. In the symbols below, \(\Omega\) is the finite state space of the chain, \(K\) its Markov kernel, and the other components of \(\mathsf H_{\mathrm{PICA}}\) are the finite stores of the implementation.

\subsubsection{Name and Host}\label{sn:13.2.1}

The PICA realization name is \(\mathsf{PICAReal}_{\mathrm{fin}}\). Its host is the finite stochastic implementation context \(\mathsf H_{\mathrm{PICA}}\), with components

\[
\mathsf H_{\mathrm{PICA}}\;=\;(\,\Omega,\;K,\;\mathbb P,\;\mathsf{Bus},\;\mathsf{CellRules},\;\mathsf{WitnessStore},\;\mathsf{UpdateStore},\;\mathsf{PairStore},\;\mathsf{DepGraph},\;A,\;V\,),
\]

where \(\Omega\) is a nonempty finite stochastic state space, \(K\) a row-stochastic Markov kernel on \(\Omega\), recorded in the stores by its finite support table \(\operatorname{supp}K=\{(\omega,\omega'):K(\omega,\omega')>0\}\), together with its entries when these lie in a declared numerical type with decidable order (the trace-support rule reads only \(\operatorname{supp}K\)), \(\mathbb P=\{P_1,\ldots,P_6\}\) the six primitive labels, and the remaining fields finite stores for the bus, cell rules, witness records, actor-update records, pair observables, the producer/consumer dependency graph, the host audit record \(A\), and the host visibility record \(V\). The host instrument is \(I_{\mathrm{PICA}}\), an admissible instrument in the sense of Subsubsection \hyperref[sn:11.1.1]{11.1.1} whose visible class contains the host stores and whose audit policy is clause (ii) of Subsubsection \hyperref[sn:13.2.2]{13.2.2} and whose source policy is a relation of Subsubsection \hyperref[sn:5.3.6]{5.3.6} containing \((t,\texttt{full},\lambda)\) for every profile \(\lambda\) and the default entry, for \(t=\texttt{committed\_state}\) and, where the host fallback policy allows it, \(t=\texttt{fallback}\), together with \((\texttt{fallback},\texttt{provisional},\mu_{\mathrm{PICA}})\); \(\mathsf{SourceOfTruthOK}\) enters through source traces (Subsubsection \hyperref[sn:13.2.7]{13.2.7}).

\subsubsection{Realization Object}\label{sn:13.2.2}

A PICA realization object is the finite typed tuple

\[
\begin{aligned}
\mathcal P\;=\;(\,&\Omega,\;K,\;\mathbb P,\;\mathsf{Bus},\;\mathsf{CellRules},\;\mathsf{WitnessStore},\;\mathsf{UpdateStore},\\
&\mathsf{StatusStore},\;\mathsf{PairStore},\;\mathsf{DepGraph},\;\mathsf{Abl},\;A,\;V,\;\mathcal N\,),
\end{aligned}
\]

extending the host data with \(\mathsf{StatusStore}\) (raw and recovered statuses), \(\mathsf{Abl}\) (finite ablation/model evidence), and \(\mathcal N\) (the PICA nonclaim register of Subsubsection \hyperref[sn:13.2.10]{13.2.10}). The object is admissible, written \(\mathsf{AdmPICAReal}(\mathcal P)\), when (i) each store is finite; (ii) every entry carries an audit record judged under \(I_{\mathrm{PICA}}\) whose checks are replayed, and every such audit defect is empty except that the audit of a cell with raw status \(\texttt{Undefined}\), and the audit of a pair record with raw status \(\texttt{blocked}\), may contain the failure codes that determine its formal status under clause 7 of Subsubsection \hyperref[sn:13.1.1]{13.1.1} (for example \(\texttt{circular\_audit}\) for \(\texttt{undefined\_circular}\), \(\texttt{threshold\_failure}\) for \(\texttt{below\_threshold}\), and \(\texttt{source\_failure}\) or \(\texttt{visibility\_failure}\) for a blocked pair); (iii) \(\mathsf{CellRules}\) contains exactly one cell record for each of the thirty-six ordered pairs of primitives and \(\mathsf{PairStore}\) exactly one pair record for each of the fifteen unordered pairs; a second record named by a duplicate-pair record is kept in \(A\), not in \(\mathsf{PairStore}\); (iv) every source record referenced by any realized component, including witness and update records and their audits, is admitted by \(\mathsf{SourcePolicy}_{I_{\mathrm{PICA}}}\); (v) the directed-cell, pair-observable, source/audit, and status-recovery records satisfy the schemas used in Subsections \hyperref[sn:4.3]{4.3}, \hyperref[sn:4.4]{4.4}, \hyperref[sn:4.8]{4.8}, and \hyperref[sn:5.6]{5.6}; (vi) \(\mathcal N\) contains the three coverage nonclaims of Subsubsection \hyperref[sn:13.6.2]{13.6.2}; (vii) every witness and update record carries a trace \((\omega_0,\dots,\omega_m)\) in \(\Omega\), \(m\ge1\), with \(K(\omega_{r-1},\omega_r)>0\) for each \(r\), replayed by its audit; and (viii) every source record \(R\) whose \(\mathrm{valid}_{\mathrm{src}}\) records a nonempty dirty-since step, or that fails a conjunct of \(\mathsf{SourceOfTruthOK}(R,Q)\) for a pair record \(Q\) whose source slot or audits name \(R\), carries in \(\mathrm{trace}_{\mathrm{src}}\) a failed check of \(R\) itself; and for every source record \(R\) whose \(\mathrm{valid}_{\mathrm{src}}\) records a nonempty dirty-since step, whose validity condition fails, or whose \(\mathrm{trace}_{\mathrm{src}}\) records a failed audit or threshold check of \(R\) itself, \(\mathsf{SourceOfTruthOK}(R,Q)\) is false for every pair record \(Q\) of the instance. In addition, the visibility policy exposes every record required by the realized fragment, including referenced audit evidence, with admissible visibility bridges supplied where required, as stipulated by clause 4 of Subsubsection \hyperref[sn:13.1.1]{13.1.1}. The audit-failure exceptions do not waive this requirement. Every cell record of \(\mathsf{CellRules}\) is well formed in the sense of Subsubsection \hyperref[sn:4.3.3]{4.3.3} and every pair record of \(\mathsf{PairStore}\) in the sense of Subsubsection \hyperref[sn:4.4.1]{4.4.1}. Admissibility is a finite predicate on the stores.

The associated realization map is

\[
\mathcal R_{\mathrm{PICA}}\;:\;\mathsf{PICAReal}_{\mathrm{fin}}\longrightarrow
\mathbf{BirdInt}^{\mathrm{aud}}_{\mathrm{fin\text{-}stoch}}{}^{\,\mathrm{cell}/\mathrm{pair}/\mathrm{prov}},
\]

sending \(\mathcal P\) to the realized finite stochastic model \(\mathcal M_{\mathcal P}\). The induced package is

\[
\mathcal T_{\mathrm{PICA}}\;=\;(\,\Omega,\;f_{\mathrm{PICA}},\;\Sigma_{f_{\mathrm{PICA}}},\;E_{\mathrm{PICA}},\;\mathcal A_{\mathrm{PICA}}\,),
\]

where \(f_{\mathrm{PICA}}\) is the bus/lens projection, \(\Sigma_{f_{\mathrm{PICA}}}=\{f_{\mathrm{PICA}}^{-1}(B):B\subseteq f_{\mathrm{PICA}}(\Omega)\}\) and \(E_{\mathrm{PICA}}=E_{f_{\mathrm{PICA}}}\) are the lens-induced content and completion, with the finite visible-claim set and the commit/update closure recorded as attached records, and \(\mathcal A_{\mathrm{PICA}}\) is the inherited package-level audit functional assembled from the finite PICA audit record \(A\).

\subsubsection{Actor-Update Map}\label{sn:13.2.3}

The PICA actor-update map sends each primitive \(P_i\) to a finite update sort drawn from \(\mathsf{UpdateStore}\):

\[
P_i\;\longmapsto\;\mathsf{Upd}_{\mathrm{PICA}}(P_i),
\]

with the assignments

\begin{longtable}[]{@{}ll@{}}
\toprule
primitive & PICA actor update sort\tabularnewline
\midrule
\endhead
\(P_1\) & descent repair / closure update\tabularnewline
\(P_2\) & support gate / admissibility update\tabularnewline
\(P_3\) & route-memory / mismatch update\tabularnewline
\(P_4\) & refinement / staging update\tabularnewline
\(P_5\) & package / completion / closure-summary update\tabularnewline
\(P_6\) & audit / provenance / source update\tabularnewline
\bottomrule
\end{longtable}

Each row records a finite update sort consumed by the directed-cell record of Subsubsection \hyperref[sn:13.2.5]{13.2.5}; the assignment is a typed map from \(\mathbb P\) to the PICA update sorts, audited by \(I_{\mathrm{PICA}}\).

\subsubsection{Informant-Witness Map}\label{sn:13.2.4}

The PICA informant-witness map sends each primitive \(P_j\) to a finite witness sort drawn from \(\mathsf{WitnessStore}\):

\[
P_j\;\longmapsto\;\mathsf{Wit}_{\mathrm{PICA}}(P_j),
\]

with the assignments

\begin{longtable}[]{@{}ll@{}}
\toprule
primitive & PICA informant witness sort\tabularnewline
\midrule
\endhead
\(P_1\) & descent/closure defect\tabularnewline
\(P_2\) & gate/support/admissibility witness\tabularnewline
\(P_3\) & route mismatch / route-memory witness\tabularnewline
\(P_4\) & refinement/stage witness\tabularnewline
\(P_5\) & package/closure/completion witness\tabularnewline
\(P_6\) & audit/provenance/source witness\tabularnewline
\bottomrule
\end{longtable}

The actor-update and informant-witness maps together specify the per-cell witness/update content for the directed-cell realization of Subsubsection \hyperref[sn:13.2.5]{13.2.5}.

\subsubsection{Directed-Cell Realization}\label{sn:13.2.5}

Each PICA directed cell becomes a typed claim in the calculus:

\[
\Gamma\,;\;\mathcal T_{\mathrm{PICA}}\,;\;I_{\mathrm{PICA}}\;\vdash\;\mathsf{Cell}_{ij}^{\lambda_{\mathrm{PICA}}}\;:\;\sigma_{ij},
\]

with the cell record

\[
(P_i\leftarrow P_j,\;W_j,\;U_i,\;L,\;V,\;\Theta,\;A,\;\delta),
\]

where \(L\) is the level/lens/interface record of Subsubsection \hyperref[sn:4.3.3]{4.3.3}, \(V\) is the visibility record, \(\Theta\) the threshold record, \(A\) the audit record, and \(\delta\) the cell defect record. For each cell, \(W_j\) is a selected record from \(\mathsf{WitnessStore}\) whose signature belongs to \(\mathsf{Wit}(P_j)\), and \(U_i\) is a selected record from \(\mathsf{UpdateStore}\) whose signature belongs to \(\mathsf{Upd}(P_i)\). The sorts in Subsubsections \hyperref[sn:13.2.3]{13.2.3}--\hyperref[sn:13.2.4]{13.2.4} specify the permitted record kinds; they are not themselves the cell's witness and update fields. The cell profile is \(\lambda_{\mathrm{PICA}}=(\mathsf{Ver},\mathsf{Ver},\texttt{fine},\texttt{dynamical},\texttt{visibility},\texttt{default},\texttt{local})\), and each pair record carries \(\mu_{\mathrm{PICA}}\) with the same entries. \(I_{\mathrm{PICA}}\) has \(\texttt{directed-cell},\texttt{pair-observable}\in\mathsf{ClaimTypes}_{I_{\mathrm{PICA}}}\), \(\mathsf{Ver}\in\mathsf{Levels}_{I_{\mathrm{PICA}}}\), \(\mathsf{ModePolicy}_{I_{\mathrm{PICA}}}(\text{directed-cell})=\mathsf{ModePolicy}_{I_{\mathrm{PICA}}}(\text{pair-observable})=\texttt{default}\), \((\texttt{fine},\texttt{dynamical},\texttt{local},\mathsf H_{\mathrm{PICA}})\in\mathsf{Compat}_{I_{\mathrm{PICA}}}\), and a bridge policy admitting the visibility type. Per-cell thresholds are carried by \(\Theta\).

The raw statuses of the thirty-six pairs are imported from \citet{Tsiokos2026PICA}. The realized records attach the witnesses of the eight raw \(\texttt{Implicit}\) cells, and no other witness, to \(\mathcal T_{\mathrm{PICA}}\) (so \(W_{\mathcal T_{\mathrm{PICA}}}\) consists of these eight records, as in the proof of Model Theorem \hyperref[cl:thm:33]{33}), with no threshold evaluating them, so that step 8 computes \(\mathsf{implicitW}\) on those cells, and the audit record of the raw \(\texttt{Undefined}\) cell lists among its evidence references the verdict reference of that cell under \(I_{\mathrm{PICA}}\), recording the circularity the PICA audit reports for it, so step 8 computes a failing \(\mathsf{CircularityCheck}\) there. The two raw \(\texttt{Trivial}\) cells have recorded updates whose effect maps fix their targets, so step 8 computes \(\mathsf{noopU}=\mathrm{true}\); the other thirty-four have \(\mathsf{noopU}=\mathrm{false}\). If all thirty-six cells are in scope, have their witness and update present, satisfy their thresholds, have empty visibility and required-bridge defects, have no package-collapse entry or other failing defect, and pass every audit check except that one circularity check, and no raw \(\texttt{Action}\) cell has as its witness one of the eight attached records (for example, the thirty-six witness records are pairwise distinct), on these records the classifier of Subsubsection \hyperref[sn:5.6.2]{5.6.2} returns the PICA-profile multiplicity record

\[
25\;\texttt{action},\quad 8\;\texttt{implicit},\quad 2\;\texttt{trivial},\quad 1\;\texttt{undefined\_circular}.
\]

Since the attachments and the circularity code are chosen from the raw statuses, this count re-expresses the raw table under the recovery map; it is not an independent computation.

This agrees with the recovery map of Subsubsection \hyperref[sn:13.2.6]{13.2.6}, in which the one raw \(\texttt{Undefined}\) pair, with its circular audit, is recovered as \(\texttt{undefined\_circular}\). It records the multiplicities of the four directed-cell statuses within the PICA host under the recovery map of Subsubsection \hyperref[sn:13.2.6]{13.2.6}, and is not a universal table theorem about \(6\times 6\) tables in \(\mathbf{BirdInt}^{\mathrm{aud}}_{\mathrm{fin}}\).

\subsubsection{Status Recovery}\label{sn:13.2.6}

Raw PICA statuses recover as formal statuses through the recovery map

\[
\texttt{Action}\;\longmapsto\;\texttt{action},\qquad\texttt{Implicit}\;\longmapsto\;\texttt{implicit},\qquad\texttt{Trivial}\;\longmapsto\;\texttt{trivial},
\]

\[
\texttt{Undefined}\;\longmapsto\;\text{reason-dependent formal status},
\]

where the reason-dependent recovery of \(\texttt{Undefined}\) is one of \(\texttt{blocked}\), \(\texttt{undefined\_circular}\), \(\texttt{outside\_scope}\), \(\texttt{absent}\), \(\texttt{below\_threshold}\), or \(\texttt{collapsed}\), decided by the formal classifier of Subsection \hyperref[sn:5.6]{5.6} acting on the cell record, the threshold record, the audit record, and the source-of-truth record.

Raw PICA status never overrides the formal classifier verdict: if the recovered formal status assigned by the classifier differs from the raw status, the formal status is the authoritative cell status, and the raw/recovered discrepancy is recorded in the source/provenance correction record of \(\mathsf{StatusStore}\). The same holds for pair observables below: the formal pair status is the verdict of the pair classifier of Subsubsection \hyperref[sn:5.6.3]{5.6.3} on the realized record.

\subsubsection{Pair-Observable Realization}\label{sn:13.2.7}

PICA pair statuses recover as formal pair statuses through

\[
\texttt{real}\;\longmapsto\;\texttt{real},\qquad\texttt{provisional}\;\longmapsto\;\texttt{real\_provisional},\qquad\texttt{duplicated}\;\longmapsto\;\texttt{real\_duplicated},
\]

\[
\texttt{blocked}\;\longmapsto\;\texttt{blocked},
\]

with a PICA pair observable record

\[
\begin{aligned}\mathsf{PairObs}_{\{i,j\}}^{\mu}\;=\;(\,&\{P_i,P_j\},\;\mu,\;\mathsf{ObservableType},\;\mathsf{BranchData},\;\mathsf{SourceOfTruth},\;L,\;V,\;\Theta,\\&A,\;\delta,\;\mathcal N,\;\mathsf{forbidsSource},\;\mathsf{partialEvidence},\;\mathsf{Dup}\,).\end{aligned}
\]

Here \(\mathsf{ObservableType}\) is the joint-witness field and \(\mathsf{BranchData}\) the branch-compatibility record of Subsubsection \hyperref[sn:4.4.1]{4.4.1}, and \(\mathsf{Dup}\) is present only for a duplicated pair.

Realness of the pair observable requires the conjunction

\[
\mathsf{SourceOfTruthOK}\;\land\;\mathsf{VisibilityPasses}\;\land\;\mathsf{AuditPasses}\;\land\;\mathsf{ThresholdPasses},
\]

Here \(\mathsf{SourceOfTruthOK}\) is the source-of-truth predicate of Subsubsection \hyperref[sn:13.2.8]{13.2.8}, and the other three predicates are the host-level visibility, audit, and threshold checks of Subsection \hyperref[sn:11.3]{11.3}, Subsubsection \hyperref[sn:4.9.2]{4.9.2} and Subsection \hyperref[sn:5.2]{5.2}. Since \(\mathsf{SourcePolicy}_{I_{\mathrm{PICA}}}\) relates only tags, modes and profiles (Subsubsection \hyperref[sn:5.3.6]{5.3.6}), the PICA store discipline records the failure of any conjunct of \(\mathsf{SourceOfTruthOK}\) as a failed check of \(R\) itself in \(\mathrm{trace}_{\mathrm{src}}\). The source defect of \(R\) is then \(\texttt{dirty}\), unless an earlier clause assigns \(\texttt{unknown}\) or \(\texttt{contradictory}\); in every case it fails the source threshold, so the first \(\texttt{blocked}\) row applies, even when an upgrade bridge is present. The \(\texttt{real\_duplicated}\) row reads the duplicate-pair record together with the conditions of the \(\texttt{real}\) row, so it also requires all four checks; if a visibility, audit or threshold defect fails, the first \(\texttt{blocked}\) row applies. An admitted fallback source, or a proxy source obtained through an admitted, audited source bridge, may receive \(\texttt{real\_provisional}\); other cases that meet no earlier row receive \(\texttt{blocked}\).

\subsubsection{Provenance Discipline}\label{sn:13.2.8}

The PICA source-of-truth predicate is

\[
\mathsf{SourceOfTruthOK}(R,Q)\;\;\Longleftrightarrow\;\;\mathsf{Committed}(R)\;\land\;\mathsf{CorrectProducer}(R,Q)\;\land\;\mathsf{CorrectProvenance}(R,Q)
\]

\[
\land\;\mathsf{NotFallbackUnlessAllowed}(R,Q)\;\land\;\mathsf{Visible}(R)\;\land\;\mathsf{Audited}(R),
\]

for a record \(R\) and a query/observable \(Q\). Each conjunct is a finite check on the host stores: \(\mathsf{Committed}\) on \(\mathsf{Bus}\) and \(\mathsf{StatusStore}\); \(\mathsf{CorrectProducer}\) on \(\mathsf{DepGraph}\); \(\mathsf{CorrectProvenance}\) on \(A\); \(\mathsf{NotFallbackUnlessAllowed}\) on the host's fallback policy; \(\mathsf{Visible}\) on \(V\); \(\mathsf{Audited}\) on \(A\).

The provenance discipline imposes two non-negotiable rules:

\begin{itemize}
\tightlist
\item
  \emph{Dirty exclusion.} Clause (viii) of \(\mathsf{AdmPICAReal}\) requires every record whose \(\mathrm{valid}_{\mathrm{src}}\) records a nonempty dirty-since step to carry in \(\mathrm{trace}_{\mathrm{src}}\) a failed check of \(R\) itself. Consequently its source threshold fails: its source defect is \(\texttt{dirty}\) unless an earlier clause assigns \(\texttt{unknown}\) or \(\texttt{contradictory}\); and \(\neg\,\mathsf{SourceOfTruthOK}(R,Q)\). A dirty record cannot be a source of truth for any query, and a dirty record in a pair's source slot, or among the \(\mathsf{SourceRefs}\) of its audit or branch-compatibility audit, cannot support a \(\texttt{real}\) pair observable: the first \(\texttt{blocked}\) row applies.
\item
  \emph{Fallback-default cap.} \(R\;\text{is fallback}\;\Longrightarrow\;Q\;\text{recovers as}\;\texttt{real\_provisional}\;\text{at most}\), unless \(\mathrm{trace}_{\mathrm{src}}(R)\) names a source-upgrade bridge admitted by \(\mathsf{SourcePolicy}_{I_{\mathrm{PICA}}}\) and \(\mathsf{BridgePolicy}_{I_{\mathrm{PICA}}}\) (Subsubsection \hyperref[sn:5.3.6]{5.3.6}). Without that admitted upgrade, a fallback record supports at most a provisional pair observation.
\end{itemize}

The discipline ensures that pair-realness is grounded in committed, correct, audited, visible, non-dirty, and not-merely-fallback records; under any failure of these conditions, the pair recovers as one of the non-real statuses of Subsubsection \hyperref[sn:13.2.7]{13.2.7}.

\subsubsection{Model Theorem 33: PICA Model-Realization}\label{sn:13.2.9}

Call a PICA realization object \(\mathcal P\) \emph{\(K\)-independent} if \(K\) is named, in the reference graph of Subsubsection \hyperref[sn:5.6.5]{5.6.5}, only by the \(K\) fields of \(\mathcal P\) and its host and by trace-support audit entries, and no rule of \(\mathsf{CheckRules}_{I_{\mathrm{PICA}}}\) other than the trace-support rule, which checks \(K(\omega_{r-1},\omega_r)>0\), takes \(K\) as an argument; and every other admissibility or classifier computation, including threshold and defect computations and transitive calls to named programs, depends on trace-support checks only through their Boolean outcomes and does not otherwise read \(K\), directly or indirectly. For such \(\mathcal P\) the realized statuses are determined by the stores, and \(K\) enters only through the support of the recorded traces.

\begin{claimtheorem}{Model Theorem 33 (PICA Model-Realization)}\phantomsection\label{cl:thm:33}

Let \(\mathcal P\in\mathsf{PICAReal}_{\mathrm{fin}}\) be admissible: \(\mathsf{AdmPICAReal}(\mathcal P)\) holds. Then the realization map \(\mathcal R_{\mathrm{PICA}}\) supplies a finite, audited assignment of host data to every component of the finite stochastic cell/pair/provenance fragment, so

\[
\mathcal R_{\mathrm{PICA}}(\mathcal P)\;\models\;\mathbf{BirdInt}^{\mathrm{aud}}_{\mathrm{fin\text{-}stoch}}{}^{\,\mathrm{cell}/\mathrm{pair}/\mathrm{prov}}.
\]

Moreover: (i) for every nonempty finite \(\Omega\) and row-stochastic Markov kernel \(K\) on \(\Omega\), and every assignment of raw statuses in \(\{\texttt{Action},\texttt{Implicit},\texttt{Trivial},\texttt{Undefined}\}\) to the thirty-six ordered pairs, \(\mathsf{AdmPICAReal}\) is satisfiable over \((\Omega,K)\) by stores on which the cell classifier reproduces that assignment under the recovery map of Subsubsection \hyperref[sn:13.2.6]{13.2.6}, with \(\texttt{Undefined}\mapsto\texttt{undefined\_circular}\), with no cell discrepancy; (ii) for pairs, for every assignment of raw pair statuses in \(\{\texttt{real},\texttt{provisional},\texttt{duplicated},\texttt{blocked}\}\) to the fifteen pairs, the stores can be chosen so that the pair classifier reproduces it under the recovery map of Subsubsection \hyperref[sn:13.2.7]{13.2.7}, with no pair discrepancy. Clause (i), applied to the raw cell table of \citet{Tsiokos2026PICA}, gives \(25\) \(\texttt{action}\), \(8\) \(\texttt{implicit}\), \(2\) \(\texttt{trivial}\) and \(1\) \(\texttt{undefined\_circular}\) cells, multiplicities that are an input of the realization, not a consequence of \(\mathsf{AdmPICAReal}\). (iii) in particular, every assignment of \(\texttt{action}\), \(\texttt{implicit}\), \(\texttt{trivial}\) and \(\texttt{undefined\_circular}\) to the thirty-six cells is realized over every nonempty finite \((\Omega,K)\) by suitable stores. (iv) the stores constructed in the proof are \(K\)-independent: once their traces are fixed, only the positive-probability checks read \(K\). (v) consequently, if \(\mathsf{AdmPICAReal}(\mathcal P)\) holds over \((\Omega,K)\), \(\mathcal P\) is \(K\)-independent, and \(K'\) is a row-stochastic kernel on \(\Omega\) positive on every recorded trace step, then \(\mathcal P\) with \(K\) replaced by \(K'\) (each trace replay rerun under \(K'\)) satisfies \(\mathsf{AdmPICAReal}\), and every realized cell and pair keeps its classifier status.

\end{claimtheorem}

The paper does not verify \(\mathsf{AdmPICAReal}\) for the stores of \citet{Tsiokos2026PICA}.

\begin{proof}

By \(\mathsf{AdmPICAReal}(\mathcal P)\), \(\mathsf{CellRules}\) and \(\mathsf{PairStore}\) contain one record for each of the thirty-six cells and fifteen pair observables listed as components in Subsubsection \hyperref[sn:13.1.1]{13.1.1}, so each component of the realized fragment is assigned a finite host datum and audit record from \(\mathcal P\):

\begin{itemize}
\tightlist
\item
  \emph{Typed directed cells.} Subsubsections \hyperref[sn:13.2.3]{13.2.3} and \hyperref[sn:13.2.4]{13.2.4} supply the actor-update and informant-witness maps; Subsubsection \hyperref[sn:13.2.5]{13.2.5} assembles them into the directed-cell record under the cell profile \(\lambda_{\mathrm{PICA}}\).
\item
  \emph{Cell-status recovery.} Subsubsection \hyperref[sn:13.2.6]{13.2.6} supplies the recovery map from raw to formal statuses, with the formal classifier of Subsection \hyperref[sn:5.6]{5.6} as the authoritative verdict.
\item
  \emph{Pair-observable contracts.} Subsubsection \hyperref[sn:13.2.7]{13.2.7} supplies the pair-observable record and the four-conjunct realness predicate.
\item
  \emph{Source-of-truth and provenance discipline.} Subsubsection \hyperref[sn:13.2.8]{13.2.8} supplies the \(\mathsf{SourceOfTruthOK}\) predicate and the dirty-exclusion and fallback-default rules.
\item
  \emph{Finite dependency and ablation records.} The host stores \(\mathsf{DepGraph}\) and \(\mathsf{Abl}\) supply finite dependency and ablation records, recorded in \(A\) under \(I_{\mathrm{PICA}}\).
\end{itemize}

Each assignment is a finite host datum, audited by \(I_{\mathrm{PICA}}\) and visible under it; clause 3 holds because \(\mathsf{AdmPICAReal}(\mathcal P)\) requires every referenced source record to be admitted by \(\mathsf{SourcePolicy}_{I_{\mathrm{PICA}}}\). Clause 5 of Subsubsection \hyperref[sn:13.1.1]{13.1.1} holds by \(\mathsf{AdmPICAReal}(\mathcal P)\). For clause 6, the witness and update records selected for each cell in Subsubsection \hyperref[sn:13.2.5]{13.2.5} have signatures in \(\mathsf{Wit}(P_j)\) and \(\mathsf{Upd}(P_i)\). For clause 7, Subsubsection \hyperref[sn:13.2.6]{13.2.6} makes the classifier's verdict the status of every realized cell and Subsubsection \hyperref[sn:13.2.7]{13.2.7} does the same for every realized pair observable, with raw PICA statuses entering only through the recovery maps and the recorded discrepancies. The realization map is admissible in the sense of Subsubsection \hyperref[sn:13.1.1]{13.1.1}, so the scoped model-realization statement receives verdict \(\texttt{model\_realization}\).

\emph{Traces and instrument.} For the satisfiability clause, take any nonempty finite \(\Omega\) and row-stochastic Markov kernel \(K\) on it; attach to every witness and update record a trace \((\omega_0,\omega_1)\) with \(K(\omega_0,\omega_1)>0\), choosing any \(\omega_0\in\Omega\) and any \(\omega_1\) with \(K(\omega_0,\omega_1)>0\), which exists because \(K(\omega_0,\cdot)\) sums to \(1\), and replay it in the record's audit. Take \(I_{\mathrm{PICA}}\) at level \(\mathsf{Ver}\), with strengths \(\{\texttt{weak}\}\), every admitted pair and claim type sent to \(\texttt{weak}\), an evidence policy admitting \(\texttt{committed\_state}\), \(\mathsf{ProvisionalTypes}_{I_{\mathrm{PICA}}}=\varnothing\), threshold policy consisting of the exact thresholds that the cell and pair records below carry, nonclaim record \(\mathcal N_{I_{\mathrm{PICA}}}=\mathcal N_{\mathrm{PICA}}\), and with scope and visible set containing, besides the host stores, the verdict references \(\ulcorner\mathsf{Cell}_{ij},I_{\mathrm{PICA}}\urcorner\) of the thirty-six cells, as for the running-example instrument of Subsection \hyperref[sn:3.8]{3.8}, with \(\mathsf H_{\mathrm{PICA}}\in\mathsf{Hosts}_{I_{\mathrm{PICA}}}\) and the level tags of the cell profiles in \(\mathsf{Levels}_{I_{\mathrm{PICA}}}\).

\emph{Cells.} Take pairwise distinct finite, typed witness and update records, on \(\mathsf H_{\mathrm{PICA}}\) and tagged \(\mathsf{Ver}\), for the thirty-six cells. Give each raw \(\texttt{Action}\) cell an update with signature in \(\mathsf{Upd}(P_i)\) (Subsubsection \hyperref[sn:3.7.5]{3.7.5}) whose recorded effect map changes its target (for \(P_6\), by adding a fresh log entry), so that \(\mathsf{noopU}\) is false; give each raw \(\texttt{Implicit}\) cell the same, with its witness attached to \(\mathcal T_{\mathrm{PICA}}\) and evaluated by no threshold, so that \(\mathsf{implicitW}\) is computed true; give each raw \(\texttt{Undefined}\) cell an update whose recorded effect changes its target and each raw \(\texttt{Trivial}\) cell a recorded identity effect; and give the audit record of each raw \(\texttt{Undefined}\) cell an evidence reference to that cell's own verdict under \(I_{\mathrm{PICA}}\), so that its computed \(\mathsf{CircularityCheck}\) fails. Give every cell a level/lens/interface record with the identity lens on \(\Omega\). Give every cell \(\lambda_{\mathrm{PICA}}\) and every pair \(\mu_{\mathrm{PICA}}\); step 4 and the profile clause of Subsubsection \hyperref[sn:4.4.1]{4.4.1} then hold, and no bridge is required. Choose each update's signature from \(\mathsf{Upd}(P_i)\) (Subsubsection \hyperref[sn:3.7.5]{3.7.5}). For \(P_5\), register a fresh fixed-point record in the first three groups and re-register an existing one in the trivial group. Thus no cell has a package-collapse entry. All cells are in scope with witness and update present, and every other check passes with its computed defect. The classifier then returns \(\texttt{action}\), \(\texttt{implicit}\), \(\texttt{trivial}\) and \(\texttt{undefined\_circular}\) on these four groups, reproducing the thirty-six raw statuses with no discrepancy. The construction treats each cell by its raw status alone, so it applies to every assignment of raw statuses. Record each raw PICA status, its formal verdict and the empty discrepancy in \(\mathsf{StatusStore}\).

\emph{Sources and pairs.} Give every witness, update, cell, dependency and ablation audit a \(\texttt{committed\_state}\) source record admitted by \(\mathsf{SourcePolicy}_{I_{\mathrm{PICA}}}\). Every record built in this proof is on \(\mathsf H_{\mathrm{PICA}}\) and tagged \(\mathsf{Ver}\). Give each of the fifteen pairs a bus entry, a producer edge, a passing provenance audit, \(\mathsf{partialEvidence}=\mathrm{false}\), an empty visibility defect, no required bridge, and a visible branch-compatibility record with \(\mathsf{compatible}=\mathrm{true}\) and its own passing audit; the \(\mathsf{SourceRefs}\) of this audit and of the pair's audit are \(\{R\}\), for the pair's source record \(R\) chosen below, so clause (viii) of \(\mathsf{AdmPICAReal}\) is met by the choice of \(R\) alone. Choose these entries so that each pair's source record \(R\) satisfies \(\mathsf{SourceOfTruthOK}(R,Q)\): \(R\) is committed on \(\mathsf{Bus}\), produced along the recorded edge, correctly provenanced, visible, audited and, if fallback, allowed by the host fallback policy; hence \(\mathrm{trace}_{\mathrm{src}}(R)\) records no failed check and \(\delta_6^{\mathrm{source}}(R)\neq\texttt{dirty}\). Choose the remaining pair data by raw pair status. A raw \(\texttt{real}\) pair gets an admitted \(\texttt{committed\_state}\) source and \(\mathsf{forbidsSource}=\mathrm{false}\), so the \(\texttt{real}\) row of Subsubsection \hyperref[sn:5.6.3]{5.6.3} holds; a raw \(\texttt{duplicated}\) pair gets the same together with a duplicate-pair record \(\mathsf{Dup}=(Q',A_{\mathsf{Dup}})\), where \(Q'\) is a copy of that pair record with a fresh identifier, no \(\mathsf{Dup}\) field and its own passing audit, stored with the audit evidence of \(A\) and not in \(\mathsf{PairStore}\); \(Q'\) is well formed, classified \(\texttt{real}\), and is not one of the fifteen realized pair records, and the \(\texttt{real\_duplicated}\) row holds for the stored record. A raw \(\texttt{provisional}\) pair gets \(\mathsf{forbidsSource}=\mathrm{false}\) and a fallback source record that \(\mathsf{SourcePolicy}_{I_{\mathrm{PICA}}}\) admits for provisional use under \(\mu\), with \(\delta_6^{\mathrm{source}}=\texttt{fallback\_admitted}\) and no source-upgrade bridge; condition (a) of the \(\texttt{real}\) row then fails, the first \(\texttt{blocked}\) row does not apply, and the \(\texttt{real\_provisional}\) row holds. Also give fallback full-mode admission under every profile and the default entry, and allow it under the host fallback policy; without an upgrade, it still yields \(\texttt{real\_provisional}\). A raw \(\texttt{blocked}\) pair gets an admitted \(\texttt{committed\_state}\) source and \(\mathsf{forbidsSource}=\mathrm{true}\), so the first \(\texttt{blocked}\) row holds. Every pair audit passes, and each pair's raw status, formal verdict and empty discrepancy are recorded in \(\mathsf{StatusStore}\). Include an audited finite ablation record and the nonclaims of Subsubsections \hyperref[sn:13.2.10]{13.2.10} and \hyperref[sn:13.6.2]{13.6.2}. These stores satisfy \(\mathsf{AdmPICAReal}\). Each cell so built passes steps 1--8 of Subsubsection \hyperref[sn:6.3.1]{6.3.1}, and only the witnesses of raw \(\texttt{Implicit}\) cells lie in \(W_{\mathcal T_{\mathrm{PICA}}}\).

\emph{Thresholds and \(K\)-independence.} Take every threshold of a raw \(\texttt{Action}\), \(\texttt{Trivial}\) or \(\texttt{Undefined}\) cell to be an exact check on the values of its recorded witness or update, every threshold of a raw \(\texttt{Implicit}\) cell an exact check on the values of its recorded update, and every threshold of a pair an exact check on its source record, each chosen so that the recorded values pass it (for instance \(\Theta_{\in S}\) with \(S\) the singleton of the recorded value); these are the thresholds meant above when every other check is said to pass, and give \(I_{\mathrm{PICA}}\) no rule other than the trace-support rule that takes \(K\) as an argument. Then \(K\) is named only by the \(K\) fields of \(\mathcal P\) and its host and by trace-support audit entries, and these stores are \(K\)-independent. For \(K\)-independent \(\mathcal P\), rerunning the trace-support checks under \(K'\) preserves their passing values, while all other checks and classifier inputs remain unchanged; hence admissibility and statuses are preserved.

\end{proof}

\subsubsection{PICA Nonclaims}\label{sn:13.2.10}

The PICA nonclaim register \(\mathcal N_{\mathrm{PICA}}\) contains, in particular:

\begin{enumerate}
\def\labelenumi{\arabic{enumi}.}
\tightlist
\item
  PICA is not the universal six-bird algebra: \(\mathsf{PICAReal}_{\mathrm{fin}}\) realizes a declared finite stochastic fragment, not the whole calculus.
\item
  The recovered \(6\times 6\) multiplicity count is a PICA-profile model fact under the cell profile \(\lambda_{\mathrm{PICA}}\), not a universal \(6\times 6\) law.
\item
  Raw PICA status is not formal status without recovery and classifier check; the formal classifier verdict is the authoritative cell status.
\item
  Directed-cell status \(\texttt{action}\) does not imply pair-realness; pair-realness requires the source/visibility/audit/threshold conjunction of Subsubsection \hyperref[sn:13.2.7]{13.2.7}.
\item
  A fallback record does not by itself support a real pair source-of-truth; it supports at most \(\texttt{real\_provisional}\).
\item
  Ablation evidence in \(\mathsf{Abl}\) is model-realization evidence, not theorem-grade evidence, in the absence of an admissible bridge to a theorem-grade claim.
\item
  A \(P_3\) route mismatch is not a \(\mathrm{P6}_{\mathrm{drive}}\) assertion in the sense of Section \hyperref[sn:8]{8}: the two are connected only through an admissible drive bridge.
\item
  A dirty record cannot support a real claim under any host configuration of \(\mathsf H_{\mathrm{PICA}}\).
\item
  PICA realizes a finite stochastic fragment, and not the entirety of \(\mathbf{BirdInt}^{\mathrm{aud}}_{\mathrm{fin}}\); this restates the fragment discipline of Subsubsection \hyperref[sn:13.1.2]{13.1.2}.
\item
  Implementation correctness of any concrete PICA implementation is future work and is not certified by this realization theorem; the theorem concerns only the finite realization record specified above.
\end{enumerate}

The register \(\mathcal N_{\mathrm{PICA}}\) is finite, recorded in \(\mathcal P\), and is part of the admissibility predicate \(\mathsf{AdmPICAReal}(\mathcal P)\) in the sense of Subsubsection \hyperref[sn:13.2.2]{13.2.2}.

\subsection{Cantor Strict-Bridge Realization}\label{sn:13.3}

This subsection records the Cantor strict-bridge realization sheet and Model Theorem \hyperref[cl:thm:34]{34}: under the audited Cantor shell, the strict-bridge passage from a base cocycle-pressure package \(T_0\) to a completion-packaged target package \(T_1\) is admissibly classified \(\texttt{strict}\). The realization is host-bound to the audited shell and fragment-scoped to the strict-bridge passage; it does not generalize beyond the shell, and its pressure-gap content is recorded as a consequence defect rather than the strictness witness itself. The realization packages the base and target theories, interfaces, bridge, pressure-gap defect, audits and nonclaims in a finite typed tuple over the audited Cantor shell; its strictness input comes from the cited Cantor theorem.

\emph{Orientation.} The Cantor model is the theorem package of \citet{Tsiokos2026Cantor}. Its substrate is a continuous Cantor system whose update is generated by six interacting mechanisms that play the roles of the six primitives, and its results hold on a distinguished class \(\mathcal S_{\mathrm{aud}}\) of lawful states, the audited shell. Two theories live on the shell. The base theory \(T_0\) is built from a cocycle and its pressure. The extension \(T_1\) is built from a packaging map that evolves a distribution, forgets detail, and reinstates a prototype. The Cantor paper proves that \(T_1\) is a strict extension of \(T_0\): the packaged-object map does not factor through the cocycle-level object map. This subsection records that result as a strict promotion bridge of the calculus. In the symbols below, \(a_n(s;x)\) is the cocycle observable at depth \(n\), parameter \(s\) and shell state \(x\); \(K\) is the normalized kernel of the substrate and \(\tau>0\) a timescale, so that \(\mu K^{\tau}\) is the state \(\mu\) evolved over the timescale \(\tau\); \(Q_\ell\) is the forgetting map of the lens \(\ell\) and \(U_\ell\) the matching reinstatement map; and \(\mathcal F\) is the family of persistent packaged strata. The pressure \(P_{T_0}(s)\) is defined by a limit in \(n\). It enters this paper only as a value recorded in the finite record of \(T_0\), established in the Cantor paper; no argument here takes the limit. In this subsection \(P_{T_0}\) and any other \(P\) with a package subscript are pressures, not role labels; \(W\) is the finite carrier, not a witness; and \(T_0,T_1\) name the Cantor theories, whose packages are \(\mathcal T^{\mathrm{Cantor}}_0,\mathcal T^{\mathrm{Cantor}}_1\).

\subsubsection{Name and Host}\label{sn:13.3.1}

The Cantor realization name is \(\mathsf{CantorShell}_{\mathrm{aud}}\). Its host is the audited Cantor shell context \(\mathsf H_{\mathrm{CantorShell}}\): a scoped audited application host built over a fixed audited shell-stable comparison domain \(\mathcal S_{\mathrm{aud}}\), recording the base and target packages, the bridge data, the pressure-gap defect record, and the audit/nonclaim register. The host is not the base finite theorem host of Sections \hyperref[sn:3]{3} and \hyperref[sn:4]{4}, but a scoped audited host whose records are judged by an instrument in the sense of Subsection \hyperref[sn:11.1]{11.1}; the host instrument is \(I_{\mathrm{Cantor}}\), an admissible instrument under which the realization claim of Subsubsection \hyperref[sn:13.3.8]{13.3.8} is recorded.

The realization object is the tuple of finite records

\[
\mathsf{CantorShell}_{\mathrm{aud}}
\;=\;
(\,\mathsf{cite}_{\mathcal S},\;W,\;T_0,\;T_1,\;\pi_0,\;\pi_1,\;\mathcal F,\;\Delta,\;A,\;\mathcal N\,),
\]

where \(\mathsf{cite}_{\mathcal S}\) is a finite citation record for the shell \(\mathcal S_{\mathrm{aud}}\) (an identifier and a pointer to its definition in the Cantor paper), \(T_0\) is the base cocycle-pressure theory with package \(\mathcal T_0^{\mathrm{Cantor}}\), \(T_1\) is the completion-packaged target theory with package \(\mathcal T_1^{\mathrm{Cantor}}\), \(\mathcal F\) is the persistent packaged-strata family, \(\Delta\) is the pressure-gap consequence record, which records the values of \(\Delta_{T_0\to T_1}\) and of \(P_{T_0}\) only at a declared finite set of parameters and at points of \(W\), each with a source record; each value is recorded as a closed interval with rational endpoints certified by the cited source, and no gate, threshold or check rule compares these values; \(A\) is the finite audit and support record, and \(\mathcal N\) is the Cantor nonclaim register of Subsubsection \hyperref[sn:13.3.9]{13.3.9}.

\subsubsection{\texorpdfstring{Base Package \(T_0\)}{Base Package T\_0}}\label{sn:13.3.2}

The base package is

\[
\mathcal T_0^{\mathrm{Cantor}}\;=\;(\,W,\;\pi_0,\;\Sigma_{\pi_0}^{\mathrm{cocycle}},\;E_0^{\mathrm{cocycle}},\;\mathcal A_{T_0}\,),
\]

a candidate package in the schema of Subsection \hyperref[sn:3.3]{3.3}, with:

\begin{itemize}
\tightlist
\item
  \(W\subseteq\mathcal S_{\mathrm{aud}}\), a finite audited subset of the audited shell-stable comparison domain that contains a witness pair of the nonfactorization of Subsubsection \hyperref[sn:13.3.5]{13.3.5}. The carrier \(W\) and the restricted object maps are finite, and the completion endomaps of the two packages are finite endomaps of \(\mathcal P(W)\), as below;
\item
  \(\pi_0:W\to\mathcal O_0\), the old cocycle-level object map onto a finite collection \(\mathcal O_0\) of cocycle-level objects;
\item
  \(\Sigma_{\pi_0}^{\mathrm{cocycle}}=\{\pi_0^{-1}(B):B\subseteq\pi_0(W)\}\), the lens-induced content of the base package; the cocycle-level claims that \(T_0\) exposes are recorded in its visibility profile \(V_{T_0}\);
\item
  \(E_0^{\mathrm{cocycle}}=E_{\pi_0}\), the lens completion \(\Sigma\mapsto\pi_0^{-1}(\pi_0(\Sigma))\) on \(\mathcal P(W)\), the canonical lens completion \(E_f\) of Subsubsection \hyperref[sn:3.3.1]{3.3.1}; the closure \(\operatorname{Cl}_{P_{T_0}}\) of the base pressure object is cited from the Cantor paper, not recorded;
\item
  \(\mathcal A_{T_0}\), the audit functional of Subsubsection \hyperref[sn:3.3.1]{3.3.1}, certifying each entry of the finite audit record family \(A_{T_0}\) of \(\mathcal T_0^{\mathrm{Cantor}}\) under \(I_{\mathrm{Cantor}}\).
\end{itemize}

The base pressure object is

\[
P_{T_0}(s)\;=\;\lim_{n\to\infty}\frac{1}{n}\,\log\!\left(\sup_{x\in\mathcal S_{\mathrm{aud}}}a_n(s;x)\right),
\]

a real-valued function of the parameter \(s\) whose values at the declared finite parameter set of Subsubsection \hyperref[sn:13.3.1]{13.3.1} are recorded in \(T_0\), each with a source record; the function itself is cited, not recorded.

\subsubsection{\texorpdfstring{Target Package \(T_1\)}{Target Package T\_1}}\label{sn:13.3.3}

The target package is

\[
\mathcal T_1^{\mathrm{Cantor}}\;=\;(\,W,\;\pi_1,\;\Sigma_{\pi_1}^{\mathrm{pkg}},\;E_1^{\mathrm{pkg}},\;\mathcal A_{T_1}\,),
\]

with \(\pi_1:W\to\mathcal O_1\) the new packaged-object map onto a finite collection \(\mathcal O_1\) of completion-packaged objects, \(\Sigma_{\pi_1}^{\mathrm{pkg}}=\{\pi_1^{-1}(B):B\subseteq\pi_1(W)\}\) the lens-induced content, with the packaged-object claims recorded in \(V_{T_1}\), \(\mathcal A_{T_1}\) the audit functional certifying the entries of the audit record family \(A_{T_1}\) of \(\mathcal T_1^{\mathrm{Cantor}}\), and completion endomap \(E_1^{\mathrm{pkg}}=E_{\pi_1}\), the lens completion \(\Sigma\mapsto\pi_1^{-1}(\pi_1(\Sigma))\) on \(\mathcal P(W)\). The packaging data of the Cantor system,

\[
(\,E_{\tau,\ell},\;\mathrm{Fix}(E_{\tau,\ell}),\;\mathcal F,\;\mathrm{Sat},\;\mathrm{Force}_{4\leftarrow 5}\,),
\]

are recorded in \(T_1\) only as finite citation records, each an identifier with a pointer to its definition in the Cantor paper, as for \(\mathsf{cite}_{\mathcal S}\); the operators themselves act on distributions of the Cantor system, not on \(\mathcal P(W)\), and are not records of the calculus. The level/lens/interface record of the bridge references these citation records. Here \(E_{\tau,\ell}(\mu)=U_{\ell}\bigl(Q_{\ell}(\mu K^{\tau})\bigr)\), \(\mathcal F\) the persistent packaged-strata family, \(\mathrm{Sat}\) the saturation operator, and \(\mathrm{Force}_{4\leftarrow 5}\) the forcing operator that records the \(P_4\leftarrow P_5\) completion-packaging relation.

\subsubsection{Cantor Strict Bridge}\label{sn:13.3.4}

The Cantor strict bridge is

\[
B_{0\to 1}^{\mathrm{Cantor}}\;=\;(\,\mathcal T_0^{\mathrm{Cantor}},\;\mathcal T_1^{\mathrm{Cantor}},\;\pi_0,\;\pi_1,\;L_{\mathrm{Cantor}},\;V_{\mathrm{Cantor}},\;\Theta_{\mathrm{Cantor}},\;A_{\mathrm{Cantor}},\;\delta_{\mathrm{Cantor}},\;\mathcal N_{\mathrm{Cantor}}\,),
\]

a promotion bridge with strict object maps in the promotion-bridge schema of Subsection \hyperref[sn:10.1]{10.1}. Its threshold record is \(\Theta_{\mathrm{Cantor}}=\{\Theta_{\mathrm{strict}}\}\), the exact strictness threshold \(\Delta_{\mathrm{fact}}(\pi_0,\pi_1)\neq\varnothing\) of Subsubsection \hyperref[sn:5.4.1]{5.4.1}, so a strictness verdict is required (Subsubsection \hyperref[sn:10.2.1]{10.2.1}); it carries no threshold on \(\delta_1^{\mathrm{split}}\) or \(\delta_{\mathrm{stab}}\). Its expanded record (Subsubsection \hyperref[sn:4.5.2]{4.5.2}) has carrier \(S=W\), codomains \(\mathcal O_0,\mathcal O_1\), the gate record of Subsubsection \hyperref[sn:13.3.7]{13.3.7}, host tag \(\mathsf H_{\mathrm{CantorShell}}\) and bridge profile \(\lambda_{\mathrm{Cantor}}\); Model Theorem \hyperref[cl:thm:34]{34} constrains \(\lambda_{\mathrm{Cantor}}\) only through \(\mathsf{AdmProfile}\) and the requirement that its level tags lie in \(\mathsf{Levels}_{I_{\mathrm{Cantor}}}\), and its Moreover clause fixes one admissible value. Its level/lens/interface record \(L_{\mathrm{Cantor}}\) records the level pair \((\mathsf{Ver},\mathsf{Ver})\), the identity lens on \(W\), and the interface \((\pi_0,\pi_1)\), and it references the attached lift data \((E_{\tau,\ell},Q_\ell,U_\ell,\mathrm{Sat},\mathrm{Force}_{4\leftarrow5})\). The defect bundle \(\delta_{\mathrm{Cantor}}\) contains the computed strictness, no-smuggling, stability, descent and audit defects of the promotion-bridge schema and the typed pressure-gap consequence defect of Subsubsection \hyperref[sn:13.3.6]{13.3.6}. Citation records for the Cantor saturation, \(4\leftarrow 5\) forcing and macro defects are attached to \(T_1\) and referenced by \(L_{\mathrm{Cantor}}\); they feed no gate.

\subsubsection{Strictness Witness}\label{sn:13.3.5}

The strictness witness is the nonfactorization

\[
\pi_1\not\factor\pi_0\quad\text{on}\quad W,
\]

equivalently,

\[
\exists\,x,x'\in W:\;\pi_0(x)=\pi_0(x')\;\land\;\pi_1(x)\neq\pi_1(x').
\]

Under the hypotheses of its strict-extension theorem on the audited shell, the Cantor paper proves that no map of sets \(\phi\) from the cocycle-level objects to the packaged objects, with no regularity condition, satisfies \(\pi_1=\phi\circ\pi_0\) on \(\mathcal S_{\mathrm{aud}}\), where \(\pi_0,\pi_1\) here denote the object maps of the Cantor paper on \(\mathcal S_{\mathrm{aud}}\), whose restrictions to \(W\) are the maps of Subsubsections \hyperref[sn:13.3.2]{13.3.2}--\hyperref[sn:13.3.3]{13.3.3}; by Theorem \hyperref[cl:thm:12]{12}, which holds for arbitrary sets, such a pair exists in \(\mathcal S_{\mathrm{aud}}\), and \(W\) is chosen to contain it; conversely, nonfactorization on \(W\subseteq\mathcal S_{\mathrm{aud}}\) implies nonfactorization on \(\mathcal S_{\mathrm{aud}}\), so nothing is lost by passing to \(W\).

The factorization defect record is

\[
\Delta_{\mathrm{fact}}^{\mathrm{Cantor}}(\pi_0,\pi_1)\;=\;\bigl\{\,(x,x')\in W\times W\;:\;\pi_0(x)=\pi_0(x')\;\land\;\pi_1(x)\neq\pi_1(x')\,\bigr\},
\]

a finite witness set on the audited shell. The strictness gate of Subsection \hyperref[sn:10.2]{10.2} returns

\[
G_{\mathrm{strict}}\;=\;\texttt{strict\_pass}\quad\Longleftrightarrow\quad\Delta_{\mathrm{fact}}^{\mathrm{Cantor}}(\pi_0,\pi_1)\;\neq\;\varnothing,
\]

so the strictness witness is the nonempty factorization defect on the audited shell. The witness is host-bound: it is a property of \(\pi_0\), \(\pi_1\), and \(\mathcal S_{\mathrm{aud}}\), recorded under \(I_{\mathrm{Cantor}}\), and not a property of \(\pi_0\) and \(\pi_1\) outside the shell.

\subsubsection{Pressure-Gap Consequence Defect}\label{sn:13.3.6}

The Cantor strict bridge carries a pressure-gap consequence defect

\[
\Delta_{T_0\to T_1}(s;x)\;=\;P_{T_0}(s)\;-\;\sum_{\Sigma\in\mathcal F(x)}w_{\Sigma}(x)\,P_{\Sigma}(s),
\]

with \(\mathcal F(x)\) the persistent packaged-strata family at \(x\), \(w_{\Sigma}(x)\) the per-stratum weight, and \(P_{\Sigma}(s)\) the per-stratum pressure object. The defect record is

\[
\mathsf{Def}^{\mathrm{press}}_{0\to 1}(s;x)\;=\;(\,\delta_{\mathrm{press}},\;\Omega_{\mathrm{press}},\;\Theta_{\mathrm{press}},\;\mathrm{polarity}_{\mathrm{press}},\;W_{\mathrm{press}},\;A_{\mathrm{press}},\;V_{\mathrm{press}},\;\sigma_{\mathrm{press}}\,),
\]

with \(\delta_{\mathrm{press}}\) storing a certified closed rational interval enclosing \(\Delta_{T_0\to T_1}(s;x)\), and the remaining fields drawn from the defect-record schema of Subsection \hyperref[sn:5.2]{5.2}. The activation assertion below is imported from the cited source; it is not a numerical test of the enclosure.

The pressure-gap defect is a \emph{consequence} defect, not the strictness witness itself: it is activated when

\[
\Delta_{T_0\to T_1}(s;x)\;\not\equiv\;0
\]

holds on the audited shell or on a recorded audited witness set, and it records a quantitative pressure deficit associated with the strict bridge, without itself supplying the nonfactorization data. The fragment discipline of Subsection \hyperref[sn:13.1]{13.1} forbids reading the pressure-gap defect as the strictness witness, and Subsubsection \hyperref[sn:13.3.9]{13.3.9} records the corresponding nonclaim.

\subsubsection{Gate Results}\label{sn:13.3.7}

The bridge audit records an origin and role tag for every use of every record in \(\mathsf{Use}(B)\); its strictness comparison uses the defining object maps on the common carrier and uses no independent target-layer novelty assertion. For the Cantor strict bridge, the promotion gates of Subsection \hyperref[sn:10.2]{10.2} return the following verdicts under \(I_{\mathrm{Cantor}}\):

\begin{longtable}[]{@{}ll@{}}
\toprule
gate & result\tabularnewline
\midrule
\endhead
\(G_{\mathrm{suff}}\) & \(\texttt{pass}\)\tabularnewline
\(G_{\mathrm{desc}}\) & \(\texttt{not\_required}\)\tabularnewline
\(G_{\mathrm{stab}}\) & \(\texttt{not\_required}\)\tabularnewline
\(G_{\mathrm{ctrl}}\) & \(\texttt{pass}\)\tabularnewline
\(G_{\mathrm{nosmuggle}}\) & \(\texttt{pass}\)\tabularnewline
\(G_{\mathrm{vis}}\) & \(\texttt{pass}\)\tabularnewline
\(G_{\mathrm{audit}}\) & \(\texttt{pass}\)\tabularnewline
\(G_{\mathrm{strict}}\) & \(\texttt{strict\_pass}\)\tabularnewline
\(G_{\mathrm{locglob}}\) & \(\texttt{not\_required}\)\tabularnewline
\bottomrule
\end{longtable}

The descent gate \(G_{\mathrm{desc}}\) is \(\texttt{not\_required}\) because the bridge claims object-map strictness on the audited shell, not closed macro dynamics; the local-global gate \(G_{\mathrm{locglob}}\) is \(\texttt{not\_required}\) because the realization is bound to the shell. The stability gate \(G_{\mathrm{stab}}\) is \(\texttt{not\_required}\) because the bridge does not claim fixed-point targets; the fixed points of the Cantor packaging \(E_{\tau,\ell}\) are cited attached data, not a claim of the bridge. The strictness gate \(G_{\mathrm{strict}}\) returns \(\texttt{strict\_pass}\) from the nonempty factorization defect of Subsubsection \hyperref[sn:13.3.5]{13.3.5}. The remaining gates pass under the hypotheses of Model Theorem \hyperref[cl:thm:34]{34}, as its proof shows.

\subsubsection{Model Theorem 34: Cantor Strict-Bridge Realization}\label{sn:13.3.8}

The predicate \(\mathsf{AdmCantorShell}(\mathsf{CantorShell}_{\mathrm{aud}})\) holds when: (1) \(\mathcal T_0^{\mathrm{Cantor}}\) and \(\mathcal T_1^{\mathrm{Cantor}}\) are admissible packages in the sense of Subsection \hyperref[sn:3.3]{3.3}, with the finite carrier \(W\); (2) \(B_{0\to1}^{\mathrm{Cantor}}\) is a record in the promotion-bridge schema of Subsection \hyperref[sn:10.1]{10.1} (Subsubsection \hyperref[sn:13.3.4]{13.3.4}); (3) the factorization defect of Subsubsection \hyperref[sn:13.3.5]{13.3.5} is recorded on \(W\); (4) the gate record of Subsubsection \hyperref[sn:13.3.7]{13.3.7} is recorded; (5) the audit record \(A\) is admissible under \(I_{\mathrm{Cantor}}\) and passes \(\mathsf{AuditPolicy}_{I_{\mathrm{Cantor}}}\); (6) the nonclaim register \(\mathcal N_{\mathrm{Cantor}}\) of Subsubsection \hyperref[sn:13.3.9]{13.3.9}, containing the coverage nonclaims of Subsubsection \hyperref[sn:13.6.2]{13.6.2}, is recorded, and \(\mathsf{SourcePolicy}_{I_{\mathrm{Cantor}}}\) admits the source records of \(W\), \(T_0\) and \(T_1\); and (7) the source record of \(W\) lists, for each \(x\in W\), a finite identifier of the shell state \(x\) fixed by the cited Cantor paper and the values \(\pi_0(x),\pi_1(x)\) as defined in the Cantor paper, and \(A_{\mathrm{Cantor}}\) checks the recorded labels against these cited values.

\begin{claimtheorem}{Model Theorem 34 (Cantor Strict-Bridge Realization)}\phantomsection\label{cl:thm:34}

Suppose the Cantor realization object is admissible:

\[
\mathsf{AdmCantorShell}\bigl(\mathsf{CantorShell}_{\mathrm{aud}}\bigr)
\]

holds, the strictness witness of Subsubsection \hyperref[sn:13.3.5]{13.3.5} records a nonempty factorization defect on \(W\), \(\mathsf{Use}(B^{\mathrm{Cantor}}_{0\to1})\subseteq\mathsf{Visible}_{I_{\mathrm{Cantor}}}\), \(\mathsf{Suppressed}_{I_{\mathrm{Cantor}}}=\varnothing\), \(\mathsf H_{\mathrm{CantorShell}}\in\mathsf{Hosts}_{I_{\mathrm{Cantor}}}\), \(\texttt{promotion}\in\mathsf{ClaimTypes}_{I_{\mathrm{Cantor}}}\), \(\mathsf{Ver}\) and both level tags of \(\lambda_{\mathrm{Cantor}}\) lie in \(\mathsf{Levels}_{I_{\mathrm{Cantor}}}\), and \(\mathsf{AdmProfile}(\lambda_{\mathrm{Cantor}},I_{\mathrm{Cantor}},\mathcal T^{\mathrm{Cantor}}_j)\) for \(j=0,1\), the audit record \(A_{\mathrm{Cantor}}\) passes \(\mathsf{AuditPolicy}_{I_{\mathrm{Cantor}}}\), and the no-smuggling defect of \(B_{0\to1}^{\mathrm{Cantor}}\) (Subsubsection \hyperref[sn:5.4.2]{5.4.2}) is empty. Then the gate record of Subsubsection \hyperref[sn:13.3.7]{13.3.7} holds, and

\[
\mathsf{CantorShell}_{\mathrm{aud}}\;\models\;\mathbf{StrictBridge}_{0\to 1}^{\mathrm{audited\ shell}}.
\]

In particular, the bridge promotion under \(\mathcal T_0^{\mathrm{Cantor}}\) classifies as \(\texttt{strict}\):

\[
\Gamma\,;\;\mathcal T_0^{\mathrm{Cantor}}\,;\;I_{\mathrm{Cantor}}\;\vdash\;\mathrm{Promote}\bigl(B_{0\to 1}^{\mathrm{Cantor}}\bigr)\;:\;\texttt{strict}.
\]

Moreover, the defect hypothesis can be met whenever the hypotheses of the cited strict-extension theorem hold (Subsubsection \hyperref[sn:13.3.5]{13.3.5}): that theorem gives nonfactorization on \(\mathcal S_{\mathrm{aud}}\), so there are \(x,x'\in\mathcal S_{\mathrm{aud}}\) with \(\pi_0(x)=\pi_0(x')\) and \(\pi_1(x)\neq\pi_1(x')\), and any finite \(W\subseteq\mathcal S_{\mathrm{aud}}\) containing \(x\) and \(x'\), for example \(W=\{x,x'\}\), has \((x,x')\in\Delta_{\mathrm{fact}}^{\mathrm{Cantor}}(\pi_0,\pi_1)\). Take \(W=\{x,x'\}\) and write \(o:=\pi_0(x)=\pi_0(x')\), \(u:=\pi_1(x)\) and \(v:=\pi_1(x')\), so \(u\ne v\); record \(\mathcal O_0=\{o\}\), \(\mathcal O_1=\{u,v\}\), and \(\Sigma_{\pi_k}\), \(E_{\pi_k}\) as in Subsubsections \hyperref[sn:13.3.2]{13.3.2}--\hyperref[sn:13.3.3]{13.3.3}. Give every record of this construction (the packages, their source, audit and nonclaim records, the source record of \(W\), and the bridge with its threshold, audit, defect and gate records) host tag \(\mathsf H_{\mathrm{CantorShell}}\) and level tag \(\mathsf{Ver}\), and give the packages \(\texttt{committed\_state}\) source records citing the Cantor paper, and give \(W\) one \(\texttt{committed\_state}\) source record citing the nonfactorization theorem and naming \(x,x'\) as its witness pair. Make every field visible, let the audit record's check rules replay the comparisons of these finite labels, and tag every record of \(\mathsf{Use}(B)\) \(\texttt{other}/\texttt{other}\), except the table of \(\pi_1\) (\(\texttt{target}/\texttt{novelty}\), which kind 1 exempts); the source record of \(W\) is \(\texttt{other}/\texttt{other}\), since the theorem only selects \(W\) and the strictness comparison recomputes \(\Delta_{\mathrm{fact}}\) from the labels \(o,u,v\); the six no-smuggling tests are then empty by inspection. Then the object satisfies \(\mathsf{AdmCantorShell}\) and the remaining hypotheses. For this, record the bridge in its expanded form \((\mathrm{id},\mathcal T_0,\mathcal T_1,S=W,O_0,O_1,\pi_0,\pi_1,L,V,\Theta,A,\delta,\mathsf{GateResults},\mathcal N,\mathsf H_{\mathrm{CantorShell}},\lambda_{\mathrm{Cantor}})\) with \(\lambda_{\mathrm{Cantor}}=(\mathsf{Ver},\mathsf{Ver},\texttt{fine},\texttt{structural},\texttt{visibility},\texttt{default},\texttt{local})\), and give \(I_{\mathrm{Cantor}}\) the claim types \(\{\texttt{promotion}\}\), the mode \(\mathsf{ModePolicy}(\texttt{promotion})=\texttt{default}\), the entry \((\texttt{fine},\texttt{structural},\texttt{local},\mathsf H_{\mathrm{CantorShell}})\in\mathsf{Compat}_{I_{\mathrm{Cantor}}}\), check rules recomputing \(\pi_0\) and \(\pi_1\) on \(W\) and \(\Delta_{\mathrm{fact}}\), exact thresholds, \(\mathsf{Hosts}_{I_{\mathrm{Cantor}}}=\{\mathsf H_{\mathrm{CantorShell}}\}\) and \(\mathsf{Levels}_{I_{\mathrm{Cantor}}}=\{\mathsf{Ver}\}\), a bridge policy admitting the visibility type, and source and evidence policies admitting \(\texttt{committed\_state}\), as for \(I_{\mathrm{prom}}\) in Subsubsection \hyperref[sn:10.6.1]{10.6.1}; its level is \(\mathsf{Ver}\), its strengths are \(\{\texttt{weak}\}\) with every admitted pair and claim type sent to \(\texttt{weak}\), its audit policy requires replay of every recorded check result under its check rules, \(\mathsf{ProvisionalTypes}_{I_{\mathrm{Cantor}}}=\varnothing\), and its nonclaim record is \(\mathcal N_{\mathrm{Cantor}}\). With these data, the bridge has the two-point object-map pattern of Theorem \hyperref[cl:thm:22]{22}, with its witness pair supplied by the cited Cantor theorem. The strictness input is the cited theorem; the calculus checks only the finite record built on it.

\end{claimtheorem}

\begin{proof}

By Subsubsection \hyperref[sn:13.3.4]{13.3.4}, \(B_{0\to 1}^{\mathrm{Cantor}}\) is a finite typed strict-bridge record in the promotion-bridge schema of Subsection \hyperref[sn:10.1]{10.1}. Subsubsection \hyperref[sn:13.3.5]{13.3.5} records that the factorization defect \(\Delta_{\mathrm{fact}}^{\mathrm{Cantor}}\) is nonempty on \(W\), so \(G_{\mathrm{strict}}=\texttt{strict\_pass}\). The remaining gates are computed as follows. \(G_{\mathrm{suff}}\) passes by clauses (1) and (2) of \(\mathsf{AdmCantorShell}\) and the agreement of the recorded defect-fed gate values with those computed here, and of the required marks with \(\Theta_{\mathrm{Cantor}}=\{\Theta_{\mathrm{strict}}\}\). \(G_{\mathrm{vis}}\) passes because \(\mathsf{Suppressed}_{I_{\mathrm{Cantor}}}=\varnothing\) and every field is visible. \(G_{\mathrm{audit}}\) passes because \(A_{\mathrm{Cantor}}\) passes \(\mathsf{AuditPolicy}_{I_{\mathrm{Cantor}}}\). \(G_{\mathrm{nosmuggle}}\) passes because its defect is empty, and so does \(G_{\mathrm{ctrl}}\), whose failure is the first violation kind of that defect. \(G_{\mathrm{desc}}\), \(G_{\mathrm{stab}}\) and \(G_{\mathrm{locglob}}\) are \(\texttt{not\_required}\), since the bridge claims neither closed macro dynamics, nor fixed-point targets, nor a local-to-global packaging. The row \(\texttt{outside\_scope}\) does not apply: \(h_B=\mathsf H_{\mathrm{CantorShell}}\in\mathsf{Hosts}_{I_{\mathrm{Cantor}}}\), the level tags of \(\lambda_{\mathrm{Cantor}}\) and \(L_{\mathrm{Cantor}}\) lie in \(\mathsf{Levels}_{I_{\mathrm{Cantor}}}\), and \(\mathsf{Use}(B)\subseteq\mathsf{Visible}_{I_{\mathrm{Cantor}}}\subseteq\mathsf{Scope}_{I_{\mathrm{Cantor}}}\). The promotion classifier of Subsubsection \hyperref[sn:5.6.4]{5.6.4}, applied with the gate record of Subsubsection \hyperref[sn:13.3.7]{13.3.7}, returns \(\mathrm{Promote}(B_{0\to 1}^{\mathrm{Cantor}}):\texttt{strict}\). This is the status the realization records for the bridge, so clause 7 of Subsubsection \hyperref[sn:13.1.1]{13.1.1} holds. Clause 1 holds because the components of the \(0\to1\) fragment are the fields of \(B_{0\to1}^{\mathrm{Cantor}}\); clause 2 by the admissible audit record \(A\); clause 3 by clause (6) of \(\mathsf{AdmCantorShell}\); clause 4 because every field is visible; and clause 5 because \(\mathcal N_{\mathrm{Cantor}}\) contains the coverage nonclaims; the realized fragment contains no directed cell, so clause 6 concerns only the typed bridge record of Subsubsection \hyperref[sn:13.3.4]{13.3.4}. By the realization-map convention of Subsubsection \hyperref[sn:13.1.1]{13.1.1}, the model-realization claim therefore receives the model-grade verdict \(\texttt{model\_realization}\) on the declared structure \(\mathbf{StrictBridge}_{0\to 1}^{\mathrm{audited\ shell}}\).

\end{proof}

\subsubsection{Cantor Nonclaims}\label{sn:13.3.9}

The Cantor nonclaim register \(\mathcal N_{\mathrm{Cantor}}\) contains:

\begin{enumerate}
\def\labelenumi{\arabic{enumi}.}
\tightlist
\item
  The realization is scoped to the audited Cantor shell \(\mathcal S_{\mathrm{aud}}\); it is not a shell-general theorem and does not assert a strict-bridge passage outside the shell.
\item
  There is no shell-general strict-bridge theorem; extending the realization beyond \(\mathcal S_{\mathrm{aud}}\) requires a separately admissible bridge.
\item
  There is no broader external-class theorem; the realization does not assert structural facts about external classes of theories unless an admissible inter-theory bridge is supplied.
\item
  Strictness does not imply macro closure: the strictness witness of Subsubsection \hyperref[sn:13.3.5]{13.3.5} is object-map nonfactorization, not closure of macro dynamics.
\item
  Strictness does not imply drive: the realization does not supply a \(\mathrm{P6}_{\mathrm{drive}}\) assertion in the sense of Section \hyperref[sn:8]{8} without a separately admissible drive bridge.
\item
  The pressure-gap consequence defect of Subsubsection \hyperref[sn:13.3.6]{13.3.6} is not a Kullback--Leibler closure deficit unless an admissible bridge to the KL closure-deficit framework is recorded.
\item
  The pressure-gap defect is not the strictness witness; it is recorded alongside the witness; its values are cited from \citet{Tsiokos2026Cantor}, and this paper does not derive it from non-factorization.
\item
  The pressure-gap functional \(\Delta_{T_0\to T_1}\) is a conditional pressure disintegration on the persistent strata family \(\mathcal F\), and is not a direct stratumwise root separation.
\item
  Cantor is a model realization/application, not a universal strict-extension theory.
\item
  The target completion \(T_1\) is not merely repeated application of one fixed idempotent completion; it is the packaged completion of Subsubsection \hyperref[sn:13.3.3]{13.3.3}. The register also contains the three coverage nonclaims of Subsubsection \hyperref[sn:13.6.2]{13.6.2}.
\end{enumerate}

The register \(\mathcal N_{\mathrm{Cantor}}\) is finite, recorded in \(\mathsf{CantorShell}_{\mathrm{aud}}\), and is part of the admissibility predicate \(\mathsf{AdmCantorShell}\).

\subsection{Graph/Cohomology Realization}\label{sn:13.4}

This subsection records the graph/cohomology realization sheet \(\mathsf{CohGraph}_{\mathrm{fin}}\) and connects the cohomology theorems of Section \hyperref[sn:8]{8} to the model-realization convention of Subsection \hyperref[sn:13.1]{13.1}. The realization is host-bound to the finite graph/cohomology host and fragment-scoped to the drive-support fragment \(P_2/P_6^{H^1}\text{-drive}/P_3\text{-separation}\).

\subsubsection{Host}\label{sn:13.4.1}

The graph/cohomology realization name is \(\mathsf{CohGraph}_{\mathrm{fin}}\). Its host is the finite graph/cohomology context \(\mathsf H_{\mathrm{graph}}\), with realization object

\[
\mathsf{CohGraph}_{\mathrm{fin}}\;=\;(\,G,\;C^{0}(G;k),\;C^{1}(G;k),\;d,\;H^{1}(G;k),\;\mathsf{Cycles}(G),\;a,\;\mathsf{Gate},\;A,\;V,\;\mathcal N\,),
\]

with components:

\begin{itemize}
\tightlist
\item
  \(G=(V,E)\), a finite support graph (finite vertex set \(V\) and finite edge set \(E\));
\item
  \(C^{0}(G;k)\) and \(C^{1}(G;k)\), the finite-dimensional \(0\)- and \(1\)-cochain spaces over the coefficient field or declared free coefficient ring \(k\) of Subsubsection \hyperref[sn:8.1.1]{8.1.1};
\item
  \(d:C^{0}(G;k)\to C^{1}(G;k)\), the coboundary operator;
\item
  \(H^{1}(G;k)\), the first cohomology group of \(G\) with coefficients in \(k\);
\item
  \(\mathsf{Cycles}(G)=Z_1(G;k)\), the cycle space of \(G\), a free \(k\)-module of rank \(\beta_1(G)\) with the finite basis of fundamental cycles of a declared spanning forest (each \(z\in Z_1\) equals \(\sum_{e\notin T}z(e)\gamma_e\), since the difference is a cycle supported on the forest \(T\) and vanishes as in the proof of Theorem \hyperref[cl:thm:15]{15}; the \(\gamma_e\) are independent, \(e\) lying only in \(\gamma_e\));
\item
  \(a\in C^{1}(G;k)\), the affinity/drive cochain assigning a finite scalar to each oriented edge of \(G\);
\item
  \(\mathsf{Gate}\), a finite typed record of admissible edge or support deletions/restrictions \(G\rightsquigarrow G'\);
\item
  \(A\) and \(V\), the host audit and visibility records;
\item
  \(\mathcal N\), the cohomology nonclaim register.
\end{itemize}

The host instrument is \(I_{\mathrm{graph}}\), an admissible instrument under which all cohomology and gating data are audited. The realization object is admissible, written \(\mathsf{AdmCohGraph}(\mathsf{CohGraph}_{\mathrm{fin}})\), when \(G\), \(a\), \(\mathsf{Gate}\), \(A\), \(V\) and \(\mathcal N\) are finite records, \(C^0(G;k)\), \(C^1(G;k)\), \(H^1(G;k)\) and \(\mathsf{Cycles}(G)\) are recorded by declared finite bases (vertices, edges, fundamental cycles; \(H^1(G;k)\) through the isomorphism \(H^1(G;k)\cong k^{\beta_1(G)}\) of Theorem \hyperref[cl:thm:17]{17}), and \(A\) records the verifications of the cochain spaces, the coboundary, the cycle space, the cohomology, and the affinity cochain, each of these verifications passes, and every assigned component's audit is replayed under \(I_{\mathrm{graph}}\) with audit defect \(\texttt{audit\_passes}\).

\subsubsection{Realized Fragment}\label{sn:13.4.2}

The graph/cohomology realization realizes the \(\mathbf{BirdInt}^{\mathrm{aud}}_{\mathrm{fin}}\) fragment

\[
\mathbf{BirdInt}^{\mathrm{aud}}_{\mathrm{fin}}{}^{\,P_{2}/P_{6}^{H^{1}}\text{-drive}/P_{3}\text{-separation}},
\]

with realization map

\begin{longtable}[]{@{}ll@{}}
\toprule
\begin{minipage}[b]{0.26\columnwidth}\raggedright
primitive\strut
\end{minipage} & \begin{minipage}[b]{0.68\columnwidth}\raggedright
cohomology realization\strut
\end{minipage}\tabularnewline
\midrule
\endhead
\begin{minipage}[t]{0.26\columnwidth}\raggedright
\(P_{2}\)\strut
\end{minipage} & \begin{minipage}[t]{0.68\columnwidth}\raggedright
support gating / edge deletion / cycle-rank capacity\strut
\end{minipage}\tabularnewline
\begin{minipage}[t]{0.26\columnwidth}\raggedright
\(\mathrm{P6}_{\mathrm{drive}}^{H^1}\)\strut
\end{minipage} & \begin{minipage}[t]{0.68\columnwidth}\raggedright
cohomological drive certificate \([a]\neq 0\)\strut
\end{minipage}\tabularnewline
\begin{minipage}[t]{0.26\columnwidth}\raggedright
\(P_{6}\) audit face\strut
\end{minipage} & \begin{minipage}[t]{0.68\columnwidth}\raggedright
audit of cochains, cycles, source-of-truth, and visibility\strut
\end{minipage}\tabularnewline
\begin{minipage}[t]{0.26\columnwidth}\raggedright
\(P_{3}\)\strut
\end{minipage} & \begin{minipage}[t]{0.68\columnwidth}\raggedright
optional route/holonomy witness, kept separate from drive\strut
\end{minipage}\tabularnewline
\begin{minipage}[t]{0.26\columnwidth}\raggedright
\(P_{4}\)\strut
\end{minipage} & \begin{minipage}[t]{0.68\columnwidth}\raggedright
optional local/global cover or gluing of cochains (out of core scope)\strut
\end{minipage}\tabularnewline
\begin{minipage}[t]{0.26\columnwidth}\raggedright
\(P_{5}\)\strut
\end{minipage} & \begin{minipage}[t]{0.68\columnwidth}\raggedright
optional quotient/support package of graph state (out of core scope)\strut
\end{minipage}\tabularnewline
\bottomrule
\end{longtable}

The core sheet of the realization focuses on \(P_{2}\) and \(\mathrm{P6}_{\mathrm{drive}}^{H^1}\); the \(P_{4}\) and \(P_{5}\) assignments above are optional extensions, recorded for completeness but not part of the core fragment, and the \(P_{3}\) separation is enforced as a nonclaim of Subsubsection \hyperref[sn:13.4.4]{13.4.4}.

\subsubsection{Drive and No-Drive Certificates}\label{sn:13.4.3}

The cohomological drive certificate is the finite typed record

\[
\kappa_{6}^{\mathrm{drive}}(a)\;:\;[a]\neq 0\;\in\;H^{1}(G;k),
\]

equivalent, by Theorem \hyperref[cl:thm:17]{17}, to the existence of a cycle integral witness:

\[
[a]\;\neq\;0\quad\Longleftrightarrow\quad\exists\,\gamma\in\mathsf{Cycles}(G):\;\oint_{\gamma}a\;\neq\;0.
\]

The drive certificate is a \(\mathrm{P6}_{\mathrm{drive}}^{H^1}\) record in the sense of Section \hyperref[sn:8]{8}, distinct from a generic \(P_{6}\) audit record: it records that the affinity cochain \(a\) has a nonzero cohomology class, equivalently a nonzero cycle integral on at least one cycle.

Each of the following finite records, with the audit and source records of Subsubsection \hyperref[sn:8.1.2]{8.1.2}, yields the no-drive certificate of that subsubsection; the first two concern the cochain \(a\), the last two the graph \(G\) and hence every \(a\in C^{1}(G;k)\):

\[
[a]\;=\;0,\qquad a\;=\;d\phi\;\;\text{for some}\;\phi\in C^{0}(G;k),\qquad H^{1}(G;k)\;=\;0,\qquad\beta_{1}(G)\;=\;0.
\]

The four records are connected by finite implications recorded in Section \hyperref[sn:8]{8}: \(a=d\phi\Rightarrow\oint_{\gamma}a=0\) for every cycle \(\gamma\); \(\beta_{1}(G)=0\Rightarrow H^{1}(G;k)=0\); and \(H^{1}(G;k)=0\Rightarrow[a]=0\) for every \(a\in C^{1}(G;k)\). The gating realization

\[
G\;=\;(V,E)\;\rightsquigarrow\;G'\;=\;(V',E')
\]

records a \(P_{2}\) support restriction with cycle-rank loss defect

\[
\delta_{2}^{\mathrm{cycle}}(G,G')\;=\;\beta_{1}(G)\;-\;\beta_{1}(G')
\]

and forest threshold \(\Theta_{2}^{\mathrm{forest}}:\beta_{1}(G')=0\). When the gate sends \(G\) to a forest \(G'\), the consequence \(H^{1}(G';k)=0\) records the absence of any cohomological drive certificate on \(G'\): by Theorem \hyperref[cl:thm:18]{18}, no certificate \(\mathrm{P6}_{\mathrm{drive}}^{H^1}\) exists on \(G'\).

The cohomology theorems of Section \hyperref[sn:8]{8} --- forest no-drive (\(\beta_{1}(G)=0\Rightarrow H^{1}(G;k)=0\)), exact-form null-drive (\(a=d\phi\Rightarrow\oint_{\gamma}a=0\)), nonzero-affinity equivalence (\([a]\neq 0\iff\exists\gamma:\oint_{\gamma}a\neq 0\)), and gating-affinity suppression (gating to a forest kills \(\mathrm{P6}_{\mathrm{drive}}^{H^1}\)) --- are recorded under the realization-map convention as the host-bound model-grade record of the drive-support fragment. Their theorem-grade status is preserved from Section \hyperref[sn:8]{8}; the model-realization convention adds the host-bound record under \(\mathsf H_{\mathrm{graph}}\) and \(I_{\mathrm{graph}}\), with verdict

\[
\Gamma\,;\;\mathsf H_{\mathrm{graph}}\,;\;I_{\mathrm{graph}}\;\vdash\;\mathsf{CohGraph}_{\mathrm{fin}}\;\models\;\mathbf{BirdInt}^{\mathrm{aud}}_{\mathrm{fin}}{}^{\,P_{2}/P_{6}^{H^{1}}\text{-drive}/P_{3}\text{-separation}}\;:\;\texttt{model\_realization}.
\]

The realization claim is admissible when \(\mathsf{AdmCohGraph}\) holds and the clauses of Subsubsection \hyperref[sn:13.1.1]{13.1.1} are met: clause 1 by the components listed there for this fragment; clauses 2 and 4 by the audit record \(A\) and the visibility record \(V\); clause 3 because \(\mathsf{SourcePolicy}_{I_{\mathrm{graph}}}\) admits the source records of the graph, cochain and audit records; clause 5 by \(\mathcal N_{\mathrm{coh}}\), which contains the coverage nonclaims of Subsubsection \hyperref[sn:13.6.2]{13.6.2}; clause 6 by the signatures of the gate and audit records in \(\mathsf{Wit}(P_2)\) and \(\mathsf{Wit}(P_6)\) and of the drive or no-drive certificate record in \(\mathsf{Wit}(P_6)\) (Subsubsection \hyperref[sn:3.7.4]{3.7.4}); and clause 7 vacuously, since no directed cell, pair observable or promotion bridge is realized. The cohomology theorems of Section \hyperref[sn:8]{8} are imported as their theorem-grade \(\texttt{accepted}\) verdicts.

\subsubsection{Nonclaims}\label{sn:13.4.4}

The graph/cohomology nonclaim register \(\mathcal N_{\mathrm{coh}}\) contains:

\begin{enumerate}
\def\labelenumi{\arabic{enumi}.}
\tightlist
\item
  The drive certificate \(\mathrm{P6}_{\mathrm{drive}}^{H^1}\) with \([a]\neq 0\) in \(H^{1}(G;k)\) is one drive certificate family, not every possible drive notion.
\item
  A \(P_{3}\)-route/holonomy witness is not a \(\mathrm{P6}_{\mathrm{drive}}^{H^1}\) assertion in the sense of Section \hyperref[sn:8]{8} in the absence of an admissible bridge \(B_{3\to 6}^{\mathrm{drive}}\).
\item
  Gating to a forest kills cycle-supported drive, not all forms of activity on \(G\).
\item
  Exact \(1\)-cochains \(a=d\phi\) have zero cycle drive: \(\oint_{\gamma}a=0\) for every cycle \(\gamma\).
\item
  The condition \(H^{1}(G;k)=0\) is a no-drive certificate for cohomological drive on \(G\), and not a no-drive certificate for any other notion of activity.
\item
  Nonzero route holonomy on \(G\) can coexist with a zero drive cochain \(a\); the two are distinct records.
\item
  The graph/cohomology host \(\mathsf H_{\mathrm{graph}}\) is finite and scoped; it is not a universal cohomological framework.
\item
  Cohomology does not replace the PICA, Cantor, or abstract-interpretation realizations of Subsections \hyperref[sn:13.2]{13.2}, \hyperref[sn:13.3]{13.3}, and Section \hyperref[sn:12]{12}.
\item
  No edge-support claim about \(G\) is accepted under \(I_{\mathrm{graph}}\) without the visibility and audit conditions of \(V\) and \(A\).
\item
  Positive holonomy-to-drive claims --- that a positive route holonomy implies a positive drive --- require an admissible drive bridge \(B_{3\to 6}^{\mathrm{drive}}\) in the sense of Section \hyperref[sn:8]{8}.
\item
  The three coverage nonclaims of Subsubsection \hyperref[sn:13.6.2]{13.6.2}.
\end{enumerate}

The register \(\mathcal N_{\mathrm{coh}}\) is finite, recorded in \(\mathsf{CohGraph}_{\mathrm{fin}}\), and is part of \(\mathsf{AdmCohGraph}\).

\subsection{Combined Model-Realization Matrix}\label{sn:13.5}

This subsection collects the four realizations of Sections \hyperref[sn:12]{12} and \hyperref[sn:13]{13} into a single matrix, with one row per realization and one column per scoped attribute. The matrix records what each realization realizes, under what host, with what claim strength, and where the supporting record is located. It does not assert any cross-model fact beyond what the per-row realizations already record; in particular, it does not state a model-coverage theorem, in keeping with Subsection \hyperref[sn:13.6]{13.6}.

\subsubsection{Rows}\label{sn:13.5.1}

The rows of the combined matrix are the four model realizations of this paper:

\begin{itemize}
\tightlist
\item
  \emph{PICA}: \(\mathsf{PICAReal}_{\mathrm{fin}}\) of Subsection \hyperref[sn:13.2]{13.2};
\item
  \emph{Cantor}: \(\mathsf{CantorShell}_{\mathrm{aud}}\) of Subsection \hyperref[sn:13.3]{13.3};
\item
  \emph{Abstract interpretation}: \(\mathsf{AIFam}_{\mathrm{fin}}\) of Section \hyperref[sn:12]{12};
\item
  \emph{Graph/cohomology}: \(\mathsf{CohGraph}_{\mathrm{fin}}\) of Subsection \hyperref[sn:13.4]{13.4}.
\end{itemize}

\subsubsection{Columns}\label{sn:13.5.2}

The columns of the combined matrix are the per-realization attributes:

\begin{itemize}
\tightlist
\item
  \emph{Host}: the host context \(\mathsf H\) and instrument \(I\) under which the realization claim is recorded.
\item
  \emph{Realized fragment}: the declared structure \(\mathcal S\) realized by the model.
\item
  \emph{Main theorem}: the theorem-grade or model-grade theorem of the realization, with its number in this paper.
\item
  \emph{Blocked overclaims}: the principal nonclaims of the realization that the fragment discipline of Subsubsection \hyperref[sn:13.1.2]{13.1.2} enforces.
\item
  \emph{Appendix support}: the appendix sheet that records the long-form realization data for the model.
\end{itemize}

The combined matrix is

\begingroup\footnotesize

\begin{longtable}[]{@{}llllll@{}}
\toprule
\begin{minipage}[b]{0.08\columnwidth}\raggedright
model\strut
\end{minipage} & \begin{minipage}[b]{0.12\columnwidth}\raggedright
host / instrument\strut
\end{minipage} & \begin{minipage}[b]{0.22\columnwidth}\raggedright
realized fragment\strut
\end{minipage} & \begin{minipage}[b]{0.12\columnwidth}\raggedright
main theorem\strut
\end{minipage} & \begin{minipage}[b]{0.22\columnwidth}\raggedright
blocked overclaims\strut
\end{minipage} & \begin{minipage}[b]{0.08\columnwidth}\raggedright
appendix support\strut
\end{minipage}\tabularnewline
\midrule
\endhead
\begin{minipage}[t]{0.08\columnwidth}\raggedright
\(\mathsf{PICAReal}_{\mathrm{fin}}\)\strut
\end{minipage} & \begin{minipage}[t]{0.12\columnwidth}\raggedright
\(\mathsf H_{\mathrm{PICA}}\) / \(I_{\mathrm{PICA}}\)\strut
\end{minipage} & \begin{minipage}[t]{0.22\columnwidth}\raggedright
\(\mathbf{BirdInt}^{\mathrm{aud}}_{\mathrm{fin\text{-}stoch}}{}^{\,\mathrm{cell}/\mathrm{pair}/\mathrm{prov}}\)\strut
\end{minipage} & \begin{minipage}[t]{0.12\columnwidth}\raggedright
Model Theorem \hyperref[cl:thm:33]{33} (model-grade)\strut
\end{minipage} & \begin{minipage}[t]{0.22\columnwidth}\raggedright
no PICA universality; no \(6\times 6\) universal law; raw status is not formal status\strut
\end{minipage} & \begin{minipage}[t]{0.08\columnwidth}\raggedright
Appendix \hyperref[sn:D.1]{D.1}\strut
\end{minipage}\tabularnewline
\begin{minipage}[t]{0.08\columnwidth}\raggedright
\(\mathsf{CantorShell}_{\mathrm{aud}}\)\strut
\end{minipage} & \begin{minipage}[t]{0.12\columnwidth}\raggedright
\(\mathsf H_{\mathrm{CantorShell}}\) / \(I_{\mathrm{Cantor}}\)\strut
\end{minipage} & \begin{minipage}[t]{0.22\columnwidth}\raggedright
\(\mathbf{StrictBridge}_{0\to 1}^{\mathrm{audited\ shell}}\)\strut
\end{minipage} & \begin{minipage}[t]{0.12\columnwidth}\raggedright
Model Theorem \hyperref[cl:thm:34]{34} (model-grade)\strut
\end{minipage} & \begin{minipage}[t]{0.22\columnwidth}\raggedright
no shell-general theorem; pressure-gap is consequence, not strictness witness; strictness does not imply drive\strut
\end{minipage} & \begin{minipage}[t]{0.08\columnwidth}\raggedright
Appendix \hyperref[sn:D.2]{D.2}\strut
\end{minipage}\tabularnewline
\begin{minipage}[t]{0.08\columnwidth}\raggedright
\(\mathsf{AIFam}_{\mathrm{fin}}\)\strut
\end{minipage} & \begin{minipage}[t]{0.12\columnwidth}\raggedright
\(\mathsf H_{\mathrm{AI}}\) / \(I_{\mathrm{AI}}\)\strut
\end{minipage} & \begin{minipage}[t]{0.22\columnwidth}\raggedright
\(\mathbf{BirdInt}^{\mathrm{aud}}_{\mathrm{fin}}{}^{\,P_{1}/P_{2}/P_{5}/P_{6}}\); with \(P_4\) given a strict refinement; with \(P_3\) given two noncommuting closures\strut
\end{minipage} & \begin{minipage}[t]{0.12\columnwidth}\raggedright
Model Theorem \hyperref[cl:thm:29]{29} (model-grade)\strut
\end{minipage} & \begin{minipage}[t]{0.22\columnwidth}\raggedright
no AI totality; no soundness-as-completeness; no closure-implies-soundness\strut
\end{minipage} & \begin{minipage}[t]{0.08\columnwidth}\raggedright
Appendix \hyperref[sn:D.3]{D.3}\strut
\end{minipage}\tabularnewline
\begin{minipage}[t]{0.08\columnwidth}\raggedright
\(\mathsf{CohGraph}_{\mathrm{fin}}\)\strut
\end{minipage} & \begin{minipage}[t]{0.12\columnwidth}\raggedright
\(\mathsf H_{\mathrm{graph}}\) / \(I_{\mathrm{graph}}\)\strut
\end{minipage} & \begin{minipage}[t]{0.22\columnwidth}\raggedright
\(\mathbf{BirdInt}^{\mathrm{aud}}_{\mathrm{fin}}{}^{\,P_{2}/P_{6}^{H^{1}}\text{-drive}/P_{3}\text{-separation}}\)\strut
\end{minipage} & \begin{minipage}[t]{0.12\columnwidth}\raggedright
Section \hyperref[sn:8]{8} theorem group plus CohGraph model-realization statement (model-grade)\strut
\end{minipage} & \begin{minipage}[t]{0.22\columnwidth}\raggedright
one drive certificate family, not all drive; \(P_{3}\) holonomy separated from drive\strut
\end{minipage} & \begin{minipage}[t]{0.08\columnwidth}\raggedright
Appendix \hyperref[sn:D.4]{D.4}\strut
\end{minipage}\tabularnewline
\bottomrule
\end{longtable}

\endgroup

The four rows record the four scoped realizations recognized by the calculus in this paper. Each row is read horizontally as a host-bound, fragment-scoped record; the matrix is not read vertically as a coverage assertion, in keeping with the cross-model role-coverage observation of Subsection \hyperref[sn:13.6]{13.6}.

Each row corresponds to a \(\models\) realization statement: Model Theorem \hyperref[cl:thm:33]{33} (Subsubsection \hyperref[sn:13.2.9]{13.2.9}) for PICA, Model Theorem \hyperref[cl:thm:34]{34} (Subsubsection \hyperref[sn:13.3.8]{13.3.8}) for Cantor, Model Theorem \hyperref[cl:thm:29]{29} (Subsection \hyperref[sn:12.5]{12.5}) for abstract interpretation, and the realization statement of Subsubsection \hyperref[sn:13.4.3]{13.4.3} for graph/cohomology. The instrument-indexed forms of the three model theorems, with verdict \(\texttt{model\_realization}\), are listed in Appendix \hyperref[sn:I]{I}. The matrix collects these four claims in one place; it does not introduce a new claim. The blocked-overclaims column is a summary view of the per-realization nonclaim registers \(\mathcal N_{\mathrm{PICA}}\), \(\mathcal N_{\mathrm{Cantor}}\), \(\mathcal N_{\mathrm{AI}}\), and \(\mathcal N_{\mathrm{coh}}\), and is not a separate, free-standing nonclaim register.

\subsection{Cross-Model Role Coverage}\label{sn:13.6}

This subsection records the cross-model role-coverage observation: which primitives \(P_{1},\ldots,P_{6}\) each realization realizes within its declared fragment, and which it does not. The observation is not a coverage theorem; the calculus does not assert that the union of the four realizations covers the six-role calculus, and Subsubsection \hyperref[sn:13.6.2]{13.6.2} records the corresponding coverage nonclaim.

\subsubsection{Coverage Table}\label{sn:13.6.1}

The role-coverage table records, per primitive, what each realization realizes within its declared fragment:

\begin{longtable}[]{@{}lllll@{}}
\toprule
\begin{minipage}[b]{0.10\columnwidth}\raggedright
role\strut
\end{minipage} & \begin{minipage}[b]{0.17\columnwidth}\raggedright
\(\mathsf{PICAReal}_{\mathrm{fin}}\)\strut
\end{minipage} & \begin{minipage}[b]{0.19\columnwidth}\raggedright
\(\mathsf{CantorShell}_{\mathrm{aud}}\)\strut
\end{minipage} & \begin{minipage}[b]{0.23\columnwidth}\raggedright
\(\mathsf{AIFam}_{\mathrm{fin}}\)\strut
\end{minipage} & \begin{minipage}[b]{0.17\columnwidth}\raggedright
\(\mathsf{CohGraph}_{\mathrm{fin}}\)\strut
\end{minipage}\tabularnewline
\midrule
\endhead
\begin{minipage}[t]{0.10\columnwidth}\raggedright
\(P_{1}\)\strut
\end{minipage} & \begin{minipage}[t]{0.17\columnwidth}\raggedright
descent / closure update\strut
\end{minipage} & \begin{minipage}[t]{0.19\columnwidth}\raggedright
macro obstruction; cited, not realized\strut
\end{minipage} & \begin{minipage}[t]{0.23\columnwidth}\raggedright
sound descent \(\alpha F\le_{A}F^{\sharp}\alpha\)\strut
\end{minipage} & \begin{minipage}[t]{0.17\columnwidth}\raggedright
not primary\strut
\end{minipage}\tabularnewline
\begin{minipage}[t]{0.10\columnwidth}\raggedright
\(P_{2}\)\strut
\end{minipage} & \begin{minipage}[t]{0.17\columnwidth}\raggedright
gating / support / admissibility\strut
\end{minipage} & \begin{minipage}[t]{0.19\columnwidth}\raggedright
admissibility obstruction (cited, not realized)\strut
\end{minipage} & \begin{minipage}[t]{0.23\columnwidth}\raggedright
representability \(\rho(p)=p\)\strut
\end{minipage} & \begin{minipage}[t]{0.17\columnwidth}\raggedright
support gating\strut
\end{minipage}\tabularnewline
\begin{minipage}[t]{0.10\columnwidth}\raggedright
\(P_{3}\)\strut
\end{minipage} & \begin{minipage}[t]{0.17\columnwidth}\raggedright
route mismatch\strut
\end{minipage} & \begin{minipage}[t]{0.19\columnwidth}\raggedright
not central\strut
\end{minipage} & \begin{minipage}[t]{0.23\columnwidth}\raggedright
noncommuting closures \(\Delta_3^{\mathrm{comp}}(\rho_k,\rho_\ell)\), when two closures fail to commute\strut
\end{minipage} & \begin{minipage}[t]{0.17\columnwidth}\raggedright
separated from drive\strut
\end{minipage}\tabularnewline
\begin{minipage}[t]{0.10\columnwidth}\raggedright
\(P_{4}\)\strut
\end{minipage} & \begin{minipage}[t]{0.17\columnwidth}\raggedright
refinement / staging\strut
\end{minipage} & \begin{minipage}[t]{0.19\columnwidth}\raggedright
\(P_{4}\!\leftarrow\!P_{5}\) completion-packaging (cited, not realized)\strut
\end{minipage} & \begin{minipage}[t]{0.23\columnwidth}\raggedright
refinement \(A_{k}\rightsquigarrow A_{\ell}\)\strut
\end{minipage} & \begin{minipage}[t]{0.17\columnwidth}\raggedright
optional local/global cover\strut
\end{minipage}\tabularnewline
\begin{minipage}[t]{0.10\columnwidth}\raggedright
\(P_{5}\)\strut
\end{minipage} & \begin{minipage}[t]{0.17\columnwidth}\raggedright
package / completion\strut
\end{minipage} & \begin{minipage}[t]{0.19\columnwidth}\raggedright
completion-packaged objects\strut
\end{minipage} & \begin{minipage}[t]{0.23\columnwidth}\raggedright
closure / packaging \(\rho=\gamma\alpha\)\strut
\end{minipage} & \begin{minipage}[t]{0.17\columnwidth}\raggedright
optional graph package\strut
\end{minipage}\tabularnewline
\begin{minipage}[t]{0.10\columnwidth}\raggedright
\(P_{6}\)\strut
\end{minipage} & \begin{minipage}[t]{0.17\columnwidth}\raggedright
audit / provenance\strut
\end{minipage} & \begin{minipage}[t]{0.19\columnwidth}\raggedright
audit of shell bridge\strut
\end{minipage} & \begin{minipage}[t]{0.23\columnwidth}\raggedright
audit of maps and soundness\strut
\end{minipage} & \begin{minipage}[t]{0.17\columnwidth}\raggedright
drive certificate \([a]\neq 0\) and audit\strut
\end{minipage}\tabularnewline
\bottomrule
\end{longtable}

A cell of the table records what the column's realization carries for the row's primitive \emph{within its declared fragment}. The phrase ``not primary'' records that the primitive is not part of the realization's core declared fragment. For \(\mathsf{PICAReal}_{\mathrm{fin}}\) the entries are the actor-update sorts of Subsubsection \hyperref[sn:13.2.3]{13.2.3}, and for \(\mathsf{CantorShell}_{\mathrm{aud}}\) records or citations of the bridge; neither builds role-channel records classified \(\texttt{active\_projection}\), so these columns are not role-fragment realizations in the sense of Subsubsection \hyperref[sn:13.1.1]{13.1.1}, unlike the \(\mathsf{AIFam}_{\mathrm{fin}}\) column.

The table reads horizontally and vertically only as a record of per-cell facts; it does not assert any cross-column equivalence of realizations of the same primitive, and it does not assert any aggregation across columns.

\subsubsection{Coverage Nonclaim}\label{sn:13.6.2}

The cross-model role-coverage observation does not, by itself, support a model-coverage theorem for \(\mathbf{BirdInt}^{\mathrm{aud}}_{\mathrm{fin}}\). The framework records the following coverage nonclaim:

\begin{itemize}
\item
  \emph{No model-coverage theorem.} The calculus does not assert that the four realizations together realize the whole calculus. Each cell of the table records one realization's datum for one role, within that realization's declared fragment; the union of the cells is not lifted to a coverage statement.
\item
  \emph{No cross-model identity.} Two cells in the same row carrying entries for the same primitive \(P_{i}\) are not identified: the PICA realization of \(P_{6}\) as audit/provenance and the cohomology realization of \(P_{6}\) as a drive certificate are distinct realizations of distinct facets of the audit role, and the matrix does not assert their identity. Cross-model identification of any per-primitive realization requires an admissible instrument-transfer bridge of Subsection \hyperref[sn:11.9]{11.9}.
\item
  \emph{Realization is not theorem grade.} The coverage table is a model-grade summary view. The \(\texttt{accepted}\) verdicts on theorem-grade claims of the calculus are recorded in Sections \hyperref[sn:6]{6}, \hyperref[sn:7]{7}, \hyperref[sn:8]{8}, \hyperref[sn:10]{10}, and \hyperref[sn:11]{11}; they are not displaced or summarized by the coverage table.
\end{itemize}

The coverage nonclaim is an instance of the fragment discipline of Subsubsection \hyperref[sn:13.1.2]{13.1.2} and the model-realization nonclaim of Subsection \hyperref[sn:2.4]{2.4}: the calculus records the per-row \(\texttt{model\_realization}\) verdicts of Subsection \hyperref[sn:13.5]{13.5}, and does not promote them, jointly or severally, to a coverage statement about \(\mathbf{BirdInt}^{\mathrm{aud}}_{\mathrm{fin}}\).
 \section{Secondary High-Structure Semantics}\label{sn:14}

This section defines, in the setting of double categories \citep{Ehresmann1963Doubles,GrandisPare1999}, a finite decorated-square semantics for \(\mathbf{BirdInt}^{\mathrm{aud}}_{\mathrm{fin}}\) and proves three high-structure recovery theorems: descent data appear as decorated squares (Theorem \hyperref[cl:thm:30]{30}), noncommuting completions give a completion-square defect (Theorem \hyperref[cl:thm:31]{31}), and strict extensions appear as no-filler conditions (Theorem \hyperref[cl:thm:32]{32}). The semantics is \emph{secondary}: it is an organizational layer over the core calculus, recording how finite categorical decorations recover descent, completion, and strict-extension data, and is not a replacement for the core calculus or its status families. Subsection \hyperref[sn:14.1]{14.1} fixes the finite double-category fragment, Subsections \hyperref[sn:14.2]{14.2}--\hyperref[sn:14.4]{14.4} prove the three theorems, and Subsection \hyperref[sn:14.5]{14.5} records the high-structure nonclaims.

The discipline of the section is the secondary-organization discipline: the decorated-square fragment supplies a finite, audit-checkable representation of selected core constructions, and does not assert that the core calculus is reducible to the decorated-square fragment. Status families remain separate, in keeping with the no-total-algebra and status-family-separation results of Section \hyperref[sn:6]{6}.

\subsection{Finite Double-Category Fragment}\label{sn:14.1}

This subsection records the finite double-category data used by the section: package objects as objects, two classes of finite maps as horizontal and vertical arrows, and the decorated squares carrying descent, quotient, completion, promotion, visibility, audit, and factorization data. The fragment is finite throughout: the object class, horizontal class, vertical class, and square class are all finite typed records audited by the host instrument \(I_{\mathrm{DC}}\).

\subsubsection{Objects, Horizontal Arrows, Vertical Arrows}\label{sn:14.1.1}

The finite double-category fragment \(\mathsf H_{\mathrm{DC}}\) is the finite typed record

\[
\mathsf H_{\mathrm{DC}}\;=\;(\,\mathrm{Ob},\;\mathrm{Hor},\;\mathrm{Ver},\;\mathrm{Sq},\;A_{\mathrm{DC}},\;V_{\mathrm{DC}},\;\mathcal N_{\mathrm{DC}}\,),
\]

with:

\begin{itemize}
\tightlist
\item
  \(\mathrm{Ob}\), the finite class of \emph{package objects}: each object is a finite typed package \(\mathcal T\) in the schema of Subsection \hyperref[sn:3.3]{3.3}, recording carrier, object map, visible interface, closure, and audit record;
\item
  \(\mathrm{Hor}\), the finite class of \emph{horizontal arrows}: each horizontal arrow is a finite typed map within or between packages, such as an update, a refinement, a lens, a package morphism in the sense of Subsection \hyperref[sn:3.3]{3.3}, or a promotion bridge when the square type requires one;
\item
  \(\mathrm{Ver}\), the finite class of \emph{vertical arrows}: each vertical arrow is a finite typed map within or between packages, such as a quotient, a completion or a transformation, with decorations such as audit, source and visibility drawn from the witness/update discipline of Subsection \hyperref[sn:3.7]{3.7}. Here \(\mathrm{Ver}\) names the vertical arrows, not the verification level \(\mathsf{Ver}\).
\item
  \(\mathrm{Sq}\), the finite class of \emph{decorated squares}, defined in Subsubsection \hyperref[sn:14.1.2]{14.1.2};
\item
  \(A_{\mathrm{DC}}\), the host audit record, recording verifications that \(\mathrm{Ob}\), \(\mathrm{Hor}\), and \(\mathrm{Ver}\) are finite and that their boundary and composition data are admissible;
\item
  \(V_{\mathrm{DC}}\), the host visibility record;
\item
  \(\mathcal N_{\mathrm{DC}}\), the host nonclaim register of Subsection \hyperref[sn:14.5]{14.5}.
\end{itemize}

Which maps sit on which side depends on the square: in the descent square of Theorem \hyperref[cl:thm:30]{30} the update \(F\) is horizontal and the quotient \(q\) vertical; in the completion square of Theorem \hyperref[cl:thm:31]{31} one completion is horizontal and the other vertical; in the factorization square of Theorem \hyperref[cl:thm:32]{32} the two object maps are horizontal.

An arrow at a package object \(X\) is a map on the carrier of \(X\); \(F:X\to X\) abbreviates such a map. A boundary set that is not the carrier of an object of \(\mathrm{Ob}\), such as \(\mathrm{im}(q)\) in Theorem \hyperref[cl:thm:30]{30}, is adjoined to \(\mathrm{Ob}\) as the carrier of an identity-lens package, as Theorem \hyperref[cl:thm:32]{32} does for \(\mathrm{im}(\pi_0)\) and \(O_1\).

Composition of horizontal arrows is along shared targets/sources, and composition of vertical arrows is along shared targets/sources, with both compositions admitted only when the resulting boundary records are finite and audited under \(I_{\mathrm{DC}}\). \(\mathsf H_{\mathrm{DC}}\) is not required to satisfy the double-category axioms (identities, associativity, interchange); the name refers only to the shape of its squares.

\subsubsection{Decorated Squares}\label{sn:14.1.2}

A decorated square \(\Xi\in\mathrm{Sq}\) is a finite typed record

\[
\Xi\;=\;(\,X,\;F,\;q,\;F^{\sharp},\;\mathrm{src}_{\Xi},\;\mathrm{tgt}_{\Xi},\;\mathrm{type}_{\Xi},\;\mathsf{Decor}_{\Xi},\;A_{\Xi},\;V_{\Xi},\;\delta_{\Xi},\;\mathcal N_{\Xi},\;\mathsf{partialEvidence}_{\Xi}\,),
\]

with:

\begin{itemize}
\tightlist
\item
  \(X\), the source object of the square (an element of \(\mathrm{Ob}\));
\item
  \(F\) and \(q\), the top horizontal edge and the vertical edge data of the square's boundary, drawn from \(\mathrm{Hor}\) and \(\mathrm{Ver}\): the vertical data is either a single arrow \(q\), used on both sides, or a pair \(q=(v_{\mathrm{left}},v_{\mathrm{right}})\) of vertical arrows;
\item
  \(F^{\sharp}\), the second horizontal arrow forming the bottom edge, drawn from \(\mathrm{Hor}\);
\item
  \(\mathrm{src}_{\Xi}\) and \(\mathrm{tgt}_{\Xi}\), the source and target boundary tuples;
\item
  \(\mathrm{type}_{\Xi}\in\{\,\texttt{descent},\;\texttt{quotient},\;\texttt{completion},\;\texttt{promotion},\;\texttt{visibility},\;\texttt{audit},\;\texttt{factorization}\,\}\), the square type;
\item
  \(\mathsf{Decor}_{\Xi}\), the decoration record carrying the type-specific data of the square: descent data for \(\texttt{descent}\), quotient data for \(\texttt{quotient}\), completion data for \(\texttt{completion}\), promotion data for \(\texttt{promotion}\), visibility-bridge data for \(\texttt{visibility}\), audit-bridge data for \(\texttt{audit}\), and factorization data for \(\texttt{factorization}\);
\item
  \(A_{\Xi}\) and \(V_{\Xi}\), the per-square audit and visibility records, with \(\mathsf{ClaimRef}(A_{\Xi})=\Xi\);
\item
  \(\delta_{\Xi}\), the per-square defect record drawn from the general defect-record schema of Subsection \hyperref[sn:5.2]{5.2} and classified in the square-status family of Subsection \hyperref[sn:5.1]{5.1};
\item
  \(\mathcal N_{\Xi}\), the per-square nonclaim register;
\item
  \(\mathsf{partialEvidence}_{\Xi}\in\{\mathrm{true},\mathrm{false}\}\), the instance's declaration that the evidence for the square is incomplete (as for claims, Subsubsection \hyperref[sn:11.1.2]{11.1.2}).
\end{itemize}

The four boundary edges must match: \(F\) runs from the source of \(v_{\mathrm{left}}\) to the source of \(v_{\mathrm{right}}\), and \(F^{\sharp}\) from the target of \(v_{\mathrm{left}}\) to the target of \(v_{\mathrm{right}}\). The square \emph{commutes} when \(v_{\mathrm{right}}\circ F=F^{\sharp}\circ v_{\mathrm{left}}\). A \emph{frame} is such a boundary with one edge left open, and a \emph{filler} of a frame is an arrow of the appropriate kind that, placed on the open edge, makes the square exact. A decorated square is (mathematically) \emph{exact} when the per-type exactness condition recorded in \(\mathsf{Decor}_{\Xi}\) holds and \(\delta_{\Xi}=\varnothing\). The exactness condition is type-specific: for \(\texttt{descent}\), it is the descent commutation equality of Subsection \hyperref[sn:14.2]{14.2}; for \(\texttt{completion}\), it is the completion-commutation equality of Subsection \hyperref[sn:14.3]{14.3}; for \(\texttt{factorization}\), it is commutation of the square of Subsection \hyperref[sn:14.4]{14.4}, with defect \(\delta_{\Xi}(\phi)=\{\,s:\phi(\pi_0(s))\ne\pi_1(s)\,\}\); for \(\texttt{visibility}\), it is \(\mathsf{AdmVisBridge}(L,I_{\mathrm{DC}},Y)\) of Subsubsection \hyperref[sn:5.5.2]{5.5.2}, where \(\mathsf{Decor}_\Xi\) records the bridge \(L\) and its target record \(Y\), of a claim, directed-cell, pair-observable or promotion-bridge schema; for \(\texttt{audit}\), it is \(\mathsf{AdmRotAuditBridge}\) of Subsubsection \hyperref[sn:11.7.3]{11.7.3} for the bridge recorded in \(\mathsf{Decor}_\Xi\); for these two types \(F\), \(q\) and \(F^{\sharp}\) are the identity map of the carrier of \(X\), adjoined to \(\mathrm{Hor}\) and \(\mathrm{Ver}\) for this purpose, so the boundary commutes and exactness is decided by the decoration alone; for \(\texttt{quotient}\) and \(\texttt{promotion}\), it is commutation \(v_{\mathrm{right}}\circ F=F^{\sharp}\circ v_{\mathrm{left}}\), with \(\delta_{\Xi}\) the set of points of the source object where it fails. The pasting law of Subsubsection \hyperref[sn:14.1.3]{14.1.3} applies to squares whose exactness condition is commutation. A decorated square is \emph{well formed} when its fields are finite and typed, its boundary edges match as above, its references resolve, and \(\delta_\Xi\) and its audit defect equal their recomputed values (Subsubsection \hyperref[sn:5.5.2]{5.5.2}); the square classifier of Subsubsection \hyperref[sn:5.6.6]{5.6.6} applies only to well-formed squares.

\subsubsection{Pasting}\label{sn:14.1.3}

The decorated-square fragment supports finite \emph{strict pasting}: given two decorated squares \(\Xi_{1}\) and \(\Xi_{2}\) that share a horizontal or vertical edge, their pasted square \(\Xi_{1}\cdot\Xi_{2}\) is recorded as a single decorated square with concatenated boundary and concatenated decoration. The pasted square is exact when its outer boundary commutes, and \(\delta_{\Xi_{1}\cdot\Xi_{2}}\) is the set of points where it fails to commute. For a horizontal paste with top edges \(F_1,F_2\),

\[
\delta_{\Xi_{1}\cdot\Xi_{2}}\;\subseteq\;\delta_{\Xi_{1}}\;\cup\;F_1^{-1}(\delta_{\Xi_{2}}),
\]

since at a point \(x\) with \(x\notin\delta_{\Xi_1}\) and \(F_1(x)\notin\delta_{\Xi_2}\) both squares commute along the path. For a vertical paste, with \(\Xi_1\) above \(\Xi_2\), left edges \(v_1,v_2\), right edges \(w_1,w_2\) and shared edge \(F_1^{\sharp}\), \(\delta_{\Xi_1\cdot\Xi_2}\subseteq\delta_{\Xi_1}\cup v_1^{-1}(\delta_{\Xi_2})\): if \(x\notin\delta_{\Xi_1}\) and \(v_1(x)\notin\delta_{\Xi_2}\), then \(w_2w_1F_1(x)=w_2F_1^{\sharp}v_1(x)=F_2^{\sharp}v_2v_1(x)\). The paste defect \(\delta_{\mathrm{paste}}(\Xi_{1},\Xi_{2})=\{e:\text{the shared edge }e\text{ is recorded differently in }\Xi_{1}\text{ and }\Xi_{2}\}\) is empty whenever pasting is admissible, and a nonempty paste defect makes the paste inadmissible. Pasting is admissible when the shared edge data agree on \(\mathrm{Hor}\) or \(\mathrm{Ver}\) and the per-square admissibility conditions hold for both factors. Strict pasting is the only pasting form admitted in this section: pseudo-pasting and weak pasting are out of scope, and the paste-defect record is the locus where any pasting failure is recorded.

The fragment \(\mathsf H_{\mathrm{DC}}\) together with the strict-pasting law is the finite decorated double-category fragment used by Theorems \hyperref[cl:thm:30]{30}, \hyperref[cl:thm:31]{31}, and \hyperref[cl:thm:32]{32} of Subsections \hyperref[sn:14.2]{14.2}, \hyperref[sn:14.3]{14.3}, and \hyperref[sn:14.4]{14.4}.

\subsection{Theorem 30: Descent-Square Recovery}\label{sn:14.2}

This subsection proves Theorem \hyperref[cl:thm:30]{30}: in the finite decorated double-category fragment of Subsection \hyperref[sn:14.1]{14.1}, descent data are recovered as decorated squares of type \(\texttt{descent}\), with exactness recorded as the commutation equality between the concrete map \(F\) and the quotient map \(q\) followed by the descended map \(F^{\sharp}\).

Let \(X\in\mathrm{Ob}\) be a finite package object, \(F\in\mathrm{Hor}\) a horizontal map \(F:X\to X\), \(q\in\mathrm{Ver}\) a quotient map \(q:X\to Y\), drawn through its corestriction \(q_{\mathrm{im}}:X\to\mathrm{im}(q)\), and \(F^{\sharp}\in\mathrm{Hor}\) the candidate descended horizontal map \(F^{\sharp}:\mathrm{im}(q)\to\mathrm{im}(q)\). Form the decorated square

\[
\Xi_{\mathrm{desc}}\;:\;
\begin{array}{ccc}
X & \xrightarrow{F} & X\\
{\scriptstyle q}\downarrow & \Downarrow\,\Xi_{\mathrm{desc}} & \downarrow{\scriptstyle q}\\
\mathrm{im}(q) & \xrightarrow{F^{\sharp}} & \mathrm{im}(q)
\end{array}
\]

with \(\mathrm{type}_{\Xi}=\texttt{descent}\) and \(\mathsf{Decor}_{\Xi}\) recording the descent data \((F,q,F^{\sharp})\).

\begin{claimtheorem}{Theorem 30 (DescentSquareRecovery)}\phantomsection\label{cl:thm:30}

Let \(\Xi_{\mathrm{desc}}\) be well formed (Subsubsection \hyperref[sn:14.1.2]{14.1.2}, with \(\delta_\Xi\) recomputed as in the square classifier of Subsubsection \hyperref[sn:5.6.6]{5.6.6}), so that \(\delta_\Xi=\delta_1^{\mathrm{cand}}(q,F,F^\sharp)\). The decorated square \(\Xi_{\mathrm{desc}}\) is exact iff the descent commutation equality

\[
F^{\sharp}\circ q_{\mathrm{im}}\;=\;q_{\mathrm{im}}\circ F
\]

holds on \(X\), where \(q_{\mathrm{im}}:X\to\mathrm{im}(q)\) is the corestriction of \(q\) to its image. Moreover, a map \(F^{\sharp}:\mathrm{im}(q)\to\mathrm{im}(q)\) with \(F^{\sharp}\circ q_{\mathrm{im}}=q_{\mathrm{im}}\circ F\), a \emph{set-map filler} of \(\Xi_{\mathrm{desc}}\) (a filler in the sense of Subsubsection \hyperref[sn:14.1.2]{14.1.2} when it also belongs to \(\mathrm{Hor}\)), exists iff the well-definedness condition

\[
q(x)\;=\;q(x')\quad\Longrightarrow\quad qF(x)\;=\;qF(x')
\]

holds for all \(x,x'\in X\). In the finite completion of \(\mathsf H_{\mathrm{DC}}\) that adjoins every map \(\mathrm{im}(q)\to\mathrm{im}(q)\) to \(\mathrm{Hor}\), with the corresponding candidate descent squares and their boundary, defect and audit records, every set-map filler is a filler, so a filler of \(\Xi_{\mathrm{desc}}\) exists there iff this condition holds.

\end{claimtheorem}

\begin{proof}

The proof has two parts: exactness of \(\Xi_{\mathrm{desc}}\) is equivalent to the commutation equality for the given \(F^{\sharp}\), and the existence of a set-map filler is equivalent to the well-definedness condition. Both parts use only the finite set-theoretic structure of \(X\) and \(\mathrm{im}(q)\).

\emph{Exactness.} The descent square is exact when its per-type condition, the commutation equality \(F^{\sharp}\circ q_{\mathrm{im}}=q_{\mathrm{im}}\circ F\), holds and its defect record is empty; for a descent square that record is the candidate defect \(\delta_1^{\mathrm{cand}}(q,F,F^\sharp)=\{x:F^\sharp(q_{\mathrm{im}}(x))\neq q_{\mathrm{im}}F(x)\}\) of Subsubsection \hyperref[sn:5.3.1]{5.3.1}, which is empty exactly when the commutation equality holds. So the square is exact iff the commutation equality holds. The split defect \(\delta_1^{\mathrm{split}}(q,F)\) (Subsubsection \hyperref[sn:5.3.1]{5.3.1}) is the obstruction in the filler clause: some set-map filler exists iff \(\delta_1^{\mathrm{split}}(q,F)=\varnothing\); in the completed host of the theorem, some filler exists iff \(\delta_1^{\mathrm{split}}(q,F)=\varnothing\).

\emph{Set-map filler from well-definedness.} Suppose \(q(x)=q(x')\Rightarrow qF(x)=qF(x')\). Define \(G:\mathrm{im}(q)\to\mathrm{im}(q)\) by \(G(y)=qF(x)\) for any \(x\in X\) with \(q(x)=y\). The well-definedness condition guarantees that the value does not depend on the choice of representative \(x\), so \(G\) is a well-defined function, and by construction \(G\circ q_{\mathrm{im}}=q_{\mathrm{im}}\circ F\): \(G\) is a set-map filler.

\emph{Well-definedness from a set-map filler.} Suppose \(G\circ q_{\mathrm{im}}=q_{\mathrm{im}}\circ F\) for some map \(G\), and suppose \(q(x)=q(x')\). Then \(q_{\mathrm{im}}(x)=q_{\mathrm{im}}(x')\), so \(G(q_{\mathrm{im}}(x))=G(q_{\mathrm{im}}(x'))\). Applying the equality on both sides gives \(q_{\mathrm{im}}F(x)=q_{\mathrm{im}}F(x')\), that is, \(qF(x)=qF(x')\).

The two parts together give both clauses of the theorem. Each step is a finite check on the finite host data of \(\mathsf H_{\mathrm{DC}}\), recorded in \(A_{\mathrm{DC}}\) and \(A_{\Xi}\).

The descent square thus recovers descent data inside the decorated double-category fragment: a descent record \((F,q,F^{\sharp})\) is exact in \(\mathsf H_{\mathrm{DC}}\) iff \(F^{\sharp}\) is the well-defined descended map of \(F\) along \(q\). The recovery is a representational result; it does not assert that every descent fact about the core calculus is reducible to a descent square, in keeping with the secondary-organization discipline of the section.

\end{proof}

\subsection{Theorem 31: Completion-Pasting Recovery}\label{sn:14.3}

This subsection proves Theorem \hyperref[cl:thm:31]{31}: in the finite decorated double-category fragment of Subsection \hyperref[sn:14.1]{14.1}, the noncommutation of two finite completions is recorded as the square defect \(\delta_\Xi=\Delta_3^{\mathrm{comp}}(E_1,E_2)\) of Subsubsection \hyperref[sn:5.3.3]{5.3.3}, and the completion square is exact iff the two completions commute.

Let \(C\in\mathrm{Ob}\) be a finite package object and let \(E_{1},E_{2}:C\to C\) be two endomaps of \(C\) (for instance the completions of Subsection \hyperref[sn:7.4]{7.4}; Theorem \hyperref[cl:thm:31]{31} uses no further property of them), with \(E_{2}\) used as the horizontal edge and \(E_{1}\) used as the vertical edge of the square. Form the decorated square

\[
\Xi_{\mathrm{comp}}\;:\;
\begin{array}{ccc}
C & \xrightarrow{E_{2}} & C\\
{\scriptstyle E_{1}}\downarrow & \Downarrow\,\Xi_{\mathrm{comp}} & \downarrow{\scriptstyle E_{1}}\\
C & \xrightarrow{E_{2}} & C
\end{array}
\]

with \(\mathrm{type}_{\Xi}=\texttt{completion}\), \(\mathsf{Decor}_{\Xi}\) recording the completion pair \((E_{1},E_{2})\), and \(\delta_{\Xi}:=\Delta_3^{\mathrm{comp}}(E_1,E_2)\).

\begin{claimtheorem}{Theorem 31 (CompletionPastingRecovery)}\phantomsection\label{cl:thm:31}

The decorated square \(\Xi_{\mathrm{comp}}\) is exact, in the sense of Subsubsection \hyperref[sn:14.1.2]{14.1.2}, iff the two completions commute, that is,

\[
E_{1}\,E_{2}\;=\;E_{2}\,E_{1}\;\;\text{on}\;\;C,
\]

and its defect record is \(\delta_\Xi=\Delta_3^{\mathrm{comp}}(E_1,E_2)\). If moreover \(\Xi_{\mathrm{comp}}\) is well formed (Subsubsection \hyperref[sn:14.1.2]{14.1.2}, with \(\delta_\Xi\) recomputed as in the square classifier of Subsubsection \hyperref[sn:5.6.6]{5.6.6}) and in scope under \(I_{\mathrm{DC}}\), its visibility defect is empty, its audit defect is \(\texttt{audit\_passes}\), no partial evidence is declared and its decoration declares no order comparison, then the square classifier of Subsubsection \hyperref[sn:5.6.6]{5.6.6} assigns \(\texttt{obstructed}\) iff there exists \(c\in C\) with \(E_{1}E_{2}(c)\neq E_{2}E_{1}(c)\), and \(\texttt{exact}\) otherwise. In the noncommuting case the noncommutation defect

\[
\Delta_3^{\mathrm{comp}}\bigl(E_{1},E_{2}\bigr)\;=\;\bigl\{\,c\in C\;:\;E_{1}E_{2}(c)\neq E_{2}E_{1}(c)\,\bigr\}
\]

is a finite typed defect record under \(I_{\mathrm{DC}}\), and the noncommutation \(E_{1}E_{2}\neq E_{2}E_{1}\) is recorded as the square defect \(\delta_\Xi=\Delta_3^{\mathrm{comp}}(E_1,E_2)\), the \(P_3\) completion-commutator defect of Subsubsection \hyperref[sn:5.3.3]{5.3.3}. For a third completion \(E_3\) on \(C\), the horizontal paste \(\Xi_{\mathrm{comp}}(E_1,E_2)\cdot\Xi_{\mathrm{comp}}(E_1,E_3)\), in which the right vertical edge \(E_1\) of the first square is the left vertical edge of the second, so that the outer square has top and bottom edges \(E_3E_2\) and both vertical edges \(E_1\), has defect

\[
\{\,c\in C:E_1E_3E_2(c)\neq E_3E_2E_1(c)\,\}\;\subseteq\;\Delta_3^{\mathrm{comp}}(E_1,E_2)\;\cup\;E_2^{-1}\bigl(\Delta_3^{\mathrm{comp}}(E_1,E_3)\bigr),
\]

and this inclusion can be strict.

\end{claimtheorem}

\begin{proof}

The decorated square \(\Xi_{\mathrm{comp}}\) has both horizontal edges equal to \(E_{2}\) and both vertical edges equal to \(E_{1}\). Exactness in the sense of Subsubsection \hyperref[sn:14.1.2]{14.1.2} is the commutation equality of the boundary composition: the upper-right path \(E_{1}\circ E_{2}:C\to C\) equals the lower-left path \(E_{2}\circ E_{1}:C\to C\). By pointwise equality of maps on the finite carrier \(C\), this is the equation \(E_{1}E_{2}=E_{2}E_{1}\) on \(C\).

\emph{Forward direction.} Suppose \(\Xi_{\mathrm{comp}}\) is exact. Then by definition \(E_{1}E_{2}(c)=E_{2}E_{1}(c)\) for every \(c\in C\), hence \(E_{1}E_{2}=E_{2}E_{1}\) on \(C\).

\emph{Reverse direction.} Suppose \(E_{1}E_{2}=E_{2}E_{1}\) on \(C\). Then the boundary composition of \(\Xi_{\mathrm{comp}}\) commutes pointwise on the finite carrier \(C\), so by the exactness criterion of Subsubsection \hyperref[sn:14.1.2]{14.1.2} the square is exact.

For the obstruction direction, suppose there exists \(c_{0}\in C\) with \(E_{1}E_{2}(c_{0})\neq E_{2}E_{1}(c_{0})\). The set

\[
\Delta_3^{\mathrm{comp}}(E_{1},E_{2})\;=\;\bigl\{\,c\in C:E_{1}E_{2}(c)\neq E_{2}E_{1}(c)\,\bigr\}
\]

is a finite subset of \(C\) containing \(c_{0}\), and is recorded as the finite typed defect record \(\delta_{\Xi}\) of the square. Since \(\delta_{\Xi}\neq\varnothing\), the completion square is not exact. Under the additional hypotheses of the status clause, the scope, visibility, audit and circularity rows of the square classifier do not apply, so it assigns \(\texttt{obstructed}\) exactly when \(\delta_\Xi\ne\varnothing\), and otherwise, with no order comparison and no partial evidence, \(\texttt{exact}\). As in the completion-pasting defect of Subsection \hyperref[sn:7.5]{7.5}, this noncommutation of two finite completions is the route/pasting witness for \(P_{3}\). Hence \(E_{1}E_{2}\neq E_{2}E_{1}\) is recovered as a \(P_{3}\)-pasting defect.

Each step is a finite check on the finite host data of \(\mathsf H_{\mathrm{DC}}\), recorded in \(A_{\mathrm{DC}}\) and \(A_{\Xi}\).

The completion square records noncommutation through \(\delta_\Xi=\Delta_3^{\mathrm{comp}}(E_1,E_2)\); commuting completions yield exact squares.
For the paste, if \(s\notin\Delta_3^{\mathrm{comp}}(E_1,E_2)\) and \(E_2s\notin\Delta_3^{\mathrm{comp}}(E_1,E_3)\), then \(E_1E_3E_2s=E_3E_1E_2s=E_3E_2E_1s\). The inclusion is strict for the running completions \(E_1,E_2\) with \(E_3\) the constant map to \(\{a,b,c\}\): \(E_3\) is idempotent and commutes with \(E_1\), both outer composites are the constant map to \(\{a,b,c\}\), so the paste has empty defect, while \(\Delta_3^{\mathrm{comp}}(E_1,E_2)\) contains \(\{a\}\).

\end{proof}

\subsection{Theorem 32: Strict Extension as No-Filler}\label{sn:14.4}

This subsection proves Theorem \hyperref[cl:thm:32]{32}: in the finite decorated double-category fragment of Subsection \hyperref[sn:14.1]{14.1}, a strict extension between two object maps \(\pi_{0},\pi_{1}\) on a common carrier \(S\) is recovered as the absence of a factorization filler in a decorated square of type \(\texttt{factorization}\). The theorem connects the strict-bridge nonfactorization witness of Section \hyperref[sn:7]{7} (and its Cantor instance in Subsection \hyperref[sn:13.3]{13.3}) to the no-filler condition in \(\mathsf H_{\mathrm{DC}}\).

Let \(S\) be a finite carrier object and let \(\pi_{0}\in\mathrm{Hor}\) and \(\pi_{1}\in\mathrm{Hor}\) be two horizontal object maps \(\pi_{0}:S\to O_{0}\) and \(\pi_{1}:S\to O_{1}\) with finite codomains. Write \(\mathrm{im}(\pi_{0})\) for the effective old interface, and let \(\phi:\mathrm{im}(\pi_{0})\to O_{1}\) be a candidate vertical filler map. For this factorization construction, take the finite completion of \(\mathsf H_{\mathrm{DC}}\) that adjoins interface package objects for \(\mathrm{im}(\pi_0)\) and \(O_1\), the corestriction \(S\to\mathrm{im}(\pi_0)\) to \(\mathrm{Hor}\), \(\mathrm{id}_S\) and every map \(\phi:\mathrm{im}(\pi_0)\to O_1\) to \(\mathrm{Ver}\), and the corresponding candidate factorization squares, with their boundary, defect and audit records, to \(\mathrm{Sq}\). In \(\Xi_{\mathrm{fact}}(\phi)\) the top edge labelled \(\pi_0\) is this corestriction. This completion is finite because the carriers are finite. Every \(\Xi_{\mathrm{fact}}(\phi)\) in Theorem \hyperref[cl:thm:32]{32} is evaluated in this completed host. Form the decorated factorization square

\[
\Xi_{\mathrm{fact}}(\phi)\;:\;
\begin{array}{ccc}
S & \xrightarrow{\pi_{0}} & \mathrm{im}(\pi_{0})\\
{\scriptstyle\mathrm{id}_{S}}\downarrow & \Downarrow\,\Xi_{\mathrm{fact}}(\phi) & \downarrow{\scriptstyle\phi}\\
S & \xrightarrow{\pi_{1}} & O_{1}
\end{array}
\]

with \(\mathrm{type}_{\Xi}=\texttt{factorization}\) and \(\mathsf{Decor}_{\Xi}\) recording the factorization candidate \((\pi_{0},\pi_{1},\phi)\).

\begin{claimtheorem}{Theorem 32 (StrictExtensionAsNoFiller)}\phantomsection\label{cl:thm:32}

For every \(\phi:\mathrm{im}(\pi_{0})\to O_{1}\), with \(\Xi_{\mathrm{fact}}(\phi)\) well formed (Subsubsection \hyperref[sn:14.1.2]{14.1.2}, with \(\delta_\Xi\) recomputed as in the square classifier of Subsubsection \hyperref[sn:5.6.6]{5.6.6}), so that \(\delta_\Xi(\phi)=\{s:\phi(\pi_0(s))\neq\pi_1(s)\}\), the factorization square \(\Xi_{\mathrm{fact}}(\phi)\) is exact, that is, \(\phi\) is a filler of its frame, iff

\[
\pi_{1}\;=\;\phi\circ\pi_{0}\;\;\text{on}\;\;S.
\]

Hence

\[
\pi_{1}\not\factor\pi_{0}\quad\Longleftrightarrow\quad\text{no exact factorization filler exists for}\;(\pi_{0},\pi_{1}),
\]

with \(\mathrm{Fill}(\pi_{0},\pi_{1})\;=\;\bigl\{\,\phi:\mathrm{im}(\pi_{0})\to O_{1}\;:\;\pi_{1}=\phi\circ\pi_{0}\,\bigr\}\) and

\[
\mathrm{Fill}(\pi_{0},\pi_{1})\;=\;\varnothing\quad\Longleftrightarrow\quad\exists\,s,s'\in S:\;\pi_{0}(s)=\pi_{0}(s')\;\land\;\pi_{1}(s)\neq\pi_{1}(s').
\]

\end{claimtheorem}

\begin{proof}

The proof translates the nonfactorization witness \(\pi_{1}\not\factor\pi_{0}\) of Subsection \hyperref[sn:7.7]{7.7} into the no-filler condition for \(\Xi_{\mathrm{fact}}\).

\emph{Filler iff factorization.} Suppose \(\phi:\mathrm{im}(\pi_{0})\to O_{1}\) satisfies \(\pi_{1}=\phi\circ\pi_{0}\) on \(S\). Then the boundary composition of \(\Xi_{\mathrm{fact}}(\phi)\) commutes pointwise: for every \(s\in S\), the upper-right path returns \(\phi(\pi_{0}(s))=\pi_{1}(s)\), and the lower-left path returns \(\pi_{1}(\mathrm{id}_{S}(s))=\pi_{1}(s)\). By the exactness criterion of Subsubsection \hyperref[sn:14.1.2]{14.1.2}, \(\Xi_{\mathrm{fact}}(\phi)\) is exact, that is, \(\phi\) is a filler. Conversely, if \(\phi\) is a filler, then by definition the boundary composition commutes pointwise, so \(\pi_{1}(s)=\phi(\pi_{0}(s))\) for every \(s\in S\), that is, \(\pi_{1}=\phi\circ\pi_{0}\) on \(S\).

\emph{No filler iff nonfactorization witness.} Suppose \(\mathrm{Fill}(\pi_{0},\pi_{1})=\varnothing\). By the equivalence above, no \(\phi\) satisfies \(\pi_{1}=\phi\circ\pi_{0}\). Equivalently, the assignment \(\pi_{0}(s)\mapsto\pi_{1}(s)\) is not a well-defined function on \(\mathrm{im}(\pi_{0})\), so there exist \(s,s'\in S\) with \(\pi_{0}(s)=\pi_{0}(s')\) and \(\pi_{1}(s)\neq\pi_{1}(s')\). Conversely, suppose there exist such \(s,s'\in S\). Then any candidate \(\phi\) with \(\pi_{1}=\phi\circ\pi_{0}\) would require \(\phi(\pi_{0}(s))=\pi_{1}(s)\) and \(\phi(\pi_{0}(s))=\phi(\pi_{0}(s'))=\pi_{1}(s')\), contradicting \(\pi_{1}(s)\neq\pi_{1}(s')\). Hence \(\mathrm{Fill}(\pi_{0},\pi_{1})=\varnothing\).

The two equivalences together yield the theorem: \(\pi_{1}\not\factor\pi_{0}\) on \(S\) in the sense of Subsection \hyperref[sn:7.7]{7.7} holds iff no exact filler exists for \(\Xi_{\mathrm{fact}}\), iff there exist \(s,s'\in S\) with \(\pi_{0}(s)=\pi_{0}(s')\) and \(\pi_{1}(s)\neq\pi_{1}(s')\). Each step is a finite check on the finite host data of \(\mathsf H_{\mathrm{DC}}\), recorded in \(A_{\mathrm{DC}}\) and \(A_{\Xi}\).

The factorization square thus records the strict-extension nonfactorization witness inside the decorated double-category fragment: a strict extension is exactly the absence of an exact filler for the factorization square. The Cantor strict-bridge realization of Subsection \hyperref[sn:13.3]{13.3} supplies a concrete instance: the finite witness set \(W\) of the audited shell carries \(\pi_{1}\not\factor\pi_{0}\) on \(W\), so \(\mathrm{Fill}(\pi_{0},\pi_{1})=\varnothing\) on \(W\), and \(\Xi_{\mathrm{fact}}\) has no filler there.

\end{proof}

\subsection{High-Structure Nonclaims}\label{sn:14.5}

This subsection records the high-structure nonclaims that fix the boundaries of Section \hyperref[sn:14]{14}: the decorated-square fragment is secondary, the status family discipline is preserved, and the general propagation of pasting defects is future work.

\subsubsection{No Double-Category Completeness}\label{sn:14.5.1}

The finite decorated double-category fragment \(\mathsf H_{\mathrm{DC}}\) of Subsection \hyperref[sn:14.1]{14.1} does not replace the core calculus \(\mathbf{BirdInt}^{\mathrm{aud}}_{\mathrm{fin}}\). Theorems \hyperref[cl:thm:30]{30}, \hyperref[cl:thm:31]{31}, and \hyperref[cl:thm:32]{32} record three specific recovery results --- descent data as descent squares, completion noncommutation as completion squares, and strict-extension nonfactorization as factorization squares --- and not a general claim that every theorem-grade fact about the calculus is reducible to a decorated-square fact.

In particular, the section does not claim:

\begin{itemize}
\tightlist
\item
  that \(\mathbf{BirdInt}^{\mathrm{aud}}_{\mathrm{fin}}\) embeds into \(\mathsf H_{\mathrm{DC}}\) as a sub-double-category;
\item
  that every directed-cell judgment of Subsection \hyperref[sn:4.3]{4.3} is recovered by a decorated square;
\item
  that the priority order of the classifier of Subsubsection \hyperref[sn:11.2.1]{11.2.1} is recovered by the per-square exactness conditions of Subsubsection \hyperref[sn:14.1.2]{14.1.2}.
\end{itemize}

The decorated-square fragment is a secondary organizational layer over the core calculus; it supplies finite, audit-checkable representations of selected core constructions where those representations are useful, and does not promote those representations to a completeness statement about the calculus.

\subsubsection{No Status Erasure}\label{sn:14.5.2}

The decorated-square fragment does not collapse the status families of Subsection \hyperref[sn:5.1]{5.1} into a single square-status family. The role-status, cell-status, pair-status, promotion-status, claim-status, gate-status, and square-status families remain separate judgment-indexed families, in keeping with the no-total-algebra and status-family-separation results of Section \hyperref[sn:6]{6}.

In particular, the per-square exactness conditions of Subsubsection \hyperref[sn:14.1.2]{14.1.2} do not assert that:

\begin{itemize}
\tightlist
\item
  every directed-cell status of Subsubsection \hyperref[sn:5.6.2]{5.6.2} is determined by the exactness of a decorated square that contains the cell;
\item
  every pair-observable status of Subsubsection \hyperref[sn:5.6.3]{5.6.3} is determined by the exactness of a decorated square that contains the pair;
\item
  every claim status of Subsubsection \hyperref[sn:11.2.1]{11.2.1} is determined by the exactness of a decorated square.
\end{itemize}

A decorated square carries its own square-status verdict only as part of a full decorated/audited record: witness and update data, language, visibility, threshold, audit, defect, square status, and nonclaim data must be present when the square is used as a \(\mathbf{BirdInt}^{\mathrm{aud}}_{\mathrm{fin}}\) claim. A bare commutative diagram is not an admissible claim. The square-status verdict does not export to the other status families, and the status-separation discipline of the calculus is preserved by Section \hyperref[sn:14]{14}.

\subsubsection{No General Obstruction Propagation Theorem}\label{sn:14.5.3}

The pasting law of Subsubsection \hyperref[sn:14.1.3]{14.1.3} bounds the defect of an admissible horizontal paste: \(\delta_{\Xi_1\cdot\Xi_2}\subseteq\delta_{\Xi_1}\cup F_1^{-1}(\delta_{\Xi_2})\), with \(\delta_{\mathrm{paste}}(\Xi_1,\Xi_2)=\varnothing\); the inclusion can be strict (Theorem \hyperref[cl:thm:31]{31}). This bound is a computation on the finite paste, not a propagation theorem.

In particular, Section \hyperref[sn:14]{14} does not claim a general theorem of the form

\begin{quote}
a defect at one square in a finite pasting diagram propagates to a defect at every square downstream of it,
\end{quote}

nor does it claim a general theorem of the form

\begin{quote}
the absence of a defect at one square in a finite pasting diagram suffices for the absence of a defect at every square in the diagram.
\end{quote}

Both directions of general defect propagation are future work. Specific propagation facts under specific pasting diagrams may be admissible under the per-diagram audit data of \(A_{\Xi_{1}\cdot\Xi_{2}}\), but no general propagation theorem is asserted by the calculus in this paper.

The three nonclaims above together fix the local boundary of the secondary high-structure semantics within \(\mathbf{BirdInt}^{\mathrm{aud}}_{\mathrm{fin}}\): the decorated-square fragment supplies finite recoveries of descent, completion, and strict-extension data, preserves the separation of the status families, and records pasted defects as finite aggregated records without a general propagation theorem.
 \section{Limitations, Nonclaims, and Future Work}\label{sn:15}

This section collects the final scope discipline of \(\mathbf{BirdInt}^{\mathrm{aud}}_{\mathrm{fin}}\) and the future-work boundaries of the paper. Subsection \hyperref[sn:15.1]{15.1} restates the positive contribution of the paper across six layers --- finite typed interaction calculus, statused profile-relative interaction, the theorem-grade finite core, the countermodel/no-go atlas, the instrument and self-reference scope, and the model-realization scope. Subsection \hyperref[sn:15.2]{15.2} records the nonclaims, numbered \(\mathrm{NC}\text{-}1\) through \(\mathrm{NC}\text{-}36\) and organized by the six topical groupings of the calculus. Subsection \hyperref[sn:15.3]{15.3} records the six classes of future work that the paper does not undertake.

The discipline of the section is uniform: the paper records exactly what it asserts and exactly what it does not assert, and treats the nonclaims and future-work entries as part of the paper's content rather than as exclusions of content. The section does not introduce any new theorem or model-realization claim; it collects, restates, and frames what the earlier sections have already recorded.

\subsection{Positive Scope Summary}\label{sn:15.1}

This subsection restates the positive contribution of the paper across the six structural layers of \(\mathbf{BirdInt}^{\mathrm{aud}}_{\mathrm{fin}}\).

\subsubsection{Finite Typed Interaction Calculus}\label{sn:15.1.1}

This paper constructs \(\mathbf{BirdInt}^{\mathrm{aud}}_{\mathrm{fin}}\), a finite typed interaction calculus over six primitive labels \(\mathbb P=\{P_{1},\ldots,P_{6}\}\), with a finite typed schema of directed-cell records, pair-observable records, promotion-bridge records, claim records, instrument records, source-of-truth records, audit records, and visibility records. The calculus is defined in Sections \hyperref[sn:3]{3} and \hyperref[sn:4]{4} as a finite typed schema with finite carriers, finite stores, and finite admissibility predicates on each component. Every record of the calculus is finite and the core calculus uses finite carriers; the macro-kernel class of Theorem \hyperref[cl:thm:8]{8} and the cochain and cycle spaces of Theorems \hyperref[cl:thm:15]{15}--\hyperref[cl:thm:18]{18}, over which the body statements quantify, may be infinite (Subsubsection \hyperref[sn:2.1.1]{2.1.1}). The Cantor realization uses finite records over a shell that is not asserted finite (Subsubsection \hyperref[sn:2.1.2]{2.1.2}). Every record is interpreted under an explicit host instrument and audit record.

\subsubsection{Statused Profile-Relative Interaction}\label{sn:15.1.2}

The second positive layer is the statused profile-relative interaction discipline of Section \hyperref[sn:5]{5}. Each judgment carries a status family: the role status of Subsubsection \hyperref[sn:5.1.1]{5.1.1}, the directed-cell status of Subsubsection \hyperref[sn:5.1.2]{5.1.2}, the pair-observable status of Subsubsection \hyperref[sn:5.1.3]{5.1.3}, the promotion status of Subsubsection \hyperref[sn:5.1.4]{5.1.4}, the claim status of Subsubsection \hyperref[sn:5.1.5]{5.1.5}, the gate status of Subsubsection \hyperref[sn:5.1.6]{5.1.6}, and the square status of Subsubsection \hyperref[sn:5.1.7]{5.1.7}. The status families are typed-separately judgment-indexed: a directed-cell status is not a pair status, a pair status is not a promotion status, and the priority orders within each family are frozen by the classifier of Subsection \hyperref[sn:5.6]{5.6}.

The same primitive pair can receive different statuses on different profiles: Theorem \hyperref[cl:thm:1]{1} exhibits this for \((P_6,P_3)\), and the regime variant of Subsubsection \hyperref[sn:3.6.3]{3.6.3} isolates the profile. Profile relativity is enforced by the per-cell profile \(\lambda\) of Subsection \hyperref[sn:4.3]{4.3} and reflected in Theorem \hyperref[cl:thm:1]{1}'s status non-factorization through the primitive pair (Section \hyperref[sn:6]{6}).

\subsubsection{Numbered Results}\label{sn:15.1.3}

The third positive layer is the numbered theorem and proposition spine of the paper. The numbered results are distributed as follows; Model Theorems \hyperref[cl:thm:29]{29}, \hyperref[cl:thm:33]{33} and \hyperref[cl:thm:34]{34} are model-grade and Proposition \hyperref[cl:prop:35]{35} is audit-grade (Appendix \hyperref[sn:I]{I}).

\begin{itemize}
\tightlist
\item
  Section \hyperref[sn:6]{6}: the no-total-algebra theorem (Theorem \hyperref[cl:thm:1]{1}) and the four core finite-interaction theorems (Theorems \hyperref[cl:thm:2]{2}--\hyperref[cl:thm:5]{5}).
\item
  Section \hyperref[sn:7]{7}: nine finite-domain theorems on descent, closure deficit, completion, saturation, strictness, and definability (Theorems \hyperref[cl:thm:6]{6}--\hyperref[cl:thm:14]{14}).
\item
  Section \hyperref[sn:8]{8}: four cohomological theorems on graph-supported drive (Theorems \hyperref[cl:thm:15]{15}--\hyperref[cl:thm:18]{18}).
\item
  Section \hyperref[sn:10]{10}: four promotion theorems (Theorems \hyperref[cl:thm:19]{19}--\hyperref[cl:thm:22]{22}).
\item
  Section \hyperref[sn:11]{11}: three self-reference and visibility theorems (Theorems \hyperref[cl:thm:23]{23}--\hyperref[cl:thm:25]{25}).
\item
  Section \hyperref[sn:12]{12}: the abstract-interpretation theorem trio (Theorems \hyperref[cl:thm:26]{26}--\hyperref[cl:thm:28]{28}) and the abstract-interpretation model-family theorem (Model Theorem \hyperref[cl:thm:29]{29}).
\item
  Section \hyperref[sn:14]{14}: the high-structure recovery theorems (Theorems \hyperref[cl:thm:30]{30}--\hyperref[cl:thm:32]{32}).
\item
  Section \hyperref[sn:13]{13}: the PICA and Cantor model-realization theorems (Model Theorems \hyperref[cl:thm:33]{33}--\hyperref[cl:thm:34]{34}), together with the graph/cohomology model-realization statement of Subsection \hyperref[sn:13.4]{13.4}.
\item
  Appendix \hyperref[sn:F]{F}: the proof-assistant validation audit report (Proposition \hyperref[cl:prop:35]{35}).
\end{itemize}

Each numbered statement has an instrument-indexed claim form and verdict, listed in Appendix \hyperref[sn:I]{I}: \(\texttt{accepted}\) for theorem-grade results, \(\texttt{model\_realization}\) for scoped model realizations, and \(\texttt{audit\_report}\) for the manifest-scoped Lean validation artifact.

\subsubsection{Countermodel/No-Go Scope}\label{sn:15.1.4}

The fourth positive layer is the countermodel atlas of Section \hyperref[sn:9]{9} and Appendix \hyperref[sn:C]{C}. The paper records twelve finite countermodels, each blocking a specific overclaim about the calculus: active package without macro closure (\(\mathsf{CM}_{1}\)), active route/holonomy without drive (\(\mathsf{CM}_{2}\)), local packages failing globally (\(\mathsf{CM}_{3}\)), directed cell active while the pair observable is blocked (\(\mathsf{CM}_{4}\)), refinement worsening closure (\(\mathsf{CM}_{5}\)), same-level self-audit circularity (\(\mathsf{CM}_{6}\)), same primitive pair with different statuses (\(\mathsf{CM}_{7}\)), idempotent completions that do not commute (\(\mathsf{CM}_{8}\)), strict extension without macro closure (\(\mathsf{CM}_{9}\)), strict extension without drive (\(\mathsf{CM}_{10}\)), fallback source unable to support a real pair observable (\(\mathsf{CM}_{11}\)), and gating that destroys cycle-supported drive (\(\mathsf{CM}_{12}\)).

Each countermodel is a finite typed structure under an instrument whose visible content blocks the corresponding overclaim under \(I\). The atlas is finite, audited, and host-bound, in keeping with the no-go scope discipline of Subsection \hyperref[sn:9.1]{9.1}.

\subsubsection{Instrument/Self-Reference Scope}\label{sn:15.1.5}

The fifth positive layer is the instrument and self-reference scope of Section \hyperref[sn:11]{11}. The paper records the instrument-indexed claim semantics of Subsection \hyperref[sn:11.1]{11.1}, the claim classifier priority of Subsection \hyperref[sn:11.2]{11.2}, the visibility and suppression discipline of Subsection \hyperref[sn:11.3]{11.3}, the no-overreading suppression theorem (Theorem \hyperref[cl:thm:23]{23}), the same-level self-audit failure theorem (Theorem \hyperref[cl:thm:24]{24}), the \(P_{6}\leftarrow P_{6}\) status separation of Subsection \hyperref[sn:11.6]{11.6}, the rotating-audit chain and bridge of Subsection \hyperref[sn:11.7]{11.7}, the finite rotating-audit theorem (Theorem \hyperref[cl:thm:25]{25}), and the instrument-transfer rule of Subsection \hyperref[sn:11.9]{11.9}.

The discipline records that no claim accepted under an instrument \(I\) uses suppressed content that no admissible bridge connects to visible content, that the claim classifier accepts no instrument's claim of its own complete self-soundness without a level-shifted audit, that finite rotating-audit chains are admissible but non-final, and that instrument transfer is bridge-mediated: acceptance under one instrument does not imply acceptance under another.

\subsubsection{Model-Realization Scope}\label{sn:15.1.6}

The sixth positive layer is the model-realization scope of Section \hyperref[sn:13]{13} and Section \hyperref[sn:12]{12}. The paper records four scoped model realizations: the PICA realization \(\mathsf{PICAReal}_{\mathrm{fin}}\) of Subsection \hyperref[sn:13.2]{13.2} with Model Theorem \hyperref[cl:thm:33]{33}, the Cantor strict-bridge realization \(\mathsf{CantorShell}_{\mathrm{aud}}\) of Subsection \hyperref[sn:13.3]{13.3} with Model Theorem \hyperref[cl:thm:34]{34}, the abstract-interpretation family \(\mathsf{AIFam}_{\mathrm{fin}}\) of Section \hyperref[sn:12]{12} with Model Theorem \hyperref[cl:thm:29]{29}, and the graph/cohomology realization \(\mathsf{CohGraph}_{\mathrm{fin}}\) of Subsection \hyperref[sn:13.4]{13.4}. Each realization is host-bound, fragment-scoped, and model-grade, with verdict \(\texttt{model\_realization}\), in keeping with the model-realization convention of Subsection \hyperref[sn:13.1]{13.1}.

\subsection{Required Nonclaims}\label{sn:15.2}

This subsection records the nonclaims of the paper, organized into six topical groupings: algebra and universality, status, package and strictness, route/refinement/locality, audit, and model. They are the thirty-six numbered nonclaims \(\mathrm{NC}\text{-}1\) through \(\mathrm{NC}\text{-}36\), the top-down nonclaim \(\mathrm{NC}\text{-}\mathrm{TD}\), and two unnumbered nonclaims on exact-six representation and unique decomposition. Each nonclaim is a scoped negative statement; together they fix the boundary of \(\mathbf{BirdInt}^{\mathrm{aud}}_{\mathrm{fin}}\) as a finite typed audited interaction calculus and prevent the positive results of Subsection \hyperref[sn:15.1]{15.1} from being read as universal foundations. Every nonclaim recorded in the body of the paper (per-section nonclaim subsections, per-realization nonclaim registers, per-bridge nonclaim records) is covered by, or refines, one of these entries, and each entry is carried by the finite nonclaim records \(\mathcal N\) of instruments, judgments, bridges, and model-realization sheets. Where a theorem, a countermodel, or a nonclaim register supports an entry, the entry names it.

\subsubsection{Algebra and Universality Nonclaims}\label{sn:15.2.1}

\begin{itemize}
\tightlist
\item
  \(\mathrm{NC}\text{-}1\). No universal six-symbol algebra: there is no faithful total operation \(\ast:\mathbb P^{2}\to\mathbb P\) whose composition with a status-decoding map recovers typed, statused, profile-relative directed-cell status across all profiles. The non-factorization of cell status through the primitive pair (Theorem \hyperref[cl:thm:1]{1}, Section \hyperref[sn:6]{6}) blocks any such operation as a status-decoding map.
\item
  \(\mathrm{NC}\text{-}2\). No ``all mathematics is six birds'' claim: \(\mathbf{BirdInt}^{\mathrm{aud}}_{\mathrm{fin}}\) is a finite typed audited interaction calculus, not a foundation for all mathematics, and no mathematical object is thereby a \(\mathbf{BirdInt}^{\mathrm{aud}}_{\mathrm{fin}}\) object under any reading.
\item
  \emph{No unrestricted exact-six theorem.} The paper does not assert that every description, theory or phenomenon decomposes into exactly the six roles (the Foundations II result is scoped to \(\mathsf{FATCD}\)), nor that every typed interaction is exactly representable as one of the thirty-six directed cells \(P_{i}\leftarrow P_{j}\). The exact-six structure is recorded as a finite typed schema with admissibility conditions, not as an unrestricted decomposition theorem.
\item
  \emph{No unique decomposition.} The paper does not assert that a typed interaction has a unique decomposition into primitive directed cells. Decomposition records are admissible only under a declared profile and audit, and different profiles can yield different admissible decompositions.
\end{itemize}

\subsubsection{Status Nonclaims}\label{sn:15.2.2}

\begin{itemize}
\tightlist
\item
  \(\mathrm{NC}\text{-}3\). No primitive-pair status determination: the directed-cell status does not factor through the primitive pair \((P_{i},P_{j})\) alone. The status non-factorization theorem (Theorem \hyperref[cl:thm:1]{1}, Section \hyperref[sn:6]{6}) records the formal counterexample.
\item
  \(\mathrm{NC}\text{-}4\). No status-family collapse: \(\mathrm{RoleStatus}\), \(\mathrm{CellStatus}\), \(\mathrm{PairStatus}\), \(\mathrm{PromotionStatus}\), \(\mathrm{ClaimStatus}\), \(\mathrm{GateStatus}\), and \(\mathrm{SqStatus}\) remain separately typed, judgment-indexed families; Theorem \hyperref[cl:thm:5]{5} of Section \hyperref[sn:6]{6} records representative separations.
\item
  \(\mathrm{NC}\text{-}5\). No directed-cell-to-pair collapse: a directed-cell status does not determine the pair-observable status. \(\mathsf{CM}_{4}\) of Section \hyperref[sn:9]{9} records a directed cell with status \(\texttt{action}\) whose corresponding pair observable is \(\texttt{blocked}\).
\item
  \(\mathrm{NC}\text{-}6\). No role-channel-to-cell collapse: the role status of a primitive does not determine the directed-cell status of any cell carrying that primitive. Role status is a separate judgment family from cell status; Theorem \hyperref[cl:thm:5]{5} of Section \hyperref[sn:6]{6} includes the channel-not-cell separation.
\end{itemize}

\subsubsection{Package and Strictness Nonclaims}\label{sn:15.2.3}

\begin{itemize}
\tightlist
\item
  \(\mathrm{NC}\text{-}7\). No package-implies-closure: a \(P_{5}\)-packaging operator does not imply macro-dynamics closure of an associated transformer. \(\mathsf{CM}_{1}\) of Section \hyperref[sn:9]{9} records an active package whose macro dynamics are not closed under \(I\).
\item
  \(\mathrm{NC}\text{-}8\). No idempotence-implies-commutation: two idempotent finite completions \(E_{1},E_{2}:C\to C\) may fail to commute. \(\mathsf{CM}_{8}\) of Appendix \hyperref[sn:C]{C}, named in Subsection \hyperref[sn:9.8]{9.8}, records two idempotent completions with \(E_{1}E_{2}\neq E_{2}E_{1}\).
\item
  \(\mathrm{NC}\text{-}9\). No repeated fixed completion as theory growth: iterating a fixed idempotent completion yields no new objects (\(E^n=E\), Theorem \hyperref[cl:thm:11]{11}); strict growth needs a new object map, bridge, refinement or nonfactorization witness. The Cantor target has a new packaged-object map \(\pi_1\); \(E_{\tau,\ell}=U_\ell\circ Q_\ell\circ K^\tau\) is attached completion data acting on distributions, rather than repetition of one fixed completion.
\item
  \(\mathrm{NC}\text{-}10\). No strictness-implies-closure: a strict-bridge passage \(\mathrm{Promote}(B):\texttt{strict}\) does not imply macro-dynamics closure on the target side of the bridge. \(\mathsf{CM}_{9}\) of Appendix \hyperref[sn:C]{C}, named in Subsection \hyperref[sn:9.8]{9.8}, records a strict bridge without closure.
\item
  \(\mathrm{NC}\text{-}11\). No strictness-implies-drive: a strict-bridge passage does not imply a \(\mathrm{P6}_{\mathrm{drive}}\) certificate on the support graph \(G_B\) recorded in the bridge's level/lens/interface record. \(\mathsf{CM}_{10}\) of Appendix \hyperref[sn:C]{C}, named in Subsection \hyperref[sn:9.8]{9.8}, records a strict bridge without drive.
\end{itemize}

\subsubsection{Route, Refinement, and Locality Nonclaims}\label{sn:15.2.4}

\begin{itemize}
\tightlist
\item
  \(\mathrm{NC}\text{-}12\). No route-mismatch-implies-drive: a \(P_{3}\) route-mismatch witness does not imply a \(\mathrm{P6}_{\mathrm{drive}}\) certificate. \(\mathsf{CM}_{2}\) of Section \hyperref[sn:9]{9} records a route mismatch without drive, and the cohomology nonclaims of Subsubsection \hyperref[sn:13.4.4]{13.4.4} record the corresponding model-level form.
\item
  \(\mathrm{NC}\text{-}13\). No refinement-implies-improvement: a lens refinement \(q_0\rightsquigarrow q_1\) does not improve every defect record. \(\mathsf{CM}_{5}\) of Section \hyperref[sn:9]{9} records a refinement that worsens closure under a chosen defect measure.
\item
  \(\mathrm{NC}\text{-}14\). No local-implies-global: local admissibility of a finite family of packages does not imply global admissibility. \(\mathsf{CM}_{3}\) of Section \hyperref[sn:9]{9} records local packages that fail globally.
\item
  \(\mathrm{NC}\text{-}\mathrm{TD}\). No structural-downward-implies-channel claim: a completion map, feasibility gate, macro profile, phase record, promoted object, or context-indexed update supplies at most a structural downward path or constraint. It certifies top-down causal difference-making only when a finite \(\mathsf{TopDownChannelRecord}\) supplies admissible macro interventions, matched controls, source-of-truth response records, visibility, audit, no-smuggling, and a passed host-specific effect threshold. The third proposition of Subsection \hyperref[sn:4.7]{4.7} proves it.
\end{itemize}

\subsubsection{Audit Nonclaims}\label{sn:15.2.5}

\begin{itemize}
\tightlist
\item
  \(\mathrm{NC}\text{-}15\). No audit-is-truth: a passing audit record or claim status \(\texttt{accepted}\) under an instrument is not instrument-free truth. The instrument-indexed claim form of Subsubsection \hyperref[sn:11.1.3]{11.1.3} fixes acceptance as relative to \(I\) and the surrounding context \(\Gamma;\mathcal T\).
\item
  \(\mathrm{NC}\text{-}16\). No same-level complete self-certification: the claim classifier does not accept, under an instrument \(I_0\), the claim \(\mathrm{Sound}(I_0)\) that \(I_0\) is completely sound. Theorem \hyperref[cl:thm:24]{24} of Section \hyperref[sn:11]{11} records the classifier routing: the claim classifier assigns same-level complete self-soundness under \(I_{0}\) the status \(\texttt{outside\_scope}\), \(\texttt{undefined\_circular}\), or a status of higher priority, and never \(\texttt{accepted}\).
\item
  \(\mathrm{NC}\text{-}17\). No rotating-audit finality: a finite rotating-audit chain is non-final. Theorem \hyperref[cl:thm:25]{25} of Section \hyperref[sn:11]{11} and the non-finality clause of Subsubsection \hyperref[sn:11.8.1]{11.8.1} record that a finite rotating-audit chain audits the lower stack but does not certify its top.
\item
  \(\mathrm{NC}\text{-}18\). No overreading: an \(I\)-accepted claim has no \(I\)-suppressed use-set element without an admissible bridge (Theorem \hyperref[cl:thm:23]{23} of Section \hyperref[sn:11]{11}); acceptance under \(I\) asserts nothing about suppressed content that no admissible bridge connects to visible content.
\item
  \(\mathrm{NC}\text{-}19\). No model-realization-as-theorem overclaim: a \(\texttt{model\_realization}\) verdict does not promote to theorem-grade \(\texttt{accepted}\) without an admissible bridge, and it does not prove the declared structure universally. The fragment discipline of Subsubsection \hyperref[sn:13.1.2]{13.1.2} records the corresponding rule.
\end{itemize}

\subsubsection{Model Nonclaims}\label{sn:15.2.6}

\begin{itemize}
\tightlist
\item
  \(\mathrm{NC}\text{-}20\). No PICA universality: \(\mathsf{PICAReal}_{\mathrm{fin}}\) realizes a finite stochastic fragment of the calculus, not its entirety. The PICA realization is host-bound and fragment-scoped to the finite stochastic cell/pair/provenance fragment.
\item
  \(\mathrm{NC}\text{-}21\). No PICA \(6\times 6\) universal algebra: the PICA \(6\times 6\) multiplicity count, recovered under the record conditions of Subsubsection \hyperref[sn:13.2.5]{13.2.5}, is a profile fact, not a universal law. It is a PICA-profile model fact under the cell profile \(\lambda_{\mathrm{PICA}}\).
\item
  \(\mathrm{NC}\text{-}22\). No raw implementation status as formal status: raw PICA status is not formal directed-cell status without recovery and classifier check.
\item
  \(\mathrm{NC}\text{-}23\). No fallback-as-real-source: a fallback record supports at most \(\texttt{real\_provisional}\) unless an admissible upgrade bridge is recorded. The fallback-default cap of Subsubsection \hyperref[sn:13.2.8]{13.2.8} records that fallback supports at most \(\texttt{real\_provisional}\) unless an explicit upgrade bridge is recorded, and \(\mathsf{CM}_{11}\) of Appendix \hyperref[sn:C]{C} gives the corresponding finite source-policy countermodel.
\item
  \(\mathrm{NC}\text{-}24\). No dirty-state realness: a dirty record in a pair's source slot, or among the \(\mathsf{SourceRefs}\) of its audit or branch-compatibility audit, cannot support a real pair observable: the first \(\texttt{blocked}\) row applies.
\item
  \(\mathrm{NC}\text{-}25\). No Cantor generalization beyond audited shell: the Cantor strict-bridge realization is shell-bound. Subsubsection \hyperref[sn:13.3.9]{13.3.9} records the corresponding nonclaim register.
\item
  \(\mathrm{NC}\text{-}26\). No Cantor pressure gap as KL closure deficit: the pressure-gap consequence defect is not a Kullback--Leibler closure-deficit without a separately admissible bridge. Cantor nonclaim 6 of Subsubsection \hyperref[sn:13.3.9]{13.3.9}.
\item
  \(\mathrm{NC}\text{-}27\). No Cantor pressure gap as strictness witness: the pressure-gap defect is a separately recorded consequence record, not the strictness witness. Cantor nonclaim 7 of Subsubsection \hyperref[sn:13.3.9]{13.3.9}.
\item
  \(\mathrm{NC}\text{-}28\). No direct autonomous stratum theorem: the persistent strata family \(\mathcal F\) does not support a direct stratumwise root-separation theorem in the present paper. Cantor nonclaim 8 of Subsubsection \hyperref[sn:13.3.9]{13.3.9}.
\item
  \(\mathrm{NC}\text{-}29\). No AI totality: a nonempty admissible AI host family sharing one concrete domain \(C\) and one concrete transformer \(F\) realizes \(P_{1}/P_{2}/P_{5}/P_{6}\), and \(P_{1}/P_{2}/P_{4}/P_{5}/P_{6}\) when its refinement records contain a strict refinement (Model Theorem \hyperref[cl:thm:29]{29}); the realization map of Model Theorem \hyperref[cl:thm:29]{29} assigns role fragments only, adds \(P_{3}\) when two closures fail to commute, and assigns no cell, pair or promotion component, so it does not realize the whole calculus.
\item
  \(\mathrm{NC}\text{-}30\). No AI soundness-as-completeness: sound abstraction does not imply complete representation. Subsubsection \hyperref[sn:12.6.2]{12.6.2}.
\item
  \(\mathrm{NC}\text{-}31\). No AI closure-implies-soundness: the closure operator \(\rho=\gamma\alpha\) does not imply transformer soundness. Subsubsection \hyperref[sn:12.6.3]{12.6.3}.
\item
  \(\mathrm{NC}\text{-}32\). No graph cohomology as all drive: the cohomological drive certificate \([a]\neq 0\) is one drive certificate family, not all drive. The cohomology nonclaims of Subsubsection \hyperref[sn:13.4.4]{13.4.4} record \(\mathrm{P6}_{\mathrm{drive}}^{H^1}\) as one finite certificate family, not all drive.
\item
  \(\mathrm{NC}\text{-}33\). No gating-preserves-drive: \(P_{2}\)-gating may suppress cycle-supported drive on \(G\). Theorem \hyperref[cl:thm:18]{18} and \(\mathsf{CM}_{12}\) of Appendix \hyperref[sn:C]{C} record it.
\item
  \(\mathrm{NC}\text{-}34\). No double-category completeness: the finite decorated-square fragment of Section \hyperref[sn:14]{14} does not replace the core calculus. Subsubsection \hyperref[sn:14.5.1]{14.5.1} records the corresponding nonclaim.
\item
  \(\mathrm{NC}\text{-}35\). No high-structure status erasure: the decorated-square fragment does not replace statuses, defects, audits, or instruments, and does not collapse the seven status families into a single square-status family. Subsubsection \hyperref[sn:14.5.2]{14.5.2}.
\item
  \(\mathrm{NC}\text{-}36\). No mechanization beyond manifest scope: the proof-assistant artifact validates the finite audited declarations listed in Appendix \hyperref[sn:F]{F}, and does not certify unlisted statements, empirical implementation correctness, a universal exact-six theorem, or a universal six-symbol algebra. The audit-grade verdict of Proposition \hyperref[cl:prop:35]{35} records the completed validation artifact for the finite manifest scope, not a certificate for unlisted statements, empirical implementations, or unrestricted host extensions.
\end{itemize}

\subsection{Future Work}\label{sn:15.3}

This subsection records the six classes of future work that the present paper does not undertake. Each class is a scoped pointer to admissible extension targets for \(\mathbf{BirdInt}^{\mathrm{aud}}_{\mathrm{fin}}\); none of the entries below is a claim about results in this paper.

\subsubsection{Classifier and Dependency Completeness}\label{sn:15.3.1}

Future work on the claim classifier of Subsection \hyperref[sn:11.2]{11.2} includes:

\begin{itemize}
\tightlist
\item
  a complete status-classifier theorem, characterizing the priority order over \(\mathrm{ClaimStatus}\) under stated soundness and admissibility conditions for the claim classifier of Subsubsection \hyperref[sn:5.6.5]{5.6.5};
\item
  a defect-to-status soundness theorem, characterizing the routing from the defect families of Subsections \hyperref[sn:5.2]{5.2}--\hyperref[sn:5.5]{5.5} through the classifier rules of Subsection \hyperref[sn:5.6]{5.6};
\item
  a dependency-lattice theorem, organizing the per-component admissibility predicates of Section \hyperref[sn:4]{4} and Section \hyperref[sn:5]{5} into a lattice of dependent admissibility checks with a finite reduction theorem.
\end{itemize}

The present paper records the priority order of Subsubsection \hyperref[sn:11.2.1]{11.2.1} and the per-status routing rules of Subsection \hyperref[sn:5.6]{5.6}, but does not establish completeness or soundness in these stronger forms.

\subsubsection{Promotion and Stacking Theory}\label{sn:15.3.2}

Future work on promotion and stacking includes:

\begin{itemize}
\tightlist
\item
  soundness and completeness conditions for the promotion classifier of Subsubsection \hyperref[sn:5.6.4]{5.6.4} over \(\mathrm{PromotionStatus}\), relative to an external account of correct promotion;
\item
  a stacking composition theorem for composed bridges whose middle interface refines the source, extending the object-map result of Subsubsection \hyperref[sn:10.7.1]{10.7.1} to the joint gate and nonclaim semantics of the composed bridge;
\item
  a local-to-global promotion theorem, identifying the conditions under which a finite family of locally admissible promotion bridges supports a global promotion verdict.
\end{itemize}

The present paper records Theorems \hyperref[cl:thm:19]{19}--\hyperref[cl:thm:22]{22} of Section \hyperref[sn:10]{10}, the eight-gate machinery, and the per-bridge classifier; the three extension targets above are recorded as future work.

\subsubsection{Verified Model Realizations}\label{sn:15.3.3}

Future work on the model-realization sheets includes:

\begin{itemize}
\tightlist
\item
  a deeper PICA realization audit, extending Model Theorem \hyperref[cl:thm:33]{33} of Subsection \hyperref[sn:13.2]{13.2} toward implementation-correctness records for the actor-update map, informant-witness map, status recovery, and pair-observable realness predicate;
\item
  a deeper Cantor strict-bridge realization audit, extending Model Theorem \hyperref[cl:thm:34]{34} of Subsection \hyperref[sn:13.3]{13.3} toward implementation-correctness records for the shell-bound strictness witness and the gate-result table;
\item
  a deeper abstract-interpretation realization audit, extending Model Theorem \hyperref[cl:thm:29]{29} of Section \hyperref[sn:12]{12} beyond the manifest-scoped Lean support toward richer concrete domains and transformers;
\item
  an implementation-correctness layer for any concrete PICA implementation, audited under a separate proof-assistant-bound instrument.
\end{itemize}

The present paper records the four realization sheets at model-grade verdict \(\texttt{model\_realization}\) and records manifest-scoped proof-assistant declarations for the required model claims. Implementation correctness and richer external model audits remain future work.

\subsubsection{Cantor Shell Expansion}\label{sn:15.3.4}

Future work on the Cantor strict-bridge realization includes:

\begin{itemize}
\tightlist
\item
  a full Cantor strictness proof, extending the audited-shell strict-bridge realization beyond \(\mathcal S_{\mathrm{aud}}\) to a wider scope under an admissible expansion bridge;
\item
  a pressure-gap-to-KL bridge, recording an admissible bridge from the pressure-gap consequence defect of Subsubsection \hyperref[sn:13.3.6]{13.3.6} to the Kullback--Leibler closure-deficit framework;
\item
  a stratumwise root-separation theorem, characterizing the persistent strata family \(\mathcal F\) via a separation property that supports stronger consequences than the conditional pressure disintegration of Subsubsection \hyperref[sn:13.3.6]{13.3.6}.
\end{itemize}

The present paper records the audited-shell realization with verdict \(\texttt{model\_realization}\) on \(\mathbf{StrictBridge}_{0\to 1}^{\mathrm{audited\ shell}}\); the three extensions above are recorded as future work.

\subsubsection{Drive and High-Structure Semantics}\label{sn:15.3.5}

Future work on the drive and high-structure layers includes:

\begin{itemize}
\tightlist
\item
  a general drive classification, extending the cohomological drive certificate \([a]\neq 0\) of Subsubsection \hyperref[sn:13.4.3]{13.4.3} to a finite family of drive certificates over multiple drive notions;
\item
  a holonomy-to-drive bridge \(B_{3\to 6}^{\mathrm{drive}}\), recording an admissible bridge from a \(P_{3}\) route/holonomy witness to a \(\mathrm{P6}_{\mathrm{drive}}\) certificate, in the sense of Subsubsection \hyperref[sn:13.4.4]{13.4.4} nonclaim 2;
\item
  gating beyond forests, characterizing \(P_{2}\)-gating effects on drive on graphs with non-trivial cycle structure under additional gating policies;
\item
  weak double-category coherence, extending the strict pasting law of Subsubsection \hyperref[sn:14.1.3]{14.1.3} to weak coherence relations on finite decorated squares;
\item
  a general obstruction propagation theorem, characterizing the propagation of square defects under finite strict pasting in the sense of Subsubsection \hyperref[sn:14.5.3]{14.5.3}.
\end{itemize}

The present paper records the cohomological theorems of Section \hyperref[sn:8]{8} and the high-structure recovery theorems of Section \hyperref[sn:14]{14}; the extensions above are recorded as future work.

\subsubsection{Self-Reference and Instrument Transfer}\label{sn:15.3.6}

Future work on self-reference and instrument transfer includes:

\begin{itemize}
\tightlist
\item
  a full instrument-transfer classifier, fixing the per-claim-type mapping from source-side acceptance and bridge defect to target-side \(\chi\in\mathrm{ClaimStatus}\) under \(J\), in the sense of Subsubsection \hyperref[sn:11.9.2]{11.9.2};
\item
  an infinite rotating-audit semantics, extending the finite rotating-audit chain of Subsection \hyperref[sn:11.7]{11.7} to an inverse limit or coalgebraic semantics admissible under a stronger meta-instrument;
\item
  a complete self-reference classification, extending Theorem \hyperref[cl:thm:24]{24} of Subsection \hyperref[sn:11.5]{11.5} to a comprehensive classification of same-level, level-shifted, and rotating self-reference patterns.
\end{itemize}

The present paper records the instrument-transfer rule of Subsubsection \hyperref[sn:11.9.1]{11.9.1}, the rotating-audit chain of Subsection \hyperref[sn:11.7]{11.7}, and Theorems \hyperref[cl:thm:24]{24} and \hyperref[cl:thm:25]{25}; the extensions above are recorded as future work.
 \section{Conclusion}\label{sn:16}

This section closes the paper by restating the contribution: the constructed object \(\mathbf{BirdInt}^{\mathrm{aud}}_{\mathrm{fin}}\), the layered theorem stack it admits, the scope discipline it imposes, and the final boundary it does not cross.

\subsection{The Constructed Object}\label{sn:16.1}

The paper constructs \(\mathbf{BirdInt}^{\mathrm{aud}}_{\mathrm{fin}}\), a finite audited typed interaction calculus for six-bird theory. The core calculus uses finite carriers. The Cantor realization uses finite records over a shell that is not asserted finite (Subsubsection \hyperref[sn:2.1.2]{2.1.2}). Every store is finite, every admissibility predicate is decidable on finite records, under the numerical decidability and computability hypotheses of Theorem \hyperref[cl:thm:3]{3}, together with the conditions the instance declares (Subsubsection \hyperref[sn:5.6.6]{5.6.6}), and every claim is recorded under an explicit host instrument. Its components are six primitive labels \(\mathbb P=\{P_{1},\ldots,P_{6}\}\), a finite typed schema of directed-cell, pair-observable, promotion-bridge, claim, instrument, source-of-truth, audit, and visibility records, seven finite status families with their classifier or per-value rules where applicable, a finite typed family of defect records that route judgments and claims to their status consequences, and a finite nonclaim register that records the negative boundary of every admissible claim.

The calculus is not a universal algebra over six symbols, an all-mathematics decomposition, or a foundation for empirical phenomena. It is a finite typed audited interaction calculus, defined by Sections \hyperref[sn:3]{3}, \hyperref[sn:4]{4}, and \hyperref[sn:5]{5}, and admitting the theorem-grade, model-grade, and audit-grade results recorded in Sections \hyperref[sn:6]{6} through \hyperref[sn:14]{14} and Appendix \hyperref[sn:F]{F}.

\subsection{The Theorem Stack}\label{sn:16.2}

The paper records its content as a layered theorem stack:

\begin{itemize}
\tightlist
\item
  \emph{Finite typed calculus and core theorems.} Sections \hyperref[sn:6]{6} and \hyperref[sn:7]{7} record the no-total-algebra theorem (Theorem \hyperref[cl:thm:1]{1}), the four core finite-interaction theorems (Theorems \hyperref[cl:thm:2]{2}--\hyperref[cl:thm:5]{5}), and the nine descent/packaging/completion/strictness theorems (Theorems \hyperref[cl:thm:6]{6}--\hyperref[cl:thm:14]{14}), all theorem-grade with verdict \(\texttt{accepted}\).
\item
  \emph{No-go and cohomological theorems.} Section \hyperref[sn:8]{8} records the four cohomological no-drive theorems (Theorems \hyperref[cl:thm:15]{15}--\hyperref[cl:thm:18]{18}) on graph-supported drive.
\item
  \emph{Countermodel/no-go atlas.} Section \hyperref[sn:9]{9} and Appendix \hyperref[sn:C]{C} record twelve countermodels (\(\mathsf{CM}_{1}\)--\(\mathsf{CM}_{12}\)), each blocking a specific overclaim about the calculus under a finite typed structure.
\item
  \emph{Promotion and instrument-indexed semantics.} Section \hyperref[sn:10]{10} records the promotion gate and strictness theorems (Theorems \hyperref[cl:thm:19]{19}--\hyperref[cl:thm:21]{21}) together with a two-point strict-bridge example (Theorem \hyperref[cl:thm:22]{22}), and Section \hyperref[sn:11]{11} records the no-overreading suppression theorem (Theorem \hyperref[cl:thm:23]{23}), the same-level self-audit failure theorem (Theorem \hyperref[cl:thm:24]{24}), and the finite rotating-audit theorem (Theorem \hyperref[cl:thm:25]{25}); Appendix \hyperref[sn:I]{I} records the instrument-indexed form of each.
\item
  \emph{Scoped model realizations.} Section \hyperref[sn:12]{12} records theorem-grade abstract-interpretation host results (Theorems \hyperref[cl:thm:26]{26}--\hyperref[cl:thm:28]{28}) and the model-grade abstract-interpretation family result (Model Theorem \hyperref[cl:thm:29]{29}). Section \hyperref[sn:13]{13} records the PICA realization (Model Theorem \hyperref[cl:thm:33]{33}), the Cantor strict-bridge realization (Model Theorem \hyperref[cl:thm:34]{34}), and the graph/cohomology realization, all model-grade with verdict \(\texttt{model\_realization}\) on declared fragments.
\item
  \emph{High-structure semantics as secondary organization.} Section \hyperref[sn:14]{14} records three decorated-square recovery theorems (Theorems \hyperref[cl:thm:30]{30}--\hyperref[cl:thm:32]{32}) over a finite double-category fragment, presented as secondary organization rather than as a replacement for the core calculus.
\item
  \emph{Mechanization audit.} Appendix \hyperref[sn:F]{F} records the manifest-scoped proof-assistant validation artifact and Proposition \hyperref[cl:prop:35]{35}, with the audit-grade verdict \(\texttt{audit\_report}\).
\end{itemize}

Together these record the paper's content. The theorem and model-realization claims carry the verdicts of their claim strengths (Subsubsection \hyperref[sn:2.3.4]{2.3.4}), while the countermodels and Lean validation proposition are recorded with their corresponding finite dependency or audit records.

\subsection{The Scope Discipline}\label{sn:16.3}

The paper replaces a naive total algebra over six primitive symbols with a typed, finite, statused, audited interaction calculus. The replacement is structural: cell status does not factor through the primitive pair (Theorem \hyperref[cl:thm:1]{1}); five non-implications separate the judgment-attached status families, and the gate and square families remain schema-separate (Theorem \hyperref[cl:thm:5]{5} and Section \hyperref[sn:5]{5}); descent is well-definedness on a finite quotient (Theorem \hyperref[cl:thm:6]{6}), strict extension is fiber mismatch on a finite carrier (Theorem \hyperref[cl:thm:12]{12}), accepted claims do not overread suppressed content (Theorem \hyperref[cl:thm:23]{23}), the claim classifier does not accept an instrument's claim of its own complete soundness at the same level (Theorem \hyperref[cl:thm:24]{24}), and finite rotating audits are admissible, given a higher level, but non-final (Theorem \hyperref[cl:thm:25]{25}). Each replacement is recorded as a finite typed result under explicit host, instrument, admissibility, and nonclaim assumptions, with defect schemas explicit where they drive status.

The scope discipline is uniform: every theorem-grade claim, model-realization claim, promotion judgment, and audit-report statement is accompanied by an instrument-indexed or host-indexed form, a finite admissibility predicate, and a nonclaim record that fixes the boundary of the claim. The paper records exactly what it asserts and exactly what it does not assert; the negative scope, recorded in Section \hyperref[sn:15]{15} and in the nonclaim subsections of each section, is part of the calculus itself.

\subsection{Final Boundary}\label{sn:16.4}

The final boundary of the paper is the finite audited scope. The calculus \(\mathbf{BirdInt}^{\mathrm{aud}}_{\mathrm{fin}}\) is a finite typed audited interaction calculus, admitting the theorem-grade, model-grade, and audit-grade results recorded in this paper, and not admitting an unrestricted exact-six theorem, a universal six-symbol algebra, a unique decomposition theorem, complete model-coverage, complete instrument-family completeness, or final self-certification.

The paper closes with this boundary in hand: the constructed object, the layered theorem stack, the scope discipline, and the explicit non-extension of any of the above into a universal foundation claim. The contribution is the calculus and what it admits --- finite, typed, statused, audited, and bounded by an explicit nonclaim register.
 
\appendix
\section{Full Definitions}\label{sn:A}

This appendix records the long-form definitions and tables of \(\mathbf{BirdInt}^{\mathrm{aud}}_{\mathrm{fin}}\), supplementing the body of the paper. Subsection \hyperref[sn:A.1]{A.1} records the full domain tuple. Subsection \hyperref[sn:A.2]{A.2} records the full primitive witness table. Subsection \hyperref[sn:A.3]{A.3} records the full primitive update table. Subsection \hyperref[sn:A.4]{A.4} records the full profile grammar and admissibility table. Subsection \hyperref[sn:A.5]{A.5} records the full status-family tables. Subsection \hyperref[sn:A.6]{A.6} records the full primitive and promotion defect tables. Each subsection collects, in one place, definitions made in the body: A.1 the tuple of Subsubsection \hyperref[sn:3.1.1]{3.1.1}, \hyperref[sn:A.2]{A.2} and \hyperref[sn:A.3]{A.3} the tables of Subsubsections \hyperref[sn:3.7.4]{3.7.4} and \hyperref[sn:3.7.5]{3.7.5}, \hyperref[sn:A.4]{A.4} the profiles of Subsection \hyperref[sn:3.6]{3.6}, \hyperref[sn:A.5]{A.5} the status sets of Sections \hyperref[sn:4]{4} and \hyperref[sn:5]{5}, and \hyperref[sn:A.6]{A.6} the defect schemas of Subsections \hyperref[sn:5.2]{5.2} through \hyperref[sn:5.4]{5.4}. The body's definition is the reference.

\subsection{Full Domain Tuple}\label{sn:A.1}

The full domain tuple of \(\mathbf{BirdInt}^{\mathrm{aud}}_{\mathrm{fin}}\) is the nineteen-component finite typed record displayed in Subsubsection \hyperref[sn:3.1.1]{3.1.1}, with components:

\begin{itemize}
\tightlist
\item
  \(\mathsf{Host}\), the finite and scoped host inventory, recording finite carriers, finite maps, finite records, finite relations, finite diagrams, finite graphs, and finite audit records (Subsection \hyperref[sn:2.1]{2.1});
\item
  \(\mathsf{Cont}\), the finite content set whose records may be visible, suppressed, outside-scope, or unknown under an instrument (Subsections \hyperref[sn:3.2]{3.2} and \hyperref[sn:3.5]{3.5});
\item
  \(\mathbb P=\{P_{1},P_{2},P_{3},P_{4},P_{5},P_{6}\}\), the finite primitive labels with stable role glosses (Subsection \hyperref[sn:2.2]{2.2});
\item
  \(\mathsf{Lev}\), the finite level set \(\{\mathsf{Beh},\mathsf{Ver},\mathsf{Str}\}\) used by the calculus (Subsubsection \hyperref[sn:2.2.3]{2.2.3});
\item
  \(\mathsf{Instr}\), the finite class of instrument records, with components introduced in Subsection \hyperref[sn:3.4]{3.4} and used in the claim semantics of Subsection \hyperref[sn:11.1]{11.1};
\item
  \(\mathsf{Pkg}\), the finite class of theory packages with components \((Z,f,\Sigma_{f},E,\mathcal A)\) (Subsection \hyperref[sn:3.3]{3.3});
\item
  \(\mathsf{Prof}\), the finite typed profile family \(\lambda\) with admissibility predicates (Subsection \hyperref[sn:3.6]{3.6});
\item
  \(\mathsf{Wit}\), the finite typed witness family for primitive informants (Subsection \hyperref[sn:3.7]{3.7} and Table A.2);
\item
  \(\mathsf{Upd}\), the finite typed update family for primitive actors (Subsection \hyperref[sn:3.7]{3.7} and Table A.3);
\item
  \(\mathsf{Def}\), the finite typed defect family with per-component schemas (Subsections \hyperref[sn:5.2]{5.2} through \hyperref[sn:5.5]{5.5} and Table A.6);
\item
  \(\mathsf{Judg}\), the finite typed judgment family --- role-channel, directed-cell, pair-observable, promotion, dependency, instrument-indexed claim and model-realization judgments; gate and square records are classified by \(\mathrm{GateStatus}\) and \(\mathrm{SqStatus}\) but are not judgment forms --- with per-judgment admissibility predicates (Subsections \hyperref[sn:4.2]{4.2} through \hyperref[sn:4.7]{4.7}, \hyperref[sn:10.2]{10.2}, \hyperref[sn:11.1]{11.1} and \hyperref[sn:14.1]{14.1}, and Subsubsection \hyperref[sn:13.1.1]{13.1.1});
\item
  \(\mathsf{Status}\), the seven finite status families, with their family-specific classifier and priority rules where defined (Subsection \hyperref[sn:5.1]{5.1} and Table A.5);
\item
  \(\mathsf{Gate}\), the finite typed gate family --- the eight promotion gates, including the strictness gate, and bridge-admissibility gates (Subsection \hyperref[sn:10.2]{10.2});
\item
  \(\mathsf{Dep}\), the finite typed dependency relation between content classes under an instrument (Subsection \hyperref[sn:4.6]{4.6});
\item
  \(\mathsf{Audit}\), the finite typed audit family --- host audit, claim audit, bridge audit, rotating-audit (Subsubsection \hyperref[sn:4.9.2]{4.9.2} and Section \hyperref[sn:11]{11});
\item
  \(\mathsf{Vis}\), the finite family of visibility maps and visibility bridges (Subsection \hyperref[sn:3.5]{3.5});
\item
  \(\mathsf{Source}\), the finite family of source records (Subsubsection \hyperref[sn:4.9.1]{4.9.1});
\item
  \(\mathsf{Nonclaim}\), the finite nonclaim register \(\mathcal N\) (Subsection \hyperref[sn:15.2]{15.2} and per-record nonclaim subfields throughout);
\item
  \(\mathsf{Real}\), the finite family of model-realization modules (Section \hyperref[sn:13]{13}).
\end{itemize}

Every strong claim under the calculus is instrument-indexed:

\[
\Gamma\,;\;\mathcal T\,;\;I\;\vdash\;\varphi\;:\;\chi,
\]

with \(\chi\in\mathrm{ClaimStatus}\) as recorded in Subsubsection \hyperref[sn:5.1.5]{5.1.5}. The shorthand \(I\vdash\varphi(\mathcal T)\) abbreviates \(\Gamma;\mathcal T;I\vdash\varphi:\texttt{accepted}\) under the acceptance rule of Subsubsection \hyperref[sn:11.1.3]{11.1.3} and never abbreviates instrument-free truth.

The domain is \emph{finite}: every component is finite, every admissibility predicate is a finite check on finite records, decidable when numerical fields have decidable equality and order and threshold data are computable (Theorem \hyperref[cl:thm:3]{3}), relative to the conditions the instance declares (Subsubsection \hyperref[sn:5.6.6]{5.6.6}), and every claim under the domain is recorded with a finite typed instrument-indexed claim form. The design principle of the domain is that interaction is typed, statused, profile-relative, instrument-indexed, and audited --- and that the primitive-pair collapse \(P_{i}\ast P_{j}=P_{k}\) is rejected as a universal representation of interaction in keeping with Theorem \hyperref[cl:thm:1]{1} of Section \hyperref[sn:6]{6}.

\subsection{Full Primitive Witness Table}\label{sn:A.2}

The witness families are typed by the informant primitive: \(\mathsf{Wit}:\mathbb P\to\mathrm{Type}\), with \(W_{j}\in\mathsf{Wit}(P_{j})\) for the informant \(P_{j}\) of any directed cell \(P_{i}\leftarrow P_{j}\). The full witness families per primitive are:

\begingroup
\setlength{\tabcolsep}{4pt}
\renewcommand{\arraystretch}{1.08}
\begin{center}
\begin{tabularx}{0.96\linewidth}{@{}>{\raggedright\arraybackslash}p{0.11\linewidth}>{\raggedright\arraybackslash}X@{}}
\toprule
primitive & witness family \\
\midrule
$P_{1}$ & descent success/failure, rewrite obstruction, closure deficit \\
$P_{2}$ & representability, nonrepresentability, gate, support, capacity \\
$P_{3}$ & route mismatch, holonomy, commutator, critical pair, pasting mismatch \\
$P_{4}$ & stage, refinement, filtration, cover, local/global obstruction \\
$P_{5}$ & package map, quotient, completion, idempotence, fixed point \\
$P_{6}$ & audit, provenance, source-of-truth, threshold check, nonclaim, drive certificate \\
\bottomrule
\end{tabularx}
\end{center}
\endgroup

The typed actor/informant rule of Subsection \hyperref[sn:4.3]{4.3} fixes the admissibility of any directed-cell judgment: a cell record \(P_{i}\leftarrow P_{j}\) is admissible only when its witness slot holds \(W_{j}\in\mathsf{Wit}(P_{j})\), of a type drawn from the row of \(P_{j}\) above, or a typed, audited absence marker (Subsubsection \hyperref[sn:3.7.3]{3.7.3}), and its update slot holds an actor update of a type drawn from the row of \(P_{i}\) or such a marker in Subsection \hyperref[sn:A.3]{A.3} below. The witness families are finite for each primitive, and admissibility of the typed witness assignment is a decidable check on the finite records of the cell.

\subsection{Full Primitive Update Table}\label{sn:A.3}

The update families are typed by the actor primitive: \(\mathsf{Upd}:\mathbb P\to\mathrm{Type}\), with \(U_{i}\in\mathsf{Upd}(P_{i})\) for the actor \(P_{i}\) of any directed cell \(P_{i}\leftarrow P_{j}\). The full update/effect families per primitive are:

\begingroup
\setlength{\tabcolsep}{4pt}
\renewcommand{\arraystretch}{1.08}
\begin{center}
\begin{tabularx}{0.96\linewidth}{@{}>{\raggedright\arraybackslash}p{0.11\linewidth}>{\raggedright\arraybackslash}X@{}}
\toprule
primitive & update/effect family \\
\midrule
$P_{1}$ & mark descent failure, install descended update, rewrite, repair \\
$P_{2}$ & mark representable/nonrepresentable, gate support, update admissibility \\
$P_{3}$ & record route mismatch, holonomy, commutator, critical pair \\
$P_{4}$ & refine lens, advance stage, record local/global obstruction \\
$P_{5}$ & install package, apply completion, record fixed points, saturate \\
$P_{6}$ & add audit, check provenance, commit source, record nonclaim, certify drive/no-drive \\
\bottomrule
\end{tabularx}
\end{center}
\endgroup

The actor update is the second half of the typed actor/informant rule of Subsection \hyperref[sn:4.3]{4.3}: a cell record \(P_{i}\leftarrow P_{j}\) is admissible only when both \(W_{j}\in\mathsf{Wit}(P_{j})\) and \(U_{i}\in\mathsf{Upd}(P_{i})\) hold for the witness and update assigned to the cell. The update families are finite for each primitive, and admissibility of the typed update assignment is a decidable check on the finite records of the cell.

\subsection{Full Profile Grammar}\label{sn:A.4}

The profile grammar of \(\mathbf{BirdInt}^{\mathrm{aud}}_{\mathrm{fin}}\) is the finite typed record schema

\[
\lambda\;=\;(\,\ell_{i},\;\ell_{j},\;\mathsf{Scale},\;\mathsf{Regime},\;\mathsf{BridgeType},\;\mathsf{InstrumentMode},\;\mathsf{Locality}\,),
\]

with components:

\begin{longtable}[]{@{}ll@{}}
\toprule
\begin{minipage}[b]{0.47\columnwidth}\raggedright
component\strut
\end{minipage} & \begin{minipage}[b]{0.47\columnwidth}\raggedright
admissible values / role\strut
\end{minipage}\tabularnewline
\midrule
\endhead
\begin{minipage}[t]{0.47\columnwidth}\raggedright
\(\ell_i\)\strut
\end{minipage} & \begin{minipage}[t]{0.47\columnwidth}\raggedright
actor-side level tag in \(\mathsf{Lev}=\{\mathsf{Beh},\mathsf{Ver},\mathsf{Str}\}\)\strut
\end{minipage}\tabularnewline
\begin{minipage}[t]{0.47\columnwidth}\raggedright
\(\ell_j\)\strut
\end{minipage} & \begin{minipage}[t]{0.47\columnwidth}\raggedright
informant-side level tag in \(\mathsf{Lev}=\{\mathsf{Beh},\mathsf{Ver},\mathsf{Str}\}\)\strut
\end{minipage}\tabularnewline
\begin{minipage}[t]{0.47\columnwidth}\raggedright
\(\mathsf{Scale}\)\strut
\end{minipage} & \begin{minipage}[t]{0.47\columnwidth}\raggedright
finite scale label for the carrier, partition, kernel, graph, or other hosted data used by the judgment\strut
\end{minipage}\tabularnewline
\begin{minipage}[t]{0.47\columnwidth}\raggedright
\(\mathsf{Regime}\)\strut
\end{minipage} & \begin{minipage}[t]{0.47\columnwidth}\raggedright
finite regime label, such as dynamical, structural, or audit, as fixed by the surrounding judgment family\strut
\end{minipage}\tabularnewline
\begin{minipage}[t]{0.47\columnwidth}\raggedright
\(\mathsf{BridgeType}\)\strut
\end{minipage} & \begin{minipage}[t]{0.47\columnwidth}\raggedright
bridge type admitted by \(\mathsf{BridgePolicy}_I\), such as visibility, level, instrument-transfer, drive, or model-realization\strut
\end{minipage}\tabularnewline
\begin{minipage}[t]{0.47\columnwidth}\raggedright
\(\mathsf{InstrumentMode}\)\strut
\end{minipage} & \begin{minipage}[t]{0.47\columnwidth}\raggedright
finite instrument-mode label, such as default, audit-only, transfer, or rotating-audit\strut
\end{minipage}\tabularnewline
\begin{minipage}[t]{0.47\columnwidth}\raggedright
\(\mathsf{Locality}\)\strut
\end{minipage} & \begin{minipage}[t]{0.47\columnwidth}\raggedright
finite locality label, such as local, semi-global, or global use of the judgment data\strut
\end{minipage}\tabularnewline
\bottomrule
\end{longtable}

When individual fields of a profile must be named, the paper writes them as \(\mathsf{Scale}_{\lambda}\), \(\mathsf{Regime}_{\lambda}\), \(\mathsf{BridgeType}_{\lambda}\), \(\mathsf{InstrumentMode}_{\lambda}\), and \(\mathsf{Locality}_{\lambda}\).

Specific named profiles and their defining locations are listed in Subsubsection \hyperref[sn:3.6.1]{3.6.1}; each has the seven components displayed above.

Profile admissibility is the predicate

\[
\mathsf{AdmProfile}(\lambda,I,\mathcal T)
\]

on a profile \(\lambda\), an instrument \(I\), and a target package \(\mathcal T\). The admissibility predicate is finite and decidable on the finite components of \(\lambda\), \(I\), and \(\mathcal T\).

\begin{longtable}[]{@{}ll@{}}
\toprule
\begin{minipage}[b]{0.47\columnwidth}\raggedright
check\strut
\end{minipage} & \begin{minipage}[b]{0.47\columnwidth}\raggedright
admissibility condition\strut
\end{minipage}\tabularnewline
\midrule
\endhead
\begin{minipage}[t]{0.47\columnwidth}\raggedright
level\strut
\end{minipage} & \begin{minipage}[t]{0.47\columnwidth}\raggedright
\(\ell_i,\ell_j\in\mathsf{Lev}\) (a cross-level bridge is checked on the judgment's \(L\), not here)\strut
\end{minipage}\tabularnewline
\begin{minipage}[t]{0.47\columnwidth}\raggedright
scale, regime, locality\strut
\end{minipage} & \begin{minipage}[t]{0.47\columnwidth}\raggedright
\((\mathsf{Scale},\mathsf{Regime},\mathsf{Locality},h_{\mathcal T})\in\mathsf{Compat}_I\)\strut
\end{minipage}\tabularnewline
\begin{minipage}[t]{0.47\columnwidth}\raggedright
bridge type\strut
\end{minipage} & \begin{minipage}[t]{0.47\columnwidth}\raggedright
\(\mathsf{BridgeType}\) is admitted by \(\mathsf{BridgePolicy}_I\)\strut
\end{minipage}\tabularnewline
\begin{minipage}[t]{0.47\columnwidth}\raggedright
instrument mode\strut
\end{minipage} & \begin{minipage}[t]{0.47\columnwidth}\raggedright
\(\mathsf{InstrumentMode}=\mathsf{ModePolicy}_I(\text{surrounding judgment family})\)\strut
\end{minipage}\tabularnewline
\bottomrule
\end{longtable}

Profile mutation is not allowed without an explicit profile update record: a judgment under profile \(\lambda\) does not become a judgment under a different profile \(\lambda'\) unless a finite typed update record for that profile change is present in \(\mathsf{Cont}\) and passes the relevant visibility, source, and audit checks. This rule records the profile-relative discipline used later by the no-total-algebra theorem; it does not prove that theorem here.

\subsection{Full Status Tables}\label{sn:A.5}

The tables below give a brief gloss of each value of the seven status families; for the directed-cell, pair and claim families the glosses restate the classifier conditions of Subsubsection \hyperref[sn:5.6.6]{5.6.6}. The value sets are displayed in Subsubsections \hyperref[sn:4.2.2]{4.2.2}, \hyperref[sn:4.3.4]{4.3.4}, \hyperref[sn:4.5.3]{4.5.3} and \hyperref[sn:4.8.2]{4.8.2}, Subsection \hyperref[sn:4.4]{4.4}, and Subsubsections \hyperref[sn:5.1.6]{5.1.6} and \hyperref[sn:5.1.7]{5.1.7}; the family-specific classifier and priority rules are fixed in Section \hyperref[sn:5]{5} and Section \hyperref[sn:11]{11}. The rows below are not in priority order; for role, directed-cell, pair, promotion and claim records use Subsubsections \hyperref[sn:5.6.1]{5.6.1}--\hyperref[sn:5.6.5]{5.6.5}, respectively, and for square records use Subsubsection \hyperref[sn:5.6.6]{5.6.6}.

\subsubsection{Role Status}\label{sn:A.5.1}

\begin{longtable}[]{@{}ll@{}}
\toprule
\begin{minipage}[b]{0.33\columnwidth}\raggedright
value\strut
\end{minipage} & \begin{minipage}[b]{0.61\columnwidth}\raggedright
meaning\strut
\end{minipage}\tabularnewline
\midrule
\endhead
\begin{minipage}[t]{0.33\columnwidth}\raggedright
\(\texttt{active\_projection}\)\strut
\end{minipage} & \begin{minipage}[t]{0.61\columnwidth}\raggedright
role evidence \(R_i\) is present, in scope, visible and audited, \(\mathsf{collapse}_i\) is false and \(\Theta_i^{\mathrm{act}}\) is met; no profile is read\strut
\end{minipage}\tabularnewline
\begin{minipage}[t]{0.33\columnwidth}\raggedright
\(\texttt{below\_threshold}\)\strut
\end{minipage} & \begin{minipage}[t]{0.61\columnwidth}\raggedright
role evidence exists but does not meet activation threshold\strut
\end{minipage}\tabularnewline
\begin{minipage}[t]{0.33\columnwidth}\raggedright
\(\texttt{collapsed\_boundary\_case}\)\strut
\end{minipage} & \begin{minipage}[t]{0.61\columnwidth}\raggedright
the role-channel record's declared Boolean field \(\mathsf{collapse}_i\) is true\strut
\end{minipage}\tabularnewline
\begin{minipage}[t]{0.33\columnwidth}\raggedright
\(\texttt{absent\_with\_record}\)\strut
\end{minipage} & \begin{minipage}[t]{0.61\columnwidth}\raggedright
absence of the role/witness is audited\strut
\end{minipage}\tabularnewline
\begin{minipage}[t]{0.33\columnwidth}\raggedright
\(\texttt{inapplicable\_outside\_scope}\)\strut
\end{minipage} & \begin{minipage}[t]{0.61\columnwidth}\raggedright
instrument or host does not adjudicate this role\strut
\end{minipage}\tabularnewline
\bottomrule
\end{longtable}

The role/cell separation rule of Section \hyperref[sn:5]{5} records that an active role channel does not imply \(\texttt{action}\) on every cell carrying that role.

\subsubsection{Cell Status}\label{sn:A.5.2}

\begin{longtable}[]{@{}ll@{}}
\toprule
\begin{minipage}[b]{0.22\columnwidth}\raggedright
value\strut
\end{minipage} & \begin{minipage}[b]{0.72\columnwidth}\raggedright
meaning\strut
\end{minipage}\tabularnewline
\midrule
\endhead
\begin{minipage}[t]{0.22\columnwidth}\raggedright
\(\texttt{action}\)\strut
\end{minipage} & \begin{minipage}[t]{0.72\columnwidth}\raggedright
no earlier row applies; \(W\) and \(U\) are present; the cell is well-formed, every defect entry meets its threshold, the visibility defect is empty, and \(\delta_6^{\mathrm{audit}}=\texttt{audit\_passes}\)\strut
\end{minipage}\tabularnewline
\begin{minipage}[t]{0.22\columnwidth}\raggedright
\(\texttt{implicit}\)\strut
\end{minipage} & \begin{minipage}[t]{0.72\columnwidth}\raggedright
\(\mathsf{implicitW}\): \(W\) is attached to \(\mathcal T\) and no threshold of \(\Theta\) evaluates \(W\) (a computed condition)\strut
\end{minipage}\tabularnewline
\begin{minipage}[t]{0.22\columnwidth}\raggedright
\(\texttt{trivial}\)\strut
\end{minipage} & \begin{minipage}[t]{0.72\columnwidth}\raggedright
\(\mathsf{noopU}\) holds, or \(\delta\) contains \(\texttt{identity\_package}\)\strut
\end{minipage}\tabularnewline
\begin{minipage}[t]{0.22\columnwidth}\raggedright
\(\texttt{blocked}\)\strut
\end{minipage} & \begin{minipage}[t]{0.72\columnwidth}\raggedright
a required bridge is missing or weak, an overread, weak-bridge or unknown visibility defect is nonempty (an outside-scope use is caught by the first row), or the audit defect contains \(\texttt{source\_failure}\), \(\texttt{visibility\_failure}\), \(\texttt{missing\_audit}\), \(\texttt{failed\_audit}\) or \(\texttt{nonclaim\_failure}\)\strut
\end{minipage}\tabularnewline
\begin{minipage}[t]{0.22\columnwidth}\raggedright
\(\texttt{undefined\_circular}\)\strut
\end{minipage} & \begin{minipage}[t]{0.72\columnwidth}\raggedright
\(\texttt{circular\_audit}\in\delta_6^{\mathrm{audit}}\)\strut
\end{minipage}\tabularnewline
\begin{minipage}[t]{0.22\columnwidth}\raggedright
\(\texttt{below\_threshold}\)\strut
\end{minipage} & \begin{minipage}[t]{0.72\columnwidth}\raggedright
\(W\) and \(U\) are present and \(\texttt{threshold\_failure}\in\delta_6^{\mathrm{audit}}\)\strut
\end{minipage}\tabularnewline
\begin{minipage}[t]{0.22\columnwidth}\raggedright
\(\texttt{collapsed}\)\strut
\end{minipage} & \begin{minipage}[t]{0.72\columnwidth}\raggedright
\(\delta\) contains \(\texttt{total\_collapse}\) or \(\texttt{singleton\_collapse}\)\strut
\end{minipage}\tabularnewline
\begin{minipage}[t]{0.22\columnwidth}\raggedright
\(\texttt{absent}\)\strut
\end{minipage} & \begin{minipage}[t]{0.72\columnwidth}\raggedright
\(W\) or \(U\) is an audited absence marker, and \(\delta_6^{\mathrm{audit}}\) contains neither \(\texttt{missing\_audit}\) nor \(\texttt{failed\_audit}\)\strut
\end{minipage}\tabularnewline
\begin{minipage}[t]{0.22\columnwidth}\raggedright
\(\texttt{outside\_scope}\)\strut
\end{minipage} & \begin{minipage}[t]{0.72\columnwidth}\raggedright
the host tag is not in \(\mathsf{Hosts}_I\), a level tag is not in \(\mathsf{Levels}_I\), or \(\delta_{\mathrm{outside}}(I,C)=\mathsf{Use}(C)\cap\mathsf{Outside}_I\neq\varnothing\)\strut
\end{minipage}\tabularnewline
\bottomrule
\end{longtable}

\subsubsection{Pair Status}\label{sn:A.5.3}

\begin{longtable}[]{@{}ll@{}}
\toprule
\begin{minipage}[b]{0.19\columnwidth}\raggedright
value\strut
\end{minipage} & \begin{minipage}[b]{0.75\columnwidth}\raggedright
meaning\strut
\end{minipage}\tabularnewline
\midrule
\endhead
\begin{minipage}[t]{0.19\columnwidth}\raggedright
\(\texttt{real}\)\strut
\end{minipage} & \begin{minipage}[t]{0.75\columnwidth}\raggedright
pair observable has source-of-truth, compatibility, visibility, threshold, and audit\strut
\end{minipage}\tabularnewline
\begin{minipage}[t]{0.19\columnwidth}\raggedright
\(\texttt{real\_provisional}\)\strut
\end{minipage} & \begin{minipage}[t]{0.75\columnwidth}\raggedright
the source is a fallback or proxy source that \(\mathsf{SourcePolicy}_I\) admits for provisional use under \(\mu\), the visibility defect is empty, \(\mathsf{partialEvidence}\) is false, and every bridge this row needs is admissible (Subsubsection \hyperref[sn:5.6.3]{5.6.3})\strut
\end{minipage}\tabularnewline
\begin{minipage}[t]{0.19\columnwidth}\raggedright
\(\texttt{real\_duplicated}\)\strut
\end{minipage} & \begin{minipage}[t]{0.75\columnwidth}\raggedright
the conditions of \(\texttt{real}\) hold and the record carries an explicit duplicate-pair record\strut
\end{minipage}\tabularnewline
\begin{minipage}[t]{0.19\columnwidth}\raggedright
\(\texttt{blocked}\)\strut
\end{minipage} & \begin{minipage}[t]{0.75\columnwidth}\raggedright
an earlier blocking condition holds, including incompatible branches or a branch audit other than \(\texttt{audit\_passes}\), or no preceding pair row applies\strut
\end{minipage}\tabularnewline
\bottomrule
\end{longtable}

The pair-realness rule of Section \hyperref[sn:5]{5} requires \(\delta_6^{\mathrm{source}}=\texttt{committed}\), or \(\texttt{fallback\_admitted}\) together with an audited source-upgrade bridge; on the PICA host failure of \(\mathsf{SourceOfTruthOK}\) fails the source threshold and blocks the pair.

\subsubsection{Promotion Status}\label{sn:A.5.4}

The values \(\texttt{strict}\) and \(\texttt{non\_strict}\) are accepted-status refinements: a \(\texttt{strict}\) verdict implies the accepted core and \(\pi_{j+1}\not\factor\pi_{j}\), and a \(\texttt{non\_strict}\) verdict implies the accepted core and \(\pi_{j+1}\factor\pi_{j}\), in keeping with Theorems \hyperref[cl:thm:19]{19}, \hyperref[cl:thm:20]{20}, \hyperref[cl:thm:21]{21} of Section \hyperref[sn:10]{10}.

\subsubsection{Claim Status}\label{sn:A.5.5}

\begin{longtable}[]{@{}ll@{}}
\toprule
\begin{minipage}[b]{0.26\columnwidth}\raggedright
value\strut
\end{minipage} & \begin{minipage}[b]{0.68\columnwidth}\raggedright
meaning\strut
\end{minipage}\tabularnewline
\midrule
\endhead
\begin{minipage}[t]{0.26\columnwidth}\raggedright
\(\texttt{accepted}\)\strut
\end{minipage} & \begin{minipage}[t]{0.68\columnwidth}\raggedright
all required checks pass\strut
\end{minipage}\tabularnewline
\begin{minipage}[t]{0.26\columnwidth}\raggedright
\(\texttt{rejected}\)\strut
\end{minipage} & \begin{minipage}[t]{0.68\columnwidth}\raggedright
the claim is well-formed and the check rules of the instrument assign a rejection verdict\strut
\end{minipage}\tabularnewline
\begin{minipage}[t]{0.26\columnwidth}\raggedright
\(\texttt{provisional}\)\strut
\end{minipage} & \begin{minipage}[t]{0.68\columnwidth}\raggedright
\(\mathsf{partialEvidence}\) is true and the claim type lies in \(\mathsf{ProvisionalTypes}_I\)\strut
\end{minipage}\tabularnewline
\begin{minipage}[t]{0.26\columnwidth}\raggedright
\(\texttt{blocked}\)\strut
\end{minipage} & \begin{minipage}[t]{0.68\columnwidth}\raggedright
in scope, with an unbridged suppressed or unknown use, an evidence entry that is not an admitted source passing \(\Theta_{\mathrm{source}}\), or a failed audit source or visibility check; also the final default\strut
\end{minipage}\tabularnewline
\begin{minipage}[t]{0.26\columnwidth}\raggedright
\(\texttt{outside\_scope}\)\strut
\end{minipage} & \begin{minipage}[t]{0.68\columnwidth}\raggedright
instrument or host does not adjudicate claim\strut
\end{minipage}\tabularnewline
\begin{minipage}[t]{0.26\columnwidth}\raggedright
\(\texttt{absent\_with\_record}\)\strut
\end{minipage} & \begin{minipage}[t]{0.68\columnwidth}\raggedright
a record referred to by the content of \(\varphi\), other than an evidence source, is absent and the absence is audited (an absent evidence source gives \(\texttt{blocked}\) unless the earlier \(\texttt{outside\_scope}\) row applies)\strut
\end{minipage}\tabularnewline
\begin{minipage}[t]{0.26\columnwidth}\raggedright
\(\texttt{below\_threshold}\)\strut
\end{minipage} & \begin{minipage}[t]{0.68\columnwidth}\raggedright
evidence exists but threshold fails\strut
\end{minipage}\tabularnewline
\begin{minipage}[t]{0.26\columnwidth}\raggedright
\(\texttt{failed\_audit}\)\strut
\end{minipage} & \begin{minipage}[t]{0.68\columnwidth}\raggedright
required audit fails\strut
\end{minipage}\tabularnewline
\begin{minipage}[t]{0.26\columnwidth}\raggedright
\(\texttt{undefined\_circular}\)\strut
\end{minipage} & \begin{minipage}[t]{0.68\columnwidth}\raggedright
the use-set of the claim reaches its own verdict reference by resolved references, or \(\texttt{circular\_audit}\in\delta_6^{\mathrm{audit}}\), or the claim type is \(\texttt{soundness}\) and some \(\ulcorner\psi,I'\urcorner\in\mathsf{Use}(\varphi)\) has \(\mathsf{Level}(I')\ge\mathsf{Level}(I)\)\strut
\end{minipage}\tabularnewline
\bottomrule
\end{longtable}

The full priority order of \(\mathrm{ClaimStatus}\) is fixed in Subsubsection \hyperref[sn:5.6.5]{5.6.5}.

\subsubsection{Gate Status}\label{sn:A.5.6}

The gate status set \(\mathrm{GateStatus}\) and its refinement \(\mathrm{StrictGateStatus}\) for the strictness gate are displayed in Subsubsection \hyperref[sn:5.1.6]{5.1.6}. The core promotion gates are \(\mathsf{CoreGates}=\{G_{\mathrm{suff}},G_{\mathrm{desc}},G_{\mathrm{stab}},G_{\mathrm{ctrl}},G_{\mathrm{nosmuggle}},G_{\mathrm{vis}},G_{\mathrm{audit}}\}\), with the additional gate \(G_{\mathrm{strict}}\) and the optional gate \(G_{\mathrm{locglob}}\) (Subsection \hyperref[sn:10.2]{10.2}).

\subsubsection{Square Status}\label{sn:A.5.7}

\begin{longtable}[]{@{}ll@{}}
\toprule
\begin{minipage}[b]{0.28\columnwidth}\raggedright
value\strut
\end{minipage} & \begin{minipage}[b]{0.66\columnwidth}\raggedright
meaning\strut
\end{minipage}\tabularnewline
\midrule
\endhead
\begin{minipage}[t]{0.28\columnwidth}\raggedright
\(\texttt{exact}\)\strut
\end{minipage} & \begin{minipage}[t]{0.66\columnwidth}\raggedright
square defect is empty and exact comparison holds\strut
\end{minipage}\tabularnewline
\begin{minipage}[t]{0.28\columnwidth}\raggedright
\(\texttt{lax}\)\strut
\end{minipage} & \begin{minipage}[t]{0.66\columnwidth}\raggedright
the decoration declares an order comparison, \(\mathsf{partialEvidence}_\Xi\) is false, \(\delta_\Xi=\varnothing\), and equality of the two boundary composites fails\strut
\end{minipage}\tabularnewline
\begin{minipage}[t]{0.28\columnwidth}\raggedright
\(\texttt{obstructed}\)\strut
\end{minipage} & \begin{minipage}[t]{0.66\columnwidth}\raggedright
square is in scope and audited but defect is nonempty\strut
\end{minipage}\tabularnewline
\begin{minipage}[t]{0.28\columnwidth}\raggedright
\(\texttt{blocked}\)\strut
\end{minipage} & \begin{minipage}[t]{0.66\columnwidth}\raggedright
nonempty visibility defect, or audit defect containing \(\texttt{source\_failure}\) or \(\texttt{visibility\_failure}\), or no earlier row completes a per-type check\strut
\end{minipage}\tabularnewline
\begin{minipage}[t]{0.28\columnwidth}\raggedright
\(\texttt{outside\_scope}\)\strut
\end{minipage} & \begin{minipage}[t]{0.66\columnwidth}\raggedright
host/instrument does not adjudicate square\strut
\end{minipage}\tabularnewline
\begin{minipage}[t]{0.28\columnwidth}\raggedright
\(\texttt{failed\_audit}\)\strut
\end{minipage} & \begin{minipage}[t]{0.66\columnwidth}\raggedright
square audit fails\strut
\end{minipage}\tabularnewline
\begin{minipage}[t]{0.28\columnwidth}\raggedright
\(\texttt{undefined\_circular}\)\strut
\end{minipage} & \begin{minipage}[t]{0.66\columnwidth}\raggedright
square support/audit is circular\strut
\end{minipage}\tabularnewline
\begin{minipage}[t]{0.28\columnwidth}\raggedright
\(\texttt{provisional}\)\strut
\end{minipage} & \begin{minipage}[t]{0.66\columnwidth}\raggedright
\(\mathsf{partialEvidence}_\Xi\) is true and \(\mathrm{type}_\Xi\in\mathsf{ProvisionalTypes}_I\)\strut
\end{minipage}\tabularnewline
\bottomrule
\end{longtable}

The seven status families are typed-separately judgment-indexed and remain non-collapsing, in keeping with Theorem \hyperref[cl:thm:5]{5} of Section \hyperref[sn:6]{6} and the status-discipline of Section \hyperref[sn:5]{5}.

\subsection{Full Defect Tables}\label{sn:A.6}

The general defect record schema and the per-primitive and promotion defect tables are recorded below.

\subsubsection{General Defect Record Schema}\label{sn:A.6.1}

A typed defect record under instrument \(I\) on a target package \(\mathcal T\) has the schema

\[
\mathsf{Def}_{i}^{\lambda}(\mathcal T,I)\;=\;(\,\delta_{i},\;\Omega_{i},\;\Theta_{i},\;\mathrm{polarity}_{i},\;W_{i},\;A_{i},\;V_{i},\;\sigma_{i}\,),
\]

with components:

\begingroup
\setlength{\tabcolsep}{4pt}
\renewcommand{\arraystretch}{1.08}
\begin{center}
\begin{tabularx}{0.96\linewidth}{@{}>{\raggedright\arraybackslash}p{0.12\linewidth}>{\raggedright\arraybackslash}X@{}}
\toprule
field & meaning \\
\midrule
$\delta_{i}$ & defect value (set, function, scalar, or other finite typed value) \\
$\Omega_{i}$ & value object / defect type \\
$\Theta_{i}$ & threshold or pass condition \\
$\mathrm{polarity}_{i}$ & obstruction, certificate, gap, residual, capacity loss, or other polarity tag \\
$W_{i}$ & witness, certificate, or counterexample for $\delta_{i}$ \\
$A_{i}$ & audit record for $\delta_{i}$ \\
$V_{i}$ & visibility record for $\delta_{i}$ \\
$\sigma_{i}$ & induced status or status annotation \\
\bottomrule
\end{tabularx}
\end{center}
\endgroup

The general threshold rule on a defect record is

\[
\mathsf{Eval}_{\Theta_{i}}(\delta_{i})\;\in\;\{\,\texttt{pass},\;\texttt{fail},\;\texttt{not\_checked},\;\texttt{outside\_scope}\,\}.
\]

\subsubsection{Threshold Types}\label{sn:A.6.2}

The admissible thresholds form the finite typed family

\[
\Theta\;\in\;\{\,\Theta_{=0},\;\Theta_{\neq0},\;\Theta_{\le\varepsilon},\;\Theta_{\ge\varepsilon},\;\Theta_{\in S},\;\Theta_{\notin S},\;\Theta_{\mathrm{pass}},\;\Theta_{\mathrm{source}},\;\Theta_{\mathrm{bridge}},\;\Theta_{\mathrm{audit}},\;\Theta_{\mathrm{vis}}\,\},
\]

with intended readings recorded in Subsection \hyperref[sn:5.2]{5.2}.

\subsubsection{Primitive Defect Families}\label{sn:A.6.3}

The per-primitive defect records are as follows. Each row is a finite typed instance of the general defect schema of Subsubsection \hyperref[sn:A.6.1]{A.6.1}.

\begingroup
\setlength{\tabcolsep}{4pt}
\renewcommand{\arraystretch}{1.08}
\setlength{\LTleft}{0pt}
\setlength{\LTright}{0pt}
\begin{longtable}{@{}>{\raggedright\arraybackslash}p{0.11\linewidth}>{\raggedright\arraybackslash}p{0.49\linewidth}>{\raggedright\arraybackslash}p{0.34\linewidth}@{}}
\toprule
primitive & defect record & threshold / status use \\
\midrule
\endhead
$P_1$ & \(\begin{aligned}[t]\delta_1^{\mathrm{split}}(q,F)&=\{\,(x,x'):\ q(x)=q(x'),\\ &\quad qF(x)\neq qF(x')\,\}\end{aligned}\) & $\Theta_1^{\mathrm{desc}}:\delta_1^{\mathrm{split}}=\varnothing$; nonempty records an active $P_1$ obstruction \\
$P_1$ & \(\begin{aligned}[t]\delta_1^{\mathrm{cand}}(q,F,F^\sharp)&=\{\,x:\ F^\sharp(q_{\mathrm{im}}(x))\\ &\quad \neq q_{\mathrm{im}}F(x)\,\}\end{aligned}\) & candidate descended-map defect \\
$P_1$ & \(\begin{aligned}[t]\mathrm{CD}_{\tau}^{\mu}(\Pi)&=I(X_t;\Pi(X_{t+\tau})\mid\Pi(X_t))\end{aligned}\) & $\Theta_1^{\mathrm{closure}}:\mathrm{CD}_{\tau}^{\mu}(\Pi)=0$ for exact closure, or $\mathrm{CD}_{\tau}^{\mu}(\Pi)\le\varepsilon$ for approximate closure \\
$P_2$ & \(\delta_2^{\mathrm{rep}}(s;f)=\begin{cases}0,&s\in\Sigma_f,\\ 1,&s\notin\Sigma_f,\end{cases}\) & in full finite-lens mode, $s\in\Sigma_f$ iff $s=f^{-1}(B)$ for some $B\subseteq\mathrm{im}(f)$; nonrepresentability is witnessed by $z,z'$ with $f(z)=f(z')$ and $\mathbf 1_s(z)\neq\mathbf 1_s(z')$ \\
$P_2$ & $\delta_2^{\mathrm{cycle}}(G,G')=\beta_1(G)-\beta_1(G')$ & drive-suppression threshold $\Theta_2^{\mathrm{forest}}:\beta_1(G')=0$ \\
$P_3$ & \(\Delta_3^{\mathrm{comp}}(E_1,E_2)=\{\,c:\ E_1E_2(c)\neq E_2E_1(c)\,\}\) & route mismatch is active when $\Delta_3^{\mathrm{comp}}\neq\varnothing$; route coherence holds when $\Delta_3^{\mathrm{comp}}=\varnothing$ \\
$P_3$ & $\delta_3^{\mathrm{hol}}(\gamma)=\{\,\phi\in\Phi:\operatorname{Hol}(\gamma)(\phi)\neq\phi\,\}$ & active when $\operatorname{Hol}(\gamma)\neq\mathrm{id}$; no $P6_{\mathrm{drive}}$ consequence follows without a drive bridge \\
$P_4$ & \(\delta_4^{\mathrm{ref}}(f_0,f_1,\rho)=\{\,z:\ f_0(z)\neq\rho f_1(z)\,\}\) & refinement passes iff $\delta_4^{\mathrm{ref}}=\varnothing$ \\
$P_4$ & \(\begin{aligned}[t]\delta_4^{\mathrm{glue}}=\{\,x\in U\cap V:\ &\mathrm{res}_{U,U\cap V}(q_U(x))\\ &\neq\mathrm{res}_{V,U\cap V}(q_V(x))\,\}\end{aligned}\) & gluing passes iff $\delta_4^{\mathrm{glue}}=\varnothing$; failed gluing may induce $\texttt{globally\_obstructed}$ \\
$P_5$ & \(\delta_5^{\mathrm{idem}}(E)=\{\,x:\ E^2(x)\neq E(x)\,\}\) & $\Theta_5^{\mathrm{idem}}:\delta_5^{\mathrm{idem}}=\varnothing$; if $U$ applies idempotent $E$ to a recorded $s\in\operatorname{im}E$, then $e_U(s)=s$, so the cell is $\texttt{trivial}$ unless an earlier row applies \\
$P_5$ & $\delta_5^{\mathrm{fix}}(E)=\{\,o:\ E(o)\neq o\,\}$; on a host with a declared norm, $\|E(o)-o\|$ is a pointwise residual & fixed-point residual; an approximate threshold uses a declared scalar residual \\
$P_5$ & package collapse values such as $\texttt{identity\_package}$, $\texttt{total\_collapse}$, and $\texttt{singleton\_collapse}$ & $\texttt{total\_collapse}$ and $\texttt{singleton\_collapse}$ induce $\texttt{collapsed}$; $\texttt{identity\_package}$ induces $\texttt{trivial}$ \\
$P_6$ & $\delta_6^{\mathrm{audit}}\subseteq\{\,\texttt{missing\_audit},\allowbreak\,\texttt{failed\_audit},\allowbreak\,\texttt{source\_failure},\allowbreak\,\texttt{visibility\_failure},\allowbreak\,\texttt{threshold\_failure},\allowbreak\,\texttt{nonclaim\_failure},\allowbreak\,\texttt{circular\_audit}\,\}$ & $\Theta_6^{\mathrm{audit}}$: the $\mathsf{CheckRules}$ field of $A$ is drawn from $\mathsf{CheckRules}_I$ and $\delta_6^{\mathrm{audit}}(A)=\varnothing$ (written $\texttt{audit\_passes}$); classifier priority decides the status when several codes occur, and circular audit induces $\texttt{undefined\_circular}$ \\
$P_6$ & $\delta_6^{\mathrm{source}}\in\{\,\texttt{committed},\allowbreak\,\texttt{fallback\_admitted},\allowbreak\,\texttt{proxy\_admitted},\allowbreak\,\texttt{dirty},\allowbreak\,\texttt{fallback\_forbidden},\allowbreak\,\texttt{unknown},\allowbreak\,\texttt{contradictory},\allowbreak\,\texttt{missing}\,\}$ & $\Theta_{\mathrm{source}}$ passes at $\texttt{committed}$, $\texttt{fallback\_admitted}$ and $\texttt{proxy\_admitted}$; real pair claims require $\texttt{committed}$, or $\texttt{fallback\_admitted}$ with an audited source-upgrade bridge \\
$P_6$ & drive certificate $\kappa_6^{\mathrm{drive}}$, kept separate from $\delta_6^{\mathrm{audit}}$ & cohomological drive uses $[a]\neq0\in H^1(G)$, equivalently $\exists\gamma:\oint_{\gamma}a\neq0$; each of $[a]=0$, $a=d\phi$, $H^1(G)=0$ and $\beta_1(G)=0$ yields the no-drive certificate of Subsubsection \hyperref[sn:8.1.2]{8.1.2} \\
\bottomrule
\end{longtable}
\endgroup

\subsubsection{Promotion-Specific Defects}\label{sn:A.6.4}

The promotion bridge \(B_{j\to j+1}\) carries the promotion-specific defects recorded in Subsection \hyperref[sn:5.4]{5.4}, in addition to any per-primitive defect contributions of its fields.

\begingroup
\setlength{\tabcolsep}{4pt}
\renewcommand{\arraystretch}{1.08}
\begin{center}
\begin{tabularx}{0.96\linewidth}{@{}>{\raggedright\arraybackslash}p{0.44\linewidth}>{\raggedright\arraybackslash}X@{}}
\toprule
defect & threshold / status use \\
\midrule
\(\begin{aligned}[t]\Delta_{\mathrm{fact}}(\pi_j,\pi_{j+1})&=\{\,(s,s'):\pi_j(s)=\pi_j(s'),\\ &\quad \pi_{j+1}(s)\neq\pi_{j+1}(s')\,\}\end{aligned}\) & $\Theta_{\mathrm{strict}}:\Delta_{\mathrm{fact}}\neq\varnothing$ gives $G_{\mathrm{strict}}=\texttt{strict\_pass}$; $\Theta_{\mathrm{nonstrict}}:\Delta_{\mathrm{fact}}=\varnothing$ gives $G_{\mathrm{strict}}=\texttt{non\_strict\_pass}$ \\
$\delta_{\mathrm{nosmuggle}}(B)$ & contains violations: target-layer data used to prove target novelty; suppressed content used without bridge; simulation used as theorem; local evidence used as global without gluing; $P_3$-witness used as $P6_{\mathrm{drive}}$ without drive bridge; lossy summary used as exact data. Threshold $\Theta_{\mathrm{nosmuggle}}:\delta_{\mathrm{nosmuggle}}=\varnothing$; failure induces $\mathrm{Promote}(B):\texttt{failed\_no\_smuggling}$ \\
$\delta_{\mathrm{stab}}=\{o\in\mathrm{im}(\pi_{j+1}):E(\pi_{j+1}^{-1}(o))\neq\pi_{j+1}^{-1}(o)\}$, or $\{o:E(\iota o)\neq\iota o\}$ via the recorded embedding $\iota$, or its approximate residual analogue & if required and failed, induces $\mathrm{Promote}(B):\texttt{failed\_stability}$ \\
descent defect inside promotion & if macro closure is claimed, the promotion requires $G_{\mathrm{desc}}=\texttt{pass}$; failure induces $\mathrm{Promote}(B):\texttt{failed\_descent}$; if no dynamics or closure claim is made, $G_{\mathrm{desc}}=\texttt{not\_required}$ \\
\bottomrule
\end{tabularx}
\end{center}
\endgroup

Each defect record above conforms to the general schema of Subsubsection \hyperref[sn:A.6.1]{A.6.1}, with the polarity tag, threshold, and induced promotion status drawn from the schemas of Subsection \hyperref[sn:5.4]{5.4} and the promotion classifier of Subsubsection \hyperref[sn:5.6.4]{5.6.4}.
 \section{Long Proofs}\label{sn:B}

This appendix records long-form proofs that are referenced from the main text. The proofs of Sections \hyperref[sn:6]{6}, \hyperref[sn:7]{7}, \hyperref[sn:8]{8}, \hyperref[sn:10]{10}, \hyperref[sn:11]{11}, \hyperref[sn:12]{12}, and \hyperref[sn:14]{14} are written in full where feasible inside the main text; this appendix supplements them with the longer technical content of the closure-deficit theorem, the cohomological theorems of Section \hyperref[sn:8]{8} and the fixed-interface definability bound of Section \hyperref[sn:7]{7}; Subsection \hyperref[sn:B.3]{B.3} points to the main-text proofs of Theorems \hyperref[cl:thm:9]{9} and \hyperref[cl:thm:31]{31}.

Each of Subsections \hyperref[sn:B.1]{B.1}, \hyperref[sn:B.2]{B.2} and \hyperref[sn:B.4]{B.4} records the long-form proof of a main-text result, ends with the instrument-indexed form of that result (listed for all numbered results in Appendix \hyperref[sn:I]{I}), and adds technical lemmas only where the main-text proof would otherwise become unreadable. The full statements remain at their main-text locations; this appendix is a pointer-and-extension document, not a re-statement of theorems.

\subsection{Closure-Deficit Proof}\label{sn:B.1}

The closure-deficit theorem of Subsection \hyperref[sn:7.3]{7.3} (Theorem \hyperref[cl:thm:8]{8}) is the finite Markov/KL result on the host \(\mathsf H_{\mathrm{prob}}\). The main-text proof gives the decomposition; this subsection records the same calculation in long form and identifies the associated \(P_1\) closure-deficit record. It does not introduce a new theorem.

\emph{Statement (recall).} Fix a finite state space \(X\), a finite stochastic kernel on \(X\), a fixed distribution, a finite abstraction \(\Pi:X\to Y\), a time lag \(\tau\), the declared KL predictive closure loss \(\mathcal L_\tau\), and an admissible macro-kernel class \(\mathcal K_{\mathrm{adm}}\) containing the conditional macro kernel

\[
K^\ast(\cdot\mid y)=P(\Pi(X_{t+\tau})\mid\Pi(X_t)=y)
\]

on the support of \(\Pi(X_t)\). Then

\[
\min_{K\in\mathcal K_{\mathrm{adm}}}\mathcal L_\tau(K)
\;=\;
I\bigl(X_t;\Pi(X_{t+\tau})\mid\Pi(X_t)\bigr).
\]

\emph{Proof.} Write

\[
Y_t=\Pi(X_t),\qquad Y'=\Pi(X_{t+\tau}).
\]

For a macro kernel \(K:Y\to\Delta(Y)\), the declared loss is

\[
\mathcal L_\tau(K)
\;=\;
\mathbb E\!\left[
D_{\mathrm{KL}}\bigl(P(Y'\mid X_t)\,\|\,K(\cdot\mid Y_t)\bigr)
\right].
\]

Because the state space is finite, the expectation decomposes over the finite support of \(Y_t\). For any \(y\) with \(P(Y_t=y)>0\), expand the conditional loss:

\[
\mathbb E\!\left[
D_{\mathrm{KL}}\bigl(P(Y'\mid X_t)\,\|\,K(\cdot\mid y)\bigr)
\mid Y_t=y
\right].
\]

Insert and subtract the conditional macro law \(P(Y'\mid Y_t=y)\). The finite KL projection identity gives

\[
\begin{aligned}
&\mathbb E\!\left[
D_{\mathrm{KL}}\bigl(P(Y'\mid X_t)\,\|\,K(\cdot\mid y)\bigr)
\mid Y_t=y
\right] \\
&\quad =
\mathbb E\!\left[
D_{\mathrm{KL}}\bigl(P(Y'\mid X_t)\,\|\,P(Y'\mid Y_t=y)\bigr)
\mid Y_t=y
\right]
+D_{\mathrm{KL}}\bigl(P(Y'\mid Y_t=y)\,\|\,K(\cdot\mid y)\bigr).
\end{aligned}
\]

The first term is independent of \(K\). The second term is non-negative and is minimized, with value zero, by

\[
K^\ast(\cdot\mid y)=P(Y'\mid Y_t=y),
\]

which belongs to \(\mathcal K_{\mathrm{adm}}\) by hypothesis. Therefore the minimizing loss is the expectation of the first term:

\[
\min_{K\in\mathcal K_{\mathrm{adm}}}\mathcal L_\tau(K)
\;=\;
\mathbb E\!\left[
D_{\mathrm{KL}}\bigl(P(Y'\mid X_t)\,\|\,P(Y'\mid Y_t)\bigr)
\right].
\]

For finite variables this expectation is the conditional mutual information \(I(X_t;Y'\mid Y_t)\). Substituting \(Y_t=\Pi(X_t)\) and \(Y'=\Pi(X_{t+\tau})\) proves the equality.

Support-zero cases are ignored by restricting the calculation to the finite support of \(Y_t\), equivalently by the standard convention that zero-probability conditioning events do not contribute to the expectation. The finite-state hypothesis makes every sum a finite sum of terms in \([0,+\infty]\) (as in Theorem \hyperref[cl:thm:8]{8}), and the minimizer \(K^\ast\) a finite macro kernel with finite loss.

The associated \(P_1\) closure-deficit value is

\[
\mathrm{CD}_{\tau}^{\mu}(\Pi)
\;=\;
I\bigl(X_t;\Pi(X_{t+\tau})\mid\Pi(X_t)\bigr),
\]

with exact-closure threshold \(\Theta_1^{\mathrm{closure}}:\mathrm{CD}_{\tau}^{\mu}(\Pi)=0\) and approximate threshold \(\mathrm{CD}_{\tau}^{\mu}(\Pi)\le\varepsilon\), as in the defect schema of Subsubsection \hyperref[sn:5.3.1]{5.3.1}. Written as a defect record,

\[
\delta_{\mathrm{cl}}^{\mathrm{prob}}
\;=\;
(\,\mathrm{CD}_{\tau}^{\mu}(\Pi),\;\Omega_{\mathrm{cl}},\;\Theta_1^{\mathrm{closure}},\;\mathrm{polarity}_{\mathrm{gap}},\;W_{\mathrm{cl}},\;A_{\mathrm{prob}},\;V_{\mathrm{prob}},\;\sigma_{\mathrm{cl}}\,),
\]

where \(W_{\mathrm{cl}}\) is the finite collection of conditional-distribution entries used in the KL calculation, \(A_{\mathrm{prob}}\) is the host audit record, \(V_{\mathrm{prob}}\) is the visibility record, and \(\sigma_{\mathrm{cl}}\) is the status annotation assigned by the later classifier. This record is a supporting defect record, not an additional theorem statement.

The instrument-indexed verdict remains the main-text theorem line

\[
\Gamma\,;\;\mathsf H_{\mathrm{prob}}\,;\;I_{\mathrm{prob}}\;\vdash\;\textsc{FiniteMarkovClosureDeficit}\;:\;\texttt{accepted}.
\]

\subsection{Graph/Cohomology Proofs}\label{sn:B.2}

This subsection records long-form proofs of the four graph/cohomology theorems of Section \hyperref[sn:8]{8}: forest no-drive (Theorem \hyperref[cl:thm:15]{15}), exact-form null-drive (Theorem \hyperref[cl:thm:16]{16}), nonzero-affinity equivalence (Theorem \hyperref[cl:thm:17]{17}), and gating-affinity suppression (Theorem \hyperref[cl:thm:18]{18}). The main-text proofs in Section \hyperref[sn:8]{8} are sufficient for the headline statements; this subsection records the long-form arguments and the associated \(P_2\) and \(\mathrm{P6}_{\mathrm{drive}}^{H^{1}}\) defect bookkeeping.

\subsubsection{Forest No-Drive (Theorem 15)}\label{sn:B.2.1}

\emph{Statement (recall).} Let \(G=(V,E)\) be a finite graph with first Betti number \(\beta_{1}(G)=0\), that is, \(G\) is a forest. Let \(C^{0}(G;k)\), \(C^{1}(G;k)\), and \(d:C^{0}(G;k)\to C^{1}(G;k)\) be the finite-dimensional (finitely generated free) cochain complex over the declared coefficient field or free coefficient host \(k\). Then \(Z_{1}(G;k)=0\) and \(H^{1}(G;k)=0\).

\emph{Proof.} For a finite graph over the declared coefficient field, or over a declared free coefficient host with the corresponding free-rank interpretation, finite graph cohomology satisfies

\[
H^{1}(G;k)\cong\mathrm{Hom}_k(H_{1}(G;k),k),
\]

by the universal coefficient theorem: the Ext term vanishes because \(H_0(G;\mathbb Z)\) is free, and \(H_1(G;k)=H_1(G;\mathbb Z)\otimes k\) is free since \(H_1(G;\mathbb Z)\) is a free abelian group of rank \(\beta_1(G)\) (the proof of Theorem \hyperref[cl:thm:17]{17} gives a direct argument by spanning forests).

The first Betti number is the finite dimension, or free rank, of \(H_1(G;k)\):

\[
\dim H_1(G;k)=\beta_1(G).
\]

If \(\beta_1(G)=0\), then \(H_1(G;k)=0\); since \(G\) has no 2-cells, \(H_1(G;k)=Z_1(G;k)\), so \(Z_1(G;k)=0\). Therefore \(\mathrm{Hom}(H_1(G;k),k)=0\), and hence \(H^1(G;k)=0\). Equivalently, a forest has no nontrivial cycle space, so there is no nonzero cycle-supported first cohomology.

The associated no-drive defect record is

\[
\delta_{6}^{\mathrm{forest}}\;=\;(\,H^{1}(G;k),\;\Omega_{\mathrm{drive}},\;\Theta_{=0},\;\mathrm{polarity}_{\mathrm{no\text{-}drive}},\;W_{\mathrm{forest}},\;A_{\mathrm{graph}},\;V_{\mathrm{graph}},\;\sigma_{\mathrm{no\text{-}drive}}\,),
\]

with \(W_{\mathrm{forest}}\) the finite forest-witness data (the absence of cycles in \(G\)).

The instrument-indexed verdict is

\[
\Gamma\,;\;\mathsf H_{\mathrm{graph}}\,;\;I_{\mathrm{graph}}\;\vdash\;\textsc{ForestNoDrive}\;:\;\texttt{accepted}.\quad\square
\]

\subsubsection{Exact-Form Null-Drive (Theorem 16)}\label{sn:B.2.2}

\emph{Statement (recall).} Let \(a\in C^{1}(G;k)\) be a \(1\)-cochain with \(a=d\phi\) for some \(\phi\in C^{0}(G;k)\). Then for every cycle \(\gamma\in\mathsf{Cycles}(G)\), the cycle integral satisfies \(\oint_{\gamma}a=0\).

\emph{Proof.} By definition, the cycle integral is the linear pairing \(\oint_{\gamma}a=\langle a,\gamma\rangle=\sum_{e}\gamma(e)\,a(e)\) of the cochain \(a\) with the coefficients \(\gamma(e)\in k\) of \(\gamma\in Z_1(G;k)\). Substituting \(a=d\phi\), we have \(a(e)=\phi(e^{+})-\phi(e^{-})\) for edge \(e\) with head \(e^{+}\) and tail \(e^{-}\), so

\[
\oint_{\gamma}d\phi\;=\;\sum_{e}\gamma(e)\bigl(\phi(e^{+})-\phi(e^{-})\bigr)\;=\;\sum_{v}\phi(v)\,(\partial\gamma)(v)\;=\;0,
\]

since \(\partial\gamma=0\) for \(\gamma\in Z_1(G;k)\).

The associated no-drive defect record is

\[
\begin{aligned}\delta_{6}^{\mathrm{exact}}=\bigl(&\{\oint_{\gamma}a:\gamma\in\mathsf{Cycles}(G)\},\Omega_{\mathrm{drive}},\Theta_{=0},\mathrm{polarity}_{\mathrm{no\text{-}drive}},\\&W_{\mathrm{exact}}=\phi,A_{\mathrm{graph}},V_{\mathrm{graph}},\sigma_{\mathrm{no\text{-}drive}}\bigr).\end{aligned}
\]

The instrument-indexed verdict is \(\Gamma;\mathsf H_{\mathrm{graph}};I_{\mathrm{graph}}\vdash\textsc{ExactFormNullDrive}:\texttt{accepted}\). \(\square\)

\subsubsection{Nonzero-Affinity Equivalence (Theorem 17)}\label{sn:B.2.3}

\emph{Statement (recall).} For \(a\in C^{1}(G;k)\), the cohomology class \([a]\neq 0\) in \(H^{1}(G;k)\) iff there exists a cycle \(\gamma\in\mathsf{Cycles}(G)\) with \(\oint_{\gamma}a\neq 0\).

\emph{Proof.} Forward direction, by contraposition. Suppose \(\oint_{\gamma}a=0\) for every cycle \(\gamma\). Fix a spanning forest \(F\) of \(G\) with a root in each component, and let \(\phi(v)\) be the signed sum of \(a\) along the \(F\)-path from the root to \(v\). Then \(a(e)=\phi(\mathrm{head}(e))-\phi(\mathrm{tail}(e))\) for every \(e\in F\) by construction, and for every \(e\notin F\) because the fundamental cycle of \(e\) with respect to \(F\) has zero integral. So \(a=d^{0}\phi\) and \([a]=0\). The argument uses only the additive group of \(k\) (Subsection \hyperref[sn:8.4]{8.4} gives the details).

Reverse direction. Suppose there exists \(\gamma\in\mathsf{Cycles}(G)\) with \(\oint_{\gamma}a\neq 0\). By the exact-form null-drive theorem (Subsubsection \hyperref[sn:B.2.2]{B.2.2}), an exact \(1\)-cochain has zero cycle integral on every cycle. Hence \(a\) is not exact, so \([a]\neq 0\) in \(H^{1}(G;k)\).

The associated drive certificate record is

\[
\kappa_{6}^{\mathrm{drive}}(a)\;=\;(\,[a]\neq 0,\;\Omega_{\mathrm{drive}},\;\Theta_{\neq0},\;\mathrm{polarity}_{\mathrm{drive}},\;W=\gamma,\;A_{\mathrm{graph}},\;V_{\mathrm{graph}},\;\sigma_{\mathrm{drive}}\,).
\]

The instrument-indexed verdict is \(\Gamma;\mathsf H_{\mathrm{graph}};I_{\mathrm{graph}}\vdash\textsc{NonzeroAffinityEquivalence}:\texttt{accepted}\). \(\square\)

\subsubsection{Gating-Affinity Suppression (Theorem 18)}\label{sn:B.2.4}

\emph{Statement (recall).} Let \(G\rightsquigarrow G'=(V',E')\) be a \(P_{2}\)-gating that sends \(G\) to a forest, that is, \(\beta_{1}(G')=0\). Then \(H^{1}(G';k)=0\), and no positive cycle-supported cohomological drive certificate \(\mathrm{P6}_{\mathrm{drive}}^{H^1}\) exists on \(G'\); for every \(a\in C^{1}(G';k)\) every cycle integral vanishes, which is a no-drive certificate once the source and audit records required by the surrounding instrument are attached.

\emph{Proof.} By the forest no-drive theorem (Subsubsection \hyperref[sn:B.2.1]{B.2.1}), \(\beta_{1}(G')=0\) implies \(H^{1}(G';k)=0\). Hence every cohomology class on \(G'\) is zero; in particular, the restriction \([a|_{G'}]\) is the zero class in \(H^{1}(G';k)\), so the drive certificate \(\kappa_{6}^{\mathrm{drive}}(a)\) does not survive the gating \(G\rightsquigarrow G'\). Since \(\beta_{1}(G')=0\), the graph \(G'\) has no nonzero cycle, so every cycle integral on \(G'\) vanishes for every \(a\in C^{1}(G';k)\); with the source and audit records required by the surrounding instrument, this is the no-drive certificate.

The cycle-rank loss defect \(\delta_{2}^{\mathrm{cycle}}(G,G')=\beta_{1}(G)-\beta_{1}(G')\) records the number of cycles destroyed by the gating; the forest threshold \(\Theta_{2}^{\mathrm{forest}}:\beta_{1}(G')=0\) records the gating-to-forest condition.

The instrument-indexed verdict is \(\Gamma;\mathsf H_{\mathrm{graph}};I_{\mathrm{graph}}\vdash\textsc{GatingAffinitySuppression}:\texttt{accepted}\). \(\square\)

The four graph/cohomology theorems together record the no-drive boundary of the graph/cohomology realization of Subsection \hyperref[sn:13.4]{13.4}: forest hosts have no cycle-supported cohomological drive, exact \(1\)-cochains have no cycle drive, nonzero cohomology is equivalent to a nonzero cycle integral, and gating to a forest kills cycle-supported drive.

\subsection{Completion and Pasting Proofs}\label{sn:B.3}

The proofs of Theorems \hyperref[cl:thm:9]{9} and \hyperref[cl:thm:31]{31} are complete in Subsubsection \hyperref[sn:7.4.1]{7.4.1} and Subsection \hyperref[sn:14.3]{14.3} and are not repeated here.

\subsection{Definability Bound Proof}\label{sn:B.4}

This subsection records the long-form proof of the fixed-interface definability bound (Theorem \hyperref[cl:thm:14]{14}) of Section \hyperref[sn:7]{7}. The main-text proof is in Subsection \hyperref[sn:7.9]{7.9}; this subsection records the explicit counting argument and the definability set on a fixed visible interface.

\emph{Statement (recall).} Let \(f:X\to Y\) be a finite map on \(\mathsf H_{\mathrm{fin}}\). The full finite-lens definability of \(f\) is

\[
\mathrm{Definable}(f)\;=\;\{\,f^{-1}(B):B\subseteq\mathrm{im}(f)\,\}.
\]

Then \(\mathrm{Definable}(f)\) has cardinality

\[
|\mathrm{Definable}(f)|\;=\;2^{|\mathrm{im}(f)|}.
\]

\emph{Proof.} The set \(\mathrm{Definable}(f)\) is in bijection with the powerset of \(\mathrm{im}(f)\). Define

\[
\Phi:\mathcal P(\mathrm{im}(f))\longrightarrow\mathrm{Definable}(f),
\qquad
\Phi(B)=f^{-1}(B).
\]

The map is well-defined because each \(B\subseteq\mathrm{im}(f)\) gives a subset \(f^{-1}(B)\subseteq X\), and by definition this subset lies in \(\mathrm{Definable}(f)\).

\emph{Injectivity.} Suppose \(B,B'\subseteq\mathrm{im}(f)\) and \(B\neq B'\). Choose \(y\in B\triangle B'\). Without loss of generality, let \(y\in B\setminus B'\). Since \(y\in\mathrm{im}(f)\), there exists \(x\in X\) with \(f(x)=y\). Then \(x\in f^{-1}(B)\) but \(x\notin f^{-1}(B')\), so \(f^{-1}(B)\neq f^{-1}(B')\). Hence \(\Phi(B)\neq\Phi(B')\).

\emph{Surjectivity.} By definition, every element of \(\mathrm{Definable}(f)\) is of the form \(f^{-1}(B)\) for some \(B\subseteq\mathrm{im}(f)\). Hence every definable predicate lies in the image of \(\Phi\).

The map \(\Phi\) is therefore a bijection, and

\[
|\mathrm{Definable}(f)|
\;=\;
|\mathcal P(\mathrm{im}(f))|
\;=\;
2^{|\mathrm{im}(f)|}.
\]

The fixed-interface definability bound records that the number of package distinctions definable through the fixed lens \(f\) depends only on the finite interface \(\mathrm{im}(f)\). It does not count open-ended definability after changing the interface, the package, the instrument, or the host.

The instrument-indexed verdict is

\[
\Gamma\,;\;\mathsf H_{\mathrm{fin}}\,;\;I_{\mathrm{fin}}\;\vdash\;\textsc{FixedInterfaceDefinabilityBound}\;:\;\texttt{accepted}.\quad\square
\]

The bound is finite and decidable on the finite records. It supports the fixed-interface discipline used in Section \hyperref[sn:10]{10}: strict promotion is recorded as an interface change rather than as theory growth at the same fixed interface.
 \section{Full Countermodel Atlas}\label{sn:C}

This appendix records the full twelve-countermodel atlas. Sheets \(\mathsf{CM}_1\) through \(\mathsf{CM}_6\) summarize the constructions of Section \hyperref[sn:9]{9} with their dependency judgments and point back to the corresponding subsections of Section \hyperref[sn:9]{9}. Sheets \(\mathsf{CM}_7\) through \(\mathsf{CM}_{12}\) supply the full constructions for the appendix-only countermodels referenced in Subsection \hyperref[sn:9.8]{9.8}.

\subsection{Atlas Convention}\label{sn:C.1}

Each sheet uses the form fixed in Subsubsection \hyperref[sn:9.1.1]{9.1.1}: construction, supported no-go (with dependency judgment), what it blocks, what it does not block, and required positive bridge or repair where applicable. The predicates \(S\) and \(X\) of every displayed judgment are read by the convention \emph{Reading the predicates} in the \emph{Supported no-go} item of Subsubsection \hyperref[sn:9.1.1]{9.1.1}, and the relation types by the truth conditions of Subsubsection \hyperref[sn:4.6.1]{4.6.1}. The tuple form \((\mathrm{id},H,\mathsf{Data},\mathsf{ActiveClaim},\mathsf{FailedClaim},W,A,V,\mathcal N)\) of Subsubsection \hyperref[sn:9.1.1]{9.1.1} is the underlying representation of each sheet; the prose fields below unpack that representation for the reader.

Where a sheet repeats a Section \hyperref[sn:9]{9} entry, the sheet is brief and points back to the corresponding main-text subsection. Where the sheet is appendix-only, the construction is given in full with the data of the calculus made explicit.

\subsection{\texorpdfstring{\(\mathsf{CM}_1\): Active Package but No Macro Closure}{C.2 \textbackslash mathsf\{CM\}\_1: Active Package but No Macro Closure}}\label{sn:C.2}

Construction, blocked overclaim, and required bridge are recorded in Subsection \hyperref[sn:9.2]{9.2}: a finite package \(q:\{a,b,c\}\to\{A,B\}\) with \(q(a)=q(b)=A\), \(q(c)=B\), together with \(F:X\to X\) given by \(F(a)=a\), \(F(b)=c\), \(F(c)=c\). The split-pair \((a,b)\) witnesses \(\delta_1^{\mathrm{split}}(q,F)\neq\varnothing\), so no descended update through \(q\) exists.

The dependency judgment is

\[
\Gamma\,;\;\mathsf H_{\mathrm{fin}}\,;\;I_{\mathrm{fin}}\;\vdash\;P_5(q)\;\mathrel{\mathsf{dep}_{\texttt{insufficient}}}\;P_1(qF=F^\sharp q),
\]

supporting the no-go \(P_5\not\Rightarrow P_1\text{ macro closure}\) (NC-7). The required positive bridge is the descent record \(\delta_1^{\mathrm{split}}(q,F)=\varnothing\) together with the audit and visibility data of the surrounding judgment, with Theorem \hyperref[cl:thm:7]{7} supplying the equivalence between the empty descent defect and the existence of \(F^\sharp\).

\subsection{\texorpdfstring{\(\mathsf{CM}_2\): Active Route / Holonomy but No Drive}{C.3 \textbackslash mathsf\{CM\}\_2: Active Route / Holonomy but No Drive}}\label{sn:C.3}

Construction and blocked overclaim are recorded in Subsection \hyperref[sn:9.3]{9.3}: a finite cycle graph \(G\) on \(\mathsf H_{\mathrm{graph}}\) (for instance the triangle \(0\to 1\to 2\to 0\)) with active \(P_3\)-route data, the transport record with fiber \(\{u,v\}\) whose only nonidentity edge permutation is the swap on \(2\to0\), so that \(\operatorname{Hol}(\gamma)\) is the swap, and the log-ratio cochain \(a(x\to y)=\log\bigl(M(x,y)/M(y,x)\bigr)\) of the symmetric random walk \(M(x,y)=1/2\) on \(G\), which is \(a=0\). Every cycle integral vanishes, so \([a]=0\in H^1(G)\) and no positive cycle-supported cohomological drive certificate is admissible on \(G\) for this \(a\).

The dependency judgment is

\[
\Gamma\,;\;\mathsf H_{\mathrm{graph}}\,;\;I_{\mathrm{graph}}\;\vdash\;P_3^{\mathrm{holonomy}}\;\mathrel{\mathsf{dep}_{\texttt{requires\_bridge}}}\;\mathrm{P6}_{\mathrm{drive}}^{H^1},
\]

supporting NC-12 (no route-mismatch-implies-drive). The required positive bridge is the drive bridge \(B_{3\to 6}^{\mathrm{drive}}\) in \(\mathsf{BridgePolicy}_I\) that connects the route/holonomy data to a specific cochain \(a\) with \([a]\neq 0\).

\subsection{\texorpdfstring{\(\mathsf{CM}_3\): Local Packages Fail Globally}{C.4 \textbackslash mathsf\{CM\}\_3: Local Packages Fail Globally}}\label{sn:C.4}

Construction and blocked overclaim are recorded in Subsection \hyperref[sn:9.4]{9.4}: a finite carrier \(Z=\{a,b,c\}\) with cover \(U=\{a,b\}\), \(V=\{b,c\}\) and overlap \(U\cap V=\{b\}\), with all package targets equal to \(\{0,1\}\), identity restriction maps, and local packages \(q_U,q_V\) that assign incompatible values \(0,1\) on the overlap \(\{b\}\). The gluing defect \(\delta_4^{\mathrm{glue}}\) is non-empty, with \(b\) as a witness, and no global package extends both local data.

The dependency judgment is

\[
\Gamma\,;\;\mathsf H_{\mathrm{fin}}\,;\;I_{\mathrm{fin}}\;\vdash\;\bigl\{P_5(q_U),\,P_5(q_V)\bigr\}\;\mathrel{\mathsf{dep}_{\texttt{insufficient}}}\;P_5(q_{\mathrm{global}}),
\]

supporting NC-14 (no local-implies-global). The required positive bridge is a \(P_4\)-gluing bridge \(B_{\mathrm{glue}}\) with \(\delta_4^{\mathrm{glue}}=\varnothing\) together with a \(P_6\)-audit compatibility record on the overlaps.

\subsection{\texorpdfstring{\(\mathsf{CM}_4\): Directed Cell Active but Pair Observable Blocked}{C.5 \textbackslash mathsf\{CM\}\_4: Directed Cell Active but Pair Observable Blocked}}\label{sn:C.5}

Construction and blocked overclaim are recorded in Subsection \hyperref[sn:9.5]{9.5}: a directed cell \(\mathsf{Cell}_{63}^{\lambda_{\mathrm{audit}}}\) with \(W_3=\mathsf{RouteMismatchWit}\), \(U_6=\mathsf{AddAuditRecord}\) and audit-only profile \(\lambda_{\mathrm{audit}}\), classified \(\texttt{action}\) by the cell classifier of Subsubsection \hyperref[sn:5.6.2]{5.6.2}; together with a pair-observable record \(\mathsf{PairObs}_{\{3,6\}}^{\mu_{\mathrm{drive}}}\) on the same primitives with profile \(\mu_{\mathrm{drive}}\) asserting drive content, but with no admissible drive bridge and no admissible drive source-of-truth, classified \(\texttt{blocked}\) by the pair classifier of Subsubsection \hyperref[sn:5.6.3]{5.6.3}.

The dependency judgment is

\[
\Gamma\,;\;\mathsf H_{\mathrm{fin}}\,;\;I_{\mathrm{fin}}\;\vdash\;\mathsf{Cell}_{63}^{\lambda_{\mathrm{audit}}}:\texttt{action}\;\mathrel{\mathsf{dep}_{\texttt{insufficient}}}\;\mathsf{PairObs}_{\{3,6\}}^{\mu_{\mathrm{drive}}}:\texttt{real},
\]

supporting NC-5 (no cell-to-pair collapse). The required positive bridge is a real or fallback source under \(\mathsf{SourcePolicy}_I\) together with the drive bridge \(B_{3\to 6}^{\mathrm{drive}}\) when the pair claims drive content; Theorem \hyperref[cl:thm:5]{5} of Section \hyperref[sn:6]{6} records the status-family separation in theorem-grade form.

\subsection{\texorpdfstring{\(\mathsf{CM}_5\): Refinement Worsens Closure}{C.6 \textbackslash mathsf\{CM\}\_5: Refinement Worsens Closure}}\label{sn:C.6}

Construction and blocked overclaim are recorded in Subsection \hyperref[sn:9.6]{9.6}: a finite carrier \(X=\{a,b,c,d\}\) with coarse abstraction \(q_0:X\to\{*\}\) and refinement \(q_1:X\to\{A,B,C\}\) given by \(q_1(a)=q_1(b)=A\), \(q_1(c)=B\), \(q_1(d)=C\), with \(q_0\) factoring through \(q_1\) via the forgetting map \(r:\{A,B,C\}\to\{*\}\). The finite update \(F(a)=c\), \(F(b)=d\), \(F(c)=c\), \(F(d)=d\) descends through \(q_0\) trivially but creates a \(q_1\)-split at \((a,b)\), so \(\delta_1^{\mathrm{split}}(q_1,F)\neq\varnothing\).

The dependency judgment is

\[
\Gamma\,;\;\mathsf H_{\mathrm{fin}}\,;\;I_{\mathrm{fin}}\;\vdash\;P_4(q_0\leadsto q_1)\;\mathrel{\mathsf{dep}_{\texttt{insufficient}}}\;\bigl(\delta_1^{\mathrm{split}}(q_1,F)\;\subseteq\;\delta_1^{\mathrm{split}}(q_0,F)\bigr),
\]

supporting NC-13 (no refinement-implies-improvement). A positive refinement-improvement theorem requires a \(P_4\)-monotonicity bridge with explicit hypotheses on the refinement, the update, and the closure functional, not asserted in this paper (Subsubsection \hyperref[sn:9.6.5]{9.6.5}).

\subsection{\texorpdfstring{\(\mathsf{CM}_6\): Same-Level Self-Audit Circularity}{C.7 \textbackslash mathsf\{CM\}\_6: Same-Level Self-Audit Circularity}}\label{sn:C.7}

Construction and blocked overclaim are recorded in Subsection \hyperref[sn:9.7]{9.7}: a finite admissible instrument \(I_0\) under package \(\mathcal T_0\) asked to certify \(\mathrm{Sound}(I_0)\) at the same level; the instance contains the audit record \(A_0\) of Subsubsection \hyperref[sn:9.7.1]{9.7.1}. Two cases arise --- outside-scope when the claim's host tag is not in \(\mathsf{Hosts}_{I_0}\), a level tag of its profile is not in \(\mathsf{Levels}_{I_0}\), \(\texttt{soundness}\notin\mathsf{ClaimTypes}(I_0)\) or \(\mathsf{Use}(\mathrm{Sound}(I_0))\not\subseteq\mathsf{Scope}_{I_0}\), and undefined-circular (or a status of higher priority) when none of these four conditions holds and the use-set of \(\mathrm{Sound}(I_0)\) contains the verdict reference \(\ulcorner\mathrm{Sound}(I_0),I_0\urcorner\) --- and in neither case does the claim classifier accept \(\mathrm{Sound}(I_0)\) under \(I_0\) at the same level.

The dependency judgment is

\[
\Gamma\,;\;\mathsf H_{\mathrm{instr}}\,;\;I_{\mathrm{instr}}\;\vdash\;I_0\text{ audits }\mathcal T_0\;\mathrel{\mathsf{dep}_{\texttt{blocks}}}\;\bigl(I_0\vdash\mathrm{Sound}(I_0):\texttt{accepted}\bigr),
\]

supporting NC-16 (no same-level complete self-certification). The required positive repair is a level shift via a higher instrument \(I_1\) that audits \(I_0\) through an admissible audit bridge \(B^{\mathrm{audit}}_{0\to 1}\); Theorem \hyperref[cl:thm:24]{24} of Section \hyperref[sn:11]{11} states the classifier routing of the same-level self-soundness claim, and Theorem \hyperref[cl:thm:25]{25} records the rotating-audit theorem with NC-17 (no rotating-audit finality).

\subsection{\texorpdfstring{\(\mathsf{CM}_7\): Same Primitive Pair, Different Statuses}{C.8 \textbackslash mathsf\{CM\}\_7: Same Primitive Pair, Different Statuses}}\label{sn:C.8}

\subsubsection{Construction}\label{sn:C.8.1}

Take the primitive pair \(P_6\leftarrow P_3\) in \(\mathbf{BirdInt}^{\mathrm{aud}}_{\mathrm{fin}}\).

\textbf{Cell 1 --- audit-only profile \(\lambda_{\mathrm{audit}}\).} Take witness \(W_3=\mathsf{RouteMismatchWit}\in\mathsf{Wit}(P_3)\) and update \(U_6=\mathsf{AddAuditRecord}\in\mathsf{Upd}(P_6)\), with the witness visible and audited under the surrounding instrument \(I\) and a profile \(\lambda_{\mathrm{audit}}\) that records the audit-only instrument mode for an audit update. The cell

\[
\mathsf{Cell}_{63}^{\lambda_{\mathrm{audit}}}\;=\;(P_6\leftarrow P_3,\,W_3,\,U_6,\,L,\,V,\,\Theta_{\ge1},\,A,\,\delta=\varnothing)
\]

is well-formed, and the classifier of Subsubsection \hyperref[sn:5.6.2]{5.6.2} assigns

\[
\mathsf{Cell}_{63}^{\lambda_{\mathrm{audit}}}\;:\;\texttt{action}.
\]

\textbf{Cell 2 --- drive-certification profile.} Use the same primitive pair \(P_6\leftarrow P_3\), but with update \(U_6'=\mathsf{CertifyDrive}\in\mathsf{Upd}(P_6)\) and a profile \(\lambda_{\mathrm{drive}}\) that requires a drive bridge. Suppose the only member of \(\mathcal L\) bridging the drive datum is a registered drive bridge \(L_{\mathrm{drive}}\) whose strength is below that required by \(\lambda_{\mathrm{drive}}\), so that \(\delta_{\mathrm{weakbridge}}\) is nonempty (the record is written out in Subsubsection \hyperref[sn:6.1.1]{6.1.1}). The cell record

\[
\mathsf{Cell}_{63}^{\lambda_{\mathrm{drive}}}
\]

fails the bridge check; its classifier verdict is

\[
\mathsf{Cell}_{63}^{\lambda_{\mathrm{drive}}}\;:\;\texttt{blocked}.
\]

The same primitive pair \((P_6,P_3)\) thus carries two cell records with different profile, update, and bridge requirements, and these records receive distinct cell statuses.

\subsubsection{Supported No-Go}\label{sn:C.8.2}

The construction supports the dependency judgment

\[
\Gamma\,;\;\mathsf H_{\mathrm{fin}}\,;\;I_{\mathrm{fin}}\;\vdash\;\mathsf{Cell}_{63}^{\lambda_{\mathrm{audit}}}:\texttt{action}\;\mathrel{\mathsf{dep}_{\texttt{insufficient}}}\;\mathsf{Cell}_{63}^{\lambda_{\mathrm{drive}}}:\texttt{action},
\]

witnessed by the instance of Theorem \hyperref[cl:thm:1]{1}; hence cell status does not factor through the primitive pair: there is no map \(g:\mathbb P\times\mathbb P\to\mathrm{CellStatus}\) recovering both verdicts from the same pair. In particular, since \(d\circ *\) would be such a \(g\), no total operation \(*:\mathbb P\times\mathbb P\to\mathbb P\) together with a decoder \(d:\mathbb P\to\mathrm{CellStatus}\) can faithfully encode typed, statused, profile-relative interaction. The construction is the working content of NC-1 and NC-3 of Subsection \hyperref[sn:15.2]{15.2}, and Theorem \hyperref[cl:thm:1]{1} of Section \hyperref[sn:6]{6} makes the non-factorization formal.

\subsubsection{What It Blocks}\label{sn:C.8.3}

The overclaim ``the \(6\times 6\) primitive-pair table is a universal interaction algebra'' is rejected. NC-1 and NC-3 are the corresponding nonclaims.

\subsubsection{What It Does Not Block}\label{sn:C.8.4}

The construction does not say cell records on \((P_6,P_3)\) are unattainable. Both are finite covered records for the directed-cell classifier; what fails is the inference from the primitive pair alone to a cell status.

\subsubsection{Required Positive Bridge}\label{sn:C.8.5}

A positive cell-status claim on \((P_6,P_3)\) requires the surrounding profile, witness, update, threshold, audit, visibility, and bridge data of the cell record; the primitive pair alone is not sufficient input.

\subsection{\texorpdfstring{\(\mathsf{CM}_8\): Idempotent Completions Do Not Commute}{C.9 \textbackslash mathsf\{CM\}\_8: Idempotent Completions Do Not Commute}}\label{sn:C.9}

\subsubsection{Construction}\label{sn:C.9.1}

Take a finite three-element set \(U=\{a,b,c\}\) and let \(C=\mathcal P(U)\) be its finite power set on \(\mathsf H_{\mathrm{fin}}\). Define two finite endomaps \(E_1,E_2:C\to C\) by

\[
E_1(S)\;=\;\begin{cases}\,S\cup\{b\}, & a\in S,\\[2pt] S, & a\notin S,\end{cases}\qquad E_2(S)\;=\;\begin{cases}\,S\cup\{c\}, & b\in S,\\[2pt] S, & b\notin S.\end{cases}
\]

\textbf{Idempotence.} For \(E_1\): if \(a\in S\), then \(E_1(S)=S\cup\{b\}\) still contains \(a\), so \(E_1(E_1(S))=(S\cup\{b\})\cup\{b\}=S\cup\{b\}=E_1(S)\). If \(a\notin S\), then \(E_1(S)=S\), hence \(E_1(E_1(S))=E_1(S)\). So \(E_1^2=E_1\). The argument for \(E_2\) is symmetric: \(E_2^2=E_2\).

\textbf{Noncommutation.} Take \(S_0=\{a\}\). Then

\[
E_1\,E_2(\{a\})\;=\;E_1(\{a\})\;=\;\{a,b\},
\]

since \(E_2\) leaves \(\{a\}\) unchanged (\(b\notin\{a\}\)). On the other hand,

\[
E_2\,E_1(\{a\})\;=\;E_2(\{a,b\})\;=\;\{a,b,c\},
\]

since \(E_1\) adds \(b\) (because \(a\in\{a\}\)) and then \(E_2\) adds \(c\) (because \(b\in\{a,b\}\)). Hence \(E_1E_2(\{a\})\neq E_2E_1(\{a\})\), and \(\Delta_3^{\mathrm{comp}}(E_1,E_2)\neq\varnothing\) by Theorem \hyperref[cl:thm:10]{10}.

\subsubsection{Supported No-Go}\label{sn:C.9.2}

The construction supports the dependency judgment

\[
\Gamma\,;\;\mathsf H_{\mathrm{fin}}\,;\;I_{\mathrm{fin}}\;\vdash\;\bigl(E_1^2=E_1\,\wedge\,E_2^2=E_2\bigr)\;\mathrel{\mathsf{dep}_{\texttt{insufficient}}}\;\bigl(E_1E_2=E_2E_1\bigr),
\]

recording that individual idempotence does not imply commutation, equivalently that \(P_5\)-individual exactness does not imply \(P_3\)-route coherence.

\subsubsection{What It Blocks}\label{sn:C.9.3}

The overclaim ``valid completions always commute'' is rejected. NC-8 of Subsection \hyperref[sn:15.2]{15.2} records the corresponding nonclaim, and Theorem \hyperref[cl:thm:9]{9} of Section \hyperref[sn:7]{7} makes the noncommuting-completions result theorem-grade.

\subsubsection{What It Does Not Block}\label{sn:C.9.4}

The construction does not say completions are always non-commuting, or that route coherence is unattainable. Specific pairs of idempotent completions on a finite carrier may commute --- for instance, when the two maps are identical or when one of them is the identity --- but commutation is not implied by individual idempotence.

\subsubsection{Required Positive Bridge}\label{sn:C.9.5}

A positive commutation claim \(E_1E_2=E_2E_1\) requires an explicit route-coherence record: the empty defect \(\Delta_3^{\mathrm{comp}}(E_1,E_2)=\varnothing\) together with the audit and visibility data of the surrounding judgment. Theorem \hyperref[cl:thm:10]{10} makes the equivalence between the empty defect and commutation precise.

\subsection{\texorpdfstring{\(\mathsf{CM}_9\): Strict Extension Without Macro Closure}{C.10 \textbackslash mathsf\{CM\}\_9: Strict Extension Without Macro Closure}}\label{sn:C.10}

\subsubsection{Construction}\label{sn:C.10.1}

Take a finite carrier \(S=\{a,b,c\}\) on \(\mathsf H_{\mathrm{fin}}\). Define an old object map \(\pi_0:S\to O_0=\{*\}\) collapsing every element to a single value:

\[
\pi_0(a)\;=\;\pi_0(b)\;=\;\pi_0(c)\;=\;*.
\]

Define a new object map \(\pi_1:S\to O_1=\{U,V\}\) with

\[
\pi_1(a)\;=\;U,\qquad \pi_1(b)\;=\;U,\qquad \pi_1(c)\;=\;V.
\]

Since \(\pi_0\) collapses everything but \(\pi_1\) separates \(c\), there exist \(s,s'\in S\) --- for example \(s=a\) and \(s'=c\) --- with \(\pi_0(s)=\pi_0(s')=*\) but \(\pi_1(s)=U\neq V=\pi_1(s')\). By Theorem \hyperref[cl:thm:12]{12}, \(\pi_1\not\factor\pi_0\), so the bridge is strict at the object-map level: the strictness gate \(G_{\mathrm{strict}}\) records \(\texttt{strict\_pass}\) from \(\Delta_{\mathrm{fact}}(\pi_0,\pi_1)\neq\varnothing\).

Now define a finite update \(F:S\to S\) by

\[
F(a)\;=\;a,\qquad F(b)\;=\;c,\qquad F(c)\;=\;c.
\]

Then \(\pi_1(a)=\pi_1(b)=U\), but \(\pi_1F(a)=\pi_1(a)=U\) and \(\pi_1F(b)=\pi_1(c)=V\), so \(\pi_1F(a)\neq\pi_1F(b)\). The descent defect \(\delta_1^{\mathrm{split}}(\pi_1,F)\) is non-empty, with \((a,b)\) as a witness, and by Theorem \hyperref[cl:thm:7]{7} no descended dynamics exists for \(F\) through \(\pi_1\). The object map is strict, but macro closure on the target package fails.

\subsubsection{Supported No-Go}\label{sn:C.10.2}

The construction supports the dependency judgment

\[
\Gamma\,;\;\mathsf H_{\mathrm{fin}}\,;\;I_{\mathrm{prom}}\;\vdash\;\Delta_{\mathrm{fact}}(\pi_0,\pi_1)\neq\varnothing\;\mathrel{\mathsf{dep}_{\texttt{insufficient}}}\;P_1(\pi_1F=F^\sharp\pi_1),
\]

recording that a strict object map does not by itself imply macro closure on the target. The bridge carrying these object maps is completed as in the instance of Theorem \hyperref[cl:thm:22]{22} (Subsubsection \hyperref[sn:10.6.1]{10.6.1}), with \(S=\{a,b,c\}\), \(O_0=\{*\}\) and \(O_1=\{U,V\}\): packages \(\mathcal T_k=(S,\pi_k,\Sigma_{\pi_k},\mathrm{id},\mathcal A_k)\), the instrument \(I_{\mathrm{prom}}\), every field visible, and \(G_{\mathrm{desc}}=\texttt{not\_required}\) because the bridge makes no closure claim. The gates pass as in Subsubsection \hyperref[sn:10.6.2]{10.6.2}, so the promotion classifier assigns \(\texttt{strict}\), while \((a,b)\in\delta_1^{\mathrm{split}}(\pi_1,F)\) for the update \(F\) recorded in the target instance.

\subsubsection{What It Blocks}\label{sn:C.10.3}

The overclaim ``strict extension automatically gives closed dynamics'' is rejected. NC-10 of Subsection \hyperref[sn:15.2]{15.2} records the corresponding nonclaim.

\subsubsection{What It Does Not Block}\label{sn:C.10.4}

The construction does not say a strict object map is incompatible with macro closure. A different finite update \(F'\) --- one whose action is consistent with the fibers of \(\pi_1\) --- would descend through \(\pi_1\) and produce an admissible descended dynamics on the target package; the construction shows only that descent is independent data from strictness.

\subsubsection{Required Positive Bridge}\label{sn:C.10.5}

A positive macro-closure claim on a promotion bridge with a strict object map requires an explicit descent record: the empty defect \(\delta_1^{\mathrm{split}}(\pi_1,F)=\varnothing\) recorded by the descent gate \(G_{\mathrm{desc}}\) of Subsubsection \hyperref[sn:5.4.4]{5.4.4}, together with the other audit, visibility, and required gate data of the bridge.

\subsection{\texorpdfstring{\(\mathsf{CM}_{10}\): Strict Extension Without Drive}{C.11 \textbackslash mathsf\{CM\}\_\{10\}: Strict Extension Without Drive}}\label{sn:C.11}

\subsubsection{Construction}\label{sn:C.11.1}

Use the same strict object maps from \(\mathsf{CM}_9\): \(S=\{a,b,c\}\), \(\pi_0:S\to\{*\}\) collapsing all elements, and \(\pi_1:S\to\{U,V\}\) separating \(c\). By Theorem \hyperref[cl:thm:12]{12}, \(\pi_1\not\factor\pi_0\), so the object map is strict.

Record in the bridge's level/lens/interface record \(L\) a finite support graph \(G_B\) that is a forest, for example the path

\[
a\;\text{---}\;b\;\text{---}\;c
\]

with vertex set \(V(G)=S=\{a,b,c\}\), the support of the bridge's carrier, and edge set \(E(G)=\{(a,b),(b,c)\}\). Then \(|E(G)|=2\), \(|V(G)|=3\) and \(G\) is connected, so \(\beta_1(G)=2-3+1=0\). By Theorem \hyperref[cl:thm:15]{15},

\[
H^1(G)\;=\;0,
\]

and every \(1\)-cochain \(\omega\) on \(G\) has \([\omega]=0\in H^1(G)\), so no admissible cycle-supported cohomological drive certificate \(\mathrm{P6}_{\mathrm{drive}}^{H^1}\) is available on \(G\). The object map is strict, but the drive face is absent on the support graph.

\subsubsection{Supported No-Go}\label{sn:C.11.2}

The construction supports the dependency judgment

\[
\Gamma\,;\;\mathsf H_{\mathrm{graph}}\,;\;I_{\mathrm{graph}}\;\vdash\;\Delta_{\mathrm{fact}}(\pi_0,\pi_1)\neq\varnothing\;\mathrel{\mathsf{dep}_{\texttt{insufficient}}}\;\mathrm{P6}_{\mathrm{drive}}^{H^1},
\]

recording that a strict object map does not by itself imply a cycle-supported cohomological drive certificate. In this instance, form a new bridge \(B^{\mathrm{graph}}_{0\to1}\) from the same object-map tables as \(\mathsf{CM}_9\), with both packages, the path support \(a\)--\(b\)--\(c\), and the bridge tagged by \(\mathsf H_{\mathrm{graph}}\). Tag both packages and the bridge with host \(\mathsf H_{\mathrm{graph}}\), give the bridge the profile \(\lambda_{\mathrm{prom}}\) of Subsubsection \hyperref[sn:10.6.1]{10.6.1}, and give \(I_{\mathrm{graph}}\), besides \(\texttt{promotion}\in\mathsf{ClaimTypes}_{I_{\mathrm{graph}}}\) and the finite source, visibility, audit and gate checks for these records, \(\mathsf{Hosts}\ni\mathsf H_{\mathrm{graph}}\), \(\mathsf{Levels}\ni\mathsf{Ver}\), \(\mathsf{ModePolicy}(\texttt{promotion})=\texttt{default}\), \((\texttt{fine},\texttt{structural},\texttt{local},\mathsf H_{\mathrm{graph}})\in\mathsf{Compat}_{I_{\mathrm{graph}}}\), a bridge policy admitting the visibility type, and source and evidence policies admitting \(\texttt{committed\_state}\) in the sense of Subsubsection \hyperref[sn:5.3.6]{5.3.6}, so that \(\mathsf{AdmProfile}(\lambda_{\mathrm{prom}},I_{\mathrm{graph}},\mathcal T_j)\) holds for \(j=0,1\). Recompute the gate results under \(I_{\mathrm{graph}}\); the required core gates pass, \(G_{\mathrm{desc}}=\texttt{not\_required}\), and the nonempty factorization defect gives \(G_{\mathrm{strict}}=\texttt{strict\_pass}\). The promotion classifier therefore assigns this bridge \(\texttt{strict}\).

\subsubsection{What It Blocks}\label{sn:C.11.3}

The overclaim ``strictness implies drive'' is rejected. NC-11 of Subsection \hyperref[sn:15.2]{15.2} records the corresponding nonclaim.

\subsubsection{What It Does Not Block}\label{sn:C.11.4}

The construction does not say a strict object map is incompatible with drive. A different support graph with \(\beta_1\ge 1\) --- for example, a triangle --- could carry a positive drive certificate; the construction shows only that the choice of support graph is independent data from the strictness of the object map.

\subsubsection{Required Positive Bridge}\label{sn:C.11.5}

A positive drive claim on a promotion bridge with a strict object map requires an explicit drive bridge in \(\mathsf{BridgePolicy}_I\) together with a support graph \(G'\) on which \(H^1(G')\neq 0\) and an admissible cochain \(\omega\) with \([\omega]\neq 0\in H^1(G')\) (written \(\omega\) here, since \(a\) names a vertex of this construction). Theorem \hyperref[cl:thm:17]{17} makes the equivalence between \([\omega]\neq 0\) and existence of a cycle with \(\oint_\gamma \omega\neq 0\) precise.

\subsection{\texorpdfstring{\(\mathsf{CM}_{11}\): Fallback Source Cannot Support Real Pair}{C.12 \textbackslash mathsf\{CM\}\_\{11\}: Fallback Source Cannot Support Real Pair}}\label{sn:C.12}

\subsubsection{Construction}\label{sn:C.12.1}

Take a pair-observable record \(Q=\mathsf{PairObs}_{\{i,j\}}^{\mu}\) on a primitive pair \(\{i,j\}\) on the host \(\mathsf H_{\mathrm{PICA}}\), under a local instrument \(I_{\mathrm{PICA}}\) (written \(I\) below) with \((t,\texttt{full},\mu)\in\mathsf{SourcePolicy}_I\), among the five candidate tags of Subsubsection \hyperref[sn:4.9.1]{4.9.1}, exactly for \(t\in\{\texttt{committed\_state},\texttt{independent\_pair\_witness}\}\), with or without entries \((\texttt{fallback},m,\mu)\).

Suppose the only available source record has type \(\texttt{fallback}\):

\[
\mathrm{src}\;=\;(\mathrm{id}_{\mathrm{src}},\;\texttt{fallback},\;\mathrm{trace}_{\mathrm{src}},\;\mathrm{valid}_{\mathrm{src}}),
\]

and that the record carries no audited source-upgrade bridge admitted by \(\mathsf{SourcePolicy}_I\) and \(\mathsf{BridgePolicy}_I\). Assume also that the source is within its validity condition, its trace records no failed audit or threshold check, and it is not obtained through a typed source bridge. Then \(\delta_6^{\mathrm{source}}\) is \(\texttt{fallback\_admitted}\) if \((\texttt{fallback},m,\mu)\in\mathsf{SourcePolicy}_I\) for some mode \(m\), and \(\texttt{fallback\_forbidden}\) otherwise (rules 5--6 of Subsubsection \hyperref[sn:5.3.6]{5.3.6}). In the first case condition (a) of the \(\texttt{real}\) row fails, so, when no condition of the first \(\texttt{blocked}\) row holds, \((\texttt{fallback},\texttt{provisional},\mu)\in\mathsf{SourcePolicy}_I\) and that row's other conditions hold, the classifier assigns

\[
Q\;:\;\texttt{real\_provisional},
\]

and otherwise \(\texttt{blocked}\); in the second the source threshold fails and the first \(\texttt{blocked}\) row applies. In neither case is \(Q\) \(\texttt{real}\).

\subsubsection{Supported No-Go}\label{sn:C.12.2}

The construction supports the dependency judgment

\[
\Gamma\,;\;\mathsf H_{\mathrm{PICA}}\,;\;I_{\mathrm{PICA}}\;\vdash\;\mathrm{src}:\texttt{fallback}\;\mathrel{\mathsf{dep}_{\texttt{insufficient}}}\;Q:\texttt{real},
\]

recording that fallback evidence does not by itself support a real pair-observable verdict.

\subsubsection{What It Blocks}\label{sn:C.12.3}

The overclaim ``placeholder or fallback source counts as a real pair source-of-truth'' is rejected. The pair-observable source-of-truth discipline of Subsubsection \hyperref[sn:4.4.3]{4.4.3} records the schema-level statement; this countermodel records the failure case at the dependency level.

\subsubsection{What It Does Not Block}\label{sn:C.12.4}

The construction does not say fallback sources are useless. A fallback source can support \(\texttt{real\_provisional}\), which is a legitimate pair-observable verdict, and, with an audited source-upgrade bridge admitted by the source and bridge policies, \(\texttt{real}\). The construction shows only that a fallback source does not by itself license the strongest verdict \(\texttt{real}\).

\subsubsection{Required Positive Bridge}\label{sn:C.12.5}

A positive \(Q:\texttt{real}\) verdict requires a real source under \(\mathsf{SourcePolicy}_I\) --- one of the five real-source classes from Subsubsection \hyperref[sn:4.9.1]{4.9.1}: \(\texttt{committed\_state}\), \(\texttt{audited\_cell\_records}\), \(\texttt{independent\_pair\_witness}\), \(\texttt{simulation\_trace}\), or \(\texttt{ablation\_record}\), admitted by \(\mu\). A fallback source can support \(\texttt{real}\) only when the record carries an audited source-upgrade bridge admitted by \(\mathsf{SourcePolicy}_I\) and \(\mathsf{BridgePolicy}_I\); absent that bridge, the fallback source supports at most \(\texttt{real\_provisional}\), or \(\texttt{blocked}\) if fallback is forbidden.

\subsection{\texorpdfstring{\(\mathsf{CM}_{12}\): Gating Destroys Cycle-Supported Drive}{C.13 \textbackslash mathsf\{CM\}\_\{12\}: Gating Destroys Cycle-Supported Drive}}\label{sn:C.13}

\subsubsection{Construction}\label{sn:C.13.1}

Take the triangle graph \(G\) on \(\mathsf H_{\mathrm{graph}}\) with vertex set \(V(G)=\{0,1,2\}\) and oriented edge set \(E(G)=\{(0,1),(1,2),(2,0)\}\). Then \(|E(G)|=3\), \(|V(G)|=3\), \(c(G)=1\), so

\[
\beta_1(G)\;=\;3\;-\;3\;+\;1\;=\;1,
\]

and the cycle space \(Z_1(G)\) is one-dimensional, spanned by the triangle cycle \(\gamma=(0,1)(1,2)(2,0)\).

Choose a finite \(1\)-cochain \(a\in C^1(G)\) with nonzero cycle integral, for example \(a((0,1))=1\), \(a((1,2))=0\), \(a((2,0))=0\) over a free coefficient host. Then

\[
\oint_\gamma a\;=\;a((0,1))\;+\;a((1,2))\;+\;a((2,0))\;=\;1\;\neq\;0,
\]

so \([a]\neq 0\in H^1(G)\) and \(\mathrm{P6}_{\mathrm{drive}}^{H^1}\) is active on \(G\) for this \(a\).

Now apply a \(P_2\)-style gate that deletes the edge \((2,0)\): the resulting graph \(G'\) has vertex set \(V(G')=\{0,1,2\}\) and edge set \(E(G')=\{(0,1),(1,2)\}\), that is, the path \(0\to 1\to 2\). Then \(|E(G')|=2\), \(|V(G')|=3\), \(c(G')=1\), so

\[
\beta_1(G')\;=\;2\;-\;3\;+\;1\;=\;0.
\]

By Theorem \hyperref[cl:thm:15]{15}, \(H^1(G')=0\), and by Theorem \hyperref[cl:thm:18]{18} no cohomological drive certificate \(\mathrm{P6}_{\mathrm{drive}}^{H^1}\) exists on \(G'\). The gate has destroyed the cycle-supported drive that was active on \(G\).

\subsubsection{Supported No-Go}\label{sn:C.13.2}

The construction supports the dependency judgment

\[
\Gamma\,;\;\mathsf H_{\mathrm{graph}}\,;\;I_{\mathrm{graph}}\;\vdash\;P_2(\text{edge deletion to forest})\;\mathrel{\mathsf{dep}_{\texttt{destroys}}}\;\mathrm{P6}_{\mathrm{drive}}^{H^1},
\]

recording that \(P_2\)-style gating that produces a forest destroys the cycle-supported cohomological drive certificate. The universal clause of \texttt{destroys} holds by Theorem \hyperref[cl:thm:18]{18}: every gate output \(G'\) with \(\beta_1(G')=0\) has \(H^1(G')=0\), so no recorded cochain has a nonzero class on \(G'\) (only cochains recorded on the gated graph count, Subsubsection \hyperref[sn:9.1.1]{9.1.1}).

\subsubsection{What It Blocks}\label{sn:C.13.3}

The overclaim ``gating only removes irrelevant support'' is rejected. The gating-suppression discipline of Section \hyperref[sn:8]{8} records the corresponding theorem (Theorem \hyperref[cl:thm:18]{18}), and this sheet gives the finite countermodel instance: gating to a forest can destroy a previously active drive certificate.

\subsubsection{What It Does Not Block}\label{sn:C.13.4}

The construction does not say all gates destroy drive. A \(P_2\)-style gate whose output retains nontrivial cycles (\(\beta_1(G')\ge 1\)) may leave drive certificates intact; the construction shows only that gates producing forests are sufficient to suppress cycle-supported cohomological drive.

\subsubsection{Required Positive Repair}\label{sn:C.13.5}

A positive drive-preservation claim under gating requires that the output graph \(G'\) have \(\beta_1(G')\ge 1\) and that the class \([a|_{G'}]\) of the restricted cochain be nonzero in \(H^1(G')\) --- equivalently, by Theorem \hyperref[cl:thm:17]{17}, that some cycle \(\gamma'\in Z_1(G')\) have \(\oint_{\gamma'}a\neq 0\). Without these conditions, the drive certificate is absent on \(G'\).
 \section{Model-Realization Sheets}\label{sn:D}

This appendix records the four scoped model-realization sheets in long form: PICA (D.1), Cantor strict-bridge (D.2), abstract interpretation (D.3), and graph/cohomology (D.4). Subsection \hyperref[sn:D.5]{D.5} records the combined coverage matrix across the four realizations. The sheets supplement the main-text records of Section \hyperref[sn:13]{13} and Section \hyperref[sn:12]{12} with per-realization reference data.

Each sheet declares the host, the realization object, the realization map, the realized fragment, the actor-update map and informant-witness map (where applicable), the per-judgment realizations, the source-of-truth and provenance discipline (where applicable), the realization theorem, and the realization-specific nonclaim register. Each sheet is host-bound and fragment-scoped: the main-text claim that the corresponding model realizes a declared fragment of \(\mathbf{BirdInt}^{\mathrm{aud}}_{\mathrm{fin}}\) is recorded in Section \hyperref[sn:13]{13} (or Section \hyperref[sn:12]{12} for AI), and this appendix does not introduce a new realization claim.

\subsection{PICA Sheet}\label{sn:D.1}

The PICA realization sheet \(\mathsf{PICAReal}_{\mathrm{fin}}\) records the realization of the finite stochastic cell/pair/provenance fragment of \(\mathbf{BirdInt}^{\mathrm{aud}}_{\mathrm{fin}}\) on the finite stochastic implementation host \(\mathsf H_{\mathrm{PICA}}\). The main-text record is in Subsection \hyperref[sn:13.2]{13.2}; this subsection records the long-form sheet.

\subsubsection{Host and Realization Object}\label{sn:D.1.1}

The host \(\mathsf H_{\mathrm{PICA}}\) has components

\[
\mathsf H_{\mathrm{PICA}}\;=\;(\,\Omega,\;K,\;\mathbb P,\;\mathsf{Bus},\;\mathsf{CellRules},\;\mathsf{WitnessStore},\;\mathsf{UpdateStore},\;\mathsf{PairStore},\;\mathsf{DepGraph},\;A,\;V\,),
\]

with \(\Omega\) a finite stochastic state space, \(K\) a finite Markov kernel on \(\Omega\), \(\mathbb P=\{P_{1},\ldots,P_{6}\}\), and the remaining fields finite stores. The realization object \(\mathcal P\) extends the host with \(\mathsf{StatusStore}\), \(\mathsf{Abl}\), and \(\mathcal N\):

\[
\begin{aligned}
\mathcal P\;=\;(\,&\Omega,\;K,\;\mathbb P,\;\mathsf{Bus},\;\mathsf{CellRules},\;\mathsf{WitnessStore},\;\mathsf{UpdateStore},\\
&\mathsf{StatusStore},\;\mathsf{PairStore},\;\mathsf{DepGraph},\;\mathsf{Abl},\;A,\;V,\;\mathcal N\,).
\end{aligned}
\]

The host instrument is \(I_{\mathrm{PICA}}\).

\subsubsection{Realization Map and Induced Package}\label{sn:D.1.2}

The realization map \(\mathcal R_{\mathrm{PICA}}:\mathsf{PICAReal}_{\mathrm{fin}}\to\mathbf{BirdInt}^{\mathrm{aud}}_{\mathrm{fin\text{-}stoch}}{}^{\mathrm{cell}/\mathrm{pair}/\mathrm{prov}}\) sends the realization object \(\mathcal P\) to the induced package

\[
\mathcal T_{\mathrm{PICA}}\;=\;(\,\Omega,\;f_{\mathrm{PICA}},\;\Sigma_{f_{\mathrm{PICA}}},\;E_{\mathrm{PICA}},\;\mathcal A_{\mathrm{PICA}}\,),
\]

with PICA-field interpretation: \(\Omega\) is the finite stochastic carrier, \(K\) provides finite Markov dynamics, the bus/lens projection provides \(f_{\mathrm{PICA}}\), \(\Sigma_{f_{\mathrm{PICA}}}=\{f_{\mathrm{PICA}}^{-1}(B):B\subseteq f_{\mathrm{PICA}}(\Omega)\}\) and \(E_{\mathrm{PICA}}=E_{f_{\mathrm{PICA}}}\) are the lens-induced content and completion, the visible PICA claims and the commit/update closure are attached records, and the finite audit/provenance ledger \(A\) assembles the inherited package-level audit functional \(\mathcal A_{\mathrm{PICA}}\). The audit record \(A\) remains distinct from the audit functional \(\mathcal A_{\mathrm{PICA}}\).

\subsubsection{Actor-Update and Informant-Witness Maps}\label{sn:D.1.3}

The actor-update map sends each primitive to its PICA update sort, and the informant-witness map sends each primitive to its PICA witness sort:

\begin{longtable}[]{@{}lll@{}}
\toprule
\begin{minipage}[b]{0.08\columnwidth}\raggedright
primitive\strut
\end{minipage} & \begin{minipage}[b]{0.46\columnwidth}\raggedright
actor update sort\strut
\end{minipage} & \begin{minipage}[b]{0.37\columnwidth}\raggedright
informant witness sort\strut
\end{minipage}\tabularnewline
\midrule
\endhead
\begin{minipage}[t]{0.08\columnwidth}\raggedright
\(P_{1}\)\strut
\end{minipage} & \begin{minipage}[t]{0.46\columnwidth}\raggedright
descent repair / closure update\strut
\end{minipage} & \begin{minipage}[t]{0.37\columnwidth}\raggedright
descent/closure defect\strut
\end{minipage}\tabularnewline
\begin{minipage}[t]{0.08\columnwidth}\raggedright
\(P_{2}\)\strut
\end{minipage} & \begin{minipage}[t]{0.46\columnwidth}\raggedright
support gate / admissibility update\strut
\end{minipage} & \begin{minipage}[t]{0.37\columnwidth}\raggedright
gate/support/admissibility witness\strut
\end{minipage}\tabularnewline
\begin{minipage}[t]{0.08\columnwidth}\raggedright
\(P_{3}\)\strut
\end{minipage} & \begin{minipage}[t]{0.46\columnwidth}\raggedright
route-memory / mismatch update\strut
\end{minipage} & \begin{minipage}[t]{0.37\columnwidth}\raggedright
route mismatch / route-memory witness\strut
\end{minipage}\tabularnewline
\begin{minipage}[t]{0.08\columnwidth}\raggedright
\(P_{4}\)\strut
\end{minipage} & \begin{minipage}[t]{0.46\columnwidth}\raggedright
refinement / staging update\strut
\end{minipage} & \begin{minipage}[t]{0.37\columnwidth}\raggedright
refinement/stage witness\strut
\end{minipage}\tabularnewline
\begin{minipage}[t]{0.08\columnwidth}\raggedright
\(P_{5}\)\strut
\end{minipage} & \begin{minipage}[t]{0.46\columnwidth}\raggedright
package / completion / closure-summary update\strut
\end{minipage} & \begin{minipage}[t]{0.37\columnwidth}\raggedright
package/closure/completion witness\strut
\end{minipage}\tabularnewline
\begin{minipage}[t]{0.08\columnwidth}\raggedright
\(P_{6}\)\strut
\end{minipage} & \begin{minipage}[t]{0.46\columnwidth}\raggedright
audit / provenance / source update\strut
\end{minipage} & \begin{minipage}[t]{0.37\columnwidth}\raggedright
audit/provenance/source witness\strut
\end{minipage}\tabularnewline
\bottomrule
\end{longtable}

\subsubsection{Directed-Cell and Status Recovery}\label{sn:D.1.4}

Each PICA cell is realized as \(\Gamma;\mathcal T_{\mathrm{PICA}};I_{\mathrm{PICA}}\vdash\mathsf{Cell}_{ij}^{\lambda_{\mathrm{PICA}}}:\sigma_{ij}\) with cell record \((P_{i}\leftarrow P_{j},W_{j},U_{i},L,V,\Theta,A,\delta)\) and PICA cell profile \(\lambda_{\mathrm{PICA}}\). If the realized records meet the conditions of Subsubsection \hyperref[sn:13.2.5]{13.2.5}, the recovered \(6\times 6\) multiplicity count under the formal classifier is \(25\;\texttt{action},\;8\;\texttt{implicit},\;2\;\texttt{trivial},\;1\;\texttt{undefined\_circular}\), recorded as a PICA-profile model fact rather than a universal \(6\times 6\) law (PICA nonclaim 2 of Subsubsection \hyperref[sn:13.2.10]{13.2.10}).

The raw-to-formal status recovery map is

\[
\texttt{Action}\;\to\;\texttt{action},\quad\texttt{Implicit}\;\to\;\texttt{implicit},\quad\texttt{Trivial}\;\to\;\texttt{trivial},
\]

with \(\texttt{Undefined}\) recovering to a reason-dependent formal status in \(\texttt{blocked}\), \(\texttt{undefined\_circular}\), \(\texttt{outside\_scope}\), \(\texttt{absent}\), \(\texttt{below\_threshold}\), or \(\texttt{collapsed}\), decided by the formal classifier of Subsection \hyperref[sn:5.6]{5.6}. Raw PICA status never overrides the formal verdict.

\subsubsection{Pair-Observable Realization and Provenance}\label{sn:D.1.5}

PICA pair-status recovery: \(\texttt{real}\to\texttt{real}\), \(\texttt{provisional}\to\texttt{real\_provisional}\), \(\texttt{duplicated}\to\texttt{real\_duplicated}\), \(\texttt{blocked}\to\texttt{blocked}\).

A PICA pair observable record is

\[
\begin{aligned}\mathsf{PairObs}_{\{i,j\}}^{\mu}\;=\;(\,&\{P_i,P_j\},\;\mu,\;\mathsf{ObservableType},\;\mathsf{BranchData},\;\mathsf{SourceOfTruth},\;L,\;V,\;\Theta,\\&A,\;\delta,\;\mathcal N,\;\mathsf{forbidsSource},\;\mathsf{partialEvidence},\;\mathsf{Dup}\,),\end{aligned}
\]

with realness condition \(\mathsf{SourceOfTruthOK}\land\mathsf{VisibilityPasses}\land\mathsf{AuditPasses}\land\mathsf{ThresholdPasses}\). The source-of-truth predicate is

\[
\begin{aligned}
\mathsf{SourceOfTruthOK}(R,Q)\;\Leftrightarrow\;&\mathsf{Committed}(R)\land\mathsf{CorrectProducer}(R,Q)\\
&\land\mathsf{CorrectProvenance}(R,Q)\land\mathsf{NotFallbackUnlessAllowed}(R,Q)\\
&\land\mathsf{Visible}(R)\land\mathsf{Audited}(R).
\end{aligned}
\]

Dirty exclusion: \(R.\mathsf{dirty\_since\_step}\neq\varnothing\Rightarrow\neg\mathsf{SourceOfTruthOK}(R,Q)\). Fallback default: a fallback record supports at most \(\texttt{real\_provisional}\) unless an explicit upgrade bridge is recorded.

\subsubsection{Realization Theorem and Nonclaims}\label{sn:D.1.6}

The sheet records the model-grade realization statement corresponding to Model Theorem \hyperref[cl:thm:33]{33}. If \(\mathcal P\in\mathsf{PICAReal}_{\mathrm{fin}}\) is admissible, then

\[
\Gamma;\mathsf H_{\mathrm{PICA}};I_{\mathrm{PICA}}\vdash\mathsf{PICAReal}_{\mathrm{fin}}\models\mathbf{BirdInt}^{\mathrm{aud}}_{\mathrm{fin\text{-}stoch}}{}^{\mathrm{cell}/\mathrm{pair}/\mathrm{prov}}:\texttt{model\_realization}.
\]

The PICA nonclaim register \(\mathcal N_{\mathrm{PICA}}\) contains the ten nonclaims of Subsubsection \hyperref[sn:13.2.10]{13.2.10}: no PICA universality, no \(6\times 6\) universal law, raw status not formal status, \(\texttt{action}\) not pair-realness, fallback without an admitted upgrade not a real source, ablation evidence is model-grade, \(P_{3}\) not \(\mathrm{P6}_{\mathrm{drive}}\) without bridge, dirty records not real, finite stochastic fragment only, implementation correctness future work.

\subsection{Cantor Sheet}\label{sn:D.2}

The Cantor strict-bridge realization sheet \(\mathsf{CantorShell}_{\mathrm{aud}}\) records the realization of \(\mathbf{StrictBridge}_{0\to 1}^{\mathrm{audited\ shell}}\) on the audited Cantor shell host \(\mathsf H_{\mathrm{CantorShell}}\). The main-text record is in Subsection \hyperref[sn:13.3]{13.3}.

\subsubsection{Host and Realization Object}\label{sn:D.2.1}

The host \(\mathsf H_{\mathrm{CantorShell}}\) is a scoped audited application host built over a fixed audited shell-stable comparison domain \(\mathcal S_{\mathrm{aud}}\). The realization object is

\[
\mathsf{CantorShell}_{\mathrm{aud}}\;=\;(\,\mathsf{cite}_{\mathcal S},\;W,\;T_{0},\;T_{1},\;\pi_{0},\;\pi_{1},\;\mathcal F,\;\Delta,\;A,\;\mathcal N\,),
\]

with \(\mathsf{cite}_{\mathcal S}\) the finite citation record for the audited shell-stable comparison domain \(\mathcal S_{\mathrm{aud}}\), \(W\subseteq\mathcal S_{\mathrm{aud}}\) the finite audited subset of Subsubsection \hyperref[sn:13.3.2]{13.3.2} containing a nonfactorization pair, \(T_{0}\) the base cocycle-pressure theory, \(T_{1}\) the completion-packaged theory, \(\pi_{0}\) the old cocycle-level object map, \(\pi_{1}\) the new packaged-object map, \(\mathcal F\) the persistent packaged-strata family, \(\Delta\) the pressure-gap consequence, \(A\) the audit/theoremlet support record, and \(\mathcal N\) the Cantor nonclaim register. The host instrument is \(I_{\mathrm{Cantor}}\).

\subsubsection{\texorpdfstring{Base Package \(T_{0}\) and Target Package \(T_{1}\)}{D.2.2 Base Package T\_\{0\} and Target Package T\_\{1\}}}\label{sn:D.2.2}

The base package is \(\mathcal T_{0}^{\mathrm{Cantor}}=(W,\pi_{0}|_W,\Sigma_{\pi_{0}}^{\mathrm{cocycle}},E_{0}^{\mathrm{cocycle}},\mathcal A_{T_{0}})\), with \(\pi_{0}|_W:W\to\mathcal O_{0}\), \(E_{0}^{\mathrm{cocycle}}=E_{\pi_0}\) the lens completion on \(\mathcal P(W)\), the closure \(\operatorname{Cl}_{P_{T_{0}}}\) cited from the Cantor paper, not recorded, and base pressure object

\[
P_{T_{0}}(s)\;=\;\lim_{n\to\infty}\frac{1}{n}\,\log\!\left(\sup_{x\in\mathcal S_{\mathrm{aud}}}a_{n}(s;x)\right).
\]

The target package is \(\mathcal T_{1}^{\mathrm{Cantor}}=(W,\pi_{1}|_W,\Sigma_{\pi_{1}}^{\mathrm{pkg}},E_{1}^{\mathrm{pkg}},\mathcal A_{T_{1}})\), with \(\pi_{1}|_W:W\to\mathcal O_{1}\) completion endomap \(E_{1}^{\mathrm{pkg}}=E_{\pi_1}\) on \(\mathcal P(W)\), and attached packaging citation records \((E_{\tau,\ell},\mathrm{Fix}(E_{\tau,\ell}),\mathcal F,\mathrm{Sat},\mathrm{Force}_{4\leftarrow 5})\) with \(E_{\tau,\ell}(\mu)=U_{\ell}(Q_{\ell}(\mu K^{\tau}))\).

For \(k=0,1\), \(\mathcal A_{T_k}\) is the finite audit functional on the attached audit record family \(A_{T_k}\), with \(\mathcal A_{T_k}(e)=\texttt{pass}\) for every entry \(e\in A_{T_k}\).

\subsubsection{Cantor Strict Bridge}\label{sn:D.2.3}

The Cantor strict bridge is

\[
B_{0\to 1}^{\mathrm{Cantor}}\;=\;(\,\mathcal T_{0}^{\mathrm{Cantor}},\;\mathcal T_{1}^{\mathrm{Cantor}},\;\pi_{0},\;\pi_{1},\;L_{\mathrm{Cantor}},\;V_{\mathrm{Cantor}},\;\Theta_{\mathrm{Cantor}},\;A_{\mathrm{Cantor}},\;\delta_{\mathrm{Cantor}},\;\mathcal N_{\mathrm{Cantor}}\,),
\]

where \(L_{\mathrm{Cantor}}\) records the level pair \((\mathsf{Ver},\mathsf{Ver})\), the identity lens on \(W\) and the interface \((\pi_0,\pi_1)\), and references the attached lift data \((E_{\tau,\ell},Q_\ell,U_\ell,\mathrm{Sat},\mathrm{Force}_{4\leftarrow5})\); the defect bundle \(\delta_{\mathrm{Cantor}}\) contains the computed promotion-bridge defects and the typed pressure-gap consequence defect; citation records for the Cantor saturation, forcing and macro defects are attached to \(T_1\) and referenced by \(L_{\mathrm{Cantor}}\), and they feed no gate.

\subsubsection{Strictness Witness and Pressure-Gap Defect}\label{sn:D.2.4}

The strictness witness is \(\pi_1|_W\not\factor\pi_0|_W\), equivalently \(\exists x,x'\in W:\pi_0(x)=\pi_0(x')\land\pi_1(x)\ne\pi_1(x')\). Define

\[
\Delta_{\mathrm{fact}}^{\mathrm{Cantor}}\;=\;\{\,(x,x')\in W\times W:\pi_0(x)=\pi_0(x')\land\pi_1(x)\ne\pi_1(x')\,\}.
\]

The strictness gate \(G_{\mathrm{strict}}=\texttt{strict\_pass}\Leftrightarrow\Delta_{\mathrm{fact}}^{\mathrm{Cantor}}\neq\varnothing\).

The pressure-gap consequence defect is

\[
\Delta_{T_{0}\to T_{1}}(s;x)\;=\;P_{T_{0}}(s)\;-\;\sum_{\Sigma\in\mathcal F(x)}w_{\Sigma}(x)\,P_{\Sigma}(s).
\]

It is a \emph{consequence} defect, not the strictness witness itself.

\subsubsection{Gate Results}\label{sn:D.2.5}

\begin{longtable}[]{@{}ll@{}}
\toprule
gate & result\tabularnewline
\midrule
\endhead
\(G_{\mathrm{suff}}\) & \(\texttt{pass}\)\tabularnewline
\(G_{\mathrm{desc}}\) & \(\texttt{not\_required}\)\tabularnewline
\(G_{\mathrm{stab}}\) & \(\texttt{not\_required}\)\tabularnewline
\(G_{\mathrm{ctrl}}\) & \(\texttt{pass}\)\tabularnewline
\(G_{\mathrm{nosmuggle}}\) & \(\texttt{pass}\)\tabularnewline
\(G_{\mathrm{vis}}\) & \(\texttt{pass}\)\tabularnewline
\(G_{\mathrm{audit}}\) & \(\texttt{pass}\)\tabularnewline
\(G_{\mathrm{strict}}\) & \(\texttt{strict\_pass}\)\tabularnewline
\(G_{\mathrm{locglob}}\) & \(\texttt{not\_required}\)\tabularnewline
\bottomrule
\end{longtable}

\subsubsection{Realization Theorem and Nonclaims}\label{sn:D.2.6}

The sheet records the model-grade realization statement corresponding to Model Theorem \hyperref[cl:thm:34]{34}. If \(\mathsf{AdmCantorShell}(\mathsf{CantorShell}_{\mathrm{aud}})\) holds, the strictness witness records a nonempty \(\Delta_{\mathrm{fact}}^{\mathrm{Cantor}}\) on \(W\), and the remaining hypotheses of Model Theorem \hyperref[cl:thm:34]{34} hold (so that, by that theorem, the gate record of D.2.5 holds), then

\[
\Gamma;\mathsf H_{\mathrm{CantorShell}};I_{\mathrm{Cantor}}\vdash\mathsf{CantorShell}_{\mathrm{aud}}\models\mathbf{StrictBridge}_{0\to 1}^{\mathrm{audited\ shell}}:\texttt{model\_realization},
\]

and the associated promotion judgment is \(\Gamma;\mathcal T_{0}^{\mathrm{Cantor}};I_{\mathrm{Cantor}}\vdash\mathrm{Promote}(B_{0\to 1}^{\mathrm{Cantor}}):\texttt{strict}\).

The Cantor nonclaim register \(\mathcal N_{\mathrm{Cantor}}\) contains the ten nonclaims of Subsubsection \hyperref[sn:13.3.9]{13.3.9}: shell-bound only, no shell-general theorem, no broader external-class theorem, strictness not macro closure, strictness not drive, pressure gap not KL closure deficit without bridge, pressure gap not strictness witness, \(\Delta_{T_{0}\to T_{1}}\) is conditional disintegration not stratumwise root separation, model realization not universal strict-extension theory, \(T_{1}\) not single repeated idempotent.

\subsection{Abstract-Interpretation Sheet}\label{sn:D.3}

The abstract-interpretation realization sheet \(\mathsf{AIFam}_{\mathrm{fin}}\) records the realization of the role fragment \(P_{1}/P_{2}/P_{5}/P_{6}\), and of \(P_{1}/P_{2}/P_{4}/P_{5}/P_{6}\) when the refinement records contain a pair with \(\rho_k\ne\rho_\ell\), on the finite abstract-interpretation host \(\mathsf H_{\mathrm{AI}}\), supported by the theorem-grade theorems 26--28 of Section \hyperref[sn:12]{12} and the model-grade Model Theorem \hyperref[cl:thm:29]{29}.

\subsubsection{Host}\label{sn:D.3.1}

A finite abstract-interpretation host is the eight-component finite typed record

\[
H\;=\;(\,C,\;A,\;\alpha,\;\gamma,\;\rho,\;F,\;F^{\sharp},\;A_{\mathrm{AI}}\,),
\]

with \((C,\le_{C})\) a finite concrete ordered domain, \((A,\le_{A})\) a finite abstract ordered domain, \(\alpha\dashv\gamma\) a Galois connection between \(C\) and \(A\), \(\rho=\gamma\alpha:C\to C\) the host closure operator, \(F:C\to C\) and \(F^{\sharp}:A\to A\) monotone transformers, and \(A_{\mathrm{AI}}\) the host audit record. The host instrument is \(I_{\mathrm{AI}}\).

\subsubsection{Single-Host Record and Field Glosses}\label{sn:D.3.2}

\begin{longtable}[]{@{}ll@{}}
\toprule
field & meaning\tabularnewline
\midrule
\endhead
\(C\) & finite concrete ordered domain\tabularnewline
\(A\) & finite abstract ordered domain\tabularnewline
\(\alpha:C\to A\) & abstraction map\tabularnewline
\(\gamma:A\to C\) & concretization map\tabularnewline
\(\rho=\gamma\alpha\) & host closure operator on \(C\)\tabularnewline
\(F:C\to C\) & monotone concrete transformer\tabularnewline
\(F^{\sharp}:A\to A\) & monotone abstract transformer\tabularnewline
\(A_{\mathrm{AI}}\) & host audit record\tabularnewline
\bottomrule
\end{longtable}

For the model-family realization, the finite family is recorded as

\[
\mathsf{AIFam}_{\mathrm{fin}}\;=\;(\,\{H_k\}_{k\in\mathcal J},\;\mathsf{Refine},\;\mathsf{Audit},\;\mathcal N_{\mathrm{AI}}\,),
\]

with finite index set \(\mathcal J\) and

\[
H_k=(\,C,\;A_k,\;\alpha_k,\;\gamma_k,\;\rho_k,\;F,\;F_k^{\sharp},\;A_{\mathrm{AI},k}\,).
\]

A refinement edge \(A_k\leadsto A_\ell\) records \(\rho_k\le_C\rho_\ell\) pointwise (the certificate \(R_{k\to\ell}\) of Subsection \hyperref[sn:12.5]{12.5}); a forgetting map \(r_{k\ell}:A_k\to A_\ell\) with \(\alpha_\ell=r_{k\ell}\alpha_k\) supplies it, since \(\alpha_k\gamma_k\alpha_k=\alpha_k\) (by the argument for \(\gamma\alpha\gamma=\gamma\) in the proof of Theorem \hyperref[cl:thm:28]{28}, with unit and counit exchanged), \(\alpha_\ell\rho_k(c)=r_{k\ell}\alpha_k\gamma_k\alpha_k(c)=\alpha_\ell(c)\), so \(\rho_k(c)\le_C\rho_\ell(c)\) by the adjunction. The canonical AI package for a host is \(\mathcal T_{\mathrm{AI}}=(C,\alpha,\Sigma_\alpha,\rho,\mathcal A_{\mathrm{AI}})\), with \(\Sigma_\alpha=\{\alpha^{-1}(B):B\subseteq A\}\) in full finite-lens mode; the audit functional \(\mathcal A_{\mathrm{AI}}\) is assembled from the finite host audit record \(A_{\mathrm{AI}}\).

\subsubsection{Realized Roles and Theorem Trio}\label{sn:D.3.3}

The realization map sends each realized primitive to its AI-host witness:

\begin{longtable}[]{@{}lll@{}}
\toprule
primitive & AI realization & core condition\tabularnewline
\midrule
\endhead
\(P_{1}\) & descent / sound transformer & \(\alpha F\le_{A}F^{\sharp}\alpha\)\tabularnewline
\(P_{2}\) & representability & \(p\;\text{representable}\iff\rho(p)=p\)\tabularnewline
\(P_{4}\) & refinement of abstract domain & \(A_{k}\rightsquigarrow A_{\ell}\)\tabularnewline
\(P_{5}\) & closure / packaging & \(\rho=\gamma\alpha,\;\rho^{2}=\rho\)\tabularnewline
\(P_{6}\) & audit / provenance & \(A_{\mathrm{AI}}\) checks maps and defects\tabularnewline
\bottomrule
\end{longtable}

The supporting theorems are:

\begin{itemize}
\tightlist
\item
  \textbf{Theorem \hyperref[cl:thm:26]{26} (AIClosureAsPackaging).} \(\rho=\gamma\alpha\) is extensive, monotone, and idempotent.
\item
  \textbf{Theorem \hyperref[cl:thm:27]{27} (AIDescentAsSoundTransformer).} \(\alpha F\le_{A}F^{\sharp}\alpha\iff F\gamma\le_{C}\gamma F^{\sharp}\).
\item
  \textbf{Theorem \hyperref[cl:thm:28]{28} (AIRepresentabilityAsFixedness).} \(p\,\text{representable}\iff\rho(p)=p\), with \(\mathrm{Rep}_{\gamma}(C)=\mathrm{Fix}(\rho)=\mathrm{im}(\gamma)\).
\end{itemize}

\subsubsection{\texorpdfstring{Optional \(P_{3}\) Extension}{D.3.4 Optional P\_\{3\} Extension}}\label{sn:D.3.4}

A \(P_{3}\) extension may be added by recording two AI host closures \(\rho_{1}=\gamma_{1}\alpha_{1}\) and \(\rho_{2}=\gamma_{2}\alpha_{2}\) and the noncommutation witness \(\rho_{1}\rho_{2}\neq\rho_{2}\rho_{1}\), in which case Model Theorem \hyperref[cl:thm:29]{29} also realizes \(P_3\), with witness the route-mismatch record \((\rho_1,\rho_2,c_{12})\), where \(c_{12}\in\Delta_3^{\mathrm{comp}}(\rho_1,\rho_2)\); no directed cell is realized.

\subsubsection{Realization Theorem and Nonclaims}\label{sn:D.3.5}

The sheet records the model-grade realization statement corresponding to Model Theorem \hyperref[cl:thm:29]{29}. If \(\mathsf{AIFam}_{\mathrm{fin}}\) is an admissible finite family of AI hosts whose refinement records contain a pair \((k,\ell)\) with \(\rho_k\neq\rho_\ell\), then

\[
\Gamma;\mathsf H_{\mathrm{AI}};I_{\mathrm{AI}}\vdash\mathsf{AIFam}_{\mathrm{fin}}\models\mathbf{BirdInt}^{\mathrm{aud}}_{\mathrm{fin}}{}^{\,P_{1}/P_{2}/P_{4}/P_{5}/P_{6}}:\texttt{model\_realization}.
\]

The AI nonclaim register \(\mathcal N_{\mathrm{AI}}\) contains: AI is not the whole calculus, soundness is not completeness, closure does not imply descent, representability does not imply sound transformer, refinement does not automatically improve every defect, AI realization includes \(P_{3}\) only when two host closures fail to commute, abstract claims cannot overread concrete distinctions invisible to \(\alpha\), audit of soundness is not instrument-free truth, AI model-family claims are finite-host scoped, and \(P_{5}\)-closure does not imply drive, and the three coverage nonclaims of Subsubsection \hyperref[sn:13.6.2]{13.6.2}.

\subsection{Graph/Cohomology Sheet}\label{sn:D.4}

The graph/cohomology realization sheet \(\mathsf{CohGraph}_{\mathrm{fin}}\) records the realization of the drive-support fragment \(P_{2}/\mathrm{P6}_{\mathrm{drive}}^{H^{1}}/P_{3}\text{-separation}\) on the finite graph/cohomology host \(\mathsf H_{\mathrm{graph}}\). The main-text record is in Subsection \hyperref[sn:13.4]{13.4}.

\subsubsection{Host and Realization Object}\label{sn:D.4.1}

The realization object is

\[
\mathsf{CohGraph}_{\mathrm{fin}}\;=\;(\,G,\;C^{0}(G;k),\;C^{1}(G;k),\;d,\;H^{1}(G;k),\;\mathsf{Cycles}(G),\;a,\;\mathsf{Gate},\;A,\;V,\;\mathcal N\,),
\]

with \(G=(V,E)\) a finite support graph, \(C^{0}(G;k)\) and \(C^{1}(G;k)\) the finite-dimensional cochain spaces, \(d:C^{0}\to C^{1}\) the coboundary, \(H^{1}(G;k)\) the first cohomology, \(\mathsf{Cycles}(G)\) the finite-dimensional cycle space, \(a\in C^{1}(G;k)\) the affinity/drive cochain, \(\mathsf{Gate}\) the gating record, \(A\) and \(V\) host audit and visibility records, and \(\mathcal N\) the cohomology nonclaim register. The host instrument is \(I_{\mathrm{graph}}\).

\subsubsection{Realized Fragment and Realization Map}\label{sn:D.4.2}

The realization fragment is \(\mathbf{BirdInt}^{\mathrm{aud}}_{\mathrm{fin}}{}^{\,P_{2}/\mathrm{P6}_{\mathrm{drive}}^{H^{1}}/P_{3}\text{-separation}}\), with realization map

\begin{longtable}[]{@{}ll@{}}
\toprule
\begin{minipage}[b]{0.34\columnwidth}\raggedright
primitive\strut
\end{minipage} & \begin{minipage}[b]{0.60\columnwidth}\raggedright
cohomology realization\strut
\end{minipage}\tabularnewline
\midrule
\endhead
\begin{minipage}[t]{0.34\columnwidth}\raggedright
\(P_{2}\)\strut
\end{minipage} & \begin{minipage}[t]{0.60\columnwidth}\raggedright
support gating / edge deletion / cycle-rank capacity\strut
\end{minipage}\tabularnewline
\begin{minipage}[t]{0.34\columnwidth}\raggedright
\(\mathrm{P6}_{\mathrm{drive}}^{H^{1}}\)\strut
\end{minipage} & \begin{minipage}[t]{0.60\columnwidth}\raggedright
cohomological drive certificate \([a]\neq 0\)\strut
\end{minipage}\tabularnewline
\begin{minipage}[t]{0.34\columnwidth}\raggedright
\(P_{6}^{\mathrm{audit}}\)\strut
\end{minipage} & \begin{minipage}[t]{0.60\columnwidth}\raggedright
audit of cochains, cycles, source-of-truth, and visibility\strut
\end{minipage}\tabularnewline
\begin{minipage}[t]{0.34\columnwidth}\raggedright
\(P_{3}\)\strut
\end{minipage} & \begin{minipage}[t]{0.60\columnwidth}\raggedright
optional route/holonomy witness, kept separate from drive\strut
\end{minipage}\tabularnewline
\begin{minipage}[t]{0.34\columnwidth}\raggedright
\(P_{4}\)\strut
\end{minipage} & \begin{minipage}[t]{0.60\columnwidth}\raggedright
optional local/global cover or gluing of cochains\strut
\end{minipage}\tabularnewline
\begin{minipage}[t]{0.34\columnwidth}\raggedright
\(P_{5}\)\strut
\end{minipage} & \begin{minipage}[t]{0.60\columnwidth}\raggedright
optional quotient/support package of graph state\strut
\end{minipage}\tabularnewline
\bottomrule
\end{longtable}

The core sheet focuses on \(P_{2}\) and \(\mathrm{P6}_{\mathrm{drive}}^{H^{1}}\).

\subsubsection{Drive and No-Drive Certificates}\label{sn:D.4.3}

The drive certificate is \(\kappa_{6}^{\mathrm{drive}}(a):[a]\neq 0\in H^{1}(G;k)\), equivalent to \(\exists\,\gamma\in\mathsf{Cycles}(G):\oint_{\gamma}a\neq 0\). No-drive certificates: \([a]=0\), \(a=d\phi\), \(H^{1}(G;k)=0\), \(\beta_{1}(G)=0\).

\subsubsection{Gating Realization}\label{sn:D.4.4}

A \(P_{2}\)-gate \(G\rightsquigarrow G'\) has cycle-rank loss defect \(\delta_{2}^{\mathrm{cycle}}(G,G')=\beta_{1}(G)-\beta_{1}(G')\) and forest threshold \(\Theta_{2}^{\mathrm{forest}}:\beta_{1}(G')=0\). When \(G'\) is a forest, \(H^{1}(G';k)=0\) and no certificate \(\mathrm{P6}_{\mathrm{drive}}^{H^{1}}\) exists on \(G'\).

\subsubsection{Core Theorems}\label{sn:D.4.5}

The core graph/cohomology theorems of Section \hyperref[sn:8]{8} (with long-form proofs in Appendix \hyperref[sn:B.2]{B.2}) are: forest no-drive (Theorem \hyperref[cl:thm:15]{15}: \(\beta_{1}(G)=0\Rightarrow H^{1}(G;k)=0\)), exact-form null-drive (Theorem \hyperref[cl:thm:16]{16}: \(a=d\phi\Rightarrow\oint_{\gamma}a=0\)), nonzero-affinity equivalence (Theorem \hyperref[cl:thm:17]{17}: \([a]\neq 0\iff\exists\gamma:\oint_{\gamma}a\neq 0\)), and gating-affinity suppression (Theorem \hyperref[cl:thm:18]{18}: gating to a forest kills cycle-supported drive).

\subsubsection{Realization Theorem and Nonclaims}\label{sn:D.4.6}

\textbf{Cohomology Support Realization (model-grade).} Every admissible \(\mathsf{CohGraph}_{\mathrm{fin}}\) realizes the fragment containing \(P_{2}\)-support gating and cycle-rank loss, \(\mathrm{P6}_{\mathrm{drive}}^{H^{1}}\) certificates, no-drive certificates, and \(P_{3}\)/drive separation:

\[
\Gamma;\mathsf H_{\mathrm{graph}};I_{\mathrm{graph}}\vdash\mathsf{CohGraph}_{\mathrm{fin}}\models\mathbf{BirdInt}^{\mathrm{aud}}_{\mathrm{fin}}{}^{\,P_{2}/\mathrm{P6}_{\mathrm{drive}}^{H^{1}}/P_{3}\text{-separation}}:\texttt{model\_realization}.
\]

The cohomology nonclaim register \(\mathcal N_{\mathrm{coh}}\) contains the three coverage nonclaims of Subsubsection \hyperref[sn:13.6.2]{13.6.2} and the entries of Subsubsection \hyperref[sn:13.4.4]{13.4.4}: \(\mathrm{P6}_{\mathrm{drive}}^{H^{1}}\) is one drive certificate family not all drive, \(P_{3}\) holonomy not \(\mathrm{P6}_{\mathrm{drive}}\) without a bridge, gating to a forest kills cycle-supported drive only, exact \(1\)-cochains have zero cycle drive, \(H^{1}(G;k)=0\) is a no-drive certificate for cohomological drive only, nonzero holonomy can coexist with zero drive cochain, host is finite and scoped, cohomology does not replace PICA/Cantor/AI, no edge-support claim accepted without visibility/audit, positive holonomy-to-drive claims require \(B_{3\to 6}^{\mathrm{drive}}\).

\subsection{Combined Coverage Matrix}\label{sn:D.5}

This subsection records the combined coverage matrix across the four model realizations. The matrix has one row per realization and one column per scoped attribute (host/instrument, realized fragment, claim strength); Subsubsection \hyperref[sn:D.5.2]{D.5.2} separately records per-primitive role coverage. The matrix collects the per-realization data of Subsections \hyperref[sn:D.1]{D.1}--\hyperref[sn:D.4]{D.4} in one place; it is the long-form analogue of the matrix recorded in Subsection \hyperref[sn:13.5]{13.5}.

\subsubsection{Realization Matrix}\label{sn:D.5.1}

\begingroup\footnotesize

\begin{longtable}[]{@{}llll@{}}
\toprule
\begin{minipage}[b]{0.07\columnwidth}\raggedright
model\strut
\end{minipage} & \begin{minipage}[b]{0.22\columnwidth}\raggedright
host / instrument\strut
\end{minipage} & \begin{minipage}[b]{0.41\columnwidth}\raggedright
realized fragment\strut
\end{minipage} & \begin{minipage}[b]{0.19\columnwidth}\raggedright
claim strength\strut
\end{minipage}\tabularnewline
\midrule
\endhead
\begin{minipage}[t]{0.07\columnwidth}\raggedright
\(\mathsf{PICAReal}_{\mathrm{fin}}\)\strut
\end{minipage} & \begin{minipage}[t]{0.22\columnwidth}\raggedright
\(\mathsf H_{\mathrm{PICA}}\) / \(I_{\mathrm{PICA}}\)\strut
\end{minipage} & \begin{minipage}[t]{0.41\columnwidth}\raggedright
\(\mathbf{BirdInt}^{\mathrm{aud}}_{\mathrm{fin\text{-}stoch}}{}^{\mathrm{cell}/\mathrm{pair}/\mathrm{prov}}\)\strut
\end{minipage} & \begin{minipage}[t]{0.19\columnwidth}\raggedright
\(\texttt{model\_realization}\)\strut
\end{minipage}\tabularnewline
\begin{minipage}[t]{0.07\columnwidth}\raggedright
\(\mathsf{CantorShell}_{\mathrm{aud}}\)\strut
\end{minipage} & \begin{minipage}[t]{0.22\columnwidth}\raggedright
\(\mathsf H_{\mathrm{CantorShell}}\) / \(I_{\mathrm{Cantor}}\)\strut
\end{minipage} & \begin{minipage}[t]{0.41\columnwidth}\raggedright
\(\mathbf{StrictBridge}_{0\to 1}^{\mathrm{audited\ shell}}\)\strut
\end{minipage} & \begin{minipage}[t]{0.19\columnwidth}\raggedright
\(\texttt{model\_realization}\) (application-scoped)\strut
\end{minipage}\tabularnewline
\begin{minipage}[t]{0.07\columnwidth}\raggedright
\(\mathsf{AIFam}_{\mathrm{fin}}\)\strut
\end{minipage} & \begin{minipage}[t]{0.22\columnwidth}\raggedright
\(\mathsf H_{\mathrm{AI}}\) / \(I_{\mathrm{AI}}\)\strut
\end{minipage} & \begin{minipage}[t]{0.41\columnwidth}\raggedright
\(\mathbf{BirdInt}^{\mathrm{aud}}_{\mathrm{fin}}{}^{\,P_{1}/P_{2}/P_{5}/P_{6}}\); with \(P_4\) given a strict refinement; with \(P_3\) given two noncommuting closures\strut
\end{minipage} & \begin{minipage}[t]{0.19\columnwidth}\raggedright
\(\texttt{model\_realization}\)\strut
\end{minipage}\tabularnewline
\begin{minipage}[t]{0.07\columnwidth}\raggedright
\(\mathsf{CohGraph}_{\mathrm{fin}}\)\strut
\end{minipage} & \begin{minipage}[t]{0.22\columnwidth}\raggedright
\(\mathsf H_{\mathrm{graph}}\) / \(I_{\mathrm{graph}}\)\strut
\end{minipage} & \begin{minipage}[t]{0.41\columnwidth}\raggedright
\(\mathbf{BirdInt}^{\mathrm{aud}}_{\mathrm{fin}}{}^{\,P_{2}/\mathrm{P6}_{\mathrm{drive}}^{H^{1}}/P_{3}\text{-separation}}\)\strut
\end{minipage} & \begin{minipage}[t]{0.19\columnwidth}\raggedright
\(\texttt{model\_realization}\)\strut
\end{minipage}\tabularnewline
\bottomrule
\end{longtable}

\endgroup

\subsubsection{Cross-Model Role Coverage}\label{sn:D.5.2}

The cross-model role-coverage table records, per primitive, what each realization carries within its declared fragment:

\begin{longtable}[]{@{}lllll@{}}
\toprule
\begin{minipage}[b]{0.05\columnwidth}\raggedright
role\strut
\end{minipage} & \begin{minipage}[b]{0.18\columnwidth}\raggedright
\(\mathsf{PICA}\)\strut
\end{minipage} & \begin{minipage}[b]{0.20\columnwidth}\raggedright
\(\mathsf{Cantor}\)\strut
\end{minipage} & \begin{minipage}[b]{0.25\columnwidth}\raggedright
\(\mathsf{AI}\)\strut
\end{minipage} & \begin{minipage}[b]{0.19\columnwidth}\raggedright
\(\mathsf{CohGraph}\)\strut
\end{minipage}\tabularnewline
\midrule
\endhead
\begin{minipage}[t]{0.05\columnwidth}\raggedright
\(P_{1}\)\strut
\end{minipage} & \begin{minipage}[t]{0.18\columnwidth}\raggedright
descent / closure update\strut
\end{minipage} & \begin{minipage}[t]{0.20\columnwidth}\raggedright
macro obstruction; cited, not realized\strut
\end{minipage} & \begin{minipage}[t]{0.25\columnwidth}\raggedright
sound descent \(\alpha F\le_{A}F^{\sharp}\alpha\)\strut
\end{minipage} & \begin{minipage}[t]{0.19\columnwidth}\raggedright
not primary\strut
\end{minipage}\tabularnewline
\begin{minipage}[t]{0.05\columnwidth}\raggedright
\(P_{2}\)\strut
\end{minipage} & \begin{minipage}[t]{0.18\columnwidth}\raggedright
gating / support / admissibility\strut
\end{minipage} & \begin{minipage}[t]{0.20\columnwidth}\raggedright
admissibility obstruction (cited, not realized)\strut
\end{minipage} & \begin{minipage}[t]{0.25\columnwidth}\raggedright
representability \(\rho(p)=p\)\strut
\end{minipage} & \begin{minipage}[t]{0.19\columnwidth}\raggedright
support gating\strut
\end{minipage}\tabularnewline
\begin{minipage}[t]{0.05\columnwidth}\raggedright
\(P_{3}\)\strut
\end{minipage} & \begin{minipage}[t]{0.18\columnwidth}\raggedright
route mismatch\strut
\end{minipage} & \begin{minipage}[t]{0.20\columnwidth}\raggedright
not central\strut
\end{minipage} & \begin{minipage}[t]{0.25\columnwidth}\raggedright
noncommuting closures \(\Delta_3^{\mathrm{comp}}(\rho_k,\rho_\ell)\), when two closures fail to commute\strut
\end{minipage} & \begin{minipage}[t]{0.19\columnwidth}\raggedright
separated from drive\strut
\end{minipage}\tabularnewline
\begin{minipage}[t]{0.05\columnwidth}\raggedright
\(P_{4}\)\strut
\end{minipage} & \begin{minipage}[t]{0.18\columnwidth}\raggedright
refinement / staging\strut
\end{minipage} & \begin{minipage}[t]{0.20\columnwidth}\raggedright
\(P_{4}\!\leftarrow\!P_{5}\) completion-packaging (cited, not realized)\strut
\end{minipage} & \begin{minipage}[t]{0.25\columnwidth}\raggedright
refinement \(A_{k}\rightsquigarrow A_{\ell}\)\strut
\end{minipage} & \begin{minipage}[t]{0.19\columnwidth}\raggedright
optional local/global cover\strut
\end{minipage}\tabularnewline
\begin{minipage}[t]{0.05\columnwidth}\raggedright
\(P_{5}\)\strut
\end{minipage} & \begin{minipage}[t]{0.18\columnwidth}\raggedright
package / completion\strut
\end{minipage} & \begin{minipage}[t]{0.20\columnwidth}\raggedright
completion-packaged objects\strut
\end{minipage} & \begin{minipage}[t]{0.25\columnwidth}\raggedright
closure / packaging \(\rho=\gamma\alpha\)\strut
\end{minipage} & \begin{minipage}[t]{0.19\columnwidth}\raggedright
optional graph package\strut
\end{minipage}\tabularnewline
\begin{minipage}[t]{0.05\columnwidth}\raggedright
\(P_{6}\)\strut
\end{minipage} & \begin{minipage}[t]{0.18\columnwidth}\raggedright
audit / provenance\strut
\end{minipage} & \begin{minipage}[t]{0.20\columnwidth}\raggedright
audit of shell bridge\strut
\end{minipage} & \begin{minipage}[t]{0.25\columnwidth}\raggedright
audit of maps and soundness\strut
\end{minipage} & \begin{minipage}[t]{0.19\columnwidth}\raggedright
drive certificate \([a]\neq 0\) and audit\strut
\end{minipage}\tabularnewline
\bottomrule
\end{longtable}

As in Subsubsection \hyperref[sn:13.6.1]{13.6.1}, the PICA and Cantor columns record update sorts and citations, not role-channel records classified \(\texttt{active\_projection}\).

\subsubsection{Combined Nonclaims}\label{sn:D.5.3}

The combined nonclaim register records ten cross-model nonclaims, each attached to the realization or realizations it names:

\begin{enumerate}
\def\labelenumi{\arabic{enumi}.}
\tightlist
\item
  No model realization is the whole of \(\mathbf{BirdInt}^{\mathrm{aud}}_{\mathrm{fin}}\).
\item
  PICA is not the universal six-symbol algebra.
\item
  Cantor is audited-shell scoped.
\item
  The abstract-interpretation family realizes \(P_1/P_2/P_5/P_6\), adds \(P_4\) when its refinement records contain a strict refinement, and adds \(P_3\) when two host closures fail to commute.
\item
  The graph/cohomology host realizes cohomological drive, not all drive notions.
\item
  Model-realization evidence is not theorem-grade unless an admissible bridge is supplied.
\item
  Ablations, simulations, and fallbacks require source-of-truth contracts.
\item
  Strictness does not imply macro closure or drive.
\item
  Audit is not final self-certification.
\item
  No model may overread suppressed content without an admissible bridge.
\end{enumerate}

The combined matrix records four scoped realizations and the per-primitive role coverage they admit. The matrix is not read vertically as a coverage assertion: the framework does not assert that the union of the four realizations covers the whole calculus, in keeping with the model-coverage nonclaim of Subsubsection \hyperref[sn:13.6.2]{13.6.2}.
 \section{High-Structure Semantics}\label{sn:E}

This appendix records the long-form definitions of the finite decorated-square fragment of Section \hyperref[sn:14]{14}. Subsection \hyperref[sn:E.1]{E.1} records the fragment in full, and Subsection \hyperref[sn:E.2]{E.2} points to the recovery theorems and nonclaims in Section \hyperref[sn:14]{14}.

The appendix is host-bound to the finite decorated double-category fragment \(\mathsf H_{\mathrm{DC}}\) and supplements Section \hyperref[sn:14]{14}; it does not introduce new theorems beyond Theorems \hyperref[cl:thm:30]{30}, \hyperref[cl:thm:31]{31}, and \hyperref[cl:thm:32]{32}.

\subsection{Finite Decorated-Square Fragment}\label{sn:E.1}

This subsection expands the schema of Subsection \hyperref[sn:14.1]{14.1}. The decorated-square fragment is a secondary organizational layer over the core calculus: it does not replace the directed-cell, pair-observable, promotion, claim, gate, or square status families, and it does not assert a completeness theorem for the calculus.

\subsubsection{Objects, Horizontal Arrows, Vertical Arrows}\label{sn:E.1.1}

The finite decorated double-category fragment \(\mathsf H_{\mathrm{DC}}\) is the finite typed record

\[
\mathsf H_{\mathrm{DC}}\;=\;(\,\mathrm{Ob},\;\mathrm{Hor},\;\mathrm{Ver},\;\mathrm{Sq},\;A_{\mathrm{DC}},\;V_{\mathrm{DC}},\;\mathcal N_{\mathrm{DC}}\,),
\]

with components fixed in Subsubsection \hyperref[sn:14.1.1]{14.1.1}: \(\mathrm{Ob}\) the finite class of package objects \(\mathcal T\) of Subsection \hyperref[sn:3.3]{3.3}; \(\mathrm{Hor}\) and \(\mathrm{Ver}\) the classes of Subsubsection \hyperref[sn:14.1.1]{14.1.1}, which maps sit on which side being fixed by the square type; \(\mathrm{Sq}\) the finite class of decorated squares; \(A_{\mathrm{DC}}\) the host audit record; \(V_{\mathrm{DC}}\) the host visibility record; and \(\mathcal N_{\mathrm{DC}}\) the host nonclaim register.

Compositions of horizontal and vertical arrows are admitted only when the resulting boundary records are finite and audited under \(I_{\mathrm{DC}}\).

\clearpage

\subsubsection{Decorated Square Schema}\label{sn:E.1.2}

A decorated square \(\Xi\in\mathrm{Sq}\) is a finite typed record

\[
\Xi\;=\;(\,X,\;F,\;q,\;F^{\sharp},\;\mathrm{src}_{\Xi},\;\mathrm{tgt}_{\Xi},\;\mathrm{type}_{\Xi},\;\mathsf{Decor}_{\Xi},\;A_{\Xi},\;V_{\Xi},\;\delta_{\Xi},\;\mathcal N_{\Xi},\;\mathsf{partialEvidence}_{\Xi}\,),
\]

with \(\mathrm{type}_{\Xi}\in\{\,\texttt{descent},\;\texttt{quotient},\;\texttt{completion},\;\texttt{promotion},\;\texttt{visibility},\;\texttt{audit},\;\texttt{factorization}\,\}\) and \(\mathsf{Decor}_{\Xi}\) recording the type-specific decoration data.

\begingroup
\setlength{\tabcolsep}{4pt}
\renewcommand{\arraystretch}{1.06}
\begin{center}
\begin{tabularx}{0.96\linewidth}{@{}>{\raggedright\arraybackslash}p{0.19\linewidth}>{\raggedright\arraybackslash}X@{}}
\toprule
field & meaning \\
\midrule
$X$ & source object of the square, an element of $\mathrm{Ob}$ \\
$F$ & horizontal arrow component of the boundary \\
$q$ & vertical arrow component of the boundary \\
$F^{\sharp}$ & bottom horizontal arrow or candidate descended/completed arrow \\
$\mathrm{src}_{\Xi},\mathrm{tgt}_{\Xi}$ & source and target boundary tuples \\
$\mathrm{type}_{\Xi}$ & one of the seven square types listed above \\
$\mathsf{Decor}_{\Xi}$ & type-specific decoration: descent, quotient, completion, promotion, visibility, audit, or factorization data \\
$A_{\Xi},V_{\Xi}$ & per-square audit and visibility records \\
$\delta_{\Xi}$ & per-square defect record, classified in the square-status family \\
$\mathcal N_{\Xi}$ & per-square nonclaim register \\
$\mathsf{partialEvidence}_{\Xi}$ & declared Boolean: the evidence for the square is incomplete \\
\bottomrule
\end{tabularx}
\end{center}
\endgroup

The exactness condition for each square type is recorded in Subsubsection \hyperref[sn:14.1.2]{14.1.2}: descent commutation equality for \(\texttt{descent}\), completion-commutation equality for \(\texttt{completion}\), factorization-filler condition for \(\texttt{factorization}\), and bridge-admissibility conditions for \(\texttt{visibility}\) and \(\texttt{audit}\).

\subsubsection{Strict Pasting}\label{sn:E.1.3}

The strict-pasting law of Subsubsection \hyperref[sn:14.1.3]{14.1.3} bounds the defect of an admissible horizontal paste of two decorated squares \(\Xi_{1}\) and \(\Xi_{2}\) with top edges \(F_1,F_2\):

\[
\delta_{\Xi_{1}\cdot\Xi_{2}}\;\subseteq\;\delta_{\Xi_{1}}\;\cup\;F_1^{-1}(\delta_{\Xi_{2}}),
\]

and the inclusion can be strict (Theorem \hyperref[cl:thm:31]{31}). The paste defect \(\delta_{\mathrm{paste}}(\Xi_{1},\Xi_{2})\) is the finite record of shared edges recorded differently in the two squares; it is empty whenever pasting is admissible. Pasting is admissible when the shared edge data agree on \(\mathrm{Hor}\) or \(\mathrm{Ver}\) and the per-square admissibility conditions hold for both factors. Pseudo-pasting and weak pasting are out of scope.

\subsection{Recovery Theorems and Nonclaims}\label{sn:E.2}

The descent-square, completion-pasting and strict-extension recoveries (Theorems \hyperref[cl:thm:30]{30}, \hyperref[cl:thm:31]{31} and \hyperref[cl:thm:32]{32}) are stated and proved in Subsections \hyperref[sn:14.2]{14.2}, \hyperref[sn:14.3]{14.3} and \hyperref[sn:14.4]{14.4}, and the high-structure nonclaims are recorded in Subsection \hyperref[sn:14.5]{14.5}; this appendix does not repeat them.
 \section{Mechanization Disclosure and Audit Table}\label{sn:F}

This appendix is the canonical disclosure venue for the proof-assistant artifact attached to this paper. The body gives the mathematical statements and proofs; this appendix records the full paper-to-Lean alignment, the wording used to describe Lean coverage, and the validation boundary.

\subsection{Coverage Classes}\label{sn:F.1}

The table of Subsection \hyperref[sn:F.2]{F.2} lists each numbered result and definition with the Lean declarations that support it and one of the following coverage classes.

\begin{itemize}
\tightlist
\item
  \textbf{Direct proof.} The Lean declaration proves the paper statement for the stated finite objects.
\item
  \textbf{Direct proof under stated hypotheses.} The Lean declaration proves the paper statement from hypotheses that the table names in its scope column.
\item
  \textbf{Partial proof.} The Lean declarations prove part of the paper statement, or prove it for a restricted class of objects; the scope column names the part or the restriction, and the paper's proof covers the rest.
\item
  \textbf{Record projection.} The Lean declaration states the result for a finite record whose fields carry the paper's content. The declaration checks the shape of the result; the mathematics is in the paper's proof.
\item
  \textbf{Schema mirror.} A definition-level typed mirror of a paper definition.
\end{itemize}

A model-realization statement is covered at most as a record projection: the Lean records do not verify a concrete implementation of the model host.

\subsection{Paper-to-Lean Disclosure Table}\label{sn:F.2}

\begin{landscape}
\begingroup\scriptsize
\setlength{\tabcolsep}{2pt}
\setlength{\LTleft}{0pt}
\setlength{\LTright}{0pt}
\begin{longtable}{@{}>{\raggedright\arraybackslash}p{0.12\linewidth}>{\raggedright\arraybackslash}p{0.17\linewidth}>{\raggedright\arraybackslash}p{0.08\linewidth}>{\raggedright\arraybackslash}p{0.18\linewidth}>{\raggedright\arraybackslash}p{0.09\linewidth}>{\raggedright\arraybackslash}p{0.16\linewidth}>{\raggedright\arraybackslash}p{0.09\linewidth}@{}}
\toprule
Paper label & Paper name & Body location & Lean declarations & Claim strength & Coverage class & Lean scope\tabularnewline
\midrule
\endhead
\path{def:primitive-labels} & Finite primitive labels P1 through P6. & Section 2 & \path{SixBirdsIII.Primitive} & definition & schema mirror & Typed schema only.\tabularnewline
\path{def:levels} & Behavioral, verification, and structural levels. & Section 2 & \path{SixBirdsIII.Level} & definition & schema mirror & Typed schema only.\tabularnewline
\path{def:hosts} & Finite permitted host inventory. & Section 2 & \path{SixBirdsIII.Host} & definition & schema mirror & Typed schema only.\tabularnewline
\path{def:birdint-domain} & Finite audited interaction calculus domain tuple. & Section 3 & \path{SixBirdsIII.BirdIntDomain} & definition & schema mirror & Typed schema only.\tabularnewline
\path{def:role-status} & Role-channel status family. & Section 5 & \path{SixBirdsIII.RoleStatus} & definition & schema mirror & Typed schema only.\tabularnewline
\path{def:cell-status} & Directed-cell status family. & Section 5 & \path{SixBirdsIII.CellStatus} & definition & schema mirror & Typed schema only.\tabularnewline
\path{def:pair-status} & Pair-observable status family. & Section 5 & \path{SixBirdsIII.PairStatus} & definition & schema mirror & Typed schema only.\tabularnewline
\path{def:promotion-status} & Promotion status family. & Section 5 & \path{SixBirdsIII.PromotionStatus} & definition & schema mirror & Typed schema only.\tabularnewline
\path{def:claim-status} & Instrument-indexed claim status family. & Section 5 & \path{SixBirdsIII.ClaimStatus} & definition & schema mirror & Typed schema only.\tabularnewline
\path{def:gate-status} & Promotion and claim gate status family. & Section 5 & \path{SixBirdsIII.GateStatus} & definition & schema mirror & Typed schema only.\tabularnewline
\path{def:strict-gate-status} & Strictness-gate status family. & Section 5 & \path{SixBirdsIII.StrictGateStatus} & definition & schema mirror & Typed schema only.\tabularnewline
\path{def:square-status} & High-structure square status family. & Section 5 & \path{SixBirdsIII.SqStatus} & definition & schema mirror & Typed schema only.\tabularnewline
\path{def:threshold-tags} & Threshold tag vocabulary. & Section 5 & \path{SixBirdsIII.ThresholdTag} & definition & schema mirror & Typed schema only.\tabularnewline
\path{def:visibility-tags} & Visibility partition tags. & Section 3 & \path{SixBirdsIII.VisibilityTag} & definition & schema mirror & Typed schema only.\tabularnewline
\path{def:profile-record} & Profile record skeleton. & Section 3 & \path{SixBirdsIII.Profile} & definition & schema mirror & Typed schema only.\tabularnewline
\path{def:instrument-record} & Instrument record skeleton. & Section 3 & \path{SixBirdsIII.Instrument} & definition & schema mirror & Typed schema only.\tabularnewline
\path{def:directed-cell-record} & Directed-cell record skeleton. & Section 4 & \path{SixBirdsIII.DirectedCellRecord} & definition & schema mirror & Typed schema only.\tabularnewline
\path{def:pair-observable-record} & Pair-observable record skeleton. & Section 4 & \path{SixBirdsIII.PairObservableRecord} & definition & schema mirror & Typed schema only.\tabularnewline
\path{def:promotion-bridge-record} & Promotion-bridge record skeleton. & Section 4 & \path{SixBirdsIII.PromotionBridgeRecord} & definition & schema mirror & Typed schema only.\tabularnewline
\path{def:defect-record} & Typed defect record skeleton. & Section 5 & \path{SixBirdsIII.DefectRecord} & definition & schema mirror & Typed schema only.\tabularnewline
\path{def:claim-record} & Instrument-indexed claim record skeleton. & Section 4 & \path{SixBirdsIII.ClaimRecord} & definition & schema mirror & Typed schema only.\tabularnewline
\path{def:top-down-channel-record} & Finite top-down channel record schema. & Subsection 4.7 & \path{SixBirdsIII.TopDownChannelRecord} & definition & schema mirror & Typed schema only.\tabularnewline
\path{def:top-down-gates} & Required top-down channel gate vocabulary. & Subsection 4.7 & \path{SixBirdsIII.TopDownGate} & definition & schema mirror & Typed schema only.\tabularnewline
\path{def:struct-down} & Structural downward influence predicate. & Subsection 4.7 & \path{SixBirdsIII.StructDown} & definition & schema mirror & Typed schema only.\tabularnewline
\path{thm:no-total-algebra} & No total primitive-pair algebra recovers directed-cell status. & Section 6 & \path{SixBirdsIII.no_total_algebra}\newline \path{SixBirdsIII.cm4_same_pair_classifier_countermodel} & theorem-grade & partial proof & Lean witness: an action cell and a missing-witness cell on P6/P3 (not well formed in Lean); the paper's action and blocked cells of Theorem 1 are not formalized.\tabularnewline
\path{thm:typed-non-collapse} & Active typed interaction does not identify actor and informant primitives. & Section 6 & \path{SixBirdsIII.typed_non_collapse} & theorem-grade & partial proof & One prefilled typed cell; full admissibility is unproved.\tabularnewline
\path{thm:cell-well-formedness-decidability} & Directed-cell well-formedness is decidable on finite records. & Section 6 & \path{SixBirdsIII.cell_well_formedness_decidable} & theorem-grade & record projection & Seven Boolean fields; no eight-step admissibility checker.\tabularnewline
\path{thm:covered-cell-unique-status} & A covered well-formed directed cell receives exactly one status. & Section 6 & \path{SixBirdsIII.covered_cell_unique_status} & theorem-grade & record projection & Total classifier; coveredness is unused.\tabularnewline
\path{thm:status-family-separation} & Judgment-attached status families remain separated. & Section 6 & \path{SixBirdsIII.status_family_separation}\newline \path{SixBirdsIII.status_family_separation_same_pair}\newline \path{SixBirdsIII.cm4_cell_pair_classifier_separation} & theorem-grade & partial proof & P6/P3 cell-pair classifier branch; other families remain record examples.\tabularnewline
\path{thm:descent-through-quotient} & Finite update descent through a quotient is equivalent to fiber constancy. & Section 7 & \path{SixBirdsIII.descent_through_quotient} & theorem-grade & direct proof & Image descent for finite-map data.\tabularnewline
\path{thm:descent-defect-equivalence} & The descent defect vanishes exactly when descent exists. & Section 7 & \path{SixBirdsIII.descent_defect_equivalence} & theorem-grade & direct proof & Witness form of the finite descent defect.\tabularnewline
\path{thm:finite-markov-closure-deficit} & Finite Markov predictive closure loss equals conditional mutual information under assumptions. & Section 7 & \path{SixBirdsIII.finite_markov_closure_deficit}\newline \path{SixBirdsIII.finite_markov_kl_decomposition}\newline \path{SixBirdsIII.finite_rat_kl_nonnegative}\newline \path{SixBirdsIII.finite_excess_loss_nonnegative}\newline \path{SixBirdsIII.finite_conditional_excess_zero}\newline \path{SixBirdsIII.finite_markov_closure_deficit_from_kl}\newline \path{SixBirdsIII.ratFinSum_add}\newline \path{SixBirdsIII.ratFinSum_sub}\newline \path{SixBirdsIII.ratFinSum_mul}\newline \path{SixBirdsIII.ratFinSum_congr}\newline \path{SixBirdsIII.ratFinSum_zero}\newline \path{SixBirdsIII.ratFinSum_nonneg}\newline \path{SixBirdsIII.ratFinSum_mono}\newline \path{SixBirdsIII.scalarFinSum_add}\newline \path{SixBirdsIII.scalarFinSum_scale}\newline \path{SixBirdsIII.scalarFinSum_congr}\newline \path{SixBirdsIII.scalarFinSum_zero}\newline \path{SixBirdsIII.scalarFinSum_nonneg}\newline \path{SixBirdsIII.scalarFinSum_mono}\newline \path{SixBirdsIII.scalarFinSum_ratCast} & theorem-grade & partial proof & Ordered log laws (real ln); positive kernels; conditional disintegration. The Lean decomposition quantifies over strictly positive rational kernels. The paper extends the identity to kernels with zero entries using extended-real KL; its lower bound and attainment criterion are proved in the paper.\tabularnewline
\path{thm:noncommuting-completions} & Finite idempotent completions need not commute. & Section 7 & \path{SixBirdsIII.noncommuting_completions} & theorem-grade & direct proof & Explicit finite noncommuting idempotents.\tabularnewline
\path{thm:completion-pasting-defect} & Completion-pasting defect is nonempty exactly when completions do not commute. & Section 7 & \path{SixBirdsIII.completion_pasting_defect} & theorem-grade & direct proof & Defect predicate for completion mismatch.\tabularnewline
\path{thm:idempotent-saturation} & An idempotent finite self-map saturates after one positive iterate. & Section 7 & \path{SixBirdsIII.idempotent_saturation} & theorem-grade & direct proof & Idempotent completion equations.\tabularnewline
\path{thm:strict-extension-nonfactorization} & Strict extension nonfactorization is equivalent to a fiber mismatch. & Section 7 & \path{SixBirdsIII.strict_extension_nonfactorization} & theorem-grade & direct proof & Factorization through the image.\tabularnewline
\path{thm:factorization-defect-equivalence} & The factorization defect records exactly nonfactorization witnesses. & Section 7 & \path{SixBirdsIII.factorization_defect_equivalence} & theorem-grade & direct proof & Explicit finite factorization defect.\tabularnewline
\path{thm:fixed-interface-definability-bound} & Finite-lens definability over a fixed interface has cardinality 2 to the image size. & Section 7 & \path{SixBirdsIII.fixed_interface_definability_bound} & theorem-grade & partial proof & Predicate bijection only; the numerical cardinality step is absent.\tabularnewline
\path{thm:forest-no-drive} & A finite forest has no first-cohomology drive. & Section 8 & \path{SixBirdsIII.forest_no_drive}\newline \path{SixBirdsIII.finite_graph_forest_no_drive} & theorem-grade & partial proof & Rooted parent/rank forests; beta-one-to-cohomology bridge absent.\tabularnewline
\path{thm:exact-form-null-drive} & Exact finite one-forms have zero integral around every cycle. & Section 8 & \path{SixBirdsIII.exact_form_null_drive}\newline \path{SixBirdsIII.finite_graph_exact_form_null_drive}\newline \path{SixBirdsIII.finiteWalkIntegral_coboundary} & theorem-grade & partial proof & Int-valued cochains and finite-walk telescoping only.\tabularnewline
\path{thm:nonzero-affinity-equivalence} & Nonzero cohomology class is equivalent to nonzero cycle affinity. & Section 8 & \path{SixBirdsIII.nonzero_affinity_equivalence}\newline \path{SixBirdsIII.finite_graph_nonzero_affinity_equivalence} & theorem-grade & partial proof & Connected finite graphs with Int-valued cochains.\tabularnewline
\path{thm:gating-affinity-suppression} & Gating to a forest suppresses cycle-supported drive. & Section 8 & \path{SixBirdsIII.gating_affinity_suppression}\newline \path{SixBirdsIII.finite_graph_gating_affinity_suppression} & theorem-grade & partial proof & Target-forest integrals only; gate data supplied; absence status unproved.\tabularnewline
\path{thm:promotion-gate-soundness} & Accepted-family promotion verdicts require every required core gate to pass. & Section 10 & \path{SixBirdsIII.promotion_gate_soundness} & theorem-grade & record projection & Predeclared Boolean core-gate flags.\tabularnewline
\path{thm:strict-promotion-nonfactorization} & Strict promotion implies a nonfactorization witness. & Section 10 & \path{SixBirdsIII.strict_promotion_nonfactorization}\newline \path{SixBirdsIII.promotion_strict_gate_iff_factorization_defect}\newline \path{SixBirdsIII.finite_strict_promotion_nonfactorization} & theorem-grade & direct proof under stated hypotheses & Finite maps and Boolean core gates; strictness computed from defect.\tabularnewline
\path{thm:nonstrict-promotion-factorization} & Non-strict promotion implies a factorization filler. & Section 10 & \path{SixBirdsIII.nonstrict_promotion_factorization}\newline \path{SixBirdsIII.promotion_nonstrict_gate_iff_factorization}\newline \path{SixBirdsIII.finite_nonstrict_promotion_factorization} & theorem-grade & direct proof under stated hypotheses & Finite maps and Boolean core gates; non-strictness computed from factorization.\tabularnewline
\path{thm:toy22-strict-bridge} & A finite two-layer Toy22 bridge witnesses strict promotion. & Section 10 & \path{SixBirdsIII.toy22_strict_bridge} & theorem-grade & direct proof & Explicit two-point bridge witness.\tabularnewline
\path{thm:no-overreading-suppression} & No accepted claim depends on unbridged suppressed content. & Section 11 & \path{SixBirdsIII.no_overreading_suppression} & theorem-grade & record projection & Suppressed-use Boolean classifier only.\tabularnewline
\path{thm:same-level-self-audit-failure} & Same-level complete self-audit does not receive accepted status. & Section 11 & \path{SixBirdsIII.same_level_self_audit_failure} & theorem-grade & record projection & Same-level Boolean classifier only.\tabularnewline
\path{thm:finite-rotating-audit} & A finite lower instrument stack admits a one-level audit extension. & Section 11 & \path{SixBirdsIII.finite_rotating_audit}\newline \path{SixBirdsIII.finite_rotating_audit_of_higher_level}\newline \path{SixBirdsIII.no_level_above_structural}\newline \path{SixBirdsIII.no_higher_level_for_top_stack} & theorem-grade & partial proof & Higher Level only; bridge, report, compliance, and nonfinality remain flags.\tabularnewline
\path{thm:ai-closure-as-packaging} & The abstract-interpretation map is a closure operator. & Section 12 & \path{SixBirdsIII.ai_closure_as_packaging} & theorem-grade & direct proof under stated hypotheses & Adjunction and order laws supplied by AIHost.\tabularnewline
\path{thm:ai-descent-as-sound-transformer} & Abstract interpretation soundness inequalities give a sound transformer. & Section 12 & \path{SixBirdsIII.ai_descent_as_sound_transformer} & theorem-grade & direct proof under stated hypotheses & Adjunction and monotonicity supplied by AIHost.\tabularnewline
\path{thm:ai-representability-as-fixedness} & Representability is fixedness under the abstract-interpretation closure. & Section 12 & \path{SixBirdsIII.ai_representability_as_fixedness} & theorem-grade & record projection & Forward direction assumes gamma\_fixed in AIHost.\tabularnewline
\path{model:ai-model-family} & Finite abstract-interpretation family realizes the named role fragment. & Section 12 & \path{SixBirdsIII.ai_model_family} & model-grade & record projection & Role proofs supplied by AIModelFamily.\tabularnewline
\path{thm:descent-square-recovery} & Decorated descent square exactness recovers descent commutation. & Section 14 & \path{SixBirdsIII.descent_square_recovery} & theorem-grade & partial proof & Image-descent equivalence only; decorated squares absent.\tabularnewline
\path{thm:completion-pasting-recovery} & Decorated completion square exactness recovers completion commutation. & Section 14 & \path{SixBirdsIII.completion_pasting_recovery} & theorem-grade & partial proof & Commutation-defect equivalence only; decorated squares absent.\tabularnewline
\path{thm:strict-extension-as-no-filler} & Strict extension is equivalent to absence of an exact filler. & Section 14 & \path{SixBirdsIII.strict_extension_as_no_filler} & theorem-grade & partial proof & Factorization-defect equivalence only; decorated fillers absent.\tabularnewline
\path{model:pica-model-realization} & Finite PICA realizes the cell, pair, and provenance fragment. & Section 13 & \path{SixBirdsIII.pica_model_realization} & model-grade & record projection & Success record supplied without PICA store semantics.\tabularnewline
\path{model:cantor-strict-bridge-realization} & Audited Cantor shell realizes the strict-bridge fragment. & Section 13 & \path{SixBirdsIII.cantor_strict_bridge_realization} & model-grade & record projection & Verdict supplied without concrete Cantor-shell maps.\tabularnewline
\path{prop:mechanization-target-feasibility} & Dated strict validation run of the Lean artifact. & Appendix F & \path{SixBirdsIII.mechanization_target_feasibility} & audit-grade & record projection & Preset validation flags; the run is reported in Proposition 35.\tabularnewline
\path{prop:top-down-channel-gate-soundness} & Accepted top-down channel records pass every required top-down gate. & Subsection 4.7 & \path{SixBirdsIII.top_down_channel_gate_soundness} & proposition-grade & record projection & Boolean gate flags; no intervention or effect data.\tabularnewline
\path{prop:top-down-channel-implies-struct-down} & Accepted top-down channel records supply structural downward influence. & Subsection 4.7 & \path{SixBirdsIII.top_down_channel_implies_struct_down} & proposition-grade & record projection & Accepted Boolean status includes the structural-path flag.\tabularnewline
\path{prop:struct-down-insufficient-for-top-down-channel} & A structural downward path can exist while the top-down channel claim is not accepted. & Subsection 4.7 & \path{SixBirdsIII.structural_downward_influence_not_top_down_channel} & proposition-grade & direct proof & Explicit schema record with structural path and blocked status.\tabularnewline
\bottomrule
\end{longtable}
\endgroup
\end{landscape}
 
\subsection{Reproducibility and Trust Base}\label{sn:F.3}

The Lean artifact is the repository's finite validation track for the manifest scope listed in the table above. It is checked under the pinned toolchain recorded in the repository, and the strict validation command checks the build, the required manifest entries, declaration existence, declaration kind, theorem axiom dependencies, placeholder scans, and semantic probes.

The allowed non-project axioms for theorem audits are:

{[}
\texttt{Classical.choice},\quad
\texttt{propext},\quad
\texttt{Quot.sound},\quad
\texttt{choice},\quad
\texttt{funext}.
{]}

The first three are Lean 4's standard axioms; the entries \texttt{choice} and \texttt{funext} name no further axiom (\texttt{funext} is a theorem derived from \texttt{Quot.sound}). All theorem-backed manifest entries are audited against this trust base. The trust-base statement is a proof-assistant audit statement; it is not an additional mathematical theorem about the calculus.

\subsection{Compressed Module Map}\label{sn:F.4}

The root import collects the finite support for primitive labels, levels, hosts, status families, record schemas, finite-map and completion theorems, graph/cohomology theorems, promotion and instrument theorems, abstract-interpretation theorems, high-structure recovery theorems, model-realization records, top-down channel records, semantic probes, and the validation-harness record.

This module map is an audit map. It records where the finite support for each manifest entry lives; it does not certify statements outside the manifest.

\subsection{Proposition 35: Mechanization Target Feasibility}\label{sn:F.5}

\begin{claimproposition}{Proposition 35 (MechanizationTargetFeasibility)}\phantomsection\label{cl:prop:35}

On 2026-09-27, at revision \texttt{c0ad1eb} of the repository \url{https://github.com/ioannist/six-birds-foundations-iii}, with the Lean toolchain \texttt{leanprover/lean4:v4.28.0}, the strict validation command \texttt{scripts/validate\_lean.sh --tier full --require-full} reported the following for the manifest of Subsection \hyperref[sn:F.2]{F.2}: the project builds; all 24 definitions and all 38 required claims of the manifest (the 35 numbered results and the three unnumbered propositions of Subsection \hyperref[sn:4.7]{4.7}) have a Lean declaration; all 96 primary and supporting declarations exist and have the declaration kind the manifest expects; all 72 axiom audits of theorem declarations stay within the trust base of Subsection \hyperref[sn:F.3]{F.3}; no placeholder markers and no semantic-audit findings are present; and the validator's overall verdict is \texttt{full}: every manifest row has a checked declaration of the expected kind. Of the 38 claims, 9 are direct proofs, 4 direct proofs under stated hypotheses, 13 partial proofs and 12 record projections (Subsection \hyperref[sn:F.2]{F.2}). The verdict certifies that each manifest entry is represented by a checked Lean declaration. What each declaration proves is given by its coverage class and scope in the table of Subsection \hyperref[sn:F.2]{F.2}.

\end{claimproposition}

\begin{proof}

The validation artifact consists of the finite manifest, the Lean declarations it names, the pinned toolchain, the trust-base list, and the validator. The validator checks that the project builds under the pinned toolchain, that each required manifest row is present, that each named declaration exists and has the expected declaration kind, that the non-project axiom dependencies of each theorem declaration lie in the trust base, that the sources contain no placeholder markers, and that the semantic probe module exposes the required finite checks. The run of the statement passed each of these finite checks, with the counts stated. The Lean declaration \texttt{SixBirdsIII.mechanization\_target\_feasibility} records the categories of these checks as a finite record (a record projection in the table of Subsection \hyperref[sn:F.2]{F.2}); the run itself is the evidence for the statement.

\end{proof}

\subsection{Mechanization Nonclaims}\label{sn:F.6}

The validation artifact is manifest-scoped. It does not assert proof-assistant certification for unlisted statements, empirical correctness of any concrete implementation, a universal exact-six theorem, a universal six-symbol algebra, a unique decomposition theorem, complete model coverage, complete same-level self-certification, or host extensions outside the finite audited manifest scope. For a result whose coverage class is record projection, the Lean declaration does not certify the mathematics of the result; the paper's proof does.
 \section{Dependency Graph}\label{sn:G}

This appendix records the long-form dependency graph of \(\mathbf{BirdInt}^{\mathrm{aud}}_{\mathrm{fin}}\). Subsection \hyperref[sn:G.1]{G.1} records the definition nodes. Subsection \hyperref[sn:G.2]{G.2} records the theorem build order in topological form. Subsection \hyperref[sn:G.3]{G.3} records the critical proof paths.

The dependency graph is a finite typed record of the producer/consumer relations among definitions, lemmas, theorems, countermodels, model-realization records, and high-structure recovery theorems in the paper. It is a bookkeeping appendix, not an additional theorem.

\subsection{Definition Nodes}\label{sn:G.1}

The definition nodes of \(\mathbf{BirdInt}^{\mathrm{aud}}_{\mathrm{fin}}\) are the finite typed schemas that the calculus introduces and uses throughout the paper. Each node is a finite typed record schema; the dependency graph records which schemas consume the records of which other schemas as part of their admissibility predicate.

The eleven definition-node families are:

\begin{enumerate}
\def\labelenumi{\arabic{enumi}.}
\tightlist
\item
  \textbf{Host} --- finite host schema (Subsection \hyperref[sn:2.1]{2.1}).
\item
  \textbf{Content} --- finite content universe with visibility and suppression records (Subsections \hyperref[sn:3.2]{3.2} and \hyperref[sn:3.5]{3.5}).
\item
  \textbf{Theory packages} --- \(\mathcal T=(Z,f,\Sigma_{f},E,\mathcal A)\) (Subsection \hyperref[sn:3.3]{3.3}).
\item
  \textbf{Instruments} --- instrument records \(I\) (Subsection \hyperref[sn:3.4]{3.4} and Subsubsection \hyperref[sn:11.1.1]{11.1.1}).
\item
  \textbf{Profiles} --- typed profile \(\lambda\) (Subsection \hyperref[sn:3.6]{3.6} and Appendix \hyperref[sn:A.4]{A.4}).
\item
  \textbf{Witness/update maps} --- \(\mathsf{Wit}:\mathbb P\to\mathrm{Type}\) and \(\mathsf{Upd}:\mathbb P\to\mathrm{Type}\) (Subsection \hyperref[sn:3.7]{3.7} and Tables A.2, A.3).
\item
  \textbf{Statuses} --- seven typed status families (Subsection \hyperref[sn:5.1]{5.1} and Table A.5).
\item
  \textbf{Defects} --- typed defect families per primitive and per promotion (Subsections \hyperref[sn:5.2]{5.2}--\hyperref[sn:5.5]{5.5} and Table A.6).
\item
  \textbf{Gates} --- promotion gates and bridge-admissibility gates (Subsection \hyperref[sn:10.2]{10.2}).
\item
  \textbf{Audits} --- host audit, claim audit, bridge audit, rotating-audit (Subsubsection \hyperref[sn:4.9.2]{4.9.2} and Section \hyperref[sn:11]{11}).
\item
  \textbf{Model-realization maps} --- realization maps \(\mathrm{real}_{\mathcal M\to\mathcal S}\) (Subsubsection \hyperref[sn:13.1.1]{13.1.1}).
\end{enumerate}

The producer/consumer relations among these eleven node families are recorded as a directed dependency graph: hosts support content and packages; packages, instruments, profiles, witness/update maps, visibility records, thresholds, audits, and defects feed the judgment schemas; classifiers consume judgment and defect data to assign statuses; gates feed promotion-status verdicts; audits feed visibility, source, and instrument-bound claim verdicts; model-realization maps consume the host, package, judgment, status, and audit nodes to produce model-grade \(\texttt{model\_realization}\) verdicts on declared fragments.

The dependency graph is a finite typed record on the eleven node families; it is admissible as a bookkeeping object when every node is finite and every dependency edge is recorded, including cross-stage edges.

\subsection{Theorem Build Order}\label{sn:G.2}

The nine stages list the paper's order of exposition. This is not a topological ordering of the full proof graph: Theorem \hyperref[cl:thm:5]{5} cites the later countermodel \(\mathsf{CM}_4\) (Theorem \hyperref[cl:thm:1]{1} builds its two records in Subsubsection \hyperref[sn:6.1.1]{6.1.1}, and \(\mathsf{CM}_7\) restates them), and Theorem \hyperref[cl:thm:5]{5} uses Theorems \hyperref[cl:thm:22]{22}, \hyperref[cl:thm:24]{24} and \hyperref[cl:thm:25]{25}. These cross-stage edges are listed at the end of this subsection.

The nine-stage build order is:

\begin{enumerate}
\def\labelenumi{\arabic{enumi}.}
\tightlist
\item
  \textbf{Formal host and records.} Sections \hyperref[sn:2]{2}--\hyperref[sn:5]{5}: finite host mathematics, primitive labels, claim strengths, status discipline, theory packages, instruments, source-of-truth records, audit records, visibility/suppression records, directed-cell records, pair-observable records, promotion bridges, profiles, status families, and defect schemas.
\item
  \textbf{Finite map/completion lemmas.} Sections \hyperref[sn:6]{6} and \hyperref[sn:7]{7}: finite-cell admissibility predicates, effective-image representative arguments, idempotent-iteration arguments, and finite fiber-mismatch arguments.
\item
  \textbf{Core finite theorem spine.} Section \hyperref[sn:6]{6} (Theorems \hyperref[cl:thm:1]{1}--\hyperref[cl:thm:5]{5}) and Section \hyperref[sn:7]{7} (Theorems \hyperref[cl:thm:6]{6}--\hyperref[cl:thm:14]{14}): no-total-algebra, typed non-collapse, well-formedness decidability, covered-cell status uniqueness, status-family separation, descent through quotient, descent defect equivalence, Markov closure deficit, noncommuting completions, completion-pasting defect, idempotent saturation, strict-extension nonfactorization, factorization defect equivalence, fixed-interface definability bound.
\item
  \textbf{No-go and countermodels.} Section \hyperref[sn:8]{8} (Theorems \hyperref[cl:thm:15]{15}--\hyperref[cl:thm:18]{18}) and Section \hyperref[sn:9]{9} plus Appendix \hyperref[sn:C]{C} (\(\mathsf{CM}_{1}\)--\(\mathsf{CM}_{12}\)): forest no-drive, exact-form null-drive, nonzero-affinity equivalence, gating-affinity suppression, plus the twelve finite countermodels.
\item
  \textbf{Promotion.} Section \hyperref[sn:10]{10} (Theorems \hyperref[cl:thm:19]{19}--\hyperref[cl:thm:22]{22}): promotion gate soundness, strict promotion implies nonfactorization, non-strict promotion implies factorization, Toy22 strict-bridge example.
\item
  \textbf{Instruments and self-reference.} Section \hyperref[sn:11]{11} (Theorems \hyperref[cl:thm:23]{23}--\hyperref[cl:thm:25]{25}): no-overreading suppression, same-level self-audit failure, finite rotating-audit theorem.
\item
  \textbf{Model families.} Section \hyperref[sn:12]{12} (Theorems \hyperref[cl:thm:26]{26}--\hyperref[cl:thm:28]{28} and Model Theorem \hyperref[cl:thm:29]{29}) and Section \hyperref[sn:13]{13} (Model Theorems \hyperref[cl:thm:33]{33}, \hyperref[cl:thm:34]{34}, plus the graph/cohomology realization): closure as packaging, descent as sound transformer, representability as fixedness, AI model-family theorem, PICA realization, Cantor strict-bridge realization, graph/cohomology realization. The realization statements in this stage are model-grade, not theorem-grade.
\item
  \textbf{High-structure semantics.} Section \hyperref[sn:14]{14} (Theorems \hyperref[cl:thm:30]{30}--\hyperref[cl:thm:32]{32}): descent-square recovery, completion-pasting recovery, strict extension as no-filler.
\item
  \textbf{Mechanization disclosure appendix.} Appendix \hyperref[sn:F]{F}: paper-to-Lean disclosure table, validation command, toolchain, trust base, Proposition \hyperref[cl:prop:35]{35}, and manifest-scoped mechanization nonclaims.
\end{enumerate}

The nine stages give the order of exposition, not a topological ordering of every dependency. The edges of the theorem spine that point from an earlier stage to a later one are:

\begin{longtable}[]{@{}lll@{}}
\toprule
\begin{minipage}[b]{0.30\columnwidth}\raggedright
consumer (stage)\strut
\end{minipage} & \begin{minipage}[b]{0.30\columnwidth}\raggedright
producer (stage)\strut
\end{minipage} & \begin{minipage}[b]{0.30\columnwidth}\raggedright
what is used\strut
\end{minipage}\tabularnewline
\midrule
\endhead
\begin{minipage}[t]{0.30\columnwidth}\raggedright
Theorem \hyperref[cl:thm:1]{1} (3)\strut
\end{minipage} & \begin{minipage}[t]{0.30\columnwidth}\raggedright
\(\mathsf{CM}_7\) (4)\strut
\end{minipage} & \begin{minipage}[t]{0.30\columnwidth}\raggedright
none: \(\mathsf{CM}_7\) restates the two records built in Subsubsection \hyperref[sn:6.1.1]{6.1.1}, so this entry is a cross-reference, not a dependency\strut
\end{minipage}\tabularnewline
\begin{minipage}[t]{0.30\columnwidth}\raggedright
Theorem \hyperref[cl:thm:5]{5} (3)\strut
\end{minipage} & \begin{minipage}[t]{0.30\columnwidth}\raggedright
\(\mathsf{CM}_4\) (4)\strut
\end{minipage} & \begin{minipage}[t]{0.30\columnwidth}\raggedright
the action cell and blocked pair observable of case (b)\strut
\end{minipage}\tabularnewline
\begin{minipage}[t]{0.30\columnwidth}\raggedright
Theorem \hyperref[cl:thm:5]{5} (3)\strut
\end{minipage} & \begin{minipage}[t]{0.30\columnwidth}\raggedright
Theorem \hyperref[cl:thm:22]{22} (5)\strut
\end{minipage} & \begin{minipage}[t]{0.30\columnwidth}\raggedright
the strict-bridge instance of cases (c), (d) and (e)\strut
\end{minipage}\tabularnewline
\begin{minipage}[t]{0.30\columnwidth}\raggedright
Theorem \hyperref[cl:thm:5]{5} (3)\strut
\end{minipage} & \begin{minipage}[t]{0.30\columnwidth}\raggedright
Theorem \hyperref[cl:thm:24]{24} (6)\strut
\end{minipage} & \begin{minipage}[t]{0.30\columnwidth}\raggedright
non-acceptance of \(\mathrm{Sound}(I_{\mathrm{prom}})\) and \(\mathrm{Sound}(I^{+}_{\mathrm{prom}})\)\strut
\end{minipage}\tabularnewline
\begin{minipage}[t]{0.30\columnwidth}\raggedright
Theorem \hyperref[cl:thm:5]{5} (3)\strut
\end{minipage} & \begin{minipage}[t]{0.30\columnwidth}\raggedright
Theorem \hyperref[cl:thm:25]{25} (6)\strut
\end{minipage} & \begin{minipage}[t]{0.30\columnwidth}\raggedright
lifted-soundness acceptance and rejection in case (e)\strut
\end{minipage}\tabularnewline
\begin{minipage}[t]{0.30\columnwidth}\raggedright
promotion classifier, Subsubsection \hyperref[sn:5.6.4]{5.6.4} (1)\strut
\end{minipage} & \begin{minipage}[t]{0.30\columnwidth}\raggedright
gates, Subsection \hyperref[sn:10.2]{10.2} (5)\strut
\end{minipage} & \begin{minipage}[t]{0.30\columnwidth}\raggedright
the gate record \(\mathsf{GateResults}(B)\)\strut
\end{minipage}\tabularnewline
\bottomrule
\end{longtable}

\subsection{Critical Proof Paths}\label{sn:G.3}

This subsection records the five critical proof paths in the dependency graph. Each path is the chain of definitions and theorems that supports a specific theorem-grade or model-grade conclusion.

\subsubsection{Strict Promotion Path}\label{sn:G.3.1}

The strict promotion path is the chain of records and theorems supporting the strict-promotion verdict \(\mathrm{Promote}(B):\texttt{strict}\):

Theory packages and promotion bridges supply the gate records of Subsection \hyperref[sn:10.2]{10.2}, which the promotion classifier of Subsubsection \hyperref[sn:5.6.4]{5.6.4} uses to assign a strict verdict. Theorem \hyperref[cl:thm:19]{19} derives core-gate soundness from accepted-family verdicts. Separately, Theorem \hyperref[cl:thm:12]{12} supplies the finite factorization criterion, Theorem \hyperref[cl:thm:13]{13} identifies nonfactorization with a nonempty \(\Delta_{\mathrm{fact}}\), and Theorem \hyperref[cl:thm:20]{20} combines Theorem \hyperref[cl:thm:13]{13} with the strict classifier rule. Theorem \hyperref[cl:thm:20]{20} does not use Theorem \hyperref[cl:thm:19]{19}.

The path is consumed by the Cantor strict-bridge realization (Model Theorem \hyperref[cl:thm:34]{34}) and, through Theorems \hyperref[cl:thm:12]{12} and \hyperref[cl:thm:13]{13}, by the strict-extension recovery in the high-structure fragment (Theorem \hyperref[cl:thm:32]{32}).

\subsubsection{Package versus Closure Path}\label{sn:G.3.2}

The package-versus-closure path is the chain supporting the \(P_{5}\)-packaging vs macro-closure separation:

theory packages \(\to\) quotient/package map \(q\) and update \(F\) \(\to\) effective-image descent setup \(\to\) Theorem \hyperref[cl:thm:6]{6} (descent through quotient) \(\to\) Theorem \hyperref[cl:thm:7]{7} (descent defect equivalence) \(\to\) \(\mathsf{CM}_{1}\) (active package without macro closure).

The path is consumed by package and closure nonclaims throughout the paper, including the AI host realization (Theorem \hyperref[cl:thm:26]{26} closure as packaging, paired with Theorem \hyperref[cl:thm:27]{27} descent as sound transformer) and the Cantor strict-bridge realization (Model Theorem \hyperref[cl:thm:34]{34}), where macro-dynamics closure is \emph{not} asserted without the required descent data.

\subsubsection{Completion versus Route Coherence Path}\label{sn:G.3.3}

The completion-versus-route path is the chain supporting the completion-noncommutation vs \(P_{3}\)-pasting recovery:

The finite completions \(E_1,E_2\) give the noncommuting instance of Theorem \hyperref[cl:thm:9]{9}, also recorded as \(\mathsf{CM}_8\). Theorem \hyperref[cl:thm:10]{10} independently identifies noncommutation with a nonempty completion-pasting defect; Theorem \hyperref[cl:thm:31]{31} uses that defect in the decorated-square recovery. \(\mathsf{CM}_8\) is an instance, not a premise of Theorem \hyperref[cl:thm:10]{10}.

The path is consumed by Theorem \hyperref[cl:thm:31]{31} of Section \hyperref[sn:14]{14} in Subsection \hyperref[sn:14.3]{14.3}.

\subsubsection{Self-Reference Honesty Path}\label{sn:G.3.4}

The self-reference honesty path is the chain supporting the no-overreading and same-level self-audit failure theorems:

Instrument records, visibility and suppression rules, and classifier priority support Theorem \hyperref[cl:thm:23]{23}. The claim classifier's self-reference rule separately supports Theorem \hyperref[cl:thm:24]{24}. The rotating-audit chain and Theorem \hyperref[cl:thm:24]{24} then support Theorem \hyperref[cl:thm:25]{25}, including its non-finality clause.

The path is consumed by every model-realization sheet of Section \hyperref[sn:13]{13} (whose claim-grade verdicts are scoped under instruments) and by the manifest-scope mechanization boundary of Appendix \hyperref[sn:F]{F}.

\subsubsection{Drive Separation Path}\label{sn:G.3.5}

The drive separation path is the chain supporting the cohomological-drive vs holonomy separation:

finite graph \(G\) and cycle space \(\mathsf{Cycles}(G)\) \(\to\) first cohomology \(H^{1}(G;k)\) \(\to\) Theorem \hyperref[cl:thm:15]{15} (forest no-drive), Theorem \hyperref[cl:thm:16]{16} (exact-form null-drive), and Theorem \hyperref[cl:thm:17]{17} (nonzero-affinity equivalence) \(\to\) \(\mathsf{CM}_{2}\) (active route/holonomy without drive); and, on the same graph data, Theorem \hyperref[cl:thm:18]{18} (gating-affinity suppression) \(\to\) \(\mathsf{CM}_{12}\) (gating destroys cycle-supported drive).

The path is consumed by the graph/cohomology realization (Subsection \hyperref[sn:13.4]{13.4}) and by the cohomology nonclaim register \(\mathcal N_{\mathrm{coh}}\) of Subsubsection \hyperref[sn:13.4.4]{13.4.4}.

The five critical paths are the principal supporting chains for the paper's theorem-grade and model-grade conclusions. They do not exhaust the dependency graph; they record the load-bearing paths that the paper relies on most heavily.
 \section{Deferred and Excluded Claims}\label{sn:H}

This appendix records the deferred and excluded claims of \(\mathbf{BirdInt}^{\mathrm{aud}}_{\mathrm{fin}}\). Subsection \hyperref[sn:H.1]{H.1} records the deferred claims, organized by future-paper group. Subsection \hyperref[sn:H.2]{H.2} records the excluded claims, which lie outside the calculus's positive scope.

The appendix is a finite typed register that consolidates the future-work items of Subsection \hyperref[sn:15.3]{15.3} and the negative-scope items of Subsection \hyperref[sn:15.2]{15.2} (and the per-section nonclaim subsections) in one place. It does not introduce new theorems; it records what is \emph{not} in this paper and \emph{not} a target of this paper.

\subsection{Deferred Claims}\label{sn:H.1}

This subsection records the deferred claims by future-paper group, mirroring the six classes of future work in Subsection \hyperref[sn:15.3]{15.3}.

\subsubsection{Classifier and Dependency Completeness}\label{sn:H.1.1}

Deferred items:

\begin{itemize}
\tightlist
\item
  a complete status-classifier theorem characterizing the priority order over \(\mathrm{ClaimStatus}\);
\item
  a defect-to-status soundness theorem characterizing the routing from the defect families of Subsections \hyperref[sn:5.2]{5.2}--\hyperref[sn:5.5]{5.5} through the classifier rules of Subsection \hyperref[sn:5.6]{5.6};
\item
  a dependency-lattice theorem organizing the per-component admissibility predicates of Sections \hyperref[sn:4]{4} and \hyperref[sn:5]{5} into a lattice with a finite reduction theorem.
\end{itemize}

\subsubsection{Promotion and Stacking Theory}\label{sn:H.1.2}

Deferred items:

\begin{itemize}
\tightlist
\item
  soundness and completeness conditions for the promotion classifier of Subsubsection \hyperref[sn:5.6.4]{5.6.4} over \(\mathrm{PromotionStatus}\), relative to an external account of correct promotion;
\item
  a stacking composition theorem for composed bridges whose middle interface refines the source, extending the object-map result of Subsubsection \hyperref[sn:10.7.1]{10.7.1} to the joint gate and nonclaim semantics of the composed bridge;
\item
  a local-to-global promotion theorem identifying conditions under which a finite family of locally admissible promotion bridges supports a global promotion verdict.
\end{itemize}

\subsubsection{Verified Model Realizations}\label{sn:H.1.3}

Deferred items:

\begin{itemize}
\tightlist
\item
  a deeper PICA realization audit, extending Model Theorem \hyperref[cl:thm:33]{33} of Subsection \hyperref[sn:13.2]{13.2} toward implementation-correctness records for the actor-update map, informant-witness map, status recovery, and pair-observable realness predicate;
\item
  a deeper Cantor strict-bridge realization audit, extending Model Theorem \hyperref[cl:thm:34]{34} of Subsection \hyperref[sn:13.3]{13.3} toward implementation-correctness records for the shell-bound strictness witness and the gate-result table;
\item
  a deeper abstract-interpretation realization audit, extending Model Theorem \hyperref[cl:thm:29]{29} of Section \hyperref[sn:12]{12} beyond the manifest-scoped Lean support toward richer concrete domains and transformers;
\item
  an implementation-correctness layer for any concrete PICA implementation, audited under a separate proof-assistant-bound instrument.
\end{itemize}

\subsubsection{Cantor Shell Expansion}\label{sn:H.1.4}

Deferred items:

\begin{itemize}
\tightlist
\item
  a full Cantor strictness proof extending the audited-shell strict-bridge realization beyond \(\mathcal S_{\mathrm{aud}}\) via an admissible expansion bridge;
\item
  a pressure-gap-to-KL bridge recording an admissible bridge from the pressure-gap consequence defect of Subsubsection \hyperref[sn:13.3.6]{13.3.6} to the Kullback--Leibler closure-deficit framework;
\item
  a stratumwise root-separation theorem characterizing the persistent strata family \(\mathcal F\) via a separation property stronger than the conditional pressure disintegration of Subsubsection \hyperref[sn:13.3.6]{13.3.6}.
\end{itemize}

\subsubsection{Drive and High-Structure Semantics}\label{sn:H.1.5}

Deferred items:

\begin{itemize}
\tightlist
\item
  a general drive classification extending the cohomological drive certificate \([a]\neq 0\) of Subsubsection \hyperref[sn:13.4.3]{13.4.3} to a finite family of drive certificates over multiple drive notions;
\item
  a holonomy-to-drive bridge \(B_{3\to 6}^{\mathrm{drive}}\) from a \(P_{3}\) route/holonomy witness to a \(\mathrm{P6}_{\mathrm{drive}}\) certificate;
\item
  gating beyond forests, characterizing \(P_{2}\)-gating effects on drive on graphs with non-trivial cycle structure under additional gating policies;
\item
  weak double-category coherence extending the strict pasting law of Subsubsection \hyperref[sn:14.1.3]{14.1.3} to weak coherence relations on finite decorated squares;
\item
  a general obstruction propagation theorem characterizing the propagation of square defects under finite strict pasting, in the sense of Subsubsection \hyperref[sn:14.5.3]{14.5.3}.
\end{itemize}

\subsubsection{Self-Reference and Instrument Transfer}\label{sn:H.1.6}

Deferred items:

\begin{itemize}
\tightlist
\item
  a full instrument-transfer classifier fixing the per-claim-type mapping from source-side acceptance and bridge defect to target-side \(\chi\in\mathrm{ClaimStatus}\) under \(J\), in the sense of Subsubsection \hyperref[sn:11.9.2]{11.9.2};
\item
  an infinite rotating-audit semantics extending the finite rotating-audit chain of Subsection \hyperref[sn:11.7]{11.7} to an inverse limit or coalgebraic semantics admissible under a stronger meta-instrument;
\item
  a complete self-reference classification extending Theorem \hyperref[cl:thm:24]{24} of Subsection \hyperref[sn:11.5]{11.5} to a comprehensive classification of same-level, level-shifted, and rotating self-reference patterns.
\end{itemize}

The six groups above record the principal future-work items of the paper. Each item is a scoped pointer to an admissible extension target, in keeping with the future-work boundary of Subsection \hyperref[sn:15.3]{15.3}.

\subsection{Excluded Claims}\label{sn:H.2}

This subsection records claims that lie outside the calculus's positive scope. Each item below is \emph{not} deferred to future work but excluded as a positive claim of \(\mathbf{BirdInt}^{\mathrm{aud}}_{\mathrm{fin}}\).

The seven excluded claims are:

\begin{enumerate}
\def\labelenumi{\arabic{enumi}.}
\tightlist
\item
  \textbf{Universal six-symbol algebra.} No total operation \(\ast:\mathbb P^{2}\to\mathbb P\) recovers cell status across profiles. Theorem \hyperref[cl:thm:1]{1} of Section \hyperref[sn:6]{6} records the formal counterexample, and \(\mathrm{NC}\text{-}1\) of Subsection \hyperref[sn:15.2]{15.2} records the corresponding nonclaim.
\item
  \textbf{All mathematics is six birds.} The claim that every mathematical structure decomposes into the six primitives is excluded; \(\mathbf{BirdInt}^{\mathrm{aud}}_{\mathrm{fin}}\) is a finite typed audited interaction calculus, not a foundation for all mathematics. \(\mathrm{NC}\text{-}2\) of Subsection \hyperref[sn:15.2]{15.2} records the corresponding nonclaim.
\item
  \textbf{Unrestricted exact-six.} The unrestricted decomposition claim, that every typed interaction is exactly representable as one of the thirty-six directed cells, is excluded. The exact-six structure is a finite typed schema with admissibility conditions, not an unrestricted decomposition theorem. The excluded claim is recorded in Subsubsection \hyperref[sn:2.4.3]{2.4.3} and the corresponding nonclaim is the unnumbered entry \emph{No unrestricted exact-six theorem} of Subsubsection \hyperref[sn:15.2.1]{15.2.1}.
\item
  \textbf{PICA as canonical algebra.} The claim that PICA defines the canonical \(6\times 6\) algebra of the calculus is excluded. PICA is a finite stochastic model realization, scoped to the cell/pair/provenance fragment, and the recovered \(6\times 6\) multiplicity count is a profile fact, not a universal law. PICA nonclaims 1 and 2 of Subsubsection \hyperref[sn:13.2.10]{13.2.10} and \(\mathrm{NC}\text{-}20\), \(\mathrm{NC}\text{-}21\) of Subsection \hyperref[sn:15.2]{15.2} record the corresponding nonclaims.
\item
  \textbf{Cantor pressure gap as unbridged strictness witness.} The claim that the pressure-gap defect \(\Delta_{T_{0}\to T_{1}}\) is itself the strictness witness is excluded. The strictness witness is the factorization defect \(\Delta_{\mathrm{fact}}^{\mathrm{Cantor}}\), and the pressure-gap defect is a separately recorded consequence record whose values are cited from the Cantor paper, not derived here from the witness. Cantor nonclaim 7 of Subsubsection \hyperref[sn:13.3.9]{13.3.9} and \(\mathrm{NC}\text{-}27\) of Subsection \hyperref[sn:15.2]{15.2} record the corresponding nonclaim.
\item
  \textbf{Provenance as proof.} The claim that provenance records constitute proof of any underlying mathematical fact is excluded. Provenance records are finite typed records of producer, source, audit, and visibility data; they support pair-realness and source-of-truth predicates but are not theorem-grade proofs of mathematical claims. The audit nonclaim \(\mathrm{NC}\text{-}15\) of Subsection \hyperref[sn:15.2]{15.2} (no audit-is-truth) records the corresponding boundary.
\item
  \textbf{High-structure semantics as replacement calculus.} The claim that the finite decorated double-category fragment of Section \hyperref[sn:14]{14} replaces or supersedes the core calculus is excluded. The decorated-square fragment is a secondary organizational layer; it records three specific recoveries (Theorems \hyperref[cl:thm:30]{30}--\hyperref[cl:thm:32]{32}) and not a general claim that every theorem-grade fact reduces to a decorated-square fact. \(\mathrm{NC}\text{-}34\), \(\mathrm{NC}\text{-}35\) of Subsection \hyperref[sn:15.2]{15.2} record the corresponding nonclaims.
\end{enumerate}

The seven excluded claims above are outside the positive scope of the paper. Each is recorded as a negative scope item in Section \hyperref[sn:15]{15}, and each is supported by a theorem, a countermodel, or a per-realization nonclaim register elsewhere in the paper. The appendix records them in one place to make the paper's negative scope register explicit.

The seven excluded claims and the six groups of deferred claims in Subsection \hyperref[sn:H.1]{H.1} together record the boundary of \(\mathbf{BirdInt}^{\mathrm{aud}}_{\mathrm{fin}}\): the calculus admits the theorem-grade, model-grade, and audit-grade results recorded in this paper; defers the items of Subsection \hyperref[sn:H.1]{H.1} to future work; and excludes the items of Subsection \hyperref[sn:H.2]{H.2} from any positive claim.
 \section{Claim Strengths and Instrument-Indexed Forms}\label{sn:I}

The body states and proves each numbered result in the mathematical vocabulary of its section. The calculus also records each result as an instrument-indexed judgment \(\Gamma\,;\;\mathsf H\,;\;I\;\vdash\;X\;:\;v\) under the host \(\mathsf H\) and instrument \(I\) of its section, in the instrument-relative judgment form of Subsubsection \hyperref[sn:3.4.2]{3.4.2}. This appendix lists the claim strength of each numbered result (Subsection \hyperref[sn:I.1]{I.1}) and its instrument-indexed form (Subsection \hyperref[sn:I.2]{I.2}).

\subsection{Claim Strengths}\label{sn:I.1}

The claim strengths of Subsubsection \hyperref[sn:2.3.1]{2.3.1} have the following meanings.

\begingroup
\setlength{\tabcolsep}{4pt}
\renewcommand{\arraystretch}{1.08}
\begin{center}
\begin{tabularx}{0.96\linewidth}{@{}>{\raggedright\arraybackslash}p{0.24\linewidth}>{\raggedright\arraybackslash}X@{}}
\toprule
Strength & Meaning \\
\midrule
theorem-grade & A statement with a proof in the calculus, stated under explicit host, instrument, package, profile, visibility, threshold, audit, and nonclaim assumptions. \\
proposition-grade & A statement that follows directly from definitions and prior theorems and is proved here. \\
schema-grade & A definition, judgment form, classifier specification, gate list, or discipline rule. Not a standalone theorem; later theorems may quote it. \\
model-only & A statement valid only as a scoped realization or application in a named model host. Not a universal theorem of the calculus. \\
countermodel-grade & A finite construction together with the no-go statement it supports. The construction and the no-go are both required. \\
future-work & A statement that is plausible from the calculus but not proved in this paper. Stated only in the limitations and future-work section. \\
excluded & A statement that this paper does not assert as a theorem, model realization, schema, or future-work claim. Recorded in Appendix \hyperref[sn:H.2]{H.2} and in the limitations discussion, with the theorem, countermodel, or nonclaim that supports the exclusion. \\
\bottomrule
\end{tabularx}
\end{center}
\endgroup

A claim strength scopes a statement; it does not rank it. Subsection \hyperref[sn:I.2]{I.2} uses model-grade, the synonym of model-only for a model theorem, and three refined tags: audit-grade for the validation record of Proposition \hyperref[cl:prop:35]{35}, theorem-grade with assumptions for Theorem \hyperref[cl:thm:8]{8}, whose statement carries finite probabilistic assumptions, and theorem-grade example for Theorem \hyperref[cl:thm:22]{22}, which is an explicit instance.

\subsection{Instrument-Indexed Forms of the Numbered Results}\label{sn:I.2}

For a theorem the verdict is \(\texttt{accepted}\). For a model theorem the verdict is \(\texttt{model\_realization}\): it records that the host supplies an admissible model-realization map for the declared fragment, in the sense of the model-realization convention of Subsection \hyperref[sn:13.1]{13.1}, and it is not the \(\texttt{accepted}\) verdict of a theorem-grade claim. For Proposition \hyperref[cl:prop:35]{35} the verdict is \(\texttt{audit\_report}\). Each entry gives the result, its claim strength, and its judgment.

\textbf{Theorem \hyperref[cl:thm:1]{1} (NoTotalAlgebra)}, theorem-grade:

\[
\Gamma\,;\;\mathcal T\,;\;I_{\mathrm{fin}}\;\vdash\;\textsc{NoTotalAlgebra}\;:\;\texttt{accepted},
\]

where \(\mathcal T\) is the finite admissible package carried by the instance and all cell records are interpreted in \(\mathbf{BirdInt}^{\mathrm{aud}}_{\mathrm{fin}}\).

\textbf{Theorem \hyperref[cl:thm:2]{2} (TypedNonCollapse)}, theorem-grade:

\[
\Gamma\,;\;\mathcal T\,;\;I_{\mathrm{fin}}\;\vdash\;\textsc{TypedNonCollapse}\;:\;\texttt{accepted}.
\]

\textbf{Theorem \hyperref[cl:thm:3]{3} (CellWellFormednessDecidability)}, theorem-grade:

\[
\Gamma\,;\;\mathcal T\,;\;I_{\mathrm{fin}}\;\vdash\;\forall\,C,\;(\mathsf{WFCell}(C)=\texttt{wf}\iff C\text{ well formed})\;:\;\texttt{accepted}.
\]

\textbf{Theorem \hyperref[cl:thm:4]{4} (CoveredCellUniqueStatus)}, theorem-grade:

\[
\Gamma\,;\;\mathcal T\,;\;I_{\mathrm{fin}}\;\vdash\;\textsc{CoveredCellUniqueStatus}\;:\;\texttt{accepted}.
\]

\textbf{Theorem \hyperref[cl:thm:5]{5} (StatusFamilySeparation)}, theorem-grade:

\[
\Gamma\,;\;\mathcal T\,;\;I_{\mathrm{fin}}\;\vdash\;\textsc{StatusFamilySeparation}\;:\;\texttt{accepted}.
\]

\textbf{Theorem \hyperref[cl:thm:6]{6} (DescentThroughQuotient)}, theorem-grade:

\[
\Gamma\,;\;\mathsf H_{\mathrm{fin}}\,;\;I_{\mathrm{fin}}\;\vdash\;\textsc{DescentThroughQuotient}\;:\;\texttt{accepted}.
\]

\textbf{Theorem \hyperref[cl:thm:7]{7} (DescentDefectEquivalence)}, theorem-grade:

\[
\Gamma\,;\;\mathsf H_{\mathrm{fin}}\,;\;I_{\mathrm{fin}}\;\vdash\;\textsc{DescentDefectEquivalence}\;:\;\texttt{accepted}.
\]

\textbf{Theorem \hyperref[cl:thm:8]{8} (FiniteMarkovClosureDeficit)}, theorem-grade with assumptions:

\[
\Gamma\,;\;\mathsf H_{\mathrm{prob}}\,;\;I_{\mathrm{prob}}\;\vdash\;\textsc{FiniteMarkovClosureDeficit}\;:\;\texttt{accepted}.
\]

\textbf{Theorem \hyperref[cl:thm:9]{9} (NoncommutingCompletions)}, theorem-grade:

\[
\Gamma\,;\;\mathsf H_{\mathrm{fin}}\,;\;I_{\mathrm{fin}}\;\vdash\;\textsc{NoncommutingCompletions}\;:\;\texttt{accepted}.
\]

\textbf{Theorem \hyperref[cl:thm:10]{10} (CompletionPastingDefect)}, theorem-grade:

\[
\Gamma\,;\;\mathsf H_{\mathrm{fin}}\,;\;I_{\mathrm{fin}}\;\vdash\;\textsc{CompletionPastingDefect}\;:\;\texttt{accepted}.
\]

\textbf{Theorem \hyperref[cl:thm:11]{11} (IdempotentSaturation)}, theorem-grade:

\[
\Gamma\,;\;\mathsf H_{\mathrm{fin}}\,;\;I_{\mathrm{fin}}\;\vdash\;\textsc{IdempotentSaturation}\;:\;\texttt{accepted}.
\]

\textbf{Theorem \hyperref[cl:thm:12]{12} (StrictExtensionNonfactorization)}, theorem-grade:

\[
\Gamma\,;\;\mathsf H_{\mathrm{fin}}\,;\;I_{\mathrm{fin}}\;\vdash\;\textsc{StrictExtensionNonfactorization}\;:\;\texttt{accepted}.
\]

\textbf{Theorem \hyperref[cl:thm:13]{13} (FactorizationDefectEquivalence)}, theorem-grade:

\[
\Gamma\,;\;\mathsf H_{\mathrm{fin}}\,;\;I_{\mathrm{fin}}\;\vdash\;\textsc{FactorizationDefectEquivalence}\;:\;\texttt{accepted}.
\]

\textbf{Theorem \hyperref[cl:thm:14]{14} (FixedInterfaceDefinabilityBound)}, theorem-grade:

\[
\Gamma\,;\;\mathsf H_{\mathrm{fin}}\,;\;I_{\mathrm{fin}}\;\vdash\;\textsc{FixedInterfaceDefinabilityBound}\;:\;\texttt{accepted}.
\]

\textbf{Theorem \hyperref[cl:thm:15]{15} (ForestNoDrive)}, theorem-grade:

\[
\Gamma\,;\;\mathsf H_{\mathrm{graph}}\,;\;I_{\mathrm{graph}}\;\vdash\;\textsc{ForestNoDrive}\;:\;\texttt{accepted}.
\]

\textbf{Theorem \hyperref[cl:thm:16]{16} (ExactFormNullDrive)}, theorem-grade:

\[
\Gamma\,;\;\mathsf H_{\mathrm{graph}}\,;\;I_{\mathrm{graph}}\;\vdash\;\textsc{ExactFormNullDrive}\;:\;\texttt{accepted}.
\]

\textbf{Theorem \hyperref[cl:thm:17]{17} (NonzeroAffinityEquivalence)}, theorem-grade:

\[
\Gamma\,;\;\mathsf H_{\mathrm{graph}}\,;\;I_{\mathrm{graph}}\;\vdash\;\textsc{NonzeroAffinityEquivalence}\;:\;\texttt{accepted}.
\]

\textbf{Theorem \hyperref[cl:thm:18]{18} (GatingAffinitySuppression)}, theorem-grade:

\[
\Gamma\,;\;\mathsf H_{\mathrm{graph}}\,;\;I_{\mathrm{graph}}\;\vdash\;\textsc{GatingAffinitySuppression}\;:\;\texttt{accepted}.
\]

\textbf{Theorem \hyperref[cl:thm:19]{19} (PromotionGateSoundness)}, theorem-grade:

\[
\Gamma\,;\;\mathsf H_{\mathrm{fin}}\,;\;I_{\mathrm{fin}}\;\vdash\;\textsc{PromotionGateSoundness}\;:\;\texttt{accepted}.
\]

\textbf{Theorem \hyperref[cl:thm:20]{20} (StrictPromotionImpliesNonfactorization)}, theorem-grade:

\[
\Gamma\,;\;\mathsf H_{\mathrm{fin}}\,;\;I_{\mathrm{fin}}\;\vdash\;\textsc{StrictPromotionImpliesNonfactorization}\;:\;\texttt{accepted}.
\]

\textbf{Theorem \hyperref[cl:thm:21]{21} (NonStrictPromotionImpliesFactorization)}, theorem-grade:

\[
\Gamma\,;\;\mathsf H_{\mathrm{fin}}\,;\;I_{\mathrm{fin}}\;\vdash\;\textsc{NonStrictPromotionImpliesFactorization}\;:\;\texttt{accepted}.
\]

\textbf{Theorem \hyperref[cl:thm:22]{22} (Toy22StrictBridge)}, theorem-grade example:

\[
\Gamma\,;\;\mathsf H_{\mathrm{fin}}\,;\;I_{\mathrm{fin}}\;\vdash\;\textsc{Toy22StrictBridge}\;:\;\texttt{accepted},
\]

where \(I_{\mathrm{fin}}\) meets conditions 1--9 of Subsubsection \hyperref[sn:2.4.1]{2.4.1} in the instance of Theorem \hyperref[cl:thm:22]{22} extended by this claim record, its audit and its evidence source. The promotion instrument \(I_{\mathrm{prom}}\) of Subsubsection \hyperref[sn:10.6.1]{10.6.1} classifies the bridge, but would assign this claim \(\texttt{outside\_scope}\): its scope contains only the records of \(\mathcal T_0\), \(\mathcal T_1\) and the bridge.

\textbf{Theorem \hyperref[cl:thm:23]{23} (NoOverreadingSuppression)}, theorem-grade:

\[
\Gamma\,;\;\mathsf H_{\mathrm{instr}}\,;\;I_{\mathrm{instr}}\;\vdash\;\textsc{NoOverreadingSuppression}\;:\;\texttt{accepted}.
\]

\textbf{Theorem \hyperref[cl:thm:24]{24} (SameLevelSelfAuditFailure)}, theorem-grade:

\[
\Gamma\,;\;\mathsf H_{\mathrm{instr}}\,;\;I_{\mathrm{instr}}\;\vdash\;\textsc{SameLevelSelfAuditFailure}\;:\;\texttt{accepted}.
\]

\textbf{Theorem \hyperref[cl:thm:25]{25} (FiniteRotatingAudit)}, theorem-grade:

\[
\Gamma\,;\;\mathsf H_{\mathrm{instr}}\,;\;I_{\mathrm{instr}}\;\vdash\;\textsc{FiniteRotatingAudit}\;:\;\texttt{accepted}.
\]

\textbf{Theorem \hyperref[cl:thm:26]{26} (AIClosureAsPackaging)}, theorem-grade:

\[
\Gamma\,;\;\mathsf H_{\mathrm{AI}}\,;\;I_{\mathrm{AI}}\;\vdash\;\textsc{AIClosureAsPackaging}\;:\;\texttt{accepted}.
\]

\textbf{Theorem \hyperref[cl:thm:27]{27} (AIDescentAsSoundTransformer)}, theorem-grade:

\[
\Gamma\,;\;\mathsf H_{\mathrm{AI}}\,;\;I_{\mathrm{AI}}\;\vdash\;\textsc{AIDescentAsSoundTransformer}\;:\;\texttt{accepted}.
\]

\textbf{Theorem \hyperref[cl:thm:28]{28} (AIRepresentabilityAsFixedness)}, theorem-grade:

\[
\Gamma\,;\;\mathsf H_{\mathrm{AI}}\,;\;I_{\mathrm{AI}}\;\vdash\;\textsc{AIRepresentabilityAsFixedness}\;:\;\texttt{accepted}.
\]

\textbf{Model Theorem \hyperref[cl:thm:29]{29} (AIModelFamily)}, model-grade, for a family whose refinement records contain a pair with \(\rho_k\neq\rho_\ell\):

\[
\Gamma\,;\;\mathsf H_{\mathrm{AI}}\,;\;I_{\mathrm{AI}}\;\vdash\;\mathsf{AIFam}_{\mathrm{fin}}\;\models\;\mathbf{BirdInt}^{\mathrm{aud}}_{\mathrm{fin}}{}^{P_1/P_2/P_4/P_5/P_6}\;:\;\texttt{model\_realization}.
\]

Otherwise the same form holds with the fragment \(P_1/P_2/P_5/P_6\). The family additionally realizes \(P_3\) when two host closures fail to commute.

\textbf{Theorem \hyperref[cl:thm:30]{30} (DescentSquareRecovery)}, theorem-grade:

\[
\Gamma\,;\;\mathsf H_{\mathrm{DC}}\,;\;I_{\mathrm{DC}}\;\vdash\;\textsc{DescentSquareRecovery}\;:\;\texttt{accepted}.
\]

\textbf{Theorem \hyperref[cl:thm:31]{31} (CompletionPastingRecovery)}, theorem-grade:

\[
\Gamma\,;\;\mathsf H_{\mathrm{DC}}\,;\;I_{\mathrm{DC}}\;\vdash\;\textsc{CompletionPastingRecovery}\;:\;\texttt{accepted}.
\]

\textbf{Theorem \hyperref[cl:thm:32]{32} (StrictExtensionAsNoFiller)}, theorem-grade:

\[
\Gamma\,;\;\mathsf H_{\mathrm{DC}}\,;\;I_{\mathrm{DC}}\;\vdash\;\textsc{StrictExtensionAsNoFiller}\;:\;\texttt{accepted}.
\]

\textbf{Model Theorem \hyperref[cl:thm:33]{33} (PICA Model-Realization)}, model-grade:

\[
\Gamma\,;\;\mathsf H_{\mathrm{PICA}}\,;\;I_{\mathrm{PICA}}\;\vdash\;\mathsf{PICAReal}_{\mathrm{fin}}\;\models\;\mathbf{BirdInt}^{\mathrm{aud}}_{\mathrm{fin\text{-}stoch}}{}^{\,\mathrm{cell}/\mathrm{pair}/\mathrm{prov}}\;:\;\texttt{model\_realization}.
\]

\textbf{Model Theorem \hyperref[cl:thm:34]{34} (Cantor Strict-Bridge Realization)}, model-grade:

\[
\Gamma\,;\;\mathsf H_{\mathrm{CantorShell}}\,;\;I_{\mathrm{Cantor}}\;\vdash\;\mathsf{CantorShell}_{\mathrm{aud}}\;\models\;\mathbf{StrictBridge}_{0\to 1}^{\mathrm{audited\ shell}}\;:\;\texttt{model\_realization}.
\]

\textbf{Proposition \hyperref[cl:prop:35]{35} (MechanizationTargetFeasibility)}, audit-grade:

\[
\begin{aligned}
\Gamma\,;\;\mathsf H_{\mathrm{fin}}\,;\;I_{\mathrm{meta}}\;\vdash\;&\text{the manifest-scoped validation artifact checks}\\
&\text{the finite audited paper declarations}\;:\;\texttt{audit\_report}.
\end{aligned}
\]

The verdict \(\texttt{audit\_report}\) is not a judgment verdict: it names the recorded outcome of the validator run of Proposition \hyperref[cl:prop:35]{35}, a repository-level validation audit. It does not change theorem-grade results into model-realization results, and it does not promote model-grade statements to universal theorem-grade claims.
 \section{Version History}\label{sn:J}

Version 1 appeared on 11 May 2026 and Version 2 on 2 October 2026. Version 2 keeps every numbered result of Version 1 with its number and name, in the same order; Proposition \hyperref[cl:prop:35]{35} moved from Section \hyperref[sn:15]{15} of Version 1 to Appendix \hyperref[sn:F.5]{F.5}, and the limitations and conclusion became Sections \hyperref[sn:15]{15} and \hyperref[sn:16]{16}.

\emph{Corrected statements.} Theorem \hyperref[cl:thm:2]{2} asserted a non-implication for every admissible cell; it now exhibits a well-formed \(\texttt{action}\) cell with \(i\neq j\) and every cell status on \((P_6,P_3)\). Theorem \hyperref[cl:thm:18]{18} concluded an \(\texttt{absent\_with\_record}\) certificate that its hypotheses do not supply; it now states that no positive certificate exists. Theorem \hyperref[cl:thm:23]{23} now requires the bridge family to contain the claim's own bridges. Theorem \hyperref[cl:thm:24]{24} stated a dichotomy that ignores higher-priority claim statuses; it now allows them. Theorem \hyperref[cl:thm:25]{25} now assumes a level above the stack, since the level set is finite, and requires every threshold evaluation reachable from its stack-defect computations to be computed by an \(\mathcal R\)-program. Theorem \hyperref[cl:thm:30]{30} separates its fixed-filler and existence clauses and assumes a well-formed square. Theorem \hyperref[cl:thm:32]{32} distinguishes exactness of a given candidate from filler existence and adds square well-formedness. Model Theorem \hyperref[cl:thm:29]{29} realized \(P_4\) unconditionally; it now requires a recorded refinement at the chosen host whose endpoint closures are distinct. In \(\mathsf{CM}_9\) and \(\mathsf{CM}_{10}\) the antecedent is a nonempty factorization defect with a complete strict bridge.

\emph{Added hypotheses and strengthened results.} Theorem \hyperref[cl:thm:3]{3} assumes decidable equality and order on numerical fields and computable threshold data, and Theorem \hyperref[cl:thm:31]{31} a well-formed, in-scope, visible and audited square. Theorems \hyperref[cl:thm:1]{1}, \hyperref[cl:thm:4]{4}, \hyperref[cl:thm:5]{5}, \hyperref[cl:thm:8]{8}, \hyperref[cl:thm:15]{15}, \hyperref[cl:thm:17]{17}, \hyperref[cl:thm:19]{19} and \hyperref[cl:thm:22]{22} have stronger or more explicit conclusions (for example five non-implications in Theorem \hyperref[cl:thm:5]{5}, the lumpability criterion in Theorem \hyperref[cl:thm:8]{8}, and \(Z_1(G)=0\) in Theorem \hyperref[cl:thm:15]{15}); Theorems \hyperref[cl:thm:24]{24}, \hyperref[cl:thm:25]{25} and Model Theorems \hyperref[cl:thm:29]{29}, \hyperref[cl:thm:33]{33} and \hyperref[cl:thm:34]{34} drop or weaken hypotheses or add clauses; Theorems \hyperref[cl:thm:26]{26}--\hyperref[cl:thm:28]{28} hold for pre-hosts and arbitrary posets, and Theorems \hyperref[cl:thm:6]{6}, \hyperref[cl:thm:7]{7} and \hyperref[cl:thm:10]{10}--\hyperref[cl:thm:14]{14} for arbitrary sets and maps. \(\mathsf{CM}_6\) now blocks rather than being insufficient. Proposition \hyperref[cl:prop:35]{35} is a dated validation report with its coverage counts, and the unnumbered statements of Subsection \hyperref[sn:4.7]{4.7} are three propositions with proofs.

\emph{Definitions made precise.} Cell well-formedness (eight conditions), coverage, and the cell, claim, pair and promotion classifiers are stated as finite tables with totality rows; bridges, visibility and source defects, instrument admissibility, the model-realization admissibility clauses of Subsubsection \hyperref[sn:13.1.1]{13.1.1}, the AI host and the level order are given explicit conditions.

\emph{Additions and removals.} Version 2 expands the reader's map (Subsection \hyperref[sn:1.4]{1.4}) and adds the Map of Results (Subsection \hyperref[sn:1.5]{1.5}), the relation to existing work (Subsection \hyperref[sn:1.6]{1.6}), the vocabulary and symbol tables (Subsections \hyperref[sn:2.3.4]{2.3.4}, \hyperref[sn:2.5]{2.5}), a running example (Subsection \hyperref[sn:3.8]{3.8}), the classifier tables of Subsubsection \hyperref[sn:5.6.6]{5.6.6}, Appendix \hyperref[sn:I]{I}, and a rebuilt Appendix \hyperref[sn:F]{F}. The claim-strength tags and instrument-indexed judgments were moved from the statements to Appendix \hyperref[sn:I]{I}, the negative-scope list of Subsection \hyperref[sn:1.3]{1.3} to Subsections \hyperref[sn:2.4.3]{2.4.3} and \hyperref[sn:15.2]{15.2}, and the final nonclaim block into Subsection \hyperref[sn:15.2]{15.2}, keeping NC-1 to NC-36. The strength \(\texttt{dropped}\) is now \(\texttt{excluded}\), and the notations \(\mathrm{Def}(f)\), \(\mathrm{Rep}_\alpha\) and \(\delta_{\mathrm{comp}}\) are now \(\mathrm{Definable}(f)\), \(\mathrm{Rep}_\gamma\) and \(\Delta_3^{\mathrm{comp}}\).
 
\clearpage

\bibliographystyle{plainnat}
\bibliography{references}

\end{document}
